\documentclass[10pt,oneside]{book}
\usepackage[a4paper,margin=24mm,headheight=15pt]{geometry}
\usepackage{fontspec}
\providecommand{\KomorebiFontDirectory}{/usr/share/fonts/opentype/urw-base35/}

\setmainfont{P052-Roman.otf}[Path=\KomorebiFontDirectory,BoldFont=P052-Bold.otf,ItalicFont=P052-Italic.otf,BoldItalicFont=P052-BoldItalic.otf,SmallCapsFont=P052-Roman.otf]
\setsansfont{Nimbus Sans}
\setmonofont{DejaVu Sans Mono}[Scale=0.80]
\usepackage{amsmath,amssymb,mathtools,longtable,booktabs,array,xcolor,fancyhdr,fvextra,xurl,needspace}
\usepackage[unicode,hidelinks]{hyperref}
\hypersetup{pdftitle={Paper XI: Lean Guide},pdfauthor={Komorebi P.},pdfsubject={Packet realization, conditional comparison and source-pinned formal scope}}
\pagestyle{fancy}\fancyhf{}\fancyhead[L]{\small Paper XI\quad Lean Guide}\fancyhead[R]{\small Komorebi P.}\fancyfoot[C]{\thepage}
\setlength{\parindent}{0pt}\setlength{\parskip}{6pt plus 1pt}\emergencystretch=2em
\setcounter{tocdepth}{0}

% Operation Komorebi: shared publication profile.
\usepackage{komorebi-publication}
\renewcommand{\KomorebiTitle}{Tame and Wild Packet Geometry}
\renewcommand{\KomorebiSubtitle}{For local Langlands}
\renewcommand{\KomorebiSeries}{Core series XI}
\renewcommand{\KomorebiKind}{Lean and evidence guide}
\renewcommand{\KomorebiEdition}{Revised edition 1.1.0 / 18 October 2026}
\renewcommand{\KomorebiShortTitle}{Paper XI}
\renewcommand{\KomorebiTitleSize}{38}
\renewcommand{\KomorebiTitleLeading}{43}
\renewcommand{\KomorebiProfile}{classical}
\renewcommand{\KomorebiOrientation}{portrait}
\KomorebiSetup
% Shared publication layout. Mathematical definitions and proofs are unchanged.
\RequirePackage{setspace,caption,needspace,etoolbox}
\definecolor{pubNavy}{HTML}{3C263F}
\definecolor{pubTeal}{HTML}{7B4165}
\definecolor{pubCopper}{HTML}{A68A56}
\definecolor{pubInk}{HTML}{22282C}
\definecolor{pubMuted}{HTML}{55616A}
\definecolor{pubWash}{HTML}{F5F0F4}
\newcommand{\CoreLayout}{%
  \geometry{letterpaper,left=.87in,right=.87in,top=.88in,bottom=.86in,headheight=14pt,headsep=18pt,footskip=26pt,ignorehead,ignorefoot}%
  \color{pubInk}\setstretch{1.06}%
  \setlength{\parindent}{1.05em}\setlength{\parskip}{2.5pt plus .5pt minus .3pt}%
  \setlength{\emergencystretch}{3em}\raggedbottom
  \setlength{\abovedisplayskip}{9pt plus 2pt minus 2pt}%
  \setlength{\belowdisplayskip}{9pt plus 2pt minus 2pt}%
  \setlength{\abovedisplayshortskip}{5pt plus 2pt minus 1pt}%
  \setlength{\belowdisplayshortskip}{6pt plus 2pt minus 1pt}%
  \pagestyle{empty}\fancypagestyle{plain}{\fancyhf{}\renewcommand{\headrulewidth}{0pt}\renewcommand{\footrulewidth}{0pt}}%
  \titleformat{\section}{\normalfont\bfseries\fontsize{14}{18}\selectfont\color{pubNavy}\raggedright}{\textcolor{pubTeal}{\thesection}}{.65em}{}%
  \titleformat{\subsection}{\normalfont\bfseries\fontsize{11.8}{15}\selectfont\color{pubNavy}\raggedright}{\thesubsection}{.65em}{}%
  \titleformat{\subsubsection}{\normalfont\itshape\fontsize{11.3}{14}\selectfont\color{pubNavy}\raggedright}{\thesubsubsection}{.7em}{}%
  \titleformat{\part}{\normalfont\bfseries\fontsize{18}{23}\selectfont\color{pubNavy}\raggedright}{\textcolor{pubTeal}{Part \thepart}}{.7em}{}%
  \titlespacing*{\part}{0pt}{16pt}{18pt}%
  \titlespacing*{\section}{0pt}{18pt plus 3pt minus 2pt}{8pt}%
  \titlespacing*{\subsection}{0pt}{13pt plus 2pt minus 1pt}{6pt}%
  \titlespacing*{\subsubsection}{0pt}{11pt}{4pt}%
  \ifcsname chapter\endcsname
    \titleformat{\chapter}[display]{\normalfont\color{pubNavy}}{\sffamily\fontsize{8.5}{11}\selectfont\color{pubTeal}\MakeUppercase{\chaptertitlename}\enspace\thechapter}{9pt}{\fontsize{24}{29}\selectfont\bfseries\raggedright}%
    \titlespacing*{\chapter}{0pt}{6pt}{22pt}%
  \fi
  \captionsetup{format=hang,justification=raggedright,singlelinecheck=false,font={small,sf,stretch=1.06},labelfont={bf,sf,color=pubInk},textfont={color=pubMuted},labelsep=period,skip=8pt}%
  \hypersetup{colorlinks=true,linkcolor=pubTeal,citecolor=pubTeal,urlcolor=pubTeal,breaklinks=true}%
  \clubpenalty=10000\widowpenalty=10000\displaywidowpenalty=10000
  \AtBeginEnvironment{longtable}{\ifdim\csname f@size\endcsname pt>10pt\relax\small\fi}
  \frenchspacing
}
\newcommand{\CoreOpening}[3]{%
  \thispagestyle{empty}\begingroup\setlength{\parindent}{0pt}%
  {\sffamily\fontsize{8.5}{11}\selectfont\color{pubTeal}\MakeUppercase{#1}\par}%
  \vspace{18pt}%
  {\fontsize{28}{34}\selectfont\color{pubNavy}\raggedright#2\par}%
  \vspace{12pt}%
  {\sffamily\fontsize{11}{15.5}\selectfont\color{pubMuted}\raggedright#3\par}%
  \vspace{19pt}\noindent\textcolor{pubCopper}{\rule{42pt}{.8pt}}\par
  \vspace{16pt}\endgroup\normalfont\color{pubInk}%
}

\newcommand{\CoreContext}[2]{%
 \par\begingroup\small\sffamily\color{pubMuted}\setlength{\parindent}{0pt}%
 \textbf{#1.}\enspace #2\par\endgroup\vspace{11pt}%
}

% Core-series presentation additions, applied after the inherited setup.
\CoreLayout
\renewcommand{\familydefault}{\rmdefault}
\definecolor{KomorebiInk}{HTML}{3C263F}
\definecolor{KomorebiPine}{HTML}{7B4165}
\definecolor{KomorebiOchre}{HTML}{A68A56}
\definecolor{KomorebiWarm}{HTML}{FFFFFF}
\definecolor{KomorebiMuted}{HTML}{655D68}
\colorlet{ink}{KomorebiInk}\colorlet{accent}{KomorebiPine}
\colorlet{teal}{KomorebiPine}\colorlet{Teal}{KomorebiPine}
\colorlet{Navy}{KomorebiInk}\colorlet{navy}{KomorebiInk}
\fancypagestyle{fancy}{\fancyhf{}\renewcommand{\headrulewidth}{0pt}\renewcommand{\footrulewidth}{0pt}}
\renewcommand{\KomorebiCover}{%
 \clearpage\hypersetup{pageanchor=false}\thispagestyle{empty}\begingroup
 \setlength{\parindent}{0pt}\setlength{\parskip}{0pt}
 \vspace*{22pt}
 {\sffamily\fontsize{8.5}{11}\selectfont\color{pubTeal}\MakeUppercase{\KomorebiSeries}\par}
 \vspace{30pt}
 {\rmfamily\fontsize{30}{36}\selectfont\color{pubNavy}\raggedright\KomorebiTitle\par}
 \vspace{18pt}
 {\rmfamily\fontsize{19}{25}\selectfont\color{pubNavy}\raggedright\KomorebiSubtitle\par}
 \vspace{24pt}\textcolor{pubCopper}{\rule{42pt}{.8pt}}\par\vspace{18pt}
 {\sffamily\fontsize{10}{14}\selectfont\color{pubMuted}\KomorebiKind\par}
 \vfill
 {\rmfamily\fontsize{13}{17}\selectfont Komorebi P.\par}\vspace{9pt}
 {\sffamily\fontsize{8.5}{12}\selectfont\color{pubMuted}\KomorebiEdition\par}
 \vspace{20pt}\endgroup\clearpage\setcounter{page}{1}\hypersetup{pageanchor=true}\color{pubInk}\normalfont
}

\begin{document}
\KomorebiCover
\KomorebiOpeningNote{A source-pinned guide to the proved formal mechanisms and their arithmetic
interfaces. The general realization endpoint is Conjecture 66. The release
includes Markdown proofs, construction notes, audits and research records.}
\frontmatter

\tableofcontents\mainmatter
\chapter{A reader contract for this edition}
Paper XI develops packet decomposition, finite representation comparisons, actual arithmetic/type adapters, and a conditional comparison interface. This guide makes the logical boundaries visible in the source language. The traditional paper gives the mathematical exposition. The visual guide illustrates the constructions. The retained Lean sources provide the formal statements and proofs at their printed hypotheses.

The general arithmetic realization endpoint is presented as \textbf{Conjecture 66}. Conjecture 66 isolates the missing realization claim from historical Claim 16.2, indexed as original obligation 66. This separation preserves the proved finite comparison and the proved conditional implication. It records that finite assembly and coherent transport have not yet supplied an actual supercuspidal realization in the original generality.

The accompanying release includes Markdown (.md) files containing detailed proofs, construction notes, proof audits, and research records. These documents distinguish established results, conditional results, and the remaining obligations associated with Conjecture 66.

\section{What Lean certifies}
Acceptance by Lean checks that the declared conclusion follows from its binders, imported declarations, and foundational axioms. It does not silently instantiate a local field, an arithmetic type, a reference correspondence, or a normalized local-factor convention. Those objects must appear in the source or in a separate mathematical attachment.

A structure may contain an existence field. A theorem that consumes that structure proves an implication using the field. It is still a substantive theorem when the comparison map, descent or uniqueness is constructed. It is not a proof of the structure's arithmetic existence. In particular, \texttt{Realization.exists\_atom} is an existence input in the formal comparison interface.

\section{Evidence available for this guide}
The source catalog was generated for the inherited edition from retained production source bytes. This edition preserves its selected modules, source hashes, import edges and named lexical records. A separately dated fresh execution replays the selected source in the pinned environment and audits its actual declaration dependencies. The current result is summarized below. Historical compiler and same-kernel records later in this guide retain their original dates and scopes. The preserved catalog's historical evidence flags describe the edition in which that catalog was generated; the dated current receipts determine this release's execution status. The PDF build and the formal execution are separate operations.

\begin{quote}
\textbf{Fresh verification, 18 October 2026: all required gates passed.}
The 423 selected publication modules were compiled with their retained source
hashes unchanged. All 4,342 catalogued names resolved to their recorded owners.
Four source names---two private theorems and two named local instances---are
matched uniquely to their kernel names using Lean's own private-name mapping;
the source and kernel names are both recorded.
The recovered terminal and the new comparison supplement were checked as
separately counted targets.
\end{quote}

\begin{center}\small
\begin{tabular}{@{}lrrr@{}}
\toprule
Population & Modules & Catalogued names & Logical declarations\\
\midrule
Selected sources & 423 & 4,342 & 7,672\\
Recovered terminal & 1 & 0 & 51\\
One-candidate supplement & 1 & 1 & 1\\
\bottomrule
\end{tabular}
\end{center}

The last column is the complete population of logical declarations owned by
the named modules in the imported private kernel environment. It includes
generated and private logical declarations. It is distinct from the retained
lexical catalog, and is not a count of independent mathematical discoveries.
The audit separately lists 1,636 compiler-only auxiliary names for
the selected sources; these are absent from the kernel constant map.

Every audited declaration's transitive axiom dependencies belong to
\texttt{Classical.choice}, \texttt{Quot.sound}, and \texttt{propext}.
No forbidden dependency was found. All three aggregate imports passed.
Two isolated negative controls were rejected as intended: uncatalogued public
and private custom axioms in a module with a safely catalogued theorem, and a
deliberately unresolved catalog name. These fixtures are not Paper XI modules.
The supplement's observed axioms were \texttt{Quot.sound} and \texttt{propext}.

\textbf{Fresh-source lineage and unsuccessful attempts.}
The first replay started from an empty Paper XI object directory and passed
421 selected modules plus the supplement. Its unfinished tower compilation
was manually stopped early after a longer historical duration was found.
The first run remains recorded as \texttt{PARTIAL\_OR\_FAIL}; it was not a
source rejection, timeout, or out-of-memory event.
The first continuation authenticated every completed source/object hash and
dependency closure, reused only those 422 successful objects from this
session, and compiled the remaining modules with a 3,600-second per-module
limit. The tower and its base-change dependent completed in
1,066.8 and 827.6 seconds, respectively, completing
all 423 selected modules. The terminal then emitted substantive elaboration
errors, including a heartbeat-limit diagnostic. After inspection of its log,
that known-failed process was deliberately terminated. This second run also
remains \texttt{PARTIAL\_OR\_FAIL}; its stop was diagnostic-driven.

The source-revision continuation authenticated the successful source/object
closure and invalidated only the revised failed terminal. Its terminal process
was then killed at the environment's 8\,GiB memory limit; the recorded
out-of-memory counters increased. This third replay also remains failed.
An unchanged-source continuation with synchronous elaboration,
\texttt{Elab.async=false}, also reached the memory limit. That fourth full
replay remains failed. Bounded diagnostics isolated an expensive conversion
in the lower-inclusion witness. The final continuation authenticated the
successful closure and compiled separately reviewed local proof repairs.
No historical Paper XI object was used. Pinned official external Mathlib cache
objects were imported; all Mathlib sources were not rebuilt.

\textbf{Separate additions and audit history.}
The terminal required more than the missing explicit
\nolinkurl{PaperXI.RamificationPolyphase} import. Its repair also corrects
implicit arguments, required standing-hypothesis bindings, declaration syntax
and local proofs. Five private component lemmas factor the comparison proof;
the lower-inclusion witness uses the retained equality in its forward direction.
A separate review checks the intended conclusions and
already-declared hypotheses. The original source, failed intermediate and
exact diffs are retained; the absent historical receipt is not recreated.
The supplement proves the supplied-candidate implication printed below.
The exploratory axiom harness needed explicit counter types, a distinction
between compiler auxiliaries and kernel declarations, a larger recursion
bound for the catalogue literal, and private-name resolution. The failed
pilot records remain alongside the successful final audits. No
selected Lean theorem statement or proof body was changed.

The complete results are in
\nolinkurl{verification/lean/RELEASE_VERIFICATION.json} and its prose companion.
The portable project supplies source, exact pins and replay scripts, with
downloaded toolchains and generated objects omitted. This release ran fresh
source compilation and complete ownership/axiom audits. It did not run a new
additional full \texttt{leanchecker} pass or use an independently implemented
kernel. These results establish the formal declarations at their stated
hypotheses; the four open arithmetic tasks retain their separate statuses.


The recorded environment is Lean 4.34.0-rc2, commit \nolinkurl{6a10ac8c22beadecabdbb0919c2b50214762f91d}, and Mathlib revision \nolinkurl{aeb72f1e9981b6d35b426b16dca7759be3e7ba06}. The reported allowed transitive axiom set is \texttt{propext}, \texttt{Classical.choice}, and \texttt{Quot.sound}. Same-kernel replay is an additional invocation of Lean's kernel, not an independently implemented checker.

\section{How the catalog is selected}
The selected union contains 423 distinct module paths and 4,342 explicitly named lexical declaration commands. It starts with the 302-module v0.44 source tree, adds three checked normalized-Frobenius modules, then retains the arithmetic/type branch through authenticated base108 and checked continuations 110--113. Identical repeated module bytes are collapsed. Layer-specific provenance remains in the JSON catalog.

This is a source census, not a count of independent mathematical discoveries or kernel-generated declarations. Constructors, recursors and unnamed instances are outside this lexical count. The historical ``65/66'' number concerns original paper obligations, whereas the appendix includes many supporting definitions and lemmas. The conjecture receives no Lean theorem-count promotion.

\clearpage
\section{Supplement: the comparison of one supplied candidate}
The publication comparison fixes one admissible candidate and one admissible reference. A separating factor system identifies their isomorphism classes if their declared factor profiles agree. The additional module \nolinkurl{PaperXI.SingleCandidateComparison} applies \nolinkurl{LocalFactorIdentification.class_eq_iff} to exactly these supplied objects.

The expected factor profile is a parameter. H5 is assumed only for this candidate. The proof does not consume \texttt{Realization.exists\_atom}, does not quantify H5 over every realizer, and does not instantiate an arithmetic correspondence. The separating factor system and admissibility conditions remain explicit inputs. This is a formal presentation of the paper's conditional implication, not an additional arithmetic existence result.

The complete supplementary module is printed below. It is counted separately from the unchanged 423-module catalog and the forty inherited teaching excerpts. Its own compiler and axiom receipt is included with the release.
\VerbatimInput[fontsize=\small,breaklines=true,breakanywhere=true]{supplements/SingleCandidateComparison.lean}

\section{The corrected moving-space phase domain}
In odd residue characteristic, the monograph's breakwise formula now applies to \(M=\operatorname{im}(g-I)\) and its break summands, with \(g\) explicitly symplectic and fixed point free on \(M\). The selected theorem \nolinkurl{PaperXI.BreakOddPhase.source} already proves this product of moving phases. \nolinkurl{BreakMoving} supplies the complement and nondegeneracy facts. A zero moving space contributes one. This source alignment does not add an order-prime-to-\(p\) hypothesis and does not modify the retained Lean module.

\section{Reading literal excerpts}
A displayed declaration excerpt retains the exact source text through a top-level \texttt{:=}, omitting its subsequent proof or value. Named-argument assignments inside binders do not terminate the excerpt. Structures retain all their fields. The complete declaration is available in the source file. A proof excerpt is a contiguous source slice. Ambient assumptions are explained beside it and retained as raw setup lines in \texttt{EXCERPTS.json}; imported definitions remain necessary.

The display is not a separately compiled replacement module. Unicode supported by the monospace font appears directly. Unsupported glyphs use reversible \texttt{[U+XXXX]} tokens. The raw excerpt, display file, conversion map and SHA-256 hashes are included.

\clearpage
\section*{Navigation through the five readers}
\phantomsection
\addcontentsline{toc}{chapter}{Navigation through the five readers}
The formal core, the written arithmetic constructions and the research-progress record have different evidence types. This edition preserves the selected formal source and provides direct reading routes into the other four readers. The selected sources remain unchanged. The terminal replay and a one-candidate supplementary lemma are counted separately; neither supplies an arithmetic-existence axiom or a new discovery count.

\begin{tabular}{@{}>{\raggedright\arraybackslash}p{.29\linewidth}>{\raggedright\arraybackslash}p{.66\linewidth}@{}}\toprule
Question & Reading route and formal boundary\\\midrule
The finite comparison & \href{https://komorebi.invalid/b07/lean_finite}{Monograph finite tame package}. Its finite statement and its arithmetic interpretation retain distinct hypotheses.\\[9pt]
The existence field & \href{https://komorebi.invalid/b07/lean_contract}{Monograph Conjecture 66}, \href{https://komorebi.invalid/b07/lean_research_contract}{progress Section 1} and \href{https://komorebi.invalid/b07/lean_v01}{atlas V01}. The source field \texttt{Realization.exists\_atom} remains an input.\\[9pt]
The pointwise comparison & \href{https://komorebi.invalid/b07/lean_comparison}{Comparison of one realized candidate}. H5 is about that candidate. The more general formal interface's all-realizers binder is not a stronger publication conjecture.\\[9pt]
Ordinary arithmetic progress & \href{https://komorebi.invalid/b07/lean_arithmetic}{Monograph Section 19} and \href{https://komorebi.invalid/b07/lean_progress}{the separate progress paper}. The explicit parameter, character, conductor and marking arguments are not newly Lean-certified here.\\[9pt]
Worked interpretation & \href{https://komorebi.invalid/b07/lean_core_visual}{Core visual guide} and \href{https://komorebi.invalid/b07/lean_atlas}{research atlas}. Abstract finite examples are not certificates for an unspecified arithmetic type.\\\bottomrule
\end{tabular}

\subsection*{The original-obligation crosswalk}
The existing 66-row historical table now links its printed-number column to the actual current monograph statements. The final row points to Conjecture 66, while the separate comparison proposition has its own link. These links are navigational. They do not assert that every printed theorem is definitionally identical to its selected entry module.

All 423 selected module paths, the 4,342 named lexical declaration records, and the 40 teaching excerpts retain their inherited source identities. The current source audit checks their correspondence and preservation. Fresh compiler and axiom receipts identify the executed targets separately from that inherited census.
\clearpage
\section*{Evidence boundaries for this revision}
\phantomsection
\addcontentsline{toc}{chapter}{Evidence boundaries for this revision}
\textbf{Compilation and publication are separate.} This revision builds and reviews the PDF readers and records a separate fresh Lean execution. The preceding execution summary gives the current result. Historical compiler and same-kernel receipts remain attributable to their exact source revisions. No independently implemented proof checker is claimed.

\textbf{The terminal has a separately dated repair and replay.} The original successful v0.45 receipt remains unavailable. Adding the missing explicit \nolinkurl{PaperXI.RamificationPolyphase} import exposed further elaboration errors. The repaired source also corrects implicit-argument calls, required standing-hypothesis bindings, declaration syntax and local proofs. The original, failed intermediate, exact diffs and hypothesis review are retained. Its compiler status is reported separately from the unchanged 423-module selection. This target is a restricted tame Weil-side attachment; it does not construct a general wild automorphic realizer.

\textbf{Block 2 is a preserved-reader baseline.} The 33-page first part of the research companion survives in its PDF and records. Its complete editable working archive has not been recovered. Later continuations are independently editable. This edition does not claim a unified rebuild of the missing source.

\textbf{The selected classical-source gap has been closed.} The Henniart--Lemaire induction interface has now been checked against the full official text, including direct visual inspection of the relevant normalization formulas. The bounded review preserves its characteristic, rank, genericity and orbit hypotheses. Other Bushnell--Kutzko, conductor, tower and descent inputs retain their own evidence. This source check is not H5 or corrected-character occurrence.

\subsection*{Four open obligations, not four missing formal lemmas}
General arithmetic realization requires an actual common host with the declared actions and transports. H5 requires the characterizing factors for that same realizer. The corrected-character branch still needs occurrence in its fixed candidate. Where the stronger native contract is prescribed, the actual square, section and honest-phase certificate must be computed on the chosen type. The \href{https://komorebi.invalid/b07/lean_native}{research certificate plate} deliberately leaves those arithmetic fields uncomputed.

The positive transported-marking contract remains separate. Its hypotheses are not silently strengthened by the native-marking task. Likewise, a failed original carrier does not exclude every permitted enlarged inducing carrier.

\subsection*{Reproducible reading routes}
The delivered builder resolves the local companion destinations after compiling all three original readers. Keep the two preserved research companions and the revised originals together. The source catalogs remain unchanged. The new navigation JSON records the current targets without overwriting historical page numbers or evidence dates.

\chapter{Fourier projectors and Frobenius transport}
The cyclic packet is constructed by averaging against an additive character. Idempotence, orthogonality and completeness are consequences of the representation law and character orthogonality. The packet space is the actual range of the computed projector, and Frobenius transports that range by its conjugation law.

\subsection{The averaging operator}
\label{guide:E01}
The ambient cyclic representation and additive-character hypotheses are declared in the source module. The index is a residue label in \(\mathbb Z/M\mathbb Z\). The scalar field is \(\mathbb C\), so the finite order is invertible.
The definition sums actual operators. Later results compute their products rather than postulating a family of orthogonal idempotents.

\begin{quote}\footnotesize\textbf{E01}: \nolinkurl{PaperXI.FourierProjectors}, source lines 13--15.\par Contiguous source proof excerpt.\par Source SHA-256: {\scriptsize\texttt{02af62a8c8e5753ef7a695629d44c71db407211d18a6bd6b6a48865ba731de52}}.\end{quote}

\VerbatimInput[fontsize=\scriptsize,numbers=left,firstnumber=13,breaklines=true,breakanywhere=true,numbersep=5pt]{excerpts/E01.render.txt}

\subsection{Completeness of the packet resolution}
\label{guide:E02}
The finite cyclic group and primitive character remain fixed. The sum ranges over all residue labels.
Completeness says that every vector is the sum of its packet components. A vanishing packet is permitted, and the theorem does not assert that every residue label occurs.

\begin{quote}\footnotesize\textbf{E02}: \nolinkurl{PaperXI.FourierProjectors}, source lines 68--77.\par Contiguous source proof excerpt.\par Source SHA-256: {\scriptsize\texttt{02af62a8c8e5753ef7a695629d44c71db407211d18a6bd6b6a48865ba731de52}}.\end{quote}

\VerbatimInput[fontsize=\scriptsize,numbers=left,firstnumber=68,breaklines=true,breakanywhere=true,numbersep=5pt]{excerpts/E02.render.txt}

\subsection{An actual packet equivalence}
\label{guide:E03}
The supplied invertible Frobenius operator satisfies the declared conjugation relation with the cyclic inertia action. The multiplication direction in the residue label is fixed by that relation.
The operator and its inverse restrict to the two computed packet ranges. Their inverse identities give a linear equivalence, so equal packet dimensions follow from a constructed transport.

\begin{quote}\footnotesize\textbf{E03}: \nolinkurl{PaperXI.PacketTransport}, source lines 62--66.\par Literal declaration excerpt. A proof/value after a top-level assignment is omitted; structure fields are retained.\par Source SHA-256: {\scriptsize\texttt{dc7282e77c935ca61b8807fe87f4fe2fa372d720a28d0a1c8982c6c749e1a910}}.\end{quote}

\VerbatimInput[fontsize=\scriptsize,numbers=left,firstnumber=62,breaklines=true,breakanywhere=true,numbersep=5pt]{excerpts/E03.render.txt}

\section{Geometric Frobenius convention and the publication repair}
The geometric convention in this edition is
\[
 \Phi^{-1}t\Phi=t^q,
 \qquad \Phi P_j\Phi^{-1}=P_{qj}.
\]
Under the other relation \(\Phi t\Phi^{-1}=t^q\), the correct label is \(q^{-1}j\). The original Proposition 3.1 proof paired that other relation with the geometric label \(qj\). The final traditional paper repairs the displayed relation to the geometric convention, preserving its intended packet transport. The checked Lean module already states both conventions separately. No new Lean proof is claimed for this prose convention repair.

\subsection{The existing geometric covariance theorem}
\label{guide:E04}
The cyclic order is nonzero. The element \(q\) is a unit modulo that order, and the source hypothesis uses the conjugation by \(F^{-1}\).
This exact signature and proof support the publication correction. They distinguish the Frobenius direction from its action on projector labels, rather than changing both directions simultaneously.

\begin{quote}\footnotesize\textbf{E04}: \nolinkurl{PaperXI.TameProjectors}, source lines 45--52.\par Contiguous source proof excerpt.\par Source SHA-256: {\scriptsize\texttt{d5881299706cbe60a5cbefc9dc60bb9bdffbc1365228529c181b9a5120b71d8e}}.\end{quote}

\VerbatimInput[fontsize=\scriptsize,numbers=left,firstnumber=45,breaklines=true,breakanywhere=true,numbersep=5pt]{excerpts/E04.render.txt}

\chapter{Induction, return scalars and tensor interactions}
The induced representation branch uses a degree homomorphism and its degree subgroups. A normalized degree-one element gives coset representatives. Its cyclic action has a closing scalar determined by the inducing character. Tensor interactions are constructed as actual representation equivalences indexed by greatest-common-divisor classes.

\subsection{The closing scalar is derived}
\label{guide:E05}
An actual character of the degree subgroup and a degree-one Frobenius element are supplied. The degree \(n\) is positive.
The scalar evaluates the character at the full return element. It is not an independent certificate whose agreement with the action is assumed.

\begin{quote}\footnotesize\textbf{E05}: \nolinkurl{PaperXI.InducedCycleInvariants}, source lines 33--33.\par Contiguous source proof excerpt.\par Source SHA-256: {\scriptsize\texttt{2faaeb89211dca9d88df443a27ad69ab7e310783d4e07a3ad54d85c3c5bcae1a}}.\end{quote}

\VerbatimInput[fontsize=\scriptsize,numbers=left,firstnumber=33,breaklines=true,breakanywhere=true,numbersep=5pt]{excerpts/E05.render.txt}

\subsection{Dimension and characteristic polynomial}
\label{guide:E06}
The basis of the genuine degree-induced space has already been constructed. Finite dimensionality and the nonzero degree hypotheses occur in the module's setup.
The resulting dimension and polynomial record the entire cyclic block, including its closing scalar. Their interpretation as arithmetic local data needs the corresponding group and normalization attachments.

\begin{quote}\footnotesize\textbf{E06}: \nolinkurl{PaperXI.InducedCycleInvariants}, source lines 75--77.\par Literal declaration excerpt. A proof/value after a top-level assignment is omitted; structure fields are retained.\par Source SHA-256: {\scriptsize\texttt{2faaeb89211dca9d88df443a27ad69ab7e310783d4e07a3ad54d85c3c5bcae1a}}.\end{quote}

\VerbatimInput[fontsize=\scriptsize,numbers=left,firstnumber=75,breaklines=true,breakanywhere=true,numbersep=5pt]{excerpts/E06.render.txt}

\subsection{The tensor decomposition is an equivalence}
\label{guide:E07}
Two positive degrees and two characters of their degree subgroups are supplied. Representatives use \(\varphi^{-i}\), fixing the direction of the shifted interaction character.
There are \(\gcd(n,m)\) interaction blocks of degree \(\operatorname{lcm}(n,m)\). The source builds an intertwining equivalence of representations, not merely the dimension identity \(\gcd(n,m)\operatorname{lcm}(n,m)=nm\).

\begin{quote}\footnotesize\textbf{E07}: \nolinkurl{PaperXI.TameTensorSource}, source lines 66--69.\par Literal declaration excerpt. A proof/value after a top-level assignment is omitted; structure fields are retained.\par Source SHA-256: {\scriptsize\texttt{b52fef0aac236edae7a22103dcf267e147fa13606783529fa64e3c4c4e954d3e}}.\end{quote}

\VerbatimInput[fontsize=\scriptsize,numbers=left,firstnumber=66,breaklines=true,breakanywhere=true,numbersep=5pt]{excerpts/E07.render.txt}

\chapter{Finite conductors and the tame comparison package}
The finite conductors are sums of codimensions of fixed spaces for a declared subgroup family. Arithmetic Swan and Artin conductors require that this family and its weights match the arithmetic ramification conventions. The finite residue Gauss phase is also a declared character sum before its arithmetic interpretation.

\subsection{A declared finite ramification sum}
\label{guide:E08}
The representation, a finite set of depths, a subgroup at each depth and rational weights are supplied. The source does not construct an absolute Weil ramification filtration here.
These inputs determine the exact finite sum. The tame vanishing theorem applies when the selected positive-depth groups act trivially.

\begin{quote}\footnotesize\textbf{E08}: \nolinkurl{PaperXI.TameConductor}, source lines 56--58.\par Contiguous source proof excerpt.\par Source SHA-256: {\scriptsize\texttt{808eb443d2a292fc79cd70ce323c369acd62c38e9393ec4c19269dc06d575dcf}}.\end{quote}

\VerbatimInput[fontsize=\scriptsize,numbers=left,firstnumber=56,breaklines=true,breakanywhere=true,numbersep=5pt]{excerpts/E08.render.txt}

\subsection{Proposition 16.1 at finite scope}
\label{guide:E09}
The exact induced spaces, degree character, cyclic relation, inertia label, finite-depth subgroup data and residue characters are supplied. The additive character is nontrivial.
The package derives observations, a tensor equivalence, the total tensor dimension and the Gauss-phase magnitude. Its magnitude formula retains the trivial multiplicative-character exception. This finite statement does not certify standard local-Weil factors or an arithmetic epsilon-factor normalization.

\begin{quote}\footnotesize\textbf{E09}: \nolinkurl{PaperXI.FiniteTameComparison}, source lines 93--107.\par Literal declaration excerpt. A proof/value after a top-level assignment is omitted; structure fields are retained.\par Source SHA-256: {\scriptsize\texttt{1c4357da32eaade5be0c88670eaa605f47fbc4269253888bbf9f72509d0e428d}}.\end{quote}

\VerbatimInput[fontsize=\scriptsize,numbers=left,firstnumber=93,breaklines=true,breakanywhere=true,numbersep=5pt]{excerpts/E09.render.txt}

\chapter{Heisenberg actions, central products and phase comparison}
A finite Heisenberg action combines translations and modulation operators. Irreducibility is proved from those actual operators. Break assembly uses a genuine quotient central product and a tensor carrier. To extend the bridge through a stabilizer, the retained action on the whole central extension and the scalar discrepancy must both be controlled.

\subsection{The finite Schrodinger action}
\label{guide:E10}
The finite Heisenberg group and its central character are fixed by the source's preceding constructions. The function-space carrier is literal.
The multiplication law is checked for the displayed operators. Invertibility, central action and irreducibility are subsequent consequences.

\begin{quote}\footnotesize\textbf{E10}: \nolinkurl{PaperXI.HeisenbergRepresentation}, source lines 14--16.\par Literal declaration excerpt. A proof/value after a top-level assignment is omitted; structure fields are retained.\par Source SHA-256: {\scriptsize\texttt{386317cf254a374627df755dc02465986fc602151303b4e2b3413f2b82b46492}}.\end{quote}

\VerbatimInput[fontsize=\scriptsize,numbers=left,firstnumber=14,breaklines=true,breakanywhere=true,numbersep=5pt]{excerpts/E10.render.txt}

\subsection{Irreducibility from actual translations and modulations}
\label{guide:E11}
The source supplies the required nondegenerate character pairing. Its invariant-subspace proof uses coordinate projections and spans the full function space.
This finite representation theorem supplies a building block. It does not by itself embed the block in a canonical arithmetic simple type.

\begin{quote}\footnotesize\textbf{E11}: \nolinkurl{PaperXI.HeisenbergRepresentation}, source lines 58--59.\par Literal declaration excerpt. A proof/value after a top-level assignment is omitted; structure fields are retained.\par Source SHA-256: {\scriptsize\texttt{386317cf254a374627df755dc02465986fc602151303b4e2b3413f2b82b46492}}.\end{quote}

\VerbatimInput[fontsize=\scriptsize,numbers=left,firstnumber=58,breaklines=true,breakanywhere=true,numbersep=5pt]{excerpts/E11.render.txt}

\subsection{The assembled group is a quotient}
\label{guide:E12}
A finite indexed family of central-extension data shares a declared center. The balanced central subgroup is defined earlier in the module.
The quotient construction retains an explicit center even for an empty family. The quotient's representation must descend through the balancing relation, and that descent is proved in the supporting modules.

\begin{quote}\footnotesize\textbf{E12}: \nolinkurl{PaperXI.CentralProduct}, source lines 39--41.\par Contiguous source proof excerpt.\par Source SHA-256: {\scriptsize\texttt{9d1e6a4ece84a862ef75383b66b97af694c9ac2298fc2047a6c4844bc6aa34af}}.\end{quote}

\VerbatimInput[fontsize=\scriptsize,numbers=left,firstnumber=39,breaklines=true,breakanywhere=true,numbersep=5pt]{excerpts/E12.render.txt}

\subsection{Precisely what break assembly consumes}
\label{guide:E13}
The finite family includes the actual groups, quotient groups and complex carriers. Its record keeps central extensions, retained lifts, central characters and relevant pairings.
An initial tensor intertwiner and a final representation equivalence are not input fields. They are constructed. Dyadic square and retained-group data remain part of the input and cannot be recovered from an alternating commutator form alone.

\begin{quote}\footnotesize\textbf{E13}: \nolinkurl{PaperXI.FiniteBreakAssembly}, source lines 22--37.\par Literal declaration excerpt. A proof/value after a top-level assignment is omitted; structure fields are retained.\par Source SHA-256: {\scriptsize\texttt{ecca3693f7733c7ceef597b6e821b111610c82d612fd66abfe754af0e5be434c}}.\end{quote}

\VerbatimInput[fontsize=\scriptsize,numbers=left,firstnumber=22,breaklines=true,breakanywhere=true,numbersep=5pt]{excerpts/E13.render.txt}

\subsection{Closing the honest stabilizer action}
\label{guide:E14}
Honest source and target stabilizer actions are supplied together with covariance. The discrepancy is tested on a generating set.
The construction kills the scalar discrepancy through the generator phase hypothesis. Matching quotient actions alone would leave this scalar issue unresolved.

\begin{quote}\footnotesize\textbf{E14}: \nolinkurl{PaperXI.FiniteBreakAssembly}, source lines 111--115.\par Literal declaration excerpt. A proof/value after a top-level assignment is omitted; structure fields are retained.\par Source SHA-256: {\scriptsize\texttt{ecca3693f7733c7ceef597b6e821b111610c82d612fd66abfe754af0e5be434c}}.\end{quote}

\VerbatimInput[fontsize=\scriptsize,numbers=left,firstnumber=111,breaklines=true,breakanywhere=true,numbersep=5pt]{excerpts/E14.render.txt}

\chapter{Odd phases and dyadic quadratic refinements}
The paper separates odd phases defined by a determinant from dyadic phases defined by a quadratic refinement. The formal certificate computes the declared determinant phase through its characteristic polynomial. Identifying that value with an arithmetic Weil trace or a canonical beta extension requires an additional argument. In residue characteristic two, an alternating form alone does not retain the quadratic data needed for the phase.

\subsection{The actual polynomial phase certificate}
\label{guide:E15}
The residue characteristic is an odd prime. A finite-dimensional vector space carries an actual nondegenerate alternating form. The supplied invertible linear operator need not be fixed-point-free for this algebraic identity.
The form derives even dimension. Evaluation of the actual characteristic polynomial at one gives the determinant of the difference operator with the declared sign convention. This certifies the defined phase, without assuming that it equals a separately normalized arithmetic phase.

\begin{quote}\footnotesize\textbf{E15}: \nolinkurl{PaperXI.OddPhaseCertificate}, source lines 46--50.\par Literal declaration excerpt. A proof/value after a top-level assignment is omitted; structure fields are retained.\par Source SHA-256: {\scriptsize\texttt{783d84147c4f09cbca432ad799b1e57d9bd47da13f3841e92d4f071d49709f2e}}.\end{quote}

\VerbatimInput[fontsize=\scriptsize,numbers=left,firstnumber=46,breaklines=true,breakanywhere=true,numbersep=5pt]{excerpts/E15.render.txt}

\subsection{Refinements form an affine space}
\label{guide:E16}
Two quadratic refinements have the same bilinear form over the field with two elements. Nondegeneracy is not required for the affine-difference statement.
Their difference is uniquely linear. Hence the bilinear form does not choose one refinement. When the arithmetic contract prescribes square data or a whole marked finite extension, that datum must be retained.

\begin{quote}\footnotesize\textbf{E16}: \nolinkurl{PaperXI.DyadicAffine}, source lines 74--75.\par Literal declaration excerpt. A proof/value after a top-level assignment is omitted; structure fields are retained.\par Source SHA-256: {\scriptsize\texttt{06fd5b5a57b17060abd22a11fb29592fcd53b06aa09388b92314fe2fd23692e3}}.\end{quote}

\VerbatimInput[fontsize=\scriptsize,numbers=left,firstnumber=74,breaklines=true,breakanywhere=true,numbersep=5pt]{excerpts/E16.render.txt}

\subsection{The exact Arf Gauss sign}
\label{guide:E17}
The finite index set parametrizes symplectic planes. Each plane has its explicit quadratic-refinement coefficients in the field with two elements.
The integer character sum is \( (-1)^{\operatorname{Arf}(q)}2^d\). The proof factors the complete finite sum into plane sums. It is an exact theorem, not an asymptotic or fitted numerical rule.

\begin{quote}\footnotesize\textbf{E17}: \nolinkurl{PaperXI.DyadicGauss}, source lines 51--58.\par Contiguous source proof excerpt.\par Source SHA-256: {\scriptsize\texttt{634fe67a476f515f627455fa43a1ec870a4ccc66b87867cc5d40ee1227c91dea}}.\end{quote}

\VerbatimInput[fontsize=\scriptsize,numbers=left,firstnumber=51,breaklines=true,breakanywhere=true,numbersep=5pt]{excerpts/E17.render.txt}

\subsection{Counting refinements with a prescribed Arf invariant}
\label{guide:E18}
The supplied symplectic basis identifies the finite-dimensional nondegenerate alternating space. The source separately retains dimension zero and the odd-Arf count.
The formulas count the complete refinement space, not a sampled collection. An isometry preserving only the bilinear form need not preserve the prescribed Arf invariant or its phase.

\begin{quote}\footnotesize\textbf{E18}: \nolinkurl{PaperXI.DyadicRefinementCounts}, source lines 42--44.\par Literal declaration excerpt. A proof/value after a top-level assignment is omitted; structure fields are retained.\par Source SHA-256: {\scriptsize\texttt{5d9cdc46b7bde9d4f7ceed89b2a7a68ec38f36fee3142300f40b6a8d5301838b}}.\end{quote}

\VerbatimInput[fontsize=\scriptsize,numbers=left,firstnumber=42,breaklines=true,breakanywhere=true,numbersep=5pt]{excerpts/E18.render.txt}

\chapter{Matrix holonomy and global section equivalence}
Local packet comparisons can be transported around a graph. A rooted spanning tree determines root paths. The remaining original edges become chords. Full matrix compatibility on the resulting fundamental cycles is the necessary and sufficient obstruction to a based global gauge.

\subsection{The rooted chord criterion}
\label{guide:E19}
The graph is a connected quiver with invertible edge transports. A root comparison is fixed. The cycle condition is an equality of the full transported linear maps.
Trace, determinant and projective cycle tests are weaker than this matrix identity. The source constructs a gauge precisely when the chord tests hold and derives compatibility on every edge.

\begin{quote}\footnotesize\textbf{E19}: \nolinkurl{PaperXI.GraphGauge}, source lines 155--159.\par Literal declaration excerpt. A proof/value after a top-level assignment is omitted; structure fields are retained.\par Source SHA-256: {\scriptsize\texttt{1750d323d264d4760717c62ab3aa154cbebefcdd24e80cc76ef22d7d8f5a572b}}.\end{quote}

\VerbatimInput[fontsize=\scriptsize,numbers=left,firstnumber=155,breaklines=true,breakanywhere=true,numbersep=5pt]{excerpts/E19.render.txt}

\subsection{The space of compatible sections}
\label{guide:E20}
Each vertex has a vector-space fiber and each directed edge carries a linear equivalence. A section satisfies every edge equation.
The space is an actual linear subspace of the product of fibers. It can be zero. A nonzero section space, an irreducible section representation, and an arithmetic supercuspidal atom are separate assertions.

\begin{quote}\footnotesize\textbf{E20}: \nolinkurl{PaperXI.GluedPacketSections}, source lines 22--26.\par Contiguous source proof excerpt.\par Source SHA-256: {\scriptsize\texttt{5bf8047d59226ed80f93c4668b031aedab1df38bd5ee39456a9335f2685372bf}}.\end{quote}

\VerbatimInput[fontsize=\scriptsize,numbers=left,firstnumber=22,breaklines=true,breakanywhere=true,numbersep=5pt]{excerpts/E20.render.txt}

\subsection{Gluing produces a representation equivalence}
\label{guide:E21}
Covariance of the source and target edge systems is supplied. A root intertwiner and the full fundamental-cycle compatibility equations are supplied.
The resulting equivalence acts on the global compatible-section spaces. Its existence is proved by transporting the root intertwiner, rather than stored as a field in a comparison record.

\begin{quote}\footnotesize\textbf{E21}: \nolinkurl{PaperXI.GluedPacketSections}, source lines 138--144.\par Literal declaration excerpt. A proof/value after a top-level assignment is omitted; structure fields are retained.\par Source SHA-256: {\scriptsize\texttt{5bf8047d59226ed80f93c4668b031aedab1df38bd5ee39456a9335f2685372bf}}.\end{quote}

\VerbatimInput[fontsize=\scriptsize,numbers=left,firstnumber=138,breaklines=true,breakanywhere=true,numbersep=5pt]{excerpts/E21.render.txt}

\chapter{Monodromy and conditional atomic reconstruction}
Finite atomic results can propagate through a supplied decomposition and a supplied ordered-quotient interface. The reconstruction theorem constructs sorting, atomic replacement and descent. The external decomposition and ordered Langlands-quotient existence statements remain printed hypotheses.

\subsection{The monodromy normalization}
\label{guide:E22}
A twisting character \(\nu\) and a representation are fixed. Copy \(i\) has the twisted action \(\nu(g)^i\rho(g)\), and the nilpotent operator shifts copy index.
The declared relation is \(R(g)NR(g)^{-1}=\nu(g)^{-1}N\). The inverse on the twisting character fixes the convention; changing it would change the normalization.

\begin{quote}\footnotesize\textbf{E22}: \nolinkurl{PaperXI.MonodromyAtom}, source lines 64--66.\par Literal declaration excerpt. A proof/value after a top-level assignment is omitted; structure fields are retained.\par Source SHA-256: {\scriptsize\texttt{f53d79ddc74c43995955bc45c8bce600d821bb6d118e9426d85e8478c0829bf2}}.\end{quote}

\VerbatimInput[fontsize=\scriptsize,numbers=left,firstnumber=64,breaklines=true,breakanywhere=true,numbersep=5pt]{excerpts/E22.render.txt}

\subsection{Source decomposition is an explicit interface}
\label{guide:E23}
An isomorphism relation on source blocks is fixed. The record specifies the allowed displays and their intrinsic classwise uniqueness.
The formal reconstruction does not derive classification of all arithmetic source objects from arbitrary finite data. It uses this interface to make replacement independent of display choices.

\begin{quote}\footnotesize\textbf{E23}: \nolinkurl{PaperXI.AtomicReconstruction}, source lines 17--21.\par Literal declaration excerpt. A proof/value after a top-level assignment is omitted; structure fields are retained.\par Source SHA-256: {\scriptsize\texttt{e848f4ab5b8717eda8465737b7a726fbb877d3e6effc3728eaeab08698da52cb}}.\end{quote}

\VerbatimInput[fontsize=\scriptsize,numbers=left,firstnumber=17,breaklines=true,breakanywhere=true,numbersep=5pt]{excerpts/E23.render.txt}

\subsection{The reconstructed map and its uniqueness}
\label{guide:E24}
The source decomposition interface, ordered target-quotient interface and atomic replacement map are supplied. Ordering uses the specified target data.
The candidate and choice independence are constructed. This conditional theorem remains useful when each interface is instantiated by separate mathematics, but it is not an unrestricted correspondence existence theorem.

\begin{quote}\footnotesize\textbf{E24}: \nolinkurl{PaperXI.AtomicReconstruction}, source lines 143--144.\par Literal declaration excerpt. A proof/value after a top-level assignment is omitted; structure fields are retained.\par Source SHA-256: {\scriptsize\texttt{e848f4ab5b8717eda8465737b7a726fbb877d3e6effc3728eaeab08698da52cb}}.\end{quote}

\VerbatimInput[fontsize=\scriptsize,numbers=left,firstnumber=143,breaklines=true,breakanywhere=true,numbersep=5pt]{excerpts/E24.render.txt}

\chapter{Local rings, ramification projectors and normalized Frobenius}
The later local-field branch attaches finite residue extensions and normalized geometric Frobenius to actual algebraic objects. Lower ramification projectors average genuine finite groups. Their successive differences form depth-detail channels. A degree-only quotient forgets inertia, and the source proves exactly which detail information then vanishes.

\subsection{Lower congruence groups}
\label{guide:E25}
The action is on a supplied local ring, with its maximal-ideal powers. In a finite Galois DVR application these agree with the usual nonnegative lower groups under the stated dictionary.
The generic definition is based on congruences of the actual action. It is not asserted to be an absolute Weil ramification filtration without the arithmetic attachment.

\begin{quote}\footnotesize\textbf{E25}: \nolinkurl{PaperXI.ValuationRamification}, source lines 15--28.\par Contiguous source proof excerpt.\par Source SHA-256: {\scriptsize\texttt{92ddb509dc327389c2168f44ddfd5f7e16f4315af16835a54398077fe81146ca}}.\end{quote}

\VerbatimInput[fontsize=\scriptsize,numbers=left,firstnumber=15,breaklines=true,breakanywhere=true,numbersep=5pt]{excerpts/E25.render.txt}

\subsection{An exact depth-detail resolution}
\label{guide:E26}
The finite-group ramification filtration, representation and stabilization hypothesis are supplied. Averaging is over actual finite subgroups.
Nested fixed spaces give nested projectors. Their successive differences are orthogonal detail projectors, and the telescope reconstructs the identity at the stated stabilized depth.

\begin{quote}\footnotesize\textbf{E26}: \nolinkurl{PaperXI.RamificationPolyphase}, source lines 311--313.\par Literal declaration excerpt. A proof/value after a top-level assignment is omitted; structure fields are retained.\par Source SHA-256: {\scriptsize\texttt{2f34bd694def457221d6f00b874c8b54547dcc95ee249e3afcf34d71b5b64549}}.\end{quote}

\VerbatimInput[fontsize=\scriptsize,numbers=left,firstnumber=311,breaklines=true,breakanywhere=true,numbersep=5pt]{excerpts/E26.render.txt}

\subsection{The unconditional quotient target is the degree range}
\label{guide:E27}
The residue--Frobenius pullback has a degree homomorphism whose kernel is inertia. No degree-one lift is needed for this first-isomorphism statement.
The quotient is equivalent to the image of degree. To replace that image by all integer degrees, the source adds a degree-one Frobenius-lift hypothesis or a subsequent construction that supplies it.

\begin{quote}\footnotesize\textbf{E27}: \nolinkurl{PaperXI.RamifiedUnramifiedQuotient}, source lines 42--43.\par Literal declaration excerpt. A proof/value after a top-level assignment is omitted; structure fields are retained.\par Source SHA-256: {\scriptsize\texttt{755d457b0f11dfbac94eaee10326fa290bd5b6a331d8883bb8ce26811f00a747}}.\end{quote}

\VerbatimInput[fontsize=\scriptsize,numbers=left,firstnumber=42,breaklines=true,breakanywhere=true,numbersep=5pt]{excerpts/E27.render.txt}

\subsection{The constructed normalized quotient}
\label{guide:E28}
The normalized-Frobenius continuation uses the geometric lift constructed in its preceding modules. Its exact local-field and residue hypotheses are in the source setup.
This checked continuation fills the bounded normalized quotient step. Its retained report explicitly leaves full arithmetic/automorphic integration unestablished.

\begin{quote}\footnotesize\textbf{E28}: \nolinkurl{PaperXI.NormalizedFrobeniusQuotient}, source lines 79--79.\par Literal declaration excerpt. A proof/value after a top-level assignment is omitted; structure fields are retained.\par Source SHA-256: {\scriptsize\texttt{917cfa62d1b914d7d47a3c31f682a52d3338dc040a1606450f1cff71ca075b1f}}.\end{quote}

\VerbatimInput[fontsize=\scriptsize,numbers=left,firstnumber=79,breaklines=true,breakanywhere=true,numbersep=5pt]{excerpts/E28.render.txt}

\chapter{Arithmetic type and induction adapters}
The arithmetic branch uses actual local fields, continuous coordinate equivalences, literal subgroups, original characters, and induction carriers. Generic support and intertwining arguments derive irreducibility once the type dictionary is supplied. Later continuations reduce some dictionary leaves to quotient forms and actual lattice algebra.

\subsection{A candidate from three arithmetic profile leaves}
\label{guide:E29}
The ambient setup fixes local fields \(F,E\), coordinates, the literal subgroup \(D\), its character \(\theta\), the prescribed \(\kappa\), residue representation \(\sigma\), and the original tensor carrier. The displayed hypotheses keep trace support, character covariance and the oriented ambient index identity explicit.
The proof derives the required Heisenberg intersection dimension and feeds it into the inherited candidate constructor. It preserves the original carrier and representations. The existence of canonical type data and the actual truth of its retained source hypotheses are separate mathematical attachments.

\begin{quote}\footnotesize\textbf{E29}: \nolinkurl{PaperXI.ActualProfileCandidate110}, source lines 34--88.\par Literal declaration excerpt. A proof/value after a top-level assignment is omitted; structure fields are retained.\par Source SHA-256: {\scriptsize\texttt{dd434e182eec32397e64a3e6a43193cc56a3c965eff65c4e36e14f4bd0b2c27b}}.\end{quote}

\VerbatimInput[fontsize=\scriptsize,numbers=left,firstnumber=34,breaklines=true,breakanywhere=true,numbersep=5pt]{excerpts/E29.render.txt}

\subsection{The residue-rank-one specialization}
\label{guide:E30}
The residue rank is \(\mathrm{Fin}\,1\). The original character, carrier, support dictionary, beta-extension intertwining and occurrence input remain fixed.
The finite cuspidal residue factor and several normalization/support leaves are derived in this branch. It still does not manufacture every canonical simple character or establish the entire original claim.

\begin{quote}\footnotesize\textbf{E30}: \nolinkurl{PaperXI.RankOneSourceCandidate110}, source lines 29--60.\par Literal declaration excerpt. A proof/value after a top-level assignment is omitted; structure fields are retained.\par Source SHA-256: {\scriptsize\texttt{7f57ede141636d625efea3d9b67e905b1f34bcd5dd6c475f31a7e2ec2a7b2928}}.\end{quote}

\VerbatimInput[fontsize=\scriptsize,numbers=left,firstnumber=29,breaklines=true,breakanywhere=true,numbersep=5pt]{excerpts/E30.render.txt}

\section{Why orientation is load-bearing}
For the literal pair group, \(\iota(b)=g\iota(a)g^{-1}\). The first coordinate therefore belongs to
\[
 I_g=\{a:g\iota(a)g^{-1}\in\iota(K)\},
 \qquad \iota(I_g)=\iota(K)\cap g^{-1}\iota(K)g.
\]
The scalar intersection uses the same orientation. Under whole-action irreducibility and the original scalar/trace profile,
\[
 d^2=[K:H].
\]
The positive oriented index identity \([K:I_g]=[H:C_g]\) then yields \([I_g:C_g]=d^2\), giving the one-dimensional intertwining space. Schur is applied to the whole irreducible action, not to arbitrary intersection restrictions. Infinite subgroup index is encoded by zero in the relevant natural-valued index API, so positivity is essential to cancellation.

\chapter{Morita lattices and exact one-plus comparisons}
The centralizer coefficient ring acts inside a potentially noncommutative endomorphism ring. It must not be treated as central in the full ambient algebra. Morita corner arguments produce bases, coordinate quotient equivalences and actual additive indices. Multiplicative one-plus indices require genuine group closure and a common deep ideal.

\subsection{The actual Cartan index}
\label{guide:E31}
Actual ring homomorphisms embed centralizer matrices in the endomorphism ring. The finite/full lattice and radical-commutation conditions are explicit.
The index follows from actual coordinate conjugation and the preceding quotient computation. A numerical index formula is not inserted as a premise, but the arithmetic lattice instance must satisfy the printed hypotheses.

\begin{quote}\footnotesize\textbf{E31}: \nolinkurl{PaperXI.MoritaActualCartan110}, source lines 137--152.\par Literal declaration excerpt. A proof/value after a top-level assignment is omitted; structure fields are retained.\par Source SHA-256: {\scriptsize\texttt{dd18d3a60333f0e737ef4ff22c34db295f49460de701a9fddf9fcd3dc22d293a}}.\end{quote}

\VerbatimInput[fontsize=\scriptsize,numbers=left,firstnumber=137,breaklines=true,breakanywhere=true,numbersep=5pt]{excerpts/E31.render.txt}

\subsection{From additive lattices to unit groups}
\label{guide:E32}
The source keeps the original integral-order ring and its actual one-plus subgroups. A common ideal and a finite quotient allow the exact coset comparison.
An arbitrary additive lattice is not automatically a subgroup under \(1+x\), and a conjugate lattice need not be an ideal of the full field algebra. The theorem's carrier and group hypotheses prevent either substitution.

\begin{quote}\footnotesize\textbf{E32}: \nolinkurl{PaperXI.OnePlusIntersection110}, source lines 179--189.\par Literal declaration excerpt. A proof/value after a top-level assignment is omitted; structure fields are retained.\par Source SHA-256: {\scriptsize\texttt{512e3b7a894d7db7474f8586c4c652dddd45bfd98197f637f7be8e5a8fec3e7e}}.\end{quote}

\VerbatimInput[fontsize=\scriptsize,numbers=left,firstnumber=179,breaklines=true,breakanywhere=true,numbersep=5pt]{excerpts/E32.render.txt}

\subsection{The original carrier after dictionary reductions}
\label{guide:E33}
The literal type quotient/form, actual one-plus groups, Morita lattices and their coordinate matches are supplied. The source fixes the original character and beta extension.
Trace support, covariance, finite quotient, degree-square and index equality are derived from those lower-level instances. The canonical type dictionary and arithmetic category/factor interpretation still require an actual instance.

\begin{quote}\footnotesize\textbf{E33}: \nolinkurl{PaperXI.ActualSourceCandidate111}, source lines 40--84.\par Literal declaration excerpt. A proof/value after a top-level assignment is omitted; structure fields are retained.\par Source SHA-256: {\scriptsize\texttt{c802c3b3d1be72ca3d961a4edaa42cca066eb81cbc902a0a0abe9eee58d60005}}.\end{quote}

\VerbatimInput[fontsize=\scriptsize,numbers=left,firstnumber=40,breaklines=true,breakanywhere=true,numbersep=5pt]{excerpts/E33.render.txt}

\chapter{Complex scalar cochains and honest cyclic actions}
Equal scalar commutators determine a complex scalar cocycle class under the stated divisible-coefficient hypotheses. This can construct an actual representation of a supplied group. The original retained finite-group lift and the arithmetic source attachment remain separate when more structure is prescribed.

\subsection{What equal commutators supply over complex units}
\label{guide:E34}
The two normalized cocycles have coefficients in \(\mathbb C^\times\). Their commutators agree.
A cochain is constructed whose coboundary gives the cocycle ratio. This scalar-pushout result does not identify distinct marked finite central extensions such as the dyadic retained-group examples.

\begin{quote}\footnotesize\textbf{E34}: \nolinkurl{PaperXI.ScalarCocycle112}, source lines 171--174.\par Literal declaration excerpt. A proof/value after a top-level assignment is omitted; structure fields are retained.\par Source SHA-256: {\scriptsize\texttt{351ac3d1e21afeb0a8897f6eea1ba99f83bb63eaeee390eebefa44c1dfbb0f4e}}.\end{quote}

\VerbatimInput[fontsize=\scriptsize,numbers=left,firstnumber=171,breaklines=true,breakanywhere=true,numbersep=5pt]{excerpts/E34.render.txt}

\subsection{The equivariant actual-group action}
\label{guide:E35}
The actual groups and their full automorphisms are supplied. The scalar bridge has been corrected using the explicit cyclic coprimality conditions.
The result retains the original theta restriction and intertwines the prescribed packet generator. It requires full actions, not only an action on the symplectic quotient.

\begin{quote}\footnotesize\textbf{E35}: \nolinkurl{PaperXI.EquivariantRepresentation113}, source lines 67--91.\par Literal declaration excerpt. A proof/value after a top-level assignment is omitted; structure fields are retained.\par Source SHA-256: {\scriptsize\texttt{f1cf8ff3343f2647c0994270f6de5bc3478925a25406ea3a221861807b08b6c9}}.\end{quote}

\VerbatimInput[fontsize=\scriptsize,numbers=left,firstnumber=67,breaklines=true,breakanywhere=true,numbersep=5pt]{excerpts/E35.render.txt}

\subsection{Determinant-one normalization}
\label{guide:E36}
The generator implements the supplied full automorphism on the same irreducible carrier. Its cyclic order and carrier dimension satisfy the required coprimality.
Scalar rescaling can impose order and determinant normalization with uniqueness at the stated scope. The theorem does not identify the implementer with the independently prescribed arithmetic beta extension.

\begin{quote}\footnotesize\textbf{E36}: \nolinkurl{PaperXI.CyclicNormalization113}, source lines 167--173.\par Literal declaration excerpt. A proof/value after a top-level assignment is omitted; structure fields are retained.\par Source SHA-256: {\scriptsize\texttt{607587d94a544d2c4e014d81b21e4c9241d852feddfaa826df57f9715d396ce7}}.\end{quote}

\VerbatimInput[fontsize=\scriptsize,numbers=left,firstnumber=167,breaklines=true,breakanywhere=true,numbersep=5pt]{excerpts/E36.render.txt}

\chapter{Conjecture 66 and the surviving conditional comparison}
The finite bridge, the coherent graph equivalence, arithmetic realization and local-factor characterization are four different assertions. The final edition retains their order. Conjecture 66 concerns the general arithmetic realization step. The conditional comparison remains a theorem once the actual realization and the original factor characterization are supplied.

\subsection{The factor profile must characterize the declared range}
\label{guide:E37}
An atom type, its isomorphism setoid, a test type and value type are fixed. The record contains admissibility, invariant profiles and separation on that admissible range.
The separation field is the external characterization input. In an arithmetic application, all normalizations and the original characterizing family must be instantiated. A finite-group character instance is an example of the abstract interface, not a local-Weil converse theorem.

\begin{quote}\footnotesize\textbf{E37}: \nolinkurl{PaperXI.LocalFactorIdentification}, source lines 14--18.\par Literal declaration excerpt. A proof/value after a top-level assignment is omitted; structure fields are retained.\par Source SHA-256: {\scriptsize\texttt{3e81e1ce1c7b50d2771432e73d46f02de2a2ce4910e0a3420aa67e6b02a514f3}}.\end{quote}

\VerbatimInput[fontsize=\scriptsize,numbers=left,firstnumber=14,breaklines=true,breakanywhere=true,numbersep=5pt]{excerpts/E37.render.txt}

\subsection{Existence is an explicit arithmetic field}
\label{guide:E38}
The factor system and allowed datum type are fixed. A realizes relation connects each datum to possible actual atoms.
The record's \texttt{exists\_atom} supplies existence. Neither atom uniqueness nor reference agreement is stored here. General arithmetic existence of this record is the unresolved endpoint.

\begin{quote}\footnotesize\textbf{E38}: \nolinkurl{PaperXI.LocalFactorIdentification}, source lines 25--28.\par Literal declaration excerpt. A proof/value after a top-level assignment is omitted; structure fields are retained.\par Source SHA-256: {\scriptsize\texttt{3e81e1ce1c7b50d2771432e73d46f02de2a2ce4910e0a3420aa67e6b02a514f3}}.\end{quote}

\VerbatimInput[fontsize=\scriptsize,numbers=left,firstnumber=25,breaklines=true,breakanywhere=true,numbersep=5pt]{excerpts/E38.render.txt}

\subsection{H5 acts on the same realized candidate}
\label{guide:E39}
An actual realization is supplied. The realized atom's profile and the reference atom's profile agree with the same expected profile on the original test family.
Separation then proves equality of atom classes. H5 must refer to this same chosen completion. Replacing it with a factor identity for another candidate would not discharge the displayed hypothesis.

\begin{quote}\footnotesize\textbf{E39}: \nolinkurl{PaperXI.LocalFactorIdentification}, source lines 49--55.\par Contiguous source proof excerpt.\par Source SHA-256: {\scriptsize\texttt{3e81e1ce1c7b50d2771432e73d46f02de2a2ce4910e0a3420aa67e6b02a514f3}}.\end{quote}

\VerbatimInput[fontsize=\scriptsize,numbers=left,firstnumber=49,breaklines=true,breakanywhere=true,numbersep=5pt]{excerpts/E39.render.txt}

\subsection{The full formal conditional signature}
\label{guide:E40}
The finite source data construct a tensor bridge. Honest actions, phase compatibility and matrix holonomy supply the coherent finite comparison. The allowed target diagram, an arithmetic realization, admissibility and factor profiles are then supplied explicitly.
The conclusion contains a concrete compatible gauge, an equivalence of global section representations, and equality of the realized atom class with the reference. The source does not assert realization of every arbitrary finite datum. Its theorem name and successful Lean status must be read together with these hypotheses.

\begin{quote}\footnotesize\textbf{E40}: \nolinkurl{PaperXI.ConditionalAtomicComparison}, source lines 127--143.\par Literal declaration excerpt. A proof/value after a top-level assignment is omitted; structure fields are retained.\par Source SHA-256: {\scriptsize\texttt{ce8d1ea0f8658d4d968aeb6665d5f1d3e4cc6bb1d472ebfbfbce6215f1103a58}}.\end{quote}

\VerbatimInput[fontsize=\scriptsize,numbers=left,firstnumber=127,breaklines=true,breakanywhere=true,numbersep=5pt]{excerpts/E40.render.txt}

\section{The publication statement}
Fix the declared class \(\mathcal A\) of wild atoms and the paper's normalizations. For every \(\rho\in\mathcal A\) with supplied H1--H4 assembled finite packet datum \(d_\rho\), Conjecture 66 asserts that there exists one actual admissible supercuspidal representation \(\pi\) of \(\mathrm{GL}_{\dim\rho}(F)\) together with an arithmetic realization witness \(R(d_\rho,\pi)\), recovering the prescribed packet and preserving the declared stabilizer action, central-extension/phase data, and compatible transports. This witness uses actual arithmetic inducing data and explicit comparison/recovery. It does not require the whole assembled packet to be the original maximal-simple-type quotient.

Given that same \(\pi\) and realization witness, the conditional comparison proposition assumes original H5 for \(\pi\) on the stipulated characterizing family with the fixed normalization. The supplied characterization then gives \(\pi\simeq\mathrm{LLC}(\rho)\). This proposition has no claim that H5 produces a candidate, and no requirement that H5 hold for all possible realizers.

The retained general Lean interface has an \emph{all-realizers} profile binder because its \texttt{Realization} relation may have several realizers. The publication proposition needs the factor equality only for its one supplied candidate. The pointwise separation lemma \texttt{class\_eq\_iff} gives that implication directly. The broader interface should not be misread as strengthening the publication conjecture.

The traditional paper gives the realization definition, the declared source class with its supplied finite hypothesis data, and the retained hypothesis numbering. The release records the broad arithmetic realization ambition as a conjecture, the factor comparison as a conditional theorem, and the restricted constructions of Blocks 166--167 as ordinary mathematical realization results within their printed hypotheses. The guide does not add a Lean axiom named after the conjecture and does not substitute a \texttt{sorry} declaration for a proof.

The restricted constructions are not Lean-certified local-field/simple-character developments. Their proof notes and audits accompany the release. They retain the same-candidate H5 boundary. The dyadic low-rank enlargement changes the arithmetic assembly carrier while preserving the original candidate through induction in stages; it must not be described as enlarging a fixed maximal-type quotient that the dimension calculation obstructs.

\chapter{Verification lineage and reproducibility}
The retained verification reports certify specific source revisions and specific dependency closures. A full project-source scan, a focused compiler run, a transitive axiom audit, and a kernel replay are distinct checks. The evidence remains attributable to its original layer rather than collapsed into one global PASS flag.

\section{Recorded scopes}
The v0.44 record checks the new ramified/unramified quotient module, its 28 public declarations, focused imports and same-kernel replay. Its source-only reconstruction freshly builds the 20-module required project closure, reusing pinned Lean and Mathlib objects. The cumulative source scan covers 302 production modules.

The normalized-Frobenius record checks three new modules and 50 public declarations in two source states. Its clean project closure has 29 modules and agrees byte-for-byte with the primary object closure. It explicitly records that the full Block45 arithmetic/automorphic contract is not established.

Continuation110 records 20 source-set modules and 223 audited public exports. Continuation111 records nine new modules and 78 audited exports. Continuation112 records three new modules and 60 exports. Continuation113 records four modules and 76 exports, with a separately recorded internal compile/kernel audit. In each continuation, inherited modules and Mathlib are authenticated or reused at the documented scope rather than rebuilt in full.

\section{Failures remain part of the evidence}
Continuation110 retains an out-of-memory full-baseline attempt under an 8 GiB virtual-address-space cap and an interrupted larger attempt. The passing bounded new-module replay is not a claim that those failed full-baseline attempts succeeded.

Continuation113 retains the rejected theta0 arithmetic attachment, a corrected sign convention, failed harness/parser attempts and an overwritten early arithmetic receipt disclosed as unavailable. Its final corrected replay is a new accepted run, not a reconstruction of that missing receipt. The finite formal results and the failed arithmetic candidate are different claims.

\section{Portable runtime}
The preserved portable-runtime README describes an offline Linux x86\_64 kit for Lean 4.34.0-rc2 and the pinned Mathlib revision. Its full executable distribution is an external dependency where the source handoff does not include it. Native Windows launchers and the Linux/WSL path have different prerequisites. The release documentation gives the exact retained resources and availability.

A reader should first authenticate the runtime, Mathlib revision and branch archive hashes, then replay in a fresh output directory. Do not overwrite original receipts. The documented Continuation113 verifier checks the baseline110--112 archives and inherited objects, compiles only its four new modules, audits named exports and runs \texttt{leanchecker}. It does not claim to rebuild every historical Paper XI module or all of Mathlib.

\begin{Verbatim}[fontsize=\small,breaklines=true,breakanywhere=true]
python scripts/verify_new113.py \
  --kit /absolute/path/to/Komorebi_Portable_Lean_4.34.0-rc2 \
  --archives /absolute/path/to/baseline_archives \
  --out evidence/replay_on_new_machine
\end{Verbatim}

\section{Rebuilding this guide}
The editable folder contains the LaTeX source, displayed and raw excerpts, complete selected Lean source copies, source catalog and this PDF builder. From the checkpoint root, run the following command with a new output directory outside the checkpoint. XeLaTeX and the declared fonts/packages must already be installed. The builder installs nothing, uses no network and fixes the source date epoch for reproducibility.

\begin{Verbatim}[fontsize=\small,breaklines=true]
python3 tools/build_readers.py --output /absolute/path/new-output
\end{Verbatim}

\appendix\chapter{Complete selected named-source catalog}\label{guide:catalog}

Every selected module below is copied unchanged into the editable guide's \texttt{sources/PaperXI} folder. The catalog gives module identity, selected layer, SHA-256, direct project imports, module documentation and every named lexical declaration with its raw source line. These line pins locate declarations in the complete source, not independently compiled snippets. Earlier repeated source-layer identities remain in \texttt{SOURCE\_VERSIONS.json}.

\section{\texorpdfstring{\nolinkurl{PaperXI.AbstractHeisenberg}}{PaperXI.AbstractHeisenberg}}\label{module:PaperXI.AbstractHeisenberg}

\textbf{Selected layer:} v0.44 parent.\quad\textbf{Source lines:} 184.\par\textbf{SHA-256:} {\scriptsize\nolinkurl{d83045df228867234c49fb9dd9e838bbf810394614dd22dce2401c6744e3dacd}}.

\VerbatimInput[fontsize=\scriptsize,breaklines=true,breakanywhere=true]{excerpts/D_AbstractHeisenberg.txt}

\textbf{Direct project imports:} \nolinkurl{PaperXI.CentralExtension}.

\begin{longtable}{@{}p{0.07\textwidth}p{0.12\textwidth}p{0.75\textwidth}@{}}\toprule Line&Kind&Name\\\midrule\endhead

18&structure&\nolinkurl{PaperXI.AbstractHeisenberg.Element}\\

24&instance&\nolinkurl{PaperXI.AbstractHeisenberg.instMul}\\

28&instance&\nolinkurl{PaperXI.AbstractHeisenberg.instOne}\\

33&instance&\nolinkurl{PaperXI.AbstractHeisenberg.instInv}\\

36&theorem&\nolinkurl{PaperXI.AbstractHeisenberg.mul_vector}\\

37&theorem&\nolinkurl{PaperXI.AbstractHeisenberg.mul_central}\\

39&theorem&\nolinkurl{PaperXI.AbstractHeisenberg.one_vector}\\

40&theorem&\nolinkurl{PaperXI.AbstractHeisenberg.one_central}\\

41&theorem&\nolinkurl{PaperXI.AbstractHeisenberg.inv_vector}\\

42&theorem&\nolinkurl{PaperXI.AbstractHeisenberg.inv_central}\\

45&instance&\nolinkurl{PaperXI.AbstractHeisenberg.instGroup}\\

55&def&\nolinkurl{PaperXI.AbstractHeisenberg.coordinates}\\

61&instance&\nolinkurl{PaperXI.AbstractHeisenberg.instFinite}\\

65&theorem&\nolinkurl{PaperXI.AbstractHeisenberg.finite_of_finite_field}\\

70&def&\nolinkurl{PaperXI.AbstractHeisenberg.centralElement}\\

72&theorem&\nolinkurl{PaperXI.AbstractHeisenberg.centralElement_vector}\\

73&theorem&\nolinkurl{PaperXI.AbstractHeisenberg.centralElement_central}\\

74&theorem&\nolinkurl{PaperXI.AbstractHeisenberg.centralElement_zero}\\

75&theorem&\nolinkurl{PaperXI.AbstractHeisenberg.centralElement_mul}\\

79&theorem&\nolinkurl{PaperXI.AbstractHeisenberg.centralElement_commute}\\

84&def&\nolinkurl{PaperXI.AbstractHeisenberg.inclusion}\\

90&def&\nolinkurl{PaperXI.AbstractHeisenberg.quotient}\\

95&theorem&\nolinkurl{PaperXI.AbstractHeisenberg.inclusion_apply}\\

97&theorem&\nolinkurl{PaperXI.AbstractHeisenberg.quotient_apply}\\

101&def&\nolinkurl{PaperXI.AbstractHeisenberg.datum}\\

121&theorem&\nolinkurl{PaperXI.AbstractHeisenberg.datum_i}\\

122&theorem&\nolinkurl{PaperXI.AbstractHeisenberg.datum_q}\\

125&theorem&\nolinkurl{PaperXI.AbstractHeisenberg.inv_central_of_isAlt}\\

129&theorem&\nolinkurl{PaperXI.AbstractHeisenberg.mul_swap}\\

140&theorem&\nolinkurl{PaperXI.AbstractHeisenberg.commutator}\\

151&def&\nolinkurl{PaperXI.AbstractHeisenberg.lift}\\

163&theorem&\nolinkurl{PaperXI.AbstractHeisenberg.lift_vector}\\

167&theorem&\nolinkurl{PaperXI.AbstractHeisenberg.lift_central}\\

171&theorem&\nolinkurl{PaperXI.AbstractHeisenberg.lift_inclusion}\\

177&theorem&\nolinkurl{PaperXI.AbstractHeisenberg.lift_quotient}\\

\bottomrule\end{longtable}

\section{\texorpdfstring{\nolinkurl{PaperXI.AbstractHeisenbergExtension}}{PaperXI.AbstractHeisenbergExtension}}\label{module:PaperXI.AbstractHeisenbergExtension}

\textbf{Selected layer:} v0.44 parent.\quad\textbf{Source lines:} 58.\par\textbf{SHA-256:} {\scriptsize\nolinkurl{64c993ffe68d2705763bdb7173dcdb8096ddb3540a82363620d5ac0e820f8d06}}.

\VerbatimInput[fontsize=\scriptsize,breaklines=true,breakanywhere=true]{excerpts/D_AbstractHeisenbergExtension.txt}

\textbf{Direct project imports:} \nolinkurl{PaperXI.AbstractHeisenberg}, \nolinkurl{PaperXI.CentralExtensionCocycle}, \nolinkurl{PaperXI.SymplecticCentralExtension}.

\begin{longtable}{@{}p{0.07\textwidth}p{0.12\textwidth}p{0.75\textwidth}@{}}\toprule Line&Kind&Name\\\midrule\endhead

15&theorem&\nolinkurl{PaperXI.AbstractHeisenberg.centralCommutator_eq}\\

23&theorem&\nolinkurl{PaperXI.AbstractHeisenberg.scalar_commutator}\\

33&theorem&\nolinkurl{PaperXI.AbstractHeisenberg.scalar_radical_eq_bot}\\

43&theorem&\nolinkurl{PaperXI.AbstractHeisenberg.scalar_radical_eq_bot_scalarCharacter}\\

52&theorem&\nolinkurl{PaperXI.AbstractHeisenberg.center_eq_inclusion_range}\\

\bottomrule\end{longtable}

\section{\texorpdfstring{\nolinkurl{PaperXI.ActualCartanCandidate106}}{PaperXI.ActualCartanCandidate106}}\label{module:PaperXI.ActualCartanCandidate106}

\textbf{Selected layer:} base108 arithmetic branch.\quad\textbf{Source lines:} 79.\par\textbf{SHA-256:} {\scriptsize\nolinkurl{e87ad679e0338ba7b5ffda904079152531d758cd29a472151aa21431bc543d06}}.

\VerbatimInput[fontsize=\scriptsize,breaklines=true,breakanywhere=true]{excerpts/D_ActualCartanCandidate106.txt}

\textbf{Direct project imports:} \nolinkurl{PaperXI.ActualLocalCartan106}, \nolinkurl{PaperXI.ActualCuspidalCandidate105}.

\begin{longtable}{@{}p{0.07\textwidth}p{0.12\textwidth}p{0.75\textwidth}@{}}\toprule Line&Kind&Name\\\midrule\endhead

31&theorem&\nolinkurl{PaperXI.ActualCartanCandidate106.candidate_exists}\\

\bottomrule\end{longtable}

\section{\texorpdfstring{\nolinkurl{PaperXI.ActualCharacterCandidate108}}{PaperXI.ActualCharacterCandidate108}}\label{module:PaperXI.ActualCharacterCandidate108}

\textbf{Selected layer:} base108 arithmetic branch.\quad\textbf{Source lines:} 89.\par\textbf{SHA-256:} {\scriptsize\nolinkurl{a311e3796f48d7d7a6c2cafd95e7ff962b8fb7543a0cdb7a81d04d9d12dbe4f2}}.

\VerbatimInput[fontsize=\scriptsize,breaklines=true,breakanywhere=true]{excerpts/D_ActualCharacterCandidate108.txt}

\textbf{Direct project imports:} \nolinkurl{PaperXI.ActualCharacterData108}.

\begin{longtable}{@{}p{0.07\textwidth}p{0.12\textwidth}p{0.75\textwidth}@{}}\toprule Line&Kind&Name\\\midrule\endhead

31&theorem&\nolinkurl{PaperXI.ActualCharacterCandidate108.candidate_exists}\\

\bottomrule\end{longtable}

\section{\texorpdfstring{\nolinkurl{PaperXI.ActualCharacterData108}}{PaperXI.ActualCharacterData108}}\label{module:PaperXI.ActualCharacterData108}

\textbf{Selected layer:} base108 arithmetic branch.\quad\textbf{Source lines:} 49.\par\textbf{SHA-256:} {\scriptsize\nolinkurl{adcb8f31671a6ce4587678730fd689cdda4c87ebd3f1c8c67391cf42ee921bc4}}.

\VerbatimInput[fontsize=\scriptsize,breaklines=true,breakanywhere=true]{excerpts/D_ActualCharacterData108.txt}

\textbf{Direct project imports:} \nolinkurl{PaperXI.ActualHeisenbergCandidate107}, \nolinkurl{PaperXI.CharacterSupport108}.

\begin{longtable}{@{}p{0.07\textwidth}p{0.12\textwidth}p{0.75\textwidth}@{}}\toprule Line&Kind&Name\\\midrule\endhead

22&theorem&\nolinkurl{PaperXI.ActualCharacterData108.domain_le_compact}\\

26&def&\nolinkurl{PaperXI.ActualCharacterData108.intoCompact}\\

30&theorem&\nolinkurl{PaperXI.ActualCharacterData108.intoCompact_mem_first}\\

37&def&\nolinkurl{PaperXI.ActualCharacterData108.Contains}\\

41&theorem&\nolinkurl{PaperXI.ActualCharacterData108.scalar_restriction}\\

\bottomrule\end{longtable}

\section{\texorpdfstring{\nolinkurl{PaperXI.ActualCharacterIdentification61}}{PaperXI.ActualCharacterIdentification61}}\label{module:PaperXI.ActualCharacterIdentification61}

\textbf{Selected layer:} base108 arithmetic branch.\quad\textbf{Source lines:} 100.\par\textbf{SHA-256:} {\scriptsize\nolinkurl{8507c7ccc941069745d0352c942ee3b600e8fb68441dc2ecd4d458f31ce367fc}}.

\VerbatimInput[fontsize=\scriptsize,breaklines=true,breakanywhere=true]{excerpts/D_ActualCharacterIdentification61.txt}

\textbf{Direct project imports:} \nolinkurl{PaperXI.ActualCharacterType61}.

\begin{longtable}{@{}p{0.07\textwidth}p{0.12\textwidth}p{0.75\textwidth}@{}}\toprule Line&Kind&Name\\\midrule\endhead

27&theorem&\nolinkurl{PaperXI.ActualCharacterIdentification61.compactGroup_eq_of_full_normalizer}\\

59&theorem&\nolinkurl{PaperXI.ActualCharacterIdentification61.uniformizer_mem_full_normalizer}\\

69&theorem&\nolinkurl{PaperXI.ActualCharacterIdentification61.extendedSubgroup_eq_full_normalizer}\\

\bottomrule\end{longtable}

\section{\texorpdfstring{\nolinkurl{PaperXI.ActualCharacterResidue61}}{PaperXI.ActualCharacterResidue61}}\label{module:PaperXI.ActualCharacterResidue61}

\textbf{Selected layer:} base108 arithmetic branch.\quad\textbf{Source lines:} 133.\par\textbf{SHA-256:} {\scriptsize\nolinkurl{401e10d5858bab78cc6a399bac1bd8c3d594a891773157e9ff975c2a7f0b805f}}.

\VerbatimInput[fontsize=\scriptsize,breaklines=true,breakanywhere=true]{excerpts/D_ActualCharacterResidue61.txt}

\textbf{Direct project imports:} \nolinkurl{PaperXI.ActualCharacterType61}, \nolinkurl{PaperXI.ActualELatticeFieldJoin59}.

\begin{longtable}{@{}p{0.07\textwidth}p{0.12\textwidth}p{0.75\textwidth}@{}}\toprule Line&Kind&Name\\\midrule\endhead

29&def&\nolinkurl{PaperXI.ActualCharacterResidue61.intoModel}\\

35&def&\nolinkurl{PaperXI.ActualCharacterResidue61.reduction}\\

40&theorem&\nolinkurl{PaperXI.ActualCharacterResidue61.reduction_kernel}\\

48&theorem&\nolinkurl{PaperXI.ActualCharacterResidue61.integralMatrix_mem_compactGroup}\\

60&theorem&\nolinkurl{PaperXI.ActualCharacterResidue61.integralUnit_mem_compactGroup}\\

76&theorem&\nolinkurl{PaperXI.ActualCharacterResidue61.reduction_surjective}\\

110&theorem&\nolinkurl{PaperXI.ActualCharacterResidue61.reduction_conjugation}\\

\bottomrule\end{longtable}

\section{\texorpdfstring{\nolinkurl{PaperXI.ActualCharacterType61}}{PaperXI.ActualCharacterType61}}\label{module:PaperXI.ActualCharacterType61}

\textbf{Selected layer:} base108 arithmetic branch.\quad\textbf{Source lines:} 163.\par\textbf{SHA-256:} {\scriptsize\nolinkurl{31cbff49441e1406a497ce67cbba0965509582286af1f2fc8337525909e25c1f}}.

\VerbatimInput[fontsize=\scriptsize,breaklines=true,breakanywhere=true]{excerpts/D_ActualCharacterType61.txt}

\textbf{Direct project imports:} \nolinkurl{PaperXI.CharacterNormalizer61}, \nolinkurl{PaperXI.ActualELatticeTopology60}.

\begin{longtable}{@{}p{0.07\textwidth}p{0.12\textwidth}p{0.75\textwidth}@{}}\toprule Line&Kind&Name\\\midrule\endhead

23&theorem&\nolinkurl{PaperXI.ActualCharacterType61.power_normalizes_chainStabilizer}\\

53&def&\nolinkurl{PaperXI.ActualCharacterType61.chainGroup}\\

58&theorem&\nolinkurl{PaperXI.ActualCharacterType61.matrixCongruence_le_chainGroup}\\

62&theorem&\nolinkurl{PaperXI.ActualCharacterType61.chainGroup_open}\\

68&theorem&\nolinkurl{PaperXI.ActualCharacterType61.chainGroup_compact}\\

83&theorem&\nolinkurl{PaperXI.ActualCharacterType61.coordinateAut_unitsEmbedding}\\

96&theorem&\nolinkurl{PaperXI.ActualCharacterType61.coordinateAut_embeddedUniformizer}\\

110&theorem&\nolinkurl{PaperXI.ActualCharacterType61.uniformizer_normalizes_chainGroup}\\

126&def&\nolinkurl{PaperXI.ActualCharacterType61.compactGroup}\\

130&def&\nolinkurl{PaperXI.ActualCharacterType61.firstSubgroup}\\

135&theorem&\nolinkurl{PaperXI.ActualCharacterType61.domain_le_compactGroup}\\

140&theorem&\nolinkurl{PaperXI.ActualCharacterType61.compactGroup_compact_open}\\

148&theorem&\nolinkurl{PaperXI.ActualCharacterType61.firstSubgroup_open}\\

154&theorem&\nolinkurl{PaperXI.ActualCharacterType61.uniformizer_normalizes_compactGroup}\\

\bottomrule\end{longtable}

\section{\texorpdfstring{\nolinkurl{PaperXI.ActualCompactTypeCompletion58}}{PaperXI.ActualCompactTypeCompletion58}}\label{module:PaperXI.ActualCompactTypeCompletion58}

\textbf{Selected layer:} base108 arithmetic branch.\quad\textbf{Source lines:} 97.\par\textbf{SHA-256:} {\scriptsize\nolinkurl{9c0e0c63cd7aa7f6c5ce5841cf9dbd3b97709ee0300314cb808dae8caef66b2a}}.

\VerbatimInput[fontsize=\scriptsize,breaklines=true,breakanywhere=true]{excerpts/D_ActualCompactTypeCompletion58.txt}

\textbf{Direct project imports:} \nolinkurl{PaperXI.ActualNormalizerTypeExtension58}, \nolinkurl{PaperXI.CompactTypeTensorNormalizer58}.

\begin{longtable}{@{}p{0.07\textwidth}p{0.12\textwidth}p{0.75\textwidth}@{}}\toprule Line&Kind&Name\\\midrule\endhead

25&theorem&\nolinkurl{PaperXI.ActualCompactTypeCompletion58.firstFactor_irreducible}\\

53&theorem&\nolinkurl{PaperXI.ActualCompactTypeCompletion58.completed}\\

\bottomrule\end{longtable}

\section{\texorpdfstring{\nolinkurl{PaperXI.ActualCoordinateDeepIdeal111}}{PaperXI.ActualCoordinateDeepIdeal111}}\label{module:PaperXI.ActualCoordinateDeepIdeal111}

\textbf{Selected layer:} Continuation111.\quad\textbf{Source lines:} 80.\par\textbf{SHA-256:} {\scriptsize\nolinkurl{1478145467bf2f189641de88b717440f1439dc649d84164ad972b32966afeb91}}.

\VerbatimInput[fontsize=\scriptsize,breaklines=true,breakanywhere=true]{excerpts/D_ActualCoordinateDeepIdeal111.txt}

\textbf{Direct project imports:} \nolinkurl{PaperXI.CoordinateDeepIdeal111}.

\begin{longtable}{@{}p{0.07\textwidth}p{0.12\textwidth}p{0.75\textwidth}@{}}\toprule Line&Kind&Name\\\midrule\endhead

14&def&\nolinkurl{PaperXI.ActualCoordinateDeepIdeal111.innerHomeomorph}\\

26&def&\nolinkurl{PaperXI.ActualCoordinateDeepIdeal111.integralEmbedding}\\

29&theorem&\nolinkurl{PaperXI.ActualCoordinateDeepIdeal111.integralEmbedding_open}\\

37&theorem&\nolinkurl{PaperXI.ActualCoordinateDeepIdeal111.conjugate_lattice_open}\\

49&theorem&\nolinkurl{PaperXI.ActualCoordinateDeepIdeal111.exists_conjugate_common_ideal}\\

66&theorem&\nolinkurl{PaperXI.ActualCoordinateDeepIdeal111.onePlus_conjugate_relIndex}\\

\bottomrule\end{longtable}

\section{\texorpdfstring{\nolinkurl{PaperXI.ActualCuspidalCandidate105}}{PaperXI.ActualCuspidalCandidate105}}\label{module:PaperXI.ActualCuspidalCandidate105}

\textbf{Selected layer:} base108 arithmetic branch.\quad\textbf{Source lines:} 118.\par\textbf{SHA-256:} {\scriptsize\nolinkurl{b37fd23618fec1d17f60baca5392c3d0d4e8c255da68dfff131b5ca746840100}}.

\VerbatimInput[fontsize=\scriptsize,breaklines=true,breakanywhere=true]{excerpts/D_ActualCuspidalCandidate105.txt}

\textbf{Direct project imports:} \nolinkurl{PaperXI.FullTypeIntertwining105}, \nolinkurl{PaperXI.ActualPacketTypeRealization70}.

\begin{longtable}{@{}p{0.07\textwidth}p{0.12\textwidth}p{0.75\textwidth}@{}}\toprule Line&Kind&Name\\\midrule\endhead

26&def&\nolinkurl{PaperXI.ActualCuspidalCandidate105.centralizerMap}\\

29&abbrev&\nolinkurl{PaperXI.ActualCuspidalCandidate105.compact}\\

30&abbrev&\nolinkurl{PaperXI.ActualCuspidalCandidate105.extended}\\

36&theorem&\nolinkurl{PaperXI.ActualCuspidalCandidate105.centralizer_scalar}\\

54&theorem&\nolinkurl{PaperXI.ActualCuspidalCandidate105.candidate_exists}\\

\bottomrule\end{longtable}

\section{\texorpdfstring{\nolinkurl{PaperXI.ActualDepthZeroExtension}}{PaperXI.ActualDepthZeroExtension}}\label{module:PaperXI.ActualDepthZeroExtension}

\textbf{Selected layer:} base108 arithmetic branch.\quad\textbf{Source lines:} 155.\par\textbf{SHA-256:} {\scriptsize\nolinkurl{ff5750ec70106f6aa0ce85b66fc16095d1d17973b914f7b37b2e2c1efdcda17c}}.

\VerbatimInput[fontsize=\scriptsize,breaklines=true,breakanywhere=true]{excerpts/D_ActualDepthZeroExtension.txt}

\textbf{Direct project imports:} \nolinkurl{PaperXI.ActualIntegralGeneralLinear}, \nolinkurl{PaperXI.ActualLocalUnitDecomposition}, \nolinkurl{PaperXI.CentralExtensionInductionConsequences}.

\begin{longtable}{@{}p{0.07\textwidth}p{0.12\textwidth}p{0.75\textwidth}@{}}\toprule Line&Kind&Name\\\midrule\endhead

25&def&\nolinkurl{PaperXI.ActualDepthZeroExtension.determinantDegree}\\

29&def&\nolinkurl{PaperXI.ActualDepthZeroExtension.scalarUniformizer}\\

32&theorem&\nolinkurl{PaperXI.ActualDepthZeroExtension.scalarUniformizer_central}\\

35&theorem&\nolinkurl{PaperXI.ActualDepthZeroExtension.determinantDegree_embedding}\\

41&theorem&\nolinkurl{PaperXI.ActualDepthZeroExtension.integralSubgroup_le_degree_kernel}\\

46&theorem&\nolinkurl{PaperXI.ActualDepthZeroExtension.determinantDegree_scalarUniformizer}\\

56&theorem&\nolinkurl{PaperXI.ActualDepthZeroExtension.scalarUniformizer_free}\\

66&abbrev&\nolinkurl{PaperXI.ActualDepthZeroExtension.extendedSubgroup}\\

71&theorem&\nolinkurl{PaperXI.ActualDepthZeroExtension.integral_le_extended}\\

76&theorem&\nolinkurl{PaperXI.ActualDepthZeroExtension.scalar_unit_mem_integral}\\

83&theorem&\nolinkurl{PaperXI.ActualDepthZeroExtension.extendedSubgroup_eq_scalar_sup}\\

106&theorem&\nolinkurl{PaperXI.ActualDepthZeroExtension.extendedSubgroup_open}\\

111&abbrev&\nolinkurl{PaperXI.ActualDepthZeroExtension.inclusion}\\

118&abbrev&\nolinkurl{PaperXI.ActualDepthZeroExtension.extendedRepresentation}\\

122&theorem&\nolinkurl{PaperXI.ActualDepthZeroExtension.extendedRepresentation_restrict}\\

126&theorem&\nolinkurl{PaperXI.ActualDepthZeroExtension.extendedRepresentation_irreducible}\\

131&theorem&\nolinkurl{PaperXI.ActualDepthZeroExtension.extendedRepresentation_smooth}\\

137&abbrev&\nolinkurl{PaperXI.ActualDepthZeroExtension.inducedRepresentation}\\

142&theorem&\nolinkurl{PaperXI.ActualDepthZeroExtension.induced_smooth}\\

146&theorem&\nolinkurl{PaperXI.ActualDepthZeroExtension.induced_nontrivial}\\

151&theorem&\nolinkurl{PaperXI.ActualDepthZeroExtension.induced_scalar}\\

\bottomrule\end{longtable}

\section{\texorpdfstring{\nolinkurl{PaperXI.ActualELatticeFieldJoin59}}{PaperXI.ActualELatticeFieldJoin59}}\label{module:PaperXI.ActualELatticeFieldJoin59}

\textbf{Selected layer:} base108 arithmetic branch.\quad\textbf{Source lines:} 80.\par\textbf{SHA-256:} {\scriptsize\nolinkurl{cd8801f07f19c6130f8e0cb4b83972ac28c399481be71c0f8cee79f7e9bbd660}}.

\VerbatimInput[fontsize=\scriptsize,breaklines=true,breakanywhere=true]{excerpts/D_ActualELatticeFieldJoin59.txt}

\textbf{Direct project imports:} \nolinkurl{PaperXI.ActualELatticeQuotient59}, \nolinkurl{PaperXI.ActualELatticeFreeness59}.

\begin{longtable}{@{}p{0.07\textwidth}p{0.12\textwidth}p{0.75\textwidth}@{}}\toprule Line&Kind&Name\\\midrule\endhead

24&theorem&\nolinkurl{PaperXI.ActualELatticeFieldJoin59.integralMatrix_scalar}\\

34&theorem&\nolinkurl{PaperXI.ActualELatticeFieldJoin59.integralUnit_mem_typeGroup}\\

40&theorem&\nolinkurl{PaperXI.ActualELatticeFieldJoin59.extendedSubgroup_eq_field_join}\\

60&theorem&\nolinkurl{PaperXI.ActualELatticeFieldJoin59.field_normalizes_typeGroup}\\

71&theorem&\nolinkurl{PaperXI.ActualELatticeFieldJoin59.extendedSubgroup_coe_eq_field_mul}\\

\bottomrule\end{longtable}

\section{\texorpdfstring{\nolinkurl{PaperXI.ActualELatticeFreeness59}}{PaperXI.ActualELatticeFreeness59}}\label{module:PaperXI.ActualELatticeFreeness59}

\textbf{Selected layer:} base108 arithmetic branch.\quad\textbf{Source lines:} 73.\par\textbf{SHA-256:} {\scriptsize\nolinkurl{b9c837f72cdde593dc7a1bf64f366e0e5453e390448eedf225b0024251c84a47}}.

\VerbatimInput[fontsize=\scriptsize,breaklines=true,breakanywhere=true]{excerpts/D_ActualELatticeFreeness59.txt}

\textbf{Direct project imports:} \nolinkurl{PaperXI.ActualELatticeTypeGroup59}.

\begin{longtable}{@{}p{0.07\textwidth}p{0.12\textwidth}p{0.75\textwidth}@{}}\toprule Line&Kind&Name\\\midrule\endhead

21&def&\nolinkurl{PaperXI.ActualELatticeFreeness59.scalarAction}\\

28&theorem&\nolinkurl{PaperXI.ActualELatticeFreeness59.scalarAction_apply}\\

33&theorem&\nolinkurl{PaperXI.ActualELatticeFreeness59.uniformizerAction_zpow_apply}\\

41&theorem&\nolinkurl{PaperXI.ActualELatticeFreeness59.typeGroup_le_stabilizer}\\

49&theorem&\nolinkurl{PaperXI.ActualELatticeFreeness59.uniformizerAction_exponent_nonneg}\\

66&theorem&\nolinkurl{PaperXI.ActualELatticeFreeness59.uniformizerAction_free}\\

\bottomrule\end{longtable}

\section{\texorpdfstring{\nolinkurl{PaperXI.ActualELatticeNormalizer59}}{PaperXI.ActualELatticeNormalizer59}}\label{module:PaperXI.ActualELatticeNormalizer59}

\textbf{Selected layer:} base108 arithmetic branch.\quad\textbf{Source lines:} 129.\par\textbf{SHA-256:} {\scriptsize\nolinkurl{7328521c5eb95fe067aaae8fbb771a976d75965543e6cb1619720f6df9d68c17}}.

\VerbatimInput[fontsize=\scriptsize,breaklines=true,breakanywhere=true]{excerpts/D_ActualELatticeNormalizer59.txt}

\textbf{Direct project imports:} \nolinkurl{PaperXI.ActualIntegralGeneralLinear}, \nolinkurl{PaperXI.ActualLocalUnitDecomposition}, \nolinkurl{PaperXI.NormalizerExtensionInducing58}.

\begin{longtable}{@{}p{0.07\textwidth}p{0.12\textwidth}p{0.75\textwidth}@{}}\toprule Line&Kind&Name\\\midrule\endhead

24&def&\nolinkurl{PaperXI.ActualELatticeNormalizer59.chainStabilizer}\\

35&def&\nolinkurl{PaperXI.ActualELatticeNormalizer59.firstCongruence}\\

59&theorem&\nolinkurl{PaperXI.ActualELatticeNormalizer59.firstCongruence_le}\\

63&theorem&\nolinkurl{PaperXI.ActualELatticeNormalizer59.stabilizer_normalizes_firstCongruence}\\

84&def&\nolinkurl{PaperXI.ActualELatticeNormalizer59.periodicChain}\\

88&theorem&\nolinkurl{PaperXI.ActualELatticeNormalizer59.periodicChain_shift}\\

96&theorem&\nolinkurl{PaperXI.ActualELatticeNormalizer59.power_normalizes_firstCongruence}\\

\bottomrule\end{longtable}

\section{\texorpdfstring{\nolinkurl{PaperXI.ActualELatticeQuotient59}}{PaperXI.ActualELatticeQuotient59}}\label{module:PaperXI.ActualELatticeQuotient59}

\textbf{Selected layer:} base108 arithmetic branch.\quad\textbf{Source lines:} 89.\par\textbf{SHA-256:} {\scriptsize\nolinkurl{a49f9630deecf93df96b4f831fb702cd48b0154c5563d3fe4f919ad2c9f47d24}}.

\VerbatimInput[fontsize=\scriptsize,breaklines=true,breakanywhere=true]{excerpts/D_ActualELatticeQuotient59.txt}

\textbf{Direct project imports:} \nolinkurl{PaperXI.ActualELatticeTypeGroup59}, \nolinkurl{PaperXI.ProductNormalizerQuotient59}.

\begin{longtable}{@{}p{0.07\textwidth}p{0.12\textwidth}p{0.75\textwidth}@{}}\toprule Line&Kind&Name\\\midrule\endhead

22&def&\nolinkurl{PaperXI.ActualELatticeQuotient59.typeCongruence}\\

25&instance&\nolinkurl{PaperXI.ActualELatticeQuotient59.typeCongruence_normal}\\

29&def&\nolinkurl{PaperXI.ActualELatticeQuotient59.residueEquiv}\\

37&def&\nolinkurl{PaperXI.ActualELatticeQuotient59.reduction}\\

40&theorem&\nolinkurl{PaperXI.ActualELatticeQuotient59.reduction_surjective}\\

43&theorem&\nolinkurl{PaperXI.ActualELatticeQuotient59.reduction_kernel}\\

50&theorem&\nolinkurl{PaperXI.ActualELatticeQuotient59.finite_quotient}\\

54&theorem&\nolinkurl{PaperXI.ActualELatticeQuotient59.uniformizer_centralizes}\\

62&theorem&\nolinkurl{PaperXI.ActualELatticeQuotient59.uniformizer_normalizes_congruence}\\

68&theorem&\nolinkurl{PaperXI.ActualELatticeQuotient59.uniformizer_normalizes}\\

75&def&\nolinkurl{PaperXI.ActualELatticeQuotient59.uniformizerConjugation}\\

81&theorem&\nolinkurl{PaperXI.ActualELatticeQuotient59.reduction_conjugation}\\

\bottomrule\end{longtable}

\section{\texorpdfstring{\nolinkurl{PaperXI.ActualELatticeResidue59}}{PaperXI.ActualELatticeResidue59}}\label{module:PaperXI.ActualELatticeResidue59}

\textbf{Selected layer:} base108 arithmetic branch.\quad\textbf{Source lines:} 157.\par\textbf{SHA-256:} {\scriptsize\nolinkurl{7adc864addc36f19b3d94431dcbaa62f118e45633f9bc356692c1c8c0698d5a3}}.

\VerbatimInput[fontsize=\scriptsize,breaklines=true,breakanywhere=true]{excerpts/D_ActualELatticeResidue59.txt}

\textbf{Direct project imports:} \nolinkurl{PaperXI.ActualIntegralGeneralLinear}, \nolinkurl{PaperXI.ActualLocalUnitDecomposition}.

\begin{longtable}{@{}p{0.07\textwidth}p{0.12\textwidth}p{0.75\textwidth}@{}}\toprule Line&Kind&Name\\\midrule\endhead

20&def&\nolinkurl{PaperXI.ActualELatticeResidue59.restrictScalarMatrix}\\

25&theorem&\nolinkurl{PaperXI.ActualELatticeResidue59.restrictScalarMatrix_apply}\\

28&theorem&\nolinkurl{PaperXI.ActualELatticeResidue59.restrictScalarMatrix_injective}\\

37&def&\nolinkurl{PaperXI.ActualELatticeResidue59.integralMatrix}\\

40&theorem&\nolinkurl{PaperXI.ActualELatticeResidue59.integralMatrix_injective}\\

44&def&\nolinkurl{PaperXI.ActualELatticeResidue59.integralLattice}\\

51&theorem&\nolinkurl{PaperXI.ActualELatticeResidue59.mem_integralLattice}\\

55&theorem&\nolinkurl{PaperXI.ActualELatticeResidue59.integralLattice_coe}\\

58&theorem&\nolinkurl{PaperXI.ActualELatticeResidue59.integralMatrix_coe}\\

64&theorem&\nolinkurl{PaperXI.ActualELatticeResidue59.integralMatrix_preserves}\\

74&def&\nolinkurl{PaperXI.ActualELatticeResidue59.nextLattice}\\

79&theorem&\nolinkurl{PaperXI.ActualELatticeResidue59.mem_nextLattice}\\

84&theorem&\nolinkurl{PaperXI.ActualELatticeResidue59.coe_mem_nextLattice_iff}\\

111&theorem&\nolinkurl{PaperXI.ActualELatticeResidue59.reduction_kernel_iff}\\

139&def&\nolinkurl{PaperXI.ActualELatticeResidue59.integralImage}\\

141&def&\nolinkurl{PaperXI.ActualELatticeResidue59.integralImageEquiv}\\

144&def&\nolinkurl{PaperXI.ActualELatticeResidue59.residueAction}\\

147&theorem&\nolinkurl{PaperXI.ActualELatticeResidue59.residueAction_surjective}\\

152&def&\nolinkurl{PaperXI.ActualELatticeResidue59.residueQuotientEquiv}\\

\bottomrule\end{longtable}

\section{\texorpdfstring{\nolinkurl{PaperXI.ActualELatticeScalarPeriod60}}{PaperXI.ActualELatticeScalarPeriod60}}\label{module:PaperXI.ActualELatticeScalarPeriod60}

\textbf{Selected layer:} base108 arithmetic branch.\quad\textbf{Source lines:} 87.\par\textbf{SHA-256:} {\scriptsize\nolinkurl{10bfffe3207aaacbc9291830df80f34c502d7383c79a8a9e1b7b723e7e3021f0}}.

\VerbatimInput[fontsize=\scriptsize,breaklines=true,breakanywhere=true]{excerpts/D_ActualELatticeScalarPeriod60.txt}

\textbf{Direct project imports:} \nolinkurl{PaperXI.ActualELatticeFieldJoin59}, \nolinkurl{PaperXI.ActualEmbeddedUniformizer58}.

\begin{longtable}{@{}p{0.07\textwidth}p{0.12\textwidth}p{0.75\textwidth}@{}}\toprule Line&Kind&Name\\\midrule\endhead

26&theorem&\nolinkurl{PaperXI.ActualELatticeScalarPeriod60.base_uniformizer_degree_ne_one}\\

42&theorem&\nolinkurl{PaperXI.ActualELatticeScalarPeriod60.exists_positive_base_normalForm}\\

67&theorem&\nolinkurl{PaperXI.ActualELatticeScalarPeriod60.chain_period_of_normalForm}\\

81&theorem&\nolinkurl{PaperXI.ActualELatticeScalarPeriod60.exists_chain_period}\\

\bottomrule\end{longtable}

\section{\texorpdfstring{\nolinkurl{PaperXI.ActualELatticeTopology60}}{PaperXI.ActualELatticeTopology60}}\label{module:PaperXI.ActualELatticeTopology60}

\textbf{Selected layer:} base108 arithmetic branch.\quad\textbf{Source lines:} 170.\par\textbf{SHA-256:} {\scriptsize\nolinkurl{4936918f1210fcc205c45cc2b2e787c6dfa51e1d1399c1c9e3988e6d9e791981}}.

\VerbatimInput[fontsize=\scriptsize,breaklines=true,breakanywhere=true]{excerpts/D_ActualELatticeTopology60.txt}

\textbf{Direct project imports:} \nolinkurl{PaperXI.LatticeCongruenceTopology60}, \nolinkurl{PaperXI.ActualELatticeScalarPeriod60}.

\begin{longtable}{@{}p{0.07\textwidth}p{0.12\textwidth}p{0.75\textwidth}@{}}\toprule Line&Kind&Name\\\midrule\endhead

29&theorem&\nolinkurl{PaperXI.ActualELatticeTopology60.field_nhdsNE_neBot}\\

41&def&\nolinkurl{PaperXI.ActualELatticeTopology60.coordinatesOfBasis}\\

56&def&\nolinkurl{PaperXI.ActualELatticeTopology60.scalarMatrixEmbedding}\\

62&theorem&\nolinkurl{PaperXI.ActualELatticeTopology60.integralLattice_compact}\\

67&theorem&\nolinkurl{PaperXI.ActualELatticeTopology60.integralLattice_open}\\

74&theorem&\nolinkurl{PaperXI.ActualELatticeTopology60.latticeChain_compact}\\

86&theorem&\nolinkurl{PaperXI.ActualELatticeTopology60.latticeChain_open}\\

95&def&\nolinkurl{PaperXI.ActualELatticeTopology60.matrixType}\\

99&def&\nolinkurl{PaperXI.ActualELatticeTopology60.matrixCongruence}\\

105&theorem&\nolinkurl{PaperXI.ActualELatticeTopology60.matrixCongruence_open}\\

115&theorem&\nolinkurl{PaperXI.ActualELatticeTopology60.matrixType_open}\\

124&theorem&\nolinkurl{PaperXI.ActualELatticeTopology60.matrixType_compact}\\

139&theorem&\nolinkurl{PaperXI.ActualELatticeTopology60.matrixCongruence_compact}\\

148&theorem&\nolinkurl{PaperXI.ActualELatticeTopology60.compact_open_of_basis}\\

155&def&\nolinkurl{PaperXI.ActualELatticeTopology60.finiteExtensionCoordinates}\\

161&theorem&\nolinkurl{PaperXI.ActualELatticeTopology60.finiteExtension_compact_open}\\

\bottomrule\end{longtable}

\section{\texorpdfstring{\nolinkurl{PaperXI.ActualELatticeTypeGroup59}}{PaperXI.ActualELatticeTypeGroup59}}\label{module:PaperXI.ActualELatticeTypeGroup59}

\textbf{Selected layer:} base108 arithmetic branch.\quad\textbf{Source lines:} 136.\par\textbf{SHA-256:} {\scriptsize\nolinkurl{985e4b5cffb7f0083ef712004a55d0dd95409019eac2f83d57c77cd971ef0455}}.

\VerbatimInput[fontsize=\scriptsize,breaklines=true,breakanywhere=true]{excerpts/D_ActualELatticeTypeGroup59.txt}

\textbf{Direct project imports:} \nolinkurl{PaperXI.ActualELatticeNormalizer59}, \nolinkurl{PaperXI.ActualELatticeResidue59}.

\begin{longtable}{@{}p{0.07\textwidth}p{0.12\textwidth}p{0.75\textwidth}@{}}\toprule Line&Kind&Name\\\midrule\endhead

22&def&\nolinkurl{PaperXI.ActualELatticeTypeGroup59.uniformizerAction}\\

27&theorem&\nolinkurl{PaperXI.ActualELatticeTypeGroup59.uniformizerAction_apply}\\

31&def&\nolinkurl{PaperXI.ActualELatticeTypeGroup59.latticeChain}\\

35&theorem&\nolinkurl{PaperXI.ActualELatticeTypeGroup59.latticeChain_zero}\\

40&theorem&\nolinkurl{PaperXI.ActualELatticeTypeGroup59.latticeChain_one}\\

56&theorem&\nolinkurl{PaperXI.ActualELatticeTypeGroup59.integralMatrix_preserves_iff}\\

65&theorem&\nolinkurl{PaperXI.ActualELatticeTypeGroup59.integralMatrix_commute}\\

74&theorem&\nolinkurl{PaperXI.ActualELatticeTypeGroup59.integralMatrix_mem_stabilizer}\\

83&def&\nolinkurl{PaperXI.ActualELatticeTypeGroup59.congruence}\\

87&theorem&\nolinkurl{PaperXI.ActualELatticeTypeGroup59.integralImage_normalizes}\\

94&theorem&\nolinkurl{PaperXI.ActualELatticeTypeGroup59.integralMatrix_mem_congruence_iff}\\

124&theorem&\nolinkurl{PaperXI.ActualELatticeTypeGroup59.congruence_subgroupOf_integralImage}\\

133&def&\nolinkurl{PaperXI.ActualELatticeTypeGroup59.typeGroup}\\

\bottomrule\end{longtable}

\section{\texorpdfstring{\nolinkurl{PaperXI.ActualEmbeddedUniformizer58}}{PaperXI.ActualEmbeddedUniformizer58}}\label{module:PaperXI.ActualEmbeddedUniformizer58}

\textbf{Selected layer:} base108 arithmetic branch.\quad\textbf{Source lines:} 177.\par\textbf{SHA-256:} {\scriptsize\nolinkurl{3f8d9f340d7a6f79785e6b82d82a800344788509d9b97b51c77aa96429018a33}}.

\VerbatimInput[fontsize=\scriptsize,breaklines=true,breakanywhere=true]{excerpts/D_ActualEmbeddedUniformizer58.txt}

\textbf{Direct project imports:} \nolinkurl{PaperXI.ActualDepthZeroExtension}, \nolinkurl{PaperXI.CompactDegreeSeparation}.

\begin{longtable}{@{}p{0.07\textwidth}p{0.12\textwidth}p{0.75\textwidth}@{}}\toprule Line&Kind&Name\\\midrule\endhead

30&def&\nolinkurl{PaperXI.ActualEmbeddedUniformizer58.unitsEmbedding}\\

34&def&\nolinkurl{PaperXI.ActualEmbeddedUniformizer58.embeddedUniformizer}\\

39&theorem&\nolinkurl{PaperXI.ActualEmbeddedUniformizer58.unitsEmbedding_continuous}\\

42&theorem&\nolinkurl{PaperXI.ActualEmbeddedUniformizer58.determinantDegree_continuous}\\

48&theorem&\nolinkurl{PaperXI.ActualEmbeddedUniformizer58.unitsEmbedding_baseScalar}\\

57&theorem&\nolinkurl{PaperXI.ActualEmbeddedUniformizer58.determinantDegree_integral_unit}\\

71&theorem&\nolinkurl{PaperXI.ActualEmbeddedUniformizer58.algebraEmbedding_continuous}\\

89&theorem&\nolinkurl{PaperXI.ActualEmbeddedUniformizer58.embeddedUniformizer_degree_ne_one}\\

107&theorem&\nolinkurl{PaperXI.ActualEmbeddedUniformizer58.embeddedUniformizer_free}\\

117&theorem&\nolinkurl{PaperXI.ActualEmbeddedUniformizer58.embeddedUniformizer_free_of_algebraMap}\\

126&theorem&\nolinkurl{PaperXI.ActualEmbeddedUniformizer58.normalForm_injective}\\

152&theorem&\nolinkurl{PaperXI.ActualEmbeddedUniformizer58.field_product_normalForm}\\

166&theorem&\nolinkurl{PaperXI.ActualEmbeddedUniformizer58.field_product_unique_normalForm}\\

\bottomrule\end{longtable}

\section{\texorpdfstring{\nolinkurl{PaperXI.ActualFieldTypeQuotient69}}{PaperXI.ActualFieldTypeQuotient69}}\label{module:PaperXI.ActualFieldTypeQuotient69}

\textbf{Selected layer:} base108 arithmetic branch.\quad\textbf{Source lines:} 157.\par\textbf{SHA-256:} {\scriptsize\nolinkurl{fd2a01056fe19cf73010c917683c3bcdac14d241c1d9f3bc3accb3584b709bd6}}.

\VerbatimInput[fontsize=\scriptsize,breaklines=true,breakanywhere=true]{excerpts/D_ActualFieldTypeQuotient69.txt}

\textbf{Direct project imports:} \nolinkurl{PaperXI.ActualNormalizerTypeExtension58}, \nolinkurl{PaperXI.CompactQuotientPeriod69}.

\begin{longtable}{@{}p{0.07\textwidth}p{0.12\textwidth}p{0.75\textwidth}@{}}\toprule Line&Kind&Name\\\midrule\endhead

21&abbrev&\nolinkurl{PaperXI.ActualFieldTypeQuotient69.fieldGroup}\\

25&def&\nolinkurl{PaperXI.ActualFieldTypeQuotient69.scalarHom}\\

31&def&\nolinkurl{PaperXI.ActualFieldTypeQuotient69.scalarSubgroup}\\

36&theorem&\nolinkurl{PaperXI.ActualFieldTypeQuotient69.scalar_central}\\

43&instance&\nolinkurl{PaperXI.ActualFieldTypeQuotient69.scalar_normal}\\

51&def&\nolinkurl{PaperXI.ActualFieldTypeQuotient69.compactPart}\\

55&theorem&\nolinkurl{PaperXI.ActualFieldTypeQuotient69.degree_kernel}\\

98&theorem&\nolinkurl{PaperXI.ActualFieldTypeQuotient69.quotient_compact}\\

\bottomrule\end{longtable}

\section{\texorpdfstring{\nolinkurl{PaperXI.ActualGroupIndex111}}{PaperXI.ActualGroupIndex111}}\label{module:PaperXI.ActualGroupIndex111}

\textbf{Selected layer:} Continuation111.\quad\textbf{Source lines:} 273.\par\textbf{SHA-256:} {\scriptsize\nolinkurl{569cdbec7bc0fcd987fb8fc614db2c4369221abc0007eb98bfdff0896e0800c2}}.

\VerbatimInput[fontsize=\scriptsize,breaklines=true,breakanywhere=true]{excerpts/D_ActualGroupIndex111.txt}

\textbf{Direct project imports:} \nolinkurl{PaperXI.ActualCoordinateDeepIdeal111}, \nolinkurl{PaperXI.ActualOnePlusTopology111}, \nolinkurl{PaperXI.MoritaEndomorphismCoordinates111}, \nolinkurl{PaperXI.MoritaPairIndex110}, \nolinkurl{PaperXI.MoritaRadicalDictionary110}, \nolinkurl{PaperXI.DegreeIndex110}.

\begin{longtable}{@{}p{0.07\textwidth}p{0.12\textwidth}p{0.75\textwidth}@{}}\toprule Line&Kind&Name\\\midrule\endhead

24&theorem&\nolinkurl{PaperXI.ActualGroupIndex111.units_innerRing_inv}\\

33&theorem&\nolinkurl{PaperXI.ActualGroupIndex111.map_inner_inv_eq_comap}\\

47&theorem&\nolinkurl{PaperXI.ActualGroupIndex111.integralUnitMap_eq_actualEmbedding}\\

53&theorem&\nolinkurl{PaperXI.ActualGroupIndex111.onePlus_backward_relIndex}\\

93&theorem&\nolinkurl{PaperXI.ActualGroupIndex111.onePlus_index_of_Morita_image}\\

125&theorem&\nolinkurl{PaperXI.ActualGroupIndex111.ambient_pair_index}\\

163&theorem&\nolinkurl{PaperXI.ActualGroupIndex111.ambient_pair_index_of_radical_products}\\

202&structure&\nolinkurl{PaperXI.ActualGroupIndex111.ActualLatticeData}\\

244&theorem&\nolinkurl{PaperXI.ActualGroupIndex111.ActualLatticeData.actual_index}\\

\bottomrule\end{longtable}

\section{\texorpdfstring{\nolinkurl{PaperXI.ActualHeisenbergCandidate107}}{PaperXI.ActualHeisenbergCandidate107}}\label{module:PaperXI.ActualHeisenbergCandidate107}

\textbf{Selected layer:} base108 arithmetic branch.\quad\textbf{Source lines:} 89.\par\textbf{SHA-256:} {\scriptsize\nolinkurl{7dd358423f2d83b2201f0e8b9eed0484c2e577f041d171fe447c59806e11d971}}.

\VerbatimInput[fontsize=\scriptsize,breaklines=true,breakanywhere=true]{excerpts/D_ActualHeisenbergCandidate107.txt}

\textbf{Direct project imports:} \nolinkurl{PaperXI.ActualCartanCandidate106}, \nolinkurl{PaperXI.ResidueRestriction107}, \nolinkurl{PaperXI.HeisenbergReduction107}, \nolinkurl{PaperXI.IntertwiningSpace107}.

\begin{longtable}{@{}p{0.07\textwidth}p{0.12\textwidth}p{0.75\textwidth}@{}}\toprule Line&Kind&Name\\\midrule\endhead

34&theorem&\nolinkurl{PaperXI.ActualHeisenbergCandidate107.candidate_exists}\\

\bottomrule\end{longtable}

\section{\texorpdfstring{\nolinkurl{PaperXI.ActualIntegerTower111}}{PaperXI.ActualIntegerTower111}}\label{module:PaperXI.ActualIntegerTower111}

\textbf{Selected layer:} Continuation111.\quad\textbf{Source lines:} 178.\par\textbf{SHA-256:} {\scriptsize\nolinkurl{6effc42145758945edec994d27d6de01bfae48a6752ba7b8f53ac9f8c9ee0300}}.

\VerbatimInput[fontsize=\scriptsize,breaklines=true,breakanywhere=true]{excerpts/D_ActualIntegerTower111.txt}

\textbf{Direct project imports:} \nolinkurl{PaperXI.MoritaActualCartan110}.

\begin{longtable}{@{}p{0.07\textwidth}p{0.12\textwidth}p{0.75\textwidth}@{}}\toprule Line&Kind&Name\\\midrule\endhead

20&theorem&\nolinkurl{PaperXI.ActualIntegerTower111.map_integer_of_continuous}\\

43&def&\nolinkurl{PaperXI.ActualIntegerTower111.integerMap}\\

47&theorem&\nolinkurl{PaperXI.ActualIntegerTower111.integerMap_coe}\\

51&theorem&\nolinkurl{PaperXI.ActualIntegerTower111.integer_of_map_integer}\\

79&theorem&\nolinkurl{PaperXI.ActualIntegerTower111.integer_comap_eq}\\

88&theorem&\nolinkurl{PaperXI.ActualIntegerTower111.valuationExtension}\\

93&def&\nolinkurl{PaperXI.ActualIntegerTower111.integerAlgebra}\\

97&theorem&\nolinkurl{PaperXI.ActualIntegerTower111.integerAlgebra_map_coe}\\

102&theorem&\nolinkurl{PaperXI.ActualIntegerTower111.integer_tower}\\

109&theorem&\nolinkurl{PaperXI.ActualIntegerTower111.integerAlgebra_map_eq}\\

117&theorem&\nolinkurl{PaperXI.ActualIntegerTower111.map_valuation_le_iff}\\

126&theorem&\nolinkurl{PaperXI.ActualIntegerTower111.integer_end_smul_apply}\\

131&theorem&\nolinkurl{PaperXI.ActualIntegerTower111.integer_end_tower}\\

143&theorem&\nolinkurl{PaperXI.ActualIntegerTower111.integer_matrix_smul_apply}\\

149&theorem&\nolinkurl{PaperXI.ActualIntegerTower111.integer_matrix_tower}\\

162&theorem&\nolinkurl{PaperXI.ActualIntegerTower111.exists_actual_cartan_index_of_integer_lattice}\\

\bottomrule\end{longtable}

\section{\texorpdfstring{\nolinkurl{PaperXI.ActualIntegralGeneralLinear}}{PaperXI.ActualIntegralGeneralLinear}}\label{module:PaperXI.ActualIntegralGeneralLinear}

\textbf{Selected layer:} base108 arithmetic branch.\quad\textbf{Source lines:} 219.\par\textbf{SHA-256:} {\scriptsize\nolinkurl{37fadf54dfd99b1032584b0bd3b561824077296be369ecc43f74deffb3847996}}.

\VerbatimInput[fontsize=\scriptsize,breaklines=true,breakanywhere=true]{excerpts/D_ActualIntegralGeneralLinear.txt}

\textbf{Direct project imports:} \nolinkurl{PaperXI.ActualResidueGeneralLinear}, \nolinkurl{PaperXI.QuotientRepresentation}, \nolinkurl{PaperXI.SmoothDegreeQuotient}, \nolinkurl{PaperXI.FiniteRepresentation}.

\begin{longtable}{@{}p{0.07\textwidth}p{0.12\textwidth}p{0.75\textwidth}@{}}\toprule Line&Kind&Name\\\midrule\endhead

27&def&\nolinkurl{PaperXI.ActualIntegralGeneralLinear.embedding}\\

30&theorem&\nolinkurl{PaperXI.ActualIntegralGeneralLinear.embedding_apply}\\

33&theorem&\nolinkurl{PaperXI.ActualIntegralGeneralLinear.embedding_injective}\\

41&def&\nolinkurl{PaperXI.ActualIntegralGeneralLinear.integralSubgroup}\\

44&def&\nolinkurl{PaperXI.ActualIntegralGeneralLinear.integralEquiv}\\

47&theorem&\nolinkurl{PaperXI.ActualIntegralGeneralLinear.integralEquiv_coe}\\

50&theorem&\nolinkurl{PaperXI.ActualIntegralGeneralLinear.embedding_integralEquiv_symm}\\

55&theorem&\nolinkurl{PaperXI.ActualIntegralGeneralLinear.mem_integralSubgroup}\\

84&theorem&\nolinkurl{PaperXI.ActualIntegralGeneralLinear.field_t2Space}\\

92&theorem&\nolinkurl{PaperXI.ActualIntegralGeneralLinear.embedding_isEmbedding}\\

96&theorem&\nolinkurl{PaperXI.ActualIntegralGeneralLinear.embedding_continuous}\\

100&theorem&\nolinkurl{PaperXI.ActualIntegralGeneralLinear.integralEquiv_symm_continuous}\\

111&theorem&\nolinkurl{PaperXI.ActualIntegralGeneralLinear.integralSubgroup_compact}\\

120&theorem&\nolinkurl{PaperXI.ActualIntegralGeneralLinear.integralSubgroup_open}\\

131&def&\nolinkurl{PaperXI.ActualIntegralGeneralLinear.reductionOnSubgroup}\\

134&theorem&\nolinkurl{PaperXI.ActualIntegralGeneralLinear.reductionOnSubgroup_integralEquiv}\\

138&theorem&\nolinkurl{PaperXI.ActualIntegralGeneralLinear.reductionOnSubgroup_surjective}\\

142&theorem&\nolinkurl{PaperXI.ActualIntegralGeneralLinear.mem_reduction_kernel}\\

156&theorem&\nolinkurl{PaperXI.ActualIntegralGeneralLinear.reduction_kernel_open}\\

171&theorem&\nolinkurl{PaperXI.ActualIntegralGeneralLinear.reductionOnSubgroup_kernel_open}\\

177&def&\nolinkurl{PaperXI.ActualIntegralGeneralLinear.residueQuotientEquiv}\\

182&theorem&\nolinkurl{PaperXI.ActualIntegralGeneralLinear.finite_residue_quotient}\\

191&def&\nolinkurl{PaperXI.ActualIntegralGeneralLinear.inflated}\\

194&theorem&\nolinkurl{PaperXI.ActualIntegralGeneralLinear.inflated_integralEquiv}\\

199&theorem&\nolinkurl{PaperXI.ActualIntegralGeneralLinear.inflated_irreducible}\\

205&theorem&\nolinkurl{PaperXI.ActualIntegralGeneralLinear.inflated_isSmooth}\\

216&theorem&\nolinkurl{PaperXI.ActualIntegralGeneralLinear.inflated_finiteDimensional}\\

\bottomrule\end{longtable}

\section{\texorpdfstring{\nolinkurl{PaperXI.ActualIntersection108}}{PaperXI.ActualIntersection108}}\label{module:PaperXI.ActualIntersection108}

\textbf{Selected layer:} base108 arithmetic branch.\quad\textbf{Source lines:} 59.\par\textbf{SHA-256:} {\scriptsize\nolinkurl{fede966a6ccc302109f79282ebd6c8a1b66b3c9dcbafc6b1cccd28b1160bd34b}}.

\VerbatimInput[fontsize=\scriptsize,breaklines=true,breakanywhere=true]{excerpts/D_ActualIntersection108.txt}

\textbf{Direct project imports:} \nolinkurl{PaperXI.ActualCharacterCandidate108}, \nolinkurl{PaperXI.FiniteIntersectionCriterion108}.

\begin{longtable}{@{}p{0.07\textwidth}p{0.12\textwidth}p{0.75\textwidth}@{}}\toprule Line&Kind&Name\\\midrule\endhead

20&abbrev&\nolinkurl{PaperXI.ActualIntersection108.embedding}\\

23&abbrev&\nolinkurl{PaperXI.ActualIntersection108.restriction}\\

25&abbrev&\nolinkurl{PaperXI.ActualIntersection108.quotient}\\

30&theorem&\nolinkurl{PaperXI.ActualIntersection108.finite}\\

47&theorem&\nolinkurl{PaperXI.ActualIntersection108.dimension}\\

\bottomrule\end{longtable}

\section{\texorpdfstring{\nolinkurl{PaperXI.ActualLocalCartan106}}{PaperXI.ActualLocalCartan106}}\label{module:PaperXI.ActualLocalCartan106}

\textbf{Selected layer:} base108 arithmetic branch.\quad\textbf{Source lines:} 44.\par\textbf{SHA-256:} {\scriptsize\nolinkurl{ad8205ed61e0d305b7023ba61e8f17b266407a62ec64e846e7993a90d62882b2}}.

\VerbatimInput[fontsize=\scriptsize,breaklines=true,breakanywhere=true]{excerpts/D_ActualLocalCartan106.txt}

\textbf{Direct project imports:} \nolinkurl{PaperXI.DVRCartan106}, \nolinkurl{PaperXI.ActualLocalUnitDecomposition}, \nolinkurl{PaperXI.ActualIntegralGeneralLinear}, \nolinkurl{PaperXI.ActualResidueGeneralLinear}.

\begin{longtable}{@{}p{0.07\textwidth}p{0.12\textwidth}p{0.75\textwidth}@{}}\toprule Line&Kind&Name\\\midrule\endhead

16&theorem&\nolinkurl{PaperXI.ActualLocalCartan106.uniformizer_residue_zero}\\

22&theorem&\nolinkurl{PaperXI.ActualLocalCartan106.hasCartan}\\

31&theorem&\nolinkurl{PaperXI.ActualLocalCartan106.intertwining_support}\\

\bottomrule\end{longtable}

\section{\texorpdfstring{\nolinkurl{PaperXI.ActualLocalUnitDecomposition}}{PaperXI.ActualLocalUnitDecomposition}}\label{module:PaperXI.ActualLocalUnitDecomposition}

\textbf{Selected layer:} base108 arithmetic branch.\quad\textbf{Source lines:} 128.\par\textbf{SHA-256:} {\scriptsize\nolinkurl{a8395974d47aefc05e157314ff5c0f41b5ec8a33c9bfb98579de8b62ab5e273d}}.

\VerbatimInput[fontsize=\scriptsize,breaklines=true,breakanywhere=true]{excerpts/D_ActualLocalUnitDecomposition.txt}

\textbf{Direct project imports:} \nolinkurl{PaperXI.ActualLocalUnitTopology}.

\begin{longtable}{@{}p{0.07\textwidth}p{0.12\textwidth}p{0.75\textwidth}@{}}\toprule Line&Kind&Name\\\midrule\endhead

22&def&\nolinkurl{PaperXI.ActualLocalUnitDecomposition.uniformizer}\\

24&theorem&\nolinkurl{PaperXI.ActualLocalUnitDecomposition.uniformizer_irreducible}\\

27&theorem&\nolinkurl{PaperXI.ActualLocalUnitDecomposition.uniformizer_ne_zero}\\

31&def&\nolinkurl{PaperXI.ActualLocalUnitDecomposition.uniformizerUnit}\\

33&theorem&\nolinkurl{PaperXI.ActualLocalUnitDecomposition.uniformizerUnit_coe}\\

35&theorem&\nolinkurl{PaperXI.ActualLocalUnitDecomposition.uniformizer_value_pos}\\

38&theorem&\nolinkurl{PaperXI.ActualLocalUnitDecomposition.uniformizer_value_lt_one}\\

42&def&\nolinkurl{PaperXI.ActualLocalUnitDecomposition.normalFormHom}\\

50&theorem&\nolinkurl{PaperXI.ActualLocalUnitDecomposition.normalFormHom_apply}\\

54&theorem&\nolinkurl{PaperXI.ActualLocalUnitDecomposition.normalForm_surjective}\\

65&theorem&\nolinkurl{PaperXI.ActualLocalUnitDecomposition.normalForm_injective}\\

82&def&\nolinkurl{PaperXI.ActualLocalUnitDecomposition.normalFormEquiv}\\

85&theorem&\nolinkurl{PaperXI.ActualLocalUnitDecomposition.normalFormEquiv_apply}\\

89&def&\nolinkurl{PaperXI.ActualLocalUnitDecomposition.degree}\\

92&theorem&\nolinkurl{PaperXI.ActualLocalUnitDecomposition.degree_normalForm}\\

97&theorem&\nolinkurl{PaperXI.ActualLocalUnitDecomposition.degree_unit}\\

100&theorem&\nolinkurl{PaperXI.ActualLocalUnitDecomposition.degree_unitEmbedding}\\

103&theorem&\nolinkurl{PaperXI.ActualLocalUnitDecomposition.degree_uniformizer}\\

107&theorem&\nolinkurl{PaperXI.ActualLocalUnitDecomposition.ker_degree}\\

119&theorem&\nolinkurl{PaperXI.ActualLocalUnitDecomposition.degree_continuous}\\

\bottomrule\end{longtable}

\section{\texorpdfstring{\nolinkurl{PaperXI.ActualLocalUnitTopology}}{PaperXI.ActualLocalUnitTopology}}\label{module:PaperXI.ActualLocalUnitTopology}

\textbf{Selected layer:} base108 arithmetic branch.\quad\textbf{Source lines:} 102.\par\textbf{SHA-256:} {\scriptsize\nolinkurl{2d87df54c86bef03b5bf3277d309f08733cde2c39d287c2226b9d04bcbe388ea}}.

\VerbatimInput[fontsize=\scriptsize,breaklines=true,breakanywhere=true]{excerpts/D_ActualLocalUnitTopology.txt}

\textbf{Direct project imports:} none.

\begin{longtable}{@{}p{0.07\textwidth}p{0.12\textwidth}p{0.75\textwidth}@{}}\toprule Line&Kind&Name\\\midrule\endhead

23&def&\nolinkurl{PaperXI.ActualLocalUnitTopology.unitEmbedding}\\

27&def&\nolinkurl{PaperXI.ActualLocalUnitTopology.compactUnitSubgroup}\\

29&theorem&\nolinkurl{PaperXI.ActualLocalUnitTopology.coe_unitEmbedding}\\

32&theorem&\nolinkurl{PaperXI.ActualLocalUnitTopology.unitEmbedding_injective}\\

35&theorem&\nolinkurl{PaperXI.ActualLocalUnitTopology.mem_compactUnitSubgroup_iff}\\

50&theorem&\nolinkurl{PaperXI.ActualLocalUnitTopology.compactUnitSubgroup_eq_units}\\

62&theorem&\nolinkurl{PaperXI.ActualLocalUnitTopology.continuous_unitEmbedding}\\

67&instance&\nolinkurl{PaperXI.ActualLocalUnitTopology.localFieldT2Space}\\

73&theorem&\nolinkurl{PaperXI.ActualLocalUnitTopology.isCompact_compactUnitSubgroup}\\

78&theorem&\nolinkurl{PaperXI.ActualLocalUnitTopology.isOpen_compactUnitSubgroup}\\

84&def&\nolinkurl{PaperXI.ActualLocalUnitTopology.unitEmbeddingEquiv}\\

92&theorem&\nolinkurl{PaperXI.ActualLocalUnitTopology.coe_unitEmbeddingEquiv}\\

96&theorem&\nolinkurl{PaperXI.ActualLocalUnitTopology.unitEmbedding_isOpenMap}\\

\bottomrule\end{longtable}

\section{\texorpdfstring{\nolinkurl{PaperXI.ActualNormalizerTypeExtension58}}{PaperXI.ActualNormalizerTypeExtension58}}\label{module:PaperXI.ActualNormalizerTypeExtension58}

\textbf{Selected layer:} base108 arithmetic branch.\quad\textbf{Source lines:} 82.\par\textbf{SHA-256:} {\scriptsize\nolinkurl{eed55102f4bdce2532e92a4655f4da687a9fa494d89917d61f2a8cef613b46f1}}.

\VerbatimInput[fontsize=\scriptsize,breaklines=true,breakanywhere=true]{excerpts/D_ActualNormalizerTypeExtension58.txt}

\textbf{Direct project imports:} \nolinkurl{PaperXI.ActualEmbeddedUniformizer58}, \nolinkurl{PaperXI.NormalizerExtensionInducing58}.

\begin{longtable}{@{}p{0.07\textwidth}p{0.12\textwidth}p{0.75\textwidth}@{}}\toprule Line&Kind&Name\\\midrule\endhead

32&theorem&\nolinkurl{PaperXI.ActualNormalizerTypeExtension58.extendedSubgroup_eq_field_join}\\

60&theorem&\nolinkurl{PaperXI.ActualNormalizerTypeExtension58.exists_smooth_irreducible_extension}\\

\bottomrule\end{longtable}

\section{\texorpdfstring{\nolinkurl{PaperXI.ActualOnePlusTopology111}}{PaperXI.ActualOnePlusTopology111}}\label{module:PaperXI.ActualOnePlusTopology111}

\textbf{Selected layer:} Continuation111.\quad\textbf{Source lines:} 114.\par\textbf{SHA-256:} {\scriptsize\nolinkurl{cdef6d0ce5c93ad213cc8af1fef105ca1375f9566fdf088cdaad7677198cb9ac}}.

\VerbatimInput[fontsize=\scriptsize,breaklines=true,breakanywhere=true]{excerpts/D_ActualOnePlusTopology111.txt}

\textbf{Direct project imports:} \nolinkurl{PaperXI.ActualCoordinateDeepIdeal111}, \nolinkurl{PaperXI.ActualIntegralGeneralLinear}.

\begin{longtable}{@{}p{0.07\textwidth}p{0.12\textwidth}p{0.75\textwidth}@{}}\toprule Line&Kind&Name\\\midrule\endhead

19&def&\nolinkurl{PaperXI.ActualOnePlusTopology111.matrixUnitsHomeomorph}\\

38&theorem&\nolinkurl{PaperXI.ActualOnePlusTopology111.matrix_isOpenUnits}\\

53&theorem&\nolinkurl{PaperXI.ActualOnePlusTopology111.actual_matrix_isOpenUnits}\\

59&theorem&\nolinkurl{PaperXI.ActualOnePlusTopology111.additive_open_of_image_group_open}\\

100&theorem&\nolinkurl{PaperXI.ActualOnePlusTopology111.onePlus_conjugate_relIndex_of_group_open}\\

\bottomrule\end{longtable}

\section{\texorpdfstring{\nolinkurl{PaperXI.ActualPacketTypeRealization70}}{PaperXI.ActualPacketTypeRealization70}}\label{module:PaperXI.ActualPacketTypeRealization70}

\textbf{Selected layer:} base108 arithmetic branch.\quad\textbf{Source lines:} 109.\par\textbf{SHA-256:} {\scriptsize\nolinkurl{079fe5114db1fe95bd51e370b7932a6394bce0b44bad32671fa820774e806037}}.

\VerbatimInput[fontsize=\scriptsize,breaklines=true,breakanywhere=true]{excerpts/D_ActualPacketTypeRealization70.txt}

\textbf{Direct project imports:} \nolinkurl{PaperXI.PacketTypeAttachment70}, \nolinkurl{PaperXI.PrescribedFieldType70}, \nolinkurl{PaperXI.ArithmeticTypeCertificate69}, \nolinkurl{PaperXI.ActualCharacterIdentification61}, \nolinkurl{PaperXI.ActualCharacterResidue61}.

\begin{longtable}{@{}p{0.07\textwidth}p{0.12\textwidth}p{0.75\textwidth}@{}}\toprule Line&Kind&Name\\\midrule\endhead

36&theorem&\nolinkurl{PaperXI.ActualPacketTypeRealization70.candidate_exists}\\

\bottomrule\end{longtable}

\section{\texorpdfstring{\nolinkurl{PaperXI.ActualProfileCandidate110}}{PaperXI.ActualProfileCandidate110}}\label{module:PaperXI.ActualProfileCandidate110}

\textbf{Selected layer:} Continuation110.\quad\textbf{Source lines:} 94.\par\textbf{SHA-256:} {\scriptsize\nolinkurl{dd434e182eec32397e64a3e6a43193cc56a3c965eff65c4e36e14f4bd0b2c27b}}.

\VerbatimInput[fontsize=\scriptsize,breaklines=true,breakanywhere=true]{excerpts/D_ActualProfileCandidate110.txt}

\textbf{Direct project imports:} \nolinkurl{PaperXI.DegreeIndex110}.

\begin{longtable}{@{}p{0.07\textwidth}p{0.12\textwidth}p{0.75\textwidth}@{}}\toprule Line&Kind&Name\\\midrule\endhead

34&theorem&\nolinkurl{PaperXI.ActualProfileCandidate110.candidate_exists_from_original_profiles}\\

\bottomrule\end{longtable}

\section{\texorpdfstring{\nolinkurl{PaperXI.ActualResidueGeneralLinear}}{PaperXI.ActualResidueGeneralLinear}}\label{module:PaperXI.ActualResidueGeneralLinear}

\textbf{Selected layer:} base108 arithmetic branch.\quad\textbf{Source lines:} 69.\par\textbf{SHA-256:} {\scriptsize\nolinkurl{33f125bdc926fb2a521cd18f1d46ba7c0d70d277e4d96bf0827d0fc610cd3459}}.

\VerbatimInput[fontsize=\scriptsize,breaklines=true,breakanywhere=true]{excerpts/D_ActualResidueGeneralLinear.txt}

\textbf{Direct project imports:} none.

\begin{longtable}{@{}p{0.07\textwidth}p{0.12\textwidth}p{0.75\textwidth}@{}}\toprule Line&Kind&Name\\\midrule\endhead

21&def&\nolinkurl{PaperXI.ActualResidueGeneralLinear.residueMap}\\

25&theorem&\nolinkurl{PaperXI.ActualResidueGeneralLinear.residueMap_apply}\\

30&theorem&\nolinkurl{PaperXI.ActualResidueGeneralLinear.residueMap_surjective}\\

52&def&\nolinkurl{PaperXI.ActualResidueGeneralLinear.reduction}\\

55&theorem&\nolinkurl{PaperXI.ActualResidueGeneralLinear.reduction_apply}\\

59&theorem&\nolinkurl{PaperXI.ActualResidueGeneralLinear.reduction_surjective}\\

65&theorem&\nolinkurl{PaperXI.ActualResidueGeneralLinear.finite_reduction_target}\\

\bottomrule\end{longtable}

\section{\texorpdfstring{\nolinkurl{PaperXI.ActualScalarLift113}}{PaperXI.ActualScalarLift113}}\label{module:PaperXI.ActualScalarLift113}

\textbf{Selected layer:} Continuation113.\quad\textbf{Source lines:} 221.\par\textbf{SHA-256:} {\scriptsize\nolinkurl{9dca2abac60d683f19c4b2bd2bd257b8f18e5b5f459cf39df9948de4ccf41ce0}}.

\VerbatimInput[fontsize=\scriptsize,breaklines=true,breakanywhere=true]{excerpts/D_ActualScalarLift113.txt}

\textbf{Direct project imports:} \nolinkurl{PaperXI.ScalarRepresentation112}, \nolinkurl{PaperXI.CyclicDefect113}.

\begin{longtable}{@{}p{0.07\textwidth}p{0.12\textwidth}p{0.75\textwidth}@{}}\toprule Line&Kind&Name\\\midrule\endhead

19&def&\nolinkurl{PaperXI.ActualScalarLift113.actualPushout}\\

30&theorem&\nolinkurl{PaperXI.ActualScalarLift113.actualPushout_section}\\

37&theorem&\nolinkurl{PaperXI.ActualScalarLift113.actualPushout_subgroup}\\

50&theorem&\nolinkurl{PaperXI.ActualScalarLift113.quotient_action_power}\\

62&theorem&\nolinkurl{PaperXI.ActualScalarLift113.quotient_power_eq_one}\\

72&def&\nolinkurl{PaperXI.ActualScalarLift113.subgroupMap}\\

88&def&\nolinkurl{PaperXI.ActualScalarLift113.actualPhase}\\

92&theorem&\nolinkurl{PaperXI.ActualScalarLift113.actualPhase_one}\\

98&theorem&\nolinkurl{PaperXI.ActualScalarLift113.sectionScalar_action}\\

127&theorem&\nolinkurl{PaperXI.ActualScalarLift113.actualPhase_mul}\\

137&def&\nolinkurl{PaperXI.ActualScalarLift113.actualLift}\\

161&theorem&\nolinkurl{PaperXI.ActualScalarLift113.actualLift_phase}\\

167&theorem&\nolinkurl{PaperXI.ActualScalarLift113.actualLift_intertwines}\\

175&theorem&\nolinkurl{PaperXI.ActualScalarLift113.actualLift_power_intertwines}\\

191&theorem&\nolinkurl{PaperXI.ActualScalarLift113.actualLift_power_eq_one}\\

\bottomrule\end{longtable}

\section{\texorpdfstring{\nolinkurl{PaperXI.ActualSourceCandidate111}}{PaperXI.ActualSourceCandidate111}}\label{module:PaperXI.ActualSourceCandidate111}

\textbf{Selected layer:} Continuation111.\quad\textbf{Source lines:} 102.\par\textbf{SHA-256:} {\scriptsize\nolinkurl{c802c3b3d1be72ca3d961a4edaa42cca066eb81cbc902a0a0abe9eee58d60005}}.

\VerbatimInput[fontsize=\scriptsize,breaklines=true,breakanywhere=true]{excerpts/D_ActualSourceCandidate111.txt}

\textbf{Direct project imports:} \nolinkurl{PaperXI.ActualProfileCandidate110}, \nolinkurl{PaperXI.HeisenbergCharacterProfile111}, \nolinkurl{PaperXI.ActualGroupIndex111}.

\begin{longtable}{@{}p{0.07\textwidth}p{0.12\textwidth}p{0.75\textwidth}@{}}\toprule Line&Kind&Name\\\midrule\endhead

40&theorem&\nolinkurl{PaperXI.ActualSourceCandidate111.candidate_exists_from_source_dictionary}\\

\bottomrule\end{longtable}

\section{\texorpdfstring{\nolinkurl{PaperXI.ActualUnramifiedBaseChange}}{PaperXI.ActualUnramifiedBaseChange}}\label{module:PaperXI.ActualUnramifiedBaseChange}

\textbf{Selected layer:} v0.44 parent.\quad\textbf{Source lines:} 115.\par\textbf{SHA-256:} {\scriptsize\nolinkurl{8f4e3c5117f404df075772a6e9b3782e2463b7e84973d380f248d9c9c5c4f87f}}.

\VerbatimInput[fontsize=\scriptsize,breaklines=true,breakanywhere=true]{excerpts/D_ActualUnramifiedBaseChange.txt}

\textbf{Direct project imports:} \nolinkurl{PaperXI.UnramifiedFieldTower}.

\begin{longtable}{@{}p{0.07\textwidth}p{0.12\textwidth}p{0.75\textwidth}@{}}\toprule Line&Kind&Name\\\midrule\endhead

43&theorem&\nolinkurl{PaperXI.ActualUnramifiedBaseChange.frobeniusPolynomial_eq_modelBaseChange}\\

52&theorem&\nolinkurl{PaperXI.ActualUnramifiedBaseChange.localFactor_eq_modelBaseChange}\\

62&theorem&\nolinkurl{PaperXI.ActualUnramifiedBaseChange.restricted_inertiaCodimension_eq_zero}\\

71&theorem&\nolinkurl{PaperXI.ActualUnramifiedBaseChange.restricted_swanConductorExponent_eq_zero}\\

79&theorem&\nolinkurl{PaperXI.ActualUnramifiedBaseChange.restricted_artinConductorExponent_eq_zero}\\

87&theorem&\nolinkurl{PaperXI.ActualUnramifiedBaseChange.artinConductorExponent_restriction}\\

98&theorem&\nolinkurl{PaperXI.ActualUnramifiedBaseChange.irreducible_restriction_frobeniusPolynomial}\\

107&theorem&\nolinkurl{PaperXI.ActualUnramifiedBaseChange.irreducible_restriction_localFactor}\\

\bottomrule\end{longtable}

\section{\texorpdfstring{\nolinkurl{PaperXI.AlgebraActionModule}}{PaperXI.AlgebraActionModule}}\label{module:PaperXI.AlgebraActionModule}

\textbf{Selected layer:} v0.44 parent.\quad\textbf{Source lines:} 250.\par\textbf{SHA-256:} {\scriptsize\nolinkurl{5ceb841756c4ca3a00e4fe21b5969470175a83df02d2c417316405eced836820}}.

\VerbatimInput[fontsize=\scriptsize,breaklines=true,breakanywhere=true]{excerpts/D_AlgebraActionModule.txt}

\textbf{Direct project imports:} none.

\begin{longtable}{@{}p{0.07\textwidth}p{0.12\textwidth}p{0.75\textwidth}@{}}\toprule Line&Kind&Name\\\midrule\endhead

19&def&\nolinkurl{PaperXI.AlgebraActionModule.Space}\\

25&def&\nolinkurl{PaperXI.AlgebraActionModule.equiv}\\

27&instance&\nolinkurl{PaperXI.AlgebraActionModule.instModule}\\

29&theorem&\nolinkurl{PaperXI.AlgebraActionModule.equiv_smul}\\

32&theorem&\nolinkurl{PaperXI.AlgebraActionModule.equiv_symm_apply}\\

34&instance&\nolinkurl{PaperXI.AlgebraActionModule.instScalarTower}\\

41&def&\nolinkurl{PaperXI.AlgebraActionModule.Stable}\\

45&def&\nolinkurl{PaperXI.AlgebraActionModule.toSubspace}\\

47&theorem&\nolinkurl{PaperXI.AlgebraActionModule.mem_toSubspace}\\

50&theorem&\nolinkurl{PaperXI.AlgebraActionModule.stable_toSubspace}\\

55&def&\nolinkurl{PaperXI.AlgebraActionModule.ofSubspace}\\

61&theorem&\nolinkurl{PaperXI.AlgebraActionModule.mem_ofSubspace}\\

64&theorem&\nolinkurl{PaperXI.AlgebraActionModule.toSubspace_ofSubspace}\\

67&theorem&\nolinkurl{PaperXI.AlgebraActionModule.ofSubspace_toSubspace}\\

71&def&\nolinkurl{PaperXI.AlgebraActionModule.submoduleOrderIso}\\

78&theorem&\nolinkurl{PaperXI.AlgebraActionModule.toSubspace_bot}\\

79&theorem&\nolinkurl{PaperXI.AlgebraActionModule.toSubspace_top}\\

82&theorem&\nolinkurl{PaperXI.AlgebraActionModule.simple_iff}\\

119&def&\nolinkurl{PaperXI.AlgebraActionModule.Intertwines}\\

123&def&\nolinkurl{PaperXI.AlgebraActionModule.toModuleMap}\\

130&def&\nolinkurl{PaperXI.AlgebraActionModule.toComplexMap}\\

133&theorem&\nolinkurl{PaperXI.AlgebraActionModule.toModuleMap_apply}\\

136&theorem&\nolinkurl{PaperXI.AlgebraActionModule.toComplexMap_apply}\\

139&theorem&\nolinkurl{PaperXI.AlgebraActionModule.intertwines_toComplexMap}\\

143&def&\nolinkurl{PaperXI.AlgebraActionModule.intertwinerEquiv}\\

150&theorem&\nolinkurl{PaperXI.AlgebraActionModule.toComplexMap_toModuleMap}\\

153&theorem&\nolinkurl{PaperXI.AlgebraActionModule.toModuleMap_toComplexMap}\\

159&theorem&\nolinkurl{PaperXI.AlgebraActionModule.toComplexMap_id}\\

165&theorem&\nolinkurl{PaperXI.AlgebraActionModule.toComplexMap_comp}\\

173&def&\nolinkurl{PaperXI.AlgebraActionModule.toModuleEquiv}\\

183&def&\nolinkurl{PaperXI.AlgebraActionModule.toComplexEquiv}\\

186&theorem&\nolinkurl{PaperXI.AlgebraActionModule.toModuleEquiv_apply}\\

190&theorem&\nolinkurl{PaperXI.AlgebraActionModule.toComplexEquiv_apply}\\

193&theorem&\nolinkurl{PaperXI.AlgebraActionModule.intertwines_toComplexEquiv}\\

197&def&\nolinkurl{PaperXI.AlgebraActionModule.equivalenceEquiv}\\

204&theorem&\nolinkurl{PaperXI.AlgebraActionModule.toComplexEquiv_toModuleEquiv}\\

208&theorem&\nolinkurl{PaperXI.AlgebraActionModule.toModuleEquiv_toComplexEquiv}\\

213&theorem&\nolinkurl{PaperXI.AlgebraActionModule.toComplexEquiv_refl}\\

218&theorem&\nolinkurl{PaperXI.AlgebraActionModule.toComplexEquiv_trans}\\

232&def&\nolinkurl{PaperXI.AlgebraActionModule.ofModuleEquiv}\\

240&theorem&\nolinkurl{PaperXI.AlgebraActionModule.ofModuleEquiv_apply}\\

244&theorem&\nolinkurl{PaperXI.AlgebraActionModule.ofModule_simple_iff}\\

\bottomrule\end{longtable}

\section{\texorpdfstring{\nolinkurl{PaperXI.ArithmeticAttachment110}}{PaperXI.ArithmeticAttachment110}}\label{module:PaperXI.ArithmeticAttachment110}

\textbf{Selected layer:} Continuation110.\quad\textbf{Source lines:} 161.\par\textbf{SHA-256:} {\scriptsize\nolinkurl{e67a6b640cb79d7ea347d00621361e950ea439519e53e1f472ef6758460a8693}}.

\VerbatimInput[fontsize=\scriptsize,breaklines=true,breakanywhere=true]{excerpts/D_ArithmeticAttachment110.txt}

\textbf{Direct project imports:} \nolinkurl{PaperXI.ActualCharacterCandidate108}.

\begin{longtable}{@{}p{0.07\textwidth}p{0.12\textwidth}p{0.75\textwidth}@{}}\toprule Line&Kind&Name\\\midrule\endhead

18&def&\nolinkurl{PaperXI.ArithmeticAttachment110.residueOne}\\

21&theorem&\nolinkurl{PaperXI.ArithmeticAttachment110.residueOne_irreducible}\\

27&theorem&\nolinkurl{PaperXI.ArithmeticAttachment110.residueOne_cuspidal}\\

37&theorem&\nolinkurl{PaperXI.ArithmeticAttachment110.gl_one_scalar}\\

45&def&\nolinkurl{PaperXI.ArithmeticAttachment110.tensorUnitEquiv}\\

66&theorem&\nolinkurl{PaperXI.ArithmeticAttachment110.integral_normalization_of_full_rank_one}\\

88&theorem&\nolinkurl{PaperXI.ArithmeticAttachment110.rank_one_candidate_exists}\\

\bottomrule\end{longtable}

\section{\texorpdfstring{\nolinkurl{PaperXI.ArithmeticFieldInduction69}}{PaperXI.ArithmeticFieldInduction69}}\label{module:PaperXI.ArithmeticFieldInduction69}

\textbf{Selected layer:} base108 arithmetic branch.\quad\textbf{Source lines:} 83.\par\textbf{SHA-256:} {\scriptsize\nolinkurl{84d9fdbbdab35cebcb07e0b60f4363cd4e176ac7b0483a850db3af4ebe56588e}}.

\VerbatimInput[fontsize=\scriptsize,breaklines=true,breakanywhere=true]{excerpts/D_ArithmeticFieldInduction69.txt}

\textbf{Direct project imports:} \nolinkurl{PaperXI.ActualFieldTypeQuotient69}, \nolinkurl{PaperXI.DerivedCentralCharacter69}, \nolinkurl{PaperXI.CompactModCenterInduction68}.

\begin{longtable}{@{}p{0.07\textwidth}p{0.12\textwidth}p{0.75\textwidth}@{}}\toprule Line&Kind&Name\\\midrule\endhead

27&theorem&\nolinkurl{PaperXI.ArithmeticFieldInduction69.fieldGroup_open}\\

32&def&\nolinkurl{PaperXI.ArithmeticFieldInduction69.scalarCharacter}\\

38&theorem&\nolinkurl{PaperXI.ArithmeticFieldInduction69.scalarCharacter_action}\\

46&theorem&\nolinkurl{PaperXI.ArithmeticFieldInduction69.induction_package}\\

72&theorem&\nolinkurl{PaperXI.ArithmeticFieldInduction69.induced_scalar}\\

\bottomrule\end{longtable}

\section{\texorpdfstring{\nolinkurl{PaperXI.ArithmeticTypeCertificate69}}{PaperXI.ArithmeticTypeCertificate69}}\label{module:PaperXI.ArithmeticTypeCertificate69}

\textbf{Selected layer:} base108 arithmetic branch.\quad\textbf{Source lines:} 106.\par\textbf{SHA-256:} {\scriptsize\nolinkurl{a360d478cc7c93440aa4806bf37892a9cbc4c847d393a8be3079c7fc9da439ac}}.

\VerbatimInput[fontsize=\scriptsize,breaklines=true,breakanywhere=true]{excerpts/D_ArithmeticTypeCertificate69.txt}

\textbf{Direct project imports:} \nolinkurl{PaperXI.ArithmeticFieldInduction69}, \nolinkurl{PaperXI.FullIntertwiningLocus69}.

\begin{longtable}{@{}p{0.07\textwidth}p{0.12\textwidth}p{0.75\textwidth}@{}}\toprule Line&Kind&Name\\\midrule\endhead

24&def&\nolinkurl{PaperXI.ArithmeticTypeCertificate69.centralCharacter}\\

29&theorem&\nolinkurl{PaperXI.ArithmeticTypeCertificate69.centralCharacter_action}\\

35&theorem&\nolinkurl{PaperXI.ArithmeticTypeCertificate69.induced_centralCharacter}\\

41&theorem&\nolinkurl{PaperXI.ArithmeticTypeCertificate69.scalarHom_continuous}\\

56&theorem&\nolinkurl{PaperXI.ArithmeticTypeCertificate69.centralCharacter_locallyConstant}\\

87&theorem&\nolinkurl{PaperXI.ArithmeticTypeCertificate69.from_full_intertwining_locus}\\

\bottomrule\end{longtable}

\section{\texorpdfstring{\nolinkurl{PaperXI.AtomicMultiset}}{PaperXI.AtomicMultiset}}\label{module:PaperXI.AtomicMultiset}

\textbf{Selected layer:} v0.44 parent.\quad\textbf{Source lines:} 122.\par\textbf{SHA-256:} {\scriptsize\nolinkurl{e9b0e86f94713a76638c5bde3c53ba97140aec5c543aad941c9bba61ede4d3f6}}.

\VerbatimInput[fontsize=\scriptsize,breaklines=true,breakanywhere=true]{excerpts/D_AtomicMultiset.txt}

\textbf{Direct project imports:} none.

\begin{longtable}{@{}p{0.07\textwidth}p{0.12\textwidth}p{0.75\textwidth}@{}}\toprule Line&Kind&Name\\\midrule\endhead

14&structure&\nolinkurl{PaperXI.AtomicMultiset.Block}\\

21&def&\nolinkurl{PaperXI.AtomicMultiset.blockMap}\\

24&theorem&\nolinkurl{PaperXI.AtomicMultiset.blockMap_atom}\\

27&theorem&\nolinkurl{PaperXI.AtomicMultiset.blockMap_length}\\

30&theorem&\nolinkurl{PaperXI.AtomicMultiset.blockMap_id}\\

33&theorem&\nolinkurl{PaperXI.AtomicMultiset.blockMap_comp}\\

36&def&\nolinkurl{PaperXI.AtomicMultiset.extend}\\

41&theorem&\nolinkurl{PaperXI.AtomicMultiset.extend_singleton}\\

44&theorem&\nolinkurl{PaperXI.AtomicMultiset.extend_cons}\\

47&theorem&\nolinkurl{PaperXI.AtomicMultiset.extend_card}\\

50&theorem&\nolinkurl{PaperXI.AtomicMultiset.extend_id}\\

54&theorem&\nolinkurl{PaperXI.AtomicMultiset.extend_comp}\\

59&theorem&\nolinkurl{PaperXI.AtomicMultiset.extend_unique}\\

70&theorem&\nolinkurl{PaperXI.AtomicMultiset.extend_injective}\\

80&def&\nolinkurl{PaperXI.AtomicMultiset.extendEquiv}\\

94&def&\nolinkurl{PaperXI.AtomicMultiset.dimension}\\

97&theorem&\nolinkurl{PaperXI.AtomicMultiset.dimension_extend}\\

107&def&\nolinkurl{PaperXI.AtomicMultiset.descend}\\

111&theorem&\nolinkurl{PaperXI.AtomicMultiset.descend_mk}\\

116&theorem&\nolinkurl{PaperXI.AtomicMultiset.extend_natural}\\

\bottomrule\end{longtable}

\section{\texorpdfstring{\nolinkurl{PaperXI.AtomicNaturality}}{PaperXI.AtomicNaturality}}\label{module:PaperXI.AtomicNaturality}

\textbf{Selected layer:} v0.44 parent.\quad\textbf{Source lines:} 106.\par\textbf{SHA-256:} {\scriptsize\nolinkurl{ba747e172fb8b6398cc4ce4076c4f8dc2c3f5fcff98a6a24f1e3f1b4a4584873}}.

\VerbatimInput[fontsize=\scriptsize,breaklines=true,breakanywhere=true]{excerpts/D_AtomicNaturality.txt}

\textbf{Direct project imports:} \nolinkurl{PaperXI.AtomicReconstruction}.

\begin{longtable}{@{}p{0.07\textwidth}p{0.12\textwidth}p{0.75\textwidth}@{}}\toprule Line&Kind&Name\\\midrule\endhead

15&theorem&\nolinkurl{PaperXI.AtomicNaturality.reconstruction_from_representatives}\\

22&theorem&\nolinkurl{PaperXI.AtomicNaturality.blocks_twist}\\

36&theorem&\nolinkurl{PaperXI.AtomicNaturality.realize_twist}\\

53&theorem&\nolinkurl{PaperXI.AtomicNaturality.candidate_twist}\\

67&theorem&\nolinkurl{PaperXI.AtomicNaturality.extend_injective_iff}\\

84&theorem&\nolinkurl{PaperXI.AtomicNaturality.candidate_injective}\\

92&theorem&\nolinkurl{PaperXI.AtomicNaturality.decomposition_independence}\\

\bottomrule\end{longtable}

\section{\texorpdfstring{\nolinkurl{PaperXI.AtomicReconstruction}}{PaperXI.AtomicReconstruction}}\label{module:PaperXI.AtomicReconstruction}

\textbf{Selected layer:} v0.44 parent.\quad\textbf{Source lines:} 153.\par\textbf{SHA-256:} {\scriptsize\nolinkurl{e848f4ab5b8717eda8465737b7a726fbb877d3e6effc3728eaeab08698da52cb}}.

\VerbatimInput[fontsize=\scriptsize,breaklines=true,breakanywhere=true]{excerpts/D_AtomicReconstruction.txt}

\textbf{Direct project imports:} \nolinkurl{PaperXI.AtomicMultiset}.

\begin{longtable}{@{}p{0.07\textwidth}p{0.12\textwidth}p{0.75\textwidth}@{}}\toprule Line&Kind&Name\\\midrule\endhead

17&structure&\nolinkurl{PaperXI.AtomicReconstruction.Decomposition}\\

25&def&\nolinkurl{PaperXI.AtomicReconstruction.rawBlocks}\\

28&theorem&\nolinkurl{PaperXI.AtomicReconstruction.rawBlocks_eq}\\

33&theorem&\nolinkurl{PaperXI.AtomicReconstruction.rawBlocks_iso}\\

41&def&\nolinkurl{PaperXI.AtomicReconstruction.blocks}\\

44&theorem&\nolinkurl{PaperXI.AtomicReconstruction.blocks_of_display}\\

52&structure&\nolinkurl{PaperXI.AtomicReconstruction.OrderedQuotient}\\

63&def&\nolinkurl{PaperXI.AtomicReconstruction.ordered}\\

67&theorem&\nolinkurl{PaperXI.AtomicReconstruction.ordered_pairwise}\\

74&theorem&\nolinkurl{PaperXI.AtomicReconstruction.ordered_multiset}\\

83&def&\nolinkurl{PaperXI.AtomicReconstruction.realize}\\

86&theorem&\nolinkurl{PaperXI.AtomicReconstruction.realize_eq}\\

95&theorem&\nolinkurl{PaperXI.AtomicReconstruction.realize_spec}\\

103&def&\nolinkurl{PaperXI.AtomicReconstruction.candidate}\\

108&def&\nolinkurl{PaperXI.AtomicReconstruction.Reconstructs}\\

114&theorem&\nolinkurl{PaperXI.AtomicReconstruction.candidate_spec}\\

121&theorem&\nolinkurl{PaperXI.AtomicReconstruction.candidate_eq}\\

130&theorem&\nolinkurl{PaperXI.AtomicReconstruction.candidate_unique}\\

143&theorem&\nolinkurl{PaperXI.AtomicReconstruction.formal_atomic_reconstruction}\\

148&theorem&\nolinkurl{PaperXI.AtomicReconstruction.candidate_isomorphism}\\

\bottomrule\end{longtable}

\section{\texorpdfstring{\nolinkurl{PaperXI.BilinearHeisenberg}}{PaperXI.BilinearHeisenberg}}\label{module:PaperXI.BilinearHeisenberg}

\textbf{Selected layer:} v0.44 parent.\quad\textbf{Source lines:} 152.\par\textbf{SHA-256:} {\scriptsize\nolinkurl{d4615e543af67b750697c7fdbf5f7d3b790aa0b4d6992aa2246a3aa00399e7d9}}.

\VerbatimInput[fontsize=\scriptsize,breaklines=true,breakanywhere=true]{excerpts/D_BilinearHeisenberg.txt}

\textbf{Direct project imports:} \nolinkurl{PaperXI.SymplecticHeisenbergUniqueness}.

\begin{longtable}{@{}p{0.07\textwidth}p{0.12\textwidth}p{0.75\textwidth}@{}}\toprule Line&Kind&Name\\\midrule\endhead

13&structure&\nolinkurl{PaperXI.BilinearHeisenberg.Element}\\

19&instance&\nolinkurl{PaperXI.BilinearHeisenberg.instMul}\\

21&instance&\nolinkurl{PaperXI.BilinearHeisenberg.instOne}\\

22&instance&\nolinkurl{PaperXI.BilinearHeisenberg.instInv}\\

25&theorem&\nolinkurl{PaperXI.BilinearHeisenberg.mul_vector}\\

27&theorem&\nolinkurl{PaperXI.BilinearHeisenberg.mul_central}\\

29&theorem&\nolinkurl{PaperXI.BilinearHeisenberg.one_vector}\\

30&theorem&\nolinkurl{PaperXI.BilinearHeisenberg.one_central}\\

31&theorem&\nolinkurl{PaperXI.BilinearHeisenberg.inv_vector}\\

32&theorem&\nolinkurl{PaperXI.BilinearHeisenberg.inv_central}\\

35&instance&\nolinkurl{PaperXI.BilinearHeisenberg.instGroup}\\

44&def&\nolinkurl{PaperXI.BilinearHeisenberg.i}\\

49&def&\nolinkurl{PaperXI.BilinearHeisenberg.q}\\

54&theorem&\nolinkurl{PaperXI.BilinearHeisenberg.i_vector}\\

55&theorem&\nolinkurl{PaperXI.BilinearHeisenberg.i_central}\\

56&theorem&\nolinkurl{PaperXI.BilinearHeisenberg.q_apply}\\

59&def&\nolinkurl{PaperXI.BilinearHeisenberg.datum}\\

78&theorem&\nolinkurl{PaperXI.BilinearHeisenberg.commutator}\\

88&theorem&\nolinkurl{PaperXI.BilinearHeisenberg.centralCommutator}\\

96&instance&\nolinkurl{PaperXI.BilinearHeisenberg.instFinite}\\

100&theorem&\nolinkurl{PaperXI.BilinearHeisenberg.order}\\

109&theorem&\nolinkurl{PaperXI.BilinearHeisenberg.source_commutator}\\

121&theorem&\nolinkurl{PaperXI.BilinearHeisenberg.exists_irreducible}\\

136&theorem&\nolinkurl{PaperXI.BilinearHeisenberg.exists_equiv}\\

145&theorem&\nolinkurl{PaperXI.BilinearHeisenberg.finrank_sqrt}\\

\bottomrule\end{longtable}

\section{\texorpdfstring{\nolinkurl{PaperXI.BoundaryExamples}}{PaperXI.BoundaryExamples}}\label{module:PaperXI.BoundaryExamples}

\textbf{Selected layer:} v0.44 parent.\quad\textbf{Source lines:} 30.\par\textbf{SHA-256:} {\scriptsize\nolinkurl{d847590d5166021de5e552840a8b6bac45da991bbac45860b3cd96c3aaeabf89}}.

\textbf{Direct project imports:} \nolinkurl{PaperXI.GeneralChain}, \nolinkurl{PaperXI.ChainSharpness}, \nolinkurl{PaperXI.ScalarCovariance}, \nolinkurl{PaperXI.FourierCovariance}, \nolinkurl{PaperXI.HolonomyResidual}.

\begin{longtable}{@{}p{0.07\textwidth}p{0.12\textwidth}p{0.75\textwidth}@{}}\toprule Line&Kind&Name\\\midrule\endhead

11&theorem&\nolinkurl{PaperXI.BoundaryExamples.empty_chain}\\

14&theorem&\nolinkurl{PaperXI.BoundaryExamples.singleton_chain}\\

17&theorem&\nolinkurl{PaperXI.BoundaryExamples.length_three_square_nonzero}\\

21&theorem&\nolinkurl{PaperXI.BoundaryExamples.length_three_cube_zero}\\

24&theorem&\nolinkurl{PaperXI.BoundaryExamples.inverse_two_mod_five}\\

26&theorem&\nolinkurl{PaperXI.BoundaryExamples.inverse_label_differs}\\

\bottomrule\end{longtable}

\section{\texorpdfstring{\nolinkurl{PaperXI.BreakBridge}}{PaperXI.BreakBridge}}\label{module:PaperXI.BreakBridge}

\textbf{Selected layer:} v0.44 parent.\quad\textbf{Source lines:} 104.\par\textbf{SHA-256:} {\scriptsize\nolinkurl{8867bca4f28ab57140255d385f80f0d940f50f9d7bfd3de435c9b4b6385937d9}}.

\VerbatimInput[fontsize=\scriptsize,breaklines=true,breakanywhere=true]{excerpts/D_BreakBridge.txt}

\textbf{Direct project imports:} \nolinkurl{PaperXI.BreakTensorMap}, \nolinkurl{PaperXI.HeisenbergTensor}, \nolinkurl{PaperXI.CharacterGenerators}.

\begin{longtable}{@{}p{0.07\textwidth}p{0.12\textwidth}p{0.75\textwidth}@{}}\toprule Line&Kind&Name\\\midrule\endhead

18&theorem&\nolinkurl{PaperXI.BreakBridge.ambient_intertwines}\\

30&theorem&\nolinkurl{PaperXI.BreakBridge.quotient_intertwines}\\

47&theorem&\nolinkurl{PaperXI.BreakBridge.stabilizer_intertwines}\\

61&theorem&\nolinkurl{PaperXI.BreakBridge.stabilizer_intertwines_of_factor_phases}\\

75&theorem&\nolinkurl{PaperXI.BreakBridge.semidirect_intertwines}\\

\bottomrule\end{longtable}

\section{\texorpdfstring{\nolinkurl{PaperXI.BreakBridgeAlong}}{PaperXI.BreakBridgeAlong}}\label{module:PaperXI.BreakBridgeAlong}

\textbf{Selected layer:} v0.44 parent.\quad\textbf{Source lines:} 80.\par\textbf{SHA-256:} {\scriptsize\nolinkurl{cc313014e54d45de8e7391509928e831d857cc92c21b22d029f4c5541b823658}}.

\VerbatimInput[fontsize=\scriptsize,breaklines=true,breakanywhere=true]{excerpts/D_BreakBridgeAlong.txt}

\textbf{Direct project imports:} \nolinkurl{PaperXI.BreakBridge}, \nolinkurl{PaperXI.CentralProductMap}, \nolinkurl{PaperXI.HeisenbergTensorFactorization}.

\begin{longtable}{@{}p{0.07\textwidth}p{0.12\textwidth}p{0.75\textwidth}@{}}\toprule Line&Kind&Name\\\midrule\endhead

19&theorem&\nolinkurl{PaperXI.BreakBridgeAlong.quotient_intertwines}\\

45&theorem&\nolinkurl{PaperXI.BreakBridgeAlong.source}\\

\bottomrule\end{longtable}

\section{\texorpdfstring{\nolinkurl{PaperXI.BreakDyadicPhase}}{PaperXI.BreakDyadicPhase}}\label{module:PaperXI.BreakDyadicPhase}

\textbf{Selected layer:} v0.44 parent.\quad\textbf{Source lines:} 115.\par\textbf{SHA-256:} {\scriptsize\nolinkurl{6b938157b4b523f51398fb5d57351690ba2200647a98c90f9d715a12030349df}}.

\VerbatimInput[fontsize=\scriptsize,breaklines=true,breakanywhere=true]{excerpts/D_BreakDyadicPhase.txt}

\textbf{Direct project imports:} \nolinkurl{PaperXI.BreakQuadratic}.

\begin{longtable}{@{}p{0.07\textwidth}p{0.12\textwidth}p{0.75\textwidth}@{}}\toprule Line&Kind&Name\\\midrule\endhead

20&theorem&\nolinkurl{PaperXI.BreakDyadicPhase.gaussSum_product}\\

31&theorem&\nolinkurl{PaperXI.BreakDyadicPhase.normalizedGaussSign_product}\\

42&theorem&\nolinkurl{PaperXI.BreakDyadicPhase.arf_sum}\\

53&theorem&\nolinkurl{PaperXI.BreakDyadicPhase.normalizedGaussSign_empty}\\

66&theorem&\nolinkurl{PaperXI.BreakDyadicPhase.transport_refinement}\\

77&theorem&\nolinkurl{PaperXI.BreakDyadicPhase.arf_of_decomposition}\\

95&theorem&\nolinkurl{PaperXI.BreakDyadicPhase.normalizedGaussSign_of_decomposition}\\

110&theorem&\nolinkurl{PaperXI.BreakDyadicPhase.half_finrank_of_decomposition}\\

\bottomrule\end{longtable}

\section{\texorpdfstring{\nolinkurl{PaperXI.BreakHeisenbergSum}}{PaperXI.BreakHeisenbergSum}}\label{module:PaperXI.BreakHeisenbergSum}

\textbf{Selected layer:} v0.44 parent.\quad\textbf{Source lines:} 143.\par\textbf{SHA-256:} {\scriptsize\nolinkurl{8f7d2102714c520e91a760b6fe970dd48fb99ed979d8d4cd14280b6ffc18e50c}}.

\VerbatimInput[fontsize=\scriptsize,breaklines=true,breakanywhere=true]{excerpts/D_BreakHeisenbergSum.txt}

\textbf{Direct project imports:} \nolinkurl{PaperXI.AbstractHeisenberg}, \nolinkurl{PaperXI.CentralProduct}.

\begin{longtable}{@{}p{0.07\textwidth}p{0.12\textwidth}p{0.75\textwidth}@{}}\toprule Line&Kind&Name\\\midrule\endhead

19&def&\nolinkurl{PaperXI.BreakHeisenbergSum.form}\\

27&theorem&\nolinkurl{PaperXI.BreakHeisenbergSum.form_apply}\\

30&theorem&\nolinkurl{PaperXI.BreakHeisenbergSum.form_isAlt}\\

36&def&\nolinkurl{PaperXI.BreakHeisenbergSum.fromProduct}\\

49&theorem&\nolinkurl{PaperXI.BreakHeisenbergSum.fromProduct_vector}\\

53&theorem&\nolinkurl{PaperXI.BreakHeisenbergSum.fromProduct_central}\\

58&theorem&\nolinkurl{PaperXI.BreakHeisenbergSum.fromProduct_surjective}\\

65&theorem&\nolinkurl{PaperXI.BreakHeisenbergSum.fromProduct_balanced}\\

72&def&\nolinkurl{PaperXI.BreakHeisenbergSum.sumMap}\\

76&theorem&\nolinkurl{PaperXI.BreakHeisenbergSum.sumMap_mk}\\

80&theorem&\nolinkurl{PaperXI.BreakHeisenbergSum.sumMap_center}\\

84&theorem&\nolinkurl{PaperXI.BreakHeisenbergSum.sumMap_surjective}\\

91&theorem&\nolinkurl{PaperXI.BreakHeisenbergSum.sumMap_injective}\\

116&def&\nolinkurl{PaperXI.BreakHeisenbergSum.sumEquiv}\\

120&theorem&\nolinkurl{PaperXI.BreakHeisenbergSum.sumEquiv_mk}\\

124&theorem&\nolinkurl{PaperXI.BreakHeisenbergSum.sumEquiv_center}\\

130&def&\nolinkurl{PaperXI.BreakHeisenbergSum.decompositionEquiv}\\

136&theorem&\nolinkurl{PaperXI.BreakHeisenbergSum.decompositionEquiv_center}\\

\bottomrule\end{longtable}

\section{\texorpdfstring{\nolinkurl{PaperXI.BreakMoving}}{PaperXI.BreakMoving}}\label{module:PaperXI.BreakMoving}

\textbf{Selected layer:} v0.44 parent.\quad\textbf{Source lines:} 116.\par\textbf{SHA-256:} {\scriptsize\nolinkurl{e72bb01b58c7e900a66b0919ad0b78a71e7ae979c3ed55025b524b2d80519da9}}.

\VerbatimInput[fontsize=\scriptsize,breaklines=true,breakanywhere=true]{excerpts/D_BreakMoving.txt}

\textbf{Direct project imports:} \nolinkurl{PaperXI.MovingPhase}, \nolinkurl{PaperXI.OrthogonalPhase}.

\begin{longtable}{@{}p{0.07\textwidth}p{0.12\textwidth}p{0.75\textwidth}@{}}\toprule Line&Kind&Name\\\midrule\endhead

17&def&\nolinkurl{PaperXI.BreakMoving.productEquiv}\\

30&theorem&\nolinkurl{PaperXI.BreakMoving.productEquiv_coe}\\

34&theorem&\nolinkurl{PaperXI.BreakMoving.productEquiv_symm_coe}\\

39&theorem&\nolinkurl{PaperXI.BreakMoving.productEquiv_intertwines}\\

51&def&\nolinkurl{PaperXI.BreakMoving.equiv}\\

57&theorem&\nolinkurl{PaperXI.BreakMoving.equiv_coe}\\

62&theorem&\nolinkurl{PaperXI.BreakMoving.equiv_intertwines}\\

72&theorem&\nolinkurl{PaperXI.BreakMoving.equiv_preserves}\\

89&theorem&\nolinkurl{PaperXI.BreakMoving.isCompl_of_moving_fixedPointFree}\\

109&theorem&\nolinkurl{PaperXI.BreakMoving.movingForm_nondegenerate}\\

\bottomrule\end{longtable}

\section{\texorpdfstring{\nolinkurl{PaperXI.BreakOddPhase}}{PaperXI.BreakOddPhase}}\label{module:PaperXI.BreakOddPhase}

\textbf{Selected layer:} v0.44 parent.\quad\textbf{Source lines:} 164.\par\textbf{SHA-256:} {\scriptsize\nolinkurl{32575366712c4c41dec1087b72f1a84ef94ced9b0af536673412a053eec700cf}}.

\VerbatimInput[fontsize=\scriptsize,breaklines=true,breakanywhere=true]{excerpts/D_BreakOddPhase.txt}

\textbf{Direct project imports:} \nolinkurl{PaperXI.BreakMoving}.

\begin{longtable}{@{}p{0.07\textwidth}p{0.12\textwidth}p{0.75\textwidth}@{}}\toprule Line&Kind&Name\\\midrule\endhead

20&theorem&\nolinkurl{PaperXI.BreakOddPhase.component_fixedPointFree}\\

42&theorem&\nolinkurl{PaperXI.BreakOddPhase.product_of_moving_forms}\\

58&theorem&\nolinkurl{PaperXI.BreakOddPhase.product}\\

74&theorem&\nolinkurl{PaperXI.BreakOddPhase.determinants_ne_zero}\\

85&theorem&\nolinkurl{PaperXI.BreakOddPhase.empty}\\

92&theorem&\nolinkurl{PaperXI.BreakOddPhase.zero_moving_factor}\\

100&theorem&\nolinkurl{PaperXI.BreakOddPhase.internal_intertwines}\\

112&theorem&\nolinkurl{PaperXI.BreakOddPhase.internal_form}\\

126&theorem&\nolinkurl{PaperXI.BreakOddPhase.internal_moving_form}\\

147&theorem&\nolinkurl{PaperXI.BreakOddPhase.source}\\

\bottomrule\end{longtable}

\section{\texorpdfstring{\nolinkurl{PaperXI.BreakQuadratic}}{PaperXI.BreakQuadratic}}\label{module:PaperXI.BreakQuadratic}

\textbf{Selected layer:} v0.44 parent.\quad\textbf{Source lines:} 104.\par\textbf{SHA-256:} {\scriptsize\nolinkurl{37bf741c07a7c2c90e5139d5bd490d61a60eccc0d2573bcf0d51e59ce81f4f6a}}.

\VerbatimInput[fontsize=\scriptsize,breaklines=true,breakanywhere=true]{excerpts/D_BreakQuadratic.txt}

\textbf{Direct project imports:} \nolinkurl{PaperXI.DyadicAbstractOrthogonal}.

\begin{longtable}{@{}p{0.07\textwidth}p{0.12\textwidth}p{0.75\textwidth}@{}}\toprule Line&Kind&Name\\\midrule\endhead

16&def&\nolinkurl{PaperXI.BreakQuadratic.form}\\

25&theorem&\nolinkurl{PaperXI.BreakQuadratic.form_apply}\\

28&theorem&\nolinkurl{PaperXI.BreakQuadratic.form_alternating}\\

32&theorem&\nolinkurl{PaperXI.BreakQuadratic.form_single_right}\\

42&theorem&\nolinkurl{PaperXI.BreakQuadratic.form_nondegenerate}\\

51&def&\nolinkurl{PaperXI.BreakQuadratic.refinement}\\

57&theorem&\nolinkurl{PaperXI.BreakQuadratic.refinement_apply}\\

63&theorem&\nolinkurl{PaperXI.BreakQuadratic.refinement_polar}\\

71&theorem&\nolinkurl{PaperXI.BreakQuadratic.gauss_sum_product}\\

82&theorem&\nolinkurl{PaperXI.BreakQuadratic.finrank_eq_twice_half}\\

91&theorem&\nolinkurl{PaperXI.BreakQuadratic.half_finrank_product}\\

\bottomrule\end{longtable}

\section{\texorpdfstring{\nolinkurl{PaperXI.BreakTensorFactorization}}{PaperXI.BreakTensorFactorization}}\label{module:PaperXI.BreakTensorFactorization}

\textbf{Selected layer:} v0.44 parent.\quad\textbf{Source lines:} 77.\par\textbf{SHA-256:} {\scriptsize\nolinkurl{9201c286d4ce3d4c53eda5ba1a756940f87a36ce7ddfb768a007ec158398ec24}}.

\VerbatimInput[fontsize=\scriptsize,breaklines=true,breakanywhere=true]{excerpts/D_BreakTensorFactorization.txt}

\textbf{Direct project imports:} \nolinkurl{PaperXI.CentralProductComparison}, \nolinkurl{PaperXI.HeisenbergTensorComparison}, \nolinkurl{PaperXI.BreakHeisenbergSum}, \nolinkurl{PaperXI.AbstractHeisenbergExtension}.

\begin{longtable}{@{}p{0.07\textwidth}p{0.12\textwidth}p{0.75\textwidth}@{}}\toprule Line&Kind&Name\\\midrule\endhead

23&def&\nolinkurl{PaperXI.BreakTensorFactorization.source}\\

53&theorem&\nolinkurl{PaperXI.BreakTensorFactorization.odd_source}\\

\bottomrule\end{longtable}

\section{\texorpdfstring{\nolinkurl{PaperXI.BreakTensorMap}}{PaperXI.BreakTensorMap}}\label{module:PaperXI.BreakTensorMap}

\textbf{Selected layer:} v0.44 parent.\quad\textbf{Source lines:} 100.\par\textbf{SHA-256:} {\scriptsize\nolinkurl{52fc48e66b2d78e5a85d528b3d8eccdee4b0668c3dc10b4ba49f3c44482caebb}}.

\VerbatimInput[fontsize=\scriptsize,breaklines=true,breakanywhere=true]{excerpts/D_BreakTensorMap.txt}

\textbf{Direct project imports:} \nolinkurl{PaperXI.FiniteTensor}.

\begin{longtable}{@{}p{0.07\textwidth}p{0.12\textwidth}p{0.75\textwidth}@{}}\toprule Line&Kind&Name\\\midrule\endhead

14&def&\nolinkurl{PaperXI.BreakTensorMap.tensor}\\

17&theorem&\nolinkurl{PaperXI.BreakTensorMap.tensor_tprod}\\

22&theorem&\nolinkurl{PaperXI.BreakTensorMap.tensor_symm_tprod}\\

27&theorem&\nolinkurl{PaperXI.BreakTensorMap.tensor_symm}\\

35&theorem&\nolinkurl{PaperXI.BreakTensorMap.intertwines}\\

49&theorem&\nolinkurl{PaperXI.BreakTensorMap.descends}\\

62&def&\nolinkurl{PaperXI.BreakTensorMap.stabilizer}\\

66&theorem&\nolinkurl{PaperXI.BreakTensorMap.stabilizer_tprod}\\

73&def&\nolinkurl{PaperXI.BreakTensorMap.productCharacter}\\

78&theorem&\nolinkurl{PaperXI.BreakTensorMap.productCharacter_apply}\\

82&theorem&\nolinkurl{PaperXI.BreakTensorMap.scalar_intertwines}\\

\bottomrule\end{longtable}

\section{\texorpdfstring{\nolinkurl{PaperXI.CanonicalDVRValuation}}{PaperXI.CanonicalDVRValuation}}\label{module:PaperXI.CanonicalDVRValuation}

\textbf{Selected layer:} v0.44 parent.\quad\textbf{Source lines:} 225.\par\textbf{SHA-256:} {\scriptsize\nolinkurl{d30233e7e4f5fffd0c9ee97fc3e54a765389da26cf8f917bd7c158e249eae6a1}}.

\VerbatimInput[fontsize=\scriptsize,breaklines=true,breakanywhere=true]{excerpts/D_CanonicalDVRValuation.txt}

\textbf{Direct project imports:} \nolinkurl{PaperXI.LocalFieldResidueExtension}.

\begin{longtable}{@{}p{0.07\textwidth}p{0.12\textwidth}p{0.75\textwidth}@{}}\toprule Line&Kind&Name\\\midrule\endhead

28&abbrev&\nolinkurl{PaperXI.CanonicalDVRValuation.canonicalValuation}\\

32&def&\nolinkurl{PaperXI.CanonicalDVRValuation.integerRingEquiv}\\

37&theorem&\nolinkurl{PaperXI.CanonicalDVRValuation.integerRingEquiv_coe}\\

43&theorem&\nolinkurl{PaperXI.CanonicalDVRValuation.integerRing_image}\\

50&theorem&\nolinkurl{PaperXI.CanonicalDVRValuation.integer_iff_exists_lift}\\

59&theorem&\nolinkurl{PaperXI.CanonicalDVRValuation.unit_value}\\

65&theorem&\nolinkurl{PaperXI.CanonicalDVRValuation.maximalIdeal_iff_value_lt_one}\\

72&theorem&\nolinkurl{PaperXI.CanonicalDVRValuation.uniformizer_isUniformizer}\\

85&theorem&\nolinkurl{PaperXI.CanonicalDVRValuation.uniformizer_value}\\

92&theorem&\nolinkurl{PaperXI.CanonicalDVRValuation.uniformizer_value_lt_one}\\

108&abbrev&\nolinkurl{PaperXI.CanonicalDVRValuation.upperRing}\\

109&abbrev&\nolinkurl{PaperXI.CanonicalDVRValuation.upperField}\\

110&abbrev&\nolinkurl{PaperXI.CanonicalDVRValuation.upperValuation}\\

115&def&\nolinkurl{PaperXI.CanonicalDVRValuation.upperIntegerRingEquiv}\\

121&theorem&\nolinkurl{PaperXI.CanonicalDVRValuation.upperIntegerRingEquiv_coe}\\

129&theorem&\nolinkurl{PaperXI.CanonicalDVRValuation.upper_integerRing_image}\\

136&theorem&\nolinkurl{PaperXI.CanonicalDVRValuation.upper_integer_iff_exists_lift}\\

144&theorem&\nolinkurl{PaperXI.CanonicalDVRValuation.upper_unit_value}\\

152&theorem&\nolinkurl{PaperXI.CanonicalDVRValuation.base_unit_value}\\

161&theorem&\nolinkurl{PaperXI.CanonicalDVRValuation.base_uniformizer_generates_upper}\\

171&theorem&\nolinkurl{PaperXI.CanonicalDVRValuation.base_uniformizer_isUniformizer}\\

182&theorem&\nolinkurl{PaperXI.CanonicalDVRValuation.base_uniformizer_value}\\

191&def&\nolinkurl{PaperXI.CanonicalDVRValuation.chosenBaseUniformizer}\\

195&theorem&\nolinkurl{PaperXI.CanonicalDVRValuation.chosenBaseUniformizer_irreducible}\\

200&theorem&\nolinkurl{PaperXI.CanonicalDVRValuation.chosenBaseUniformizer_span}\\

207&theorem&\nolinkurl{PaperXI.CanonicalDVRValuation.chosenBaseUniformizer_isUpperUniformizer}\\

216&theorem&\nolinkurl{PaperXI.CanonicalDVRValuation.chosenBaseUniformizer_value}\\

\bottomrule\end{longtable}

\section{\texorpdfstring{\nolinkurl{PaperXI.CentralCharacterAction}}{PaperXI.CentralCharacterAction}}\label{module:PaperXI.CentralCharacterAction}

\textbf{Selected layer:} v0.44 parent.\quad\textbf{Source lines:} 111.\par\textbf{SHA-256:} {\scriptsize\nolinkurl{7d5028051c4d9345e5ff3724fa65990af1cb4900a25053f54c82f6e7dafcf04f}}.

\VerbatimInput[fontsize=\scriptsize,breaklines=true,breakanywhere=true]{excerpts/D_CentralCharacterAction.txt}

\textbf{Direct project imports:} \nolinkurl{PaperXI.CentralExtensionCocycle}, \nolinkurl{PaperXI.TwistedAction}.

\begin{longtable}{@{}p{0.07\textwidth}p{0.12\textwidth}p{0.75\textwidth}@{}}\toprule Line&Kind&Name\\\midrule\endhead

15&theorem&\nolinkurl{PaperXI.CentralCharacter.section_operators_one}\\

20&theorem&\nolinkurl{PaperXI.CentralCharacter.section_operators_mul}\\

32&abbrev&\nolinkurl{PaperXI.CentralCharacter.Algebra}\\

38&abbrev&\nolinkurl{PaperXI.CentralCharacter.basis}\\

42&def&\nolinkurl{PaperXI.CentralCharacter.toAlgebraAction}\\

50&theorem&\nolinkurl{PaperXI.CentralCharacter.toAlgebraAction_basis}\\

59&theorem&\nolinkurl{PaperXI.CentralCharacter.action_normal_form}\\

76&theorem&\nolinkurl{PaperXI.CentralCharacter.action_stable_iff}\\

93&theorem&\nolinkurl{PaperXI.CentralCharacter.action_intertwines_iff}\\

\bottomrule\end{longtable}

\section{\texorpdfstring{\nolinkurl{PaperXI.CentralCharacterInduction}}{PaperXI.CentralCharacterInduction}}\label{module:PaperXI.CentralCharacterInduction}

\textbf{Selected layer:} base108 arithmetic branch.\quad\textbf{Source lines:} 62.\par\textbf{SHA-256:} {\scriptsize\nolinkurl{26563eced498d52375622ece6a9b1081ff7d338c5f94399babdd60a81e749349}}.

\VerbatimInput[fontsize=\scriptsize,breaklines=true,breakanywhere=true]{excerpts/D_CentralCharacterInduction.txt}

\textbf{Direct project imports:} \nolinkurl{PaperXI.OpenSubgroupSmoothInduction}.

\begin{longtable}{@{}p{0.07\textwidth}p{0.12\textwidth}p{0.75\textwidth}@{}}\toprule Line&Kind&Name\\\midrule\endhead

22&theorem&\nolinkurl{PaperXI.CentralCharacterInduction.ind_central_scalar}\\

37&theorem&\nolinkurl{PaperXI.CentralCharacterInduction.intertwiner_eq_zero_of_scalar_ne}\\

51&theorem&\nolinkurl{PaperXI.CentralCharacterInduction.ind_intertwiner_eq_zero}\\

\bottomrule\end{longtable}

\section{\texorpdfstring{\nolinkurl{PaperXI.CentralCharacterIrreducible}}{PaperXI.CentralCharacterIrreducible}}\label{module:PaperXI.CentralCharacterIrreducible}

\textbf{Selected layer:} v0.44 parent.\quad\textbf{Source lines:} 49.\par\textbf{SHA-256:} {\scriptsize\nolinkurl{bb4d95bc7100a200f2b5e915893df883a1e7768c06b5e5006432e5b6c4b3d367}}.

\VerbatimInput[fontsize=\scriptsize,breaklines=true,breakanywhere=true]{excerpts/D_CentralCharacterIrreducible.txt}

\textbf{Direct project imports:} \nolinkurl{PaperXI.CentralCharacterAction}, \nolinkurl{PaperXI.CentralCharacterRepresentation}, \nolinkurl{PaperXI.AlgebraActionModule}, \nolinkurl{PaperXI.RepresentationStability}.

\begin{longtable}{@{}p{0.07\textwidth}p{0.12\textwidth}p{0.75\textwidth}@{}}\toprule Line&Kind&Name\\\midrule\endhead

18&theorem&\nolinkurl{PaperXI.CentralCharacter.irreducible_iff_simple}\\

34&theorem&\nolinkurl{PaperXI.CentralCharacter.simpleModule}\\

42&theorem&\nolinkurl{PaperXI.CentralCharacter.ofAlgebraAction_irreducible_iff}\\

\bottomrule\end{longtable}

\section{\texorpdfstring{\nolinkurl{PaperXI.CentralCharacterRepresentation}}{PaperXI.CentralCharacterRepresentation}}\label{module:PaperXI.CentralCharacterRepresentation}

\textbf{Selected layer:} v0.44 parent.\quad\textbf{Source lines:} 83.\par\textbf{SHA-256:} {\scriptsize\nolinkurl{115c7bf9c958e23517b9990efc2179bb7595d5c416e839a96ee9705c1448d3fe}}.

\VerbatimInput[fontsize=\scriptsize,breaklines=true,breakanywhere=true]{excerpts/D_CentralCharacterRepresentation.txt}

\textbf{Direct project imports:} \nolinkurl{PaperXI.CentralCharacterAction}.

\begin{longtable}{@{}p{0.07\textwidth}p{0.12\textwidth}p{0.75\textwidth}@{}}\toprule Line&Kind&Name\\\midrule\endhead

14&def&\nolinkurl{PaperXI.CentralCharacter.ofAlgebraAction}\\

27&theorem&\nolinkurl{PaperXI.CentralCharacter.ofAlgebraAction_apply}\\

33&theorem&\nolinkurl{PaperXI.CentralCharacter.ofAlgebraAction_central}\\

40&theorem&\nolinkurl{PaperXI.CentralCharacter.ofAlgebraAction_section}\\

46&theorem&\nolinkurl{PaperXI.CentralCharacter.ofAlgebraAction_isUnit}\\

53&theorem&\nolinkurl{PaperXI.CentralCharacter.of_toAlgebraAction}\\

62&theorem&\nolinkurl{PaperXI.CentralCharacter.to_ofAlgebraAction}\\

74&def&\nolinkurl{PaperXI.CentralCharacter.actionEquiv}\\

\bottomrule\end{longtable}

\section{\texorpdfstring{\nolinkurl{PaperXI.CentralExtension}}{PaperXI.CentralExtension}}\label{module:PaperXI.CentralExtension}

\textbf{Selected layer:} v0.44 parent.\quad\textbf{Source lines:} 85.\par\textbf{SHA-256:} {\scriptsize\nolinkurl{738dea99a4dea853d5968a45499d8cd3804c8aae571749faa7ecd0f2beada19a}}.

\VerbatimInput[fontsize=\scriptsize,breaklines=true,breakanywhere=true]{excerpts/D_CentralExtension.txt}

\textbf{Direct project imports:} \nolinkurl{PaperXI.CocycleCommutator}.

\begin{longtable}{@{}p{0.07\textwidth}p{0.12\textwidth}p{0.75\textwidth}@{}}\toprule Line&Kind&Name\\\midrule\endhead

11&structure&\nolinkurl{PaperXI.CentralExtension.Datum}\\

22&theorem&\nolinkurl{PaperXI.CentralExtension.q_i}\\

28&def&\nolinkurl{PaperXI.CentralExtension.sectionMap}\\

32&theorem&\nolinkurl{PaperXI.CentralExtension.section_one}\\

35&theorem&\nolinkurl{PaperXI.CentralExtension.q_section}\\

41&theorem&\nolinkurl{PaperXI.CentralExtension.exists_coordinate}\\

49&def&\nolinkurl{PaperXI.CentralExtension.coordinate}\\

51&theorem&\nolinkurl{PaperXI.CentralExtension.coordinate_spec}\\

56&theorem&\nolinkurl{PaperXI.CentralExtension.normal_form}\\

60&theorem&\nolinkurl{PaperXI.CentralExtension.coordinate_eq_of_eq}\\

67&theorem&\nolinkurl{PaperXI.CentralExtension.coordinate_one}\\

70&theorem&\nolinkurl{PaperXI.CentralExtension.coordinate_section}\\

73&theorem&\nolinkurl{PaperXI.CentralExtension.coordinate_i}\\

77&theorem&\nolinkurl{PaperXI.CentralExtension.normal_form_unique}\\

\bottomrule\end{longtable}

\section{\texorpdfstring{\nolinkurl{PaperXI.CentralExtensionCocycle}}{PaperXI.CentralExtensionCocycle}}\label{module:PaperXI.CentralExtensionCocycle}

\textbf{Selected layer:} v0.44 parent.\quad\textbf{Source lines:} 127.\par\textbf{SHA-256:} {\scriptsize\nolinkurl{20408e0fc0cb3e5f4edced01e3f8829ed5b80a905ed55a06d4545e27170151fa}}.

\VerbatimInput[fontsize=\scriptsize,breaklines=true,breakanywhere=true]{excerpts/D_CentralExtensionCocycle.txt}

\textbf{Direct project imports:} \nolinkurl{PaperXI.CentralExtension}.

\begin{longtable}{@{}p{0.07\textwidth}p{0.12\textwidth}p{0.75\textwidth}@{}}\toprule Line&Kind&Name\\\midrule\endhead

15&def&\nolinkurl{PaperXI.CentralExtension.defect}\\

17&theorem&\nolinkurl{PaperXI.CentralExtension.section_mul}\\

22&theorem&\nolinkurl{PaperXI.CentralExtension.defect_one_left}\\

25&theorem&\nolinkurl{PaperXI.CentralExtension.defect_one_right}\\

28&theorem&\nolinkurl{PaperXI.CentralExtension.normal_form_mul}\\

38&theorem&\nolinkurl{PaperXI.CentralExtension.coordinate_mul}\\

49&theorem&\nolinkurl{PaperXI.CentralExtension.defect_assoc}\\

56&def&\nolinkurl{PaperXI.CentralExtension.centralCocycle}\\

60&theorem&\nolinkurl{PaperXI.CentralExtension.normalized_centralCocycle}\\

64&def&\nolinkurl{PaperXI.CentralExtension.scalarCocycle}\\

69&theorem&\nolinkurl{PaperXI.CentralExtension.scalarCocycle_coeff}\\

72&theorem&\nolinkurl{PaperXI.CentralExtension.normalized_scalarCocycle}\\

76&theorem&\nolinkurl{PaperXI.CentralExtension.product_normal_form}\\

82&theorem&\nolinkurl{PaperXI.CentralExtension.lift_commutation}\\

100&def&\nolinkurl{PaperXI.CentralExtension.centralCommutator}\\

102&theorem&\nolinkurl{PaperXI.CentralExtension.centralCommutator_eq}\\

110&theorem&\nolinkurl{PaperXI.CentralExtension.i_centralCommutator}\\

116&theorem&\nolinkurl{PaperXI.CentralExtension.centralCommutator_lift_independent}\\

123&theorem&\nolinkurl{PaperXI.CentralExtension.scalar_commutator}\\

\bottomrule\end{longtable}

\section{\texorpdfstring{\nolinkurl{PaperXI.CentralExtensionComparison}}{PaperXI.CentralExtensionComparison}}\label{module:PaperXI.CentralExtensionComparison}

\textbf{Selected layer:} v0.44 parent.\quad\textbf{Source lines:} 61.\par\textbf{SHA-256:} {\scriptsize\nolinkurl{e36d9e9e7da551d6db27596259bc69841e32a06e82d7b1c79f3bb69942a35082}}.

\VerbatimInput[fontsize=\scriptsize,breaklines=true,breakanywhere=true]{excerpts/D_CentralExtensionComparison.txt}

\textbf{Direct project imports:} \nolinkurl{PaperXI.CentralExtensionUniqueness}, \nolinkurl{PaperXI.HeisenbergTransport}.

\begin{longtable}{@{}p{0.07\textwidth}p{0.12\textwidth}p{0.75\textwidth}@{}}\toprule Line&Kind&Name\\\midrule\endhead

15&def&\nolinkurl{PaperXI.CentralExtensionComparison.equivAlong}\\

25&theorem&\nolinkurl{PaperXI.CentralExtensionComparison.equivAlong_intertwines}\\

38&theorem&\nolinkurl{PaperXI.CentralExtensionComparison.exists_unique_along}\\

\bottomrule\end{longtable}

\section{\texorpdfstring{\nolinkurl{PaperXI.CentralExtensionExistence}}{PaperXI.CentralExtensionExistence}}\label{module:PaperXI.CentralExtensionExistence}

\textbf{Selected layer:} v0.44 parent.\quad\textbf{Source lines:} 43.\par\textbf{SHA-256:} {\scriptsize\nolinkurl{c073b326ebaed1aafad9eeaf1e4cbb0707dfa8112ae901e0b421f9d114227a2a}}.

\VerbatimInput[fontsize=\scriptsize,breaklines=true,breakanywhere=true]{excerpts/D_CentralExtensionExistence.txt}

\textbf{Direct project imports:} \nolinkurl{PaperXI.CentralExtensionUniqueness}, \nolinkurl{PaperXI.CentralCharacterRepresentation}.

\begin{longtable}{@{}p{0.07\textwidth}p{0.12\textwidth}p{0.75\textwidth}@{}}\toprule Line&Kind&Name\\\midrule\endhead

16&theorem&\nolinkurl{PaperXI.CentralExtensionExistence.exists_irreducible}\\

\bottomrule\end{longtable}

\section{\texorpdfstring{\nolinkurl{PaperXI.CentralExtensionInducing}}{PaperXI.CentralExtensionInducing}}\label{module:PaperXI.CentralExtensionInducing}

\textbf{Selected layer:} base108 arithmetic branch.\quad\textbf{Source lines:} 208.\par\textbf{SHA-256:} {\scriptsize\nolinkurl{a435de848de5be66c076c64e660fe71344453fd66a68dc1d2f842854b984317f}}.

\VerbatimInput[fontsize=\scriptsize,breaklines=true,breakanywhere=true]{excerpts/D_CentralExtensionInducing.txt}

\textbf{Direct project imports:} \nolinkurl{PaperXI.SmoothDegreeQuotient}, \nolinkurl{PaperXI.RepresentationStability}.

\begin{longtable}{@{}p{0.07\textwidth}p{0.12\textwidth}p{0.75\textwidth}@{}}\toprule Line&Kind&Name\\\midrule\endhead

21&theorem&\nolinkurl{PaperXI.CentralExtensionInducing.hfree_from_degree}\\

32&def&\nolinkurl{PaperXI.CentralExtensionInducing.normalFormHom}\\

45&theorem&\nolinkurl{PaperXI.CentralExtensionInducing.normalFormHom_apply}\\

50&def&\nolinkurl{PaperXI.CentralExtensionInducing.extendedSubgroup}\\

53&theorem&\nolinkurl{PaperXI.CentralExtensionInducing.le_extendedSubgroup}\\

58&theorem&\nolinkurl{PaperXI.CentralExtensionInducing.central_mem_extendedSubgroup}\\

63&theorem&\nolinkurl{PaperXI.CentralExtensionInducing.extendedSubgroup_eq_sup}\\

75&theorem&\nolinkurl{PaperXI.CentralExtensionInducing.normalFormHom_injective}\\

91&def&\nolinkurl{PaperXI.CentralExtensionInducing.normalFormEquiv}\\

97&theorem&\nolinkurl{PaperXI.CentralExtensionInducing.normalFormEquiv_coe}\\

104&def&\nolinkurl{PaperXI.CentralExtensionInducing.subgroupInclusion}\\

109&def&\nolinkurl{PaperXI.CentralExtensionInducing.centralElement}\\

116&def&\nolinkurl{PaperXI.CentralExtensionInducing.productRepresentation}\\

127&def&\nolinkurl{PaperXI.CentralExtensionInducing.extendedRepresentation}\\

135&theorem&\nolinkurl{PaperXI.CentralExtensionInducing.extendedRepresentation_normalForm}\\

143&theorem&\nolinkurl{PaperXI.CentralExtensionInducing.extendedRepresentation_restrict}\\

153&theorem&\nolinkurl{PaperXI.CentralExtensionInducing.extendedRepresentation_central}\\

167&theorem&\nolinkurl{PaperXI.CentralExtensionInducing.extendedRepresentation_irreducible}\\

183&theorem&\nolinkurl{PaperXI.CentralExtensionInducing.extendedSubgroup_isOpen}\\

189&theorem&\nolinkurl{PaperXI.CentralExtensionInducing.extendedRepresentation_smooth}\\

\bottomrule\end{longtable}

\section{\texorpdfstring{\nolinkurl{PaperXI.CentralExtensionInductionConsequences}}{PaperXI.CentralExtensionInductionConsequences}}\label{module:PaperXI.CentralExtensionInductionConsequences}

\textbf{Selected layer:} base108 arithmetic branch.\quad\textbf{Source lines:} 95.\par\textbf{SHA-256:} {\scriptsize\nolinkurl{6e51c61426161f53ded9cdde06049b7358cf387ee47e88de993af2d98a6838f1}}.

\VerbatimInput[fontsize=\scriptsize,breaklines=true,breakanywhere=true]{excerpts/D_CentralExtensionInductionConsequences.txt}

\textbf{Direct project imports:} \nolinkurl{PaperXI.CentralExtensionInducing}, \nolinkurl{PaperXI.CentralCharacterInduction}.

\begin{longtable}{@{}p{0.07\textwidth}p{0.12\textwidth}p{0.75\textwidth}@{}}\toprule Line&Kind&Name\\\midrule\endhead

22&def&\nolinkurl{PaperXI.CentralExtensionInductionConsequences.inducedRepresentation}\\

31&theorem&\nolinkurl{PaperXI.CentralExtensionInductionConsequences.induced_central_scalar}\\

41&theorem&\nolinkurl{PaperXI.CentralExtensionInductionConsequences.induced_nontrivial}\\

50&theorem&\nolinkurl{PaperXI.CentralExtensionInductionConsequences.induced_intertwiner_eq_zero}\\

62&theorem&\nolinkurl{PaperXI.CentralExtensionInductionConsequences.induced_not_equiv}\\

84&theorem&\nolinkurl{PaperXI.CentralExtensionInductionConsequences.induced_isSmooth}\\

\bottomrule\end{longtable}

\section{\texorpdfstring{\nolinkurl{PaperXI.CentralExtensionSectionChange}}{PaperXI.CentralExtensionSectionChange}}\label{module:PaperXI.CentralExtensionSectionChange}

\textbf{Selected layer:} v0.44 parent.\quad\textbf{Source lines:} 267.\par\textbf{SHA-256:} {\scriptsize\nolinkurl{5360a7116c397b9588d620f7febc89b633f697bd8cec091dd4ecc76468e03f4f}}.

\VerbatimInput[fontsize=\scriptsize,breaklines=true,breakanywhere=true]{excerpts/D_CentralExtensionSectionChange.txt}

\textbf{Direct project imports:} \nolinkurl{PaperXI.CentralExtensionCocycle}, \nolinkurl{PaperXI.TwistedRescaling}, \nolinkurl{PaperXI.CentralCharacterAction}.

\begin{longtable}{@{}p{0.07\textwidth}p{0.12\textwidth}p{0.75\textwidth}@{}}\toprule Line&Kind&Name\\\midrule\endhead

18&theorem&\nolinkurl{PaperXI.CentralExtension.cocycle_eq_of_coeff_eq}\\

27&def&\nolinkurl{PaperXI.CentralExtension.sectionFactor}\\

29&theorem&\nolinkurl{PaperXI.CentralExtension.section_factor}\\

33&theorem&\nolinkurl{PaperXI.CentralExtension.sectionFactor_one}\\

38&def&\nolinkurl{PaperXI.CentralExtension.sectionDefect}\\

41&theorem&\nolinkurl{PaperXI.CentralExtension.i_sectionDefect}\\

46&theorem&\nolinkurl{PaperXI.CentralExtension.sectionDefect_mul}\\

51&theorem&\nolinkurl{PaperXI.CentralExtension.sectionDefect_eq_rescale}\\

70&def&\nolinkurl{PaperXI.CentralExtension.sectionCocycle}\\

76&theorem&\nolinkurl{PaperXI.CentralExtension.sectionCocycle_eq_rescale}\\

81&theorem&\nolinkurl{PaperXI.CentralExtension.normalized_sectionCocycle}\\

88&def&\nolinkurl{PaperXI.CentralExtension.sectionScalarCocycle}\\

94&theorem&\nolinkurl{PaperXI.CentralExtension.sectionScalarCocycle_eq_rescale}\\

106&theorem&\nolinkurl{PaperXI.CentralExtension.normalized_sectionScalarCocycle}\\

114&theorem&\nolinkurl{PaperXI.CentralExtension.sectionScalarCocycle_commutator}\\

121&theorem&\nolinkurl{PaperXI.CentralExtension.sectionScalarCocycle_radical}\\

130&def&\nolinkurl{PaperXI.CentralExtension.sectionAlgebraEquiv}\\

142&theorem&\nolinkurl{PaperXI.CentralExtension.sectionAlgebraEquiv_basis}\\

152&theorem&\nolinkurl{PaperXI.CentralExtension.sectionAlgebraEquiv_symm_basis}\\

165&def&\nolinkurl{PaperXI.CentralExtension.betweenSectionsEquiv}\\

174&theorem&\nolinkurl{PaperXI.CentralExtension.betweenSectionsEquiv_basis}\\

191&theorem&\nolinkurl{PaperXI.CentralExtension.sectionOperators_factor}\\

202&theorem&\nolinkurl{PaperXI.CentralExtension.sectionOperators_mul}\\

214&def&\nolinkurl{PaperXI.CentralExtension.sectionAction}\\

224&theorem&\nolinkurl{PaperXI.CentralExtension.sectionAction_basis}\\

235&theorem&\nolinkurl{PaperXI.CentralExtension.sectionAction_eq_transport}\\

254&theorem&\nolinkurl{PaperXI.CentralExtension.sectionAction_betweenSections}\\

\bottomrule\end{longtable}

\section{\texorpdfstring{\nolinkurl{PaperXI.CentralExtensionStructure}}{PaperXI.CentralExtensionStructure}}\label{module:PaperXI.CentralExtensionStructure}

\textbf{Selected layer:} v0.44 parent.\quad\textbf{Source lines:} 117.\par\textbf{SHA-256:} {\scriptsize\nolinkurl{601bab4de431a620e5b9f2c84ea154cd63b7336382468bdf4dd6cca2a0a1ee35}}.

\VerbatimInput[fontsize=\scriptsize,breaklines=true,breakanywhere=true]{excerpts/D_CentralExtensionStructure.txt}

\textbf{Direct project imports:} \nolinkurl{PaperXI.CentralExtensionCocycle}.

\begin{longtable}{@{}p{0.07\textwidth}p{0.12\textwidth}p{0.75\textwidth}@{}}\toprule Line&Kind&Name\\\midrule\endhead

17&def&\nolinkurl{PaperXI.CentralExtension.coordinates}\\

26&theorem&\nolinkurl{PaperXI.CentralExtension.coordinates_apply}\\

29&theorem&\nolinkurl{PaperXI.CentralExtension.coordinates_symm_apply}\\

33&theorem&\nolinkurl{PaperXI.CentralExtension.finite_middle}\\

38&theorem&\nolinkurl{PaperXI.CentralExtension.natCard_eq}\\

42&theorem&\nolinkurl{PaperXI.CentralExtension.card_eq}\\

48&theorem&\nolinkurl{PaperXI.CentralExtension.radical_scalarCocycle_eq_bot_iff}\\

69&theorem&\nolinkurl{PaperXI.CentralExtension.radical_scalarCocycle_eq_bot_of_faithful}\\

81&theorem&\nolinkurl{PaperXI.CentralExtension.radical_scalarCocycle_of_injective}\\

95&theorem&\nolinkurl{PaperXI.CentralExtension.radical_scalarCocycle_eq_bot_of_central}\\

102&theorem&\nolinkurl{PaperXI.CentralExtension.center_eq_range}\\

\bottomrule\end{longtable}

\section{\texorpdfstring{\nolinkurl{PaperXI.CentralExtensionUniqueness}}{PaperXI.CentralExtensionUniqueness}}\label{module:PaperXI.CentralExtensionUniqueness}

\textbf{Selected layer:} v0.44 parent.\quad\textbf{Source lines:} 99.\par\textbf{SHA-256:} {\scriptsize\nolinkurl{2e5ae550279047849c064e92382606a085429ff1dabb04de593c14006d39a9d2}}.

\VerbatimInput[fontsize=\scriptsize,breaklines=true,breakanywhere=true]{excerpts/D_CentralExtensionUniqueness.txt}

\textbf{Direct project imports:} \nolinkurl{PaperXI.CentralCharacterIrreducible}, \nolinkurl{PaperXI.TwistedModules}, \nolinkurl{PaperXI.SchurScalar}.

\begin{longtable}{@{}p{0.07\textwidth}p{0.12\textwidth}p{0.75\textwidth}@{}}\toprule Line&Kind&Name\\\midrule\endhead

18&theorem&\nolinkurl{PaperXI.CentralExtensionUniqueness.finiteDimensional}\\

33&theorem&\nolinkurl{PaperXI.CentralExtensionUniqueness.finrank_sqrt}\\

47&theorem&\nolinkurl{PaperXI.CentralExtensionUniqueness.exists_equiv}\\

68&def&\nolinkurl{PaperXI.CentralExtensionUniqueness.equiv}\\

76&theorem&\nolinkurl{PaperXI.CentralExtensionUniqueness.equiv_intertwines}\\

87&theorem&\nolinkurl{PaperXI.CentralExtensionUniqueness.unique_scalar}\\

\bottomrule\end{longtable}

\section{\texorpdfstring{\nolinkurl{PaperXI.CentralProduct}}{PaperXI.CentralProduct}}\label{module:PaperXI.CentralProduct}

\textbf{Selected layer:} v0.44 parent.\quad\textbf{Source lines:} 159.\par\textbf{SHA-256:} {\scriptsize\nolinkurl{9d1e6a4ece84a862ef75383b66b97af694c9ac2298fc2047a6c4844bc6aa34af}}.

\VerbatimInput[fontsize=\scriptsize,breaklines=true,breakanywhere=true]{excerpts/D_CentralProduct.txt}

\textbf{Direct project imports:} \nolinkurl{PaperXI.CentralExtension}.

\begin{longtable}{@{}p{0.07\textwidth}p{0.12\textwidth}p{0.75\textwidth}@{}}\toprule Line&Kind&Name\\\midrule\endhead

15&abbrev&\nolinkurl{PaperXI.CentralProduct.Ambient}\\

18&def&\nolinkurl{PaperXI.CentralProduct.balanced}\\

23&def&\nolinkurl{PaperXI.CentralProduct.relation}\\

25&theorem&\nolinkurl{PaperXI.CentralProduct.balanced_central}\\

32&instance&\nolinkurl{PaperXI.CentralProduct.relation_normal}\\

39&def&\nolinkurl{PaperXI.CentralProduct.Group}\\

43&def&\nolinkurl{PaperXI.CentralProduct.mk}\\

45&theorem&\nolinkurl{PaperXI.CentralProduct.mk_surjective}\\

47&theorem&\nolinkurl{PaperXI.CentralProduct.mk_balanced}\\

50&def&\nolinkurl{PaperXI.CentralProduct.center}\\

52&theorem&\nolinkurl{PaperXI.CentralProduct.center_apply}\\

54&theorem&\nolinkurl{PaperXI.CentralProduct.center_injective}\\

69&def&\nolinkurl{PaperXI.CentralProduct.ambientProjection}\\

74&theorem&\nolinkurl{PaperXI.CentralProduct.relation_le_ker_projection}\\

79&def&\nolinkurl{PaperXI.CentralProduct.projection}\\

82&theorem&\nolinkurl{PaperXI.CentralProduct.projection_mk}\\

85&theorem&\nolinkurl{PaperXI.CentralProduct.projection_center}\\

89&theorem&\nolinkurl{PaperXI.CentralProduct.projection_surjective}\\

96&theorem&\nolinkurl{PaperXI.CentralProduct.mk_centers}\\

105&theorem&\nolinkurl{PaperXI.CentralProduct.range_center_eq_ker_projection}\\

122&theorem&\nolinkurl{PaperXI.CentralProduct.center_commutes}\\

131&def&\nolinkurl{PaperXI.CentralProduct.datum}\\

139&theorem&\nolinkurl{PaperXI.CentralProduct.datum_i}\\

141&theorem&\nolinkurl{PaperXI.CentralProduct.datum_q}\\

144&def&\nolinkurl{PaperXI.CentralProduct.descend}\\

148&theorem&\nolinkurl{PaperXI.CentralProduct.descend_mk}\\

153&theorem&\nolinkurl{PaperXI.CentralProduct.hom_ext}\\

\bottomrule\end{longtable}

\section{\texorpdfstring{\nolinkurl{PaperXI.CentralProductCommutator}}{PaperXI.CentralProductCommutator}}\label{module:PaperXI.CentralProductCommutator}

\textbf{Selected layer:} v0.44 parent.\quad\textbf{Source lines:} 72.\par\textbf{SHA-256:} {\scriptsize\nolinkurl{00857ad77d53b9578db693336a58114a88ad7c58f683a74132a7730fff7c4dbe}}.

\VerbatimInput[fontsize=\scriptsize,breaklines=true,breakanywhere=true]{excerpts/D_CentralProductCommutator.txt}

\textbf{Direct project imports:} \nolinkurl{PaperXI.CentralProduct}, \nolinkurl{PaperXI.CentralExtensionStructure}.

\begin{longtable}{@{}p{0.07\textwidth}p{0.12\textwidth}p{0.75\textwidth}@{}}\toprule Line&Kind&Name\\\midrule\endhead

16&theorem&\nolinkurl{PaperXI.CentralProduct.commutator_mk}\\

32&theorem&\nolinkurl{PaperXI.CentralProduct.scalar_commutator_product}\\

48&theorem&\nolinkurl{PaperXI.CentralProduct.radical_eq_bot}\\

67&theorem&\nolinkurl{PaperXI.CentralProduct.full_center}\\

\bottomrule\end{longtable}

\section{\texorpdfstring{\nolinkurl{PaperXI.CentralProductComparison}}{PaperXI.CentralProductComparison}}\label{module:PaperXI.CentralProductComparison}

\textbf{Selected layer:} v0.44 parent.\quad\textbf{Source lines:} 149.\par\textbf{SHA-256:} {\scriptsize\nolinkurl{0c1567c5e43684d459663ec1d8c0995f0cdd4d26a9b3b84b56e719e76dcf1a67}}.

\VerbatimInput[fontsize=\scriptsize,breaklines=true,breakanywhere=true]{excerpts/D_CentralProductComparison.txt}

\textbf{Direct project imports:} \nolinkurl{PaperXI.CentralProductFactors}.

\begin{longtable}{@{}p{0.07\textwidth}p{0.12\textwidth}p{0.75\textwidth}@{}}\toprule Line&Kind&Name\\\midrule\endhead

22&def&\nolinkurl{PaperXI.CentralProductComparison.factors}\\

24&theorem&\nolinkurl{PaperXI.CentralProductComparison.factors_single}\\

30&theorem&\nolinkurl{PaperXI.CentralProductComparison.factors_projection}\\

45&def&\nolinkurl{PaperXI.CentralProductComparison.centerProduct}\\

52&theorem&\nolinkurl{PaperXI.CentralProductComparison.factors_centers}\\

72&def&\nolinkurl{PaperXI.CentralProductComparison.ambient}\\

76&theorem&\nolinkurl{PaperXI.CentralProductComparison.ambient_apply}\\

81&theorem&\nolinkurl{PaperXI.CentralProductComparison.ambient_balanced}\\

87&def&\nolinkurl{PaperXI.CentralProductComparison.comparison}\\

90&theorem&\nolinkurl{PaperXI.CentralProductComparison.comparison_mk}\\

93&theorem&\nolinkurl{PaperXI.CentralProductComparison.comparison_center}\\

99&theorem&\nolinkurl{PaperXI.CentralProductComparison.comparison_projection}\\

108&theorem&\nolinkurl{PaperXI.CentralProductComparison.comparison_injective}\\

126&theorem&\nolinkurl{PaperXI.CentralProductComparison.comparison_surjective}\\

138&def&\nolinkurl{PaperXI.CentralProductComparison.equiv}\\

142&theorem&\nolinkurl{PaperXI.CentralProductComparison.equiv_center}\\

146&theorem&\nolinkurl{PaperXI.CentralProductComparison.equiv_mk}\\

\bottomrule\end{longtable}

\section{\texorpdfstring{\nolinkurl{PaperXI.CentralProductFactors}}{PaperXI.CentralProductFactors}}\label{module:PaperXI.CentralProductFactors}

\textbf{Selected layer:} v0.44 parent.\quad\textbf{Source lines:} 89.\par\textbf{SHA-256:} {\scriptsize\nolinkurl{2c019b264e9f537a7886668193c3e1777ac631b16c10cf02acbd25e8f30b3e82}}.

\VerbatimInput[fontsize=\scriptsize,breaklines=true,breakanywhere=true]{excerpts/D_CentralProductFactors.txt}

\textbf{Direct project imports:} \nolinkurl{PaperXI.CentralProduct}.

\begin{longtable}{@{}p{0.07\textwidth}p{0.12\textwidth}p{0.75\textwidth}@{}}\toprule Line&Kind&Name\\\midrule\endhead

13&def&\nolinkurl{PaperXI.CentralProduct.factor}\\

16&theorem&\nolinkurl{PaperXI.CentralProduct.factor_apply}\\

20&theorem&\nolinkurl{PaperXI.CentralProduct.projection_factor}\\

28&theorem&\nolinkurl{PaperXI.CentralProduct.factor_center}\\

43&theorem&\nolinkurl{PaperXI.CentralProduct.factor_injective}\\

61&theorem&\nolinkurl{PaperXI.CentralProduct.factor_commute}\\

72&theorem&\nolinkurl{PaperXI.CentralProduct.center_surjective_of_isEmpty}\\

81&def&\nolinkurl{PaperXI.CentralProduct.emptyEquiv}\\

86&theorem&\nolinkurl{PaperXI.CentralProduct.emptyEquiv_center}\\

\bottomrule\end{longtable}

\section{\texorpdfstring{\nolinkurl{PaperXI.CentralProductMap}}{PaperXI.CentralProductMap}}\label{module:PaperXI.CentralProductMap}

\textbf{Selected layer:} v0.44 parent.\quad\textbf{Source lines:} 64.\par\textbf{SHA-256:} {\scriptsize\nolinkurl{a5dad822167784fc5033f5187a85a8e99a172466c1f65742380c20fa01d6e4d8}}.

\VerbatimInput[fontsize=\scriptsize,breaklines=true,breakanywhere=true]{excerpts/D_CentralProductMap.txt}

\textbf{Direct project imports:} \nolinkurl{PaperXI.CentralProduct}.

\begin{longtable}{@{}p{0.07\textwidth}p{0.12\textwidth}p{0.75\textwidth}@{}}\toprule Line&Kind&Name\\\midrule\endhead

15&def&\nolinkurl{PaperXI.CentralProduct.ambientMap}\\

21&def&\nolinkurl{PaperXI.CentralProduct.map}\\

29&theorem&\nolinkurl{PaperXI.CentralProduct.map_mk}\\

34&theorem&\nolinkurl{PaperXI.CentralProduct.map_center}\\

40&def&\nolinkurl{PaperXI.CentralProduct.congr}\\

55&theorem&\nolinkurl{PaperXI.CentralProduct.congr_mk}\\

60&theorem&\nolinkurl{PaperXI.CentralProduct.congr_center}\\

\bottomrule\end{longtable}

\section{\texorpdfstring{\nolinkurl{PaperXI.ChainSharpness}}{PaperXI.ChainSharpness}}\label{module:PaperXI.ChainSharpness}

\textbf{Selected layer:} v0.44 parent.\quad\textbf{Source lines:} 25.\par\textbf{SHA-256:} {\scriptsize\nolinkurl{9dc676cc958252b1a0d2ad618ee5f85bcded1dbd8813022e05f9a32c85129ece}}.

\textbf{Direct project imports:} \nolinkurl{PaperXI.MonodromyChain}.

\begin{longtable}{@{}p{0.07\textwidth}p{0.12\textwidth}p{0.75\textwidth}@{}}\toprule Line&Kind&Name\\\midrule\endhead

9&theorem&\nolinkurl{PaperXI.MonodromyChain.identity_path}\\

16&theorem&\nolinkurl{PaperXI.MonodromyChain.identity_shift_power_ne_zero}\\

\bottomrule\end{longtable}

\section{\texorpdfstring{\nolinkurl{PaperXI.CharacterAttachment108}}{PaperXI.CharacterAttachment108}}\label{module:PaperXI.CharacterAttachment108}

\textbf{Selected layer:} base108 arithmetic branch.\quad\textbf{Source lines:} 73.\par\textbf{SHA-256:} {\scriptsize\nolinkurl{3ab5684d139d5f725153a9bc2d64d616384a28734b3839f0bc3c1f89856219fe}}.

\VerbatimInput[fontsize=\scriptsize,breaklines=true,breakanywhere=true]{excerpts/D_CharacterAttachment108.txt}

\textbf{Direct project imports:} \nolinkurl{PaperXI.TypeSideDescent}, \nolinkurl{PaperXI.CharacterNormalizer61}, \nolinkurl{PaperXI.RepresentationStability}, \nolinkurl{PaperXI.IntertwiningAlong105}.

\begin{longtable}{@{}p{0.07\textwidth}p{0.12\textwidth}p{0.75\textwidth}@{}}\toprule Line&Kind&Name\\\midrule\endhead

14&theorem&\nolinkurl{PaperXI.CharacterAttachment108.irreducible_of_restriction}\\

23&def&\nolinkurl{PaperXI.CharacterAttachment108.insideCharacter}\\

27&theorem&\nolinkurl{PaperXI.CharacterAttachment108.scalar_of_contains}\\

50&theorem&\nolinkurl{PaperXI.CharacterAttachment108.character_of_scalar}\\

\bottomrule\end{longtable}

\section{\texorpdfstring{\nolinkurl{PaperXI.CharacterGenerators}}{PaperXI.CharacterGenerators}}\label{module:PaperXI.CharacterGenerators}

\textbf{Selected layer:} v0.44 parent.\quad\textbf{Source lines:} 59.\par\textbf{SHA-256:} {\scriptsize\nolinkurl{5f499da97dbae7d5948989b633f51e14f7adac1468a62e22af5ec653fbc2b1a6}}.

\VerbatimInput[fontsize=\scriptsize,breaklines=true,breakanywhere=true]{excerpts/D_CharacterGenerators.txt}

\textbf{Direct project imports:} \nolinkurl{PaperXI.CyclicBridge}.

\begin{longtable}{@{}p{0.07\textwidth}p{0.12\textwidth}p{0.75\textwidth}@{}}\toprule Line&Kind&Name\\\midrule\endhead

12&theorem&\nolinkurl{PaperXI.CharacterGenerators.eq_one_on_closure}\\

18&theorem&\nolinkurl{PaperXI.CharacterGenerators.trivial}\\

25&theorem&\nolinkurl{PaperXI.CharacterGenerators.trivial_iff}\\

35&theorem&\nolinkurl{PaperXI.CharacterGenerators.bridge_from_actions}\\

\bottomrule\end{longtable}

\section{\texorpdfstring{\nolinkurl{PaperXI.CharacterNormalizer61}}{PaperXI.CharacterNormalizer61}}\label{module:PaperXI.CharacterNormalizer61}

\textbf{Selected layer:} base108 arithmetic branch.\quad\textbf{Source lines:} 104.\par\textbf{SHA-256:} {\scriptsize\nolinkurl{4d1fb07d5d106c49247b3900dea666ce7859054359e7328de46b86d3cf3cdf1f}}.

\VerbatimInput[fontsize=\scriptsize,breaklines=true,breakanywhere=true]{excerpts/D_CharacterNormalizer61.txt}

\textbf{Direct project imports:} \nolinkurl{PaperXI.NormalizerIntertwining60}.

\begin{longtable}{@{}p{0.07\textwidth}p{0.12\textwidth}p{0.75\textwidth}@{}}\toprule Line&Kind&Name\\\midrule\endhead

18&def&\nolinkurl{PaperXI.CharacterNormalizer61.graph}\\

21&def&\nolinkurl{PaperXI.CharacterNormalizer61.normalizer}\\

25&theorem&\nolinkurl{PaperXI.CharacterNormalizer61.conjugation_mem}\\

34&theorem&\nolinkurl{PaperXI.CharacterNormalizer61.character_covariance}\\

45&theorem&\nolinkurl{PaperXI.CharacterNormalizer61.normalizes_domain}\\

57&theorem&\nolinkurl{PaperXI.CharacterNormalizer61.domain_le_normalizer}\\

76&theorem&\nolinkurl{PaperXI.CharacterNormalizer61.character_intertwinedAt}\\

85&def&\nolinkurl{PaperXI.CharacterNormalizer61.compactPart}\\

91&theorem&\nolinkurl{PaperXI.CharacterNormalizer61.normalizer_open}\\

95&theorem&\nolinkurl{PaperXI.CharacterNormalizer61.compactPart_compact_open}\\

\bottomrule\end{longtable}

\section{\texorpdfstring{\nolinkurl{PaperXI.CharacterPairing108}}{PaperXI.CharacterPairing108}}\label{module:PaperXI.CharacterPairing108}

\textbf{Selected layer:} base108 arithmetic branch.\quad\textbf{Source lines:} 78.\par\textbf{SHA-256:} {\scriptsize\nolinkurl{e86637f400e8ad723c2b13db11b5fc48fb9110e11a1d764d1925fed4a1e06dfb}}.

\VerbatimInput[fontsize=\scriptsize,breaklines=true,breakanywhere=true]{excerpts/D_CharacterPairing108.txt}

\textbf{Direct project imports:} \nolinkurl{PaperXI.HeisenbergTrace108}.

\begin{longtable}{@{}p{0.07\textwidth}p{0.12\textwidth}p{0.75\textwidth}@{}}\toprule Line&Kind&Name\\\midrule\endhead

22&theorem&\nolinkurl{PaperXI.CharacterPairing108.dimension_formula}\\

65&theorem&\nolinkurl{PaperXI.CharacterPairing108.dimension_one_iff}\\

\bottomrule\end{longtable}

\section{\texorpdfstring{\nolinkurl{PaperXI.CharacterProjector}}{PaperXI.CharacterProjector}}\label{module:PaperXI.CharacterProjector}

\textbf{Selected layer:} v0.44 parent.\quad\textbf{Source lines:} 100.\par\textbf{SHA-256:} {\scriptsize\nolinkurl{c3d3017695299f3ef67a28bee5c519bbe37aa9c0768e52c4fdbb8c9be912d7cc}}.

\VerbatimInput[fontsize=\scriptsize,breaklines=true,breakanywhere=true]{excerpts/D_CharacterProjector.txt}

\textbf{Direct project imports:} none.

\begin{longtable}{@{}p{0.07\textwidth}p{0.12\textwidth}p{0.75\textwidth}@{}}\toprule Line&Kind&Name\\\midrule\endhead

15&def&\nolinkurl{PaperXI.CharacterProjector.element}\\

20&theorem&\nolinkurl{PaperXI.CharacterProjector.coeff_element}\\

26&theorem&\nolinkurl{PaperXI.CharacterProjector.central}\\

42&theorem&\nolinkurl{PaperXI.CharacterProjector.action}\\

50&theorem&\nolinkurl{PaperXI.CharacterProjector.central_scalar}\\

58&theorem&\nolinkurl{PaperXI.CharacterProjector.nontrivial_space}\\

66&theorem&\nolinkurl{PaperXI.CharacterProjector.degree_ne_zero}\\

72&theorem&\nolinkurl{PaperXI.CharacterProjector.irreducible_action}\\

\bottomrule\end{longtable}

\section{\texorpdfstring{\nolinkurl{PaperXI.CharacterSupport108}}{PaperXI.CharacterSupport108}}\label{module:PaperXI.CharacterSupport108}

\textbf{Selected layer:} base108 arithmetic branch.\quad\textbf{Source lines:} 71.\par\textbf{SHA-256:} {\scriptsize\nolinkurl{4a1a02bf6685a09772dbbe49ad1f29037975b2bbc299e736b1eba53df4aa8960}}.

\VerbatimInput[fontsize=\scriptsize,breaklines=true,breakanywhere=true]{excerpts/D_CharacterSupport108.txt}

\textbf{Direct project imports:} \nolinkurl{PaperXI.CharacterAttachment108}, \nolinkurl{PaperXI.FullTypeIntertwining105}, \nolinkurl{PaperXI.IntertwiningSpace107}.

\begin{longtable}{@{}p{0.07\textwidth}p{0.12\textwidth}p{0.75\textwidth}@{}}\toprule Line&Kind&Name\\\midrule\endhead

16&def&\nolinkurl{PaperXI.CharacterSupport108.intoFirst}\\

26&theorem&\nolinkurl{PaperXI.CharacterSupport108.restricted_iff_character}\\

38&theorem&\nolinkurl{PaperXI.CharacterSupport108.full_iff_character}\\

48&theorem&\nolinkurl{PaperXI.CharacterSupport108.full_iff_restricted}\\

54&theorem&\nolinkurl{PaperXI.CharacterSupport108.existsAt_of_dimension_pos}\\

64&theorem&\nolinkurl{PaperXI.CharacterSupport108.beta_of_dimension_pos}\\

\bottomrule\end{longtable}

\section{\texorpdfstring{\nolinkurl{PaperXI.CharacterTwist}}{PaperXI.CharacterTwist}}\label{module:PaperXI.CharacterTwist}

\textbf{Selected layer:} v0.44 parent.\quad\textbf{Source lines:} 54.\par\textbf{SHA-256:} {\scriptsize\nolinkurl{2f04977b3efe71419a28d3985c9ca035ba1e9f43cccb4779638ea69098b5f58e}}.

\VerbatimInput[fontsize=\scriptsize,breaklines=true,breakanywhere=true]{excerpts/D_CharacterTwist.txt}

\textbf{Direct project imports:} \nolinkurl{PaperXI.CyclicInducedCoordinates}.

\begin{longtable}{@{}p{0.07\textwidth}p{0.12\textwidth}p{0.75\textwidth}@{}}\toprule Line&Kind&Name\\\midrule\endhead

11&def&\nolinkurl{PaperXI.CharacterTwist.twist}\\

16&theorem&\nolinkurl{PaperXI.CharacterTwist.twist_apply}\\

20&def&\nolinkurl{PaperXI.CharacterTwist.tensorEquiv}\\

27&def&\nolinkurl{PaperXI.CharacterTwist.mapEquiv}\\

37&theorem&\nolinkurl{PaperXI.CharacterTwist.restrict_eq}\\

42&theorem&\nolinkurl{PaperXI.CharacterTwist.power_action}\\

48&theorem&\nolinkurl{PaperXI.CharacterTwist.closing_scalar}\\

\bottomrule\end{longtable}

\section{\texorpdfstring{\nolinkurl{PaperXI.CliffordAction}}{PaperXI.CliffordAction}}\label{module:PaperXI.CliffordAction}

\textbf{Selected layer:} v0.44 parent.\quad\textbf{Source lines:} 86.\par\textbf{SHA-256:} {\scriptsize\nolinkurl{aff52bad200146645c61ebf8647615481ee1f908eb4fa8aa5ccdd39bf79393bf}}.

\VerbatimInput[fontsize=\scriptsize,breaklines=true,breakanywhere=true]{excerpts/D_CliffordAction.txt}

\textbf{Direct project imports:} \nolinkurl{PaperXI.CliffordIrreducible}, \nolinkurl{PaperXI.CliffordChoices}.

\begin{longtable}{@{}p{0.07\textwidth}p{0.12\textwidth}p{0.75\textwidth}@{}}\toprule Line&Kind&Name\\\midrule\endhead

31&theorem&\nolinkurl{PaperXI.CliffordAction.all_choices}\\

56&theorem&\nolinkurl{PaperXI.CliffordAction.irreducible_all_choices}\\

\bottomrule\end{longtable}

\section{\texorpdfstring{\nolinkurl{PaperXI.CliffordChoices}}{PaperXI.CliffordChoices}}\label{module:PaperXI.CliffordChoices}

\textbf{Selected layer:} v0.44 parent.\quad\textbf{Source lines:} 148.\par\textbf{SHA-256:} {\scriptsize\nolinkurl{1e5599eae92df49fe01d3a846e8a21e070e97ca8aa79461ebafdfaf666c9a9e0}}.

\VerbatimInput[fontsize=\scriptsize,breaklines=true,breakanywhere=true]{excerpts/D_CliffordChoices.txt}

\textbf{Direct project imports:} \nolinkurl{PaperXI.CliffordMultiplicity}, \nolinkurl{PaperXI.ProjectiveCocycle}.

\begin{longtable}{@{}p{0.07\textwidth}p{0.12\textwidth}p{0.75\textwidth}@{}}\toprule Line&Kind&Name\\\midrule\endhead

16&theorem&\nolinkurl{PaperXI.CliffordChoices.operator_rescale}\\

32&theorem&\nolinkurl{PaperXI.CliffordChoices.existsUnique_intertwiner_scalar}\\

46&def&\nolinkurl{PaperXI.CliffordChoices.intertwinerScalar}\\

52&theorem&\nolinkurl{PaperXI.CliffordChoices.intertwinerScalar_spec}\\

61&theorem&\nolinkurl{PaperXI.CliffordChoices.operator_choice}\\

72&def&\nolinkurl{PaperXI.CliffordChoices.representativeFactor}\\

79&theorem&\nolinkurl{PaperXI.CliffordChoices.representative_factorization}\\

85&def&\nolinkurl{PaperXI.CliffordChoices.representativeIntertwiner}\\

103&theorem&\nolinkurl{PaperXI.CliffordChoices.representativeIntertwiner_apply}\\

110&theorem&\nolinkurl{PaperXI.CliffordChoices.operator_representative}\\

126&def&\nolinkurl{PaperXI.CliffordChoices.choiceCochain}\\

135&theorem&\nolinkurl{PaperXI.CliffordChoices.arbitrary_choices}\\

\bottomrule\end{longtable}

\section{\texorpdfstring{\nolinkurl{PaperXI.CliffordCocycle}}{PaperXI.CliffordCocycle}}\label{module:PaperXI.CliffordCocycle}

\textbf{Selected layer:} v0.44 parent.\quad\textbf{Source lines:} 98.\par\textbf{SHA-256:} {\scriptsize\nolinkurl{915d42076be1fed3ce0abb96e2e74367ebbd9d51d78dd6220a3cd73f37a04354}}.

\VerbatimInput[fontsize=\scriptsize,breaklines=true,breakanywhere=true]{excerpts/D_CliffordCocycle.txt}

\textbf{Direct project imports:} \nolinkurl{PaperXI.CliffordChoices}.

\begin{longtable}{@{}p{0.07\textwidth}p{0.12\textwidth}p{0.75\textwidth}@{}}\toprule Line&Kind&Name\\\midrule\endhead

16&def&\nolinkurl{PaperXI.CliffordCocycle.scalar}\\

25&theorem&\nolinkurl{PaperXI.CliffordCocycle.scalar_spec}\\

38&theorem&\nolinkurl{PaperXI.CliffordCocycle.operator_defect}\\

52&theorem&\nolinkurl{PaperXI.CliffordCocycle.sectionOperator_mul}\\

67&theorem&\nolinkurl{PaperXI.CliffordCocycle.sectionOperator_ne_zero}\\

77&def&\nolinkurl{PaperXI.CliffordCocycle.cocycle}\\

85&theorem&\nolinkurl{PaperXI.CliffordCocycle.cocycle_coeff}\\

91&theorem&\nolinkurl{PaperXI.CliffordCocycle.normalized}\\

\bottomrule\end{longtable}

\section{\texorpdfstring{\nolinkurl{PaperXI.CliffordCocycleChoices}}{PaperXI.CliffordCocycleChoices}}\label{module:PaperXI.CliffordCocycleChoices}

\textbf{Selected layer:} v0.44 parent.\quad\textbf{Source lines:} 87.\par\textbf{SHA-256:} {\scriptsize\nolinkurl{b81243b34012e08c729b15becd2058df413240d7047d7975e8f6d119746f4bd8}}.

\VerbatimInput[fontsize=\scriptsize,breaklines=true,breakanywhere=true]{excerpts/D_CliffordCocycleChoices.txt}

\textbf{Direct project imports:} \nolinkurl{PaperXI.CliffordCocycle}.

\begin{longtable}{@{}p{0.07\textwidth}p{0.12\textwidth}p{0.75\textwidth}@{}}\toprule Line&Kind&Name\\\midrule\endhead

18&def&\nolinkurl{PaperXI.CliffordCocycleChoices.choiceCocycle}\\

23&theorem&\nolinkurl{PaperXI.CliffordCocycleChoices.choiceOperator_mul}\\

35&theorem&\nolinkurl{PaperXI.CliffordCocycleChoices.choiceCocycle_unique}\\

55&theorem&\nolinkurl{PaperXI.CliffordCocycleChoices.commutator_choice}\\

62&theorem&\nolinkurl{PaperXI.CliffordCocycleChoices.radical_choice}\\

68&theorem&\nolinkurl{PaperXI.CliffordCocycleChoices.choiceCochain_one}\\

80&theorem&\nolinkurl{PaperXI.CliffordCocycleChoices.choiceCocycle_normalized}\\

\bottomrule\end{longtable}

\section{\texorpdfstring{\nolinkurl{PaperXI.CliffordCocycleClass}}{PaperXI.CliffordCocycleClass}}\label{module:PaperXI.CliffordCocycleClass}

\textbf{Selected layer:} v0.44 parent.\quad\textbf{Source lines:} 111.\par\textbf{SHA-256:} {\scriptsize\nolinkurl{88207bfc7c10216386294ac5631533d2b07f6a592136a204d91b41553bd7711c}}.

\VerbatimInput[fontsize=\scriptsize,breaklines=true,breakanywhere=true]{excerpts/D_CliffordCocycleClass.txt}

\textbf{Direct project imports:} \nolinkurl{PaperXI.CocycleClass}, \nolinkurl{PaperXI.CliffordCocycleChoices}, \nolinkurl{PaperXI.CliffordTransport}, \nolinkurl{PaperXI.CliffordSource}.

\begin{longtable}{@{}p{0.07\textwidth}p{0.12\textwidth}p{0.75\textwidth}@{}}\toprule Line&Kind&Name\\\midrule\endhead

21&def&\nolinkurl{PaperXI.CliffordCocycleClass.value}\\

24&theorem&\nolinkurl{PaperXI.CliffordCocycleClass.commutator_value}\\

28&theorem&\nolinkurl{PaperXI.CliffordCocycleClass.radical_value}\\

32&theorem&\nolinkurl{PaperXI.CliffordCocycleClass.choice_class}\\

42&theorem&\nolinkurl{PaperXI.CliffordCocycleClass.class_of_multiplication}\\

67&theorem&\nolinkurl{PaperXI.CliffordCocycleClass.value_transport}\\

82&def&\nolinkurl{PaperXI.CliffordCocycleClass.sourceValue}\\

101&theorem&\nolinkurl{PaperXI.CliffordCocycleClass.source_transport}\\

\bottomrule\end{longtable}

\section{\texorpdfstring{\nolinkurl{PaperXI.CliffordDatum}}{PaperXI.CliffordDatum}}\label{module:PaperXI.CliffordDatum}

\textbf{Selected layer:} v0.44 parent.\quad\textbf{Source lines:} 112.\par\textbf{SHA-256:} {\scriptsize\nolinkurl{7321b347ed121e9981fdb5e48e5cdcc56ffa0b1ed4c556fb547dcc4a7eab8b5b}}.

\VerbatimInput[fontsize=\scriptsize,breaklines=true,breakanywhere=true]{excerpts/D_CliffordDatum.txt}

\textbf{Direct project imports:} \nolinkurl{PaperXI.QuotientRepresentation}.

\begin{longtable}{@{}p{0.07\textwidth}p{0.12\textwidth}p{0.75\textwidth}@{}}\toprule Line&Kind&Name\\\midrule\endhead

14&abbrev&\nolinkurl{PaperXI.CliffordDatum.restriction}\\

18&def&\nolinkurl{PaperXI.CliffordDatum.conjugation}\\

23&theorem&\nolinkurl{PaperXI.CliffordDatum.conjugation_coe}\\

26&theorem&\nolinkurl{PaperXI.CliffordDatum.conjugation_one}\\

29&theorem&\nolinkurl{PaperXI.CliffordDatum.conjugation_mul}\\

34&theorem&\nolinkurl{PaperXI.CliffordDatum.conjugation_inv_cancel}\\

40&abbrev&\nolinkurl{PaperXI.CliffordDatum.conjugate}\\

44&def&\nolinkurl{PaperXI.CliffordDatum.Invariant}\\

48&def&\nolinkurl{PaperXI.CliffordDatum.sectionMap}\\

53&theorem&\nolinkurl{PaperXI.CliffordDatum.section_one}\\

56&theorem&\nolinkurl{PaperXI.CliffordDatum.quotient_section}\\

64&def&\nolinkurl{PaperXI.CliffordDatum.defect}\\

71&theorem&\nolinkurl{PaperXI.CliffordDatum.defect_coe}\\

74&theorem&\nolinkurl{PaperXI.CliffordDatum.section_mul}\\

78&theorem&\nolinkurl{PaperXI.CliffordDatum.defect_one_left}\\

81&theorem&\nolinkurl{PaperXI.CliffordDatum.defect_one_right}\\

85&def&\nolinkurl{PaperXI.CliffordDatum.coordinate}\\

92&theorem&\nolinkurl{PaperXI.CliffordDatum.normal_form}\\

97&def&\nolinkurl{PaperXI.CliffordDatum.sectionIntertwiner}\\

106&theorem&\nolinkurl{PaperXI.CliffordDatum.sectionIntertwiner_one}\\

\bottomrule\end{longtable}

\section{\texorpdfstring{\nolinkurl{PaperXI.CliffordEnvelope}}{PaperXI.CliffordEnvelope}}\label{module:PaperXI.CliffordEnvelope}

\textbf{Selected layer:} v0.44 parent.\quad\textbf{Source lines:} 147.\par\textbf{SHA-256:} {\scriptsize\nolinkurl{550b0e257a986100dd3ddf95022f7766f8366363760787b3b28b03c08cc04548}}.

\VerbatimInput[fontsize=\scriptsize,breaklines=true,breakanywhere=true]{excerpts/D_CliffordEnvelope.txt}

\textbf{Direct project imports:} \nolinkurl{PaperXI.CliffordSource}, \nolinkurl{PaperXI.CliffordTransport}, \nolinkurl{PaperXI.ElementaryEnvelopeTransport}.

\begin{longtable}{@{}p{0.07\textwidth}p{0.12\textwidth}p{0.75\textwidth}@{}}\toprule Line&Kind&Name\\\midrule\endhead

23&def&\nolinkurl{PaperXI.CliffordEnvelope.pairing}\\

26&theorem&\nolinkurl{PaperXI.CliffordEnvelope.pairing_coe}\\

31&theorem&\nolinkurl{PaperXI.CliffordEnvelope.pairing_self}\\

35&abbrev&\nolinkurl{PaperXI.CliffordEnvelope.Space}\\

40&def&\nolinkurl{PaperXI.CliffordEnvelope.form}\\

43&theorem&\nolinkurl{PaperXI.CliffordEnvelope.form_isAlt}\\

46&theorem&\nolinkurl{PaperXI.CliffordEnvelope.form_nondegenerate}\\

49&theorem&\nolinkurl{PaperXI.CliffordEnvelope.finiteDimensional}\\

54&theorem&\nolinkurl{PaperXI.CliffordEnvelope.form_on_classes}\\

77&theorem&\nolinkurl{PaperXI.CliffordEnvelope.commutator_preserved}\\

92&theorem&\nolinkurl{PaperXI.CliffordEnvelope.exists_normalized_cochain}\\

110&theorem&\nolinkurl{PaperXI.CliffordEnvelope.pairing_preserved}\\

118&def&\nolinkurl{PaperXI.CliffordEnvelope.equiv}\\

124&theorem&\nolinkurl{PaperXI.CliffordEnvelope.equiv_mk}\\

131&theorem&\nolinkurl{PaperXI.CliffordEnvelope.equiv_form}\\

139&theorem&\nolinkurl{PaperXI.CliffordEnvelope.equiv_independent}\\

\bottomrule\end{longtable}

\section{\texorpdfstring{\nolinkurl{PaperXI.CliffordEvaluation}}{PaperXI.CliffordEvaluation}}\label{module:PaperXI.CliffordEvaluation}

\textbf{Selected layer:} v0.44 parent.\quad\textbf{Source lines:} 242.\par\textbf{SHA-256:} {\scriptsize\nolinkurl{405e9bf8ba1bdf8eba00ab3d6063ec48594e3b24d4d919d5a5b862e76bd4effd}}.

\textbf{Direct project imports:} \nolinkurl{PaperXI.IsotypicEvaluation}, \nolinkurl{PaperXI.FiniteRepresentation}, \nolinkurl{PaperXI.CliffordMultiplicity}.

\begin{longtable}{@{}p{0.07\textwidth}p{0.12\textwidth}p{0.75\textwidth}@{}}\toprule Line&Kind&Name\\\midrule\endhead

19&def&\nolinkurl{PaperXI.CliffordEvaluation.map}\\

28&theorem&\nolinkurl{PaperXI.CliffordEvaluation.map_tmul}\\

32&theorem&\nolinkurl{PaperXI.CliffordEvaluation.map_action}\\

40&theorem&\nolinkurl{PaperXI.CliffordEvaluation.bijective}\\

55&def&\nolinkurl{PaperXI.CliffordEvaluation.equiv}\\

60&theorem&\nolinkurl{PaperXI.CliffordEvaluation.equiv_tmul}\\

64&theorem&\nolinkurl{PaperXI.CliffordEvaluation.equiv_action}\\

71&theorem&\nolinkurl{PaperXI.CliffordEvaluation.multiplicity_nontrivial}\\

80&def&\nolinkurl{PaperXI.CliffordEvaluation.choiceMap}\\

90&theorem&\nolinkurl{PaperXI.CliffordEvaluation.choiceMap_apply}\\

96&theorem&\nolinkurl{PaperXI.CliffordEvaluation.choiceMap_injective}\\

107&def&\nolinkurl{PaperXI.CliffordEvaluation.choiceEquiv}\\

116&theorem&\nolinkurl{PaperXI.CliffordEvaluation.choiceEquiv_apply}\\

123&theorem&\nolinkurl{PaperXI.CliffordEvaluation.exists_unique_choice}\\

146&theorem&\nolinkurl{PaperXI.CliffordEvaluation.map_operator}\\

156&theorem&\nolinkurl{PaperXI.CliffordEvaluation.operator_unique}\\

166&theorem&\nolinkurl{PaperXI.CliffordEvaluation.choice_operator_conjugation}\\

184&def&\nolinkurl{PaperXI.CliffordEvaluation.choiceOperator}\\

194&theorem&\nolinkurl{PaperXI.CliffordEvaluation.choiceOperator_tensor}\\

206&theorem&\nolinkurl{PaperXI.CliffordEvaluation.exists_unique_choice_operator}\\

227&theorem&\nolinkurl{PaperXI.CliffordEvaluation.restriction_bijective}\\

234&theorem&\nolinkurl{PaperXI.CliffordEvaluation.restriction_multiplicity_nontrivial}\\

\bottomrule\end{longtable}

\section{\texorpdfstring{\nolinkurl{PaperXI.CliffordHeisenbergAttachment}}{PaperXI.CliffordHeisenbergAttachment}}\label{module:PaperXI.CliffordHeisenbergAttachment}

\textbf{Selected layer:} v0.44 parent.\quad\textbf{Source lines:} 144.\par\textbf{SHA-256:} {\scriptsize\nolinkurl{bfc127d6190ce5ad509d5e37b1b9507b92996f620a90657c855d90516c6648e7}}.

\VerbatimInput[fontsize=\scriptsize,breaklines=true,breakanywhere=true]{excerpts/D_CliffordHeisenbergAttachment.txt}

\textbf{Direct project imports:} \nolinkurl{PaperXI.OddEnvelopeHeisenberg}, \nolinkurl{PaperXI.DyadicHeisenbergLift}, \nolinkurl{PaperXI.CliffordEnvelope}.

\begin{longtable}{@{}p{0.07\textwidth}p{0.12\textwidth}p{0.75\textwidth}@{}}\toprule Line&Kind&Name\\\midrule\endhead

27&def&\nolinkurl{PaperXI.CliffordHeisenbergAttachment.oddFactor}\\

32&theorem&\nolinkurl{PaperXI.CliffordHeisenbergAttachment.odd_commutator_on_classes}\\

50&theorem&\nolinkurl{PaperXI.CliffordHeisenbergAttachment.odd_exists}\\

72&theorem&\nolinkurl{PaperXI.CliffordHeisenbergAttachment.odd_unique}\\

91&def&\nolinkurl{PaperXI.CliffordHeisenbergAttachment.dyadicFactor}\\

98&theorem&\nolinkurl{PaperXI.CliffordHeisenbergAttachment.dyadic_square}\\

106&theorem&\nolinkurl{PaperXI.CliffordHeisenbergAttachment.dyadic_exists}\\

128&theorem&\nolinkurl{PaperXI.CliffordHeisenbergAttachment.dyadic_unique}\\

\bottomrule\end{longtable}

\section{\texorpdfstring{\nolinkurl{PaperXI.CliffordIrreducible}}{PaperXI.CliffordIrreducible}}\label{module:PaperXI.CliffordIrreducible}

\textbf{Selected layer:} v0.44 parent.\quad\textbf{Source lines:} 98.\par\textbf{SHA-256:} {\scriptsize\nolinkurl{39476ad63852c2ef28b547eed56fe8805910a7beabbdde97f4cdbef6b5b83b63}}.

\VerbatimInput[fontsize=\scriptsize,breaklines=true,breakanywhere=true]{excerpts/D_CliffordIrreducible.txt}

\textbf{Direct project imports:} \nolinkurl{PaperXI.CliffordMultiplicity}, \nolinkurl{PaperXI.CliffordEvaluation}, \nolinkurl{PaperXI.TensorSubspace}.

\begin{longtable}{@{}p{0.07\textwidth}p{0.12\textwidth}p{0.75\textwidth}@{}}\toprule Line&Kind&Name\\\midrule\endhead

21&abbrev&\nolinkurl{PaperXI.CliffordIrreducible.Isotypic}\\

25&theorem&\nolinkurl{PaperXI.CliffordIrreducible.multiplicity_finiteDimensional}\\

31&theorem&\nolinkurl{PaperXI.CliffordIrreducible.multiplicity_nontrivial}\\

37&def&\nolinkurl{PaperXI.CliffordIrreducible.evaluation}\\

41&theorem&\nolinkurl{PaperXI.CliffordIrreducible.evaluation_tmul}\\

47&theorem&\nolinkurl{PaperXI.CliffordIrreducible.section_action}\\

64&theorem&\nolinkurl{PaperXI.CliffordIrreducible.group_action}\\

80&theorem&\nolinkurl{PaperXI.CliffordIrreducible.irreducible}\\

\bottomrule\end{longtable}

\section{\texorpdfstring{\nolinkurl{PaperXI.CliffordMultiplicity}}{PaperXI.CliffordMultiplicity}}\label{module:PaperXI.CliffordMultiplicity}

\textbf{Selected layer:} v0.44 parent.\quad\textbf{Source lines:} 241.\par\textbf{SHA-256:} {\scriptsize\nolinkurl{9ad23300239f1a6cb777e10ea3a8efb66596640c0644862baf2db01c0089b1df}}.

\VerbatimInput[fontsize=\scriptsize,breaklines=true,breakanywhere=true]{excerpts/D_CliffordMultiplicity.txt}

\textbf{Direct project imports:} \nolinkurl{PaperXI.CliffordDatum}.

\begin{longtable}{@{}p{0.07\textwidth}p{0.12\textwidth}p{0.75\textwidth}@{}}\toprule Line&Kind&Name\\\midrule\endhead

13&abbrev&\nolinkurl{PaperXI.CliffordMultiplicity.Space}\\

16&theorem&\nolinkurl{PaperXI.CliffordMultiplicity.space_intertwines}\\

20&theorem&\nolinkurl{PaperXI.CliffordMultiplicity.inverse_intertwines}\\

27&theorem&\nolinkurl{PaperXI.CliffordMultiplicity.action_conjugation}\\

35&def&\nolinkurl{PaperXI.CliffordMultiplicity.operatorMap}\\

48&theorem&\nolinkurl{PaperXI.CliffordMultiplicity.operatorMap_apply}\\

54&def&\nolinkurl{PaperXI.CliffordMultiplicity.inverseMap}\\

74&theorem&\nolinkurl{PaperXI.CliffordMultiplicity.inverseMap_apply}\\

80&def&\nolinkurl{PaperXI.CliffordMultiplicity.operator}\\

100&theorem&\nolinkurl{PaperXI.CliffordMultiplicity.operator_apply}\\

105&theorem&\nolinkurl{PaperXI.CliffordMultiplicity.operator_symm_apply}\\

111&def&\nolinkurl{PaperXI.CliffordMultiplicity.sectionOperator}\\

116&theorem&\nolinkurl{PaperXI.CliffordMultiplicity.sectionOperator_one}\\

128&theorem&\nolinkurl{PaperXI.CliffordMultiplicity.intertwiner_apply}\\

134&def&\nolinkurl{PaperXI.CliffordMultiplicity.productIntertwiner}\\

145&theorem&\nolinkurl{PaperXI.CliffordMultiplicity.productIntertwiner_apply}\\

152&def&\nolinkurl{PaperXI.CliffordMultiplicity.actionEquiv}\\

165&theorem&\nolinkurl{PaperXI.CliffordMultiplicity.actionEquiv_apply}\\

169&def&\nolinkurl{PaperXI.CliffordMultiplicity.subgroupIntertwiner}\\

185&theorem&\nolinkurl{PaperXI.CliffordMultiplicity.subgroupIntertwiner_apply}\\

191&def&\nolinkurl{PaperXI.CliffordMultiplicity.defectIntertwiner}\\

210&theorem&\nolinkurl{PaperXI.CliffordMultiplicity.defectIntertwiner_apply}\\

215&theorem&\nolinkurl{PaperXI.CliffordMultiplicity.operator_product}\\

227&theorem&\nolinkurl{PaperXI.CliffordMultiplicity.operator_subgroup_mul}\\

\bottomrule\end{longtable}

\section{\texorpdfstring{\nolinkurl{PaperXI.CliffordPacket}}{PaperXI.CliffordPacket}}\label{module:PaperXI.CliffordPacket}

\textbf{Selected layer:} v0.44 parent.\quad\textbf{Source lines:} 68.\par\textbf{SHA-256:} {\scriptsize\nolinkurl{620bdf1b1eaca5bc5f86831a95fd342e72fe43e7fb697fd0590d56a03cc92533}}.

\VerbatimInput[fontsize=\scriptsize,breaklines=true,breakanywhere=true]{excerpts/D_CliffordPacket.txt}

\textbf{Direct project imports:} \nolinkurl{PaperXI.CliffordSource}, \nolinkurl{PaperXI.TwistedActionEquivalence}.

\begin{longtable}{@{}p{0.07\textwidth}p{0.12\textwidth}p{0.75\textwidth}@{}}\toprule Line&Kind&Name\\\midrule\endhead

20&def&\nolinkurl{PaperXI.CliffordPacket.action}\\

27&theorem&\nolinkurl{PaperXI.CliffordPacket.action_basis}\\

38&theorem&\nolinkurl{PaperXI.CliffordPacket.action_bijective}\\

44&def&\nolinkurl{PaperXI.CliffordPacket.actionEquiv}\\

49&theorem&\nolinkurl{PaperXI.CliffordPacket.actionEquiv_apply}\\

55&theorem&\nolinkurl{PaperXI.CliffordPacket.finrank_sq}\\

65&theorem&\nolinkurl{PaperXI.CliffordPacket.finrank_sqrt}\\

\bottomrule\end{longtable}

\section{\texorpdfstring{\nolinkurl{PaperXI.CliffordSource}}{PaperXI.CliffordSource}}\label{module:PaperXI.CliffordSource}

\textbf{Selected layer:} v0.44 parent.\quad\textbf{Source lines:} 48.\par\textbf{SHA-256:} {\scriptsize\nolinkurl{f0b150418de3d56a1fcd7799169a9937e6b9746b2dc6a75c79de625bde828c88}}.

\VerbatimInput[fontsize=\scriptsize,breaklines=true,breakanywhere=true]{excerpts/D_CliffordSource.txt}

\textbf{Direct project imports:} \nolinkurl{PaperXI.CliffordIrreducible}, \nolinkurl{PaperXI.CliffordCocycle}.

\begin{longtable}{@{}p{0.07\textwidth}p{0.12\textwidth}p{0.75\textwidth}@{}}\toprule Line&Kind&Name\\\midrule\endhead

20&def&\nolinkurl{PaperXI.CliffordSource.cocycle}\\

26&theorem&\nolinkurl{PaperXI.CliffordSource.normalized}\\

32&theorem&\nolinkurl{PaperXI.CliffordSource.operator_mul}\\

42&theorem&\nolinkurl{PaperXI.CliffordSource.projective_irreducible}\\

\bottomrule\end{longtable}

\section{\texorpdfstring{\nolinkurl{PaperXI.CliffordStabilizer}}{PaperXI.CliffordStabilizer}}\label{module:PaperXI.CliffordStabilizer}

\textbf{Selected layer:} v0.44 parent.\quad\textbf{Source lines:} 216.\par\textbf{SHA-256:} {\scriptsize\nolinkurl{6bc6cae8489f5025d88e5dfe8375db8a056c899365557a772e5188998ff689fc}}.

\VerbatimInput[fontsize=\scriptsize,breaklines=true,breakanywhere=true]{excerpts/D_CliffordStabilizer.txt}

\textbf{Direct project imports:} \nolinkurl{PaperXI.CliffordDatum}, \nolinkurl{PaperXI.CliffordSource}, \nolinkurl{PaperXI.CliffordEnvelope}, \nolinkurl{PaperXI.CliffordTransport}, \nolinkurl{PaperXI.ElementaryEnvelopeTransport}.

\begin{longtable}{@{}p{0.07\textwidth}p{0.12\textwidth}p{0.75\textwidth}@{}}\toprule Line&Kind&Name\\\midrule\endhead

20&def&\nolinkurl{PaperXI.CliffordStabilizer.quotientAction}\\

36&theorem&\nolinkurl{PaperXI.CliffordStabilizer.quotientAction_mk}\\

41&theorem&\nolinkurl{PaperXI.CliffordStabilizer.quotientAction_comp}\\

47&theorem&\nolinkurl{PaperXI.CliffordStabilizer.quotientAction_unique}\\

67&def&\nolinkurl{PaperXI.CliffordStabilizer.envelopeAction}\\

70&theorem&\nolinkurl{PaperXI.CliffordStabilizer.envelopeAction_mk}\\

76&theorem&\nolinkurl{PaperXI.CliffordStabilizer.envelopeAction_form}\\

89&theorem&\nolinkurl{PaperXI.CliffordStabilizer.envelopeAction_equivariant}\\

137&theorem&\nolinkurl{PaperXI.CliffordStabilizer.source_commutator_preserved}\\

151&def&\nolinkurl{PaperXI.CliffordStabilizer.sourceAction}\\

157&theorem&\nolinkurl{PaperXI.CliffordStabilizer.sourceAction_mk}\\

163&theorem&\nolinkurl{PaperXI.CliffordStabilizer.sourceAction_form}\\

174&theorem&\nolinkurl{PaperXI.CliffordStabilizer.sourceAction_independent}\\

200&theorem&\nolinkurl{PaperXI.CliffordStabilizer.sourceAction_equivariant}\\

\bottomrule\end{longtable}

\section{\texorpdfstring{\nolinkurl{PaperXI.CliffordTransport}}{PaperXI.CliffordTransport}}\label{module:PaperXI.CliffordTransport}

\textbf{Selected layer:} v0.44 parent.\quad\textbf{Source lines:} 280.\par\textbf{SHA-256:} {\scriptsize\nolinkurl{3b10aa513a5ba4f75eee3ca0e288cb4683b0a0044f3c22cf1bfaf031a568a995}}.

\VerbatimInput[fontsize=\scriptsize,breaklines=true,breakanywhere=true]{excerpts/D_CliffordTransport.txt}

\textbf{Direct project imports:} \nolinkurl{PaperXI.CliffordCocycle}.

\begin{longtable}{@{}p{0.07\textwidth}p{0.12\textwidth}p{0.75\textwidth}@{}}\toprule Line&Kind&Name\\\midrule\endhead

15&def&\nolinkurl{PaperXI.CliffordTransport.subgroupEquiv}\\

26&theorem&\nolinkurl{PaperXI.CliffordTransport.subgroupEquiv_coe}\\

30&theorem&\nolinkurl{PaperXI.CliffordTransport.subgroupEquiv_conjugation}\\

39&def&\nolinkurl{PaperXI.CliffordTransport.quotientEquiv}\\

43&theorem&\nolinkurl{PaperXI.CliffordTransport.quotientEquiv_mk}\\

50&theorem&\nolinkurl{PaperXI.CliffordTransport.quotientEquiv_refl}\\

59&theorem&\nolinkurl{PaperXI.CliffordTransport.quotientEquiv_comp}\\

76&theorem&\nolinkurl{PaperXI.CliffordTransport.source_apply}\\

80&theorem&\nolinkurl{PaperXI.CliffordTransport.source_inverse_apply}\\

88&theorem&\nolinkurl{PaperXI.CliffordTransport.target_apply}\\

91&theorem&\nolinkurl{PaperXI.CliffordTransport.target_inverse_apply}\\

99&def&\nolinkurl{PaperXI.CliffordTransport.homMap}\\

113&theorem&\nolinkurl{PaperXI.CliffordTransport.homMap_apply}\\

117&def&\nolinkurl{PaperXI.CliffordTransport.inverseHomMap}\\

130&def&\nolinkurl{PaperXI.CliffordTransport.homEquiv}\\

142&theorem&\nolinkurl{PaperXI.CliffordTransport.homEquiv_apply}\\

145&theorem&\nolinkurl{PaperXI.CliffordTransport.homEquiv_symm_apply}\\

151&def&\nolinkurl{PaperXI.CliffordTransport.transportedIntertwiner}\\

162&theorem&\nolinkurl{PaperXI.CliffordTransport.transportedIntertwiner_apply}\\

167&theorem&\nolinkurl{PaperXI.CliffordTransport.operator_transport}\\

178&theorem&\nolinkurl{PaperXI.CliffordTransport.operator_transport_conj}\\

189&theorem&\nolinkurl{PaperXI.CliffordTransport.transported_representative}\\

198&def&\nolinkurl{PaperXI.CliffordTransport.cochain}\\

206&theorem&\nolinkurl{PaperXI.CliffordTransport.sectionOperator_transport}\\

230&theorem&\nolinkurl{PaperXI.CliffordTransport.cochain_one}\\

243&theorem&\nolinkurl{PaperXI.CliffordTransport.cocycle_transport}\\

265&theorem&\nolinkurl{PaperXI.CliffordTransport.commutator_preserved}\\

\bottomrule\end{longtable}

\section{\texorpdfstring{\nolinkurl{PaperXI.CocycleBasis}}{PaperXI.CocycleBasis}}\label{module:PaperXI.CocycleBasis}

\textbf{Selected layer:} v0.44 parent.\quad\textbf{Source lines:} 33.\par\textbf{SHA-256:} {\scriptsize\nolinkurl{5ba62f3c665922695e62725e35c969a9b223859a7dee0cad1163de238e51213e}}.

\VerbatimInput[fontsize=\scriptsize,breaklines=true,breakanywhere=true]{excerpts/D_CocycleBasis.txt}

\textbf{Direct project imports:} \nolinkurl{PaperXI.CocycleCommutator}.

\begin{longtable}{@{}p{0.07\textwidth}p{0.12\textwidth}p{0.75\textwidth}@{}}\toprule Line&Kind&Name\\\midrule\endhead

15&theorem&\nolinkurl{PaperXI.CocycleCommutator.rescaled_basis_mul}\\

29&theorem&\nolinkurl{PaperXI.CocycleCommutator.rescaled_basis_commutator}\\

\bottomrule\end{longtable}

\section{\texorpdfstring{\nolinkurl{PaperXI.CocycleClass}}{PaperXI.CocycleClass}}\label{module:PaperXI.CocycleClass}

\textbf{Selected layer:} v0.44 parent.\quad\textbf{Source lines:} 168.\par\textbf{SHA-256:} {\scriptsize\nolinkurl{393677b19c265df8ad023c34ede0aa8a8cb38f60e1bb050ef0b4282876a302bd}}.

\VerbatimInput[fontsize=\scriptsize,breaklines=true,breakanywhere=true]{excerpts/D_CocycleClass.txt}

\textbf{Direct project imports:} \nolinkurl{PaperXI.CocycleCommutator}.

\begin{longtable}{@{}p{0.07\textwidth}p{0.12\textwidth}p{0.75\textwidth}@{}}\toprule Line&Kind&Name\\\midrule\endhead

14&theorem&\nolinkurl{PaperXI.CocycleClass.cocycle_ext}\\

21&theorem&\nolinkurl{PaperXI.CocycleClass.rescale_one}\\

27&theorem&\nolinkurl{PaperXI.CocycleClass.rescale_trans}\\

34&abbrev&\nolinkurl{PaperXI.CocycleClass.NormalizedCocycle}\\

38&def&\nolinkurl{PaperXI.CocycleClass.Related}\\

41&theorem&\nolinkurl{PaperXI.CocycleClass.related_refl}\\

44&theorem&\nolinkurl{PaperXI.CocycleClass.related_symm}\\

50&theorem&\nolinkurl{PaperXI.CocycleClass.related_trans}\\

58&def&\nolinkurl{PaperXI.CocycleClass.coboundarySetoid}\\

64&def&\nolinkurl{PaperXI.CocycleClass.Class}\\

67&def&\nolinkurl{PaperXI.CocycleClass.of}\\

71&theorem&\nolinkurl{PaperXI.CocycleClass.of_eq_iff}\\

76&theorem&\nolinkurl{PaperXI.CocycleClass.of_eq_of_rescale}\\

80&theorem&\nolinkurl{PaperXI.CocycleClass.of_rescale}\\

86&def&\nolinkurl{PaperXI.CocycleClass.commutator}\\

95&theorem&\nolinkurl{PaperXI.CocycleClass.commutator_of}\\

99&def&\nolinkurl{PaperXI.CocycleClass.radical}\\

107&theorem&\nolinkurl{PaperXI.CocycleClass.radical_of}\\

111&def&\nolinkurl{PaperXI.CocycleClass.pullbackCocycle}\\

115&theorem&\nolinkurl{PaperXI.CocycleClass.normalized_pullback}\\

119&theorem&\nolinkurl{PaperXI.CocycleClass.pullback_rescale}\\

126&def&\nolinkurl{PaperXI.CocycleClass.pullback}\\

136&theorem&\nolinkurl{PaperXI.CocycleClass.pullback_of}\\

140&def&\nolinkurl{PaperXI.CocycleClass.congrEquiv}\\

164&theorem&\nolinkurl{PaperXI.CocycleClass.congrEquiv_of}\\

\bottomrule\end{longtable}

\section{\texorpdfstring{\nolinkurl{PaperXI.CocycleCommutator}}{PaperXI.CocycleCommutator}}\label{module:PaperXI.CocycleCommutator}

\textbf{Selected layer:} v0.44 parent.\quad\textbf{Source lines:} 122.\par\textbf{SHA-256:} {\scriptsize\nolinkurl{0b8b7f046bf6d69122f05e183848102036057443192762a7e89fa4e374f6d4df}}.

\VerbatimInput[fontsize=\scriptsize,breaklines=true,breakanywhere=true]{excerpts/D_CocycleCommutator.txt}

\textbf{Direct project imports:} none.

\begin{longtable}{@{}p{0.07\textwidth}p{0.12\textwidth}p{0.75\textwidth}@{}}\toprule Line&Kind&Name\\\midrule\endhead

14&structure&\nolinkurl{PaperXI.CocycleCommutator.Cocycle}\\

18&def&\nolinkurl{PaperXI.CocycleCommutator.Normalized}\\

22&def&\nolinkurl{PaperXI.CocycleCommutator.rescale}\\

35&def&\nolinkurl{PaperXI.CocycleCommutator.commutator}\\

37&theorem&\nolinkurl{PaperXI.CocycleCommutator.commutator_self}\\

40&theorem&\nolinkurl{PaperXI.CocycleCommutator.commutator_swap}\\

43&theorem&\nolinkurl{PaperXI.CocycleCommutator.commutator_rescale}\\

48&theorem&\nolinkurl{PaperXI.CocycleCommutator.commutator_mul_left}\\

68&theorem&\nolinkurl{PaperXI.CocycleCommutator.commutator_mul_right}\\

73&theorem&\nolinkurl{PaperXI.CocycleCommutator.commutator_one_left}\\

78&theorem&\nolinkurl{PaperXI.CocycleCommutator.commutator_one_right}\\

82&def&\nolinkurl{PaperXI.CocycleCommutator.characterMap}\\

90&theorem&\nolinkurl{PaperXI.CocycleCommutator.characterMap_apply}\\

94&def&\nolinkurl{PaperXI.CocycleCommutator.radical}\\

96&theorem&\nolinkurl{PaperXI.CocycleCommutator.mem_radical_iff}\\

103&theorem&\nolinkurl{PaperXI.CocycleCommutator.radical_rescale}\\

110&theorem&\nolinkurl{PaperXI.CocycleCommutator.normalized_rescale}\\

114&theorem&\nolinkurl{PaperXI.CocycleCommutator.normalized_rescale_iff}\\

\bottomrule\end{longtable}

\section{\texorpdfstring{\nolinkurl{PaperXI.CocycleQuotient}}{PaperXI.CocycleQuotient}}\label{module:PaperXI.CocycleQuotient}

\textbf{Selected layer:} v0.44 parent.\quad\textbf{Source lines:} 85.\par\textbf{SHA-256:} {\scriptsize\nolinkurl{e4fd2024e7a2b19c99fe508d778300fc161a53761c5ed2db8f0e5c0091101b51}}.

\VerbatimInput[fontsize=\scriptsize,breaklines=true,breakanywhere=true]{excerpts/D_CocycleQuotient.txt}

\textbf{Direct project imports:} \nolinkurl{PaperXI.CocycleCommutator}.

\begin{longtable}{@{}p{0.07\textwidth}p{0.12\textwidth}p{0.75\textwidth}@{}}\toprule Line&Kind&Name\\\midrule\endhead

11&def&\nolinkurl{PaperXI.CocycleCommutator.quotientCharacter}\\

14&theorem&\nolinkurl{PaperXI.CocycleCommutator.quotientCharacter_mk}\\

17&theorem&\nolinkurl{PaperXI.CocycleCommutator.radical_le_quotientCharacter_ker}\\

25&def&\nolinkurl{PaperXI.CocycleCommutator.quotientPairing}\\

43&theorem&\nolinkurl{PaperXI.CocycleCommutator.quotientPairing_mk}\\

47&theorem&\nolinkurl{PaperXI.CocycleCommutator.quotientPairing_self}\\

52&theorem&\nolinkurl{PaperXI.CocycleCommutator.quotientPairing_swap}\\

60&theorem&\nolinkurl{PaperXI.CocycleCommutator.quotientPairing_nondegenerate_left}\\

74&theorem&\nolinkurl{PaperXI.CocycleCommutator.quotientPairing_nondegenerate_right}\\

82&theorem&\nolinkurl{PaperXI.CocycleCommutator.quotientCharacter_injective}\\

\bottomrule\end{longtable}

\section{\texorpdfstring{\nolinkurl{PaperXI.CommutatorEnvelope}}{PaperXI.CommutatorEnvelope}}\label{module:PaperXI.CommutatorEnvelope}

\textbf{Selected layer:} v0.44 parent.\quad\textbf{Source lines:} 160.\par\textbf{SHA-256:} {\scriptsize\nolinkurl{bc6e6fc84b1c63fc1bb70568bc6eb1c3137c5f9f781783990c9c1bc51a6453b8}}.

\VerbatimInput[fontsize=\scriptsize,breaklines=true,breakanywhere=true]{excerpts/D_CommutatorEnvelope.txt}

\textbf{Direct project imports:} \nolinkurl{PaperXI.CocycleQuotient}.

\begin{longtable}{@{}p{0.07\textwidth}p{0.12\textwidth}p{0.75\textwidth}@{}}\toprule Line&Kind&Name\\\midrule\endhead

14&def&\nolinkurl{PaperXI.CommutatorEnvelope.descendedCharacter}\\

17&theorem&\nolinkurl{PaperXI.CommutatorEnvelope.descendedCharacter_mk}\\

21&theorem&\nolinkurl{PaperXI.CommutatorEnvelope.le_descendedCharacter_ker}\\

29&def&\nolinkurl{PaperXI.CommutatorEnvelope.descendedPairing}\\

48&theorem&\nolinkurl{PaperXI.CommutatorEnvelope.descendedPairing_mk}\\

54&theorem&\nolinkurl{PaperXI.CommutatorEnvelope.descends_iff}\\

68&theorem&\nolinkurl{PaperXI.CommutatorEnvelope.descendedPairing_unique}\\

81&theorem&\nolinkurl{PaperXI.CommutatorEnvelope.existsUnique_descendedPairing}\\

89&theorem&\nolinkurl{PaperXI.CommutatorEnvelope.descendedPairing_self}\\

95&theorem&\nolinkurl{PaperXI.CommutatorEnvelope.nondegenerate_iff}\\

116&def&\nolinkurl{PaperXI.CommutatorEnvelope.quotientFactor}\\

120&theorem&\nolinkurl{PaperXI.CommutatorEnvelope.quotientFactor_mk}\\

125&theorem&\nolinkurl{PaperXI.CommutatorEnvelope.quotientFactor_surjective}\\

132&theorem&\nolinkurl{PaperXI.CommutatorEnvelope.quotientFactor_pairing}\\

140&theorem&\nolinkurl{PaperXI.CommutatorEnvelope.existsUnique_quotientFactor}\\

151&theorem&\nolinkurl{PaperXI.CommutatorEnvelope.existsUnique_injective_characterFactor}\\

\bottomrule\end{longtable}

\section{\texorpdfstring{\nolinkurl{PaperXI.CompactDegreeSeparation}}{PaperXI.CompactDegreeSeparation}}\label{module:PaperXI.CompactDegreeSeparation}

\textbf{Selected layer:} base108 arithmetic branch.\quad\textbf{Source lines:} 37.\par\textbf{SHA-256:} {\scriptsize\nolinkurl{c69f0838adb9db017a926ff1a66606634858bfb0aa89710e147b19b44a60d561}}.

\VerbatimInput[fontsize=\scriptsize,breaklines=true,breakanywhere=true]{excerpts/D_CompactDegreeSeparation.txt}

\textbf{Direct project imports:} \nolinkurl{PaperXI.CentralExtensionInducing}.

\begin{longtable}{@{}p{0.07\textwidth}p{0.12\textwidth}p{0.75\textwidth}@{}}\toprule Line&Kind&Name\\\midrule\endhead

20&theorem&\nolinkurl{PaperXI.CompactDegreeSeparation.compact_le_degree_ker}\\

32&theorem&\nolinkurl{PaperXI.CompactDegreeSeparation.hfree_from_compact_degree}\\

\bottomrule\end{longtable}

\section{\texorpdfstring{\nolinkurl{PaperXI.CompactIntersection108}}{PaperXI.CompactIntersection108}}\label{module:PaperXI.CompactIntersection108}

\textbf{Selected layer:} base108 arithmetic branch.\quad\textbf{Source lines:} 55.\par\textbf{SHA-256:} {\scriptsize\nolinkurl{318e23f5dfa82f3cf42743e82c7783ac44508ce10be1a856aa6fc60407b2fa37}}.

\VerbatimInput[fontsize=\scriptsize,breaklines=true,breakanywhere=true]{excerpts/D_CompactIntersection108.txt}

\textbf{Direct project imports:} \nolinkurl{PaperXI.IntersectionModel108}, \nolinkurl{PaperXI.FiniteJointQuotient108}.

\begin{longtable}{@{}p{0.07\textwidth}p{0.12\textwidth}p{0.75\textwidth}@{}}\toprule Line&Kind&Name\\\midrule\endhead

17&theorem&\nolinkurl{PaperXI.CompactIntersection108.compact_pairs}\\

26&theorem&\nolinkurl{PaperXI.CompactIntersection108.smooth_comp}\\

35&theorem&\nolinkurl{PaperXI.CompactIntersection108.finite_quotient}\\

46&theorem&\nolinkurl{PaperXI.CompactIntersection108.finrank_quotient}\\

\bottomrule\end{longtable}

\section{\texorpdfstring{\nolinkurl{PaperXI.CompactModCenterInduction68}}{PaperXI.CompactModCenterInduction68}}\label{module:PaperXI.CompactModCenterInduction68}

\textbf{Selected layer:} base108 arithmetic branch.\quad\textbf{Source lines:} 126.\par\textbf{SHA-256:} {\scriptsize\nolinkurl{b1765821b4681b903bdbaabe6a76b05200fba99f81d94cc260e9ecc3484a4065}}.

\VerbatimInput[fontsize=\scriptsize,breaklines=true,breakanywhere=true]{excerpts/D_CompactModCenterInduction68.txt}

\textbf{Direct project imports:} \nolinkurl{PaperXI.EquivariantSplitting68}, \nolinkurl{PaperXI.InducedTypeSpace68}.

\begin{longtable}{@{}p{0.07\textwidth}p{0.12\textwidth}p{0.75\textwidth}@{}}\toprule Line&Kind&Name\\\midrule\endhead

25&theorem&\nolinkurl{PaperXI.CompactModCenterInduction68.induced_central_scalar}\\

40&def&\nolinkurl{PaperXI.CompactModCenterInduction68.subInclusion}\\

52&theorem&\nolinkurl{PaperXI.CompactModCenterInduction68.irreducible_of_compact_mod_center_mackey}\\

111&theorem&\nolinkurl{PaperXI.CompactModCenterInduction68.induction_package}\\

\bottomrule\end{longtable}

\section{\texorpdfstring{\nolinkurl{PaperXI.CompactQuotientPeriod69}}{PaperXI.CompactQuotientPeriod69}}\label{module:PaperXI.CompactQuotientPeriod69}

\textbf{Selected layer:} base108 arithmetic branch.\quad\textbf{Source lines:} 52.\par\textbf{SHA-256:} {\scriptsize\nolinkurl{db609bc27e26e3964e4eaa9f6739b066d277b23dd9c0b595f4b51fc80baa975a}}.

\VerbatimInput[fontsize=\scriptsize,breaklines=true,breakanywhere=true]{excerpts/D_CompactQuotientPeriod69.txt}

\textbf{Direct project imports:} none.

\begin{longtable}{@{}p{0.07\textwidth}p{0.12\textwidth}p{0.75\textwidth}@{}}\toprule Line&Kind&Name\\\midrule\endhead

15&theorem&\nolinkurl{PaperXI.CompactQuotientPeriod69.compact_quotient_of_period}\\

\bottomrule\end{longtable}

\section{\texorpdfstring{\nolinkurl{PaperXI.CompactSmoothSemisimple}}{PaperXI.CompactSmoothSemisimple}}\label{module:PaperXI.CompactSmoothSemisimple}

\textbf{Selected layer:} base108 arithmetic branch.\quad\textbf{Source lines:} 92.\par\textbf{SHA-256:} {\scriptsize\nolinkurl{8f8329834876b5913b22197a22cc58d386a26429cb21fcb23d51947c33d93e74}}.

\VerbatimInput[fontsize=\scriptsize,breaklines=true,breakanywhere=true]{excerpts/D_CompactSmoothSemisimple.txt}

\textbf{Direct project imports:} \nolinkurl{PaperXI.SmoothDegreeQuotient}, \nolinkurl{PaperXI.LocallyFiniteSemisimple}.

\begin{longtable}{@{}p{0.07\textwidth}p{0.12\textwidth}p{0.75\textwidth}@{}}\toprule Line&Kind&Name\\\midrule\endhead

19&def&\nolinkurl{PaperXI.CompactSmoothSemisimple.orbitQuotient}\\

27&def&\nolinkurl{PaperXI.CompactSmoothSemisimple.orbitSpan}\\

40&theorem&\nolinkurl{PaperXI.CompactSmoothSemisimple.mem_orbitSpan}\\

48&theorem&\nolinkurl{PaperXI.CompactSmoothSemisimple.subrepresentation_smooth}\\

61&theorem&\nolinkurl{PaperXI.CompactSmoothSemisimple.orbit_finite}\\

70&theorem&\nolinkurl{PaperXI.CompactSmoothSemisimple.orbitSpan_finiteDimensional}\\

75&theorem&\nolinkurl{PaperXI.CompactSmoothSemisimple.subrepresentation_kernel_quotient_finite}\\

83&theorem&\nolinkurl{PaperXI.CompactSmoothSemisimple.smooth_semisimple}\\

\bottomrule\end{longtable}

\section{\texorpdfstring{\nolinkurl{PaperXI.CompactTypeTensor58}}{PaperXI.CompactTypeTensor58}}\label{module:PaperXI.CompactTypeTensor58}

\textbf{Selected layer:} base108 arithmetic branch.\quad\textbf{Source lines:} 93.\par\textbf{SHA-256:} {\scriptsize\nolinkurl{885a6eef6b9a66f86ddf833e1077fdd3de5fe94a004487d6b539d70cba73bc6b}}.

\VerbatimInput[fontsize=\scriptsize,breaklines=true,breakanywhere=true]{excerpts/D_CompactTypeTensor58.txt}

\textbf{Direct project imports:} \nolinkurl{PaperXI.GallagherTensorIrreducibility58}, \nolinkurl{PaperXI.SmoothDegreeQuotient}, \nolinkurl{PaperXI.FiniteRepresentation}.

\begin{longtable}{@{}p{0.07\textwidth}p{0.12\textwidth}p{0.75\textwidth}@{}}\toprule Line&Kind&Name\\\midrule\endhead

27&abbrev&\nolinkurl{PaperXI.CompactTypeTensor58.action}\\

32&theorem&\nolinkurl{PaperXI.CompactTypeTensor58.restriction}\\

39&theorem&\nolinkurl{PaperXI.CompactTypeTensor58.subgroup_action}\\

48&theorem&\nolinkurl{PaperXI.CompactTypeTensor58.irreducible}\\

55&theorem&\nolinkurl{PaperXI.CompactTypeTensor58.finiteDimensional}\\

66&theorem&\nolinkurl{PaperXI.CompactTypeTensor58.smooth}\\

82&theorem&\nolinkurl{PaperXI.CompactTypeTensor58.completed}\\

\bottomrule\end{longtable}

\section{\texorpdfstring{\nolinkurl{PaperXI.CompactTypeTensorNormalizer58}}{PaperXI.CompactTypeTensorNormalizer58}}\label{module:PaperXI.CompactTypeTensorNormalizer58}

\textbf{Selected layer:} base108 arithmetic branch.\quad\textbf{Source lines:} 54.\par\textbf{SHA-256:} {\scriptsize\nolinkurl{1058ac054e5f073f587c3012d307d04d2924924594cd5ec137685c0bfb3e10b5}}.

\VerbatimInput[fontsize=\scriptsize,breaklines=true,breakanywhere=true]{excerpts/D_CompactTypeTensorNormalizer58.txt}

\textbf{Direct project imports:} \nolinkurl{PaperXI.CompactTypeTensor58}.

\begin{longtable}{@{}p{0.07\textwidth}p{0.12\textwidth}p{0.75\textwidth}@{}}\toprule Line&Kind&Name\\\midrule\endhead

26&def&\nolinkurl{PaperXI.CompactTypeTensorNormalizer58.tensorOperator}\\

29&theorem&\nolinkurl{PaperXI.CompactTypeTensorNormalizer58.tensorOperator_tmul}\\

34&theorem&\nolinkurl{PaperXI.CompactTypeTensorNormalizer58.tensorOperator_covariance}\\

43&def&\nolinkurl{PaperXI.CompactTypeTensorNormalizer58.normalizerEquiv}\\

51&theorem&\nolinkurl{PaperXI.CompactTypeTensorNormalizer58.normalizerEquiv_linear}\\

\bottomrule\end{longtable}

\section{\texorpdfstring{\nolinkurl{PaperXI.ConditionalAtomicComparison}}{PaperXI.ConditionalAtomicComparison}}\label{module:PaperXI.ConditionalAtomicComparison}

\textbf{Selected layer:} v0.44 parent.\quad\textbf{Source lines:} 151.\par\textbf{SHA-256:} {\scriptsize\nolinkurl{ce8d1ea0f8658d4d968aeb6665d5f1d3e4cc6bb1d472ebfbfbce6215f1103a58}}.

\VerbatimInput[fontsize=\scriptsize,breaklines=true,breakanywhere=true]{excerpts/D_ConditionalAtomicComparison.txt}

\textbf{Direct project imports:} \nolinkurl{PaperXI.FiniteBreakAssembly}, \nolinkurl{PaperXI.GluedPacketSections}, \nolinkurl{PaperXI.LocalFactorIdentification}.

\begin{longtable}{@{}p{0.07\textwidth}p{0.12\textwidth}p{0.75\textwidth}@{}}\toprule Line&Kind&Name\\\midrule\endhead

19&structure&\nolinkurl{PaperXI.ConditionalAtomicComparison.Diagram}\\

27&structure&\nolinkurl{PaperXI.ConditionalAtomicComparison.Atlas}\\

42&def&\nolinkurl{PaperXI.ConditionalAtomicComparison.targetDiagram}\\

45&def&\nolinkurl{PaperXI.ConditionalAtomicComparison.sourceDiagram}\\

49&def&\nolinkurl{PaperXI.ConditionalAtomicComparison.atRoot}\\

56&theorem&\nolinkurl{PaperXI.ConditionalAtomicComparison.atRoot_linear}\\

60&def&\nolinkurl{PaperXI.ConditionalAtomicComparison.Holonomy}\\

65&def&\nolinkurl{PaperXI.ConditionalAtomicComparison.leftSections}\\

67&def&\nolinkurl{PaperXI.ConditionalAtomicComparison.rightSections}\\

70&def&\nolinkurl{PaperXI.ConditionalAtomicComparison.glue}\\

75&theorem&\nolinkurl{PaperXI.ConditionalAtomicComparison.glue_apply}\\

84&theorem&\nolinkurl{PaperXI.ConditionalAtomicComparison.holonomy_iff_gauge}\\

119&def&\nolinkurl{PaperXI.ConditionalAtomicComparison.finiteComparison}\\

127&theorem&\nolinkurl{PaperXI.ConditionalAtomicComparison.source}\\

\bottomrule\end{longtable}

\section{\texorpdfstring{\nolinkurl{PaperXI.ConstructedFieldFrobenius}}{PaperXI.ConstructedFieldFrobenius}}\label{module:PaperXI.ConstructedFieldFrobenius}

\textbf{Selected layer:} normalized Frobenius.\quad\textbf{Source lines:} 65.\par\textbf{SHA-256:} {\scriptsize\nolinkurl{37abb8f28a5e22530d0f0296eaef1880e5fca49629014ef7cd762d606916b5bc}}.

\VerbatimInput[fontsize=\scriptsize,breaklines=true,breakanywhere=true]{excerpts/D_ConstructedFieldFrobenius.txt}

\textbf{Direct project imports:} \nolinkurl{PaperXI.SpectralGeometricFrobenius}, \nolinkurl{PaperXI.SpectralLocalField}.

\begin{longtable}{@{}p{0.07\textwidth}p{0.12\textwidth}p{0.75\textwidth}@{}}\toprule Line&Kind&Name\\\midrule\endhead

22&abbrev&\nolinkurl{PaperXI.ConstructedFieldFrobenius.BaseField}\\

28&theorem&\nolinkurl{PaperXI.ConstructedFieldFrobenius.restricted_integer_subring}\\

36&def&\nolinkurl{PaperXI.ConstructedFieldFrobenius.baseIntegerEquiv}\\

43&def&\nolinkurl{PaperXI.ConstructedFieldFrobenius.baseResidueEquiv}\\

46&theorem&\nolinkurl{PaperXI.ConstructedFieldFrobenius.baseResidue_card}\\

54&theorem&\nolinkurl{PaperXI.ConstructedFieldFrobenius.exists_normalized_field_automorphism}\\

\bottomrule\end{longtable}

\section{\texorpdfstring{\nolinkurl{PaperXI.ContinuousHolonomyResidual}}{PaperXI.ContinuousHolonomyResidual}}\label{module:PaperXI.ContinuousHolonomyResidual}

\textbf{Selected layer:} v0.44 parent.\quad\textbf{Source lines:} 132.\par\textbf{SHA-256:} {\scriptsize\nolinkurl{ef0ef5f7cdf0b952317f838c65b6f9e0229b32e424b9f384ab96634cd955f060}}.

\VerbatimInput[fontsize=\scriptsize,breaklines=true,breakanywhere=true]{excerpts/D_ContinuousHolonomyResidual.txt}

\textbf{Direct project imports:} \nolinkurl{PaperXI.HolonomyResidual}.

\begin{longtable}{@{}p{0.07\textwidth}p{0.12\textwidth}p{0.75\textwidth}@{}}\toprule Line&Kind&Name\\\midrule\endhead

18&def&\nolinkurl{PaperXI.ContinuousHolonomyResidual.transported}\\

21&theorem&\nolinkurl{PaperXI.ContinuousHolonomyResidual.transported_apply}\\

24&theorem&\nolinkurl{PaperXI.ContinuousHolonomyResidual.transported_toLinearEquiv}\\

28&def&\nolinkurl{PaperXI.ContinuousHolonomyResidual.cycleDefect}\\

37&def&\nolinkurl{PaperXI.ContinuousHolonomyResidual.edgeResidual}\\

44&theorem&\nolinkurl{PaperXI.ContinuousHolonomyResidual.cycleDefect_toLinearMap}\\

52&theorem&\nolinkurl{PaperXI.ContinuousHolonomyResidual.edgeResidual_toLinearMap}\\

60&theorem&\nolinkurl{PaperXI.ContinuousHolonomyResidual.residual_factorization}\\

72&theorem&\nolinkurl{PaperXI.ContinuousHolonomyResidual.residual_zero_iff}\\

86&theorem&\nolinkurl{PaperXI.ContinuousHolonomyResidual.residual_norm_bound}\\

103&theorem&\nolinkurl{PaperXI.ContinuousHolonomyResidual.uniform_bound}\\

119&theorem&\nolinkurl{PaperXI.ContinuousHolonomyResidual.product_conditioning_bound}\\

\bottomrule\end{longtable}

\section{\texorpdfstring{\nolinkurl{PaperXI.CoordinateDeepIdeal111}}{PaperXI.CoordinateDeepIdeal111}}\label{module:PaperXI.CoordinateDeepIdeal111}

\textbf{Selected layer:} Continuation111.\quad\textbf{Source lines:} 79.\par\textbf{SHA-256:} {\scriptsize\nolinkurl{00a4e4aef33f78b2682431bc4370b2158cb1628f7d1c702f634861314e19cb09}}.

\VerbatimInput[fontsize=\scriptsize,breaklines=true,breakanywhere=true]{excerpts/D_CoordinateDeepIdeal111.txt}

\textbf{Direct project imports:} \nolinkurl{PaperXI.OnePlusIntersection110}.

\begin{longtable}{@{}p{0.07\textwidth}p{0.12\textwidth}p{0.75\textwidth}@{}}\toprule Line&Kind&Name\\\midrule\endhead

21&def&\nolinkurl{PaperXI.CoordinateDeepIdeal111.uncurryHomeomorph}\\

31&theorem&\nolinkurl{PaperXI.CoordinateDeepIdeal111.exists_open_matrix_ideal_subset}\\

45&theorem&\nolinkurl{PaperXI.CoordinateDeepIdeal111.matrix_ideal_open}\\

64&theorem&\nolinkurl{PaperXI.CoordinateDeepIdeal111.exists_common_ideal_finite}\\

\bottomrule\end{longtable}

\section{\texorpdfstring{\nolinkurl{PaperXI.CuspidalAnnihilation}}{PaperXI.CuspidalAnnihilation}}\label{module:PaperXI.CuspidalAnnihilation}

\textbf{Selected layer:} v0.44 parent.\quad\textbf{Source lines:} 87.\par\textbf{SHA-256:} {\scriptsize\nolinkurl{e56bdeb22b0d449550cd009852c029feed4d614bdc02a53d49648784b4a6abd2}}.

\VerbatimInput[fontsize=\scriptsize,breaklines=true,breakanywhere=true]{excerpts/D_CuspidalAnnihilation.txt}

\textbf{Direct project imports:} \nolinkurl{PaperXI.PrimitiveCharacterBlock}, \nolinkurl{PaperXI.FiniteAverage}.

\begin{longtable}{@{}p{0.07\textwidth}p{0.12\textwidth}p{0.75\textwidth}@{}}\toprule Line&Kind&Name\\\midrule\endhead

16&def&\nolinkurl{PaperXI.CuspidalAnnihilation.subgroupAverage}\\

20&theorem&\nolinkurl{PaperXI.CuspidalAnnihilation.average_action}\\

25&theorem&\nolinkurl{PaperXI.CuspidalAnnihilation.average_range}\\

31&theorem&\nolinkurl{PaperXI.CuspidalAnnihilation.average_idempotent}\\

38&theorem&\nolinkurl{PaperXI.CuspidalAnnihilation.average_zero_iff}\\

45&theorem&\nolinkurl{PaperXI.CuspidalAnnihilation.own_block_action}\\

51&theorem&\nolinkurl{PaperXI.CuspidalAnnihilation.left_annihilation_iff}\\

66&theorem&\nolinkurl{PaperXI.CuspidalAnnihilation.right_annihilation_iff}\\

73&theorem&\nolinkurl{PaperXI.CuspidalAnnihilation.parabolic_annihilation}\\

79&theorem&\nolinkurl{PaperXI.CuspidalAnnihilation.nonzero_of_fixed_vector}\\

\bottomrule\end{longtable}

\section{\texorpdfstring{\nolinkurl{PaperXI.CuspidalCartan105}}{PaperXI.CuspidalCartan105}}\label{module:PaperXI.CuspidalCartan105}

\textbf{Selected layer:} base108 arithmetic branch.\quad\textbf{Source lines:} 82.\par\textbf{SHA-256:} {\scriptsize\nolinkurl{b2194352b6debd4ae75e3489078283c8f99b6c759d153356b9853900799a8087}}.

\VerbatimInput[fontsize=\scriptsize,breaklines=true,breakanywhere=true]{excerpts/D_CuspidalCartan105.txt}

\textbf{Direct project imports:} \nolinkurl{PaperXI.CuspidalContraction105}, \nolinkurl{PaperXI.IntertwiningAlong105}, \nolinkurl{PaperXI.SplitRadical105}.

\begin{longtable}{@{}p{0.07\textwidth}p{0.12\textwidth}p{0.75\textwidth}@{}}\toprule Line&Kind&Name\\\midrule\endhead

19&def&\nolinkurl{PaperXI.CuspidalCartan105.Cuspidal}\\

23&theorem&\nolinkurl{PaperXI.CuspidalCartan105.diagonal_intertwiner_zero}\\

48&def&\nolinkurl{PaperXI.CuspidalCartan105.HasCartan}\\

56&theorem&\nolinkurl{PaperXI.CuspidalCartan105.intertwining_support}\\

\bottomrule\end{longtable}

\section{\texorpdfstring{\nolinkurl{PaperXI.CuspidalContraction105}}{PaperXI.CuspidalContraction105}}\label{module:PaperXI.CuspidalContraction105}

\textbf{Selected layer:} base108 arithmetic branch.\quad\textbf{Source lines:} 37.\par\textbf{SHA-256:} {\scriptsize\nolinkurl{cf19708af1422fb89f6f41404483d37fb54258b7a6b1b18428ac52a31995dcbd}}.

\VerbatimInput[fontsize=\scriptsize,breaklines=true,breakanywhere=true]{excerpts/D_CuspidalContraction105.txt}

\textbf{Direct project imports:} none.

\begin{longtable}{@{}p{0.07\textwidth}p{0.12\textwidth}p{0.75\textwidth}@{}}\toprule Line&Kind&Name\\\midrule\endhead

14&def&\nolinkurl{PaperXI.CuspidalContraction105.NoFixed}\\

17&theorem&\nolinkurl{PaperXI.CuspidalContraction105.sum_fixed}\\

23&theorem&\nolinkurl{PaperXI.CuspidalContraction105.sum_eq_zero}\\

26&theorem&\nolinkurl{PaperXI.CuspidalContraction105.map_eq_zero}\\

\bottomrule\end{longtable}

\section{\texorpdfstring{\nolinkurl{PaperXI.CycleCoordinates}}{PaperXI.CycleCoordinates}}\label{module:PaperXI.CycleCoordinates}

\textbf{Selected layer:} v0.44 parent.\quad\textbf{Source lines:} 97.\par\textbf{SHA-256:} {\scriptsize\nolinkurl{0e63d65fbd8eecf42832697e5cbb2b1c5e6c7fda9b165030352c34bb3ec19004}}.

\VerbatimInput[fontsize=\scriptsize,breaklines=true,breakanywhere=true]{excerpts/D_CycleCoordinates.txt}

\textbf{Direct project imports:} \nolinkurl{PaperXI.CyclicPolynomial}.

\begin{longtable}{@{}p{0.07\textwidth}p{0.12\textwidth}p{0.75\textwidth}@{}}\toprule Line&Kind&Name\\\midrule\endhead

13&theorem&\nolinkurl{PaperXI.CycleCoordinates.basis_eq_powers}\\

28&theorem&\nolinkurl{PaperXI.CycleCoordinates.closing_of_last_step}\\

45&theorem&\nolinkurl{PaperXI.CycleCoordinates.cycle_formula}\\

55&theorem&\nolinkurl{PaperXI.CycleCoordinates.rescaled_entry}\\

62&theorem&\nolinkurl{PaperXI.CycleCoordinates.rescaled_charpoly}\\

68&theorem&\nolinkurl{PaperXI.CycleCoordinates.rescaled_det}\\

74&theorem&\nolinkurl{PaperXI.CycleCoordinates.cycle_orientation}\\

77&theorem&\nolinkurl{PaperXI.CycleCoordinates.cycle_local_polynomial}\\

87&theorem&\nolinkurl{PaperXI.CycleCoordinates.rescaled_cycle_formula}\\

\bottomrule\end{longtable}

\section{\texorpdfstring{\nolinkurl{PaperXI.CycleFixed}}{PaperXI.CycleFixed}}\label{module:PaperXI.CycleFixed}

\textbf{Selected layer:} v0.44 parent.\quad\textbf{Source lines:} 138.\par\textbf{SHA-256:} {\scriptsize\nolinkurl{008caab8c32493eb96d7225cdc6ff9c2b44718f2cbf85360a0ebef91be48e3dd}}.

\VerbatimInput[fontsize=\scriptsize,breaklines=true,breakanywhere=true]{excerpts/D_CycleFixed.txt}

\textbf{Direct project imports:} \nolinkurl{PaperXI.CycleCoordinates}.

\begin{longtable}{@{}p{0.07\textwidth}p{0.12\textwidth}p{0.75\textwidth}@{}}\toprule Line&Kind&Name\\\midrule\endhead

11&def&\nolinkurl{PaperXI.CycleFixed.fixedSpace}\\

13&theorem&\nolinkurl{PaperXI.CycleFixed.mem_fixedSpace}\\

16&theorem&\nolinkurl{PaperXI.CycleFixed.coordinate_step}\\

42&theorem&\nolinkurl{PaperXI.CycleFixed.fixed_coordinates_constant}\\

59&theorem&\nolinkurl{PaperXI.CycleFixed.basis_sum_ne_zero}\\

65&theorem&\nolinkurl{PaperXI.CycleFixed.fixed_eq_smul_sum}\\

75&theorem&\nolinkurl{PaperXI.CycleFixed.map_basis_sum}\\

90&theorem&\nolinkurl{PaperXI.CycleFixed.fixedSpace_eq}\\

126&theorem&\nolinkurl{PaperXI.CycleFixed.fixedSpace_finrank}\\

\bottomrule\end{longtable}

\section{\texorpdfstring{\nolinkurl{PaperXI.CycleSpan}}{PaperXI.CycleSpan}}\label{module:PaperXI.CycleSpan}

\textbf{Selected layer:} v0.44 parent.\quad\textbf{Source lines:} 76.\par\textbf{SHA-256:} {\scriptsize\nolinkurl{a96cff77ae41d6de99c41763e94a925f42820c3e7646ee438e6fe3bbb67f8054}}.

\VerbatimInput[fontsize=\scriptsize,breaklines=true,breakanywhere=true]{excerpts/D_CycleSpan.txt}

\textbf{Direct project imports:} \nolinkurl{PaperXI.PacketReturn}.

\begin{longtable}{@{}p{0.07\textwidth}p{0.12\textwidth}p{0.75\textwidth}@{}}\toprule Line&Kind&Name\\\midrule\endhead

10&def&\nolinkurl{PaperXI.CycleSpan.cycleVector}\\

12&def&\nolinkurl{PaperXI.CycleSpan.cycleSpan}\\

15&theorem&\nolinkurl{PaperXI.CycleSpan.power_mem}\\

18&theorem&\nolinkurl{PaperXI.CycleSpan.base_mem}\\

21&theorem&\nolinkurl{PaperXI.CycleSpan.map_mem}\\

30&theorem&\nolinkurl{PaperXI.CycleSpan.forward_power}\\

33&theorem&\nolinkurl{PaperXI.CycleSpan.backward_power}\\

37&theorem&\nolinkurl{PaperXI.CycleSpan.inverse_wrap}\\

43&theorem&\nolinkurl{PaperXI.CycleSpan.forward_stable}\\

55&theorem&\nolinkurl{PaperXI.CycleSpan.backward_stable}\\

69&theorem&\nolinkurl{PaperXI.CycleSpan.span_ne_bot}\\

\bottomrule\end{longtable}

\section{\texorpdfstring{\nolinkurl{PaperXI.CyclicAction}}{PaperXI.CyclicAction}}\label{module:PaperXI.CyclicAction}

\textbf{Selected layer:} v0.44 parent.\quad\textbf{Source lines:} 55.\par\textbf{SHA-256:} {\scriptsize\nolinkurl{581f94955c014edf5c55919aa1a8a85b4fc7ec5ec07b9110e81dea08ad0648cf}}.

\VerbatimInput[fontsize=\scriptsize,breaklines=true,breakanywhere=true]{excerpts/D_CyclicAction.txt}

\textbf{Direct project imports:} \nolinkurl{PaperXI.FourierCovariance}.

\begin{longtable}{@{}p{0.07\textwidth}p{0.12\textwidth}p{0.75\textwidth}@{}}\toprule Line&Kind&Name\\\midrule\endhead

12&abbrev&\nolinkurl{PaperXI.CyclicAction.Action}\\

14&def&\nolinkurl{PaperXI.CyclicAction.operator}\\

17&theorem&\nolinkurl{PaperXI.CyclicAction.operator_zero}\\

20&theorem&\nolinkurl{PaperXI.CyclicAction.operator_add}\\

23&theorem&\nolinkurl{PaperXI.CyclicAction.operator_nsmul}\\

31&theorem&\nolinkurl{PaperXI.CyclicAction.operator_mul}\\

36&theorem&\nolinkurl{PaperXI.CyclicAction.operator_eq_generator_pow}\\

40&theorem&\nolinkurl{PaperXI.CyclicAction.generator_period}\\

46&theorem&\nolinkurl{PaperXI.CyclicAction.conjugation_from_generator}\\

\bottomrule\end{longtable}

\section{\texorpdfstring{\nolinkurl{PaperXI.CyclicBridge}}{PaperXI.CyclicBridge}}\label{module:PaperXI.CyclicBridge}

\textbf{Selected layer:} v0.44 parent.\quad\textbf{Source lines:} 111.\par\textbf{SHA-256:} {\scriptsize\nolinkurl{c24882a79bb6079e3ec2c68c8abfdd1db6c00634d6dd9045b2be41aa67cf601b}}.

\VerbatimInput[fontsize=\scriptsize,breaklines=true,breakanywhere=true]{excerpts/D_CyclicBridge.txt}

\textbf{Direct project imports:} \nolinkurl{PaperXI.DiscrepancyCharacter}.

\begin{longtable}{@{}p{0.07\textwidth}p{0.12\textwidth}p{0.75\textwidth}@{}}\toprule Line&Kind&Name\\\midrule\endhead

12&theorem&\nolinkurl{PaperXI.CyclicBridge.character_trivial}\\

18&theorem&\nolinkurl{PaperXI.CyclicBridge.semidirect_intertwining}\\

28&theorem&\nolinkurl{PaperXI.CyclicBridge.stabilizer_covariance}\\

41&theorem&\nolinkurl{PaperXI.CyclicBridge.one_phase_bridge}\\

63&def&\nolinkurl{PaperXI.CyclicBridge.semidirectRepresentation}\\

74&theorem&\nolinkurl{PaperXI.CyclicBridge.semidirectRepresentation_inl}\\

80&theorem&\nolinkurl{PaperXI.CyclicBridge.semidirectRepresentation_inr}\\

88&theorem&\nolinkurl{PaperXI.CyclicBridge.one_phase_bridge_from_actions}\\

\bottomrule\end{longtable}

\section{\texorpdfstring{\nolinkurl{PaperXI.CyclicDefect113}}{PaperXI.CyclicDefect113}}\label{module:PaperXI.CyclicDefect113}

\textbf{Selected layer:} Continuation113.\quad\textbf{Source lines:} 339.\par\textbf{SHA-256:} {\scriptsize\nolinkurl{b6fe373c44ebd8725be1683e7f851da03bcc708c896db2103e02a5a31fbb3a19}}.

\VerbatimInput[fontsize=\scriptsize,breaklines=true,breakanywhere=true]{excerpts/D_CyclicDefect113.txt}

\textbf{Direct project imports:} \nolinkurl{PaperXI.ScalarCocycle112}.

\begin{longtable}{@{}p{0.07\textwidth}p{0.12\textwidth}p{0.75\textwidth}@{}}\toprule Line&Kind&Name\\\midrule\endhead

17&def&\nolinkurl{PaperXI.CyclicDefect113.orbitNorm}\\

21&def&\nolinkurl{PaperXI.CyclicDefect113.weightedCharacter}\\

24&theorem&\nolinkurl{PaperXI.CyclicDefect113.orbitNorm_zero}\\

27&theorem&\nolinkurl{PaperXI.CyclicDefect113.orbitNorm_succ}\\

31&theorem&\nolinkurl{PaperXI.CyclicDefect113.weightedCharacter_zero}\\

34&theorem&\nolinkurl{PaperXI.CyclicDefect113.weightedCharacter_succ}\\

39&theorem&\nolinkurl{PaperXI.CyclicDefect113.power_apply_shift}\\

45&theorem&\nolinkurl{PaperXI.CyclicDefect113.weightedCharacter_shift}\\

59&theorem&\nolinkurl{PaperXI.CyclicDefect113.bezout_power}\\

71&def&\nolinkurl{PaperXI.CyclicDefect113.correctingCharacter}\\

76&theorem&\nolinkurl{PaperXI.CyclicDefect113.correctingCharacter_coboundary}\\

92&structure&\nolinkurl{PaperXI.CyclicDefect113.ScalarLift}\\

102&def&\nolinkurl{PaperXI.CyclicDefect113.ScalarLift.phase}\\

104&theorem&\nolinkurl{PaperXI.CyclicDefect113.ScalarLift.phase_one}\\

108&theorem&\nolinkurl{PaperXI.CyclicDefect113.ScalarLift.scalar_formula}\\

121&theorem&\nolinkurl{PaperXI.CyclicDefect113.ScalarLift.phase_mul}\\

131&theorem&\nolinkurl{PaperXI.CyclicDefect113.ScalarLift.power_base}\\

141&theorem&\nolinkurl{PaperXI.CyclicDefect113.ScalarLift.power_scalar}\\

152&theorem&\nolinkurl{PaperXI.CyclicDefect113.ScalarLift.norm_eq_one}\\

162&theorem&\nolinkurl{PaperXI.CyclicDefect113.orbitNorm_coboundary}\\

172&theorem&\nolinkurl{PaperXI.CyclicDefect113.orbitNorm_mul}\\

176&theorem&\nolinkurl{PaperXI.CyclicDefect113.orbitNorm_div}\\

187&def&\nolinkurl{PaperXI.CyclicDefect113.actionDefect}\\

190&theorem&\nolinkurl{PaperXI.CyclicDefect113.actionDefect_one}\\

195&theorem&\nolinkurl{PaperXI.CyclicDefect113.actionDefect_mul}\\

205&def&\nolinkurl{PaperXI.CyclicDefect113.defectCharacter}\\

211&theorem&\nolinkurl{PaperXI.CyclicDefect113.actionDefect_norm}\\

223&def&\nolinkurl{PaperXI.CyclicDefect113.correctedCochain}\\

227&theorem&\nolinkurl{PaperXI.CyclicDefect113.correctedCochain_one}\\

232&theorem&\nolinkurl{PaperXI.CyclicDefect113.correctedCochain_cocycle}\\

243&theorem&\nolinkurl{PaperXI.CyclicDefect113.correctedCochain_covariant}\\

264&def&\nolinkurl{PaperXI.CyclicDefect113.scalarBridge}\\

279&theorem&\nolinkurl{PaperXI.CyclicDefect113.scalarBridge_base}\\

284&theorem&\nolinkurl{PaperXI.CyclicDefect113.scalarBridge_fixes_scalar}\\

290&theorem&\nolinkurl{PaperXI.CyclicDefect113.scalarBridge_equivariant}\\

308&theorem&\nolinkurl{PaperXI.CyclicDefect113.exists_equivariant_cochain}\\

325&theorem&\nolinkurl{PaperXI.CyclicDefect113.exists_equivariant_scalarBridge}\\

\bottomrule\end{longtable}

\section{\texorpdfstring{\nolinkurl{PaperXI.CyclicFrobeniusCount}}{PaperXI.CyclicFrobeniusCount}}\label{module:PaperXI.CyclicFrobeniusCount}

\textbf{Selected layer:} v0.44 parent.\quad\textbf{Source lines:} 120.\par\textbf{SHA-256:} {\scriptsize\nolinkurl{b97169de1bf48a86a58bf00236bad78341081a6db9ae46f6a8e49f95c87a4f82}}.

\VerbatimInput[fontsize=\scriptsize,breaklines=true,breakanywhere=true]{excerpts/D_CyclicFrobeniusCount.txt}

\textbf{Direct project imports:} \nolinkurl{PaperXI.FrobeniusPeriodCount}.

\begin{longtable}{@{}p{0.07\textwidth}p{0.12\textwidth}p{0.75\textwidth}@{}}\toprule Line&Kind&Name\\\midrule\endhead

16&theorem&\nolinkurl{PaperXI.CyclicFrobeniusCount.iterate_power}\\

24&theorem&\nolinkurl{PaperXI.CyclicFrobeniusCount.fixed_iff_kernel}\\

40&theorem&\nolinkurl{PaperXI.CyclicFrobeniusCount.fixedCount_gcd}\\

48&theorem&\nolinkurl{PaperXI.CyclicFrobeniusCount.fixedCount_pow_gcd}\\

54&theorem&\nolinkurl{PaperXI.CyclicFrobeniusCount.fixedCount_of_dvd}\\

60&theorem&\nolinkurl{PaperXI.CyclicFrobeniusCount.exactCount_formula}\\

72&theorem&\nolinkurl{PaperXI.CyclicFrobeniusCount.card_coprime_q}\\

79&def&\nolinkurl{PaperXI.CyclicFrobeniusCount.frobeniusPerm}\\

84&theorem&\nolinkurl{PaperXI.CyclicFrobeniusCount.frobeniusPerm_apply}\\

89&theorem&\nolinkurl{PaperXI.CyclicFrobeniusCount.orbitCount_mul_formula}\\

100&theorem&\nolinkurl{PaperXI.CyclicFrobeniusCount.moebius_sum_dvd}\\

108&theorem&\nolinkurl{PaperXI.CyclicFrobeniusCount.orbitCount_formula}\\

\bottomrule\end{longtable}

\section{\texorpdfstring{\nolinkurl{PaperXI.CyclicInducedBasis}}{PaperXI.CyclicInducedBasis}}\label{module:PaperXI.CyclicInducedBasis}

\textbf{Selected layer:} v0.44 parent.\quad\textbf{Source lines:} 83.\par\textbf{SHA-256:} {\scriptsize\nolinkurl{a90ba613fe359b594d103fb19aa93529a4e4cb8584c38b335950bcb1ceaf0eb1}}.

\VerbatimInput[fontsize=\scriptsize,breaklines=true,breakanywhere=true]{excerpts/D_CyclicInducedBasis.txt}

\textbf{Direct project imports:} \nolinkurl{PaperXI.CyclicInducedCoordinates}.

\begin{longtable}{@{}p{0.07\textwidth}p{0.12\textwidth}p{0.75\textwidth}@{}}\toprule Line&Kind&Name\\\midrule\endhead

15&def&\nolinkurl{PaperXI.CyclicInduced.cyclicRep}\\

19&theorem&\nolinkurl{PaperXI.CyclicInduced.cyclicRep_degree}\\

25&theorem&\nolinkurl{PaperXI.CyclicInduced.cyclicRep_one}\\

30&def&\nolinkurl{PaperXI.CyclicInduced.cyclicBasis}\\

34&def&\nolinkurl{PaperXI.CyclicInduced.cyclicWeight}\\

37&theorem&\nolinkurl{PaperXI.CyclicInduced.cyclicBasis_apply}\\

43&theorem&\nolinkurl{PaperXI.CyclicInduced.cyclicBasis_action}\\

51&def&\nolinkurl{PaperXI.CyclicInduced.cyclicCharacter}\\

54&theorem&\nolinkurl{PaperXI.CyclicInduced.cyclicCharacter_apply}\\

64&theorem&\nolinkurl{PaperXI.CyclicInduced.cyclicBasis_action_kernel}\\

71&theorem&\nolinkurl{PaperXI.CyclicInduced.cyclicBasis_action_zero}\\

78&theorem&\nolinkurl{PaperXI.CyclicInduced.cyclicBasis_action_phi}\\

\bottomrule\end{longtable}

\section{\texorpdfstring{\nolinkurl{PaperXI.CyclicInducedCharacter}}{PaperXI.CyclicInducedCharacter}}\label{module:PaperXI.CyclicInducedCharacter}

\textbf{Selected layer:} v0.44 parent.\quad\textbf{Source lines:} 60.\par\textbf{SHA-256:} {\scriptsize\nolinkurl{15707518268a187bd5d4f486516c842c236f5385de68af3fb7ef63eb83fcd4f5}}.

\VerbatimInput[fontsize=\scriptsize,breaklines=true,breakanywhere=true]{excerpts/D_CyclicInducedCharacter.txt}

\textbf{Direct project imports:} \nolinkurl{PaperXI.CyclicInducedDegree}.

\begin{longtable}{@{}p{0.07\textwidth}p{0.12\textwidth}p{0.75\textwidth}@{}}\toprule Line&Kind&Name\\\midrule\endhead

13&theorem&\nolinkurl{PaperXI.CyclicInduced.degreeWeight_kernel}\\

19&theorem&\nolinkurl{PaperXI.CyclicInduced.degreeCharacter_apply}\\

28&theorem&\nolinkurl{PaperXI.CyclicInduced.degreeCharacter_zero}\\

36&theorem&\nolinkurl{PaperXI.CyclicInduced.degreeWeight_zero}\\

43&theorem&\nolinkurl{PaperXI.CyclicInduced.degreeWeight_restrict}\\

54&theorem&\nolinkurl{PaperXI.CyclicInduced.degreeWeight_restrict_zero}\\

\bottomrule\end{longtable}

\section{\texorpdfstring{\nolinkurl{PaperXI.CyclicInducedCoordinates}}{PaperXI.CyclicInducedCoordinates}}\label{module:PaperXI.CyclicInducedCoordinates}

\textbf{Selected layer:} v0.44 parent.\quad\textbf{Source lines:} 179.\par\textbf{SHA-256:} {\scriptsize\nolinkurl{302c38129316e682653aa44bbff88450ec77341b005fbdbce392a31a475f2d90}}.

\VerbatimInput[fontsize=\scriptsize,breaklines=true,breakanywhere=true]{excerpts/D_CyclicInducedCoordinates.txt}

\textbf{Direct project imports:} none.

\begin{longtable}{@{}p{0.07\textwidth}p{0.12\textwidth}p{0.75\textwidth}@{}}\toprule Line&Kind&Name\\\midrule\endhead

17&def&\nolinkurl{PaperXI.CyclicInduced.scalar}\\

22&theorem&\nolinkurl{PaperXI.CyclicInduced.scalar_apply}\\

26&def&\nolinkurl{PaperXI.CyclicInduced.residual}\\

29&theorem&\nolinkurl{PaperXI.CyclicInduced.residual_rep}\\

33&theorem&\nolinkurl{PaperXI.CyclicInduced.residual_mul}\\

40&abbrev&\nolinkurl{PaperXI.CyclicInduced.CoindSpace}\\

41&abbrev&\nolinkurl{PaperXI.CyclicInduced.coinduced}\\

42&abbrev&\nolinkurl{PaperXI.CyclicInduced.IndSpace}\\

43&abbrev&\nolinkurl{PaperXI.CyclicInduced.induced}\\

46&def&\nolinkurl{PaperXI.CyclicInduced.extend}\\

52&theorem&\nolinkurl{PaperXI.CyclicInduced.extend_rep}\\

57&def&\nolinkurl{PaperXI.CyclicInduced.coordinates}\\

70&theorem&\nolinkurl{PaperXI.CyclicInduced.coordinates_apply}\\

74&def&\nolinkurl{PaperXI.CyclicInduced.weight}\\

80&def&\nolinkurl{PaperXI.CyclicInduced.coindBasis}\\

83&theorem&\nolinkurl{PaperXI.CyclicInduced.coindBasis_rep}\\

94&theorem&\nolinkurl{PaperXI.CyclicInduced.coindBasis_action}\\

120&def&\nolinkurl{PaperXI.CyclicInduced.indCoind}\\

125&def&\nolinkurl{PaperXI.CyclicInduced.basis}\\

129&theorem&\nolinkurl{PaperXI.CyclicInduced.basis_action}\\

142&def&\nolinkurl{PaperXI.CyclicInduced.conjugate}\\

150&def&\nolinkurl{PaperXI.CyclicInduced.conjugateCharacter}\\

154&theorem&\nolinkurl{PaperXI.CyclicInduced.weight_kernel}\\

163&theorem&\nolinkurl{PaperXI.CyclicInduced.basis_action_kernel}\\

170&theorem&\nolinkurl{PaperXI.CyclicInduced.basis_action_origin}\\

\bottomrule\end{longtable}

\section{\texorpdfstring{\nolinkurl{PaperXI.CyclicInducedDegree}}{PaperXI.CyclicInducedDegree}}\label{module:PaperXI.CyclicInducedDegree}

\textbf{Selected layer:} v0.44 parent.\quad\textbf{Source lines:} 59.\par\textbf{SHA-256:} {\scriptsize\nolinkurl{e60c559b41a83de4e58e4a01b0633216f888976f7748d01b1b05f7aa42aaf085}}.

\VerbatimInput[fontsize=\scriptsize,breaklines=true,breakanywhere=true]{excerpts/D_CyclicInducedDegree.txt}

\textbf{Direct project imports:} \nolinkurl{PaperXI.CyclicInducedBasis}, \nolinkurl{PaperXI.TameDegree}.

\begin{longtable}{@{}p{0.07\textwidth}p{0.12\textwidth}p{0.75\textwidth}@{}}\toprule Line&Kind&Name\\\midrule\endhead

15&theorem&\nolinkurl{PaperXI.CyclicInduced.degreeMod_phi}\\

20&abbrev&\nolinkurl{PaperXI.CyclicInduced.degreeInduced}\\

21&abbrev&\nolinkurl{PaperXI.CyclicInduced.DegreeIndSpace}\\

24&def&\nolinkurl{PaperXI.CyclicInduced.degreeBasis}\\

27&def&\nolinkurl{PaperXI.CyclicInduced.degreeWeight}\\

31&def&\nolinkurl{PaperXI.CyclicInduced.degreeCharacter}\\

35&theorem&\nolinkurl{PaperXI.CyclicInduced.degreeBasis_action}\\

42&theorem&\nolinkurl{PaperXI.CyclicInduced.degreeBasis_action_kernel}\\

48&theorem&\nolinkurl{PaperXI.CyclicInduced.degreeBasis_action_zero}\\

54&theorem&\nolinkurl{PaperXI.CyclicInduced.degreeBasis_action_phi}\\

\bottomrule\end{longtable}

\section{\texorpdfstring{\nolinkurl{PaperXI.CyclicInducedWrap}}{PaperXI.CyclicInducedWrap}}\label{module:PaperXI.CyclicInducedWrap}

\textbf{Selected layer:} v0.44 parent.\quad\textbf{Source lines:} 119.\par\textbf{SHA-256:} {\scriptsize\nolinkurl{4dec20ab681bfc9225f0bf157b406ac02ea6272612eb84376cb4b647af558bb9}}.

\VerbatimInput[fontsize=\scriptsize,breaklines=true,breakanywhere=true]{excerpts/D_CyclicInducedWrap.txt}

\textbf{Direct project imports:} \nolinkurl{PaperXI.CyclicInducedDegree}.

\begin{longtable}{@{}p{0.07\textwidth}p{0.12\textwidth}p{0.75\textwidth}@{}}\toprule Line&Kind&Name\\\midrule\endhead

14&theorem&\nolinkurl{PaperXI.CyclicInduced.cyclicWeight_phi_value}\\

38&def&\nolinkurl{PaperXI.CyclicInduced.period}\\

44&theorem&\nolinkurl{PaperXI.CyclicInduced.val_succ_of_lt}\\

50&theorem&\nolinkurl{PaperXI.CyclicInduced.val_succ_of_last}\\

57&theorem&\nolinkurl{PaperXI.CyclicInduced.cyclicWeight_phi_of_lt}\\

69&theorem&\nolinkurl{PaperXI.CyclicInduced.cyclicWeight_phi_of_last}\\

85&def&\nolinkurl{PaperXI.CyclicInduced.degreePeriod}\\

89&theorem&\nolinkurl{PaperXI.CyclicInduced.degreePeriod_coe}\\

92&theorem&\nolinkurl{PaperXI.CyclicInduced.degreeWeight_phi_of_lt}\\

97&theorem&\nolinkurl{PaperXI.CyclicInduced.degreeWeight_phi_of_last}\\

102&theorem&\nolinkurl{PaperXI.CyclicInduced.degreeBasis_phi_of_lt}\\

109&theorem&\nolinkurl{PaperXI.CyclicInduced.degreeBasis_phi_of_last}\\

\bottomrule\end{longtable}

\section{\texorpdfstring{\nolinkurl{PaperXI.CyclicLinearization}}{PaperXI.CyclicLinearization}}\label{module:PaperXI.CyclicLinearization}

\textbf{Selected layer:} v0.44 parent.\quad\textbf{Source lines:} 83.\par\textbf{SHA-256:} {\scriptsize\nolinkurl{77c69c57b811f5321306490bb3bb4aef92a7c262d5d875a613ef8d63e5170d0c}}.

\VerbatimInput[fontsize=\scriptsize,breaklines=true,breakanywhere=true]{excerpts/D_CyclicLinearization.txt}

\textbf{Direct project imports:} none.

\begin{longtable}{@{}p{0.07\textwidth}p{0.12\textwidth}p{0.75\textwidth}@{}}\toprule Line&Kind&Name\\\midrule\endhead

16&theorem&\nolinkurl{PaperXI.CyclicLinearization.unit_root}\\

29&theorem&\nolinkurl{PaperXI.CyclicLinearization.rescale_generator}\\

50&theorem&\nolinkurl{PaperXI.CyclicLinearization.lift_of_generator}\\

65&theorem&\nolinkurl{PaperXI.CyclicLinearization.finite_cyclic_linearization}\\

73&theorem&\nolinkurl{PaperXI.CyclicLinearization.finite_cyclic_scalar_equivalence}\\

\bottomrule\end{longtable}

\section{\texorpdfstring{\nolinkurl{PaperXI.CyclicNormalization113}}{PaperXI.CyclicNormalization113}}\label{module:PaperXI.CyclicNormalization113}

\textbf{Selected layer:} Continuation113.\quad\textbf{Source lines:} 386.\par\textbf{SHA-256:} {\scriptsize\nolinkurl{607587d94a544d2c4e014d81b21e4c9241d852feddfaa826df57f9715d396ce7}}.

\VerbatimInput[fontsize=\scriptsize,breaklines=true,breakanywhere=true]{excerpts/D_CyclicNormalization113.txt}

\textbf{Direct project imports:} \nolinkurl{PaperXI.GlobalHeisenbergBridge112}.

\begin{longtable}{@{}p{0.07\textwidth}p{0.12\textwidth}p{0.75\textwidth}@{}}\toprule Line&Kind&Name\\\midrule\endhead

18&def&\nolinkurl{PaperXI.CyclicNormalization113.scale}\\

21&theorem&\nolinkurl{PaperXI.CyclicNormalization113.scale_apply}\\

24&theorem&\nolinkurl{PaperXI.CyclicNormalization113.scale_pow_apply}\\

33&theorem&\nolinkurl{PaperXI.CyclicNormalization113.scalar_correction_exists}\\

48&theorem&\nolinkurl{PaperXI.CyclicNormalization113.det_scale}\\

58&theorem&\nolinkurl{PaperXI.CyclicNormalization113.normalize_honest}\\

78&theorem&\nolinkurl{PaperXI.CyclicNormalization113.implements_pow}\\

90&theorem&\nolinkurl{PaperXI.CyclicNormalization113.scale_implements}\\

101&theorem&\nolinkurl{PaperXI.CyclicNormalization113.implementer_power_scalar}\\

111&theorem&\nolinkurl{PaperXI.CyclicNormalization113.unit_nth_root}\\

124&theorem&\nolinkurl{PaperXI.CyclicNormalization113.normalize_order}\\

140&theorem&\nolinkurl{PaperXI.CyclicNormalization113.normalized_generator_exists}\\

153&theorem&\nolinkurl{PaperXI.CyclicNormalization113.normalized_generator_unique}\\

167&theorem&\nolinkurl{PaperXI.CyclicNormalization113.normalized_generator_existsUnique}\\

189&theorem&\nolinkurl{PaperXI.CyclicNormalization113.implementer_of_actual_profile}\\

219&theorem&\nolinkurl{PaperXI.CyclicNormalization113.canonical_generator_of_actual_profile}\\

239&theorem&\nolinkurl{PaperXI.CyclicNormalization113.canonical_generator_of_actual_form}\\

268&def&\nolinkurl{PaperXI.CyclicNormalization113.cyclicAction}\\

276&theorem&\nolinkurl{PaperXI.CyclicNormalization113.cyclicAction_generator}\\

286&theorem&\nolinkurl{PaperXI.CyclicNormalization113.canonical_action_of_implementer}\\

325&theorem&\nolinkurl{PaperXI.CyclicNormalization113.canonical_action_of_actual_profile}\\

347&theorem&\nolinkurl{PaperXI.CyclicNormalization113.canonical_equivariant_bridge_of_actual_profile}\\

\bottomrule\end{longtable}

\section{\texorpdfstring{\nolinkurl{PaperXI.CyclicPolynomial}}{PaperXI.CyclicPolynomial}}\label{module:PaperXI.CyclicPolynomial}

\textbf{Selected layer:} v0.44 parent.\quad\textbf{Source lines:} 81.\par\textbf{SHA-256:} {\scriptsize\nolinkurl{0b7bf89b38ea2b0bc63675d4707494ec562ab9f04b8218c17fc092d9e82ea6ab}}.

\VerbatimInput[fontsize=\scriptsize,breaklines=true,breakanywhere=true]{excerpts/D_CyclicPolynomial.txt}

\textbf{Direct project imports:} none.

\begin{longtable}{@{}p{0.07\textwidth}p{0.12\textwidth}p{0.75\textwidth}@{}}\toprule Line&Kind&Name\\\midrule\endhead

13&theorem&\nolinkurl{PaperXI.CyclicPolynomial.aeval_apply_eq_sum}\\

20&theorem&\nolinkurl{PaperXI.CyclicPolynomial.no_short_annihilator}\\

31&theorem&\nolinkurl{PaperXI.CyclicPolynomial.charpoly_of_orbit_basis}\\

45&theorem&\nolinkurl{PaperXI.CyclicPolynomial.det_of_orbit_basis}\\

59&theorem&\nolinkurl{PaperXI.CyclicPolynomial.eval_charpoly_of_orbit_basis}\\

67&theorem&\nolinkurl{PaperXI.CyclicPolynomial.reverse_cycle_polynomial}\\

72&theorem&\nolinkurl{PaperXI.CyclicPolynomial.local_polynomial_of_orbit_basis}\\

\bottomrule\end{longtable}

\section{\texorpdfstring{\nolinkurl{PaperXI.DVRCartan106}}{PaperXI.DVRCartan106}}\label{module:PaperXI.DVRCartan106}

\textbf{Selected layer:} base108 arithmetic branch.\quad\textbf{Source lines:} 68.\par\textbf{SHA-256:} {\scriptsize\nolinkurl{450b28f31d92b480d5a52f75464a037a0aaf3ebd66ce7753b50f9e474951eeda}}.

\VerbatimInput[fontsize=\scriptsize,breaklines=true,breakanywhere=true]{excerpts/D_DVRCartan106.txt}

\textbf{Direct project imports:} \nolinkurl{PaperXI.FractionMatrix106}, \nolinkurl{PaperXI.CuspidalCartan105}.

\begin{longtable}{@{}p{0.07\textwidth}p{0.12\textwidth}p{0.75\textwidth}@{}}\toprule Line&Kind&Name\\\midrule\endhead

15&def&\nolinkurl{PaperXI.DVRCartan106.unitsDiagonal}\\

21&theorem&\nolinkurl{PaperXI.DVRCartan106.normalize_diagonal}\\

55&theorem&\nolinkurl{PaperXI.DVRCartan106.hasCartan}\\

\bottomrule\end{longtable}

\section{\texorpdfstring{\nolinkurl{PaperXI.DegreeIndex110}}{PaperXI.DegreeIndex110}}\label{module:PaperXI.DegreeIndex110}

\textbf{Selected layer:} Continuation110.\quad\textbf{Source lines:} 432.\par\textbf{SHA-256:} {\scriptsize\nolinkurl{0f87f63a5f0f5c28adbb1a04732fbe110b63a5f45cdde03e98830271c7c6a6c5}}.

\VerbatimInput[fontsize=\scriptsize,breaklines=true,breakanywhere=true]{excerpts/D_DegreeIndex110.txt}

\textbf{Direct project imports:} \nolinkurl{PaperXI.FiniteProfileBridge109}, \nolinkurl{PaperXI.CompactIntersection108}, \nolinkurl{PaperXI.ActualIntersection108}, \nolinkurl{PaperXI.PairIndex110}.

\begin{longtable}{@{}p{0.07\textwidth}p{0.12\textwidth}p{0.75\textwidth}@{}}\toprule Line&Kind&Name\\\midrule\endhead

20&theorem&\nolinkurl{PaperXI.DegreeIndex110.degree_square_of_finite_model}\\

44&theorem&\nolinkurl{PaperXI.DegreeIndex110.degree_square_of_finite_kernel_quotient}\\

58&theorem&\nolinkurl{PaperXI.DegreeIndex110.degree_square}\\

71&theorem&\nolinkurl{PaperXI.DegreeIndex110.intersection_index_of_oriented_identity}\\

85&theorem&\nolinkurl{PaperXI.DegreeIndex110.dimension_one_iff_original_index}\\

120&theorem&\nolinkurl{PaperXI.DegreeIndex110.intersection_dimension_one_of_index}\\

154&theorem&\nolinkurl{PaperXI.DegreeIndex110.intersection_dimension_one}\\

186&theorem&\nolinkurl{PaperXI.DegreeIndex110.intersection_index_pos}\\

199&theorem&\nolinkurl{PaperXI.DegreeIndex110.intersection_index_ambient}\\

205&theorem&\nolinkurl{PaperXI.DegreeIndex110.scalarIntersection_index_ambient}\\

214&theorem&\nolinkurl{PaperXI.DegreeIndex110.oriented_identity_of_ambient}\\

236&def&\nolinkurl{PaperXI.DegreeIndex110.actualScalarSubgroup}\\

240&def&\nolinkurl{PaperXI.DegreeIndex110.actualScalarCharacter}\\

251&theorem&\nolinkurl{PaperXI.DegreeIndex110.actual_scalar}\\

264&theorem&\nolinkurl{PaperXI.DegreeIndex110.actualScalarSubgroup_image}\\

275&theorem&\nolinkurl{PaperXI.DegreeIndex110.actual_embedding_range_open}\\

295&theorem&\nolinkurl{PaperXI.DegreeIndex110.actual_eta_dimension_one}\\

343&theorem&\nolinkurl{PaperXI.DegreeIndex110.actual_eta_dimension}\\

390&theorem&\nolinkurl{PaperXI.DegreeIndex110.actual_eta_dimension_ambient}\\

\bottomrule\end{longtable}

\section{\texorpdfstring{\nolinkurl{PaperXI.DerivedCentralCharacter69}}{PaperXI.DerivedCentralCharacter69}}\label{module:PaperXI.DerivedCentralCharacter69}

\textbf{Selected layer:} base108 arithmetic branch.\quad\textbf{Source lines:} 41.\par\textbf{SHA-256:} {\scriptsize\nolinkurl{861257cc1f0b38c19e73ec7fec5cb12d5ee9e3affca53ab38d5287c502d56d0a}}.

\VerbatimInput[fontsize=\scriptsize,breaklines=true,breakanywhere=true]{excerpts/D_DerivedCentralCharacter69.txt}

\textbf{Direct project imports:} none.

\begin{longtable}{@{}p{0.07\textwidth}p{0.12\textwidth}p{0.75\textwidth}@{}}\toprule Line&Kind&Name\\\midrule\endhead

15&def&\nolinkurl{PaperXI.DerivedCentralCharacter69.centralAction}\\

25&def&\nolinkurl{PaperXI.DerivedCentralCharacter69.scalarEquiv}\\

30&def&\nolinkurl{PaperXI.DerivedCentralCharacter69.character}\\

34&theorem&\nolinkurl{PaperXI.DerivedCentralCharacter69.character_apply}\\

\bottomrule\end{longtable}

\section{\texorpdfstring{\nolinkurl{PaperXI.DiscrepancyCharacter}}{PaperXI.DiscrepancyCharacter}}\label{module:PaperXI.DiscrepancyCharacter}

\textbf{Selected layer:} v0.44 parent.\quad\textbf{Source lines:} 50.\par\textbf{SHA-256:} {\scriptsize\nolinkurl{cf201fbd9c03f12816213b45d8e418ea5649df147dd993af122f826106f798af}}.

\VerbatimInput[fontsize=\scriptsize,breaklines=true,breakanywhere=true]{excerpts/D_DiscrepancyCharacter.txt}

\textbf{Direct project imports:} \nolinkurl{PaperXI.SchurScalar}.

\begin{longtable}{@{}p{0.07\textwidth}p{0.12\textwidth}p{0.75\textwidth}@{}}\toprule Line&Kind&Name\\\midrule\endhead

13&theorem&\nolinkurl{PaperXI.DiscrepancyCharacter.exists_character}\\

\bottomrule\end{longtable}

\section{\texorpdfstring{\nolinkurl{PaperXI.DualCharacterSums}}{PaperXI.DualCharacterSums}}\label{module:PaperXI.DualCharacterSums}

\textbf{Selected layer:} v0.44 parent.\quad\textbf{Source lines:} 55.\par\textbf{SHA-256:} {\scriptsize\nolinkurl{cd4cf6b9b43305d2e09f09566db4519b5821134b6fac03ac8e92c83fad7c0532}}.

\VerbatimInput[fontsize=\scriptsize,breaklines=true,breakanywhere=true]{excerpts/D_DualCharacterSums.txt}

\textbf{Direct project imports:} \nolinkurl{PaperXI.SchrodingerTrace}, \nolinkurl{PaperXI.SchrodingerSeparation}, \nolinkurl{PaperXI.HeisenbergDimension}.

\begin{longtable}{@{}p{0.07\textwidth}p{0.12\textwidth}p{0.75\textwidth}@{}}\toprule Line&Kind&Name\\\midrule\endhead

14&theorem&\nolinkurl{PaperXI.DualCharacters.sum_over_points}\\

22&def&\nolinkurl{PaperXI.DualCharacters.evaluationCharacter}\\

28&theorem&\nolinkurl{PaperXI.DualCharacters.evaluationCharacter_apply}\\

31&theorem&\nolinkurl{PaperXI.DualCharacters.evaluationCharacter_ne_one}\\

40&theorem&\nolinkurl{PaperXI.DualCharacters.sum_over_dual}\\

50&theorem&\nolinkurl{PaperXI.DualCharacters.card_dual}\\

\bottomrule\end{longtable}

\section{\texorpdfstring{\nolinkurl{PaperXI.DualityLinear}}{PaperXI.DualityLinear}}\label{module:PaperXI.DualityLinear}

\textbf{Selected layer:} v0.44 parent.\quad\textbf{Source lines:} 94.\par\textbf{SHA-256:} {\scriptsize\nolinkurl{0f856e352f57db7831276fa70666d08465b1a79f9ac97e0b4a31265892c47e12}}.

\VerbatimInput[fontsize=\scriptsize,breaklines=true,breakanywhere=true]{excerpts/D_DualityLinear.txt}

\textbf{Direct project imports:} \nolinkurl{PaperXI.RepresentationStability}.

\begin{longtable}{@{}p{0.07\textwidth}p{0.12\textwidth}p{0.75\textwidth}@{}}\toprule Line&Kind&Name\\\midrule\endhead

17&def&\nolinkurl{PaperXI.DualityLinear.invariantsEquiv}\\

35&def&\nolinkurl{PaperXI.DualityLinear.bidualEquiv}\\

43&theorem&\nolinkurl{PaperXI.DualityLinear.coannihilator_stable}\\

54&theorem&\nolinkurl{PaperXI.DualityLinear.dual_irreducible}\\

71&def&\nolinkurl{PaperXI.DualityLinear.tensorHomEquiv}\\

82&def&\nolinkurl{PaperXI.DualityLinear.fixedTensorHomEquiv}\\

88&theorem&\nolinkurl{PaperXI.DualityLinear.hom_finrank_eq_tensor_invariants}\\

\bottomrule\end{longtable}

\section{\texorpdfstring{\nolinkurl{PaperXI.DyadicAbstract}}{PaperXI.DyadicAbstract}}\label{module:PaperXI.DyadicAbstract}

\textbf{Selected layer:} v0.44 parent.\quad\textbf{Source lines:} 159.\par\textbf{SHA-256:} {\scriptsize\nolinkurl{1e70f115444ed1310c3bdae48b5dd98530fa9d8f2eae45684430b7cec5e836d0}}.

\VerbatimInput[fontsize=\scriptsize,breaklines=true,breakanywhere=true]{excerpts/D_DyadicAbstract.txt}

\textbf{Direct project imports:} \nolinkurl{PaperXI.SymplecticCoordinates}, \nolinkurl{PaperXI.DyadicBasisConsequences}, \nolinkurl{PaperXI.DyadicBasisCounts}.

\begin{longtable}{@{}p{0.07\textwidth}p{0.12\textwidth}p{0.75\textwidth}@{}}\toprule Line&Kind&Name\\\midrule\endhead

15&def&\nolinkurl{PaperXI.DyadicAbstract.coordinates}\\

20&theorem&\nolinkurl{PaperXI.DyadicAbstract.coordinates_form}\\

25&def&\nolinkurl{PaperXI.DyadicAbstract.arf}\\

32&theorem&\nolinkurl{PaperXI.DyadicAbstract.arf_eq_basis}\\

43&def&\nolinkurl{PaperXI.DyadicAbstract.gaussSum}\\

50&theorem&\nolinkurl{PaperXI.DyadicAbstract.gaussSum_eq_sum}\\

58&theorem&\nolinkurl{PaperXI.DyadicAbstract.gauss_sum}\\

68&def&\nolinkurl{PaperXI.DyadicAbstract.normalizedGaussSign}\\

73&theorem&\nolinkurl{PaperXI.DyadicAbstract.normalizedGaussSign_eq_sign}\\

83&theorem&\nolinkurl{PaperXI.DyadicAbstract.arf_translation}\\

92&theorem&\nolinkurl{PaperXI.DyadicAbstract.opposite_refinements}\\

113&theorem&\nolinkurl{PaperXI.DyadicAbstract.even_count}\\

123&theorem&\nolinkurl{PaperXI.DyadicAbstract.odd_count}\\

133&theorem&\nolinkurl{PaperXI.DyadicAbstract.gauss_sum_of_finrank}\\

142&theorem&\nolinkurl{PaperXI.DyadicAbstract.even_count_of_finrank}\\

151&theorem&\nolinkurl{PaperXI.DyadicAbstract.odd_count_of_finrank}\\

\bottomrule\end{longtable}

\section{\texorpdfstring{\nolinkurl{PaperXI.DyadicAbstractOrthogonal}}{PaperXI.DyadicAbstractOrthogonal}}\label{module:PaperXI.DyadicAbstractOrthogonal}

\textbf{Selected layer:} v0.44 parent.\quad\textbf{Source lines:} 141.\par\textbf{SHA-256:} {\scriptsize\nolinkurl{84c944337500a19941c2a2a1b839131b3c7ffd8b54f89c118ebb29f3416b811f}}.

\VerbatimInput[fontsize=\scriptsize,breaklines=true,breakanywhere=true]{excerpts/D_DyadicAbstractOrthogonal.txt}

\textbf{Direct project imports:} \nolinkurl{PaperXI.DyadicAbstract}, \nolinkurl{PaperXI.DyadicIsometry}, \nolinkurl{PaperXI.DyadicOrthogonal}.

\begin{longtable}{@{}p{0.07\textwidth}p{0.12\textwidth}p{0.75\textwidth}@{}}\toprule Line&Kind&Name\\\midrule\endhead

17&theorem&\nolinkurl{PaperXI.DyadicAbstractOrthogonal.form_alternating}\\

21&theorem&\nolinkurl{PaperXI.DyadicAbstractOrthogonal.form_nondegenerate}\\

38&theorem&\nolinkurl{PaperXI.DyadicAbstractOrthogonal.arf_transport}\\

52&theorem&\nolinkurl{PaperXI.DyadicAbstractOrthogonal.arf_add}\\

71&theorem&\nolinkurl{PaperXI.DyadicAbstractOrthogonal.arf_sign_product}\\

84&theorem&\nolinkurl{PaperXI.DyadicAbstractOrthogonal.arf_of_decomposition}\\

107&theorem&\nolinkurl{PaperXI.DyadicAbstractOrthogonal.arf_sign_of_decomposition}\\

116&theorem&\nolinkurl{PaperXI.DyadicAbstractOrthogonal.normalized_sign_product}\\

129&theorem&\nolinkurl{PaperXI.DyadicAbstractOrthogonal.normalized_sign_of_decomposition}\\

\bottomrule\end{longtable}

\section{\texorpdfstring{\nolinkurl{PaperXI.DyadicAffine}}{PaperXI.DyadicAffine}}\label{module:PaperXI.DyadicAffine}

\textbf{Selected layer:} v0.44 parent.\quad\textbf{Source lines:} 81.\par\textbf{SHA-256:} {\scriptsize\nolinkurl{06fd5b5a57b17060abd22a11fb29592fcd53b06aa09388b92314fe2fd23692e3}}.

\VerbatimInput[fontsize=\scriptsize,breaklines=true,breakanywhere=true]{excerpts/D_DyadicAffine.txt}

\textbf{Direct project imports:} none.

\begin{longtable}{@{}p{0.07\textwidth}p{0.12\textwidth}p{0.75\textwidth}@{}}\toprule Line&Kind&Name\\\midrule\endhead

13&structure&\nolinkurl{PaperXI.DyadicAffine.Refinement}\\

19&theorem&\nolinkurl{PaperXI.DyadicAffine.Refinement.ext}\\

25&theorem&\nolinkurl{PaperXI.DyadicAffine.Refinement.value_zero}\\

31&theorem&\nolinkurl{PaperXI.DyadicAffine.difference_add}\\

38&def&\nolinkurl{PaperXI.DyadicAffine.linearDifference}\\

43&theorem&\nolinkurl{PaperXI.DyadicAffine.linearDifference_apply}\\

47&def&\nolinkurl{PaperXI.DyadicAffine.addLinear}\\

53&theorem&\nolinkurl{PaperXI.DyadicAffine.addLinear_apply}\\

56&theorem&\nolinkurl{PaperXI.DyadicAffine.addLinear_linearDifference}\\

61&theorem&\nolinkurl{PaperXI.DyadicAffine.linearDifference_addLinear}\\

67&def&\nolinkurl{PaperXI.DyadicAffine.refinementEquiv}\\

74&theorem&\nolinkurl{PaperXI.DyadicAffine.exists_unique_linear_difference}\\

\bottomrule\end{longtable}

\section{\texorpdfstring{\nolinkurl{PaperXI.DyadicBasis}}{PaperXI.DyadicBasis}}\label{module:PaperXI.DyadicBasis}

\textbf{Selected layer:} v0.44 parent.\quad\textbf{Source lines:} 83.\par\textbf{SHA-256:} {\scriptsize\nolinkurl{18c632203cbd0c023aa7aaf6fcc10af54acb0271d09db3fa9a38d74308c5546f}}.

\VerbatimInput[fontsize=\scriptsize,breaklines=true,breakanywhere=true]{excerpts/D_DyadicBasis.txt}

\textbf{Direct project imports:} \nolinkurl{PaperXI.DyadicCoordinates}.

\begin{longtable}{@{}p{0.07\textwidth}p{0.12\textwidth}p{0.75\textwidth}@{}}\toprule Line&Kind&Name\\\midrule\endhead

14&def&\nolinkurl{PaperXI.DyadicBasis.toCoordinates}\\

22&def&\nolinkurl{PaperXI.DyadicBasis.fromCoordinates}\\

30&def&\nolinkurl{PaperXI.DyadicBasis.refinementEquiv}\\

42&theorem&\nolinkurl{PaperXI.DyadicBasis.toCoordinates_apply}\\

47&theorem&\nolinkurl{PaperXI.DyadicBasis.fromCoordinates_apply}\\

53&def&\nolinkurl{PaperXI.DyadicBasis.basisArf}\\

56&theorem&\nolinkurl{PaperXI.DyadicBasis.basisArf_eq_coordinate_arf}\\

62&theorem&\nolinkurl{PaperXI.DyadicBasis.gauss_sum}\\

74&theorem&\nolinkurl{PaperXI.DyadicBasis.basisArf_independent}\\

\bottomrule\end{longtable}

\section{\texorpdfstring{\nolinkurl{PaperXI.DyadicBasisConsequences}}{PaperXI.DyadicBasisConsequences}}\label{module:PaperXI.DyadicBasisConsequences}

\textbf{Selected layer:} v0.44 parent.\quad\textbf{Source lines:} 68.\par\textbf{SHA-256:} {\scriptsize\nolinkurl{d45fa0f60b8166b49c7962a6afc195f6e7cb92dccabb9ee3103ff488062a4031}}.

\VerbatimInput[fontsize=\scriptsize,breaklines=true,breakanywhere=true]{excerpts/D_DyadicBasisConsequences.txt}

\textbf{Direct project imports:} \nolinkurl{PaperXI.DyadicBasis}, \nolinkurl{PaperXI.DyadicRefinement}.

\begin{longtable}{@{}p{0.07\textwidth}p{0.12\textwidth}p{0.75\textwidth}@{}}\toprule Line&Kind&Name\\\midrule\endhead

16&theorem&\nolinkurl{PaperXI.DyadicBasisConsequences.toCoordinates_addLinear}\\

25&theorem&\nolinkurl{PaperXI.DyadicBasisConsequences.arf_translation}\\

42&theorem&\nolinkurl{PaperXI.DyadicBasisConsequences.basisArf_fromCoordinates}\\

48&theorem&\nolinkurl{PaperXI.DyadicBasisConsequences.gauss_fromCoordinates}\\

57&theorem&\nolinkurl{PaperXI.DyadicBasisConsequences.opposite_refinements}\\

\bottomrule\end{longtable}

\section{\texorpdfstring{\nolinkurl{PaperXI.DyadicBasisCounts}}{PaperXI.DyadicBasisCounts}}\label{module:PaperXI.DyadicBasisCounts}

\textbf{Selected layer:} v0.44 parent.\quad\textbf{Source lines:} 35.\par\textbf{SHA-256:} {\scriptsize\nolinkurl{2b9d123b889d163bf29be6dfbab3e510195d98591db8098fb4d773a4a05b6093}}.

\VerbatimInput[fontsize=\scriptsize,breaklines=true,breakanywhere=true]{excerpts/D_DyadicBasisCounts.txt}

\textbf{Direct project imports:} \nolinkurl{PaperXI.DyadicBasis}, \nolinkurl{PaperXI.DyadicRefinementCounts}.

\begin{longtable}{@{}p{0.07\textwidth}p{0.12\textwidth}p{0.75\textwidth}@{}}\toprule Line&Kind&Name\\\midrule\endhead

14&def&\nolinkurl{PaperXI.DyadicBasisCounts.arfFiberEquiv}\\

21&theorem&\nolinkurl{PaperXI.DyadicBasisCounts.even_count}\\

28&theorem&\nolinkurl{PaperXI.DyadicBasisCounts.odd_count}\\

\bottomrule\end{longtable}

\section{\texorpdfstring{\nolinkurl{PaperXI.DyadicCoordinates}}{PaperXI.DyadicCoordinates}}\label{module:PaperXI.DyadicCoordinates}

\textbf{Selected layer:} v0.44 parent.\quad\textbf{Source lines:} 162.\par\textbf{SHA-256:} {\scriptsize\nolinkurl{fc1a2a3a6569e264cfddcb51fd16e91c22ff91ed4a741634239798f18cd83004}}.

\VerbatimInput[fontsize=\scriptsize,breaklines=true,breakanywhere=true]{excerpts/D_DyadicCoordinates.txt}

\textbf{Direct project imports:} \nolinkurl{PaperXI.DyadicAffine}, \nolinkurl{PaperXI.DyadicGauss}.

\begin{longtable}{@{}p{0.07\textwidth}p{0.12\textwidth}p{0.75\textwidth}@{}}\toprule Line&Kind&Name\\\midrule\endhead

15&def&\nolinkurl{PaperXI.DyadicCoordinates.symplectic}\\

56&theorem&\nolinkurl{PaperXI.DyadicCoordinates.symplectic_apply}\\

59&theorem&\nolinkurl{PaperXI.DyadicCoordinates.quadratic_polar}\\

67&def&\nolinkurl{PaperXI.DyadicCoordinates.coordinateRefinement}\\

71&def&\nolinkurl{PaperXI.DyadicCoordinates.firstAxis}\\

72&def&\nolinkurl{PaperXI.DyadicCoordinates.secondAxis}\\

75&theorem&\nolinkurl{PaperXI.DyadicCoordinates.single_axes}\\

79&def&\nolinkurl{PaperXI.DyadicCoordinates.linearCoefficients}\\

82&theorem&\nolinkurl{PaperXI.DyadicCoordinates.linear_coordinate_formula}\\

93&theorem&\nolinkurl{PaperXI.DyadicCoordinates.quadratic_firstAxis}\\

102&theorem&\nolinkurl{PaperXI.DyadicCoordinates.quadratic_secondAxis}\\

111&def&\nolinkurl{PaperXI.DyadicCoordinates.refinementCoefficients}\\

114&theorem&\nolinkurl{PaperXI.DyadicCoordinates.refinement_coordinate_formula}\\

125&theorem&\nolinkurl{PaperXI.DyadicCoordinates.refinementCoefficients_coordinateRefinement}\\

130&def&\nolinkurl{PaperXI.DyadicCoordinates.coordinateEquiv}\\

138&theorem&\nolinkurl{PaperXI.DyadicCoordinates.refinementCoefficients_apply}\\

145&theorem&\nolinkurl{PaperXI.DyadicCoordinates.arf_basis_formula}\\

149&theorem&\nolinkurl{PaperXI.DyadicCoordinates.refinement_arf_basis_formula}\\

155&theorem&\nolinkurl{PaperXI.DyadicCoordinates.refinement_gauss_sum}\\

\bottomrule\end{longtable}

\section{\texorpdfstring{\nolinkurl{PaperXI.DyadicCounts}}{PaperXI.DyadicCounts}}\label{module:PaperXI.DyadicCounts}

\textbf{Selected layer:} v0.44 parent.\quad\textbf{Source lines:} 135.\par\textbf{SHA-256:} {\scriptsize\nolinkurl{b1155d1f3103b197ed16cea0c4c297e29da3f6e917e6006f486e47727af34b5d}}.

\VerbatimInput[fontsize=\scriptsize,breaklines=true,breakanywhere=true]{excerpts/D_DyadicCounts.txt}

\textbf{Direct project imports:} \nolinkurl{PaperXI.DyadicGauss}.

\begin{longtable}{@{}p{0.07\textwidth}p{0.12\textwidth}p{0.75\textwidth}@{}}\toprule Line&Kind&Name\\\midrule\endhead

13&theorem&\nolinkurl{PaperXI.DyadicCounts.signed_sum_eq_card_difference}\\

19&theorem&\nolinkurl{PaperXI.DyadicCounts.card_partition}\\

27&def&\nolinkurl{PaperXI.DyadicCounts.evenCount}\\

30&def&\nolinkurl{PaperXI.DyadicCounts.oddCount}\\

33&theorem&\nolinkurl{PaperXI.DyadicCounts.f2_ne_zero_iff_eq_one}\\

47&theorem&\nolinkurl{PaperXI.DyadicCounts.parameter_card}\\

50&theorem&\nolinkurl{PaperXI.DyadicCounts.even_add_odd}\\

54&theorem&\nolinkurl{PaperXI.DyadicCounts.even_sub_odd}\\

63&theorem&\nolinkurl{PaperXI.DyadicCounts.twice_even}\\

71&theorem&\nolinkurl{PaperXI.DyadicCounts.twice_odd}\\

78&theorem&\nolinkurl{PaperXI.DyadicCounts.even_factored}\\

90&theorem&\nolinkurl{PaperXI.DyadicCounts.odd_factored}\\

102&theorem&\nolinkurl{PaperXI.DyadicCounts.counts_rank_zero}\\

110&theorem&\nolinkurl{PaperXI.DyadicCounts.even_factored_rational}\\

123&theorem&\nolinkurl{PaperXI.DyadicCounts.odd_factored_rational}\\

\bottomrule\end{longtable}

\section{\texorpdfstring{\nolinkurl{PaperXI.DyadicDimension}}{PaperXI.DyadicDimension}}\label{module:PaperXI.DyadicDimension}

\textbf{Selected layer:} v0.44 parent.\quad\textbf{Source lines:} 49.\par\textbf{SHA-256:} {\scriptsize\nolinkurl{295c12f1fe30abb310b8e54856cc51fb1f11bee0bd9f7d4258386bd30271944a}}.

\VerbatimInput[fontsize=\scriptsize,breaklines=true,breakanywhere=true]{excerpts/D_DyadicDimension.txt}

\textbf{Direct project imports:} \nolinkurl{PaperXI.DyadicBasis}.

\begin{longtable}{@{}p{0.07\textwidth}p{0.12\textwidth}p{0.75\textwidth}@{}}\toprule Line&Kind&Name\\\midrule\endhead

15&theorem&\nolinkurl{PaperXI.DyadicDimension.finrank_eq_twice_card}\\

20&theorem&\nolinkurl{PaperXI.DyadicDimension.card_eq_finrank_div_two}\\

25&theorem&\nolinkurl{PaperXI.DyadicDimension.card_pos}\\

32&theorem&\nolinkurl{PaperXI.DyadicDimension.coordinate_card_eq}\\

38&theorem&\nolinkurl{PaperXI.DyadicDimension.basisArf_independent}\\

\bottomrule\end{longtable}

\section{\texorpdfstring{\nolinkurl{PaperXI.DyadicExtensionClass}}{PaperXI.DyadicExtensionClass}}\label{module:PaperXI.DyadicExtensionClass}

\textbf{Selected layer:} v0.44 parent.\quad\textbf{Source lines:} 93.\par\textbf{SHA-256:} {\scriptsize\nolinkurl{af6825148cbf10fa7e036416c6b158a32f9be04ed3a83f3767394273839b2656}}.

\VerbatimInput[fontsize=\scriptsize,breaklines=true,breakanywhere=true]{excerpts/D_DyadicExtensionClass.txt}

\textbf{Direct project imports:} \nolinkurl{PaperXI.DyadicHeisenbergLift}.

\begin{longtable}{@{}p{0.07\textwidth}p{0.12\textwidth}p{0.75\textwidth}@{}}\toprule Line&Kind&Name\\\midrule\endhead

14&def&\nolinkurl{PaperXI.DyadicExtensionClass.upper}\\

18&theorem&\nolinkurl{PaperXI.DyadicExtensionClass.upper_basis}\\

24&theorem&\nolinkurl{PaperXI.DyadicExtensionClass.upper_symmetrization}\\

45&theorem&\nolinkurl{PaperXI.DyadicExtensionClass.exists_coboundary}\\

60&theorem&\nolinkurl{PaperXI.DyadicExtensionClass.exists_equiv_of_diagonal}\\

74&theorem&\nolinkurl{PaperXI.DyadicExtensionClass.exists_equiv_iff_diagonal}\\

85&theorem&\nolinkurl{PaperXI.DyadicExtensionClass.refinement_model_independent}\\

\bottomrule\end{longtable}

\section{\texorpdfstring{\nolinkurl{PaperXI.DyadicGauss}}{PaperXI.DyadicGauss}}\label{module:PaperXI.DyadicGauss}

\textbf{Selected layer:} v0.44 parent.\quad\textbf{Source lines:} 60.\par\textbf{SHA-256:} {\scriptsize\nolinkurl{634fe67a476f515f627455fa43a1ec870a4ccc66b87867cc5d40ee1227c91dea}}.

\VerbatimInput[fontsize=\scriptsize,breaklines=true,breakanywhere=true]{excerpts/D_DyadicGauss.txt}

\textbf{Direct project imports:} none.

\begin{longtable}{@{}p{0.07\textwidth}p{0.12\textwidth}p{0.75\textwidth}@{}}\toprule Line&Kind&Name\\\midrule\endhead

10&abbrev&\nolinkurl{PaperXI.Dyadic.F2}\\

12&def&\nolinkurl{PaperXI.Dyadic.sign}\\

14&theorem&\nolinkurl{PaperXI.Dyadic.sign_zero}\\

15&theorem&\nolinkurl{PaperXI.Dyadic.sign_one}\\

17&theorem&\nolinkurl{PaperXI.Dyadic.sign_add}\\

20&theorem&\nolinkurl{PaperXI.Dyadic.sign_injective}\\

30&theorem&\nolinkurl{PaperXI.Dyadic.sign_sum}\\

37&def&\nolinkurl{PaperXI.Dyadic.plane}\\

39&theorem&\nolinkurl{PaperXI.Dyadic.plane_gauss}\\

47&def&\nolinkurl{PaperXI.Dyadic.quadratic}\\

49&def&\nolinkurl{PaperXI.Dyadic.arf}\\

51&theorem&\nolinkurl{PaperXI.Dyadic.gauss_sum}\\

\bottomrule\end{longtable}

\section{\texorpdfstring{\nolinkurl{PaperXI.DyadicHeisenbergLift}}{PaperXI.DyadicHeisenbergLift}}\label{module:PaperXI.DyadicHeisenbergLift}

\textbf{Selected layer:} v0.44 parent.\quad\textbf{Source lines:} 155.\par\textbf{SHA-256:} {\scriptsize\nolinkurl{916986a499526e1343b2b63906fe133ac8db3cb95e6cbbbe31d6ada8e17229b1}}.

\VerbatimInput[fontsize=\scriptsize,breaklines=true,breakanywhere=true]{excerpts/D_DyadicHeisenbergLift.txt}

\textbf{Direct project imports:} \nolinkurl{PaperXI.HeisenbergCochain}, \nolinkurl{PaperXI.SymplecticCoordinates}.

\begin{longtable}{@{}p{0.07\textwidth}p{0.12\textwidth}p{0.75\textwidth}@{}}\toprule Line&Kind&Name\\\midrule\endhead

15&def&\nolinkurl{PaperXI.DyadicHeisenbergLift.coordinateFactor}\\

65&theorem&\nolinkurl{PaperXI.DyadicHeisenbergLift.coordinateFactor_apply}\\

69&theorem&\nolinkurl{PaperXI.DyadicHeisenbergLift.coordinateFactor_commutator}\\

78&theorem&\nolinkurl{PaperXI.DyadicHeisenbergLift.coordinateFactor_diagonal}\\

91&def&\nolinkurl{PaperXI.DyadicHeisenbergLift.coordinateRefinementOf}\\

100&def&\nolinkurl{PaperXI.DyadicHeisenbergLift.factorOfCoordinates}\\

106&theorem&\nolinkurl{PaperXI.DyadicHeisenbergLift.factorOfCoordinates_commutator}\\

113&theorem&\nolinkurl{PaperXI.DyadicHeisenbergLift.factorOfCoordinates_diagonal}\\

125&theorem&\nolinkurl{PaperXI.DyadicHeisenbergLift.exists_factor}\\

135&def&\nolinkurl{PaperXI.DyadicHeisenbergLift.factor}\\

139&theorem&\nolinkurl{PaperXI.DyadicHeisenbergLift.factor_commutator}\\

144&theorem&\nolinkurl{PaperXI.DyadicHeisenbergLift.factor_diagonal}\\

150&theorem&\nolinkurl{PaperXI.DyadicHeisenbergLift.square_profile}\\

\bottomrule\end{longtable}

\section{\texorpdfstring{\nolinkurl{PaperXI.DyadicIsometry}}{PaperXI.DyadicIsometry}}\label{module:PaperXI.DyadicIsometry}

\textbf{Selected layer:} v0.44 parent.\quad\textbf{Source lines:} 52.\par\textbf{SHA-256:} {\scriptsize\nolinkurl{48d9678ba58d442aa451540502f01d5e1fb8a3257e8bdf4d64fef974305c651f}}.

\VerbatimInput[fontsize=\scriptsize,breaklines=true,breakanywhere=true]{excerpts/D_DyadicIsometry.txt}

\textbf{Direct project imports:} \nolinkurl{PaperXI.DyadicDimension}.

\begin{longtable}{@{}p{0.07\textwidth}p{0.12\textwidth}p{0.75\textwidth}@{}}\toprule Line&Kind&Name\\\midrule\endhead

14&def&\nolinkurl{PaperXI.DyadicIsometry.transportRefinement}\\

21&theorem&\nolinkurl{PaperXI.DyadicIsometry.transportRefinement_apply}\\

25&def&\nolinkurl{PaperXI.DyadicIsometry.transportRefinementEquiv}\\

33&theorem&\nolinkurl{PaperXI.DyadicIsometry.gauss_sum_transport}\\

41&theorem&\nolinkurl{PaperXI.DyadicIsometry.basisArf_transport}\\

\bottomrule\end{longtable}

\section{\texorpdfstring{\nolinkurl{PaperXI.DyadicOrthogonal}}{PaperXI.DyadicOrthogonal}}\label{module:PaperXI.DyadicOrthogonal}

\textbf{Selected layer:} v0.44 parent.\quad\textbf{Source lines:} 76.\par\textbf{SHA-256:} {\scriptsize\nolinkurl{3ec83d3eeac7e69ef4986a573ba818eb9cfe33282530bb3e8f85d7ed4323d671}}.

\textbf{Direct project imports:} \nolinkurl{PaperXI.DyadicBasis}, \nolinkurl{PaperXI.DyadicTranslation}.

\begin{longtable}{@{}p{0.07\textwidth}p{0.12\textwidth}p{0.75\textwidth}@{}}\toprule Line&Kind&Name\\\midrule\endhead

14&def&\nolinkurl{PaperXI.DyadicOrthogonal.form}\\

18&theorem&\nolinkurl{PaperXI.DyadicOrthogonal.form_apply}\\

20&def&\nolinkurl{PaperXI.DyadicOrthogonal.refinement}\\

28&theorem&\nolinkurl{PaperXI.DyadicOrthogonal.gauss_product}\\

35&def&\nolinkurl{PaperXI.DyadicOrthogonal.coordinates}\\

40&theorem&\nolinkurl{PaperXI.DyadicOrthogonal.coordinates_form}\\

48&theorem&\nolinkurl{PaperXI.DyadicOrthogonal.arf_add}\\

67&theorem&\nolinkurl{PaperXI.DyadicOrthogonal.normalized_sign_product}\\

\bottomrule\end{longtable}

\section{\texorpdfstring{\nolinkurl{PaperXI.DyadicRefinement}}{PaperXI.DyadicRefinement}}\label{module:PaperXI.DyadicRefinement}

\textbf{Selected layer:} v0.44 parent.\quad\textbf{Source lines:} 48.\par\textbf{SHA-256:} {\scriptsize\nolinkurl{4056c45a4a85c4fbaf44d651ce11f83f50656138b52fb6f18d95e9b8eafa01e5}}.

\VerbatimInput[fontsize=\scriptsize,breaklines=true,breakanywhere=true]{excerpts/D_DyadicRefinement.txt}

\textbf{Direct project imports:} \nolinkurl{PaperXI.DyadicCoordinates}, \nolinkurl{PaperXI.DyadicTranslation}.

\begin{longtable}{@{}p{0.07\textwidth}p{0.12\textwidth}p{0.75\textwidth}@{}}\toprule Line&Kind&Name\\\midrule\endhead

13&def&\nolinkurl{PaperXI.DyadicRefinement.basisArf}\\

16&theorem&\nolinkurl{PaperXI.DyadicRefinement.basisArf_coordinateRefinement}\\

20&theorem&\nolinkurl{PaperXI.DyadicRefinement.translation_refinement}\\

29&theorem&\nolinkurl{PaperXI.DyadicRefinement.arf_translation}\\

36&theorem&\nolinkurl{PaperXI.DyadicRefinement.opposite_refinements}\\

45&theorem&\nolinkurl{PaperXI.DyadicRefinement.sign_eq_neg_one_pow}\\

\bottomrule\end{longtable}

\section{\texorpdfstring{\nolinkurl{PaperXI.DyadicRefinementCounts}}{PaperXI.DyadicRefinementCounts}}\label{module:PaperXI.DyadicRefinementCounts}

\textbf{Selected layer:} v0.44 parent.\quad\textbf{Source lines:} 55.\par\textbf{SHA-256:} {\scriptsize\nolinkurl{5d9cdc46b7bde9d4f7ceed89b2a7a68ec38f36fee3142300f40b6a8d5301838b}}.

\VerbatimInput[fontsize=\scriptsize,breaklines=true,breakanywhere=true]{excerpts/D_DyadicRefinementCounts.txt}

\textbf{Direct project imports:} \nolinkurl{PaperXI.DyadicCounts}, \nolinkurl{PaperXI.DyadicCoordinates}.

\begin{longtable}{@{}p{0.07\textwidth}p{0.12\textwidth}p{0.75\textwidth}@{}}\toprule Line&Kind&Name\\\midrule\endhead

13&instance&\nolinkurl{PaperXI.DyadicRefinementCounts.refinementFintype}\\

16&def&\nolinkurl{PaperXI.DyadicRefinementCounts.basisArf}\\

19&theorem&\nolinkurl{PaperXI.DyadicRefinementCounts.basisArf_coordinate}\\

24&def&\nolinkurl{PaperXI.DyadicRefinementCounts.arfFiberEquiv}\\

29&theorem&\nolinkurl{PaperXI.DyadicRefinementCounts.even_fiber_card}\\

35&theorem&\nolinkurl{PaperXI.DyadicRefinementCounts.odd_fiber_card}\\

42&theorem&\nolinkurl{PaperXI.DyadicRefinementCounts.even_refinement_count}\\

49&theorem&\nolinkurl{PaperXI.DyadicRefinementCounts.odd_refinement_count}\\

\bottomrule\end{longtable}

\section{\texorpdfstring{\nolinkurl{PaperXI.DyadicTorsor}}{PaperXI.DyadicTorsor}}\label{module:PaperXI.DyadicTorsor}

\textbf{Selected layer:} v0.44 parent.\quad\textbf{Source lines:} 70.\par\textbf{SHA-256:} {\scriptsize\nolinkurl{da8d7de85105943a07b9492568783dfc1df76d3789bb3444d344259b0ed8bb6b}}.

\VerbatimInput[fontsize=\scriptsize,breaklines=true,breakanywhere=true]{excerpts/D_DyadicTorsor.txt}

\textbf{Direct project imports:} \nolinkurl{PaperXI.DyadicAffine}.

\begin{longtable}{@{}p{0.07\textwidth}p{0.12\textwidth}p{0.75\textwidth}@{}}\toprule Line&Kind&Name\\\midrule\endhead

14&theorem&\nolinkurl{PaperXI.DyadicTorsor.addLinear_zero}\\

18&theorem&\nolinkurl{PaperXI.DyadicTorsor.addLinear_add}\\

24&theorem&\nolinkurl{PaperXI.DyadicTorsor.addLinear_injective}\\

30&theorem&\nolinkurl{PaperXI.DyadicTorsor.addLinear_free}\\

39&theorem&\nolinkurl{PaperXI.DyadicTorsor.addLinear_transitive}\\

44&instance&\nolinkurl{PaperXI.DyadicTorsor.instAddAction}\\

49&instance&\nolinkurl{PaperXI.DyadicTorsor.instVSub}\\

54&instance&\nolinkurl{PaperXI.DyadicTorsor.instAddTorsor}\\

61&abbrev&\nolinkurl{PaperXI.DyadicTorsor.canonicalTorsor}\\

64&theorem&\nolinkurl{PaperXI.DyadicTorsor.vadd_apply}\\

67&theorem&\nolinkurl{PaperXI.DyadicTorsor.vsub_apply}\\

\bottomrule\end{longtable}

\section{\texorpdfstring{\nolinkurl{PaperXI.DyadicTranslation}}{PaperXI.DyadicTranslation}}\label{module:PaperXI.DyadicTranslation}

\textbf{Selected layer:} v0.44 parent.\quad\textbf{Source lines:} 89.\par\textbf{SHA-256:} {\scriptsize\nolinkurl{098a6b0ee464537aaea2cef020e535b28ba9572b8079134076fdd5ee6fb54969}}.

\textbf{Direct project imports:} \nolinkurl{PaperXI.DyadicGauss}.

\begin{longtable}{@{}p{0.07\textwidth}p{0.12\textwidth}p{0.75\textwidth}@{}}\toprule Line&Kind&Name\\\midrule\endhead

9&def&\nolinkurl{PaperXI.Dyadic.pairing}\\

12&def&\nolinkurl{PaperXI.Dyadic.translate}\\

15&theorem&\nolinkurl{PaperXI.Dyadic.quadratic_translate}\\

23&theorem&\nolinkurl{PaperXI.Dyadic.arf_translate}\\

31&theorem&\nolinkurl{PaperXI.Dyadic.gauss_translate}\\

37&theorem&\nolinkurl{PaperXI.Dyadic.arf_orthogonal}\\

41&theorem&\nolinkurl{PaperXI.Dyadic.gauss_orthogonal}\\

57&theorem&\nolinkurl{PaperXI.Dyadic.sign_orthogonal}\\

61&theorem&\nolinkurl{PaperXI.Dyadic.arf_zero}\\

64&def&\nolinkurl{PaperXI.Dyadic.anisotropic}\\

66&theorem&\nolinkurl{PaperXI.Dyadic.arf_anisotropic}\\

76&theorem&\nolinkurl{PaperXI.Dyadic.opposite_gauss}\\

82&theorem&\nolinkurl{PaperXI.Dyadic.both_arf}\\

\bottomrule\end{longtable}

\section{\texorpdfstring{\nolinkurl{PaperXI.ElementaryEnvelope}}{PaperXI.ElementaryEnvelope}}\label{module:PaperXI.ElementaryEnvelope}

\textbf{Selected layer:} v0.44 parent.\quad\textbf{Source lines:} 139.\par\textbf{SHA-256:} {\scriptsize\nolinkurl{76c36ba11aabeef9351becc4f99f59db800649c73d1dc0f0ff57d59bc7e5d6c4}}.

\VerbatimInput[fontsize=\scriptsize,breaklines=true,breakanywhere=true]{excerpts/D_ElementaryEnvelope.txt}

\textbf{Direct project imports:} \nolinkurl{PaperXI.CocycleQuotient}.

\begin{longtable}{@{}p{0.07\textwidth}p{0.12\textwidth}p{0.75\textwidth}@{}}\toprule Line&Kind&Name\\\midrule\endhead

18&theorem&\nolinkurl{PaperXI.ElementaryEnvelope.commutator_pow}\\

25&abbrev&\nolinkurl{PaperXI.ElementaryEnvelope.groupModule}\\

28&theorem&\nolinkurl{PaperXI.ElementaryEnvelope.finiteDimensional_group}\\

36&def&\nolinkurl{PaperXI.ElementaryEnvelope.Space}\\

39&theorem&\nolinkurl{PaperXI.ElementaryEnvelope.quotient_pow}\\

45&instance&\nolinkurl{PaperXI.ElementaryEnvelope.spaceModule}\\

50&instance&\nolinkurl{PaperXI.ElementaryEnvelope.finiteSpace}\\

53&instance&\nolinkurl{PaperXI.ElementaryEnvelope.finiteDimensionalSpace}\\

58&def&\nolinkurl{PaperXI.ElementaryEnvelope.rootCharacter}\\

69&def&\nolinkurl{PaperXI.ElementaryEnvelope.rootEquiv}\\

76&theorem&\nolinkurl{PaperXI.ElementaryEnvelope.rootEquiv_coe}\\

80&def&\nolinkurl{PaperXI.ElementaryEnvelope.rootLog}\\

83&theorem&\nolinkurl{PaperXI.ElementaryEnvelope.rootLog_character}\\

86&theorem&\nolinkurl{PaperXI.ElementaryEnvelope.character_rootLog}\\

90&theorem&\nolinkurl{PaperXI.ElementaryEnvelope.quotientPairing_pow}\\

97&def&\nolinkurl{PaperXI.ElementaryEnvelope.rootPairing}\\

115&theorem&\nolinkurl{PaperXI.ElementaryEnvelope.rootPairing_coe}\\

119&def&\nolinkurl{PaperXI.ElementaryEnvelope.rootCommutator}\\

133&theorem&\nolinkurl{PaperXI.ElementaryEnvelope.rootCommutator_coe}\\

136&theorem&\nolinkurl{PaperXI.ElementaryEnvelope.rootCommutator_self}\\

\bottomrule\end{longtable}

\section{\texorpdfstring{\nolinkurl{PaperXI.ElementaryEnvelopeForm}}{PaperXI.ElementaryEnvelopeForm}}\label{module:PaperXI.ElementaryEnvelopeForm}

\textbf{Selected layer:} v0.44 parent.\quad\textbf{Source lines:} 137.\par\textbf{SHA-256:} {\scriptsize\nolinkurl{8bfa652c46358cbe6ee15446b6d435521aeac0f75cc96a605a5fb41b3bf5bc6f}}.

\VerbatimInput[fontsize=\scriptsize,breaklines=true,breakanywhere=true]{excerpts/D_ElementaryEnvelopeForm.txt}

\textbf{Direct project imports:} \nolinkurl{PaperXI.ElementaryEnvelope}.

\begin{longtable}{@{}p{0.07\textwidth}p{0.12\textwidth}p{0.75\textwidth}@{}}\toprule Line&Kind&Name\\\midrule\endhead

15&def&\nolinkurl{PaperXI.ElementaryEnvelope.logarithmicPairing}\\

37&def&\nolinkurl{PaperXI.ElementaryEnvelope.formWith}\\

42&theorem&\nolinkurl{PaperXI.ElementaryEnvelope.formWith_apply}\\

47&theorem&\nolinkurl{PaperXI.ElementaryEnvelope.formWith_character}\\

54&theorem&\nolinkurl{PaperXI.ElementaryEnvelope.formWith_smul_left}\\

59&theorem&\nolinkurl{PaperXI.ElementaryEnvelope.formWith_smul_right}\\

64&theorem&\nolinkurl{PaperXI.ElementaryEnvelope.formWith_isAlt}\\

73&theorem&\nolinkurl{PaperXI.ElementaryEnvelope.formWith_eq_zero_iff}\\

80&theorem&\nolinkurl{PaperXI.ElementaryEnvelope.formWith_nondegenerate}\\

93&def&\nolinkurl{PaperXI.ElementaryEnvelope.form}\\

95&theorem&\nolinkurl{PaperXI.ElementaryEnvelope.form_isAlt}\\

97&theorem&\nolinkurl{PaperXI.ElementaryEnvelope.form_nondegenerate}\\

101&theorem&\nolinkurl{PaperXI.ElementaryEnvelope.form_stdAddChar}\\

109&theorem&\nolinkurl{PaperXI.ElementaryEnvelope.choice_scalar_ne_zero}\\

115&theorem&\nolinkurl{PaperXI.ElementaryEnvelope.choice_map_eq_mul}\\

124&theorem&\nolinkurl{PaperXI.ElementaryEnvelope.formWith_choice}\\

131&theorem&\nolinkurl{PaperXI.ElementaryEnvelope.exists_choice_scalar}\\

\bottomrule\end{longtable}

\section{\texorpdfstring{\nolinkurl{PaperXI.ElementaryEnvelopeTransport}}{PaperXI.ElementaryEnvelopeTransport}}\label{module:PaperXI.ElementaryEnvelopeTransport}

\textbf{Selected layer:} v0.44 parent.\quad\textbf{Source lines:} 148.\par\textbf{SHA-256:} {\scriptsize\nolinkurl{56b022844ede8e64c96a72cadf7393135d60c5c9601cd457d49f95322e3fcf18}}.

\VerbatimInput[fontsize=\scriptsize,breaklines=true,breakanywhere=true]{excerpts/D_ElementaryEnvelopeTransport.txt}

\textbf{Direct project imports:} \nolinkurl{PaperXI.ElementaryEnvelopeForm}.

\begin{longtable}{@{}p{0.07\textwidth}p{0.12\textwidth}p{0.75\textwidth}@{}}\toprule Line&Kind&Name\\\midrule\endhead

15&theorem&\nolinkurl{PaperXI.ElementaryEnvelope.radical_map}\\

34&def&\nolinkurl{PaperXI.ElementaryEnvelope.quotientEquiv}\\

39&theorem&\nolinkurl{PaperXI.ElementaryEnvelope.quotientEquiv_mk}\\

43&theorem&\nolinkurl{PaperXI.ElementaryEnvelope.quotientEquiv_pairing}\\

56&def&\nolinkurl{PaperXI.ElementaryEnvelope.linearEquiv}\\

62&theorem&\nolinkurl{PaperXI.ElementaryEnvelope.linearEquiv_mk}\\

68&theorem&\nolinkurl{PaperXI.ElementaryEnvelope.linearEquiv_formWith}\\

79&theorem&\nolinkurl{PaperXI.ElementaryEnvelope.linearEquiv_form}\\

85&theorem&\nolinkurl{PaperXI.ElementaryEnvelope.space_induction_on}\\

89&theorem&\nolinkurl{PaperXI.ElementaryEnvelope.linearEquiv_refl}\\

98&theorem&\nolinkurl{PaperXI.ElementaryEnvelope.linearEquiv_trans}\\

113&def&\nolinkurl{PaperXI.ElementaryEnvelope.action}\\

135&theorem&\nolinkurl{PaperXI.ElementaryEnvelope.action_mk}\\

142&theorem&\nolinkurl{PaperXI.ElementaryEnvelope.action_form}\\

\bottomrule\end{longtable}

\section{\texorpdfstring{\nolinkurl{PaperXI.EquivariantRepresentation113}}{PaperXI.EquivariantRepresentation113}}\label{module:PaperXI.EquivariantRepresentation113}

\textbf{Selected layer:} Continuation113.\quad\textbf{Source lines:} 117.\par\textbf{SHA-256:} {\scriptsize\nolinkurl{f1cf8ff3343f2647c0994270f6de5bc3478925a25406ea3a221861807b08b6c9}}.

\VerbatimInput[fontsize=\scriptsize,breaklines=true,breakanywhere=true]{excerpts/D_EquivariantRepresentation113.txt}

\textbf{Direct project imports:} \nolinkurl{PaperXI.ActualScalarLift113}.

\begin{longtable}{@{}p{0.07\textwidth}p{0.12\textwidth}p{0.75\textwidth}@{}}\toprule Line&Kind&Name\\\midrule\endhead

38&theorem&\nolinkurl{PaperXI.EquivariantRepresentation113.scalarTransport_covariant}\\

67&theorem&\nolinkurl{PaperXI.EquivariantRepresentation113.exists_equivariant_actual_action}\\

\bottomrule\end{longtable}

\section{\texorpdfstring{\nolinkurl{PaperXI.EquivariantSplitting68}}{PaperXI.EquivariantSplitting68}}\label{module:PaperXI.EquivariantSplitting68}

\textbf{Selected layer:} base108 arithmetic branch.\quad\textbf{Source lines:} 35.\par\textbf{SHA-256:} {\scriptsize\nolinkurl{561e6ca7bb701f95ce183d25a1b39b3dbf2a5d389924851f78373fa3f52e4460}}.

\VerbatimInput[fontsize=\scriptsize,breaklines=true,breakanywhere=true]{excerpts/D_EquivariantSplitting68.txt}

\textbf{Direct project imports:} \nolinkurl{PaperXI.SmoothHomOrbit68}, \nolinkurl{PaperXI.FiniteOrbitAverage68}.

\begin{longtable}{@{}p{0.07\textwidth}p{0.12\textwidth}p{0.75\textwidth}@{}}\toprule Line&Kind&Name\\\midrule\endhead

20&theorem&\nolinkurl{PaperXI.EquivariantSplitting68.exists_equivariant_section}\\

\bottomrule\end{longtable}

\section{\texorpdfstring{\nolinkurl{PaperXI.FiniteAverage}}{PaperXI.FiniteAverage}}\label{module:PaperXI.FiniteAverage}

\textbf{Selected layer:} v0.44 parent.\quad\textbf{Source lines:} 69.\par\textbf{SHA-256:} {\scriptsize\nolinkurl{e960b8793c2a7e76fef815c456c4e61db736706716741679c3d423053722c977}}.

\VerbatimInput[fontsize=\scriptsize,breaklines=true,breakanywhere=true]{excerpts/D_FiniteAverage.txt}

\textbf{Direct project imports:} none.

\begin{longtable}{@{}p{0.07\textwidth}p{0.12\textwidth}p{0.75\textwidth}@{}}\toprule Line&Kind&Name\\\midrule\endhead

11&def&\nolinkurl{PaperXI.FiniteAverage.average}\\

14&theorem&\nolinkurl{PaperXI.FiniteAverage.average_apply}\\

18&theorem&\nolinkurl{PaperXI.FiniteAverage.card_ne_zero}\\

21&theorem&\nolinkurl{PaperXI.FiniteAverage.invariant}\\

27&theorem&\nolinkurl{PaperXI.FiniteAverage.fixes}\\

34&theorem&\nolinkurl{PaperXI.FiniteAverage.idempotent}\\

39&theorem&\nolinkurl{PaperXI.FiniteAverage.range_eq}\\

51&theorem&\nolinkurl{PaperXI.FiniteAverage.self_adjoint}\\

62&theorem&\nolinkurl{PaperXI.FiniteAverage.residual_orthogonal}\\

\bottomrule\end{longtable}

\section{\texorpdfstring{\nolinkurl{PaperXI.FiniteBreakAssembly}}{PaperXI.FiniteBreakAssembly}}\label{module:PaperXI.FiniteBreakAssembly}

\textbf{Selected layer:} v0.44 parent.\quad\textbf{Source lines:} 175.\par\textbf{SHA-256:} {\scriptsize\nolinkurl{ecca3693f7733c7ceef597b6e821b111610c82d612fd66abfe754af0e5be434c}}.

\VerbatimInput[fontsize=\scriptsize,breaklines=true,breakanywhere=true]{excerpts/D_FiniteBreakAssembly.txt}

\textbf{Direct project imports:} \nolinkurl{PaperXI.BreakBridgeAlong}, \nolinkurl{PaperXI.CentralExtensionComparison}, \nolinkurl{PaperXI.TraceDiscrepancy}.

\begin{longtable}{@{}p{0.07\textwidth}p{0.12\textwidth}p{0.75\textwidth}@{}}\toprule Line&Kind&Name\\\midrule\endhead

22&structure&\nolinkurl{PaperXI.FiniteBreakAssembly.Data}\\

43&def&\nolinkurl{PaperXI.FiniteBreakAssembly.factor}\\

51&theorem&\nolinkurl{PaperXI.FiniteBreakAssembly.factor_intertwines}\\

60&def&\nolinkurl{PaperXI.FiniteBreakAssembly.tensorBridge}\\

63&def&\nolinkurl{PaperXI.FiniteBreakAssembly.sourceAction}\\

67&def&\nolinkurl{PaperXI.FiniteBreakAssembly.targetAction}\\

71&theorem&\nolinkurl{PaperXI.FiniteBreakAssembly.tensor_intertwines}\\

79&theorem&\nolinkurl{PaperXI.FiniteBreakAssembly.source_finiteDimensional}\\

84&theorem&\nolinkurl{PaperXI.FiniteBreakAssembly.source_irreducible}\\

96&def&\nolinkurl{PaperXI.FiniteBreakAssembly.discrepancy}\\

102&theorem&\nolinkurl{PaperXI.FiniteBreakAssembly.discrepancy_law}\\

111&def&\nolinkurl{PaperXI.FiniteBreakAssembly.equivOfPhase}\\

132&theorem&\nolinkurl{PaperXI.FiniteBreakAssembly.discrepancy_and_closure}\\

151&def&\nolinkurl{PaperXI.FiniteBreakAssembly.equivOfTraces}\\

169&theorem&\nolinkurl{PaperXI.FiniteBreakAssembly.equivOfTraces_map}\\

\bottomrule\end{longtable}

\section{\texorpdfstring{\nolinkurl{PaperXI.FiniteFieldAverage}}{PaperXI.FiniteFieldAverage}}\label{module:PaperXI.FiniteFieldAverage}

\textbf{Selected layer:} v0.44 parent.\quad\textbf{Source lines:} 186.\par\textbf{SHA-256:} {\scriptsize\nolinkurl{da2bd3e2ade03fa988a3b50efe5b53ef4ed47f3636c86fc5f6841cc8b868f81d}}.

\VerbatimInput[fontsize=\scriptsize,breaklines=true,breakanywhere=true]{excerpts/D_FiniteFieldAverage.txt}

\textbf{Direct project imports:} none.

\begin{longtable}{@{}p{0.07\textwidth}p{0.12\textwidth}p{0.75\textwidth}@{}}\toprule Line&Kind&Name\\\midrule\endhead

16&def&\nolinkurl{PaperXI.FiniteFieldAverage.difference}\\

19&def&\nolinkurl{PaperXI.FiniteFieldAverage.average}\\

23&def&\nolinkurl{PaperXI.FiniteFieldAverage.companion}\\

26&theorem&\nolinkurl{PaperXI.FiniteFieldAverage.difference_apply}\\

28&theorem&\nolinkurl{PaperXI.FiniteFieldAverage.average_apply}\\

32&theorem&\nolinkurl{PaperXI.FiniteFieldAverage.companion_apply}\\

37&theorem&\nolinkurl{PaperXI.FiniteFieldAverage.weighted_identity}\\

50&theorem&\nolinkurl{PaperXI.FiniteFieldAverage.power_eq_one}\\

58&theorem&\nolinkurl{PaperXI.FiniteFieldAverage.average_mul_difference}\\

65&theorem&\nolinkurl{PaperXI.FiniteFieldAverage.difference_mul_average}\\

72&theorem&\nolinkurl{PaperXI.FiniteFieldAverage.moving_identity}\\

79&theorem&\nolinkurl{PaperXI.FiniteFieldAverage.power_fixes}\\

86&theorem&\nolinkurl{PaperXI.FiniteFieldAverage.average_fixes}\\

94&theorem&\nolinkurl{PaperXI.FiniteFieldAverage.range_average}\\

104&theorem&\nolinkurl{PaperXI.FiniteFieldAverage.ker_average}\\

117&theorem&\nolinkurl{PaperXI.FiniteFieldAverage.idempotent}\\

124&theorem&\nolinkurl{PaperXI.FiniteFieldAverage.isCompl}\\

131&def&\nolinkurl{PaperXI.FiniteFieldAverage.decomposition}\\

136&theorem&\nolinkurl{PaperXI.FiniteFieldAverage.decomposition_apply}\\

142&theorem&\nolinkurl{PaperXI.FiniteFieldAverage.decomposition_symm_fst}\\

154&theorem&\nolinkurl{PaperXI.FiniteFieldAverage.decomposition_symm_snd}\\

162&theorem&\nolinkurl{PaperXI.FiniteFieldAverage.average_mul_action}\\

170&theorem&\nolinkurl{PaperXI.FiniteFieldAverage.action_mul_average}\\

178&theorem&\nolinkurl{PaperXI.FiniteFieldAverage.average_independent}\\

\bottomrule\end{longtable}

\section{\texorpdfstring{\nolinkurl{PaperXI.FiniteFieldOrbitCount}}{PaperXI.FiniteFieldOrbitCount}}\label{module:PaperXI.FiniteFieldOrbitCount}

\textbf{Selected layer:} v0.44 parent.\quad\textbf{Source lines:} 93.\par\textbf{SHA-256:} {\scriptsize\nolinkurl{624a4a0234ebd63e8062471cb1439ed52dc73ed82e5d8f6400549c16196dde42}}.

\VerbatimInput[fontsize=\scriptsize,breaklines=true,breakanywhere=true]{excerpts/D_FiniteFieldOrbitCount.txt}

\textbf{Direct project imports:} \nolinkurl{PaperXI.CyclicFrobeniusCount}.

\begin{longtable}{@{}p{0.07\textwidth}p{0.12\textwidth}p{0.75\textwidth}@{}}\toprule Line&Kind&Name\\\midrule\endhead

13&instance&\nolinkurl{PaperXI.FiniteFieldOrbitCount.complexEnoughRoots}\\

19&theorem&\nolinkurl{PaperXI.FiniteFieldOrbitCount.characterGroup_isCyclic}\\

25&theorem&\nolinkurl{PaperXI.FiniteFieldOrbitCount.character_card}\\

28&def&\nolinkurl{PaperXI.FiniteFieldOrbitCount.characterFrobenius}\\

33&theorem&\nolinkurl{PaperXI.FiniteFieldOrbitCount.characterFrobenius_apply}\\

41&theorem&\nolinkurl{PaperXI.FiniteFieldOrbitCount.fixed_character_count}\\

47&theorem&\nolinkurl{PaperXI.FiniteFieldOrbitCount.exact_character_count}\\

54&theorem&\nolinkurl{PaperXI.FiniteFieldOrbitCount.regular_character_orbits}\\

62&theorem&\nolinkurl{PaperXI.FiniteFieldOrbitCount.regular_numerator_divisible}\\

70&def&\nolinkurl{PaperXI.FiniteFieldOrbitCount.extensionFrobenius}\\

77&theorem&\nolinkurl{PaperXI.FiniteFieldOrbitCount.extension_regular_orbits}\\

88&theorem&\nolinkurl{PaperXI.FiniteFieldOrbitCount.extensionFrobenius_apply}\\

\bottomrule\end{longtable}

\section{\texorpdfstring{\nolinkurl{PaperXI.FiniteFreeCompletion}}{PaperXI.FiniteFreeCompletion}}\label{module:PaperXI.FiniteFreeCompletion}

\textbf{Selected layer:} v0.44 parent.\quad\textbf{Source lines:} 75.\par\textbf{SHA-256:} {\scriptsize\nolinkurl{218a5b9e952143a9e26b36d2bcf93529baca8ae45f986c76903854d7c18aeead}}.

\VerbatimInput[fontsize=\scriptsize,breaklines=true,breakanywhere=true]{excerpts/D_FiniteFreeCompletion.txt}

\textbf{Direct project imports:} \nolinkurl{PaperXI.ResidueExtensionRealization}.

\begin{longtable}{@{}p{0.07\textwidth}p{0.12\textwidth}p{0.75\textwidth}@{}}\toprule Line&Kind&Name\\\midrule\endhead

14&theorem&\nolinkurl{PaperXI.FiniteFreeCompletion.linear_congruence}\\

25&theorem&\nolinkurl{PaperXI.FiniteFreeCompletion.precomplete}\\

42&theorem&\nolinkurl{PaperXI.FiniteFreeCompletion.hausdorff}\\

53&theorem&\nolinkurl{PaperXI.FiniteFreeCompletion.complete}\\

65&theorem&\nolinkurl{PaperXI.CompleteResidueRealization.complete}\\

\bottomrule\end{longtable}

\section{\texorpdfstring{\nolinkurl{PaperXI.FiniteFunctionAverage}}{PaperXI.FiniteFunctionAverage}}\label{module:PaperXI.FiniteFunctionAverage}

\textbf{Selected layer:} v0.44 parent.\quad\textbf{Source lines:} 54.\par\textbf{SHA-256:} {\scriptsize\nolinkurl{a7d776b604ef43cade5a813a977a7ce80ebd96a2b7906759e3b9fc89c7420d01}}.

\VerbatimInput[fontsize=\scriptsize,breaklines=true,breakanywhere=true]{excerpts/D_FiniteFunctionAverage.txt}

\textbf{Direct project imports:} \nolinkurl{PaperXI.TranslationAverage}.

\begin{longtable}{@{}p{0.07\textwidth}p{0.12\textwidth}p{0.75\textwidth}@{}}\toprule Line&Kind&Name\\\midrule\endhead

11&theorem&\nolinkurl{PaperXI.FiniteFunctionAverage.pairing_average_left}\\

24&theorem&\nolinkurl{PaperXI.FiniteFunctionAverage.pairing_average_right}\\

36&theorem&\nolinkurl{PaperXI.FiniteFunctionAverage.self_adjoint}\\

46&theorem&\nolinkurl{PaperXI.FiniteFunctionAverage.residual_orthogonal}\\

\bottomrule\end{longtable}

\section{\texorpdfstring{\nolinkurl{PaperXI.FiniteIntersectionCriterion108}}{PaperXI.FiniteIntersectionCriterion108}}\label{module:PaperXI.FiniteIntersectionCriterion108}

\textbf{Selected layer:} base108 arithmetic branch.\quad\textbf{Source lines:} 30.\par\textbf{SHA-256:} {\scriptsize\nolinkurl{a428d271392d27755c20bb0b6dac61b71398d4cc6370fe0642c4e5ed41364f51}}.

\VerbatimInput[fontsize=\scriptsize,breaklines=true,breakanywhere=true]{excerpts/D_FiniteIntersectionCriterion108.txt}

\textbf{Direct project imports:} \nolinkurl{PaperXI.CompactIntersection108}.

\begin{longtable}{@{}p{0.07\textwidth}p{0.12\textwidth}p{0.75\textwidth}@{}}\toprule Line&Kind&Name\\\midrule\endhead

14&theorem&\nolinkurl{PaperXI.FiniteIntersectionCriterion108.dimension_one_iff}\\

\bottomrule\end{longtable}

\section{\texorpdfstring{\nolinkurl{PaperXI.FiniteJointQuotient108}}{PaperXI.FiniteJointQuotient108}}\label{module:PaperXI.FiniteJointQuotient108}

\textbf{Selected layer:} base108 arithmetic branch.\quad\textbf{Source lines:} 47.\par\textbf{SHA-256:} {\scriptsize\nolinkurl{bd7174c5b31be9c6750454626e3e5627651dbea5f5465b49f2f3cbfbedc846df}}.

\VerbatimInput[fontsize=\scriptsize,breaklines=true,breakanywhere=true]{excerpts/D_FiniteJointQuotient108.txt}

\textbf{Direct project imports:} \nolinkurl{PaperXI.SmoothDegreeQuotient}.

\begin{longtable}{@{}p{0.07\textwidth}p{0.12\textwidth}p{0.75\textwidth}@{}}\toprule Line&Kind&Name\\\midrule\endhead

13&def&\nolinkurl{PaperXI.FiniteJointQuotient108.kernel}\\

16&instance&\nolinkurl{PaperXI.FiniteJointQuotient108.kernel_normal}\\

19&abbrev&\nolinkurl{PaperXI.FiniteJointQuotient108.quotient}\\

21&def&\nolinkurl{PaperXI.FiniteJointQuotient108.first}\\

25&def&\nolinkurl{PaperXI.FiniteJointQuotient108.second}\\

29&theorem&\nolinkurl{PaperXI.FiniteJointQuotient108.first_comp}\\

32&theorem&\nolinkurl{PaperXI.FiniteJointQuotient108.second_comp}\\

40&theorem&\nolinkurl{PaperXI.FiniteJointQuotient108.finite}\\

\bottomrule\end{longtable}

\section{\texorpdfstring{\nolinkurl{PaperXI.FiniteOrbitAverage68}}{PaperXI.FiniteOrbitAverage68}}\label{module:PaperXI.FiniteOrbitAverage68}

\textbf{Selected layer:} base108 arithmetic branch.\quad\textbf{Source lines:} 101.\par\textbf{SHA-256:} {\scriptsize\nolinkurl{41b8b3e87046ef73167fbebd34c7c4a62a6bf13714ff572356c87c64a17b3de6}}.

\VerbatimInput[fontsize=\scriptsize,breaklines=true,breakanywhere=true]{excerpts/D_FiniteOrbitAverage68.txt}

\textbf{Direct project imports:} none.

\begin{longtable}{@{}p{0.07\textwidth}p{0.12\textwidth}p{0.75\textwidth}@{}}\toprule Line&Kind&Name\\\midrule\endhead

16&abbrev&\nolinkurl{PaperXI.FiniteOrbitAverage68.Orbit}\\

20&def&\nolinkurl{PaperXI.FiniteOrbitAverage68.orbitEquiv}\\

32&def&\nolinkurl{PaperXI.FiniteOrbitAverage68.average}\\

40&theorem&\nolinkurl{PaperXI.FiniteOrbitAverage68.average_invariant}\\

51&theorem&\nolinkurl{PaperXI.FiniteOrbitAverage68.average_constraint}\\

68&def&\nolinkurl{PaperXI.FiniteOrbitAverage68.invariantHom}\\

79&theorem&\nolinkurl{PaperXI.FiniteOrbitAverage68.equivariant_section_of_finite_orbit}\\

\bottomrule\end{longtable}

\section{\texorpdfstring{\nolinkurl{PaperXI.FiniteProductDeterminant}}{PaperXI.FiniteProductDeterminant}}\label{module:PaperXI.FiniteProductDeterminant}

\textbf{Selected layer:} v0.44 parent.\quad\textbf{Source lines:} 64.\par\textbf{SHA-256:} {\scriptsize\nolinkurl{07340f6176b2b67b41bba32f65bdb41e2805856a67877d7e70b73585b9c7b556}}.

\VerbatimInput[fontsize=\scriptsize,breaklines=true,breakanywhere=true]{excerpts/D_FiniteProductDeterminant.txt}

\textbf{Direct project imports:} none.

\begin{longtable}{@{}p{0.07\textwidth}p{0.12\textwidth}p{0.75\textwidth}@{}}\toprule Line&Kind&Name\\\midrule\endhead

14&def&\nolinkurl{PaperXI.FiniteProductDeterminant.map}\\

19&theorem&\nolinkurl{PaperXI.FiniteProductDeterminant.map_apply}\\

24&theorem&\nolinkurl{PaperXI.FiniteProductDeterminant.determinant}\\

\bottomrule\end{longtable}

\section{\texorpdfstring{\nolinkurl{PaperXI.FiniteProfileBridge109}}{PaperXI.FiniteProfileBridge109}}\label{module:PaperXI.FiniteProfileBridge109}

\textbf{Selected layer:} Continuation110.\quad\textbf{Source lines:} 195.\par\textbf{SHA-256:} {\scriptsize\nolinkurl{d829f6afedc938f83bdb668850c4a684874773efe807983b4ae769844588cf0a}}.

\VerbatimInput[fontsize=\scriptsize,breaklines=true,breakanywhere=true]{excerpts/D_FiniteProfileBridge109.txt}

\textbf{Direct project imports:} \nolinkurl{PaperXI.HeisenbergTrace108}, \nolinkurl{PaperXI.CharacterPairing108}, \nolinkurl{PaperXI.FiniteJointQuotient108}.

\begin{longtable}{@{}p{0.07\textwidth}p{0.12\textwidth}p{0.75\textwidth}@{}}\toprule Line&Kind&Name\\\midrule\endhead

19&theorem&\nolinkurl{PaperXI.FiniteProfileBridge109.kernel_le_of_trace_support}\\

33&theorem&\nolinkurl{PaperXI.FiniteProfileBridge109.scalar_character_eq_one_on_kernel}\\

48&theorem&\nolinkurl{PaperXI.FiniteProfileBridge109.joint_kernel_le}\\

60&def&\nolinkurl{PaperXI.FiniteProfileBridge109.imageMap}\\

65&theorem&\nolinkurl{PaperXI.FiniteProfileBridge109.imageMap_surjective}\\

73&theorem&\nolinkurl{PaperXI.FiniteProfileBridge109.imageMap_kernel_le_character_kernel}\\

90&def&\nolinkurl{PaperXI.FiniteProfileBridge109.imageCharacter}\\

94&theorem&\nolinkurl{PaperXI.FiniteProfileBridge109.imageCharacter_apply}\\

102&theorem&\nolinkurl{PaperXI.FiniteProfileBridge109.image_scalar}\\

113&theorem&\nolinkurl{PaperXI.FiniteProfileBridge109.image_trace_zero}\\

120&theorem&\nolinkurl{PaperXI.FiniteProfileBridge109.image_intersection}\\

130&theorem&\nolinkurl{PaperXI.FiniteProfileBridge109.image_index}\\

134&theorem&\nolinkurl{PaperXI.FiniteProfileBridge109.image_characters_compatible}\\

158&theorem&\nolinkurl{PaperXI.FiniteProfileBridge109.quotient_kernel_le_action_kernel}\\

168&theorem&\nolinkurl{PaperXI.FiniteProfileBridge109.dimension_one_iff_of_original_profiles}\\

\bottomrule\end{longtable}

\section{\texorpdfstring{\nolinkurl{PaperXI.FiniteRepresentation}}{PaperXI.FiniteRepresentation}}\label{module:PaperXI.FiniteRepresentation}

\textbf{Selected layer:} v0.44 parent.\quad\textbf{Source lines:} 43.\par\textbf{SHA-256:} {\scriptsize\nolinkurl{b9ac153884bbf3b17a2701a8b50ee49c6f936899b7e807fa3a1cae5006e44d47}}.

\VerbatimInput[fontsize=\scriptsize,breaklines=true,breakanywhere=true]{excerpts/D_FiniteRepresentation.txt}

\textbf{Direct project imports:} \nolinkurl{PaperXI.RepresentationStability}.

\begin{longtable}{@{}p{0.07\textwidth}p{0.12\textwidth}p{0.75\textwidth}@{}}\toprule Line&Kind&Name\\\midrule\endhead

16&theorem&\nolinkurl{PaperXI.FiniteRepresentation.simple_finiteDimensional}\\

26&theorem&\nolinkurl{PaperXI.FiniteRepresentation.finiteDimensional}\\

35&theorem&\nolinkurl{PaperXI.FiniteRepresentation.finrank_pos}\\

\bottomrule\end{longtable}

\section{\texorpdfstring{\nolinkurl{PaperXI.FiniteResidueBaseChange}}{PaperXI.FiniteResidueBaseChange}}\label{module:PaperXI.FiniteResidueBaseChange}

\textbf{Selected layer:} v0.44 parent.\quad\textbf{Source lines:} 91.\par\textbf{SHA-256:} {\scriptsize\nolinkurl{559f2e07321f7991ab4790ffc84edb044865ea96de4a2c240148b8fb39aea272}}.

\VerbatimInput[fontsize=\scriptsize,breaklines=true,breakanywhere=true]{excerpts/D_FiniteResidueBaseChange.txt}

\textbf{Direct project imports:} \nolinkurl{PaperXI.ResidueDegreeBaseChange}.

\begin{longtable}{@{}p{0.07\textwidth}p{0.12\textwidth}p{0.75\textwidth}@{}}\toprule Line&Kind&Name\\\midrule\endhead

16&def&\nolinkurl{PaperXI.FiniteResidueBaseChange.arithmetic}\\

20&def&\nolinkurl{PaperXI.FiniteResidueBaseChange.geometric}\\

22&theorem&\nolinkurl{PaperXI.FiniteResidueBaseChange.arithmetic_apply}\\

24&theorem&\nolinkurl{PaperXI.FiniteResidueBaseChange.arithmetic_pow_apply}\\

33&theorem&\nolinkurl{PaperXI.FiniteResidueBaseChange.geometric_inverse_power}\\

42&theorem&\nolinkurl{PaperXI.FiniteResidueBaseChange.arithmetic_tower}\\

50&theorem&\nolinkurl{PaperXI.FiniteResidueBaseChange.geometric_tower}\\

56&def&\nolinkurl{PaperXI.FiniteResidueBaseChange.pullbackEquiv}\\

67&def&\nolinkurl{PaperXI.FiniteResidueBaseChange.forget}\\

73&theorem&\nolinkurl{PaperXI.FiniteResidueBaseChange.forget_injective}\\

80&theorem&\nolinkurl{PaperXI.FiniteResidueBaseChange.geometric_order}\\

86&theorem&\nolinkurl{PaperXI.FiniteResidueBaseChange.geometric_zpow_eq_one}\\

\bottomrule\end{longtable}

\section{\texorpdfstring{\nolinkurl{PaperXI.FiniteResidueSubgroup}}{PaperXI.FiniteResidueSubgroup}}\label{module:PaperXI.FiniteResidueSubgroup}

\textbf{Selected layer:} v0.44 parent.\quad\textbf{Source lines:} 128.\par\textbf{SHA-256:} {\scriptsize\nolinkurl{2da76b0ad80ae46f7b05f39967fee2a9247b7c80f1a9e0ef61b0423a026f8a72}}.

\VerbatimInput[fontsize=\scriptsize,breaklines=true,breakanywhere=true]{excerpts/D_FiniteResidueSubgroup.txt}

\textbf{Direct project imports:} \nolinkurl{PaperXI.FiniteResidueBaseChange}.

\begin{longtable}{@{}p{0.07\textwidth}p{0.12\textwidth}p{0.75\textwidth}@{}}\toprule Line&Kind&Name\\\midrule\endhead

14&theorem&\nolinkurl{PaperXI.FiniteResidueSubgroup.arithmetic_map}\\

19&theorem&\nolinkurl{PaperXI.FiniteResidueSubgroup.geometric_map}\\

26&theorem&\nolinkurl{PaperXI.FiniteResidueSubgroup.geometric_pow_map}\\

38&theorem&\nolinkurl{PaperXI.FiniteResidueSubgroup.geometric_zpow_map}\\

51&theorem&\nolinkurl{PaperXI.FiniteResidueSubgroup.fixes_subfield_iff}\\

70&def&\nolinkurl{PaperXI.FiniteResidueSubgroup.fixing}\\

87&theorem&\nolinkurl{PaperXI.FiniteResidueSubgroup.kernel_le_fixing}\\

94&theorem&\nolinkurl{PaperXI.FiniteResidueSubgroup.projection_mem_iff}\\

104&theorem&\nolinkurl{PaperXI.FiniteResidueSubgroup.projection_fixes}\\

111&def&\nolinkurl{PaperXI.FiniteResidueSubgroup.fixingEquiv}\\

122&def&\nolinkurl{PaperXI.FiniteResidueSubgroup.unramifiedEquiv}\\

\bottomrule\end{longtable}

\section{\texorpdfstring{\nolinkurl{PaperXI.FiniteTameComparison}}{PaperXI.FiniteTameComparison}}\label{module:PaperXI.FiniteTameComparison}

\textbf{Selected layer:} v0.44 parent.\quad\textbf{Source lines:} 113.\par\textbf{SHA-256:} {\scriptsize\nolinkurl{1c4357da32eaade5be0c88670eaa605f47fbc4269253888bbf9f72509d0e428d}}.

\VerbatimInput[fontsize=\scriptsize,breaklines=true,breakanywhere=true]{excerpts/D_FiniteTameComparison.txt}

\textbf{Direct project imports:} \nolinkurl{PaperXI.InducedTameProfile}, \nolinkurl{PaperXI.TameGaussNaturality}, \nolinkurl{PaperXI.TameTensorSource}.

\begin{longtable}{@{}p{0.07\textwidth}p{0.12\textwidth}p{0.75\textwidth}@{}}\toprule Line&Kind&Name\\\midrule\endhead

17&structure&\nolinkurl{PaperXI.FiniteTameComparison.Observables}\\

40&def&\nolinkurl{PaperXI.FiniteTameComparison.observed}\\

53&def&\nolinkurl{PaperXI.FiniteTameComparison.fromCoordinates}\\

64&theorem&\nolinkurl{PaperXI.FiniteTameComparison.observations_eq_coordinates}\\

83&def&\nolinkurl{PaperXI.FiniteTameComparison.interactionEquiv}\\

93&theorem&\nolinkurl{PaperXI.FiniteTameComparison.finite_comparison_package}\\

\bottomrule\end{longtable}

\section{\texorpdfstring{\nolinkurl{PaperXI.FiniteTensor}}{PaperXI.FiniteTensor}}\label{module:PaperXI.FiniteTensor}

\textbf{Selected layer:} v0.44 parent.\quad\textbf{Source lines:} 69.\par\textbf{SHA-256:} {\scriptsize\nolinkurl{8a628f3f2be9e4cbc308b13abbc71ef858a82cf897523b60554829ed81ac0782}}.

\VerbatimInput[fontsize=\scriptsize,breaklines=true,breakanywhere=true]{excerpts/D_FiniteTensor.txt}

\textbf{Direct project imports:} none.

\begin{longtable}{@{}p{0.07\textwidth}p{0.12\textwidth}p{0.75\textwidth}@{}}\toprule Line&Kind&Name\\\midrule\endhead

17&def&\nolinkurl{PaperXI.FiniteTensor.representation}\\

25&theorem&\nolinkurl{PaperXI.FiniteTensor.representation_tprod}\\

32&theorem&\nolinkurl{PaperXI.FiniteTensor.finrank}\\

40&theorem&\nolinkurl{PaperXI.FiniteTensor.finrank_empty}\\

46&theorem&\nolinkurl{PaperXI.FiniteTensor.nontrivial}\\

53&theorem&\nolinkurl{PaperXI.FiniteTensor.empty_action}\\

60&theorem&\nolinkurl{PaperXI.FiniteTensor.scalar_map}\\

\bottomrule\end{longtable}

\section{\texorpdfstring{\nolinkurl{PaperXI.FiniteUnramifiedWeil}}{PaperXI.FiniteUnramifiedWeil}}\label{module:PaperXI.FiniteUnramifiedWeil}

\textbf{Selected layer:} v0.44 parent.\quad\textbf{Source lines:} 198.\par\textbf{SHA-256:} {\scriptsize\nolinkurl{73d632d598c02e6eced406cf8bd8f48422a3b79634bfb718dea82cc440d29f40}}.

\VerbatimInput[fontsize=\scriptsize,breaklines=true,breakanywhere=true]{excerpts/D_FiniteUnramifiedWeil.txt}

\textbf{Direct project imports:} \nolinkurl{PaperXI.ResidueFrobeniusAttachment}, \nolinkurl{PaperXI.ResidueDegreeBaseChange}.

\begin{longtable}{@{}p{0.07\textwidth}p{0.12\textwidth}p{0.75\textwidth}@{}}\toprule Line&Kind&Name\\\midrule\endhead

26&abbrev&\nolinkurl{PaperXI.FiniteUnramifiedWeil.UnramifiedWeil}\\

30&def&\nolinkurl{PaperXI.FiniteUnramifiedWeil.degreeEquiv}\\

44&theorem&\nolinkurl{PaperXI.FiniteUnramifiedWeil.degreeEquiv_apply}\\

50&theorem&\nolinkurl{PaperXI.FiniteUnramifiedWeil.projection_eq_geometric_zpow_degree}\\

57&theorem&\nolinkurl{PaperXI.FiniteUnramifiedWeil.degree_bijective}\\

62&theorem&\nolinkurl{PaperXI.FiniteUnramifiedWeil.inertia_eq_bot}\\

67&def&\nolinkurl{PaperXI.FiniteUnramifiedWeil.geometricLift}\\

71&theorem&\nolinkurl{PaperXI.FiniteUnramifiedWeil.degree_geometricLift}\\

77&theorem&\nolinkurl{PaperXI.FiniteUnramifiedWeil.projection_geometricLift}\\

84&def&\nolinkurl{PaperXI.FiniteUnramifiedWeil.arithmeticLift}\\

88&theorem&\nolinkurl{PaperXI.FiniteUnramifiedWeil.degree_arithmeticLift}\\

95&theorem&\nolinkurl{PaperXI.FiniteUnramifiedWeil.projection_arithmeticLift}\\

104&theorem&\nolinkurl{PaperXI.FiniteUnramifiedWeil.eq_geometricLift_zpow_degree}\\

118&theorem&\nolinkurl{PaperXI.FiniteUnramifiedWeil.geometricLift_zpowers_eq_top}\\

128&instance&\nolinkurl{PaperXI.FiniteUnramifiedWeil.unramifiedWeil_isMulCommutative}\\

136&theorem&\nolinkurl{PaperXI.FiniteUnramifiedWeil.irreducible_finrank_eq_one}\\

144&theorem&\nolinkurl{PaperXI.FiniteUnramifiedWeil.projection_surjective}\\

156&theorem&\nolinkurl{PaperXI.FiniteUnramifiedWeil.projection_ker_eq_period}\\

189&def&\nolinkurl{PaperXI.FiniteUnramifiedWeil.finiteResidueQuotientEquiv}\\

\bottomrule\end{longtable}

\section{\texorpdfstring{\nolinkurl{PaperXI.FixedMoving}}{PaperXI.FixedMoving}}\label{module:PaperXI.FixedMoving}

\textbf{Selected layer:} v0.44 parent.\quad\textbf{Source lines:} 198.\par\textbf{SHA-256:} {\scriptsize\nolinkurl{c5159f03721e3f1446641781dcf303a90ea160dafb3eed4edac5906238153c67}}.

\VerbatimInput[fontsize=\scriptsize,breaklines=true,breakanywhere=true]{excerpts/D_FixedMoving.txt}

\textbf{Direct project imports:} none.

\begin{longtable}{@{}p{0.07\textwidth}p{0.12\textwidth}p{0.75\textwidth}@{}}\toprule Line&Kind&Name\\\midrule\endhead

13&def&\nolinkurl{PaperXI.FixedMoving.fixed}\\

16&def&\nolinkurl{PaperXI.FixedMoving.moving}\\

18&theorem&\nolinkurl{PaperXI.FixedMoving.mem_fixed}\\

21&theorem&\nolinkurl{PaperXI.FixedMoving.mem_moving}\\

25&theorem&\nolinkurl{PaperXI.FixedMoving.orthogonal}\\

37&theorem&\nolinkurl{PaperXI.FixedMoving.orthogonal_reverse}\\

43&def&\nolinkurl{PaperXI.FixedMoving.fixedForm}\\

46&def&\nolinkurl{PaperXI.FixedMoving.movingForm}\\

48&theorem&\nolinkurl{PaperXI.FixedMoving.fixedForm_alt}\\

51&theorem&\nolinkurl{PaperXI.FixedMoving.movingForm_alt}\\

55&theorem&\nolinkurl{PaperXI.FixedMoving.restrict_nondegenerate}\\

77&theorem&\nolinkurl{PaperXI.FixedMoving.movingForm_nondegenerate_of_isCompl}\\

84&theorem&\nolinkurl{PaperXI.FixedMoving.fixedForm_nondegenerate_of_isCompl}\\

91&theorem&\nolinkurl{PaperXI.FixedMoving.moving_stable}\\

97&theorem&\nolinkurl{PaperXI.FixedMoving.moving_inverse_stable}\\

105&def&\nolinkurl{PaperXI.FixedMoving.movingAction}\\

113&theorem&\nolinkurl{PaperXI.FixedMoving.movingAction_coe}\\

116&theorem&\nolinkurl{PaperXI.FixedMoving.movingAction_symm_coe}\\

119&theorem&\nolinkurl{PaperXI.FixedMoving.movingAction_preserves}\\

124&theorem&\nolinkurl{PaperXI.FixedMoving.movingAction_fixedPointFree_of_isCompl}\\

133&theorem&\nolinkurl{PaperXI.FixedMoving.moving_difference_injective}\\

145&def&\nolinkurl{PaperXI.FixedMoving.movingDifferenceEquiv}\\

151&theorem&\nolinkurl{PaperXI.FixedMoving.movingDifferenceEquiv_apply}\\

159&def&\nolinkurl{PaperXI.FixedMoving.movingEquiv}\\

179&theorem&\nolinkurl{PaperXI.FixedMoving.movingEquiv_coe}\\

184&theorem&\nolinkurl{PaperXI.FixedMoving.movingEquiv_intertwines}\\

191&theorem&\nolinkurl{PaperXI.FixedMoving.movingEquiv_preserves}\\

\bottomrule\end{longtable}

\section{\texorpdfstring{\nolinkurl{PaperXI.FixedMovingSource}}{PaperXI.FixedMovingSource}}\label{module:PaperXI.FixedMovingSource}

\textbf{Selected layer:} v0.44 parent.\quad\textbf{Source lines:} 111.\par\textbf{SHA-256:} {\scriptsize\nolinkurl{2591738676c9f4830f6e41ce811376e7b6b2c17f1afe99cfbf94cffebdb7aa22}}.

\VerbatimInput[fontsize=\scriptsize,breaklines=true,breakanywhere=true]{excerpts/D_FixedMovingSource.txt}

\textbf{Direct project imports:} \nolinkurl{PaperXI.FixedMoving}, \nolinkurl{PaperXI.FiniteFieldAverage}, \nolinkurl{PaperXI.OddPhase}.

\begin{longtable}{@{}p{0.07\textwidth}p{0.12\textwidth}p{0.75\textwidth}@{}}\toprule Line&Kind&Name\\\midrule\endhead

15&theorem&\nolinkurl{PaperXI.FixedMoving.isCompl_of_period}\\

19&def&\nolinkurl{PaperXI.FixedMoving.decomposition}\\

22&theorem&\nolinkurl{PaperXI.FixedMoving.decomposition_apply}\\

27&theorem&\nolinkurl{PaperXI.FixedMoving.decomposition_fixed}\\

33&theorem&\nolinkurl{PaperXI.FixedMoving.decomposition_moving}\\

38&theorem&\nolinkurl{PaperXI.FixedMoving.movingForm_nondegenerate}\\

44&theorem&\nolinkurl{PaperXI.FixedMoving.fixedForm_nondegenerate}\\

55&theorem&\nolinkurl{PaperXI.FixedMoving.source_order_pos}\\

63&theorem&\nolinkurl{PaperXI.FixedMoving.source_order_ne_zero}\\

68&theorem&\nolinkurl{PaperXI.FixedMoving.source_isCompl}\\

73&def&\nolinkurl{PaperXI.FixedMoving.sourceDecomposition}\\

77&theorem&\nolinkurl{PaperXI.FixedMoving.sourceDecomposition_apply}\\

81&theorem&\nolinkurl{PaperXI.FixedMoving.source_moving_fixedPointFree}\\

86&def&\nolinkurl{PaperXI.FixedMoving.sourceMovingDifference}\\

89&theorem&\nolinkurl{PaperXI.FixedMoving.sourceMovingDifference_apply}\\

95&theorem&\nolinkurl{PaperXI.FixedMoving.source}\\

\bottomrule\end{longtable}

\section{\texorpdfstring{\nolinkurl{PaperXI.FourierCovariance}}{PaperXI.FourierCovariance}}\label{module:PaperXI.FourierCovariance}

\textbf{Selected layer:} v0.44 parent.\quad\textbf{Source lines:} 40.\par\textbf{SHA-256:} {\scriptsize\nolinkurl{83b273abc08207d3b23d39152863370e2345bc7a04f422578b880b39ffb6bb6c}}.

\VerbatimInput[fontsize=\scriptsize,breaklines=true,breakanywhere=true]{excerpts/D_FourierCovariance.txt}

\textbf{Direct project imports:} \nolinkurl{PaperXI.ProjectorTransport}.

\begin{longtable}{@{}p{0.07\textwidth}p{0.12\textwidth}p{0.75\textwidth}@{}}\toprule Line&Kind&Name\\\midrule\endhead

13&def&\nolinkurl{PaperXI.FourierCovariance.mulPermutation}\\

19&def&\nolinkurl{PaperXI.FourierCovariance.fourierOperator}\\

25&theorem&\nolinkurl{PaperXI.FourierCovariance.inverse_exponent_covariance}\\

\bottomrule\end{longtable}

\section{\texorpdfstring{\nolinkurl{PaperXI.FourierPackets}}{PaperXI.FourierPackets}}\label{module:PaperXI.FourierPackets}

\textbf{Selected layer:} v0.44 parent.\quad\textbf{Source lines:} 50.\par\textbf{SHA-256:} {\scriptsize\nolinkurl{3521693ff5ede2869607d8308db0f63ce2cb6a736c5d728fd43d9b18c942cade}}.

\VerbatimInput[fontsize=\scriptsize,breaklines=true,breakanywhere=true]{excerpts/D_FourierPackets.txt}

\textbf{Direct project imports:} \nolinkurl{PaperXI.FourierProjectors}.

\begin{longtable}{@{}p{0.07\textwidth}p{0.12\textwidth}p{0.75\textwidth}@{}}\toprule Line&Kind&Name\\\midrule\endhead

13&def&\nolinkurl{PaperXI.FourierPackets.packet}\\

20&theorem&\nolinkurl{PaperXI.FourierPackets.projector_mem_packet}\\

27&theorem&\nolinkurl{PaperXI.FourierPackets.projector_on_packet}\\

36&theorem&\nolinkurl{PaperXI.FourierPackets.projector_range}\\

44&theorem&\nolinkurl{PaperXI.FourierPackets.vector_reconstruction}\\

\bottomrule\end{longtable}

\section{\texorpdfstring{\nolinkurl{PaperXI.FourierProjectors}}{PaperXI.FourierProjectors}}\label{module:PaperXI.FourierProjectors}

\textbf{Selected layer:} v0.44 parent.\quad\textbf{Source lines:} 79.\par\textbf{SHA-256:} {\scriptsize\nolinkurl{02af62a8c8e5753ef7a695629d44c71db407211d18a6bd6b6a48865ba731de52}}.

\VerbatimInput[fontsize=\scriptsize,breaklines=true,breakanywhere=true]{excerpts/D_FourierProjectors.txt}

\textbf{Direct project imports:} \nolinkurl{PaperXI.CyclicAction}.

\begin{longtable}{@{}p{0.07\textwidth}p{0.12\textwidth}p{0.75\textwidth}@{}}\toprule Line&Kind&Name\\\midrule\endhead

13&def&\nolinkurl{PaperXI.FourierProjectors.projector}\\

17&theorem&\nolinkurl{PaperXI.FourierProjectors.projector_formula}\\

21&theorem&\nolinkurl{PaperXI.FourierProjectors.character_sum}\\

26&theorem&\nolinkurl{PaperXI.FourierProjectors.action_projector}\\

44&theorem&\nolinkurl{PaperXI.FourierProjectors.projector_product}\\

58&theorem&\nolinkurl{PaperXI.FourierProjectors.projector_idempotent}\\

63&theorem&\nolinkurl{PaperXI.FourierProjectors.projector_orthogonal}\\

68&theorem&\nolinkurl{PaperXI.FourierProjectors.projector_complete}\\

\bottomrule\end{longtable}

\section{\texorpdfstring{\nolinkurl{PaperXI.FractionMatrix106}}{PaperXI.FractionMatrix106}}\label{module:PaperXI.FractionMatrix106}

\textbf{Selected layer:} base108 arithmetic branch.\quad\textbf{Source lines:} 57.\par\textbf{SHA-256:} {\scriptsize\nolinkurl{37125ed47f602558238d6c6b2e9ef6abc4f83592ef58525b03697009a9ac54b6}}.

\VerbatimInput[fontsize=\scriptsize,breaklines=true,breakanywhere=true]{excerpts/D_FractionMatrix106.txt}

\textbf{Direct project imports:} \nolinkurl{PaperXI.SmithMatrix106}, \nolinkurl{PaperXI.SplitRadical105}.

\begin{longtable}{@{}p{0.07\textwidth}p{0.12\textwidth}p{0.75\textwidth}@{}}\toprule Line&Kind&Name\\\midrule\endhead

17&theorem&\nolinkurl{PaperXI.FractionMatrix106.exists_integral_multiple}\\

38&theorem&\nolinkurl{PaperXI.FractionMatrix106.exists_integral_diagonal}\\

\bottomrule\end{longtable}

\section{\texorpdfstring{\nolinkurl{PaperXI.FreeLatticeCard110}}{PaperXI.FreeLatticeCard110}}\label{module:PaperXI.FreeLatticeCard110}

\textbf{Selected layer:} Continuation110.\quad\textbf{Source lines:} 55.\par\textbf{SHA-256:} {\scriptsize\nolinkurl{07d5f21b13996d68c3d71abd71dab1052fd72ae292beba461aee4f571f6bcd08}}.

\VerbatimInput[fontsize=\scriptsize,breaklines=true,breakanywhere=true]{excerpts/D_FreeLatticeCard110.txt}

\textbf{Direct project imports:} none.

\begin{longtable}{@{}p{0.07\textwidth}p{0.12\textwidth}p{0.75\textwidth}@{}}\toprule Line&Kind&Name\\\midrule\endhead

16&theorem&\nolinkurl{PaperXI.FreeLatticeCard110.ideal_smul_top_pi}\\

34&def&\nolinkurl{PaperXI.FreeLatticeCard110.quotientEquiv}\\

41&theorem&\nolinkurl{PaperXI.FreeLatticeCard110.card_quotient}\\

46&theorem&\nolinkurl{PaperXI.FreeLatticeCard110.card_quotient_prime_pow}\\

\bottomrule\end{longtable}

\section{\texorpdfstring{\nolinkurl{PaperXI.FrobeniusOrbit}}{PaperXI.FrobeniusOrbit}}\label{module:PaperXI.FrobeniusOrbit}

\textbf{Selected layer:} v0.44 parent.\quad\textbf{Source lines:} 42.\par\textbf{SHA-256:} {\scriptsize\nolinkurl{7d1f7d09d4d67afe3a6e5721cf265b76d2ff69a13ef63bb6a753133817f2f7c3}}.

\VerbatimInput[fontsize=\scriptsize,breaklines=true,breakanywhere=true]{excerpts/D_FrobeniusOrbit.txt}

\textbf{Direct project imports:} \nolinkurl{PaperXI.PacketTransport}.

\begin{longtable}{@{}p{0.07\textwidth}p{0.12\textwidth}p{0.75\textwidth}@{}}\toprule Line&Kind&Name\\\midrule\endhead

8&def&\nolinkurl{PaperXI.FrobeniusOrbit.Related}\\

11&theorem&\nolinkurl{PaperXI.FrobeniusOrbit.refl}\\

14&theorem&\nolinkurl{PaperXI.FrobeniusOrbit.trans}\\

20&theorem&\nolinkurl{PaperXI.FrobeniusOrbit.symm}\\

26&theorem&\nolinkurl{PaperXI.FrobeniusOrbit.step}\\

29&theorem&\nolinkurl{PaperXI.FrobeniusOrbit.back}\\

32&theorem&\nolinkurl{PaperXI.FrobeniusOrbit.step_mem}\\

35&theorem&\nolinkurl{PaperXI.FrobeniusOrbit.back_mem}\\

38&theorem&\nolinkurl{PaperXI.FrobeniusOrbit.same_orbit_iff}\\

\bottomrule\end{longtable}

\section{\texorpdfstring{\nolinkurl{PaperXI.FrobeniusPeriod}}{PaperXI.FrobeniusPeriod}}\label{module:PaperXI.FrobeniusPeriod}

\textbf{Selected layer:} v0.44 parent.\quad\textbf{Source lines:} 97.\par\textbf{SHA-256:} {\scriptsize\nolinkurl{8e34bfa6b4636a8b392bdff1eb6f25eb5c9061f592ae1bff13c3179e48ec14fb}}.

\VerbatimInput[fontsize=\scriptsize,breaklines=true,breakanywhere=true]{excerpts/D_FrobeniusPeriod.txt}

\textbf{Direct project imports:} \nolinkurl{PaperXI.PacketSupport}.

\begin{longtable}{@{}p{0.07\textwidth}p{0.12\textwidth}p{0.75\textwidth}@{}}\toprule Line&Kind&Name\\\midrule\endhead

12&def&\nolinkurl{PaperXI.FrobeniusPeriod.permutation}\\

14&def&\nolinkurl{PaperXI.FrobeniusPeriod.period}\\

16&theorem&\nolinkurl{PaperXI.FrobeniusPeriod.iterate_label}\\

28&theorem&\nolinkurl{PaperXI.FrobeniusPeriod.period_pos}\\

31&theorem&\nolinkurl{PaperXI.FrobeniusPeriod.period_le}\\

35&theorem&\nolinkurl{PaperXI.FrobeniusPeriod.return_label}\\

43&theorem&\nolinkurl{PaperXI.FrobeniusPeriod.stabilizing_power_iff}\\

49&theorem&\nolinkurl{PaperXI.FrobeniusPeriod.no_short_return}\\

56&theorem&\nolinkurl{PaperXI.FrobeniusPeriod.related_iff_nat}\\

68&theorem&\nolinkurl{PaperXI.FrobeniusPeriod.related_iff_bounded}\\

79&def&\nolinkurl{PaperXI.FrobeniusPeriod.orbitEquiv}\\

93&theorem&\nolinkurl{PaperXI.FrobeniusPeriod.orbit_card}\\

\bottomrule\end{longtable}

\section{\texorpdfstring{\nolinkurl{PaperXI.FrobeniusPeriodCount}}{PaperXI.FrobeniusPeriodCount}}\label{module:PaperXI.FrobeniusPeriodCount}

\textbf{Selected layer:} v0.44 parent.\quad\textbf{Source lines:} 125.\par\textbf{SHA-256:} {\scriptsize\nolinkurl{53f45d0eb9fed4c2ad104695c3730de6b59560120f164630b865c30a3ccb0749}}.

\VerbatimInput[fontsize=\scriptsize,breaklines=true,breakanywhere=true]{excerpts/D_FrobeniusPeriodCount.txt}

\textbf{Direct project imports:} none.

\begin{longtable}{@{}p{0.07\textwidth}p{0.12\textwidth}p{0.75\textwidth}@{}}\toprule Line&Kind&Name\\\midrule\endhead

15&def&\nolinkurl{PaperXI.FrobeniusPeriodCount.exactCount}\\

18&def&\nolinkurl{PaperXI.FrobeniusPeriodCount.fixedCount}\\

22&lemma&\nolinkurl{PaperXI.FrobeniusPeriodCount.exactCount_eq_filter}\\

28&lemma&\nolinkurl{PaperXI.FrobeniusPeriodCount.fixedCount_eq_filter}\\

34&theorem&\nolinkurl{PaperXI.FrobeniusPeriodCount.sum_exactCount}\\

59&theorem&\nolinkurl{PaperXI.FrobeniusPeriodCount.moebius_exactCount}\\

75&def&\nolinkurl{PaperXI.FrobeniusPeriodCount.regularPerm}\\

80&def&\nolinkurl{PaperXI.FrobeniusPeriodCount.orbitCount}\\

84&theorem&\nolinkurl{PaperXI.FrobeniusPeriodCount.regularPerm_minimalPeriod}\\

96&theorem&\nolinkurl{PaperXI.FrobeniusPeriodCount.exactCount_eq_orbitCount_mul}\\

119&theorem&\nolinkurl{PaperXI.FrobeniusPeriodCount.period_dvd_exactCount}\\

\bottomrule\end{longtable}

\section{\texorpdfstring{\nolinkurl{PaperXI.FullBridgeTransport}}{PaperXI.FullBridgeTransport}}\label{module:PaperXI.FullBridgeTransport}

\textbf{Selected layer:} v0.44 parent.\quad\textbf{Source lines:} 50.\par\textbf{SHA-256:} {\scriptsize\nolinkurl{af018ebd737c4d182525a72c4eddc2561a4ca9a2a228aabd0bac83dae1a9d816}}.

\VerbatimInput[fontsize=\scriptsize,breaklines=true,breakanywhere=true]{excerpts/D_FullBridgeTransport.txt}

\textbf{Direct project imports:} \nolinkurl{PaperXI.CyclicBridge}.

\begin{longtable}{@{}p{0.07\textwidth}p{0.12\textwidth}p{0.75\textwidth}@{}}\toprule Line&Kind&Name\\\midrule\endhead

14&theorem&\nolinkurl{PaperXI.FullBridgeTransport.semidirect_along}\\

36&theorem&\nolinkurl{PaperXI.FullBridgeTransport.through_factorizations}\\

\bottomrule\end{longtable}

\section{\texorpdfstring{\nolinkurl{PaperXI.FullIntertwiningLocus69}}{PaperXI.FullIntertwiningLocus69}}\label{module:PaperXI.FullIntertwiningLocus69}

\textbf{Selected layer:} base108 arithmetic branch.\quad\textbf{Source lines:} 95.\par\textbf{SHA-256:} {\scriptsize\nolinkurl{d1f43873562d510630801d569413b3ad3b4d659e2b4403e1d74714196488c12a}}.

\VerbatimInput[fontsize=\scriptsize,breaklines=true,breakanywhere=true]{excerpts/D_FullIntertwiningLocus69.txt}

\textbf{Direct project imports:} \nolinkurl{PaperXI.InductionIrreducibility}.

\begin{longtable}{@{}p{0.07\textwidth}p{0.12\textwidth}p{0.75\textwidth}@{}}\toprule Line&Kind&Name\\\midrule\endhead

14&def&\nolinkurl{PaperXI.FullIntertwiningLocus69.locus}\\

19&def&\nolinkurl{PaperXI.FullIntertwiningLocus69.internalIntertwiner}\\

30&theorem&\nolinkurl{PaperXI.FullIntertwiningLocus69.internalIntertwiner_ne_zero}\\

41&theorem&\nolinkurl{PaperXI.FullIntertwiningLocus69.subgroup_subset}\\

48&theorem&\nolinkurl{PaperXI.FullIntertwiningLocus69.locus_eq_iff}\\

71&def&\nolinkurl{PaperXI.FullIntertwiningLocus69.standardLocus}\\

75&theorem&\nolinkurl{PaperXI.FullIntertwiningLocus69.standardLocus_eq_iff}\\

88&theorem&\nolinkurl{PaperXI.FullIntertwiningLocus69.off_subgroup_of_standard_locus}\\

\bottomrule\end{longtable}

\section{\texorpdfstring{\nolinkurl{PaperXI.FullTypeIntertwining105}}{PaperXI.FullTypeIntertwining105}}\label{module:PaperXI.FullTypeIntertwining105}

\textbf{Selected layer:} base108 arithmetic branch.\quad\textbf{Source lines:} 151.\par\textbf{SHA-256:} {\scriptsize\nolinkurl{77dddac0fa33899fcac3fba052d3a5bf97fd3f7c3e9443409e6fd3be1e7f7be6}}.

\VerbatimInput[fontsize=\scriptsize,breaklines=true,breakanywhere=true]{excerpts/D_FullTypeIntertwining105.txt}

\textbf{Direct project imports:} \nolinkurl{PaperXI.CuspidalCartan105}, \nolinkurl{PaperXI.FullIntertwiningLocus69}, \nolinkurl{PaperXI.CompactModCenterInduction68}.

\begin{longtable}{@{}p{0.07\textwidth}p{0.12\textwidth}p{0.75\textwidth}@{}}\toprule Line&Kind&Name\\\midrule\endhead

20&def&\nolinkurl{PaperXI.FullTypeIntertwining105.mapOfAlong}\\

30&theorem&\nolinkurl{PaperXI.FullTypeIntertwining105.alongOfMap}\\

47&theorem&\nolinkurl{PaperXI.FullTypeIntertwining105.existsAt_iff_locus}\\

63&theorem&\nolinkurl{PaperXI.FullTypeIntertwining105.restrict_existsAt}\\

72&theorem&\nolinkurl{PaperXI.FullTypeIntertwining105.extension_irreducible}\\

92&theorem&\nolinkurl{PaperXI.FullTypeIntertwining105.full_locus}\\

120&theorem&\nolinkurl{PaperXI.FullTypeIntertwining105.induction_package}\\

\bottomrule\end{longtable}

\section{\texorpdfstring{\nolinkurl{PaperXI.GallagherTensorIrreducibility58}}{PaperXI.GallagherTensorIrreducibility58}}\label{module:PaperXI.GallagherTensorIrreducibility58}

\textbf{Selected layer:} base108 arithmetic branch.\quad\textbf{Source lines:} 165.\par\textbf{SHA-256:} {\scriptsize\nolinkurl{8bf42d852c9fc0f54c67a1dbdfb96f9df652d86138597dd1db0006bcf90e552b}}.

\VerbatimInput[fontsize=\scriptsize,breaklines=true,breakanywhere=true]{excerpts/D_GallagherTensorIrreducibility58.txt}

\textbf{Direct project imports:} \nolinkurl{PaperXI.SchurScalar}, \nolinkurl{PaperXI.RepresentationStability}.

\begin{longtable}{@{}p{0.07\textwidth}p{0.12\textwidth}p{0.75\textwidth}@{}}\toprule Line&Kind&Name\\\midrule\endhead

25&theorem&\nolinkurl{PaperXI.GallagherTensorIrreducibility58.algebraAction_surjective}\\

59&theorem&\nolinkurl{PaperXI.GallagherTensorIrreducibility58.stable_tensor_algebra}\\

73&theorem&\nolinkurl{PaperXI.GallagherTensorIrreducibility58.stable_tensor_end}\\

83&theorem&\nolinkurl{PaperXI.GallagherTensorIrreducibility58.map_rankOne}\\

95&def&\nolinkurl{PaperXI.GallagherTensorIrreducibility58.multiplicitySubspace}\\

104&theorem&\nolinkurl{PaperXI.GallagherTensorIrreducibility58.tensor_irreducible}\\

158&theorem&\nolinkurl{PaperXI.GallagherTensorIrreducibility58.quotient_tensor_irreducible}\\

\bottomrule\end{longtable}

\section{\texorpdfstring{\nolinkurl{PaperXI.GaussCharacters}}{PaperXI.GaussCharacters}}\label{module:PaperXI.GaussCharacters}

\textbf{Selected layer:} v0.44 parent.\quad\textbf{Source lines:} 67.\par\textbf{SHA-256:} {\scriptsize\nolinkurl{f49f317f384804ecff8614197f0acd360d1575046498c320b2027bd7ba16d8ca}}.

\VerbatimInput[fontsize=\scriptsize,breaklines=true,breakanywhere=true]{excerpts/D_GaussCharacters.txt}

\textbf{Direct project imports:} \nolinkurl{PaperXI.GaussSums}.

\begin{longtable}{@{}p{0.07\textwidth}p{0.12\textwidth}p{0.75\textwidth}@{}}\toprule Line&Kind&Name\\\midrule\endhead

12&def&\nolinkurl{PaperXI.GaussCharacters.sum}\\

16&theorem&\nolinkurl{PaperXI.GaussCharacters.sum_eq_unitSum}\\

22&theorem&\nolinkurl{PaperXI.GaussCharacters.extension_ne_one}\\

33&theorem&\nolinkurl{PaperXI.GaussCharacters.extension_inverse}\\

38&theorem&\nolinkurl{PaperXI.GaussCharacters.magnitude}\\

45&def&\nolinkurl{PaperXI.GaussCharacters.phase}\\

49&theorem&\nolinkurl{PaperXI.GaussCharacters.phase_eq}\\

54&theorem&\nolinkurl{PaperXI.GaussCharacters.phase_norm}\\

60&theorem&\nolinkurl{PaperXI.GaussCharacters.phase_duality}\\

\bottomrule\end{longtable}

\section{\texorpdfstring{\nolinkurl{PaperXI.GaussSums}}{PaperXI.GaussSums}}\label{module:PaperXI.GaussSums}

\textbf{Selected layer:} v0.44 parent.\quad\textbf{Source lines:} 90.\par\textbf{SHA-256:} {\scriptsize\nolinkurl{ace300219a85a3337ec026bb367507f5088f887306c309e17af938fdfca02d02}}.

\VerbatimInput[fontsize=\scriptsize,breaklines=true,breakanywhere=true]{excerpts/D_GaussSums.txt}

\textbf{Direct project imports:} none.

\begin{longtable}{@{}p{0.07\textwidth}p{0.12\textwidth}p{0.75\textwidth}@{}}\toprule Line&Kind&Name\\\midrule\endhead

16&def&\nolinkurl{PaperXI.GaussSums.unitSum}\\

20&theorem&\nolinkurl{PaperXI.GaussSums.unitSum_eq_gaussSum}\\

28&theorem&\nolinkurl{PaperXI.GaussSums.sum_mul_inverse}\\

35&theorem&\nolinkurl{PaperXI.GaussSums.conjugate_sum}\\

40&theorem&\nolinkurl{PaperXI.GaussSums.norm_sq}\\

50&theorem&\nolinkurl{PaperXI.GaussSums.magnitude}\\

56&theorem&\nolinkurl{PaperXI.GaussSums.sum_duality}\\

65&def&\nolinkurl{PaperXI.GaussSums.phase}\\

69&theorem&\nolinkurl{PaperXI.GaussSums.sqrt_card_pos}\\

73&theorem&\nolinkurl{PaperXI.GaussSums.phase_norm}\\

80&theorem&\nolinkurl{PaperXI.GaussSums.phase_duality}\\

\bottomrule\end{longtable}

\section{\texorpdfstring{\nolinkurl{PaperXI.GeneralChain}}{PaperXI.GeneralChain}}\label{module:PaperXI.GeneralChain}

\textbf{Selected layer:} v0.44 parent.\quad\textbf{Source lines:} 89.\par\textbf{SHA-256:} {\scriptsize\nolinkurl{0e5a7f76765c2e3b753f9db36e3886906a9eac283a0583ebb7f07bcb404e33d7}}.

\VerbatimInput[fontsize=\scriptsize,breaklines=true,breakanywhere=true]{excerpts/D_GeneralChain.txt}

\textbf{Direct project imports:} \nolinkurl{PaperXI.MonodromyChain}.

\begin{longtable}{@{}p{0.07\textwidth}p{0.12\textwidth}p{0.75\textwidth}@{}}\toprule Line&Kind&Name\\\midrule\endhead

12&def&\nolinkurl{PaperXI.MonodromyChain.ChainSupported}\\

15&theorem&\nolinkurl{PaperXI.MonodromyChain.arbitrary_block_reconstruction}\\

29&theorem&\nolinkurl{PaperXI.MonodromyChain.projector_eq_single}\\

36&def&\nolinkurl{PaperXI.MonodromyChain.extractedTransfers}\\

42&theorem&\nolinkurl{PaperXI.MonodromyChain.coordinate_block_zero}\\

49&theorem&\nolinkurl{PaperXI.MonodromyChain.eq_extracted_shift}\\

75&theorem&\nolinkurl{PaperXI.MonodromyChain.arbitrary_nilpotent_bound}\\

80&theorem&\nolinkurl{PaperXI.MonodromyChain.arbitrary_power_paths}\\

\bottomrule\end{longtable}

\section{\texorpdfstring{\nolinkurl{PaperXI.GeneratedIrreducibility}}{PaperXI.GeneratedIrreducibility}}\label{module:PaperXI.GeneratedIrreducibility}

\textbf{Selected layer:} v0.44 parent.\quad\textbf{Source lines:} 55.\par\textbf{SHA-256:} {\scriptsize\nolinkurl{3aca9c09ed2ae98d3df1627466cdadb860be6807be1e7555d2743c1b534bd5b2}}.

\VerbatimInput[fontsize=\scriptsize,breaklines=true,breakanywhere=true]{excerpts/D_GeneratedIrreducibility.txt}

\textbf{Direct project imports:} \nolinkurl{PaperXI.CyclicAction}.

\begin{longtable}{@{}p{0.07\textwidth}p{0.12\textwidth}p{0.75\textwidth}@{}}\toprule Line&Kind&Name\\\midrule\endhead

11&theorem&\nolinkurl{PaperXI.GeneratedIrreducibility.nontrivial_space}\\

20&def&\nolinkurl{PaperXI.GeneratedIrreducibility.actionEquiv}\\

26&theorem&\nolinkurl{PaperXI.GeneratedIrreducibility.actionEquiv_apply}\\

29&theorem&\nolinkurl{PaperXI.GeneratedIrreducibility.actionEquiv_symm_apply}\\

32&theorem&\nolinkurl{PaperXI.GeneratedIrreducibility.stable_of_generators}\\

47&theorem&\nolinkurl{PaperXI.GeneratedIrreducibility.eq_bot_or_top}\\

\bottomrule\end{longtable}

\section{\texorpdfstring{\nolinkurl{PaperXI.GlobalHeisenbergBridge112}}{PaperXI.GlobalHeisenbergBridge112}}\label{module:PaperXI.GlobalHeisenbergBridge112}

\textbf{Selected layer:} Continuation112.\quad\textbf{Source lines:} 301.\par\textbf{SHA-256:} {\scriptsize\nolinkurl{b44f5fc5e576c356e04e33a27cfbcd1d0f70946a32dafba8389de23b83d83998}}.

\VerbatimInput[fontsize=\scriptsize,breaklines=true,breakanywhere=true]{excerpts/D_GlobalHeisenbergBridge112.txt}

\textbf{Direct project imports:} \nolinkurl{PaperXI.HeisenbergCharacterProfile111}, \nolinkurl{PaperXI.FiniteRepresentation}, \nolinkurl{PaperXI.TypeSideDescent}, \nolinkurl{PaperXI.SchurScalar}.

\begin{longtable}{@{}p{0.07\textwidth}p{0.12\textwidth}p{0.75\textwidth}@{}}\toprule Line&Kind&Name\\\midrule\endhead

23&theorem&\nolinkurl{PaperXI.GlobalHeisenbergBridge112.equiv_of_common_profile}\\

54&theorem&\nolinkurl{PaperXI.GlobalHeisenbergBridge112.equiv_of_common_quotient_form}\\

72&theorem&\nolinkurl{PaperXI.GlobalHeisenbergBridge112.equiv_of_global_central_lift}\\

105&theorem&\nolinkurl{PaperXI.GlobalHeisenbergBridge112.actual_eta_bridge_of_global_lift}\\

146&theorem&\nolinkurl{PaperXI.GlobalHeisenbergBridge112.generator_discrepancy}\\

168&theorem&\nolinkurl{PaperXI.GlobalHeisenbergBridge112.generator_discrepancy_pow}\\

179&theorem&\nolinkurl{PaperXI.GlobalHeisenbergBridge112.unit_eq_one_of_coprime_powers}\\

188&theorem&\nolinkurl{PaperXI.GlobalHeisenbergBridge112.generator_phase_eq_one_of_det_coprime}\\

224&theorem&\nolinkurl{PaperXI.GlobalHeisenbergBridge112.generator_equivariant_of_det_coprime}\\

247&theorem&\nolinkurl{PaperXI.GlobalHeisenbergBridge112.cyclic_equivariant_of_generator}\\

268&theorem&\nolinkurl{PaperXI.GlobalHeisenbergBridge112.equivariant_bridge_of_common_quotient_form}\\

\bottomrule\end{longtable}

\section{\texorpdfstring{\nolinkurl{PaperXI.GluedPacketSections}}{PaperXI.GluedPacketSections}}\label{module:PaperXI.GluedPacketSections}

\textbf{Selected layer:} v0.44 parent.\quad\textbf{Source lines:} 170.\par\textbf{SHA-256:} {\scriptsize\nolinkurl{5bf8047d59226ed80f93c4668b031aedab1df38bd5ee39456a9335f2685372bf}}.

\VerbatimInput[fontsize=\scriptsize,breaklines=true,breakanywhere=true]{excerpts/D_GluedPacketSections.txt}

\textbf{Direct project imports:} \nolinkurl{PaperXI.RepresentationGraph}.

\begin{longtable}{@{}p{0.07\textwidth}p{0.12\textwidth}p{0.75\textwidth}@{}}\toprule Line&Kind&Name\\\midrule\endhead

22&def&\nolinkurl{PaperXI.GluedPacketSections.sections}\\

28&theorem&\nolinkurl{PaperXI.GluedPacketSections.section_edge}\\

37&theorem&\nolinkurl{PaperXI.GluedPacketSections.section_path}\\

46&theorem&\nolinkurl{PaperXI.GluedPacketSections.root_evaluation_injective}\\

56&def&\nolinkurl{PaperXI.GluedPacketSections.loopFixed}\\

63&def&\nolinkurl{PaperXI.GluedPacketSections.extendRoot}\\

71&def&\nolinkurl{PaperXI.GluedPacketSections.evaluationEquiv}\\

90&def&\nolinkurl{PaperXI.GluedPacketSections.sectionRepresentation}\\

101&theorem&\nolinkurl{PaperXI.GluedPacketSections.sectionRepresentation_apply}\\

106&def&\nolinkurl{PaperXI.GluedPacketSections.ofGauge}\\

123&def&\nolinkurl{PaperXI.GluedPacketSections.equivOfGauge}\\

134&def&\nolinkurl{PaperXI.GluedPacketSections.gluedGauge}\\

138&def&\nolinkurl{PaperXI.GluedPacketSections.fromChords}\\

152&theorem&\nolinkurl{PaperXI.GluedPacketSections.fromChords_apply}\\

163&theorem&\nolinkurl{PaperXI.GluedPacketSections.gauge_unique_at_root}\\

\bottomrule\end{longtable}

\section{\texorpdfstring{\nolinkurl{PaperXI.GraphChordCharacters}}{PaperXI.GraphChordCharacters}}\label{module:PaperXI.GraphChordCharacters}

\textbf{Selected layer:} v0.44 parent.\quad\textbf{Source lines:} 144.\par\textbf{SHA-256:} {\scriptsize\nolinkurl{037a0bd69fdb62bbe05ce9760a8f4bcf242c6c01d167c45f6d9bf591cf5bcca0}}.

\VerbatimInput[fontsize=\scriptsize,breaklines=true,breakanywhere=true]{excerpts/D_GraphChordCharacters.txt}

\textbf{Direct project imports:} \nolinkurl{PaperXI.GraphTree}, \nolinkurl{PaperXI.GraphLoopRepresentation}, \nolinkurl{PaperXI.ScalarHolonomy}.

\begin{longtable}{@{}p{0.07\textwidth}p{0.12\textwidth}p{0.75\textwidth}@{}}\toprule Line&Kind&Name\\\midrule\endhead

24&def&\nolinkurl{PaperXI.GraphChordCharacters.weight}\\

29&theorem&\nolinkurl{PaperXI.GraphChordCharacters.weight_tree}\\

33&theorem&\nolinkurl{PaperXI.GraphChordCharacters.weight_chord}\\

39&def&\nolinkurl{PaperXI.GraphChordCharacters.edgeSystem}\\

43&theorem&\nolinkurl{PaperXI.GraphChordCharacters.edgeSystem_chord}\\

49&theorem&\nolinkurl{PaperXI.GraphChordCharacters.traversal_tree}\\

66&theorem&\nolinkurl{PaperXI.GraphChordCharacters.rootPath_transport}\\

81&def&\nolinkurl{PaperXI.GraphChordCharacters.chordLoop}\\

87&theorem&\nolinkurl{PaperXI.GraphChordCharacters.chordLoop_transport}\\

94&theorem&\nolinkurl{PaperXI.GraphChordCharacters.chordLoop_holonomy}\\

101&theorem&\nolinkurl{PaperXI.GraphChordCharacters.field_equiv_scalar}\\

111&theorem&\nolinkurl{PaperXI.GraphChordCharacters.holonomy_isScalar}\\

117&def&\nolinkurl{PaperXI.GraphChordCharacters.character}\\

121&theorem&\nolinkurl{PaperXI.GraphChordCharacters.character_chord}\\

128&theorem&\nolinkurl{PaperXI.GraphChordCharacters.character_tests_surjective}\\

135&theorem&\nolinkurl{PaperXI.GraphChordCharacters.homology_tests_surjective}\\

\bottomrule\end{longtable}

\section{\texorpdfstring{\nolinkurl{PaperXI.GraphGauge}}{PaperXI.GraphGauge}}\label{module:PaperXI.GraphGauge}

\textbf{Selected layer:} v0.44 parent.\quad\textbf{Source lines:} 190.\par\textbf{SHA-256:} {\scriptsize\nolinkurl{1750d323d264d4760717c62ab3aa154cbebefcdd24e80cc76ef22d7d8f5a572b}}.

\VerbatimInput[fontsize=\scriptsize,breaklines=true,breakanywhere=true]{excerpts/D_GraphGauge.txt}

\textbf{Direct project imports:} \nolinkurl{PaperXI.GraphTransport}, \nolinkurl{PaperXI.GraphTree}, \nolinkurl{PaperXI.GraphLoopRepresentation}, \nolinkurl{PaperXI.HolonomyResidual}.

\begin{longtable}{@{}p{0.07\textwidth}p{0.12\textwidth}p{0.75\textwidth}@{}}\toprule Line&Kind&Name\\\midrule\endhead

27&def&\nolinkurl{PaperXI.GraphGauge.cycle}\\

31&def&\nolinkurl{PaperXI.GraphGauge.rootGauge}\\

35&def&\nolinkurl{PaperXI.GraphGauge.EdgeCompatible}\\

39&def&\nolinkurl{PaperXI.GraphGauge.LoopCompatible}\\

43&def&\nolinkurl{PaperXI.GraphGauge.CycleCompatible}\\

47&theorem&\nolinkurl{PaperXI.GraphGauge.rootGauge_root}\\

53&theorem&\nolinkurl{PaperXI.GraphGauge.rootGauge_edge_iff}\\

66&theorem&\nolinkurl{PaperXI.GraphGauge.edgeCompatible_loops}\\

72&theorem&\nolinkurl{PaperXI.GraphGauge.rootGauge_compatible_of_loops}\\

79&theorem&\nolinkurl{PaperXI.GraphGauge.gauge_unique}\\

88&theorem&\nolinkurl{PaperXI.GraphGauge.based_iff_loops}\\

101&theorem&\nolinkurl{PaperXI.GraphGauge.rootGauge_parent}\\

110&theorem&\nolinkurl{PaperXI.GraphGauge.rootGauge_tree}\\

127&theorem&\nolinkurl{PaperXI.GraphGauge.rootGauge_compatible_iff_chords}\\

141&theorem&\nolinkurl{PaperXI.GraphGauge.based_iff_chords}\\

155&theorem&\nolinkurl{PaperXI.GraphGauge.connected_based_iff_chords}\\

163&theorem&\nolinkurl{PaperXI.GraphGauge.chords_iff_loops}\\

170&theorem&\nolinkurl{PaperXI.GraphGauge.tree_independence}\\

177&theorem&\nolinkurl{PaperXI.GraphGauge.unbased_iff_conjugates}\\

\bottomrule\end{longtable}

\section{\texorpdfstring{\nolinkurl{PaperXI.GraphLoop}}{PaperXI.GraphLoop}}\label{module:PaperXI.GraphLoop}

\textbf{Selected layer:} v0.44 parent.\quad\textbf{Source lines:} 106.\par\textbf{SHA-256:} {\scriptsize\nolinkurl{f38f6d71925623d713432f41101158265cdfa0733f6bd3242f335f62bf01d198}}.

\VerbatimInput[fontsize=\scriptsize,breaklines=true,breakanywhere=true]{excerpts/D_GraphLoop.txt}

\textbf{Direct project imports:} none.

\begin{longtable}{@{}p{0.07\textwidth}p{0.12\textwidth}p{0.75\textwidth}@{}}\toprule Line&Kind&Name\\\midrule\endhead

27&abbrev&\nolinkurl{PaperXI.GraphLoop.vertex}\\

31&def&\nolinkurl{PaperXI.GraphLoop.pathClass}\\

35&theorem&\nolinkurl{PaperXI.GraphLoop.pathClass_nil}\\

38&theorem&\nolinkurl{PaperXI.GraphLoop.pathClass_comp}\\

43&theorem&\nolinkurl{PaperXI.GraphLoop.pathClass_reverse}\\

48&theorem&\nolinkurl{PaperXI.GraphLoop.pathClass_surjective}\\

53&abbrev&\nolinkurl{PaperXI.GraphLoop.Loop}\\

56&def&\nolinkurl{PaperXI.GraphLoop.loopClass}\\

59&theorem&\nolinkurl{PaperXI.GraphLoop.loopClass_nil}\\

63&theorem&\nolinkurl{PaperXI.GraphLoop.loopClass_comp}\\

67&theorem&\nolinkurl{PaperXI.GraphLoop.loopClass_reverse}\\

71&theorem&\nolinkurl{PaperXI.GraphLoop.loopClass_surjective}\\

75&abbrev&\nolinkurl{PaperXI.GraphLoop.H1}\\

78&def&\nolinkurl{PaperXI.GraphLoop.homologyClass}\\

80&theorem&\nolinkurl{PaperXI.GraphLoop.homologyClass_mul}\\

84&def&\nolinkurl{PaperXI.GraphLoop.factorCharacter}\\

87&theorem&\nolinkurl{PaperXI.GraphLoop.factorCharacter_apply_class}\\

92&theorem&\nolinkurl{PaperXI.GraphLoop.factorCharacter_unique}\\

99&theorem&\nolinkurl{PaperXI.GraphLoop.character_ext_paths}\\

\bottomrule\end{longtable}

\section{\texorpdfstring{\nolinkurl{PaperXI.GraphLoopRepresentation}}{PaperXI.GraphLoopRepresentation}}\label{module:PaperXI.GraphLoopRepresentation}

\textbf{Selected layer:} v0.44 parent.\quad\textbf{Source lines:} 138.\par\textbf{SHA-256:} {\scriptsize\nolinkurl{8e6bb58b55ba016228d5ae3cd85b19e40ccde0abd116f4a13101a8d8939a4e95}}.

\VerbatimInput[fontsize=\scriptsize,breaklines=true,breakanywhere=true]{excerpts/D_GraphLoopRepresentation.txt}

\textbf{Direct project imports:} \nolinkurl{PaperXI.GraphLoop}, \nolinkurl{PaperXI.GraphTransport}.

\begin{longtable}{@{}p{0.07\textwidth}p{0.12\textwidth}p{0.75\textwidth}@{}}\toprule Line&Kind&Name\\\midrule\endhead

25&structure&\nolinkurl{PaperXI.GraphLoop.Fiber}\\

28&instance&\nolinkurl{PaperXI.GraphLoop.fiberGroupoid}\\

40&def&\nolinkurl{PaperXI.GraphLoop.edgePrefunctor}\\

46&def&\nolinkurl{PaperXI.GraphLoop.transportFunctor}\\

51&def&\nolinkurl{PaperXI.GraphLoop.transport}\\

55&theorem&\nolinkurl{PaperXI.GraphLoop.transport_id}\\

60&theorem&\nolinkurl{PaperXI.GraphLoop.transport_comp}\\

66&theorem&\nolinkurl{PaperXI.GraphLoop.transport_edge}\\

73&theorem&\nolinkurl{PaperXI.GraphLoop.transport_pathClass}\\

90&def&\nolinkurl{PaperXI.GraphLoop.holonomy}\\

96&theorem&\nolinkurl{PaperXI.GraphLoop.holonomy_loopClass}\\

105&def&\nolinkurl{PaperXI.GraphLoop.Conjugates}\\

112&theorem&\nolinkurl{PaperXI.GraphLoop.conjugates_iff_paths}\\

129&theorem&\nolinkurl{PaperXI.GraphLoop.exists_conjugates_iff_paths}\\

\bottomrule\end{longtable}

\section{\texorpdfstring{\nolinkurl{PaperXI.GraphResidual}}{PaperXI.GraphResidual}}\label{module:PaperXI.GraphResidual}

\textbf{Selected layer:} v0.44 parent.\quad\textbf{Source lines:} 274.\par\textbf{SHA-256:} {\scriptsize\nolinkurl{ffc35ff3cd92cf45144e94c501a36af8567e4bb7531e9dc21eb81401385adbf8}}.

\VerbatimInput[fontsize=\scriptsize,breaklines=true,breakanywhere=true]{excerpts/D_GraphResidual.txt}

\textbf{Direct project imports:} \nolinkurl{PaperXI.ContinuousHolonomyResidual}, \nolinkurl{PaperXI.GraphTransport}, \nolinkurl{PaperXI.GraphTree}.

\begin{longtable}{@{}p{0.07\textwidth}p{0.12\textwidth}p{0.75\textwidth}@{}}\toprule Line&Kind&Name\\\midrule\endhead

24&def&\nolinkurl{PaperXI.GraphResidual.pathL}\\

28&theorem&\nolinkurl{PaperXI.GraphResidual.pathL_toLinearEquiv}\\

33&theorem&\nolinkurl{PaperXI.GraphResidual.pathL_apply}\\

38&theorem&\nolinkurl{PaperXI.GraphResidual.pathL_symm_apply}\\

43&theorem&\nolinkurl{PaperXI.GraphResidual.pathL_reverse}\\

50&theorem&\nolinkurl{PaperXI.GraphResidual.pathL_comp}\\

60&def&\nolinkurl{PaperXI.GraphResidual.gauge}\\

67&theorem&\nolinkurl{PaperXI.GraphResidual.gauge_toLinearEquiv}\\

74&theorem&\nolinkurl{PaperXI.GraphResidual.gauge_apply}\\

81&theorem&\nolinkurl{PaperXI.GraphResidual.gauge_root}\\

90&theorem&\nolinkurl{PaperXI.GraphResidual.gauge_parent}\\

106&def&\nolinkurl{PaperXI.GraphResidual.edgeResidual}\\

116&def&\nolinkurl{PaperXI.GraphResidual.cycleDefect}\\

125&theorem&\nolinkurl{PaperXI.GraphResidual.cycleDefect_toLinearMap}\\

134&theorem&\nolinkurl{PaperXI.GraphResidual.edgeResidual_toLinearMap}\\

144&theorem&\nolinkurl{PaperXI.GraphResidual.cycleDefect_loop}\\

162&theorem&\nolinkurl{PaperXI.GraphResidual.residual_factorization}\\

173&theorem&\nolinkurl{PaperXI.GraphResidual.residual_zero_iff}\\

184&theorem&\nolinkurl{PaperXI.GraphResidual.residual_norm_bound}\\

196&theorem&\nolinkurl{PaperXI.GraphResidual.uniform_bound}\\

212&theorem&\nolinkurl{PaperXI.GraphResidual.uniform_product_bound}\\

228&theorem&\nolinkurl{PaperXI.GraphResidual.uniform_chord_bound}\\

246&theorem&\nolinkurl{PaperXI.GraphResidual.uniform_chord_product_bound}\\

264&theorem&\nolinkurl{PaperXI.GraphResidual.connected_residual_bound}\\

\bottomrule\end{longtable}

\section{\texorpdfstring{\nolinkurl{PaperXI.GraphTransport}}{PaperXI.GraphTransport}}\label{module:PaperXI.GraphTransport}

\textbf{Selected layer:} v0.44 parent.\quad\textbf{Source lines:} 150.\par\textbf{SHA-256:} {\scriptsize\nolinkurl{72c2feeee0c9f9e98eeed01a1ff511d439bdfdac6c3bdb2732290dcc1f8151fb}}.

\VerbatimInput[fontsize=\scriptsize,breaklines=true,breakanywhere=true]{excerpts/D_GraphTransport.txt}

\textbf{Direct project imports:} none.

\begin{longtable}{@{}p{0.07\textwidth}p{0.12\textwidth}p{0.75\textwidth}@{}}\toprule Line&Kind&Name\\\midrule\endhead

17&abbrev&\nolinkurl{PaperXI.GraphTransport.Traversal}\\

20&abbrev&\nolinkurl{PaperXI.GraphTransport.Walk}\\

25&def&\nolinkurl{PaperXI.GraphTransport.symmetrifyAdd}\\

29&def&\nolinkurl{PaperXI.GraphTransport.symmetrifyModule}\\

33&def&\nolinkurl{PaperXI.GraphTransport.edge}\\

39&theorem&\nolinkurl{PaperXI.GraphTransport.edge_inl}\\

42&theorem&\nolinkurl{PaperXI.GraphTransport.edge_inr}\\

45&theorem&\nolinkurl{PaperXI.GraphTransport.edge_reverse}\\

51&def&\nolinkurl{PaperXI.GraphTransport.path}\\

56&theorem&\nolinkurl{PaperXI.GraphTransport.nil}\\

59&theorem&\nolinkurl{PaperXI.GraphTransport.cons}\\

63&theorem&\nolinkurl{PaperXI.GraphTransport.toPath}\\

68&theorem&\nolinkurl{PaperXI.GraphTransport.comp}\\

77&theorem&\nolinkurl{PaperXI.GraphTransport.reverse}\\

85&theorem&\nolinkurl{PaperXI.GraphTransport.path_reverse_cancel}\\

92&theorem&\nolinkurl{PaperXI.GraphTransport.reverse_path_cancel}\\

100&theorem&\nolinkurl{PaperXI.GraphTransport.cancel}\\

110&theorem&\nolinkurl{PaperXI.GraphTransport.gauge_edge}\\

125&theorem&\nolinkurl{PaperXI.GraphTransport.gauge_path}\\

139&theorem&\nolinkurl{PaperXI.GraphTransport.gauge_path_equiv}\\

\bottomrule\end{longtable}

\section{\texorpdfstring{\nolinkurl{PaperXI.GraphTree}}{PaperXI.GraphTree}}\label{module:PaperXI.GraphTree}

\textbf{Selected layer:} v0.44 parent.\quad\textbf{Source lines:} 245.\par\textbf{SHA-256:} {\scriptsize\nolinkurl{37a1cea3da4966b2d631d1f44a997095037c88ca15a4b3001a6b8efa607dd486}}.

\VerbatimInput[fontsize=\scriptsize,breaklines=true,breakanywhere=true]{excerpts/D_GraphTree.txt}

\textbf{Direct project imports:} none.

\begin{longtable}{@{}p{0.07\textwidth}p{0.12\textwidth}p{0.75\textwidth}@{}}\toprule Line&Kind&Name\\\midrule\endhead

18&abbrev&\nolinkurl{PaperXI.GraphTree.Edge}\\

21&abbrev&\nolinkurl{PaperXI.GraphTree.Traversal}\\

24&abbrev&\nolinkurl{PaperXI.GraphTree.Path}\\

27&def&\nolinkurl{PaperXI.GraphTree.Connected}\\

30&structure&\nolinkurl{PaperXI.GraphTree.Rooted}\\

40&theorem&\nolinkurl{PaperXI.GraphTree.exists_length}\\

45&def&\nolinkurl{PaperXI.GraphTree.distance}\\

49&theorem&\nolinkurl{PaperXI.GraphTree.distance_le}\\

53&theorem&\nolinkurl{PaperXI.GraphTree.distance_root}\\

57&theorem&\nolinkurl{PaperXI.GraphTree.exists_predecessor}\\

71&def&\nolinkurl{PaperXI.GraphTree.ofConnected}\\

83&def&\nolinkurl{PaperXI.GraphTree.rootPath}\\

92&theorem&\nolinkurl{PaperXI.GraphTree.rootPath_root}\\

96&theorem&\nolinkurl{PaperXI.GraphTree.rootPath_parent}\\

102&theorem&\nolinkurl{PaperXI.GraphTree.rooted_induction}\\

112&def&\nolinkurl{PaperXI.GraphTree.underlying}\\

118&def&\nolinkurl{PaperXI.GraphTree.selected}\\

120&theorem&\nolinkurl{PaperXI.GraphTree.underlying_endpoints}\\

125&theorem&\nolinkurl{PaperXI.GraphTree.selected_injective}\\

152&def&\nolinkurl{PaperXI.GraphTree.treeEdges}\\

155&def&\nolinkurl{PaperXI.GraphTree.Chord}\\

157&theorem&\nolinkurl{PaperXI.GraphTree.mem_treeEdges}\\

160&theorem&\nolinkurl{PaperXI.GraphTree.treeEdges_induction}\\

167&theorem&\nolinkurl{PaperXI.GraphTree.selected_not_loop}\\

183&def&\nolinkurl{PaperXI.GraphTree.Supported}\\

188&theorem&\nolinkurl{PaperXI.GraphTree.rootPath_supported}\\

196&theorem&\nolinkurl{PaperXI.GraphTree.treeEdges_no_loop}\\

200&theorem&\nolinkurl{PaperXI.GraphTree.tree_card}\\

211&theorem&\nolinkurl{PaperXI.GraphTree.chord_count}\\

226&theorem&\nolinkurl{PaperXI.GraphTree.edge_count_lower_bound}\\

233&theorem&\nolinkurl{PaperXI.GraphTree.chord_card}\\

239&theorem&\nolinkurl{PaperXI.GraphTree.chord_card_int}\\

\bottomrule\end{longtable}

\section{\texorpdfstring{\nolinkurl{PaperXI.GraphTreePaths}}{PaperXI.GraphTreePaths}}\label{module:PaperXI.GraphTreePaths}

\textbf{Selected layer:} v0.44 parent.\quad\textbf{Source lines:} 87.\par\textbf{SHA-256:} {\scriptsize\nolinkurl{732095eebf2b794f7d31d503a7928bfde9b7d40da777fd5c99b3d1cccd1e2e18}}.

\VerbatimInput[fontsize=\scriptsize,breaklines=true,breakanywhere=true]{excerpts/D_GraphTreePaths.txt}

\textbf{Direct project imports:} \nolinkurl{PaperXI.GraphTree}, \nolinkurl{PaperXI.GraphTransport}.

\begin{longtable}{@{}p{0.07\textwidth}p{0.12\textwidth}p{0.75\textwidth}@{}}\toprule Line&Kind&Name\\\midrule\endhead

20&theorem&\nolinkurl{PaperXI.GraphTreePaths.parent_transport}\\

27&theorem&\nolinkurl{PaperXI.GraphTreePaths.tree_transport}\\

55&theorem&\nolinkurl{PaperXI.GraphTreePaths.traversal_transport}\\

70&theorem&\nolinkurl{PaperXI.GraphTreePaths.supported_transport}\\

82&theorem&\nolinkurl{PaperXI.GraphTreePaths.tree_loop_identity}\\

\bottomrule\end{longtable}

\section{\texorpdfstring{\nolinkurl{PaperXI.GroupFixed}}{PaperXI.GroupFixed}}\label{module:PaperXI.GroupFixed}

\textbf{Selected layer:} v0.44 parent.\quad\textbf{Source lines:} 85.\par\textbf{SHA-256:} {\scriptsize\nolinkurl{32a4e3c21abaf0e27169d70976a894f79172987de4954913d8fc607fd925e51b}}.

\VerbatimInput[fontsize=\scriptsize,breaklines=true,breakanywhere=true]{excerpts/D_GroupFixed.txt}

\textbf{Direct project imports:} \nolinkurl{PaperXI.CycleFixed}, \nolinkurl{PaperXI.InertiaQuotient}.

\begin{longtable}{@{}p{0.07\textwidth}p{0.12\textwidth}p{0.75\textwidth}@{}}\toprule Line&Kind&Name\\\midrule\endhead

12&theorem&\nolinkurl{PaperXI.GroupFixed.fixed_of_generators}\\

22&theorem&\nolinkurl{PaperXI.GroupFixed.invariants_eq_common}\\

46&theorem&\nolinkurl{PaperXI.GroupFixed.orbit_basis_steps}\\

63&theorem&\nolinkurl{PaperXI.GroupFixed.packet_finrank}\\

\bottomrule\end{longtable}

\section{\texorpdfstring{\nolinkurl{PaperXI.HeisenbergCardinality}}{PaperXI.HeisenbergCardinality}}\label{module:PaperXI.HeisenbergCardinality}

\textbf{Selected layer:} v0.44 parent.\quad\textbf{Source lines:} 65.\par\textbf{SHA-256:} {\scriptsize\nolinkurl{3325cc5f2651ba031ec6ab41f813ec4f9e63afea875a3e4cac9827997fef3225}}.

\VerbatimInput[fontsize=\scriptsize,breaklines=true,breakanywhere=true]{excerpts/D_HeisenbergCardinality.txt}

\textbf{Direct project imports:} \nolinkurl{PaperXI.HeisenbergForm}, \nolinkurl{PaperXI.HeisenbergGeometry}.

\begin{longtable}{@{}p{0.07\textwidth}p{0.12\textwidth}p{0.75\textwidth}@{}}\toprule Line&Kind&Name\\\midrule\endhead

17&theorem&\nolinkurl{PaperXI.HeisenbergCardinality.form_isAlt}\\

19&theorem&\nolinkurl{PaperXI.HeisenbergCardinality.form_separatingLeft}\\

29&theorem&\nolinkurl{PaperXI.HeisenbergCardinality.form_nondegenerate}\\

37&theorem&\nolinkurl{PaperXI.HeisenbergCardinality.form_forall_eq_zero_iff}\\

41&theorem&\nolinkurl{PaperXI.HeisenbergCardinality.maximal_orthogonal_eq}\\

48&theorem&\nolinkurl{PaperXI.HeisenbergCardinality.finrank_eq}\\

54&theorem&\nolinkurl{PaperXI.HeisenbergCardinality.natCard_eq_pow}\\

59&theorem&\nolinkurl{PaperXI.HeisenbergCardinality.natCard_eq}\\

\bottomrule\end{longtable}

\section{\texorpdfstring{\nolinkurl{PaperXI.HeisenbergCenter}}{PaperXI.HeisenbergCenter}}\label{module:PaperXI.HeisenbergCenter}

\textbf{Selected layer:} v0.44 parent.\quad\textbf{Source lines:} 64.\par\textbf{SHA-256:} {\scriptsize\nolinkurl{f2e8ae8af020e8557e7b08bb957b7145733776c6b5c83c37218c1e0e184f4966}}.

\VerbatimInput[fontsize=\scriptsize,breaklines=true,breakanywhere=true]{excerpts/D_HeisenbergCenter.txt}

\textbf{Direct project imports:} \nolinkurl{PaperXI.HeisenbergGroup}, \nolinkurl{PaperXI.HeisenbergCardinality}.

\begin{longtable}{@{}p{0.07\textwidth}p{0.12\textwidth}p{0.75\textwidth}@{}}\toprule Line&Kind&Name\\\midrule\endhead

15&theorem&\nolinkurl{PaperXI.HeisenbergCenter.commute_iff}\\

35&theorem&\nolinkurl{PaperXI.HeisenbergCenter.mem_center_iff}\\

51&theorem&\nolinkurl{PaperXI.HeisenbergCenter.center_eq_centralElement}\\

\bottomrule\end{longtable}

\section{\texorpdfstring{\nolinkurl{PaperXI.HeisenbergCharacterProfile111}}{PaperXI.HeisenbergCharacterProfile111}}\label{module:PaperXI.HeisenbergCharacterProfile111}

\textbf{Selected layer:} Continuation111.\quad\textbf{Source lines:} 317.\par\textbf{SHA-256:} {\scriptsize\nolinkurl{1b222ac6ed229dac1c3c2682d5340cc880a2699c3c0ca3c9cbc80ae11bcb712e}}.

\VerbatimInput[fontsize=\scriptsize,breaklines=true,breakanywhere=true]{excerpts/D_HeisenbergCharacterProfile111.txt}

\textbf{Direct project imports:} \nolinkurl{PaperXI.DegreeIndex110}.

\begin{longtable}{@{}p{0.07\textwidth}p{0.12\textwidth}p{0.75\textwidth}@{}}\toprule Line&Kind&Name\\\midrule\endhead

21&def&\nolinkurl{PaperXI.HeisenbergCharacterProfile111.coordinate}\\

24&theorem&\nolinkurl{PaperXI.HeisenbergCharacterProfile111.coordinate_surjective}\\

31&theorem&\nolinkurl{PaperXI.HeisenbergCharacterProfile111.coordinate_eq_zero_iff}\\

44&theorem&\nolinkurl{PaperXI.HeisenbergCharacterProfile111.commutator_mem}\\

52&def&\nolinkurl{PaperXI.HeisenbergCharacterProfile111.characterCommutator}\\

59&theorem&\nolinkurl{PaperXI.HeisenbergCharacterProfile111.commutator_detection_of_nondegenerate_form}\\

94&theorem&\nolinkurl{PaperXI.HeisenbergCharacterProfile111.trace_support_of_nondegenerate_form}\\

116&def&\nolinkurl{PaperXI.HeisenbergCharacterProfile111.transportedQuotient}\\

120&theorem&\nolinkurl{PaperXI.HeisenbergCharacterProfile111.coordinate_transport}\\

124&def&\nolinkurl{PaperXI.HeisenbergCharacterProfile111.pulledBackCharacter}\\

128&theorem&\nolinkurl{PaperXI.HeisenbergCharacterProfile111.characterCommutator_transport}\\

140&theorem&\nolinkurl{PaperXI.HeisenbergCharacterProfile111.form_transport}\\

160&theorem&\nolinkurl{PaperXI.HeisenbergCharacterProfile111.character_intertwinedAt_iff_covariance}\\

187&theorem&\nolinkurl{PaperXI.HeisenbergCharacterProfile111.character_values_of_intertwinedAt}\\

216&instance&\nolinkurl{PaperXI.HeisenbergCharacterProfile111.actualScalarSubgroup_normal}\\

237&theorem&\nolinkurl{PaperXI.HeisenbergCharacterProfile111.actual_trace_support_of_quotient_form}\\

261&theorem&\nolinkurl{PaperXI.HeisenbergCharacterProfile111.actual_covariance_of_character_intertwining}\\

283&theorem&\nolinkurl{PaperXI.HeisenbergCharacterProfile111.actual_eta_dimension_from_quotient_form}\\

\bottomrule\end{longtable}

\section{\texorpdfstring{\nolinkurl{PaperXI.HeisenbergCochain}}{PaperXI.HeisenbergCochain}}\label{module:PaperXI.HeisenbergCochain}

\textbf{Selected layer:} v0.44 parent.\quad\textbf{Source lines:} 69.\par\textbf{SHA-256:} {\scriptsize\nolinkurl{ba7c5f4ea39d1e68eadd6f852cf159e67a2b0c454eb6083d92f7d48ca9e92f8f}}.

\VerbatimInput[fontsize=\scriptsize,breaklines=true,breakanywhere=true]{excerpts/D_HeisenbergCochain.txt}

\textbf{Direct project imports:} \nolinkurl{PaperXI.BilinearHeisenberg}.

\begin{longtable}{@{}p{0.07\textwidth}p{0.12\textwidth}p{0.75\textwidth}@{}}\toprule Line&Kind&Name\\\midrule\endhead

14&def&\nolinkurl{PaperXI.HeisenbergCochain.equiv}\\

28&theorem&\nolinkurl{PaperXI.HeisenbergCochain.equiv_center}\\

33&theorem&\nolinkurl{PaperXI.HeisenbergCochain.equiv_vector}\\

41&theorem&\nolinkurl{PaperXI.HeisenbergCochain.square}\\

51&theorem&\nolinkurl{PaperXI.HeisenbergCochain.diagonal_preserved}\\

62&theorem&\nolinkurl{PaperXI.HeisenbergCochain.no_equiv_of_diagonal_mismatch}\\

\bottomrule\end{longtable}

\section{\texorpdfstring{\nolinkurl{PaperXI.HeisenbergComparison}}{PaperXI.HeisenbergComparison}}\label{module:PaperXI.HeisenbergComparison}

\textbf{Selected layer:} v0.44 parent.\quad\textbf{Source lines:} 92.\par\textbf{SHA-256:} {\scriptsize\nolinkurl{545124ccd81c40a4378e2d854940a0ddb3402c6081775d92751fa13fa0a7f96b}}.

\VerbatimInput[fontsize=\scriptsize,breaklines=true,breakanywhere=true]{excerpts/D_HeisenbergComparison.txt}

\textbf{Direct project imports:} \nolinkurl{PaperXI.HeisenbergTransport}, \nolinkurl{PaperXI.SchrodingerUniversal}.

\begin{longtable}{@{}p{0.07\textwidth}p{0.12\textwidth}p{0.75\textwidth}@{}}\toprule Line&Kind&Name\\\midrule\endhead

21&def&\nolinkurl{PaperXI.HeisenbergComparison.equivAlong}\\

32&theorem&\nolinkurl{PaperXI.HeisenbergComparison.equivAlong_intertwines}\\

46&theorem&\nolinkurl{PaperXI.HeisenbergComparison.exists_unique_along}\\

75&theorem&\nolinkurl{PaperXI.HeisenbergComparison.exists_unique_of_isometry}\\

\bottomrule\end{longtable}

\section{\texorpdfstring{\nolinkurl{PaperXI.HeisenbergDimension}}{PaperXI.HeisenbergDimension}}\label{module:PaperXI.HeisenbergDimension}

\textbf{Selected layer:} v0.44 parent.\quad\textbf{Source lines:} 41.\par\textbf{SHA-256:} {\scriptsize\nolinkurl{08109693bdd53c46678bff188224863dda5f276ce70dee638238a08c28317a02}}.

\VerbatimInput[fontsize=\scriptsize,breaklines=true,breakanywhere=true]{excerpts/D_HeisenbergDimension.txt}

\textbf{Direct project imports:} \nolinkurl{PaperXI.HeisenbergCardinality}.

\begin{longtable}{@{}p{0.07\textwidth}p{0.12\textwidth}p{0.75\textwidth}@{}}\toprule Line&Kind&Name\\\midrule\endhead

16&theorem&\nolinkurl{PaperXI.HeisenbergDimension.natCard_dual}\\

21&theorem&\nolinkurl{PaperXI.HeisenbergDimension.natCard_symplectic}\\

24&theorem&\nolinkurl{PaperXI.HeisenbergDimension.sqrt_natCard_symplectic}\\

31&theorem&\nolinkurl{PaperXI.HeisenbergDimension.finrank_function}\\

34&theorem&\nolinkurl{PaperXI.HeisenbergDimension.finrank_function_pow}\\

37&theorem&\nolinkurl{PaperXI.HeisenbergDimension.finrank_function_sqrt}\\

\bottomrule\end{longtable}

\section{\texorpdfstring{\nolinkurl{PaperXI.HeisenbergFinite}}{PaperXI.HeisenbergFinite}}\label{module:PaperXI.HeisenbergFinite}

\textbf{Selected layer:} v0.44 parent.\quad\textbf{Source lines:} 32.\par\textbf{SHA-256:} {\scriptsize\nolinkurl{d3703379c8db33d8675a342dc898f778aa75b1599529098f11b2af1df156e553}}.

\VerbatimInput[fontsize=\scriptsize,breaklines=true,breakanywhere=true]{excerpts/D_HeisenbergFinite.txt}

\textbf{Direct project imports:} \nolinkurl{PaperXI.HeisenbergGroup}, \nolinkurl{PaperXI.HeisenbergDimension}.

\begin{longtable}{@{}p{0.07\textwidth}p{0.12\textwidth}p{0.75\textwidth}@{}}\toprule Line&Kind&Name\\\midrule\endhead

16&def&\nolinkurl{PaperXI.HeisenbergFinite.coordinates}\\

22&instance&\nolinkurl{PaperXI.HeisenbergFinite.instFinite}\\

27&theorem&\nolinkurl{PaperXI.HeisenbergFinite.natCard_eq}\\

\bottomrule\end{longtable}

\section{\texorpdfstring{\nolinkurl{PaperXI.HeisenbergForm}}{PaperXI.HeisenbergForm}}\label{module:PaperXI.HeisenbergForm}

\textbf{Selected layer:} v0.44 parent.\quad\textbf{Source lines:} 50.\par\textbf{SHA-256:} {\scriptsize\nolinkurl{b9d055bf58ee5326c45c49e9267dcd9b113d6b1d2b3f0f82d5e7cd909890230c}}.

\VerbatimInput[fontsize=\scriptsize,breaklines=true,breakanywhere=true]{excerpts/D_HeisenbergForm.txt}

\textbf{Direct project imports:} none.

\begin{longtable}{@{}p{0.07\textwidth}p{0.12\textwidth}p{0.75\textwidth}@{}}\toprule Line&Kind&Name\\\midrule\endhead

13&def&\nolinkurl{PaperXI.HeisenbergForm.form}\\

41&theorem&\nolinkurl{PaperXI.HeisenbergForm.form_apply}\\

44&theorem&\nolinkurl{PaperXI.HeisenbergForm.form_self}\\

47&theorem&\nolinkurl{PaperXI.HeisenbergForm.form_skew}\\

\bottomrule\end{longtable}

\section{\texorpdfstring{\nolinkurl{PaperXI.HeisenbergFrequency}}{PaperXI.HeisenbergFrequency}}\label{module:PaperXI.HeisenbergFrequency}

\textbf{Selected layer:} v0.44 parent.\quad\textbf{Source lines:} 197.\par\textbf{SHA-256:} {\scriptsize\nolinkurl{511658b8ac1ca7f93b8002ee2a626418b076f2cd34875bd57bbad47a4deb6b59}}.

\VerbatimInput[fontsize=\scriptsize,breaklines=true,breakanywhere=true]{excerpts/D_HeisenbergFrequency.txt}

\textbf{Direct project imports:} \nolinkurl{PaperXI.DualCharacterSums}, \nolinkurl{PaperXI.HeisenbergNormalForm}, \nolinkurl{PaperXI.HeisenbergRepresentation}.

\begin{longtable}{@{}p{0.07\textwidth}p{0.12\textwidth}p{0.75\textwidth}@{}}\toprule Line&Kind&Name\\\midrule\endhead

19&def&\nolinkurl{PaperXI.HeisenbergFrequency.projector}\\

24&theorem&\nolinkurl{PaperXI.HeisenbergFrequency.projector_apply}\\

31&theorem&\nolinkurl{PaperXI.HeisenbergFrequency.eigen}\\

63&theorem&\nolinkurl{PaperXI.HeisenbergFrequency.complete}\\

92&theorem&\nolinkurl{PaperXI.HeisenbergFrequency.projector_of_eigen}\\

114&theorem&\nolinkurl{PaperXI.HeisenbergFrequency.orthogonal_idempotents}\\

127&theorem&\nolinkurl{PaperXI.HeisenbergFrequency.modulation_translation_action}\\

147&theorem&\nolinkurl{PaperXI.HeisenbergFrequency.translation_covariance}\\

169&theorem&\nolinkurl{PaperXI.HeisenbergFrequency.nonzero_zero_fiber}\\

\bottomrule\end{longtable}

\section{\texorpdfstring{\nolinkurl{PaperXI.HeisenbergGeometry}}{PaperXI.HeisenbergGeometry}}\label{module:PaperXI.HeisenbergGeometry}

\textbf{Selected layer:} v0.44 parent.\quad\textbf{Source lines:} 82.\par\textbf{SHA-256:} {\scriptsize\nolinkurl{df27cc30cbdcf1f4f8ac3ac85c90e2266dcc866bcd81ea3fbdc18d9ac3022c7c}}.

\VerbatimInput[fontsize=\scriptsize,breaklines=true,breakanywhere=true]{excerpts/D_HeisenbergGeometry.txt}

\textbf{Direct project imports:} none.

\begin{longtable}{@{}p{0.07\textwidth}p{0.12\textwidth}p{0.75\textwidth}@{}}\toprule Line&Kind&Name\\\midrule\endhead

21&def&\nolinkurl{PaperXI.HeisenbergGeometry.IsIsotropic}\\

25&def&\nolinkurl{PaperXI.HeisenbergGeometry.IsMaximalIsotropic}\\

28&theorem&\nolinkurl{PaperXI.HeisenbergGeometry.isotropic_iff_le_orthogonal}\\

37&theorem&\nolinkurl{PaperXI.HeisenbergGeometry.isotropic_sup_span}\\

50&theorem&\nolinkurl{PaperXI.HeisenbergGeometry.orthogonal_eq}\\

60&theorem&\nolinkurl{PaperXI.HeisenbergGeometry.maximal_iff_orthogonal_eq}\\

72&theorem&\nolinkurl{PaperXI.HeisenbergGeometry.two_mul_finrank}\\

\bottomrule\end{longtable}

\section{\texorpdfstring{\nolinkurl{PaperXI.HeisenbergGroup}}{PaperXI.HeisenbergGroup}}\label{module:PaperXI.HeisenbergGroup}

\textbf{Selected layer:} v0.44 parent.\quad\textbf{Source lines:} 109.\par\textbf{SHA-256:} {\scriptsize\nolinkurl{1a146672c602d8cacb7184a8bf364365676d1766a0ffd7af7a1e1c56652776dd}}.

\VerbatimInput[fontsize=\scriptsize,breaklines=true,breakanywhere=true]{excerpts/D_HeisenbergGroup.txt}

\textbf{Direct project imports:} none.

\begin{longtable}{@{}p{0.07\textwidth}p{0.12\textwidth}p{0.75\textwidth}@{}}\toprule Line&Kind&Name\\\midrule\endhead

12&structure&\nolinkurl{PaperXI.Heisenberg.Element}\\

19&instance&\nolinkurl{PaperXI.Heisenberg.instMul}\\

23&instance&\nolinkurl{PaperXI.Heisenberg.instOne}\\

26&instance&\nolinkurl{PaperXI.Heisenberg.instInv}\\

29&theorem&\nolinkurl{PaperXI.Heisenberg.mul_position}\\

32&theorem&\nolinkurl{PaperXI.Heisenberg.mul_frequency}\\

35&theorem&\nolinkurl{PaperXI.Heisenberg.mul_central}\\

39&theorem&\nolinkurl{PaperXI.Heisenberg.one_position}\\

40&theorem&\nolinkurl{PaperXI.Heisenberg.one_frequency}\\

41&theorem&\nolinkurl{PaperXI.Heisenberg.one_central}\\

42&theorem&\nolinkurl{PaperXI.Heisenberg.inv_position}\\

43&theorem&\nolinkurl{PaperXI.Heisenberg.inv_frequency}\\

44&theorem&\nolinkurl{PaperXI.Heisenberg.inv_central}\\

47&instance&\nolinkurl{PaperXI.Heisenberg.instGroup}\\

59&def&\nolinkurl{PaperXI.Heisenberg.centralElement}\\

62&def&\nolinkurl{PaperXI.Heisenberg.translationElement}\\

65&def&\nolinkurl{PaperXI.Heisenberg.modulationElement}\\

67&theorem&\nolinkurl{PaperXI.Heisenberg.centralElement_position}\\

69&theorem&\nolinkurl{PaperXI.Heisenberg.centralElement_frequency}\\

71&theorem&\nolinkurl{PaperXI.Heisenberg.centralElement_central}\\

73&theorem&\nolinkurl{PaperXI.Heisenberg.translationElement_position}\\

75&theorem&\nolinkurl{PaperXI.Heisenberg.translationElement_frequency}\\

77&theorem&\nolinkurl{PaperXI.Heisenberg.translationElement_central}\\

79&theorem&\nolinkurl{PaperXI.Heisenberg.modulationElement_position}\\

81&theorem&\nolinkurl{PaperXI.Heisenberg.modulationElement_frequency}\\

83&theorem&\nolinkurl{PaperXI.Heisenberg.modulationElement_central}\\

86&theorem&\nolinkurl{PaperXI.Heisenberg.centralElement_zero}\\

87&theorem&\nolinkurl{PaperXI.Heisenberg.translationElement_zero}\\

88&theorem&\nolinkurl{PaperXI.Heisenberg.modulationElement_zero}\\

91&theorem&\nolinkurl{PaperXI.Heisenberg.centralElement_mul}\\

95&theorem&\nolinkurl{PaperXI.Heisenberg.translationElement_mul}\\

99&theorem&\nolinkurl{PaperXI.Heisenberg.modulationElement_mul}\\

104&theorem&\nolinkurl{PaperXI.Heisenberg.centralElement_commute}\\

\bottomrule\end{longtable}

\section{\texorpdfstring{\nolinkurl{PaperXI.HeisenbergNormalForm}}{PaperXI.HeisenbergNormalForm}}\label{module:PaperXI.HeisenbergNormalForm}

\textbf{Selected layer:} v0.44 parent.\quad\textbf{Source lines:} 33.\par\textbf{SHA-256:} {\scriptsize\nolinkurl{61b84ce5c3a491de98e96f9f6edee89760c3d6516081c9372b8cc2d4cb98c1b2}}.

\VerbatimInput[fontsize=\scriptsize,breaklines=true,breakanywhere=true]{excerpts/D_HeisenbergNormalForm.txt}

\textbf{Direct project imports:} \nolinkurl{PaperXI.HeisenbergGroup}.

\begin{longtable}{@{}p{0.07\textwidth}p{0.12\textwidth}p{0.75\textwidth}@{}}\toprule Line&Kind&Name\\\midrule\endhead

10&theorem&\nolinkurl{PaperXI.Heisenberg.normal_form}\\

16&theorem&\nolinkurl{PaperXI.Heisenberg.modulation_translation}\\

25&theorem&\nolinkurl{PaperXI.Heisenberg.translation_modulation}\\

\bottomrule\end{longtable}

\section{\texorpdfstring{\nolinkurl{PaperXI.HeisenbergPrime}}{PaperXI.HeisenbergPrime}}\label{module:PaperXI.HeisenbergPrime}

\textbf{Selected layer:} v0.44 parent.\quad\textbf{Source lines:} 29.\par\textbf{SHA-256:} {\scriptsize\nolinkurl{5f481efb90ff5aa531e502311d9145cf59b1546638b05905f8dbce01b4be4dc9}}.

\VerbatimInput[fontsize=\scriptsize,breaklines=true,breakanywhere=true]{excerpts/D_HeisenbergPrime.txt}

\textbf{Direct project imports:} \nolinkurl{PaperXI.SchrodingerOperator}.

\begin{longtable}{@{}p{0.07\textwidth}p{0.12\textwidth}p{0.75\textwidth}@{}}\toprule Line&Kind&Name\\\midrule\endhead

12&theorem&\nolinkurl{PaperXI.HeisenbergPrime.two_ne_zero}\\

17&theorem&\nolinkurl{PaperXI.HeisenbergPrime.character_nontrivial}\\

24&theorem&\nolinkurl{PaperXI.HeisenbergPrime.character_formula}\\

\bottomrule\end{longtable}

\section{\texorpdfstring{\nolinkurl{PaperXI.HeisenbergProjector}}{PaperXI.HeisenbergProjector}}\label{module:PaperXI.HeisenbergProjector}

\textbf{Selected layer:} v0.44 parent.\quad\textbf{Source lines:} 230.\par\textbf{SHA-256:} {\scriptsize\nolinkurl{8d41365d8376b7d75e5d8bd5fed0875e6d650020b0341126c480b93f5bfe0d6a}}.

\VerbatimInput[fontsize=\scriptsize,breaklines=true,breakanywhere=true]{excerpts/D_HeisenbergProjector.txt}

\textbf{Direct project imports:} \nolinkurl{PaperXI.HeisenbergRepresentation}, \nolinkurl{PaperXI.HeisenbergForm}, \nolinkurl{PaperXI.HeisenbergGeometry}, \nolinkurl{PaperXI.SchrodingerUnitary}, \nolinkurl{PaperXI.FiniteFunctionAverage}.

\begin{longtable}{@{}p{0.07\textwidth}p{0.12\textwidth}p{0.75\textwidth}@{}}\toprule Line&Kind&Name\\\midrule\endhead

16&def&\nolinkurl{PaperXI.HeisenbergProjector.liftHom}\\

25&theorem&\nolinkurl{PaperXI.HeisenbergProjector.lift_position}\\

29&theorem&\nolinkurl{PaperXI.HeisenbergProjector.lift_frequency}\\

33&theorem&\nolinkurl{PaperXI.HeisenbergProjector.lift_central}\\

37&theorem&\nolinkurl{PaperXI.HeisenbergProjector.lift_injective}\\

45&def&\nolinkurl{PaperXI.HeisenbergProjector.liftedSubgroup}\\

48&theorem&\nolinkurl{PaperXI.HeisenbergProjector.mem_liftedSubgroup_iff}\\

58&def&\nolinkurl{PaperXI.HeisenbergProjector.liftEquiv}\\

67&instance&\nolinkurl{PaperXI.HeisenbergProjector.liftedFintype}\\

71&theorem&\nolinkurl{PaperXI.HeisenbergProjector.liftEquiv_coe}\\

75&def&\nolinkurl{PaperXI.HeisenbergProjector.restriction}\\

80&theorem&\nolinkurl{PaperXI.HeisenbergProjector.restriction_apply}\\

85&def&\nolinkurl{PaperXI.HeisenbergProjector.projector}\\

89&theorem&\nolinkurl{PaperXI.HeisenbergProjector.projector_apply}\\

101&def&\nolinkurl{PaperXI.HeisenbergProjector.subgroupProjector}\\

106&theorem&\nolinkurl{PaperXI.HeisenbergProjector.projector_eq_subgroupProjector}\\

121&theorem&\nolinkurl{PaperXI.HeisenbergProjector.idempotent}\\

127&theorem&\nolinkurl{PaperXI.HeisenbergProjector.range_eq_invariants}\\

136&theorem&\nolinkurl{PaperXI.HeisenbergProjector.pairing_invariant}\\

143&theorem&\nolinkurl{PaperXI.HeisenbergProjector.self_adjoint}\\

155&theorem&\nolinkurl{PaperXI.HeisenbergProjector.residual_orthogonal}\\

173&def&\nolinkurl{PaperXI.HeisenbergProjector.positionSpace}\\

179&def&\nolinkurl{PaperXI.HeisenbergProjector.positionEquiv}\\

187&instance&\nolinkurl{PaperXI.HeisenbergProjector.positionFintype}\\

190&theorem&\nolinkurl{PaperXI.HeisenbergProjector.position_isotropic}\\

196&theorem&\nolinkurl{PaperXI.HeisenbergProjector.position_projector_apply}\\

214&theorem&\nolinkurl{PaperXI.HeisenbergProjector.position_projector_eq}\\

220&theorem&\nolinkurl{PaperXI.HeisenbergProjector.position_range_eq}\\

225&theorem&\nolinkurl{PaperXI.HeisenbergProjector.position_rank_one}\\

\bottomrule\end{longtable}

\section{\texorpdfstring{\nolinkurl{PaperXI.HeisenbergReduction107}}{PaperXI.HeisenbergReduction107}}\label{module:PaperXI.HeisenbergReduction107}

\textbf{Selected layer:} base108 arithmetic branch.\quad\textbf{Source lines:} 61.\par\textbf{SHA-256:} {\scriptsize\nolinkurl{7bb3cef0f423eb861a688bbf662334632a9bfeaf015ddd1b429d554fc2c19934}}.

\VerbatimInput[fontsize=\scriptsize,breaklines=true,breakanywhere=true]{excerpts/D_HeisenbergReduction107.txt}

\textbf{Direct project imports:} \nolinkurl{PaperXI.TensorIntertwining107}.

\begin{longtable}{@{}p{0.07\textwidth}p{0.12\textwidth}p{0.75\textwidth}@{}}\toprule Line&Kind&Name\\\midrule\endhead

15&theorem&\nolinkurl{PaperXI.HeisenbergReduction107.pullback}\\

27&theorem&\nolinkurl{PaperXI.HeisenbergReduction107.double_coset_reduction}\\

43&theorem&\nolinkurl{PaperXI.HeisenbergReduction107.integral_reduction}\\

\bottomrule\end{longtable}

\section{\texorpdfstring{\nolinkurl{PaperXI.HeisenbergRepresentation}}{PaperXI.HeisenbergRepresentation}}\label{module:PaperXI.HeisenbergRepresentation}

\textbf{Selected layer:} v0.44 parent.\quad\textbf{Source lines:} 72.\par\textbf{SHA-256:} {\scriptsize\nolinkurl{386317cf254a374627df755dc02465986fc602151303b4e2b3413f2b82b46492}}.

\VerbatimInput[fontsize=\scriptsize,breaklines=true,breakanywhere=true]{excerpts/D_HeisenbergRepresentation.txt}

\textbf{Direct project imports:} \nolinkurl{PaperXI.HeisenbergGroup}, \nolinkurl{PaperXI.SchrodingerOperator}, \nolinkurl{PaperXI.SchrodingerIrreducible}.

\begin{longtable}{@{}p{0.07\textwidth}p{0.12\textwidth}p{0.75\textwidth}@{}}\toprule Line&Kind&Name\\\midrule\endhead

14&def&\nolinkurl{PaperXI.Heisenberg.representation}\\

21&theorem&\nolinkurl{PaperXI.Heisenberg.representation_apply}\\

28&def&\nolinkurl{PaperXI.Heisenberg.representationGL}\\

32&theorem&\nolinkurl{PaperXI.Heisenberg.representationGL_coe}\\

38&theorem&\nolinkurl{PaperXI.Heisenberg.representation_isUnit}\\

43&theorem&\nolinkurl{PaperXI.Heisenberg.central_action}\\

47&theorem&\nolinkurl{PaperXI.Heisenberg.translation_action}\\

51&theorem&\nolinkurl{PaperXI.Heisenberg.modulation_action}\\

58&theorem&\nolinkurl{PaperXI.Heisenberg.irreducible}\\

\bottomrule\end{longtable}

\section{\texorpdfstring{\nolinkurl{PaperXI.HeisenbergSelector}}{PaperXI.HeisenbergSelector}}\label{module:PaperXI.HeisenbergSelector}

\textbf{Selected layer:} v0.44 parent.\quad\textbf{Source lines:} 109.\par\textbf{SHA-256:} {\scriptsize\nolinkurl{89d25c89f133dcf4a56a48864b73092bcf942c81c09adf27fcd7b8f8149c6643}}.

\VerbatimInput[fontsize=\scriptsize,breaklines=true,breakanywhere=true]{excerpts/D_HeisenbergSelector.txt}

\textbf{Direct project imports:} \nolinkurl{PaperXI.FiniteAverage}, \nolinkurl{PaperXI.SchrodingerTrace}, \nolinkurl{PaperXI.HeisenbergCardinality}, \nolinkurl{PaperXI.HeisenbergProjector}.

\begin{longtable}{@{}p{0.07\textwidth}p{0.12\textwidth}p{0.75\textwidth}@{}}\toprule Line&Kind&Name\\\midrule\endhead

17&theorem&\nolinkurl{PaperXI.HeisenbergSelector.trace_average}\\

23&theorem&\nolinkurl{PaperXI.HeisenbergSelector.rank_average}\\

30&theorem&\nolinkurl{PaperXI.HeisenbergSelector.invariants_rank_one_of_trace_one}\\

42&theorem&\nolinkurl{PaperXI.HeisenbergSelector.coordinates_zero_iff}\\

57&theorem&\nolinkurl{PaperXI.HeisenbergSelector.trace_restriction}\\

66&theorem&\nolinkurl{PaperXI.HeisenbergSelector.trace_projector}\\

77&theorem&\nolinkurl{PaperXI.HeisenbergSelector.maximal_trace_one}\\

88&theorem&\nolinkurl{PaperXI.HeisenbergSelector.maximal_rank_one}\\

100&theorem&\nolinkurl{PaperXI.HeisenbergSelector.maximal_fixed_line}\\

\bottomrule\end{longtable}

\section{\texorpdfstring{\nolinkurl{PaperXI.HeisenbergSource}}{PaperXI.HeisenbergSource}}\label{module:PaperXI.HeisenbergSource}

\textbf{Selected layer:} v0.44 parent.\quad\textbf{Source lines:} 63.\par\textbf{SHA-256:} {\scriptsize\nolinkurl{ffb52975ff4f261a1791da759b032fd5c736abf36257053961a5b5b84c963686}}.

\VerbatimInput[fontsize=\scriptsize,breaklines=true,breakanywhere=true]{excerpts/D_HeisenbergSource.txt}

\textbf{Direct project imports:} \nolinkurl{PaperXI.HeisenbergPrime}, \nolinkurl{PaperXI.HeisenbergRepresentation}, \nolinkurl{PaperXI.HeisenbergDimension}, \nolinkurl{PaperXI.HeisenbergFinite}, \nolinkurl{PaperXI.HeisenbergCenter}, \nolinkurl{PaperXI.SchrodingerTrace}.

\begin{longtable}{@{}p{0.07\textwidth}p{0.12\textwidth}p{0.75\textwidth}@{}}\toprule Line&Kind&Name\\\midrule\endhead

19&def&\nolinkurl{PaperXI.HeisenbergSource.model}\\

22&theorem&\nolinkurl{PaperXI.HeisenbergSource.model_formula}\\

27&theorem&\nolinkurl{PaperXI.HeisenbergSource.center_character}\\

33&theorem&\nolinkurl{PaperXI.HeisenbergSource.model_invertible}\\

38&theorem&\nolinkurl{PaperXI.HeisenbergSource.model_irreducible}\\

44&theorem&\nolinkurl{PaperXI.HeisenbergSource.dimension_power}\\

48&theorem&\nolinkurl{PaperXI.HeisenbergSource.dimension_sqrt}\\

52&theorem&\nolinkurl{PaperXI.HeisenbergSource.character_support}\\

\bottomrule\end{longtable}

\section{\texorpdfstring{\nolinkurl{PaperXI.HeisenbergSourceProjector}}{PaperXI.HeisenbergSourceProjector}}\label{module:PaperXI.HeisenbergSourceProjector}

\textbf{Selected layer:} v0.44 parent.\quad\textbf{Source lines:} 118.\par\textbf{SHA-256:} {\scriptsize\nolinkurl{5e8ee76ad84809af2c219cb594f6ae86e0538557cd45997e05b8df888155396e}}.

\VerbatimInput[fontsize=\scriptsize,breaklines=true,breakanywhere=true]{excerpts/D_HeisenbergSourceProjector.txt}

\textbf{Direct project imports:} \nolinkurl{PaperXI.HeisenbergSource}, \nolinkurl{PaperXI.HeisenbergSelector}, \nolinkurl{PaperXI.HeisenbergProjector}.

\begin{longtable}{@{}p{0.07\textwidth}p{0.12\textwidth}p{0.75\textwidth}@{}}\toprule Line&Kind&Name\\\midrule\endhead

17&def&\nolinkurl{PaperXI.HeisenbergSourceProjector.projector}\\

24&theorem&\nolinkurl{PaperXI.HeisenbergSourceProjector.idempotent}\\

34&theorem&\nolinkurl{PaperXI.HeisenbergSourceProjector.range_eq_invariants}\\

46&theorem&\nolinkurl{PaperXI.HeisenbergSourceProjector.self_adjoint}\\

57&theorem&\nolinkurl{PaperXI.HeisenbergSourceProjector.residual_orthogonal}\\

68&theorem&\nolinkurl{PaperXI.HeisenbergSourceProjector.lift_central_iff_identity}\\

76&theorem&\nolinkurl{PaperXI.HeisenbergSourceProjector.position_range}\\

87&theorem&\nolinkurl{PaperXI.HeisenbergSourceProjector.position_rank_one}\\

97&theorem&\nolinkurl{PaperXI.HeisenbergSourceProjector.maximal_rank_one}\\

107&theorem&\nolinkurl{PaperXI.HeisenbergSourceProjector.maximal_fixed_line}\\

\bottomrule\end{longtable}

\section{\texorpdfstring{\nolinkurl{PaperXI.HeisenbergTensor}}{PaperXI.HeisenbergTensor}}\label{module:PaperXI.HeisenbergTensor}

\textbf{Selected layer:} v0.44 parent.\quad\textbf{Source lines:} 99.\par\textbf{SHA-256:} {\scriptsize\nolinkurl{c869c46b15622f6c072689647c34c28046d5f9ff18bbc8191683f0588727d4e7}}.

\VerbatimInput[fontsize=\scriptsize,breaklines=true,breakanywhere=true]{excerpts/D_HeisenbergTensor.txt}

\textbf{Direct project imports:} \nolinkurl{PaperXI.FiniteTensor}, \nolinkurl{PaperXI.CentralProduct}.

\begin{longtable}{@{}p{0.07\textwidth}p{0.12\textwidth}p{0.75\textwidth}@{}}\toprule Line&Kind&Name\\\midrule\endhead

17&def&\nolinkurl{PaperXI.HeisenbergTensor.ambient}\\

27&theorem&\nolinkurl{PaperXI.HeisenbergTensor.ambient_apply}\\

33&theorem&\nolinkurl{PaperXI.HeisenbergTensor.ambient_tprod}\\

41&theorem&\nolinkurl{PaperXI.HeisenbergTensor.ambient_balanced}\\

55&def&\nolinkurl{PaperXI.HeisenbergTensor.descend}\\

61&theorem&\nolinkurl{PaperXI.HeisenbergTensor.descend_mk}\\

68&theorem&\nolinkurl{PaperXI.HeisenbergTensor.descend_tprod}\\

77&theorem&\nolinkurl{PaperXI.HeisenbergTensor.central}\\

88&theorem&\nolinkurl{PaperXI.HeisenbergTensor.empty_central}\\

\bottomrule\end{longtable}

\section{\texorpdfstring{\nolinkurl{PaperXI.HeisenbergTensorComparison}}{PaperXI.HeisenbergTensorComparison}}\label{module:PaperXI.HeisenbergTensorComparison}

\textbf{Selected layer:} v0.44 parent.\quad\textbf{Source lines:} 70.\par\textbf{SHA-256:} {\scriptsize\nolinkurl{cc8aa01b6a54bea28296549299ad41b7873e08c01703c5aca621d700c4469661}}.

\VerbatimInput[fontsize=\scriptsize,breaklines=true,breakanywhere=true]{excerpts/D_HeisenbergTensorComparison.txt}

\textbf{Direct project imports:} \nolinkurl{PaperXI.HeisenbergTensorFactorization}, \nolinkurl{PaperXI.CentralExtensionComparison}.

\begin{longtable}{@{}p{0.07\textwidth}p{0.12\textwidth}p{0.75\textwidth}@{}}\toprule Line&Kind&Name\\\midrule\endhead

23&def&\nolinkurl{PaperXI.HeisenbergTensor.factorizationAlong}\\

33&theorem&\nolinkurl{PaperXI.HeisenbergTensor.factorizationAlong_intertwines}\\

43&def&\nolinkurl{PaperXI.HeisenbergTensor.onFullEnvelope}\\

48&def&\nolinkurl{PaperXI.HeisenbergTensor.fullEnvelopeFactorization}\\

60&theorem&\nolinkurl{PaperXI.HeisenbergTensor.fullEnvelopeFactorization_intertwines}\\

\bottomrule\end{longtable}

\section{\texorpdfstring{\nolinkurl{PaperXI.HeisenbergTensorFactorization}}{PaperXI.HeisenbergTensorFactorization}}\label{module:PaperXI.HeisenbergTensorFactorization}

\textbf{Selected layer:} v0.44 parent.\quad\textbf{Source lines:} 87.\par\textbf{SHA-256:} {\scriptsize\nolinkurl{03fec0d85721da724851da3e533aa76b616b5ba39ddfc3fe9a3ab5b513399b16}}.

\VerbatimInput[fontsize=\scriptsize,breaklines=true,breakanywhere=true]{excerpts/D_HeisenbergTensorFactorization.txt}

\textbf{Direct project imports:} \nolinkurl{PaperXI.HeisenbergTensor}, \nolinkurl{PaperXI.CentralProductCommutator}, \nolinkurl{PaperXI.TensorIrreducible}.

\begin{longtable}{@{}p{0.07\textwidth}p{0.12\textwidth}p{0.75\textwidth}@{}}\toprule Line&Kind&Name\\\midrule\endhead

27&theorem&\nolinkurl{PaperXI.HeisenbergTensor.factor_finiteDimensional}\\

31&theorem&\nolinkurl{PaperXI.HeisenbergTensor.finiteDimensional}\\

36&theorem&\nolinkurl{PaperXI.HeisenbergTensor.nontrivial}\\

42&theorem&\nolinkurl{PaperXI.HeisenbergTensor.finrank_sq}\\

51&theorem&\nolinkurl{PaperXI.HeisenbergTensor.irreducible}\\

61&theorem&\nolinkurl{PaperXI.HeisenbergTensor.exists_factorization}\\

71&def&\nolinkurl{PaperXI.HeisenbergTensor.factorization}\\

78&theorem&\nolinkurl{PaperXI.HeisenbergTensor.factorization_intertwines}\\

\bottomrule\end{longtable}

\section{\texorpdfstring{\nolinkurl{PaperXI.HeisenbergTrace108}}{PaperXI.HeisenbergTrace108}}\label{module:PaperXI.HeisenbergTrace108}

\textbf{Selected layer:} base108 arithmetic branch.\quad\textbf{Source lines:} 47.\par\textbf{SHA-256:} {\scriptsize\nolinkurl{045edb4a7a7672920cc0a4fdfc44e0b7b64d78563fe3cf3c625335b1e09240c8}}.

\VerbatimInput[fontsize=\scriptsize,breaklines=true,breakanywhere=true]{excerpts/D_HeisenbergTrace108.txt}

\textbf{Direct project imports:} none.

\begin{longtable}{@{}p{0.07\textwidth}p{0.12\textwidth}p{0.75\textwidth}@{}}\toprule Line&Kind&Name\\\midrule\endhead

16&theorem&\nolinkurl{PaperXI.HeisenbergTrace108.character_on_subgroup}\\

23&theorem&\nolinkurl{PaperXI.HeisenbergTrace108.character_zero_of_commutator}\\

40&theorem&\nolinkurl{PaperXI.HeisenbergTrace108.character_off_subgroup}\\

\bottomrule\end{longtable}

\section{\texorpdfstring{\nolinkurl{PaperXI.HeisenbergTransport}}{PaperXI.HeisenbergTransport}}\label{module:PaperXI.HeisenbergTransport}

\textbf{Selected layer:} v0.44 parent.\quad\textbf{Source lines:} 180.\par\textbf{SHA-256:} {\scriptsize\nolinkurl{7be38801f777c48277b8229987482c3b7e8c656c351a2e5201c7a067740572e2}}.

\VerbatimInput[fontsize=\scriptsize,breaklines=true,breakanywhere=true]{excerpts/D_HeisenbergTransport.txt}

\textbf{Direct project imports:} \nolinkurl{PaperXI.HeisenbergRepresentation}, \nolinkurl{PaperXI.HeisenbergForm}.

\begin{longtable}{@{}p{0.07\textwidth}p{0.12\textwidth}p{0.75\textwidth}@{}}\toprule Line&Kind&Name\\\midrule\endhead

20&def&\nolinkurl{PaperXI.HeisenbergTransport.vector}\\

22&theorem&\nolinkurl{PaperXI.HeisenbergTransport.vector_mul}\\

25&def&\nolinkurl{PaperXI.HeisenbergTransport.PreservesForm}\\

28&theorem&\nolinkurl{PaperXI.HeisenbergTransport.preservesForm_refl}\\

31&theorem&\nolinkurl{PaperXI.HeisenbergTransport.PreservesForm.symm}\\

37&theorem&\nolinkurl{PaperXI.HeisenbergTransport.PreservesForm.trans}\\

46&def&\nolinkurl{PaperXI.HeisenbergTransport.liftIsometry}\\

69&theorem&\nolinkurl{PaperXI.HeisenbergTransport.liftIsometry_vector}\\

74&theorem&\nolinkurl{PaperXI.HeisenbergTransport.liftIsometry_central}\\

79&theorem&\nolinkurl{PaperXI.HeisenbergTransport.liftIsometry_centralElement}\\

90&theorem&\nolinkurl{PaperXI.HeisenbergTransport.liftIsometry_refl}\\

95&theorem&\nolinkurl{PaperXI.HeisenbergTransport.liftIsometry_symm}\\

101&theorem&\nolinkurl{PaperXI.HeisenbergTransport.liftIsometry_trans}\\

115&def&\nolinkurl{PaperXI.HeisenbergTransport.pullback}\\

118&theorem&\nolinkurl{PaperXI.HeisenbergTransport.pullback_apply}\\

121&theorem&\nolinkurl{PaperXI.HeisenbergTransport.pullback_refl}\\

126&theorem&\nolinkurl{PaperXI.HeisenbergTransport.pullback_trans}\\

130&theorem&\nolinkurl{PaperXI.HeisenbergTransport.pullback_symm}\\

136&def&\nolinkurl{PaperXI.HeisenbergTransport.pullbackSubrepresentationEquiv}\\

145&theorem&\nolinkurl{PaperXI.HeisenbergTransport.pullback_irreducible_iff}\\

149&theorem&\nolinkurl{PaperXI.HeisenbergTransport.pullback_irreducible}\\

155&theorem&\nolinkurl{PaperXI.HeisenbergTransport.pullback_scalar_action}\\

170&theorem&\nolinkurl{PaperXI.HeisenbergTransport.pullback_central_action}\\

\bottomrule\end{longtable}

\section{\texorpdfstring{\nolinkurl{PaperXI.HenselianConductorTransport}}{PaperXI.HenselianConductorTransport}}\label{module:PaperXI.HenselianConductorTransport}

\textbf{Selected layer:} v0.44 parent.\quad\textbf{Source lines:} 120.\par\textbf{SHA-256:} {\scriptsize\nolinkurl{4195ca7a0238b26ebca438814c59c9712c09d21fd211f8468e3738fadebcb48a}}.

\VerbatimInput[fontsize=\scriptsize,breaklines=true,breakanywhere=true]{excerpts/D_HenselianConductorTransport.txt}

\textbf{Direct project imports:} \nolinkurl{PaperXI.HenselianWeilAttachment}.

\begin{longtable}{@{}p{0.07\textwidth}p{0.12\textwidth}p{0.75\textwidth}@{}}\toprule Line&Kind&Name\\\midrule\endhead

17&def&\nolinkurl{PaperXI.HenselianConductorTransport.localLowerEquiv}\\

31&def&\nolinkurl{PaperXI.HenselianConductorTransport.lowerEquiv}\\

39&theorem&\nolinkurl{PaperXI.HenselianConductorTransport.lower_card}\\

45&theorem&\nolinkurl{PaperXI.HenselianConductorTransport.embedding_lift}\\

58&theorem&\nolinkurl{PaperXI.HenselianConductorTransport.fixed_eq}\\

78&theorem&\nolinkurl{PaperXI.HenselianConductorTransport.summand_eq}\\

87&theorem&\nolinkurl{PaperXI.HenselianConductorTransport.swan_eq}\\

98&theorem&\nolinkurl{PaperXI.HenselianConductorTransport.artin_eq}\\

108&theorem&\nolinkurl{PaperXI.HenselianConductorTransport.summable_source}\\

\bottomrule\end{longtable}

\section{\texorpdfstring{\nolinkurl{PaperXI.HenselianFieldLift}}{PaperXI.HenselianFieldLift}}\label{module:PaperXI.HenselianFieldLift}

\textbf{Selected layer:} v0.44 parent.\quad\textbf{Source lines:} 95.\par\textbf{SHA-256:} {\scriptsize\nolinkurl{af945609f76e7f115f367f158515f7c5053cb2b0290dbc8baf21c20ea5f12e09}}.

\VerbatimInput[fontsize=\scriptsize,breaklines=true,breakanywhere=true]{excerpts/D_HenselianFieldLift.txt}

\textbf{Direct project imports:} \nolinkurl{PaperXI.HenselianResidueLift}.

\begin{longtable}{@{}p{0.07\textwidth}p{0.12\textwidth}p{0.75\textwidth}@{}}\toprule Line&Kind&Name\\\midrule\endhead

16&def&\nolinkurl{PaperXI.HenselianFieldLift.element}\\

19&def&\nolinkurl{PaperXI.HenselianFieldLift.field}\\

21&theorem&\nolinkurl{PaperXI.HenselianFieldLift.element_mem}\\

24&theorem&\nolinkurl{PaperXI.HenselianFieldLift.element_integral}\\

31&instance&\nolinkurl{PaperXI.HenselianFieldLift.finiteDimensional}\\

38&theorem&\nolinkurl{PaperXI.HenselianFieldLift.residue_element}\\

43&def&\nolinkurl{PaperXI.HenselianFieldLift.fixing}\\

47&theorem&\nolinkurl{PaperXI.HenselianFieldLift.mem_fixing_iff}\\

63&theorem&\nolinkurl{PaperXI.HenselianFieldLift.fixing_iff_residue}\\

81&theorem&\nolinkurl{PaperXI.HenselianFieldLift.inertia_le_fixing}\\

89&theorem&\nolinkurl{PaperXI.HenselianFieldLift.contained_in_inertia_fixed}\\

\bottomrule\end{longtable}

\section{\texorpdfstring{\nolinkurl{PaperXI.HenselianGaloisResidue}}{PaperXI.HenselianGaloisResidue}}\label{module:PaperXI.HenselianGaloisResidue}

\textbf{Selected layer:} v0.44 parent.\quad\textbf{Source lines:} 75.\par\textbf{SHA-256:} {\scriptsize\nolinkurl{4a35a55e8f7a231bdd6f1f2affbe1d44e00c6423dea5d53811260120032d96a0}}.

\VerbatimInput[fontsize=\scriptsize,breaklines=true,breakanywhere=true]{excerpts/D_HenselianGaloisResidue.txt}

\textbf{Direct project imports:} \nolinkurl{PaperXI.ValuationResidueEmbedding}.

\begin{longtable}{@{}p{0.07\textwidth}p{0.12\textwidth}p{0.75\textwidth}@{}}\toprule Line&Kind&Name\\\midrule\endhead

13&def&\nolinkurl{PaperXI.HenselianGaloisResidue.decompositionMap}\\

17&def&\nolinkurl{PaperXI.HenselianGaloisResidue.residueAction}\\

20&theorem&\nolinkurl{PaperXI.HenselianGaloisResidue.residueAction_residue}\\

23&theorem&\nolinkurl{PaperXI.HenselianGaloisResidue.fixes_base}\\

35&theorem&\nolinkurl{PaperXI.HenselianGaloisResidue.fixed_iff_base}\\

59&theorem&\nolinkurl{PaperXI.HenselianGaloisResidue.lift_automorphism}\\

\bottomrule\end{longtable}

\section{\texorpdfstring{\nolinkurl{PaperXI.HenselianLiftTower}}{PaperXI.HenselianLiftTower}}\label{module:PaperXI.HenselianLiftTower}

\textbf{Selected layer:} v0.44 parent.\quad\textbf{Source lines:} 71.\par\textbf{SHA-256:} {\scriptsize\nolinkurl{3ba24216ac193a2da00a39d6d731e8e750f4d8932c2ea761c52b6de0657f9cc7}}.

\VerbatimInput[fontsize=\scriptsize,breaklines=true,breakanywhere=true]{excerpts/D_HenselianLiftTower.txt}

\textbf{Direct project imports:} \nolinkurl{PaperXI.HenselianFieldLift}.

\begin{longtable}{@{}p{0.07\textwidth}p{0.12\textwidth}p{0.75\textwidth}@{}}\toprule Line&Kind&Name\\\midrule\endhead

12&theorem&\nolinkurl{PaperXI.HenselianLiftTower.henselian_of_complete}\\

24&theorem&\nolinkurl{PaperXI.HenselianLiftTower.root_power_tower}\\

32&theorem&\nolinkurl{PaperXI.HenselianLiftTower.lift_tower}\\

43&theorem&\nolinkurl{PaperXI.HenselianLiftTower.lift_equiv}\\

58&theorem&\nolinkurl{PaperXI.HenselianLiftTower.field_tower}\\

\bottomrule\end{longtable}

\section{\texorpdfstring{\nolinkurl{PaperXI.HenselianResidueLift}}{PaperXI.HenselianResidueLift}}\label{module:PaperXI.HenselianResidueLift}

\textbf{Selected layer:} v0.44 parent.\quad\textbf{Source lines:} 115.\par\textbf{SHA-256:} {\scriptsize\nolinkurl{52a3d30f9b5773234105b18f548446e88910a055cc08896c493da1c5bed39370}}.

\VerbatimInput[fontsize=\scriptsize,breaklines=true,breakanywhere=true]{excerpts/D_HenselianResidueLift.txt}

\textbf{Direct project imports:} none.

\begin{longtable}{@{}p{0.07\textwidth}p{0.12\textwidth}p{0.75\textwidth}@{}}\toprule Line&Kind&Name\\\midrule\endhead

15&def&\nolinkurl{PaperXI.HenselianResidueLift.poly}\\

17&theorem&\nolinkurl{PaperXI.HenselianResidueLift.poly_monic}\\

22&theorem&\nolinkurl{PaperXI.HenselianResidueLift.card_residue_zero}\\

26&theorem&\nolinkurl{PaperXI.HenselianResidueLift.derivative_unit}\\

31&theorem&\nolinkurl{PaperXI.HenselianResidueLift.exists_lift}\\

43&theorem&\nolinkurl{PaperXI.HenselianResidueLift.unique_lift}\\

53&def&\nolinkurl{PaperXI.HenselianResidueLift.lift}\\

55&theorem&\nolinkurl{PaperXI.HenselianResidueLift.lift_pow}\\

58&theorem&\nolinkurl{PaperXI.HenselianResidueLift.residue_lift}\\

61&theorem&\nolinkurl{PaperXI.HenselianResidueLift.lift_injective}\\

66&theorem&\nolinkurl{PaperXI.HenselianResidueLift.lift_zero}\\

71&theorem&\nolinkurl{PaperXI.HenselianResidueLift.lift_one}\\

76&theorem&\nolinkurl{PaperXI.HenselianResidueLift.lift_mul}\\

82&def&\nolinkurl{PaperXI.HenselianResidueLift.sectionMonoid}\\

88&theorem&\nolinkurl{PaperXI.HenselianResidueLift.lift_nat_pow}\\

91&theorem&\nolinkurl{PaperXI.HenselianResidueLift.lift_unit}\\

96&theorem&\nolinkurl{PaperXI.HenselianResidueLift.automorphism_fix_iff}\\

107&theorem&\nolinkurl{PaperXI.HenselianResidueLift.automorphism_transport}\\

\bottomrule\end{longtable}

\section{\texorpdfstring{\nolinkurl{PaperXI.HenselianWeilAttachment}}{PaperXI.HenselianWeilAttachment}}\label{module:PaperXI.HenselianWeilAttachment}

\textbf{Selected layer:} v0.44 parent.\quad\textbf{Source lines:} 140.\par\textbf{SHA-256:} {\scriptsize\nolinkurl{af3f6d521f61881ee03e54dcc204f5911df3f47dc1b0201fa77efd687608ba2d}}.

\VerbatimInput[fontsize=\scriptsize,breaklines=true,breakanywhere=true]{excerpts/D_HenselianWeilAttachment.txt}

\textbf{Direct project imports:} \nolinkurl{PaperXI.HenselianLiftTower}, \nolinkurl{PaperXI.ValuationFixingField}, \nolinkurl{PaperXI.UnramifiedValuation}.

\begin{longtable}{@{}p{0.07\textwidth}p{0.12\textwidth}p{0.75\textwidth}@{}}\toprule Line&Kind&Name\\\midrule\endhead

19&abbrev&\nolinkurl{PaperXI.HenselianWeilAttachment.liftedField}\\

23&theorem&\nolinkurl{PaperXI.HenselianWeilAttachment.field_fixing_eq}\\

30&def&\nolinkurl{PaperXI.HenselianWeilAttachment.decompositionEquiv}\\

37&theorem&\nolinkurl{PaperXI.HenselianWeilAttachment.ring_action}\\

42&theorem&\nolinkurl{PaperXI.HenselianWeilAttachment.residue_action}\\

50&theorem&\nolinkurl{PaperXI.HenselianWeilAttachment.residue_hom}\\

58&def&\nolinkurl{PaperXI.HenselianWeilAttachment.pullbackEquiv}\\

84&def&\nolinkurl{PaperXI.HenselianWeilAttachment.sourceEquiv}\\

92&def&\nolinkurl{PaperXI.HenselianWeilAttachment.embedding}\\

100&theorem&\nolinkurl{PaperXI.HenselianWeilAttachment.embedding_injective}\\

104&theorem&\nolinkurl{PaperXI.HenselianWeilAttachment.degree_formula}\\

111&theorem&\nolinkurl{PaperXI.HenselianWeilAttachment.embedding_range}\\

130&theorem&\nolinkurl{PaperXI.HenselianWeilAttachment.fixes_liftedField_iff}\\

\bottomrule\end{longtable}

\section{\texorpdfstring{\nolinkurl{PaperXI.HolonomyResidual}}{PaperXI.HolonomyResidual}}\label{module:PaperXI.HolonomyResidual}

\textbf{Selected layer:} v0.44 parent.\quad\textbf{Source lines:} 57.\par\textbf{SHA-256:} {\scriptsize\nolinkurl{d3e383de6f107b77fc33648b64a84751caa375bf3e468908ec513f79b1e11d52}}.

\VerbatimInput[fontsize=\scriptsize,breaklines=true,breakanywhere=true]{excerpts/D_HolonomyResidual.txt}

\textbf{Direct project imports:} none.

\begin{longtable}{@{}p{0.07\textwidth}p{0.12\textwidth}p{0.75\textwidth}@{}}\toprule Line&Kind&Name\\\midrule\endhead

14&def&\nolinkurl{PaperXI.HolonomyResidual.transported}\\

17&theorem&\nolinkurl{PaperXI.HolonomyResidual.transported_apply}\\

20&def&\nolinkurl{PaperXI.HolonomyResidual.cycleDefect}\\

27&def&\nolinkurl{PaperXI.HolonomyResidual.edgeResidual}\\

34&theorem&\nolinkurl{PaperXI.HolonomyResidual.residual_factorization}\\

43&theorem&\nolinkurl{PaperXI.HolonomyResidual.residual_zero_iff}\\

\bottomrule\end{longtable}

\section{\texorpdfstring{\nolinkurl{PaperXI.HyperbolicDeterminant}}{PaperXI.HyperbolicDeterminant}}\label{module:PaperXI.HyperbolicDeterminant}

\textbf{Selected layer:} v0.44 parent.\quad\textbf{Source lines:} 66.\par\textbf{SHA-256:} {\scriptsize\nolinkurl{36feef061c87d2e39afe6bfc3e8a246b942f70904fff07a77ab5dd5cff6fd275}}.

\VerbatimInput[fontsize=\scriptsize,breaklines=true,breakanywhere=true]{excerpts/D_HyperbolicDeterminant.txt}

\textbf{Direct project imports:} \nolinkurl{PaperXI.HyperbolicSpace}.

\begin{longtable}{@{}p{0.07\textwidth}p{0.12\textwidth}p{0.75\textwidth}@{}}\toprule Line&Kind&Name\\\midrule\endhead

13&theorem&\nolinkurl{PaperXI.HyperbolicDeterminant.inverse_difference}\\

21&theorem&\nolinkurl{PaperXI.HyperbolicDeterminant.dual_difference}\\

33&theorem&\nolinkurl{PaperXI.HyperbolicDeterminant.dual_difference_det}\\

43&theorem&\nolinkurl{PaperXI.HyperbolicDeterminant.action_difference}\\

51&theorem&\nolinkurl{PaperXI.HyperbolicDeterminant.total_difference_det}\\

59&theorem&\nolinkurl{PaperXI.HyperbolicDeterminant.action_det}\\

\bottomrule\end{longtable}

\section{\texorpdfstring{\nolinkurl{PaperXI.HyperbolicPhase}}{PaperXI.HyperbolicPhase}}\label{module:PaperXI.HyperbolicPhase}

\textbf{Selected layer:} v0.44 parent.\quad\textbf{Source lines:} 56.\par\textbf{SHA-256:} {\scriptsize\nolinkurl{419bca65c3efe41e14cf7793f6d4defd5c05e15201eb00e4161a916f3cb975f6}}.

\VerbatimInput[fontsize=\scriptsize,breaklines=true,breakanywhere=true]{excerpts/D_HyperbolicPhase.txt}

\textbf{Direct project imports:} \nolinkurl{PaperXI.HyperbolicDeterminant}, \nolinkurl{PaperXI.OddPhase}.

\begin{longtable}{@{}p{0.07\textwidth}p{0.12\textwidth}p{0.75\textwidth}@{}}\toprule Line&Kind&Name\\\midrule\endhead

16&theorem&\nolinkurl{PaperXI.HyperbolicPhase.fixedPointFree}\\

22&theorem&\nolinkurl{PaperXI.HyperbolicPhase.source}\\

47&theorem&\nolinkurl{PaperXI.HyperbolicPhase.symplectic_fixedPointFree}\\

\bottomrule\end{longtable}

\section{\texorpdfstring{\nolinkurl{PaperXI.HyperbolicSpace}}{PaperXI.HyperbolicSpace}}\label{module:PaperXI.HyperbolicSpace}

\textbf{Selected layer:} v0.44 parent.\quad\textbf{Source lines:} 120.\par\textbf{SHA-256:} {\scriptsize\nolinkurl{1532ea2542009051b6a09e99d9353d460de621615cf408b4dff48af6a072c501}}.

\VerbatimInput[fontsize=\scriptsize,breaklines=true,breakanywhere=true]{excerpts/D_HyperbolicSpace.txt}

\textbf{Direct project imports:} none.

\begin{longtable}{@{}p{0.07\textwidth}p{0.12\textwidth}p{0.75\textwidth}@{}}\toprule Line&Kind&Name\\\midrule\endhead

16&abbrev&\nolinkurl{PaperXI.HyperbolicSpace.Space}\\

19&def&\nolinkurl{PaperXI.HyperbolicSpace.form}\\

39&theorem&\nolinkurl{PaperXI.HyperbolicSpace.form_apply}\\

43&theorem&\nolinkurl{PaperXI.HyperbolicSpace.form_isAlt}\\

48&theorem&\nolinkurl{PaperXI.HyperbolicSpace.form_nondegenerate}\\

72&def&\nolinkurl{PaperXI.HyperbolicSpace.action}\\

75&theorem&\nolinkurl{PaperXI.HyperbolicSpace.action_apply}\\

79&theorem&\nolinkurl{PaperXI.HyperbolicSpace.action_form}\\

85&theorem&\nolinkurl{PaperXI.HyperbolicSpace.dual_fixedPointFree}\\

107&theorem&\nolinkurl{PaperXI.HyperbolicSpace.action_fixedPointFree}\\

115&theorem&\nolinkurl{PaperXI.HyperbolicSpace.finrank}\\

\bottomrule\end{longtable}

\section{\texorpdfstring{\nolinkurl{PaperXI.InducedClassification}}{PaperXI.InducedClassification}}\label{module:PaperXI.InducedClassification}

\textbf{Selected layer:} v0.44 parent.\quad\textbf{Source lines:} 136.\par\textbf{SHA-256:} {\scriptsize\nolinkurl{5700629a23c61b9f8d38d3ca69ba94c5efdaf0c847404c0c162aa3c7e54321a0}}.

\VerbatimInput[fontsize=\scriptsize,breaklines=true,breakanywhere=true]{excerpts/D_InducedClassification.txt}

\textbf{Direct project imports:} \nolinkurl{PaperXI.InducedDual}, \nolinkurl{PaperXI.CyclicInducedWrap}.

\begin{longtable}{@{}p{0.07\textwidth}p{0.12\textwidth}p{0.75\textwidth}@{}}\toprule Line&Kind&Name\\\midrule\endhead

16&theorem&\nolinkurl{PaperXI.InducedClassification.index_transitive}\\

22&theorem&\nolinkurl{PaperXI.InducedClassification.index_stabilizer}\\

29&def&\nolinkurl{PaperXI.InducedClassification.conjugateEquiv}\\

38&theorem&\nolinkurl{PaperXI.InducedClassification.coefficient_action}\\

56&theorem&\nolinkurl{PaperXI.InducedClassification.character_of_equiv}\\

84&theorem&\nolinkurl{PaperXI.InducedClassification.equiv_iff}\\

93&theorem&\nolinkurl{PaperXI.InducedClassification.selfDual_iff}\\

110&def&\nolinkurl{PaperXI.InducedClassification.degreeDualEquiv}\\

120&theorem&\nolinkurl{PaperXI.InducedClassification.degreeSelfDual_iff}\\

\bottomrule\end{longtable}

\section{\texorpdfstring{\nolinkurl{PaperXI.InducedCycleInvariants}}{PaperXI.InducedCycleInvariants}}\label{module:PaperXI.InducedCycleInvariants}

\textbf{Selected layer:} v0.44 parent.\quad\textbf{Source lines:} 107.\par\textbf{SHA-256:} {\scriptsize\nolinkurl{2faaeb89211dca9d88df443a27ad69ab7e310783d4e07a3ad54d85c3c5bcae1a}}.

\VerbatimInput[fontsize=\scriptsize,breaklines=true,breakanywhere=true]{excerpts/D_InducedCycleInvariants.txt}

\textbf{Direct project imports:} \nolinkurl{PaperXI.CyclicInducedWrap}, \nolinkurl{PaperXI.CycleCoordinates}.

\begin{longtable}{@{}p{0.07\textwidth}p{0.12\textwidth}p{0.75\textwidth}@{}}\toprule Line&Kind&Name\\\midrule\endhead

18&def&\nolinkurl{PaperXI.InducedCycleInvariants.indexEquiv}\\

25&def&\nolinkurl{PaperXI.InducedCycleInvariants.basis}\\

28&theorem&\nolinkurl{PaperXI.InducedCycleInvariants.basis_apply}\\

33&def&\nolinkurl{PaperXI.InducedCycleInvariants.returnScalar}\\

37&theorem&\nolinkurl{PaperXI.InducedCycleInvariants.dimension}\\

41&theorem&\nolinkurl{PaperXI.InducedCycleInvariants.step}\\

51&theorem&\nolinkurl{PaperXI.InducedCycleInvariants.closing_step}\\

61&theorem&\nolinkurl{PaperXI.InducedCycleInvariants.basis_powers}\\

69&theorem&\nolinkurl{PaperXI.InducedCycleInvariants.finiteSpace}\\

75&theorem&\nolinkurl{PaperXI.InducedCycleInvariants.cycle_invariants}\\

83&theorem&\nolinkurl{PaperXI.InducedCycleInvariants.returnScalar_from_charpoly}\\

89&theorem&\nolinkurl{PaperXI.InducedCycleInvariants.ambient_polynomial}\\

99&theorem&\nolinkurl{PaperXI.InducedCycleInvariants.return_operator}\\

\bottomrule\end{longtable}

\section{\texorpdfstring{\nolinkurl{PaperXI.InducedDual}}{PaperXI.InducedDual}}\label{module:PaperXI.InducedDual}

\textbf{Selected layer:} v0.44 parent.\quad\textbf{Source lines:} 42.\par\textbf{SHA-256:} {\scriptsize\nolinkurl{e915fa1d7389f923ae824272ca73b059424964f89c7e46b0aa503af596137995}}.

\VerbatimInput[fontsize=\scriptsize,breaklines=true,breakanywhere=true]{excerpts/D_InducedDual.txt}

\textbf{Direct project imports:} \nolinkurl{PaperXI.MonomialDual}, \nolinkurl{PaperXI.CyclicInducedCoordinates}, \nolinkurl{PaperXI.MonomialInductionSubgroup}.

\begin{longtable}{@{}p{0.07\textwidth}p{0.12\textwidth}p{0.75\textwidth}@{}}\toprule Line&Kind&Name\\\midrule\endhead

15&def&\nolinkurl{PaperXI.InducedDual.indexAction}\\

21&theorem&\nolinkurl{PaperXI.InducedDual.indexAction_apply}\\

26&def&\nolinkurl{PaperXI.InducedDual.inverseCharacterEquiv}\\

\bottomrule\end{longtable}

\section{\texorpdfstring{\nolinkurl{PaperXI.InducedTameProfile}}{PaperXI.InducedTameProfile}}\label{module:PaperXI.InducedTameProfile}

\textbf{Selected layer:} v0.44 parent.\quad\textbf{Source lines:} 187.\par\textbf{SHA-256:} {\scriptsize\nolinkurl{76a61c7e77722f469f734ee21b80e55786a2aec4a1916c23c4dc48a92249d576}}.

\VerbatimInput[fontsize=\scriptsize,breaklines=true,breakanywhere=true]{excerpts/D_InducedTameProfile.txt}

\textbf{Direct project imports:} \nolinkurl{PaperXI.InducedCycleInvariants}, \nolinkurl{PaperXI.TameConductor}, \nolinkurl{PaperXI.InvariantPolynomial}.

\begin{longtable}{@{}p{0.07\textwidth}p{0.12\textwidth}p{0.75\textwidth}@{}}\toprule Line&Kind&Name\\\midrule\endhead

18&theorem&\nolinkurl{PaperXI.InducedTameProfile.inertia_degree}\\

33&def&\nolinkurl{PaperXI.InducedTameProfile.inertiaInDegree}\\

51&theorem&\nolinkurl{PaperXI.InducedTameProfile.origin_mem}\\

63&theorem&\nolinkurl{PaperXI.InducedTameProfile.basis_packet}\\

77&theorem&\nolinkurl{PaperXI.InducedTameProfile.support_of_basis}\\

104&theorem&\nolinkurl{PaperXI.InducedTameProfile.inertia_support}\\

112&theorem&\nolinkurl{PaperXI.InducedTameProfile.inertia_invariants}\\

128&theorem&\nolinkurl{PaperXI.InducedTameProfile.invariant_dimension}\\

149&theorem&\nolinkurl{PaperXI.InducedTameProfile.local_polynomial}\\

173&theorem&\nolinkurl{PaperXI.InducedTameProfile.finite_conductors}\\

\bottomrule\end{longtable}

\section{\texorpdfstring{\nolinkurl{PaperXI.InducedTwist}}{PaperXI.InducedTwist}}\label{module:PaperXI.InducedTwist}

\textbf{Selected layer:} v0.44 parent.\quad\textbf{Source lines:} 103.\par\textbf{SHA-256:} {\scriptsize\nolinkurl{53e7034ad1f3dacf431f99ae0d0334e94296550244011390bb7f137f09a4ebec}}.

\VerbatimInput[fontsize=\scriptsize,breaklines=true,breakanywhere=true]{excerpts/D_InducedTwist.txt}

\textbf{Direct project imports:} \nolinkurl{PaperXI.CharacterTwist}, \nolinkurl{PaperXI.CyclicInducedWrap}.

\begin{longtable}{@{}p{0.07\textwidth}p{0.12\textwidth}p{0.75\textwidth}@{}}\toprule Line&Kind&Name\\\midrule\endhead

14&def&\nolinkurl{PaperXI.InducedTwist.character}\\

16&theorem&\nolinkurl{PaperXI.InducedTwist.character_apply}\\

20&def&\nolinkurl{PaperXI.InducedTwist.forward}\\

31&def&\nolinkurl{PaperXI.InducedTwist.backward}\\

43&def&\nolinkurl{PaperXI.InducedTwist.coindLinearEquiv}\\

52&def&\nolinkurl{PaperXI.InducedTwist.coindEquiv}\\

62&def&\nolinkurl{PaperXI.InducedTwist.inducedEquiv}\\

74&def&\nolinkurl{PaperXI.InducedTwist.degreeEquiv}\\

81&theorem&\nolinkurl{PaperXI.InducedTwist.period_value}\\

88&theorem&\nolinkurl{PaperXI.InducedTwist.closing_basis}\\

97&theorem&\nolinkurl{PaperXI.InducedTwist.inertia_unchanged}\\

\bottomrule\end{longtable}

\section{\texorpdfstring{\nolinkurl{PaperXI.InducedTypeSpace68}}{PaperXI.InducedTypeSpace68}}\label{module:PaperXI.InducedTypeSpace68}

\textbf{Selected layer:} base108 arithmetic branch.\quad\textbf{Source lines:} 142.\par\textbf{SHA-256:} {\scriptsize\nolinkurl{fe12e8899453c20460a4e73b84995d88a3551bb34ad86ab9b31fc3b83a89b409}}.

\VerbatimInput[fontsize=\scriptsize,breaklines=true,breakanywhere=true]{excerpts/D_InducedTypeSpace68.txt}

\textbf{Direct project imports:} \nolinkurl{PaperXI.InductionIrreducibility}.

\begin{longtable}{@{}p{0.07\textwidth}p{0.12\textwidth}p{0.75\textwidth}@{}}\toprule Line&Kind&Name\\\midrule\endhead

21&theorem&\nolinkurl{PaperXI.InducedTypeSpace68.factorization}\\

35&theorem&\nolinkurl{PaperXI.InducedTypeSpace68.factorization_apply}\\

45&def&\nolinkurl{PaperXI.InducedTypeSpace68.homEquiv}\\

60&theorem&\nolinkurl{PaperXI.InducedTypeSpace68.range_eq_identity}\\

85&def&\nolinkurl{PaperXI.InducedTypeSpace68.typeSpace}\\

89&def&\nolinkurl{PaperXI.InducedTypeSpace68.typeEquiv}\\

96&def&\nolinkurl{PaperXI.InducedTypeSpace68.typeOperators}\\

99&theorem&\nolinkurl{PaperXI.InducedTypeSpace68.typeOperators_mul}\\

103&theorem&\nolinkurl{PaperXI.InducedTypeSpace68.typeOperators_add}\\

107&theorem&\nolinkurl{PaperXI.InducedTypeSpace68.typeOperators_eq_zero_iff}\\

112&theorem&\nolinkurl{PaperXI.InducedTypeSpace68.typeOperators_idempotent_iff}\\

119&theorem&\nolinkurl{PaperXI.InducedTypeSpace68.type_map_scalar}\\

134&theorem&\nolinkurl{PaperXI.InducedTypeSpace68.conjugation_eq_of_scalar}\\

\bottomrule\end{longtable}

\section{\texorpdfstring{\nolinkurl{PaperXI.InductionIrreducibility}}{PaperXI.InductionIrreducibility}}\label{module:PaperXI.InductionIrreducibility}

\textbf{Selected layer:} base108 arithmetic branch.\quad\textbf{Source lines:} 237.\par\textbf{SHA-256:} {\scriptsize\nolinkurl{b099fd6ad458401e85a83ae0289bec6be2de9575eb7b73c496fafb98d13088ed}}.

\VerbatimInput[fontsize=\scriptsize,breaklines=true,breakanywhere=true]{excerpts/D_InductionIrreducibility.txt}

\textbf{Direct project imports:} \nolinkurl{PaperXI.OpenSubgroupSmoothInduction}, \nolinkurl{PaperXI.RepresentationStability}, \nolinkurl{PaperXI.SemisimpleProjection}, \nolinkurl{PaperXI.CompactSmoothSemisimple}.

\begin{longtable}{@{}p{0.07\textwidth}p{0.12\textwidth}p{0.75\textwidth}@{}}\toprule Line&Kind&Name\\\midrule\endhead

26&def&\nolinkurl{PaperXI.InductionIrreducibility.identityInclusion}\\

36&theorem&\nolinkurl{PaperXI.InductionIrreducibility.evaluate_translate}\\

42&def&\nolinkurl{PaperXI.InductionIrreducibility.identityEvaluation}\\

53&theorem&\nolinkurl{PaperXI.InductionIrreducibility.identityEvaluation_inclusion}\\

57&def&\nolinkurl{PaperXI.InductionIrreducibility.identityProjection}\\

61&theorem&\nolinkurl{PaperXI.InductionIrreducibility.exists_nonzero_coefficient}\\

72&theorem&\nolinkurl{PaperXI.InductionIrreducibility.eq_top_of_contains_identity}\\

88&theorem&\nolinkurl{PaperXI.InductionIrreducibility.irreducible_of_projection}\\

126&theorem&\nolinkurl{PaperXI.InductionIrreducibility.algebra_action_mem}\\

137&theorem&\nolinkurl{PaperXI.InductionIrreducibility.irreducible_of_local_groupAlgebra}\\

150&def&\nolinkurl{PaperXI.InductionIrreducibility.intersectionSubgroup}\\

154&def&\nolinkurl{PaperXI.InductionIrreducibility.intersectionConjugation}\\

161&def&\nolinkurl{PaperXI.InductionIrreducibility.coefficientIntertwiner}\\

182&theorem&\nolinkurl{PaperXI.InductionIrreducibility.intertwiner_eq_zero_of_mackey}\\

205&theorem&\nolinkurl{PaperXI.InductionIrreducibility.irreducible_of_mackey}\\

221&theorem&\nolinkurl{PaperXI.InductionIrreducibility.irreducible_of_compact_open_mackey}\\

\bottomrule\end{longtable}

\section{\texorpdfstring{\nolinkurl{PaperXI.InertiaFixed}}{PaperXI.InertiaFixed}}\label{module:PaperXI.InertiaFixed}

\textbf{Selected layer:} v0.44 parent.\quad\textbf{Source lines:} 90.\par\textbf{SHA-256:} {\scriptsize\nolinkurl{ac85a017c6c2462073de65b7304984cdc44ecebf7b201d52cad8f4a5e916e967}}.

\VerbatimInput[fontsize=\scriptsize,breaklines=true,breakanywhere=true]{excerpts/D_InertiaFixed.txt}

\textbf{Direct project imports:} \nolinkurl{PaperXI.PacketReturn}.

\begin{longtable}{@{}p{0.07\textwidth}p{0.12\textwidth}p{0.75\textwidth}@{}}\toprule Line&Kind&Name\\\midrule\endhead

13&theorem&\nolinkurl{PaperXI.InertiaFixed.invariants_eq_packet_zero}\\

23&theorem&\nolinkurl{PaperXI.InertiaFixed.orbit_label_eq_zero_iff}\\

32&theorem&\nolinkurl{PaperXI.InertiaFixed.invariants_eq_top_of_basis}\\

40&theorem&\nolinkurl{PaperXI.InertiaFixed.invariants_eq_bot_of_basis}\\

56&theorem&\nolinkurl{PaperXI.InertiaFixed.fixed_forward}\\

63&theorem&\nolinkurl{PaperXI.InertiaFixed.fixed_backward}\\

70&theorem&\nolinkurl{PaperXI.InertiaFixed.invariants_of_cycle_basis}\\

\bottomrule\end{longtable}

\section{\texorpdfstring{\nolinkurl{PaperXI.InertiaQuotient}}{PaperXI.InertiaQuotient}}\label{module:PaperXI.InertiaQuotient}

\textbf{Selected layer:} v0.44 parent.\quad\textbf{Source lines:} 98.\par\textbf{SHA-256:} {\scriptsize\nolinkurl{6864c2943a2417a4883fae946d64fcb26273f3e590700e9ae16b296e8d5361bf}}.

\VerbatimInput[fontsize=\scriptsize,breaklines=true,breakanywhere=true]{excerpts/D_InertiaQuotient.txt}

\textbf{Direct project imports:} \nolinkurl{PaperXI.InteractionPolynomial}.

\begin{longtable}{@{}p{0.07\textwidth}p{0.12\textwidth}p{0.75\textwidth}@{}}\toprule Line&Kind&Name\\\midrule\endhead

13&theorem&\nolinkurl{PaperXI.InertiaQuotient.invariants_comp_of_surjective}\\

23&def&\nolinkurl{PaperXI.InertiaQuotient.invariantEquiv}\\

27&theorem&\nolinkurl{PaperXI.InertiaQuotient.invariantEquiv_apply}\\

34&def&\nolinkurl{PaperXI.InertiaQuotient.factoredFrobenius}\\

42&theorem&\nolinkurl{PaperXI.InertiaQuotient.factoredFrobenius_apply}\\

49&theorem&\nolinkurl{PaperXI.InertiaQuotient.factored_polynomial_eq}\\

58&theorem&\nolinkurl{PaperXI.InertiaQuotient.factored_local_polynomial}\\

70&theorem&\nolinkurl{PaperXI.InertiaQuotient.factored_character_trivial_iff}\\

83&theorem&\nolinkurl{PaperXI.InertiaQuotient.factored_determinant}\\

\bottomrule\end{longtable}

\section{\texorpdfstring{\nolinkurl{PaperXI.IntegerRingComparison}}{PaperXI.IntegerRingComparison}}\label{module:PaperXI.IntegerRingComparison}

\textbf{Selected layer:} v0.44 parent.\quad\textbf{Source lines:} 230.\par\textbf{SHA-256:} {\scriptsize\nolinkurl{382149600512912093535bbc83e0f1dcd500caa3db84e5340ac368157c74b06d}}.

\VerbatimInput[fontsize=\scriptsize,breaklines=true,breakanywhere=true]{excerpts/D_IntegerRingComparison.txt}

\textbf{Direct project imports:} \nolinkurl{PaperXI.SpectralLocalField}.

\begin{longtable}{@{}p{0.07\textwidth}p{0.12\textwidth}p{0.75\textwidth}@{}}\toprule Line&Kind&Name\\\midrule\endhead

45&abbrev&\nolinkurl{PaperXI.IntegerRingComparison.UpperRing}\\

46&abbrev&\nolinkurl{PaperXI.IntegerRingComparison.UpperField}\\

50&theorem&\nolinkurl{PaperXI.IntegerRingComparison.norm_le_one_iff_isIntegral}\\

95&theorem&\nolinkurl{PaperXI.IntegerRingComparison.isIntegral_iff_exists_upper_lift}\\

108&theorem&\nolinkurl{PaperXI.IntegerRingComparison.upper_integerRing_eq_spectral}\\

127&def&\nolinkurl{PaperXI.IntegerRingComparison.upperRingEquivSpectralIntegers}\\

137&instance&\nolinkurl{PaperXI.IntegerRingComparison.spectralIntegerAlgebra}\\

143&instance&\nolinkurl{PaperXI.IntegerRingComparison.spectralIntegerAlgebra_isLocal}\\

155&def&\nolinkurl{PaperXI.IntegerRingComparison.upperRingAlgEquivSpectralIntegers}\\

159&instance&\nolinkurl{PaperXI.IntegerRingComparison.spectralIntegersEtale}\\

166&theorem&\nolinkurl{PaperXI.IntegerRingComparison.upperRingEquivSpectralIntegers_coe}\\

173&theorem&\nolinkurl{PaperXI.IntegerRingComparison.upper_maximalIdeal_map}\\

182&theorem&\nolinkurl{PaperXI.IntegerRingComparison.spectral_ramificationIdx_one}\\

188&def&\nolinkurl{PaperXI.IntegerRingComparison.upperResidueFieldEquivSpectral}\\

193&def&\nolinkurl{PaperXI.IntegerRingComparison.upperResidueFieldAlgEquivSpectral}\\

200&def&\nolinkurl{PaperXI.IntegerRingComparison.spectralResidueAlgEquiv}\\

207&def&\nolinkurl{PaperXI.IntegerRingComparison.spectralResidueEquiv}\\

214&theorem&\nolinkurl{PaperXI.IntegerRingComparison.spectral_residue_finrank}\\

220&theorem&\nolinkurl{PaperXI.IntegerRingComparison.upper_degree_eq_spectral_residue_finrank}\\

226&theorem&\nolinkurl{PaperXI.IntegerRingComparison.natCard_spectralResidue}\\

\bottomrule\end{longtable}

\section{\texorpdfstring{\nolinkurl{PaperXI.InteractionPolynomial}}{PaperXI.InteractionPolynomial}}\label{module:PaperXI.InteractionPolynomial}

\textbf{Selected layer:} v0.44 parent.\quad\textbf{Source lines:} 77.\par\textbf{SHA-256:} {\scriptsize\nolinkurl{e9129207250ae783e05f622a71a16a8654e22680c10c252bdb028abccec69ec2}}.

\VerbatimInput[fontsize=\scriptsize,breaklines=true,breakanywhere=true]{excerpts/D_InteractionPolynomial.txt}

\textbf{Direct project imports:} \nolinkurl{PaperXI.InvariantPolynomial}, \nolinkurl{PaperXI.PacketCycle}.

\begin{longtable}{@{}p{0.07\textwidth}p{0.12\textwidth}p{0.75\textwidth}@{}}\toprule Line&Kind&Name\\\midrule\endhead

14&theorem&\nolinkurl{PaperXI.InteractionPolynomial.character_trivial_iff}\\

24&theorem&\nolinkurl{PaperXI.InteractionPolynomial.local_polynomial}\\

47&theorem&\nolinkurl{PaperXI.InteractionPolynomial.local_polynomial_character_form}\\

62&theorem&\nolinkurl{PaperXI.InteractionPolynomial.group_local_polynomial}\\

\bottomrule\end{longtable}

\section{\texorpdfstring{\nolinkurl{PaperXI.IntersectionModel108}}{PaperXI.IntersectionModel108}}\label{module:PaperXI.IntersectionModel108}

\textbf{Selected layer:} base108 arithmetic branch.\quad\textbf{Source lines:} 62.\par\textbf{SHA-256:} {\scriptsize\nolinkurl{8f2f518c8309162145d80c2ae8d403d2c1eed20ac545a995985c18af108d6343}}.

\VerbatimInput[fontsize=\scriptsize,breaklines=true,breakanywhere=true]{excerpts/D_IntersectionModel108.txt}

\textbf{Direct project imports:} \nolinkurl{PaperXI.IntertwiningSpace107}, \nolinkurl{PaperXI.QuotientRepresentation}, \nolinkurl{PaperXI.CharacterPairing108}.

\begin{longtable}{@{}p{0.07\textwidth}p{0.12\textwidth}p{0.75\textwidth}@{}}\toprule Line&Kind&Name\\\midrule\endhead

16&def&\nolinkurl{PaperXI.IntersectionModel108.pairs}\\

20&def&\nolinkurl{PaperXI.IntersectionModel108.left}\\

23&def&\nolinkurl{PaperXI.IntersectionModel108.right}\\

27&def&\nolinkurl{PaperXI.IntersectionModel108.spaceEquiv}\\

39&def&\nolinkurl{PaperXI.IntersectionModel108.pullbackEquiv}\\

53&theorem&\nolinkurl{PaperXI.IntersectionModel108.finrank_model}\\

\bottomrule\end{longtable}

\section{\texorpdfstring{\nolinkurl{PaperXI.IntertwiningAlong105}}{PaperXI.IntertwiningAlong105}}\label{module:PaperXI.IntertwiningAlong105}

\textbf{Selected layer:} base108 arithmetic branch.\quad\textbf{Source lines:} 78.\par\textbf{SHA-256:} {\scriptsize\nolinkurl{e1af42743d7ea4d6cbdbb96cb3d355e1f6502d69367e55041090370ccbea5478}}.

\VerbatimInput[fontsize=\scriptsize,breaklines=true,breakanywhere=true]{excerpts/D_IntertwiningAlong105.txt}

\textbf{Direct project imports:} none.

\begin{longtable}{@{}p{0.07\textwidth}p{0.12\textwidth}p{0.75\textwidth}@{}}\toprule Line&Kind&Name\\\midrule\endhead

13&theorem&\nolinkurl{PaperXI.IntertwiningAlong105.act_mul}\\

16&theorem&\nolinkurl{PaperXI.IntertwiningAlong105.act_inv_cancel}\\

20&theorem&\nolinkurl{PaperXI.IntertwiningAlong105.act_mul_inv_cancel}\\

25&def&\nolinkurl{PaperXI.IntertwiningAlong105.Along}\\

29&def&\nolinkurl{PaperXI.IntertwiningAlong105.ExistsAt}\\

32&theorem&\nolinkurl{PaperXI.IntertwiningAlong105.left_right}\\

46&theorem&\nolinkurl{PaperXI.IntertwiningAlong105.left_right_ne_zero}\\

56&theorem&\nolinkurl{PaperXI.IntertwiningAlong105.existsAt_of_left_right}\\

62&theorem&\nolinkurl{PaperXI.IntertwiningAlong105.remove_central}\\

72&theorem&\nolinkurl{PaperXI.IntertwiningAlong105.existsAt_remove_central}\\

\bottomrule\end{longtable}

\section{\texorpdfstring{\nolinkurl{PaperXI.IntertwiningSpace107}}{PaperXI.IntertwiningSpace107}}\label{module:PaperXI.IntertwiningSpace107}

\textbf{Selected layer:} base108 arithmetic branch.\quad\textbf{Source lines:} 81.\par\textbf{SHA-256:} {\scriptsize\nolinkurl{2c3ea887cc1868b2c539bf9d466778b39567a4ea10aa245876ba0dfe9bef8aa7}}.

\VerbatimInput[fontsize=\scriptsize,breaklines=true,breakanywhere=true]{excerpts/D_IntertwiningSpace107.txt}

\textbf{Direct project imports:} \nolinkurl{PaperXI.TensorIntertwining107}.

\begin{longtable}{@{}p{0.07\textwidth}p{0.12\textwidth}p{0.75\textwidth}@{}}\toprule Line&Kind&Name\\\midrule\endhead

16&def&\nolinkurl{PaperXI.IntertwiningSpace107.space}\\

33&def&\nolinkurl{PaperXI.IntertwiningSpace107.tensorMap}\\

43&def&\nolinkurl{PaperXI.IntertwiningSpace107.equivalence}\\

54&theorem&\nolinkurl{PaperXI.IntertwiningSpace107.finrank_eq}\\

61&theorem&\nolinkurl{PaperXI.IntertwiningSpace107.line_of_dimension}\\

\bottomrule\end{longtable}

\section{\texorpdfstring{\nolinkurl{PaperXI.InvariantAssembly}}{PaperXI.InvariantAssembly}}\label{module:PaperXI.InvariantAssembly}

\textbf{Selected layer:} v0.44 parent.\quad\textbf{Source lines:} 48.\par\textbf{SHA-256:} {\scriptsize\nolinkurl{2fb39cacd47c40f2a53037ffb5b416bb5f6311285265e3f57aed422baa40e5b4}}.

\VerbatimInput[fontsize=\scriptsize,breaklines=true,breakanywhere=true]{excerpts/D_InvariantAssembly.txt}

\textbf{Direct project imports:} \nolinkurl{PaperXI.GroupFixed}.

\begin{longtable}{@{}p{0.07\textwidth}p{0.12\textwidth}p{0.75\textwidth}@{}}\toprule Line&Kind&Name\\\midrule\endhead

13&def&\nolinkurl{PaperXI.InvariantAssembly.invariantsEquiv}\\

28&theorem&\nolinkurl{PaperXI.InvariantAssembly.invariantsEquiv_apply}\\

32&theorem&\nolinkurl{PaperXI.InvariantAssembly.finrank_invariants}\\

39&theorem&\nolinkurl{PaperXI.InvariantAssembly.count_invariants}\\

\bottomrule\end{longtable}

\section{\texorpdfstring{\nolinkurl{PaperXI.InvariantPolynomial}}{PaperXI.InvariantPolynomial}}\label{module:PaperXI.InvariantPolynomial}

\textbf{Selected layer:} v0.44 parent.\quad\textbf{Source lines:} 90.\par\textbf{SHA-256:} {\scriptsize\nolinkurl{cd9db934f06df326729d69a24446a44a9b6d00c9e5dbf48c6ff8c666e2b36990}}.

\VerbatimInput[fontsize=\scriptsize,breaklines=true,breakanywhere=true]{excerpts/D_InvariantPolynomial.txt}

\textbf{Direct project imports:} \nolinkurl{PaperXI.InertiaFixed}, \nolinkurl{PaperXI.CyclicPolynomial}.

\begin{longtable}{@{}p{0.07\textwidth}p{0.12\textwidth}p{0.75\textwidth}@{}}\toprule Line&Kind&Name\\\midrule\endhead

13&def&\nolinkurl{PaperXI.InvariantPolynomial.operatorPolynomial}\\

16&theorem&\nolinkurl{PaperXI.InvariantPolynomial.operatorPolynomial_eq_det}\\

22&theorem&\nolinkurl{PaperXI.InvariantPolynomial.operatorPolynomial_of_subsingleton}\\

29&theorem&\nolinkurl{PaperXI.InvariantPolynomial.operatorPolynomial_conj}\\

38&def&\nolinkurl{PaperXI.InvariantPolynomial.fixedFrobenius}\\

49&theorem&\nolinkurl{PaperXI.InvariantPolynomial.fixedFrobenius_apply}\\

55&def&\nolinkurl{PaperXI.InvariantPolynomial.localPolynomial}\\

60&theorem&\nolinkurl{PaperXI.InvariantPolynomial.localPolynomial_eq_det}\\

68&theorem&\nolinkurl{PaperXI.InvariantPolynomial.fixed_conj_of_top}\\

76&theorem&\nolinkurl{PaperXI.InvariantPolynomial.localPolynomial_of_top}\\

83&theorem&\nolinkurl{PaperXI.InvariantPolynomial.localPolynomial_of_bot}\\

\bottomrule\end{longtable}

\section{\texorpdfstring{\nolinkurl{PaperXI.IsotypicEvaluation}}{PaperXI.IsotypicEvaluation}}\label{module:PaperXI.IsotypicEvaluation}

\textbf{Selected layer:} v0.44 parent.\quad\textbf{Source lines:} 97.\par\textbf{SHA-256:} {\scriptsize\nolinkurl{9fa381127f5578659025a1972b5e641f6f3000ff3e6a0d17c4f402b6ec08beaa}}.

\textbf{Direct project imports:} none.

\begin{longtable}{@{}p{0.07\textwidth}p{0.12\textwidth}p{0.75\textwidth}@{}}\toprule Line&Kind&Name\\\midrule\endhead

19&def&\nolinkurl{PaperXI.IsotypicEvaluation.map}\\

28&theorem&\nolinkurl{PaperXI.IsotypicEvaluation.map_tmul}\\

31&theorem&\nolinkurl{PaperXI.IsotypicEvaluation.bijective_of_copies}\\

80&theorem&\nolinkurl{PaperXI.IsotypicEvaluation.bijective}\\

87&def&\nolinkurl{PaperXI.IsotypicEvaluation.equiv}\\

92&theorem&\nolinkurl{PaperXI.IsotypicEvaluation.equiv_tmul}\\

\bottomrule\end{longtable}

\section{\texorpdfstring{\nolinkurl{PaperXI.LatticeCongruenceTopology60}}{PaperXI.LatticeCongruenceTopology60}}\label{module:PaperXI.LatticeCongruenceTopology60}

\textbf{Selected layer:} base108 arithmetic branch.\quad\textbf{Source lines:} 144.\par\textbf{SHA-256:} {\scriptsize\nolinkurl{2bdb8a0315e6b0015b8c69df7a1ef382ca5455f3976d5d040c507a94e2736fd0}}.

\VerbatimInput[fontsize=\scriptsize,breaklines=true,breakanywhere=true]{excerpts/D_LatticeCongruenceTopology60.txt}

\textbf{Direct project imports:} \nolinkurl{PaperXI.LatticeMatrixTopology60}, \nolinkurl{PaperXI.PeriodicLatticeReduction60}.

\begin{longtable}{@{}p{0.07\textwidth}p{0.12\textwidth}p{0.75\textwidth}@{}}\toprule Line&Kind&Name\\\midrule\endhead

21&def&\nolinkurl{PaperXI.LatticeCongruenceTopology60.coordinateAut}\\

28&theorem&\nolinkurl{PaperXI.LatticeCongruenceTopology60.coordinateAut_apply}\\

33&def&\nolinkurl{PaperXI.LatticeCongruenceTopology60.congruence}\\

38&theorem&\nolinkurl{PaperXI.LatticeCongruenceTopology60.mem_stabilizer_iff}\\

60&def&\nolinkurl{PaperXI.LatticeCongruenceTopology60.levelSet}\\

65&theorem&\nolinkurl{PaperXI.LatticeCongruenceTopology60.levelSet_eq}\\

72&theorem&\nolinkurl{PaperXI.LatticeCongruenceTopology60.levelSet_closed}\\

89&theorem&\nolinkurl{PaperXI.LatticeCongruenceTopology60.levelSet_open}\\

105&theorem&\nolinkurl{PaperXI.LatticeCongruenceTopology60.congruence_closed}\\

120&theorem&\nolinkurl{PaperXI.LatticeCongruenceTopology60.congruence_compact}\\

130&theorem&\nolinkurl{PaperXI.LatticeCongruenceTopology60.congruence_open}\\

\bottomrule\end{longtable}

\section{\texorpdfstring{\nolinkurl{PaperXI.LatticeMatrixTopology60}}{PaperXI.LatticeMatrixTopology60}}\label{module:PaperXI.LatticeMatrixTopology60}

\textbf{Selected layer:} base108 arithmetic branch.\quad\textbf{Source lines:} 152.\par\textbf{SHA-256:} {\scriptsize\nolinkurl{c924df503533e0221a502da1cdc8e886d4290c90602f3ee33287cdfb9ef98274}}.

\VerbatimInput[fontsize=\scriptsize,breaklines=true,breakanywhere=true]{excerpts/D_LatticeMatrixTopology60.txt}

\textbf{Direct project imports:} \nolinkurl{PaperXI.ActualELatticeQuotient59}.

\begin{longtable}{@{}p{0.07\textwidth}p{0.12\textwidth}p{0.75\textwidth}@{}}\toprule Line&Kind&Name\\\midrule\endhead

22&def&\nolinkurl{PaperXI.LatticeMatrixTopology60.matrixAction}\\

26&theorem&\nolinkurl{PaperXI.LatticeMatrixTopology60.matrixAction_continuous}\\

36&def&\nolinkurl{PaperXI.LatticeMatrixTopology60.preserving}\\

45&def&\nolinkurl{PaperXI.LatticeMatrixTopology60.stabilizer}\\

50&theorem&\nolinkurl{PaperXI.LatticeMatrixTopology60.preserving_closed}\\

62&theorem&\nolinkurl{PaperXI.LatticeMatrixTopology60.preserving_open}\\

73&theorem&\nolinkurl{PaperXI.LatticeMatrixTopology60.preserving_compact}\\

92&theorem&\nolinkurl{PaperXI.LatticeMatrixTopology60.exists_scaled_basis}\\

110&theorem&\nolinkurl{PaperXI.LatticeMatrixTopology60.preserving_compact_of_open}\\

131&theorem&\nolinkurl{PaperXI.LatticeMatrixTopology60.stabilizer_compact}\\

139&theorem&\nolinkurl{PaperXI.LatticeMatrixTopology60.stabilizer_open}\\

145&theorem&\nolinkurl{PaperXI.LatticeMatrixTopology60.stabilizer_compact_of_open}\\

\bottomrule\end{longtable}

\section{\texorpdfstring{\nolinkurl{PaperXI.LiftedEtale}}{PaperXI.LiftedEtale}}\label{module:PaperXI.LiftedEtale}

\textbf{Selected layer:} v0.44 parent.\quad\textbf{Source lines:} 95.\par\textbf{SHA-256:} {\scriptsize\nolinkurl{2fa34a8ba6c429ab131c13111dc85621648e62f7ef6edf662e0367ebbddce8bf}}.

\VerbatimInput[fontsize=\scriptsize,breaklines=true,breakanywhere=true]{excerpts/D_LiftedEtale.txt}

\textbf{Direct project imports:} \nolinkurl{PaperXI.LiftedValuationFinite}, \nolinkurl{PaperXI.LocalResidueRank}.

\begin{longtable}{@{}p{0.07\textwidth}p{0.12\textwidth}p{0.75\textwidth}@{}}\toprule Line&Kind&Name\\\midrule\endhead

20&theorem&\nolinkurl{PaperXI.LiftedEtale.free}\\

28&def&\nolinkurl{PaperXI.LiftedEtale.canonicalResidueLinearEquiv}\\

38&theorem&\nolinkurl{PaperXI.LiftedEtale.rank_residue}\\

45&theorem&\nolinkurl{PaperXI.LiftedEtale.maximalIdeal_map}\\

54&theorem&\nolinkurl{PaperXI.LiftedEtale.etale}\\

65&theorem&\nolinkurl{PaperXI.LiftedEtale.ramificationIdx_one}\\

72&theorem&\nolinkurl{PaperXI.LiftedEtale.uniformizer_ideal}\\

80&theorem&\nolinkurl{PaperXI.LiftedEtale.monogenic}\\

\bottomrule\end{longtable}

\section{\texorpdfstring{\nolinkurl{PaperXI.LiftedFullResidue}}{PaperXI.LiftedFullResidue}}\label{module:PaperXI.LiftedFullResidue}

\textbf{Selected layer:} v0.44 parent.\quad\textbf{Source lines:} 115.\par\textbf{SHA-256:} {\scriptsize\nolinkurl{88447e2315204103283102ff48fb70a06eb201867cbbeceaaba3f94325b44618}}.

\VerbatimInput[fontsize=\scriptsize,breaklines=true,breakanywhere=true]{excerpts/D_LiftedFullResidue.txt}

\textbf{Direct project imports:} \nolinkurl{PaperXI.LiftedResidueDegree}.

\begin{longtable}{@{}p{0.07\textwidth}p{0.12\textwidth}p{0.75\textwidth}@{}}\toprule Line&Kind&Name\\\midrule\endhead

15&def&\nolinkurl{PaperXI.LiftedFullResidue.valuation}\\

17&def&\nolinkurl{PaperXI.LiftedFullResidue.residueMap}\\

21&theorem&\nolinkurl{PaperXI.LiftedFullResidue.residueMap_contains}\\

35&theorem&\nolinkurl{PaperXI.LiftedFullResidue.residueMap_mem}\\

56&theorem&\nolinkurl{PaperXI.LiftedFullResidue.residueMap_range}\\

63&def&\nolinkurl{PaperXI.LiftedFullResidue.toSubfield}\\

71&def&\nolinkurl{PaperXI.LiftedFullResidue.residueEquiv}\\

77&theorem&\nolinkurl{PaperXI.LiftedFullResidue.residueEquiv_val}\\

81&theorem&\nolinkurl{PaperXI.LiftedFullResidue.residue_card}\\

85&def&\nolinkurl{PaperXI.LiftedFullResidue.baseResidueMap}\\

89&theorem&\nolinkurl{PaperXI.LiftedFullResidue.base_composition}\\

95&instance&\nolinkurl{PaperXI.LiftedFullResidue.actualBaseAlgebra}\\

99&def&\nolinkurl{PaperXI.LiftedFullResidue.residueAlgEquiv}\\

109&theorem&\nolinkurl{PaperXI.LiftedFullResidue.residue_degree}\\

\bottomrule\end{longtable}

\section{\texorpdfstring{\nolinkurl{PaperXI.LiftedResidueDegree}}{PaperXI.LiftedResidueDegree}}\label{module:PaperXI.LiftedResidueDegree}

\textbf{Selected layer:} v0.44 parent.\quad\textbf{Source lines:} 82.\par\textbf{SHA-256:} {\scriptsize\nolinkurl{d49acefdc58604537e7f2afe7c3ab0f46eba16fb8096a2dc28f1f8a655103bf7}}.

\VerbatimInput[fontsize=\scriptsize,breaklines=true,breakanywhere=true]{excerpts/D_LiftedResidueDegree.txt}

\textbf{Direct project imports:} \nolinkurl{PaperXI.HenselianGaloisResidue}.

\begin{longtable}{@{}p{0.07\textwidth}p{0.12\textwidth}p{0.75\textwidth}@{}}\toprule Line&Kind&Name\\\midrule\endhead

13&abbrev&\nolinkurl{PaperXI.LiftedResidueDegree.BaseResidue}\\

15&instance&\nolinkurl{PaperXI.LiftedResidueDegree.baseResidueAlgebra}\\

20&def&\nolinkurl{PaperXI.LiftedResidueDegree.residueAut}\\

26&theorem&\nolinkurl{PaperXI.LiftedResidueDegree.residueAut_apply}\\

32&theorem&\nolinkurl{PaperXI.LiftedResidueDegree.residueAut_surjective}\\

41&def&\nolinkurl{PaperXI.LiftedResidueDegree.field}\\

46&theorem&\nolinkurl{PaperXI.LiftedResidueDegree.fixingSubgroup_eq}\\

67&theorem&\nolinkurl{PaperXI.LiftedResidueDegree.degree}\\

75&theorem&\nolinkurl{PaperXI.LiftedResidueDegree.galois}\\

\bottomrule\end{longtable}

\section{\texorpdfstring{\nolinkurl{PaperXI.LiftedResidueFrobenius}}{PaperXI.LiftedResidueFrobenius}}\label{module:PaperXI.LiftedResidueFrobenius}

\textbf{Selected layer:} v0.44 parent.\quad\textbf{Source lines:} 67.\par\textbf{SHA-256:} {\scriptsize\nolinkurl{45e5ccfb6f245b6318a6d628168f8eba77834ea08b554b9f74ba407b49a9f45c}}.

\VerbatimInput[fontsize=\scriptsize,breaklines=true,breakanywhere=true]{excerpts/D_LiftedResidueFrobenius.txt}

\textbf{Direct project imports:} \nolinkurl{PaperXI.LiftedFullResidue}, \nolinkurl{PaperXI.FiniteResidueBaseChange}.

\begin{longtable}{@{}p{0.07\textwidth}p{0.12\textwidth}p{0.75\textwidth}@{}}\toprule Line&Kind&Name\\\midrule\endhead

14&instance&\nolinkurl{PaperXI.LiftedResidueFrobenius.baseFinite}\\

20&def&\nolinkurl{PaperXI.LiftedResidueFrobenius.frobenius}\\

22&theorem&\nolinkurl{PaperXI.LiftedResidueFrobenius.frobenius_inverse_power}\\

30&theorem&\nolinkurl{PaperXI.LiftedResidueFrobenius.frobenius_in_range}\\

38&theorem&\nolinkurl{PaperXI.LiftedResidueFrobenius.decomposition_frobenius_lift}\\

45&theorem&\nolinkurl{PaperXI.LiftedResidueFrobenius.full_degree}\\

50&def&\nolinkurl{PaperXI.LiftedResidueFrobenius.cyclicQuotient}\\

59&theorem&\nolinkurl{PaperXI.LiftedResidueFrobenius.quotient_order}\\

\bottomrule\end{longtable}

\section{\texorpdfstring{\nolinkurl{PaperXI.LiftedValuationFinite}}{PaperXI.LiftedValuationFinite}}\label{module:PaperXI.LiftedValuationFinite}

\textbf{Selected layer:} v0.44 parent.\quad\textbf{Source lines:} 78.\par\textbf{SHA-256:} {\scriptsize\nolinkurl{b4d29f529bdede481493a48e735a4998dd0ff3f401796cc652dbc870a82c2b0e}}.

\VerbatimInput[fontsize=\scriptsize,breaklines=true,breakanywhere=true]{excerpts/D_LiftedValuationFinite.txt}

\textbf{Direct project imports:} \nolinkurl{PaperXI.ValuationIntegralModel}, \nolinkurl{PaperXI.LiftedFullResidue}.

\begin{longtable}{@{}p{0.07\textwidth}p{0.12\textwidth}p{0.75\textwidth}@{}}\toprule Line&Kind&Name\\\midrule\endhead

17&instance&\nolinkurl{PaperXI.LiftedValuationFinite.baseAlgebra}\\

20&instance&\nolinkurl{PaperXI.LiftedValuationFinite.baseTower}\\

23&instance&\nolinkurl{PaperXI.LiftedValuationFinite.baseFieldTower}\\

26&instance&\nolinkurl{PaperXI.LiftedValuationFinite.baseLocal}\\

34&theorem&\nolinkurl{PaperXI.LiftedValuationFinite.stable_valuation}\\

51&theorem&\nolinkurl{PaperXI.LiftedValuationFinite.integralClosure}\\

58&theorem&\nolinkurl{PaperXI.LiftedValuationFinite.finite}\\

68&theorem&\nolinkurl{PaperXI.LiftedValuationFinite.rank}\\

\bottomrule\end{longtable}

\section{\texorpdfstring{\nolinkurl{PaperXI.LiftedValuationInertia}}{PaperXI.LiftedValuationInertia}}\label{module:PaperXI.LiftedValuationInertia}

\textbf{Selected layer:} v0.44 parent.\quad\textbf{Source lines:} 51.\par\textbf{SHA-256:} {\scriptsize\nolinkurl{84e9166787ac68126346acfa087e68e5b5ad325b43fe0b761ac2b001f1ba310b}}.

\VerbatimInput[fontsize=\scriptsize,breaklines=true,breakanywhere=true]{excerpts/D_LiftedValuationInertia.txt}

\textbf{Direct project imports:} \nolinkurl{PaperXI.LiftedFullResidue}.

\begin{longtable}{@{}p{0.07\textwidth}p{0.12\textwidth}p{0.75\textwidth}@{}}\toprule Line&Kind&Name\\\midrule\endhead

16&theorem&\nolinkurl{PaperXI.LiftedValuationInertia.inertia_trivial}\\

\bottomrule\end{longtable}

\section{\texorpdfstring{\nolinkurl{PaperXI.LineCharacter}}{PaperXI.LineCharacter}}\label{module:PaperXI.LineCharacter}

\textbf{Selected layer:} v0.44 parent.\quad\textbf{Source lines:} 172.\par\textbf{SHA-256:} {\scriptsize\nolinkurl{fe421a06982cd72d4f80a09db6f3ee52a249a6be9ded4bcd4aa927800c4ed007}}.

\VerbatimInput[fontsize=\scriptsize,breaklines=true,breakanywhere=true]{excerpts/D_LineCharacter.txt}

\textbf{Direct project imports:} \nolinkurl{PaperXI.PacketBasis}.

\begin{longtable}{@{}p{0.07\textwidth}p{0.12\textwidth}p{0.75\textwidth}@{}}\toprule Line&Kind&Name\\\midrule\endhead

11&theorem&\nolinkurl{PaperXI.LineCharacter.scalar_unique}\\

16&theorem&\nolinkurl{PaperXI.LineCharacter.scalar_ne_zero}\\

22&def&\nolinkurl{PaperXI.LineCharacter.lineStabilizer}\\

39&theorem&\nolinkurl{PaperXI.LineCharacter.mem_iff_preserves_line}\\

70&def&\nolinkurl{PaperXI.LineCharacter.scalar}\\

73&theorem&\nolinkurl{PaperXI.LineCharacter.scalar_spec}\\

77&def&\nolinkurl{PaperXI.LineCharacter.character}\\

94&theorem&\nolinkurl{PaperXI.LineCharacter.character_action}\\

97&theorem&\nolinkurl{PaperXI.LineCharacter.power_action}\\

109&theorem&\nolinkurl{PaperXI.LineCharacter.inertia_mem}\\

116&theorem&\nolinkurl{PaperXI.LineCharacter.return_mem}\\

123&theorem&\nolinkurl{PaperXI.LineCharacter.character_on_inertia}\\

133&theorem&\nolinkurl{PaperXI.LineCharacter.character_on_return}\\

143&theorem&\nolinkurl{PaperXI.LineCharacter.packet_stabilizer_iff}\\

159&theorem&\nolinkurl{PaperXI.LineCharacter.expected_stabilizer_le}\\

\bottomrule\end{longtable}

\section{\texorpdfstring{\nolinkurl{PaperXI.LocalFactorIdentification}}{PaperXI.LocalFactorIdentification}}\label{module:PaperXI.LocalFactorIdentification}

\textbf{Selected layer:} v0.44 parent.\quad\textbf{Source lines:} 150.\par\textbf{SHA-256:} {\scriptsize\nolinkurl{3e81e1ce1c7b50d2771432e73d46f02de2a2ce4910e0a3420aa67e6b02a514f3}}.

\VerbatimInput[fontsize=\scriptsize,breaklines=true,breakanywhere=true]{excerpts/D_LocalFactorIdentification.txt}

\textbf{Direct project imports:} \nolinkurl{PaperXI.AtomicReconstruction}, \nolinkurl{PaperXI.CharacterProjector}.

\begin{longtable}{@{}p{0.07\textwidth}p{0.12\textwidth}p{0.75\textwidth}@{}}\toprule Line&Kind&Name\\\midrule\endhead

14&structure&\nolinkurl{PaperXI.LocalFactorIdentification.FactorSystem}\\

25&structure&\nolinkurl{PaperXI.LocalFactorIdentification.Realization}\\

32&def&\nolinkurl{PaperXI.LocalFactorIdentification.chosenAtom}\\

34&theorem&\nolinkurl{PaperXI.LocalFactorIdentification.chosenAtom_realizes}\\

37&def&\nolinkurl{PaperXI.LocalFactorIdentification.candidate}\\

40&theorem&\nolinkurl{PaperXI.LocalFactorIdentification.class_eq_iff}\\

49&theorem&\nolinkurl{PaperXI.LocalFactorIdentification.candidate_eq_reference}\\

59&theorem&\nolinkurl{PaperXI.LocalFactorIdentification.unique_realized_class}\\

70&theorem&\nolinkurl{PaperXI.LocalFactorIdentification.realization_independence}\\

81&theorem&\nolinkurl{PaperXI.LocalFactorIdentification.candidate_congr}\\

89&def&\nolinkurl{PaperXI.LocalFactorIdentification.onClasses}\\

95&theorem&\nolinkurl{PaperXI.LocalFactorIdentification.constructed_comparison}\\

105&theorem&\nolinkurl{PaperXI.LocalFactorIdentification.reconstructed_comparison}\\

117&def&\nolinkurl{PaperXI.LocalFactorIdentification.representationSetoid}\\

124&theorem&\nolinkurl{PaperXI.LocalFactorIdentification.finite_character_separates}\\

140&def&\nolinkurl{PaperXI.LocalFactorIdentification.finiteCharacterSystem}\\

\bottomrule\end{longtable}

\section{\texorpdfstring{\nolinkurl{PaperXI.LocalFieldResidueExtension}}{PaperXI.LocalFieldResidueExtension}}\label{module:PaperXI.LocalFieldResidueExtension}

\textbf{Selected layer:} v0.44 parent.\quad\textbf{Source lines:} 63.\par\textbf{SHA-256:} {\scriptsize\nolinkurl{e7438b3803d39699e4ac020c70f4f4d2377d99e32af405cdea45fa2eb1c9cb9e}}.

\VerbatimInput[fontsize=\scriptsize,breaklines=true,breakanywhere=true]{excerpts/D_LocalFieldResidueExtension.txt}

\textbf{Direct project imports:} \nolinkurl{PaperXI.FiniteFreeCompletion}.

\begin{longtable}{@{}p{0.07\textwidth}p{0.12\textwidth}p{0.75\textwidth}@{}}\toprule Line&Kind&Name\\\midrule\endhead

16&abbrev&\nolinkurl{PaperXI.LocalFieldResidueExtension.ring}\\

17&abbrev&\nolinkurl{PaperXI.LocalFieldResidueExtension.field}\\

19&instance&\nolinkurl{PaperXI.LocalFieldResidueExtension.baseAlgebra}\\

22&instance&\nolinkurl{PaperXI.LocalFieldResidueExtension.baseTower}\\

28&theorem&\nolinkurl{PaperXI.LocalFieldResidueExtension.degree}\\

32&instance&\nolinkurl{PaperXI.LocalFieldResidueExtension.finiteExtension}\\

35&theorem&\nolinkurl{PaperXI.LocalFieldResidueExtension.source}\\

50&abbrev&\nolinkurl{PaperXI.LocalFieldResidueExtension.residueExtension}\\

52&theorem&\nolinkurl{PaperXI.LocalFieldResidueExtension.every_positive_degree}\\

\bottomrule\end{longtable}

\section{\texorpdfstring{\nolinkurl{PaperXI.LocalResidueRank}}{PaperXI.LocalResidueRank}}\label{module:PaperXI.LocalResidueRank}

\textbf{Selected layer:} v0.44 parent.\quad\textbf{Source lines:} 83.\par\textbf{SHA-256:} {\scriptsize\nolinkurl{04aa1ea55d4707f48f22f293dff3a05669d750c0ecaf1dca52bce090c7553fd3}}.

\VerbatimInput[fontsize=\scriptsize,breaklines=true,breakanywhere=true]{excerpts/D_LocalResidueRank.txt}

\textbf{Direct project imports:} none.

\begin{longtable}{@{}p{0.07\textwidth}p{0.12\textwidth}p{0.75\textwidth}@{}}\toprule Line&Kind&Name\\\midrule\endhead

19&theorem&\nolinkurl{PaperXI.LocalResidueRank.map_maximal_of_rank}\\

51&theorem&\nolinkurl{PaperXI.LocalResidueRank.formallyUnramified}\\

57&theorem&\nolinkurl{PaperXI.LocalResidueRank.etale}\\

66&theorem&\nolinkurl{PaperXI.LocalResidueRank.etale_iff_rank_residue}\\

77&theorem&\nolinkurl{PaperXI.LocalResidueRank.ramificationIdx_one}\\

\bottomrule\end{longtable}

\section{\texorpdfstring{\nolinkurl{PaperXI.LocallyFiniteSemisimple}}{PaperXI.LocallyFiniteSemisimple}}\label{module:PaperXI.LocallyFiniteSemisimple}

\textbf{Selected layer:} base108 arithmetic branch.\quad\textbf{Source lines:} 108.\par\textbf{SHA-256:} {\scriptsize\nolinkurl{9d51826538fb9c038407e4a3d741ccbb3ad1c1266fc6642f1e29df6b81cefdd8}}.

\VerbatimInput[fontsize=\scriptsize,breaklines=true,breakanywhere=true]{excerpts/D_LocallyFiniteSemisimple.txt}

\textbf{Direct project imports:} \nolinkurl{PaperXI.QuotientRepresentation}.

\begin{longtable}{@{}p{0.07\textwidth}p{0.12\textwidth}p{0.75\textwidth}@{}}\toprule Line&Kind&Name\\\midrule\endhead

21&theorem&\nolinkurl{PaperXI.LocallyFiniteSemisimple.module_of_locally_semisimple}\\

39&theorem&\nolinkurl{PaperXI.LocallyFiniteSemisimple.semisimple_comp_iff}\\

45&theorem&\nolinkurl{PaperXI.LocallyFiniteSemisimple.semisimple_finite_pullback}\\

52&theorem&\nolinkurl{PaperXI.LocallyFiniteSemisimple.semisimple_of_finite_quotient}\\

59&theorem&\nolinkurl{PaperXI.LocallyFiniteSemisimple.semisimple_asSubmodule}\\

81&theorem&\nolinkurl{PaperXI.LocallyFiniteSemisimple.semisimple_of_locally_semisimple}\\

94&theorem&\nolinkurl{PaperXI.LocallyFiniteSemisimple.semisimple_of_locally_finite_quotient}\\

\bottomrule\end{longtable}

\section{\texorpdfstring{\nolinkurl{PaperXI.MonicResidueAlgebra}}{PaperXI.MonicResidueAlgebra}}\label{module:PaperXI.MonicResidueAlgebra}

\textbf{Selected layer:} v0.44 parent.\quad\textbf{Source lines:} 87.\par\textbf{SHA-256:} {\scriptsize\nolinkurl{61e77b723a6fc1939e003a5cfd834cb152674a48112bee181e683903cdccdca5}}.

\VerbatimInput[fontsize=\scriptsize,breaklines=true,breakanywhere=true]{excerpts/D_MonicResidueAlgebra.txt}

\textbf{Direct project imports:} none.

\begin{longtable}{@{}p{0.07\textwidth}p{0.12\textwidth}p{0.75\textwidth}@{}}\toprule Line&Kind&Name\\\midrule\endhead

17&abbrev&\nolinkurl{PaperXI.MonicResidueAlgebra.specialIdeal}\\

20&instance&\nolinkurl{PaperXI.MonicResidueAlgebra.finite}\\

21&instance&\nolinkurl{PaperXI.MonicResidueAlgebra.free}\\

24&def&\nolinkurl{PaperXI.MonicResidueAlgebra.specialFiberEquiv}\\

29&theorem&\nolinkurl{PaperXI.MonicResidueAlgebra.specialIdeal_maximal}\\

34&instance&\nolinkurl{PaperXI.MonicResidueAlgebra.localRing}\\

46&theorem&\nolinkurl{PaperXI.MonicResidueAlgebra.maximalIdeal_eq}\\

49&def&\nolinkurl{PaperXI.MonicResidueAlgebra.residueEquiv}\\

53&theorem&\nolinkurl{PaperXI.MonicResidueAlgebra.rank}\\

56&theorem&\nolinkurl{PaperXI.MonicResidueAlgebra.degree_pos}\\

60&instance&\nolinkurl{PaperXI.MonicResidueAlgebra.localHom}\\

63&theorem&\nolinkurl{PaperXI.MonicResidueAlgebra.residueEquiv_mk}\\

72&def&\nolinkurl{PaperXI.MonicResidueAlgebra.residueAlgebraEquiv}\\

81&theorem&\nolinkurl{PaperXI.MonicResidueAlgebra.residue_rank}\\

\bottomrule\end{longtable}

\section{\texorpdfstring{\nolinkurl{PaperXI.MonodromyAtom}}{PaperXI.MonodromyAtom}}\label{module:PaperXI.MonodromyAtom}

\textbf{Selected layer:} v0.44 parent.\quad\textbf{Source lines:} 117.\par\textbf{SHA-256:} {\scriptsize\nolinkurl{f53d79ddc74c43995955bc45c8bce600d821bb6d118e9426d85e8478c0829bf2}}.

\VerbatimInput[fontsize=\scriptsize,breaklines=true,breakanywhere=true]{excerpts/D_MonodromyAtom.txt}

\textbf{Direct project imports:} \nolinkurl{PaperXI.MonodromyChain}.

\begin{longtable}{@{}p{0.07\textwidth}p{0.12\textwidth}p{0.75\textwidth}@{}}\toprule Line&Kind&Name\\\midrule\endhead

18&structure&\nolinkurl{PaperXI.MonodromyAtom.Datum}\\

25&def&\nolinkurl{PaperXI.MonodromyAtom.action}\\

40&theorem&\nolinkurl{PaperXI.MonodromyAtom.action_apply}\\

44&def&\nolinkurl{PaperXI.MonodromyAtom.monodromy}\\

46&theorem&\nolinkurl{PaperXI.MonodromyAtom.nilpotence}\\

50&theorem&\nolinkurl{PaperXI.MonodromyAtom.commutation}\\

64&theorem&\nolinkurl{PaperXI.MonodromyAtom.covariance}\\

70&def&\nolinkurl{PaperXI.MonodromyAtom.sp}\\

78&theorem&\nolinkurl{PaperXI.MonodromyAtom.dimension}\\

83&def&\nolinkurl{PaperXI.MonodromyAtom.mapEquiv}\\

97&theorem&\nolinkurl{PaperXI.MonodromyAtom.mapEquiv_monodromy}\\

108&theorem&\nolinkurl{PaperXI.MonodromyAtom.scalar_twist}\\

\bottomrule\end{longtable}

\section{\texorpdfstring{\nolinkurl{PaperXI.MonodromyChain}}{PaperXI.MonodromyChain}}\label{module:PaperXI.MonodromyChain}

\textbf{Selected layer:} v0.44 parent.\quad\textbf{Source lines:} 122.\par\textbf{SHA-256:} {\scriptsize\nolinkurl{3cf19949976f59a32001fee38e472b4bf8f9fda9d4d7c67e3819ff2e27ab101c}}.

\VerbatimInput[fontsize=\scriptsize,breaklines=true,breakanywhere=true]{excerpts/D_MonodromyChain.txt}

\textbf{Direct project imports:} none.

\begin{longtable}{@{}p{0.07\textwidth}p{0.12\textwidth}p{0.75\textwidth}@{}}\toprule Line&Kind&Name\\\midrule\endhead

16&def&\nolinkurl{PaperXI.MonodromyChain.shift}\\

21&theorem&\nolinkurl{PaperXI.MonodromyChain.shift_apply}\\

24&def&\nolinkurl{PaperXI.MonodromyChain.path}\\

28&theorem&\nolinkurl{PaperXI.MonodromyChain.path_zero}\\

29&theorem&\nolinkurl{PaperXI.MonodromyChain.path_succ}\\

33&theorem&\nolinkurl{PaperXI.MonodromyChain.pow_apply}\\

51&theorem&\nolinkurl{PaperXI.MonodromyChain.nilpotent_bound}\\

59&def&\nolinkurl{PaperXI.MonodromyChain.projector}\\

64&theorem&\nolinkurl{PaperXI.MonodromyChain.projector_apply}\\

67&theorem&\nolinkurl{PaperXI.MonodromyChain.projector_idempotent}\\

74&theorem&\nolinkurl{PaperXI.MonodromyChain.projector_orthogonal}\\

83&theorem&\nolinkurl{PaperXI.MonodromyChain.projector_sum}\\

90&theorem&\nolinkurl{PaperXI.MonodromyChain.block_vanishes}\\

109&theorem&\nolinkurl{PaperXI.MonodromyChain.block_reconstruction}\\

\bottomrule\end{longtable}

\section{\texorpdfstring{\nolinkurl{PaperXI.MonomialAction}}{PaperXI.MonomialAction}}\label{module:PaperXI.MonomialAction}

\textbf{Selected layer:} v0.44 parent.\quad\textbf{Source lines:} 92.\par\textbf{SHA-256:} {\scriptsize\nolinkurl{70169fa35998c99c6d9830b345ee188e42778dace1f244f1181413f956a8961b}}.

\VerbatimInput[fontsize=\scriptsize,breaklines=true,breakanywhere=true]{excerpts/D_MonomialAction.txt}

\textbf{Direct project imports:} none.

\begin{longtable}{@{}p{0.07\textwidth}p{0.12\textwidth}p{0.75\textwidth}@{}}\toprule Line&Kind&Name\\\midrule\endhead

16&def&\nolinkurl{PaperXI.MonomialAction.operator}\\

18&theorem&\nolinkurl{PaperXI.MonomialAction.operator_basis}\\

22&def&\nolinkurl{PaperXI.MonomialAction.representation}\\

37&theorem&\nolinkurl{PaperXI.MonomialAction.representation_basis}\\

42&theorem&\nolinkurl{PaperXI.MonomialAction.scalar_injective}\\

47&theorem&\nolinkurl{PaperXI.MonomialAction.weight_one}\\

54&theorem&\nolinkurl{PaperXI.MonomialAction.weight_mul}\\

63&theorem&\nolinkurl{PaperXI.MonomialAction.representation_eq}\\

76&def&\nolinkurl{PaperXI.MonomialAction.tensorPermutation}\\

82&theorem&\nolinkurl{PaperXI.MonomialAction.tensor_basis}\\

\bottomrule\end{longtable}

\section{\texorpdfstring{\nolinkurl{PaperXI.MonomialDual}}{PaperXI.MonomialDual}}\label{module:PaperXI.MonomialDual}

\textbf{Selected layer:} v0.44 parent.\quad\textbf{Source lines:} 46.\par\textbf{SHA-256:} {\scriptsize\nolinkurl{1f8a72858be61f32ca3a9daa195e4e3045e31db32c4be366e1a73c604753e8b8}}.

\VerbatimInput[fontsize=\scriptsize,breaklines=true,breakanywhere=true]{excerpts/D_MonomialDual.txt}

\textbf{Direct project imports:} \nolinkurl{PaperXI.MonomialAction}.

\begin{longtable}{@{}p{0.07\textwidth}p{0.12\textwidth}p{0.75\textwidth}@{}}\toprule Line&Kind&Name\\\midrule\endhead

18&theorem&\nolinkurl{PaperXI.MonomialDual.inverse_weight}\\

26&theorem&\nolinkurl{PaperXI.MonomialDual.dual_basis_action}\\

\bottomrule\end{longtable}

\section{\texorpdfstring{\nolinkurl{PaperXI.MonomialInduction}}{PaperXI.MonomialInduction}}\label{module:PaperXI.MonomialInduction}

\textbf{Selected layer:} v0.44 parent.\quad\textbf{Source lines:} 152.\par\textbf{SHA-256:} {\scriptsize\nolinkurl{9406f172392d78a75931d9a416f9e701e1dc43b326140a53ab238c582b35813b}}.

\VerbatimInput[fontsize=\scriptsize,breaklines=true,breakanywhere=true]{excerpts/D_MonomialInduction.txt}

\textbf{Direct project imports:} \nolinkurl{PaperXI.LineCharacter}.

\begin{longtable}{@{}p{0.07\textwidth}p{0.12\textwidth}p{0.75\textwidth}@{}}\toprule Line&Kind&Name\\\midrule\endhead

12&def&\nolinkurl{PaperXI.MonomialInduction.scalarRepresentation}\\

17&theorem&\nolinkurl{PaperXI.MonomialInduction.scalarRepresentation_apply}\\

22&abbrev&\nolinkurl{PaperXI.MonomialInduction.stabilizer}\\

23&abbrev&\nolinkurl{PaperXI.MonomialInduction.character}\\

24&abbrev&\nolinkurl{PaperXI.MonomialInduction.lineRepresentation}\\

25&abbrev&\nolinkurl{PaperXI.MonomialInduction.induced}\\

26&abbrev&\nolinkurl{PaperXI.MonomialInduction.InducedSpace}\\

29&def&\nolinkurl{PaperXI.MonomialInduction.lineMap}\\

40&def&\nolinkurl{PaperXI.MonomialInduction.evaluation}\\

52&theorem&\nolinkurl{PaperXI.MonomialInduction.evaluation_mk}\\

58&theorem&\nolinkurl{PaperXI.MonomialInduction.mk_stabilizer}\\

71&def&\nolinkurl{PaperXI.MonomialInduction.reconstruction}\\

75&theorem&\nolinkurl{PaperXI.MonomialInduction.reconstruction_basis}\\

83&theorem&\nolinkurl{PaperXI.MonomialInduction.evaluation_reconstruction}\\

92&theorem&\nolinkurl{PaperXI.MonomialInduction.representative_inverse}\\

100&theorem&\nolinkurl{PaperXI.MonomialInduction.reconstruction_evaluation}\\

133&theorem&\nolinkurl{PaperXI.MonomialInduction.evaluation_bijective}\\

144&def&\nolinkurl{PaperXI.MonomialInduction.equivalence}\\

147&theorem&\nolinkurl{PaperXI.MonomialInduction.equivalence_mk}\\

\bottomrule\end{longtable}

\section{\texorpdfstring{\nolinkurl{PaperXI.MonomialInductionSubgroup}}{PaperXI.MonomialInductionSubgroup}}\label{module:PaperXI.MonomialInductionSubgroup}

\textbf{Selected layer:} v0.44 parent.\quad\textbf{Source lines:} 69.\par\textbf{SHA-256:} {\scriptsize\nolinkurl{392f827d63ce2b3a4c1ede29d1d61d70de0f705e596b32f42c617d03a6a96910}}.

\VerbatimInput[fontsize=\scriptsize,breaklines=true,breakanywhere=true]{excerpts/D_MonomialInductionSubgroup.txt}

\textbf{Direct project imports:} \nolinkurl{PaperXI.MonomialInductionTransitive}.

\begin{longtable}{@{}p{0.07\textwidth}p{0.12\textwidth}p{0.75\textwidth}@{}}\toprule Line&Kind&Name\\\midrule\endhead

12&def&\nolinkurl{PaperXI.MonomialInduction.equivalenceOfSubgroup}\\

31&def&\nolinkurl{PaperXI.MonomialInduction.transitiveEquivalenceOfSubgroupAction}\\

50&def&\nolinkurl{PaperXI.MonomialInduction.transitiveEquivalenceOfSubgroup}\\

\bottomrule\end{longtable}

\section{\texorpdfstring{\nolinkurl{PaperXI.MonomialInductionTransitive}}{PaperXI.MonomialInductionTransitive}}\label{module:PaperXI.MonomialInductionTransitive}

\textbf{Selected layer:} v0.44 parent.\quad\textbf{Source lines:} 69.\par\textbf{SHA-256:} {\scriptsize\nolinkurl{371cd6e7cd417a0f86606695b9bfa11c152a0a424eb0bf642fcdcb4a0ac04140}}.

\VerbatimInput[fontsize=\scriptsize,breaklines=true,breakanywhere=true]{excerpts/D_MonomialInductionTransitive.txt}

\textbf{Direct project imports:} \nolinkurl{PaperXI.MonomialInduction}.

\begin{longtable}{@{}p{0.07\textwidth}p{0.12\textwidth}p{0.75\textwidth}@{}}\toprule Line&Kind&Name\\\midrule\endhead

15&theorem&\nolinkurl{PaperXI.MonomialInduction.mem_lineStabilizer_iff}\\

28&theorem&\nolinkurl{PaperXI.MonomialInduction.character_eq_weight}\\

38&def&\nolinkurl{PaperXI.MonomialInduction.representative}\\

40&theorem&\nolinkurl{PaperXI.MonomialInduction.representative_spec}\\

44&theorem&\nolinkurl{PaperXI.MonomialInduction.representative_action}\\

50&theorem&\nolinkurl{PaperXI.MonomialInduction.orbit_is_monomial}\\

56&def&\nolinkurl{PaperXI.MonomialInduction.transitiveEquivalence}\\

63&theorem&\nolinkurl{PaperXI.MonomialInduction.transitiveEquivalence_mk}\\

\bottomrule\end{longtable}

\section{\texorpdfstring{\nolinkurl{PaperXI.MonomialOrbit}}{PaperXI.MonomialOrbit}}\label{module:PaperXI.MonomialOrbit}

\textbf{Selected layer:} v0.44 parent.\quad\textbf{Source lines:} 82.\par\textbf{SHA-256:} {\scriptsize\nolinkurl{7552a5e8d7825191e603a7e2b8ad25bcf97b785a4fc4d3dd7e48cb7cdc0c1734}}.

\VerbatimInput[fontsize=\scriptsize,breaklines=true,breakanywhere=true]{excerpts/D_MonomialOrbit.txt}

\textbf{Direct project imports:} \nolinkurl{PaperXI.MonomialAction}.

\begin{longtable}{@{}p{0.07\textwidth}p{0.12\textwidth}p{0.75\textwidth}@{}}\toprule Line&Kind&Name\\\midrule\endhead

23&def&\nolinkurl{PaperXI.MonomialOrbit.block}\\

29&theorem&\nolinkurl{PaperXI.MonomialOrbit.block_basis}\\

34&theorem&\nolinkurl{PaperXI.MonomialOrbit.block_single}\\

40&def&\nolinkurl{PaperXI.MonomialOrbit.sumBasis}\\

43&theorem&\nolinkurl{PaperXI.MonomialOrbit.sumBasis_apply}\\

50&def&\nolinkurl{PaperXI.MonomialOrbit.linearEquiv}\\

53&theorem&\nolinkurl{PaperXI.MonomialOrbit.linearEquiv_basis}\\

59&theorem&\nolinkurl{PaperXI.MonomialOrbit.intertwines}\\

72&def&\nolinkurl{PaperXI.MonomialOrbit.representationEquiv}\\

79&theorem&\nolinkurl{PaperXI.MonomialOrbit.block_finrank}\\

\bottomrule\end{longtable}

\section{\texorpdfstring{\nolinkurl{PaperXI.MoritaActualCartan110}}{PaperXI.MoritaActualCartan110}}\label{module:PaperXI.MoritaActualCartan110}

\textbf{Selected layer:} Continuation110.\quad\textbf{Source lines:} 266.\par\textbf{SHA-256:} {\scriptsize\nolinkurl{dd18d3a60333f0e737ef4ff22c34db295f49460de701a9fddf9fcd3dc22d293a}}.

\VerbatimInput[fontsize=\scriptsize,breaklines=true,breakanywhere=true]{excerpts/D_MoritaActualCartan110.txt}

\textbf{Direct project imports:} \nolinkurl{PaperXI.MoritaIndex110}, \nolinkurl{PaperXI.ActualLocalCartan106}, \nolinkurl{PaperXI.MoritaCartanTransport110}, \nolinkurl{PaperXI.MoritaScalarUpgrade110}.

\begin{longtable}{@{}p{0.07\textwidth}p{0.12\textwidth}p{0.75\textwidth}@{}}\toprule Line&Kind&Name\\\midrule\endhead

21&def&\nolinkurl{PaperXI.MoritaActualCartan110.centralizerEmbedding}\\

24&def&\nolinkurl{PaperXI.MoritaActualCartan110.scalarEmbedding}\\

27&def&\nolinkurl{PaperXI.MoritaActualCartan110.uniformizerEndUnit}\\

32&theorem&\nolinkurl{PaperXI.MoritaActualCartan110.centralizerEmbedding_integral}\\

39&theorem&\nolinkurl{PaperXI.MoritaActualCartan110.scalarEmbedding_integralUnit}\\

54&theorem&\nolinkurl{PaperXI.MoritaActualCartan110.scalarEmbedding_uniformizer}\\

72&theorem&\nolinkurl{PaperXI.MoritaActualCartan110.scalarEmbedding_normalizes}\\

93&theorem&\nolinkurl{PaperXI.MoritaActualCartan110.centralizerEmbedding_cartan}\\

127&def&\nolinkurl{PaperXI.MoritaActualCartan110.RadicalCommutes}\\

137&theorem&\nolinkurl{PaperXI.MoritaActualCartan110.relIndex_of_actual_cartan}\\

182&theorem&\nolinkurl{PaperXI.MoritaActualCartan110.exists_actual_cartan_index}\\

206&theorem&\nolinkurl{PaperXI.MoritaActualCartan110.exists_actual_cartan_index_of_base_lattice}\\

232&theorem&\nolinkurl{PaperXI.MoritaActualCartan110.relIndex_isEmpty}\\

246&theorem&\nolinkurl{PaperXI.MoritaActualCartan110.relIndex_subsingleton}\\

\bottomrule\end{longtable}

\section{\texorpdfstring{\nolinkurl{PaperXI.MoritaCartanTransport110}}{PaperXI.MoritaCartanTransport110}}\label{module:PaperXI.MoritaCartanTransport110}

\textbf{Selected layer:} Continuation110.\quad\textbf{Source lines:} 121.\par\textbf{SHA-256:} {\scriptsize\nolinkurl{61b021d7f9fd0fecaebb1a9dc3c89a7cbe30760ebecf6120c0af691d1e0e0645}}.

\VerbatimInput[fontsize=\scriptsize,breaklines=true,breakanywhere=true]{excerpts/D_MoritaCartanTransport110.txt}

\textbf{Direct project imports:} \nolinkurl{PaperXI.MoritaLattice110}.

\begin{longtable}{@{}p{0.07\textwidth}p{0.12\textwidth}p{0.75\textwidth}@{}}\toprule Line&Kind&Name\\\midrule\endhead

16&def&\nolinkurl{PaperXI.MoritaCartanTransport110.conjugate}\\

20&theorem&\nolinkurl{PaperXI.MoritaCartanTransport110.unitConjugation_mul}\\

26&theorem&\nolinkurl{PaperXI.MoritaCartanTransport110.unitConjugation_inv_apply_apply}\\

30&theorem&\nolinkurl{PaperXI.MoritaCartanTransport110.conjugate_mul}\\

36&theorem&\nolinkurl{PaperXI.MoritaCartanTransport110.conjugate_eq_of_normalizing}\\

44&theorem&\nolinkurl{PaperXI.MoritaCartanTransport110.conjugate_inf}\\

50&theorem&\nolinkurl{PaperXI.MoritaCartanTransport110.intersection_outer_factors}\\

58&theorem&\nolinkurl{PaperXI.MoritaCartanTransport110.intersection_relIndex_outer_factors}\\

77&def&\nolinkurl{PaperXI.MoritaCartanTransport110.mappedMatrixUnit}\\

80&theorem&\nolinkurl{PaperXI.MoritaCartanTransport110.mappedMatrixUnit_val}\\

83&theorem&\nolinkurl{PaperXI.MoritaCartanTransport110.mappedMatrixUnit_inv}\\

89&theorem&\nolinkurl{PaperXI.MoritaCartanTransport110.mappedMatrixUnit_conjugation_mem}\\

98&theorem&\nolinkurl{PaperXI.MoritaCartanTransport110.mappedMatrixUnit_mem_normalizingUnits}\\

110&theorem&\nolinkurl{PaperXI.MoritaCartanTransport110.matrix_intersection_relIndex_outer_factors}\\

\bottomrule\end{longtable}

\section{\texorpdfstring{\nolinkurl{PaperXI.MoritaCornerRank110}}{PaperXI.MoritaCornerRank110}}\label{module:PaperXI.MoritaCornerRank110}

\textbf{Selected layer:} Continuation110.\quad\textbf{Source lines:} 132.\par\textbf{SHA-256:} {\scriptsize\nolinkurl{8b3f281926d04dbaa4b40d47b74138cda2b422b7f8c4fd68785b65d37586aed2}}.

\VerbatimInput[fontsize=\scriptsize,breaklines=true,breakanywhere=true]{excerpts/D_MoritaCornerRank110.txt}

\textbf{Direct project imports:} \nolinkurl{PaperXI.MoritaLattice110}.

\begin{longtable}{@{}p{0.07\textwidth}p{0.12\textwidth}p{0.75\textwidth}@{}}\toprule Line&Kind&Name\\\midrule\endhead

22&theorem&\nolinkurl{PaperXI.MoritaCornerRank110.exists_basis_full_submodule}\\

63&def&\nolinkurl{PaperXI.MoritaCornerRank110.coefficientMap}\\

67&def&\nolinkurl{PaperXI.MoritaCornerRank110.cornerSubmodule}\\

79&theorem&\nolinkurl{PaperXI.MoritaCornerRank110.cornerSubmodule_eq_map}\\

92&theorem&\nolinkurl{PaperXI.MoritaCornerRank110.corner_fg}\\

100&theorem&\nolinkurl{PaperXI.MoritaCornerRank110.corner_full}\\

120&theorem&\nolinkurl{PaperXI.MoritaCornerRank110.exists_corner_basis_degree}\\

\bottomrule\end{longtable}

\section{\texorpdfstring{\nolinkurl{PaperXI.MoritaEndomorphismCoordinates111}}{PaperXI.MoritaEndomorphismCoordinates111}}\label{module:PaperXI.MoritaEndomorphismCoordinates111}

\textbf{Selected layer:} Continuation111.\quad\textbf{Source lines:} 176.\par\textbf{SHA-256:} {\scriptsize\nolinkurl{d0023761aee69e81f53222a6e334e6ac57122a15fd64ce9017ee0bc0e27663c0}}.

\VerbatimInput[fontsize=\scriptsize,breaklines=true,breakanywhere=true]{excerpts/D_MoritaEndomorphismCoordinates111.txt}

\textbf{Direct project imports:} \nolinkurl{PaperXI.MoritaActualCartan110}, \nolinkurl{PaperXI.ActualCuspidalCandidate105}.

\begin{longtable}{@{}p{0.07\textwidth}p{0.12\textwidth}p{0.75\textwidth}@{}}\toprule Line&Kind&Name\\\midrule\endhead

20&def&\nolinkurl{PaperXI.MoritaEndomorphismCoordinates111.blockLinearEquiv}\\

24&theorem&\nolinkurl{PaperXI.MoritaEndomorphismCoordinates111.blockLinearEquiv_apply}\\

31&theorem&\nolinkurl{PaperXI.MoritaEndomorphismCoordinates111.blockLinearEquiv_mul}\\

41&def&\nolinkurl{PaperXI.MoritaEndomorphismCoordinates111.blockRingEquiv}\\

45&theorem&\nolinkurl{PaperXI.MoritaEndomorphismCoordinates111.blockRingEquiv_apply}\\

58&def&\nolinkurl{PaperXI.MoritaEndomorphismCoordinates111.coordinateRingEquiv}\\

65&def&\nolinkurl{PaperXI.MoritaEndomorphismCoordinates111.coordinateUnitMap}\\

71&theorem&\nolinkurl{PaperXI.MoritaEndomorphismCoordinates111.blockRingEquiv_centralizerEmbedding}\\

82&theorem&\nolinkurl{PaperXI.MoritaEndomorphismCoordinates111.coordinateAut_coordinateUnitMap}\\

96&theorem&\nolinkurl{PaperXI.MoritaEndomorphismCoordinates111.coordinateUnitMap_centralizer}\\

115&theorem&\nolinkurl{PaperXI.MoritaEndomorphismCoordinates111.coordinateRingEquiv_unitConjugation}\\

128&theorem&\nolinkurl{PaperXI.MoritaEndomorphismCoordinates111.coordinate_map_conjugate}\\

142&theorem&\nolinkurl{PaperXI.MoritaEndomorphismCoordinates111.coordinate_intersection_relIndex}\\

164&theorem&\nolinkurl{PaperXI.MoritaEndomorphismCoordinates111.coordinateUnitMap_integral}\\

\bottomrule\end{longtable}

\section{\texorpdfstring{\nolinkurl{PaperXI.MoritaIndex110}}{PaperXI.MoritaIndex110}}\label{module:PaperXI.MoritaIndex110}

\textbf{Selected layer:} Continuation110.\quad\textbf{Source lines:} 144.\par\textbf{SHA-256:} {\scriptsize\nolinkurl{30251babfb354b36cc6f28c150e892b0e244e9046f24d606dd8921245105a9d2}}.

\VerbatimInput[fontsize=\scriptsize,breaklines=true,breakanywhere=true]{excerpts/D_MoritaIndex110.txt}

\textbf{Direct project imports:} \nolinkurl{PaperXI.MoritaCornerRank110}, \nolinkurl{PaperXI.MoritaQuotient110}, \nolinkurl{PaperXI.FreeLatticeCard110}, \nolinkurl{PaperXI.MoritaUniformizer110}.

\begin{longtable}{@{}p{0.07\textwidth}p{0.12\textwidth}p{0.75\textwidth}@{}}\toprule Line&Kind&Name\\\midrule\endhead

20&theorem&\nolinkurl{PaperXI.MoritaIndex110.leftPower_addSubgroupOf_eq}\\

44&theorem&\nolinkurl{PaperXI.MoritaIndex110.card_leftPower_quotient}\\

69&def&\nolinkurl{PaperXI.MoritaIndex110.coefficientUnit}\\

73&theorem&\nolinkurl{PaperXI.MoritaIndex110.coefficientUnit_coe}\\

78&theorem&\nolinkurl{PaperXI.MoritaIndex110.card_diagonal_intersection}\\

122&theorem&\nolinkurl{PaperXI.MoritaIndex110.diagonal_relIndex_of_radical_commutation}\\

\bottomrule\end{longtable}

\section{\texorpdfstring{\nolinkurl{PaperXI.MoritaLattice110}}{PaperXI.MoritaLattice110}}\label{module:PaperXI.MoritaLattice110}

\textbf{Selected layer:} Continuation110.\quad\textbf{Source lines:} 354.\par\textbf{SHA-256:} {\scriptsize\nolinkurl{abdbc4e8506fb153a12c66728668a5c6ff97d94d2f2655dd1ebf49d43f4008d0}}.

\VerbatimInput[fontsize=\scriptsize,breaklines=true,breakanywhere=true]{excerpts/D_MoritaLattice110.txt}

\textbf{Direct project imports:} none.

\begin{longtable}{@{}p{0.07\textwidth}p{0.12\textwidth}p{0.75\textwidth}@{}}\toprule Line&Kind&Name\\\midrule\endhead

18&def&\nolinkurl{PaperXI.MoritaLattice110.IsMatrixBimodule}\\

23&def&\nolinkurl{PaperXI.MoritaLattice110.cornerEmbedding}\\

29&def&\nolinkurl{PaperXI.MoritaLattice110.corner}\\

33&theorem&\nolinkurl{PaperXI.MoritaLattice110.mem_corner}\\

37&theorem&\nolinkurl{PaperXI.MoritaLattice110.entry_mem_corner}\\

45&theorem&\nolinkurl{PaperXI.MoritaLattice110.single_mem_of_corner}\\

53&theorem&\nolinkurl{PaperXI.MoritaLattice110.eq_matrix_corner}\\

68&theorem&\nolinkurl{PaperXI.MoritaLattice110.corner_bimodule}\\

82&def&\nolinkurl{PaperXI.MoritaLattice110.unitConjugation}\\

89&theorem&\nolinkurl{PaperXI.MoritaLattice110.unitConjugation_apply}\\

93&def&\nolinkurl{PaperXI.MoritaLattice110.normalizingUnits}\\

106&theorem&\nolinkurl{PaperXI.MoritaLattice110.mem_normalizingUnits}\\

110&theorem&\nolinkurl{PaperXI.MoritaLattice110.normalizes_of_left_right_map_eq}\\

137&theorem&\nolinkurl{PaperXI.MoritaLattice110.zpow_normalizes}\\

142&theorem&\nolinkurl{PaperXI.MoritaLattice110.two_sided_mem_iff_left}\\

159&def&\nolinkurl{PaperXI.MoritaLattice110.leftPower}\\

162&theorem&\nolinkurl{PaperXI.MoritaLattice110.mem_leftPower}\\

177&def&\nolinkurl{PaperXI.MoritaLattice110.diagonalUnit}\\

194&theorem&\nolinkurl{PaperXI.MoritaLattice110.diagonal_conjugation_apply}\\

200&theorem&\nolinkurl{PaperXI.MoritaLattice110.scalar_conjugation_single}\\

208&theorem&\nolinkurl{PaperXI.MoritaLattice110.corner_normalizes_of_scalar_commutation}\\

221&theorem&\nolinkurl{PaperXI.MoritaLattice110.mem_conjugate}\\

232&def&\nolinkurl{PaperXI.MoritaLattice110.entrywise}\\

239&theorem&\nolinkurl{PaperXI.MoritaLattice110.mem_entrywise}\\

244&theorem&\nolinkurl{PaperXI.MoritaLattice110.conjugate_eq_entrywise_leftPower}\\

276&theorem&\nolinkurl{PaperXI.MoritaLattice110.intersection_eq_entrywise}\\

297&theorem&\nolinkurl{PaperXI.MoritaLattice110.pow_mul_mem}\\

304&theorem&\nolinkurl{PaperXI.MoritaLattice110.leftPower_zero}\\

310&theorem&\nolinkurl{PaperXI.MoritaLattice110.inf_leftPower_eq_toNat}\\

331&theorem&\nolinkurl{PaperXI.MoritaLattice110.corner_scalar_stable}\\

341&theorem&\nolinkurl{PaperXI.MoritaLattice110.intersection_eq_nonnegative_entrywise}\\

\bottomrule\end{longtable}

\section{\texorpdfstring{\nolinkurl{PaperXI.MoritaPairIndex110}}{PaperXI.MoritaPairIndex110}}\label{module:PaperXI.MoritaPairIndex110}

\textbf{Selected layer:} Continuation110.\quad\textbf{Source lines:} 39.\par\textbf{SHA-256:} {\scriptsize\nolinkurl{1a7eeb3145c18b69770ec5d2b83acd8a4b2d23bb966e4cc47a44655a47472727}}.

\VerbatimInput[fontsize=\scriptsize,breaklines=true,breakanywhere=true]{excerpts/D_MoritaPairIndex110.txt}

\textbf{Direct project imports:} \nolinkurl{PaperXI.MoritaActualCartan110}.

\begin{longtable}{@{}p{0.07\textwidth}p{0.12\textwidth}p{0.75\textwidth}@{}}\toprule Line&Kind&Name\\\midrule\endhead

20&theorem&\nolinkurl{PaperXI.MoritaPairIndex110.actual_pair_index}\\

\bottomrule\end{longtable}

\section{\texorpdfstring{\nolinkurl{PaperXI.MoritaQuotient110}}{PaperXI.MoritaQuotient110}}\label{module:PaperXI.MoritaQuotient110}

\textbf{Selected layer:} Continuation110.\quad\textbf{Source lines:} 170.\par\textbf{SHA-256:} {\scriptsize\nolinkurl{33d3a455e6a8a202c2ce041e07c1c48e5d82d632a4f4dd114458be95809369e9}}.

\VerbatimInput[fontsize=\scriptsize,breaklines=true,breakanywhere=true]{excerpts/D_MoritaQuotient110.txt}

\textbf{Direct project imports:} \nolinkurl{PaperXI.MoritaLattice110}.

\begin{longtable}{@{}p{0.07\textwidth}p{0.12\textwidth}p{0.75\textwidth}@{}}\toprule Line&Kind&Name\\\midrule\endhead

16&def&\nolinkurl{PaperXI.MoritaQuotient110.coordinateQuotientMap}\\

22&theorem&\nolinkurl{PaperXI.MoritaQuotient110.coordinateQuotientMap_apply}\\

27&theorem&\nolinkurl{PaperXI.MoritaQuotient110.coordinateQuotientMap_surjective}\\

35&theorem&\nolinkurl{PaperXI.MoritaQuotient110.coordinateQuotientMap_ker}\\

43&def&\nolinkurl{PaperXI.MoritaQuotient110.piQuotientEquiv}\\

48&theorem&\nolinkurl{PaperXI.MoritaQuotient110.piQuotientEquiv_mk}\\

54&theorem&\nolinkurl{PaperXI.MoritaQuotient110.card_pi_quotient}\\

60&def&\nolinkurl{PaperXI.MoritaQuotient110.subgroupCoordinateQuotientMap}\\

66&theorem&\nolinkurl{PaperXI.MoritaQuotient110.subgroupCoordinateQuotientMap_apply}\\

72&theorem&\nolinkurl{PaperXI.MoritaQuotient110.subgroupCoordinateQuotientMap_surjective}\\

82&theorem&\nolinkurl{PaperXI.MoritaQuotient110.subgroupCoordinateQuotientMap_ker}\\

93&def&\nolinkurl{PaperXI.MoritaQuotient110.subgroupPiQuotientEquiv}\\

102&theorem&\nolinkurl{PaperXI.MoritaQuotient110.subgroupPiQuotientEquiv_mk}\\

110&theorem&\nolinkurl{PaperXI.MoritaQuotient110.card_subgroup_pi_quotient}\\

121&def&\nolinkurl{PaperXI.MoritaQuotient110.matrixCoordinateQuotientMap}\\

127&theorem&\nolinkurl{PaperXI.MoritaQuotient110.matrixCoordinateQuotientMap_apply}\\

132&theorem&\nolinkurl{PaperXI.MoritaQuotient110.matrixCoordinateQuotientMap_surjective}\\

141&theorem&\nolinkurl{PaperXI.MoritaQuotient110.matrixCoordinateQuotientMap_ker}\\

149&def&\nolinkurl{PaperXI.MoritaQuotient110.matrixQuotientEquiv}\\

155&theorem&\nolinkurl{PaperXI.MoritaQuotient110.matrixQuotientEquiv_mk}\\

162&theorem&\nolinkurl{PaperXI.MoritaQuotient110.card_matrix_quotient}\\

\bottomrule\end{longtable}

\section{\texorpdfstring{\nolinkurl{PaperXI.MoritaRadicalDictionary110}}{PaperXI.MoritaRadicalDictionary110}}\label{module:PaperXI.MoritaRadicalDictionary110}

\textbf{Selected layer:} Continuation110.\quad\textbf{Source lines:} 132.\par\textbf{SHA-256:} {\scriptsize\nolinkurl{caf45f046cafea555ed4871a05c2baece3d153318b1fced87d701d20622f781f}}.

\VerbatimInput[fontsize=\scriptsize,breaklines=true,breakanywhere=true]{excerpts/D_MoritaRadicalDictionary110.txt}

\textbf{Direct project imports:} \nolinkurl{PaperXI.MoritaActualCartan110}.

\begin{longtable}{@{}p{0.07\textwidth}p{0.12\textwidth}p{0.75\textwidth}@{}}\toprule Line&Kind&Name\\\midrule\endhead

17&def&\nolinkurl{PaperXI.MoritaRadicalDictionary110.additiveProduct}\\

20&theorem&\nolinkurl{PaperXI.MoritaRadicalDictionary110.mul_mem_additiveProduct}\\

30&def&\nolinkurl{PaperXI.MoritaRadicalDictionary110.maximalOrderRadical}\\

34&theorem&\nolinkurl{PaperXI.MoritaRadicalDictionary110.scalar_mem_maximalOrderRadical}\\

42&theorem&\nolinkurl{PaperXI.MoritaRadicalDictionary110.scalar_commutes_mapped}\\

49&theorem&\nolinkurl{PaperXI.MoritaRadicalDictionary110.radical_product_eq_left_image}\\

69&theorem&\nolinkurl{PaperXI.MoritaRadicalDictionary110.product_radical_eq_right_image}\\

90&theorem&\nolinkurl{PaperXI.MoritaRadicalDictionary110.scalar_images_eq_of_radical_products_eq}\\

109&def&\nolinkurl{PaperXI.MoritaRadicalDictionary110.actualRadical}\\

115&theorem&\nolinkurl{PaperXI.MoritaRadicalDictionary110.RadicalCommutes_of_actual_radical_products}\\

\bottomrule\end{longtable}

\section{\texorpdfstring{\nolinkurl{PaperXI.MoritaScalarUpgrade110}}{PaperXI.MoritaScalarUpgrade110}}\label{module:PaperXI.MoritaScalarUpgrade110}

\textbf{Selected layer:} Continuation110.\quad\textbf{Source lines:} 55.\par\textbf{SHA-256:} {\scriptsize\nolinkurl{0a05325327c85f7a0de2f900d68524b1a168fd2264a2b5d8565ba75c64356291}}.

\VerbatimInput[fontsize=\scriptsize,breaklines=true,breakanywhere=true]{excerpts/D_MoritaScalarUpgrade110.txt}

\textbf{Direct project imports:} \nolinkurl{PaperXI.MoritaLattice110}.

\begin{longtable}{@{}p{0.07\textwidth}p{0.12\textwidth}p{0.75\textwidth}@{}}\toprule Line&Kind&Name\\\midrule\endhead

17&def&\nolinkurl{PaperXI.MoritaScalarUpgrade110.upgrade}\\

24&theorem&\nolinkurl{PaperXI.MoritaScalarUpgrade110.upgrade_restrictScalars}\\

30&theorem&\nolinkurl{PaperXI.MoritaScalarUpgrade110.upgrade_fg}\\

43&theorem&\nolinkurl{PaperXI.MoritaScalarUpgrade110.scalar_stable_of_matrix_bimodule}\\

\bottomrule\end{longtable}

\section{\texorpdfstring{\nolinkurl{PaperXI.MoritaUniformizer110}}{PaperXI.MoritaUniformizer110}}\label{module:PaperXI.MoritaUniformizer110}

\textbf{Selected layer:} Continuation110.\quad\textbf{Source lines:} 62.\par\textbf{SHA-256:} {\scriptsize\nolinkurl{b8f0b8cf79c2e58d57e03edd7ec29f81f51c91f1a0a837045514ee4ccf1c126e}}.

\VerbatimInput[fontsize=\scriptsize,breaklines=true,breakanywhere=true]{excerpts/D_MoritaUniformizer110.txt}

\textbf{Direct project imports:} none.

\begin{longtable}{@{}p{0.07\textwidth}p{0.12\textwidth}p{0.75\textwidth}@{}}\toprule Line&Kind&Name\\\midrule\endhead

17&def&\nolinkurl{PaperXI.MoritaUniformizer110.uniformizerIdeal}\\

20&theorem&\nolinkurl{PaperXI.MoritaUniformizer110.exists_uniformizer}\\

24&theorem&\nolinkurl{PaperXI.MoritaUniformizer110.uniformizerIdeal_eq_maximalIdeal}\\

27&theorem&\nolinkurl{PaperXI.MoritaUniformizer110.uniformizerIdeal_ne_bot}\\

32&theorem&\nolinkurl{PaperXI.MoritaUniformizer110.uniformizerIdeal_isMaximal}\\

38&theorem&\nolinkurl{PaperXI.MoritaUniformizer110.uniformizerIdeal_isPrime}\\

43&def&\nolinkurl{PaperXI.MoritaUniformizer110.uniformizerResidueEquiv}\\

47&theorem&\nolinkurl{PaperXI.MoritaUniformizer110.uniformizerResidueEquiv_mk}\\

53&theorem&\nolinkurl{PaperXI.MoritaUniformizer110.uniformizer_quotient_finite}\\

58&theorem&\nolinkurl{PaperXI.MoritaUniformizer110.card_uniformizer_quotient}\\

\bottomrule\end{longtable}

\section{\texorpdfstring{\nolinkurl{PaperXI.MovingPhase}}{PaperXI.MovingPhase}}\label{module:PaperXI.MovingPhase}

\textbf{Selected layer:} v0.44 parent.\quad\textbf{Source lines:} 87.\par\textbf{SHA-256:} {\scriptsize\nolinkurl{c01ee25e1c3e1f60b027ae5a5f09f85341892753a59e2eae1257490b80f8ec52}}.

\VerbatimInput[fontsize=\scriptsize,breaklines=true,breakanywhere=true]{excerpts/D_MovingPhase.txt}

\textbf{Direct project imports:} \nolinkurl{PaperXI.FixedMovingSource}, \nolinkurl{PaperXI.OddPhaseTransport}.

\begin{longtable}{@{}p{0.07\textwidth}p{0.12\textwidth}p{0.75\textwidth}@{}}\toprule Line&Kind&Name\\\midrule\endhead

16&def&\nolinkurl{PaperXI.MovingPhase.phase}\\

18&theorem&\nolinkurl{PaperXI.MovingPhase.phase_formula}\\

23&theorem&\nolinkurl{PaperXI.MovingPhase.determinant_ne_zero}\\

28&theorem&\nolinkurl{PaperXI.MovingPhase.phase_sq_one}\\

31&theorem&\nolinkurl{PaperXI.MovingPhase.phase_sign}\\

37&theorem&\nolinkurl{PaperXI.MovingPhase.source_valid}\\

51&theorem&\nolinkurl{PaperXI.MovingPhase.conjugacy}\\

58&theorem&\nolinkurl{PaperXI.MovingPhase.square_class}\\

66&theorem&\nolinkurl{PaperXI.MovingPhase.source_square_class}\\

80&theorem&\nolinkurl{PaperXI.MovingPhase.phase_identity}\\

\bottomrule\end{longtable}

\section{\texorpdfstring{\nolinkurl{PaperXI.NormalInertia}}{PaperXI.NormalInertia}}\label{module:PaperXI.NormalInertia}

\textbf{Selected layer:} v0.44 parent.\quad\textbf{Source lines:} 85.\par\textbf{SHA-256:} {\scriptsize\nolinkurl{bc52e7a907e034f8583c73eb69d3de4b3f0575f0ce1174605d50cce3864ab80a}}.

\VerbatimInput[fontsize=\scriptsize,breaklines=true,breakanywhere=true]{excerpts/D_NormalInertia.txt}

\textbf{Direct project imports:} \nolinkurl{PaperXI.QuotientRepresentation}.

\begin{longtable}{@{}p{0.07\textwidth}p{0.12\textwidth}p{0.75\textwidth}@{}}\toprule Line&Kind&Name\\\midrule\endhead

16&def&\nolinkurl{PaperXI.NormalInertia.fixedSubrepresentation}\\

22&theorem&\nolinkurl{PaperXI.NormalInertia.fixed_eq_bot_or_top}\\

30&theorem&\nolinkurl{PaperXI.NormalInertia.fixed_top_iff_le_kernel}\\

46&theorem&\nolinkurl{PaperXI.NormalInertia.dimension_one_of_inertia_trivial}\\

57&theorem&\nolinkurl{PaperXI.NormalInertia.fixed_eq_bot_of_dimension_gt_one}\\

68&theorem&\nolinkurl{PaperXI.NormalInertia.fixed_eq_bot_of_not_le_kernel}\\

76&theorem&\nolinkurl{PaperXI.NormalInertia.fixed_codimension}\\

\bottomrule\end{longtable}

\section{\texorpdfstring{\nolinkurl{PaperXI.NormalizedFrobeniusQuotient}}{PaperXI.NormalizedFrobeniusQuotient}}\label{module:PaperXI.NormalizedFrobeniusQuotient}

\textbf{Selected layer:} normalized Frobenius.\quad\textbf{Source lines:} 120.\par\textbf{SHA-256:} {\scriptsize\nolinkurl{917cfa62d1b914d7d47a3c31f682a52d3338dc040a1606450f1cff71ca075b1f}}.

\VerbatimInput[fontsize=\scriptsize,breaklines=true,breakanywhere=true]{excerpts/D_NormalizedFrobeniusQuotient.txt}

\textbf{Direct project imports:} \nolinkurl{PaperXI.SpectralGeometricFrobenius}, \nolinkurl{PaperXI.RamifiedUnramifiedQuotient}.

\begin{longtable}{@{}p{0.07\textwidth}p{0.12\textwidth}p{0.75\textwidth}@{}}\toprule Line&Kind&Name\\\midrule\endhead

22&abbrev&\nolinkurl{PaperXI.NormalizedFrobeniusQuotient.Decomposition}\\

23&abbrev&\nolinkurl{PaperXI.NormalizedFrobeniusQuotient.residueHom}\\

25&abbrev&\nolinkurl{PaperXI.NormalizedFrobeniusQuotient.Weil}\\

26&abbrev&\nolinkurl{PaperXI.NormalizedFrobeniusQuotient.degree}\\

28&abbrev&\nolinkurl{PaperXI.NormalizedFrobeniusQuotient.Inertia}\\

31&def&\nolinkurl{PaperXI.NormalizedFrobeniusQuotient.inertiaEquiv}\\

37&def&\nolinkurl{PaperXI.NormalizedFrobeniusQuotient.decompositionLift}\\

40&theorem&\nolinkurl{PaperXI.NormalizedFrobeniusQuotient.residue_decompositionLift}\\

45&theorem&\nolinkurl{PaperXI.NormalizedFrobeniusQuotient.frobenius_mem_range}\\

49&def&\nolinkurl{PaperXI.NormalizedFrobeniusQuotient.weilGeometricLift}\\

55&theorem&\nolinkurl{PaperXI.NormalizedFrobeniusQuotient.degree_geometricLift}\\

59&theorem&\nolinkurl{PaperXI.NormalizedFrobeniusQuotient.lift_inverse_power}\\

65&def&\nolinkurl{PaperXI.NormalizedFrobeniusQuotient.weilArithmeticLift}\\

68&theorem&\nolinkurl{PaperXI.NormalizedFrobeniusQuotient.degree_arithmeticLift}\\

74&theorem&\nolinkurl{PaperXI.NormalizedFrobeniusQuotient.degree_surjective}\\

79&def&\nolinkurl{PaperXI.NormalizedFrobeniusQuotient.quotientEquivIntegerDegree}\\

84&theorem&\nolinkurl{PaperXI.NormalizedFrobeniusQuotient.quotientEquivIntegerDegree_mk}\\

90&theorem&\nolinkurl{PaperXI.NormalizedFrobeniusQuotient.quotient_geometricLift}\\

102&theorem&\nolinkurl{PaperXI.NormalizedFrobeniusQuotient.degree_one_action_eq_on_invariants}\\

\bottomrule\end{longtable}

\section{\texorpdfstring{\nolinkurl{PaperXI.NormalizerDictionary110}}{PaperXI.NormalizerDictionary110}}\label{module:PaperXI.NormalizerDictionary110}

\textbf{Selected layer:} Continuation110.\quad\textbf{Source lines:} 432.\par\textbf{SHA-256:} {\scriptsize\nolinkurl{19839adda13ba916ac5dbe60982bf13e65ce22a746ae2b2a63ca3cc54e4fe9e4}}.

\VerbatimInput[fontsize=\scriptsize,breaklines=true,breakanywhere=true]{excerpts/D_NormalizerDictionary110.txt}

\textbf{Direct project imports:} \nolinkurl{PaperXI.ActualCharacterIdentification61}, \nolinkurl{PaperXI.ActualCharacterResidue61}, \nolinkurl{PaperXI.ActualCuspidalCandidate105}.

\begin{longtable}{@{}p{0.07\textwidth}p{0.12\textwidth}p{0.75\textwidth}@{}}\toprule Line&Kind&Name\\\midrule\endhead

16&theorem&\nolinkurl{PaperXI.NormalizerDictionary110.inf_eq_of_support}\\

37&theorem&\nolinkurl{PaperXI.NormalizerDictionary110.compact_eq_sup_of_support}\\

61&theorem&\nolinkurl{PaperXI.NormalizerDictionary110.compact_le_normalizer_inf}\\

66&theorem&\nolinkurl{PaperXI.NormalizerDictionary110.compact_le_normalizer_of_inf_eq}\\

73&theorem&\nolinkurl{PaperXI.NormalizerDictionary110.mem_compact_iff_product}\\

123&theorem&\nolinkurl{PaperXI.NormalizerDictionary110.chainGroup_le_normalizer_congruence}\\

136&theorem&\nolinkurl{PaperXI.NormalizerDictionary110.firstSubgroup_ambient_image}\\

148&theorem&\nolinkurl{PaperXI.NormalizerDictionary110.firstSubgroup_embedding_range}\\

157&theorem&\nolinkurl{PaperXI.NormalizerDictionary110.firstSubgroup_normal}\\

169&def&\nolinkurl{PaperXI.NormalizerDictionary110.integralCentralizer}\\

174&theorem&\nolinkurl{PaperXI.NormalizerDictionary110.matrixType_eq_integral_sup}\\

182&theorem&\nolinkurl{PaperXI.NormalizerDictionary110.integralCentralizer_le_chainGroup}\\

193&theorem&\nolinkurl{PaperXI.NormalizerDictionary110.integralCentralizer_le_character_normalizer}\\

209&theorem&\nolinkurl{PaperXI.NormalizerDictionary110.restrict_mem_chain_iff}\\

234&theorem&\nolinkurl{PaperXI.NormalizerDictionary110.centralizer_inf_chain}\\

271&theorem&\nolinkurl{PaperXI.NormalizerDictionary110.centralizer_inf_congruence}\\

318&theorem&\nolinkurl{PaperXI.NormalizerDictionary110.actual_first_eq}\\

329&theorem&\nolinkurl{PaperXI.NormalizerDictionary110.actual_first_subgroup_eq}\\

340&theorem&\nolinkurl{PaperXI.NormalizerDictionary110.actual_compact_eq_sup}\\

351&theorem&\nolinkurl{PaperXI.NormalizerDictionary110.actual_compact_product}\\

367&theorem&\nolinkurl{PaperXI.NormalizerDictionary110.actual_model_le}\\

378&theorem&\nolinkurl{PaperXI.NormalizerDictionary110.actual_model_from_centralizer}\\

395&theorem&\nolinkurl{PaperXI.NormalizerDictionary110.actual_first_from_residue}\\

411&theorem&\nolinkurl{PaperXI.NormalizerDictionary110.character_support_compact}\\

\bottomrule\end{longtable}

\section{\texorpdfstring{\nolinkurl{PaperXI.NormalizerExtensionInducing58}}{PaperXI.NormalizerExtensionInducing58}}\label{module:PaperXI.NormalizerExtensionInducing58}

\textbf{Selected layer:} base108 arithmetic branch.\quad\textbf{Source lines:} 272.\par\textbf{SHA-256:} {\scriptsize\nolinkurl{e30d543d65b9008eaf4b66affdcd66e745a09db8651881df7678fda5e84c4b79}}.

\VerbatimInput[fontsize=\scriptsize,breaklines=true,breakanywhere=true]{excerpts/D_NormalizerExtensionInducing58.txt}

\textbf{Direct project imports:} \nolinkurl{PaperXI.CentralExtensionInducing}, \nolinkurl{PaperXI.QuotientRepresentation}.

\begin{longtable}{@{}p{0.07\textwidth}p{0.12\textwidth}p{0.75\textwidth}@{}}\toprule Line&Kind&Name\\\midrule\endhead

22&def&\nolinkurl{PaperXI.NormalizerExtensionInducing58.compatiblePairs}\\

40&theorem&\nolinkurl{PaperXI.NormalizerExtensionInducing58.covariance_zpow}\\

53&def&\nolinkurl{PaperXI.NormalizerExtensionInducing58.conjugation}\\

56&def&\nolinkurl{PaperXI.NormalizerExtensionInducing58.conjugationAction}\\

59&theorem&\nolinkurl{PaperXI.NormalizerExtensionInducing58.conjugation_coe}\\

62&theorem&\nolinkurl{PaperXI.NormalizerExtensionInducing58.conjugationAction_coe}\\

69&def&\nolinkurl{PaperXI.NormalizerExtensionInducing58.normalFormHom}\\

75&theorem&\nolinkurl{PaperXI.NormalizerExtensionInducing58.normalFormHom_apply}\\

80&def&\nolinkurl{PaperXI.NormalizerExtensionInducing58.extendedSubgroup}\\

82&theorem&\nolinkurl{PaperXI.NormalizerExtensionInducing58.le_extendedSubgroup}\\

86&theorem&\nolinkurl{PaperXI.NormalizerExtensionInducing58.generator_mem}\\

89&theorem&\nolinkurl{PaperXI.NormalizerExtensionInducing58.extendedSubgroup_eq_sup}\\

97&theorem&\nolinkurl{PaperXI.NormalizerExtensionInducing58.normalFormHom_injective}\\

112&def&\nolinkurl{PaperXI.NormalizerExtensionInducing58.normalFormEquiv}\\

115&theorem&\nolinkurl{PaperXI.NormalizerExtensionInducing58.normalFormEquiv_coe}\\

122&def&\nolinkurl{PaperXI.NormalizerExtensionInducing58.subgroupInclusion}\\

126&def&\nolinkurl{PaperXI.NormalizerExtensionInducing58.generator}\\

133&def&\nolinkurl{PaperXI.NormalizerExtensionInducing58.semidirectRepresentation}\\

143&def&\nolinkurl{PaperXI.NormalizerExtensionInducing58.extendedRepresentation}\\

147&theorem&\nolinkurl{PaperXI.NormalizerExtensionInducing58.extendedRepresentation_normalForm}\\

155&theorem&\nolinkurl{PaperXI.NormalizerExtensionInducing58.extendedRepresentation_restrict}\\

163&theorem&\nolinkurl{PaperXI.NormalizerExtensionInducing58.extendedRepresentation_generator}\\

172&theorem&\nolinkurl{PaperXI.NormalizerExtensionInducing58.extendedRepresentation_irreducible}\\

183&theorem&\nolinkurl{PaperXI.NormalizerExtensionInducing58.equiv_unit_covariance}\\

200&theorem&\nolinkurl{PaperXI.NormalizerExtensionInducing58.intertwiner_bijective}\\

209&def&\nolinkurl{PaperXI.NormalizerExtensionInducing58.extensionOfEquiv}\\

216&theorem&\nolinkurl{PaperXI.NormalizerExtensionInducing58.extensionOfEquiv_restrict}\\

222&theorem&\nolinkurl{PaperXI.NormalizerExtensionInducing58.extensionOfEquiv_generator}\\

231&theorem&\nolinkurl{PaperXI.NormalizerExtensionInducing58.exists_extension_of_nonzero}\\

249&theorem&\nolinkurl{PaperXI.NormalizerExtensionInducing58.extendedSubgroup_isOpen}\\

254&theorem&\nolinkurl{PaperXI.NormalizerExtensionInducing58.extendedRepresentation_smooth}\\

\bottomrule\end{longtable}

\section{\texorpdfstring{\nolinkurl{PaperXI.NormalizerIntertwining60}}{PaperXI.NormalizerIntertwining60}}\label{module:PaperXI.NormalizerIntertwining60}

\textbf{Selected layer:} base108 arithmetic branch.\quad\textbf{Source lines:} 111.\par\textbf{SHA-256:} {\scriptsize\nolinkurl{5993f730521e96961b4fed40e9a77796d990e2a0cccd720874c20680873c2c7f}}.

\VerbatimInput[fontsize=\scriptsize,breaklines=true,breakanywhere=true]{excerpts/D_NormalizerIntertwining60.txt}

\textbf{Direct project imports:} \nolinkurl{PaperXI.InductionIrreducibility}, \nolinkurl{PaperXI.NormalizerExtensionInducing58}, \nolinkurl{PaperXI.MonomialInduction}.

\begin{longtable}{@{}p{0.07\textwidth}p{0.12\textwidth}p{0.75\textwidth}@{}}\toprule Line&Kind&Name\\\midrule\endhead

23&def&\nolinkurl{PaperXI.NormalizerIntertwining60.IntertwinedAt}\\

31&def&\nolinkurl{PaperXI.NormalizerIntertwining60.intersectionInclusion}\\

37&def&\nolinkurl{PaperXI.NormalizerIntertwining60.normalizerIntertwiner}\\

44&theorem&\nolinkurl{PaperXI.NormalizerIntertwining60.normalizerIntertwiner_apply}\\

49&theorem&\nolinkurl{PaperXI.NormalizerIntertwining60.normalizerIntertwiner_ne_zero}\\

60&theorem&\nolinkurl{PaperXI.NormalizerIntertwining60.exists_normalizerIntertwiner}\\

67&theorem&\nolinkurl{PaperXI.NormalizerIntertwining60.exists_normalizerEquiv}\\

76&theorem&\nolinkurl{PaperXI.NormalizerIntertwining60.character_intertwinedAt}\\

101&theorem&\nolinkurl{PaperXI.NormalizerIntertwining60.exists_normalizerIntertwiner_of_character}\\

\bottomrule\end{longtable}

\section{\texorpdfstring{\nolinkurl{PaperXI.OddEnvelopeHeisenberg}}{PaperXI.OddEnvelopeHeisenberg}}\label{module:PaperXI.OddEnvelopeHeisenberg}

\textbf{Selected layer:} v0.44 parent.\quad\textbf{Source lines:} 52.\par\textbf{SHA-256:} {\scriptsize\nolinkurl{0ff24e111c1739d834a0e471bac34b82c9da25eb040e12f9cea2f9cf3b1d1465}}.

\VerbatimInput[fontsize=\scriptsize,breaklines=true,breakanywhere=true]{excerpts/D_OddEnvelopeHeisenberg.txt}

\textbf{Direct project imports:} \nolinkurl{PaperXI.BilinearHeisenberg}, \nolinkurl{PaperXI.HeisenbergPrime}.

\begin{longtable}{@{}p{0.07\textwidth}p{0.12\textwidth}p{0.75\textwidth}@{}}\toprule Line&Kind&Name\\\midrule\endhead

13&def&\nolinkurl{PaperXI.OddEnvelopeHeisenberg.factor}\\

15&theorem&\nolinkurl{PaperXI.OddEnvelopeHeisenberg.factor_apply}\\

19&theorem&\nolinkurl{PaperXI.OddEnvelopeHeisenberg.factor_commutator}\\

33&def&\nolinkurl{PaperXI.OddEnvelopeHeisenberg.isometryEquiv}\\

46&theorem&\nolinkurl{PaperXI.OddEnvelopeHeisenberg.isometryEquiv_center}\\

\bottomrule\end{longtable}

\section{\texorpdfstring{\nolinkurl{PaperXI.OddPhase}}{PaperXI.OddPhase}}\label{module:PaperXI.OddPhase}

\textbf{Selected layer:} v0.44 parent.\quad\textbf{Source lines:} 89.\par\textbf{SHA-256:} {\scriptsize\nolinkurl{a78d2e7c2db5c75599339479441e4e0b95181ec11ced910642a239377291249c}}.

\VerbatimInput[fontsize=\scriptsize,breaklines=true,breakanywhere=true]{excerpts/D_OddPhase.txt}

\textbf{Direct project imports:} none.

\begin{longtable}{@{}p{0.07\textwidth}p{0.12\textwidth}p{0.75\textwidth}@{}}\toprule Line&Kind&Name\\\midrule\endhead

15&def&\nolinkurl{PaperXI.OddPhase.difference}\\

17&theorem&\nolinkurl{PaperXI.OddPhase.difference_apply}\\

21&def&\nolinkurl{PaperXI.OddPhase.FixedPointFree}\\

23&theorem&\nolinkurl{PaperXI.OddPhase.fixedPointFree_iff_ker}\\

28&theorem&\nolinkurl{PaperXI.OddPhase.fixedPointFree_iff_injective}\\

35&theorem&\nolinkurl{PaperXI.OddPhase.fixedPointFree_iff_bijective}\\

40&theorem&\nolinkurl{PaperXI.OddPhase.fixedPointFree_iff_det_ne_zero}\\

45&def&\nolinkurl{PaperXI.OddPhase.differenceEquiv}\\

48&theorem&\nolinkurl{PaperXI.OddPhase.differenceEquiv_apply}\\

52&def&\nolinkurl{PaperXI.OddPhase.phase}\\

56&theorem&\nolinkurl{PaperXI.OddPhase.phase_sq_one}\\

61&theorem&\nolinkurl{PaperXI.OddPhase.phase_sign}\\

64&theorem&\nolinkurl{PaperXI.OddPhase.phase_ne_zero}\\

69&theorem&\nolinkurl{PaperXI.OddPhase.phase_zero_space}\\

74&theorem&\nolinkurl{PaperXI.OddPhase.fixedPointFree_zero_space}\\

80&theorem&\nolinkurl{PaperXI.OddPhase.phase_eq_zero_iff}\\

\bottomrule\end{longtable}

\section{\texorpdfstring{\nolinkurl{PaperXI.OddPhaseCertificate}}{PaperXI.OddPhaseCertificate}}\label{module:PaperXI.OddPhaseCertificate}

\textbf{Selected layer:} v0.44 parent.\quad\textbf{Source lines:} 61.\par\textbf{SHA-256:} {\scriptsize\nolinkurl{783d84147c4f09cbca432ad799b1e57d9bd47da13f3841e92d4f071d49709f2e}}.

\VerbatimInput[fontsize=\scriptsize,breaklines=true,breakanywhere=true]{excerpts/D_OddPhaseCertificate.txt}

\textbf{Direct project imports:} \nolinkurl{PaperXI.OddPhase}, \nolinkurl{PaperXI.OddSymplectic}.

\begin{longtable}{@{}p{0.07\textwidth}p{0.12\textwidth}p{0.75\textwidth}@{}}\toprule Line&Kind&Name\\\midrule\endhead

16&theorem&\nolinkurl{PaperXI.OddPhase.charpoly_eval_one}\\

24&theorem&\nolinkurl{PaperXI.OddPhase.determinant_neg_difference}\\

30&theorem&\nolinkurl{PaperXI.OddPhase.characteristic_determinant}\\

36&theorem&\nolinkurl{PaperXI.OddPhase.characteristic_certificate_of_even}\\

46&theorem&\nolinkurl{PaperXI.OddPhase.characteristic_certificate}\\

55&theorem&\nolinkurl{PaperXI.OddPhase.characteristic_eval_one_ne_zero}\\

\bottomrule\end{longtable}

\section{\texorpdfstring{\nolinkurl{PaperXI.OddPhaseSum}}{PaperXI.OddPhaseSum}}\label{module:PaperXI.OddPhaseSum}

\textbf{Selected layer:} v0.44 parent.\quad\textbf{Source lines:} 92.\par\textbf{SHA-256:} {\scriptsize\nolinkurl{d97a3ea60f4bcb92aedcc1ccd4ba6b44c70ccfdf1ab43161b9c4194ad26867d1}}.

\VerbatimInput[fontsize=\scriptsize,breaklines=true,breakanywhere=true]{excerpts/D_OddPhaseSum.txt}

\textbf{Direct project imports:} \nolinkurl{PaperXI.FiniteProductDeterminant}, \nolinkurl{PaperXI.OddPhaseTransport}, \nolinkurl{PaperXI.OddSymplectic}.

\begin{longtable}{@{}p{0.07\textwidth}p{0.12\textwidth}p{0.75\textwidth}@{}}\toprule Line&Kind&Name\\\midrule\endhead

16&def&\nolinkurl{PaperXI.OddPhaseSum.action}\\

20&theorem&\nolinkurl{PaperXI.OddPhaseSum.action_apply}\\

24&theorem&\nolinkurl{PaperXI.OddPhaseSum.difference}\\

31&theorem&\nolinkurl{PaperXI.OddPhaseSum.determinant}\\

36&theorem&\nolinkurl{PaperXI.OddPhaseSum.fixedPointFree}\\

42&theorem&\nolinkurl{PaperXI.OddPhaseSum.half_dimension}\\

58&theorem&\nolinkurl{PaperXI.OddPhaseSum.phase_product}\\

65&theorem&\nolinkurl{PaperXI.OddPhaseSum.phase_product_symplectic}\\

75&theorem&\nolinkurl{PaperXI.OddPhaseSum.of_decomposition}\\

86&theorem&\nolinkurl{PaperXI.OddPhaseSum.fixedPointFree_of_decomposition}\\

\bottomrule\end{longtable}

\section{\texorpdfstring{\nolinkurl{PaperXI.OddPhaseTransport}}{PaperXI.OddPhaseTransport}}\label{module:PaperXI.OddPhaseTransport}

\textbf{Selected layer:} v0.44 parent.\quad\textbf{Source lines:} 72.\par\textbf{SHA-256:} {\scriptsize\nolinkurl{97c339e43d2c59d8b62995b29c986d88df5214c3df3ae472fd5249ad2ef9ba19}}.

\VerbatimInput[fontsize=\scriptsize,breaklines=true,breakanywhere=true]{excerpts/D_OddPhaseTransport.txt}

\textbf{Direct project imports:} \nolinkurl{PaperXI.OddPhase}.

\begin{longtable}{@{}p{0.07\textwidth}p{0.12\textwidth}p{0.75\textwidth}@{}}\toprule Line&Kind&Name\\\midrule\endhead

13&theorem&\nolinkurl{PaperXI.OddPhase.difference_intertwines}\\

18&theorem&\nolinkurl{PaperXI.OddPhase.difference_conjugacy}\\

26&theorem&\nolinkurl{PaperXI.OddPhase.fixedPointFree_conjugacy}\\

42&theorem&\nolinkurl{PaperXI.OddPhase.determinant_conjugacy}\\

50&theorem&\nolinkurl{PaperXI.OddPhase.conjugacy}\\

56&theorem&\nolinkurl{PaperXI.OddPhase.square_class}\\

64&theorem&\nolinkurl{PaperXI.OddPhase.square_class_of_exists}\\

\bottomrule\end{longtable}

\section{\texorpdfstring{\nolinkurl{PaperXI.OddSymplectic}}{PaperXI.OddSymplectic}}\label{module:PaperXI.OddSymplectic}

\textbf{Selected layer:} v0.44 parent.\quad\textbf{Source lines:} 70.\par\textbf{SHA-256:} {\scriptsize\nolinkurl{e1b15050070b85893ab0d094b3ffaf7aaed326673ec11418c39d52487354e063}}.

\VerbatimInput[fontsize=\scriptsize,breaklines=true,breakanywhere=true]{excerpts/D_OddSymplectic.txt}

\textbf{Direct project imports:} none.

\begin{longtable}{@{}p{0.07\textwidth}p{0.12\textwidth}p{0.75\textwidth}@{}}\toprule Line&Kind&Name\\\midrule\endhead

17&theorem&\nolinkurl{PaperXI.OddSymplectic.transpose_toMatrix}\\

26&theorem&\nolinkurl{PaperXI.OddSymplectic.even_finrank}\\

44&theorem&\nolinkurl{PaperXI.OddSymplectic.finrank_eq_two_mul_half}\\

51&theorem&\nolinkurl{PaperXI.OddSymplectic.two_ne_zero}\\

59&theorem&\nolinkurl{PaperXI.OddSymplectic.even_finrank_zmod}\\

64&theorem&\nolinkurl{PaperXI.OddSymplectic.finrank_eq_two_mul_half_zmod}\\

\bottomrule\end{longtable}

\section{\texorpdfstring{\nolinkurl{PaperXI.OnePlusIndex110}}{PaperXI.OnePlusIndex110}}\label{module:PaperXI.OnePlusIndex110}

\textbf{Selected layer:} Continuation110.\quad\textbf{Source lines:} 87.\par\textbf{SHA-256:} {\scriptsize\nolinkurl{733e4888e34cb2eec83dd8b0d29ea395d269e9caa209f00d1fb751c61563ef49}}.

\VerbatimInput[fontsize=\scriptsize,breaklines=true,breakanywhere=true]{excerpts/D_OnePlusIndex110.txt}

\textbf{Direct project imports:} none.

\begin{longtable}{@{}p{0.07\textwidth}p{0.12\textwidth}p{0.75\textwidth}@{}}\toprule Line&Kind&Name\\\midrule\endhead

15&def&\nolinkurl{PaperXI.OnePlusIndex110.congruence}\\

35&def&\nolinkurl{PaperXI.OnePlusIndex110.offsetEquiv}\\

52&theorem&\nolinkurl{PaperXI.OnePlusIndex110.inv_mul_sub_one_mem_iff}\\

67&def&\nolinkurl{PaperXI.OnePlusIndex110.cosetEquiv}\\

81&theorem&\nolinkurl{PaperXI.OnePlusIndex110.index_eq_additive_relIndex}\\

\bottomrule\end{longtable}

\section{\texorpdfstring{\nolinkurl{PaperXI.OnePlusIntersection110}}{PaperXI.OnePlusIntersection110}}\label{module:PaperXI.OnePlusIntersection110}

\textbf{Selected layer:} Continuation110.\quad\textbf{Source lines:} 197.\par\textbf{SHA-256:} {\scriptsize\nolinkurl{512e3b7a894d7db7474f8586c4c652dddd45bfd98197f637f7be8e5a8fec3e7e}}.

\VerbatimInput[fontsize=\scriptsize,breaklines=true,breakanywhere=true]{excerpts/D_OnePlusIntersection110.txt}

\textbf{Direct project imports:} \nolinkurl{PaperXI.OnePlusIndex110}.

\begin{longtable}{@{}p{0.07\textwidth}p{0.12\textwidth}p{0.75\textwidth}@{}}\toprule Line&Kind&Name\\\midrule\endhead

15&def&\nolinkurl{PaperXI.OnePlusIntersection110.idealUnits}\\

18&theorem&\nolinkurl{PaperXI.OnePlusIntersection110.mem_idealUnits}\\

21&theorem&\nolinkurl{PaperXI.OnePlusIntersection110.mem_of_offset}\\

32&theorem&\nolinkurl{PaperXI.OnePlusIntersection110.idealUnits_le}\\

38&theorem&\nolinkurl{PaperXI.OnePlusIntersection110.idealUnits_relIndex}\\

50&theorem&\nolinkurl{PaperXI.OnePlusIntersection110.offset_le_of_group_le}\\

62&theorem&\nolinkurl{PaperXI.OnePlusIntersection110.nested_relIndex_eq}\\

90&def&\nolinkurl{PaperXI.OnePlusIntersection110.innerRing}\\

99&theorem&\nolinkurl{PaperXI.OnePlusIntersection110.mapped_offsets}\\

111&theorem&\nolinkurl{PaperXI.OnePlusIntersection110.mapped_onePlus}\\

124&def&\nolinkurl{PaperXI.OnePlusIntersection110.pullbackGroup}\\

127&def&\nolinkurl{PaperXI.OnePlusIntersection110.pullbackLattice}\\

130&theorem&\nolinkurl{PaperXI.OnePlusIntersection110.pullback_offsets}\\

143&theorem&\nolinkurl{PaperXI.OnePlusIntersection110.pullback_onePlus}\\

161&theorem&\nolinkurl{PaperXI.OnePlusIntersection110.pullback_intersection_relIndex_eq}\\

179&theorem&\nolinkurl{PaperXI.OnePlusIntersection110.conjugate_intersection_relIndex_eq}\\

\bottomrule\end{longtable}

\section{\texorpdfstring{\nolinkurl{PaperXI.OpenSubgroupSmoothInduction}}{PaperXI.OpenSubgroupSmoothInduction}}\label{module:PaperXI.OpenSubgroupSmoothInduction}

\textbf{Selected layer:} base108 arithmetic branch.\quad\textbf{Source lines:} 409.\par\textbf{SHA-256:} {\scriptsize\nolinkurl{eec2e1f81ea1ed3ed0161aa0df8518115a4605b351d08127ef4282943d5978f7}}.

\VerbatimInput[fontsize=\scriptsize,breaklines=true,breakanywhere=true]{excerpts/D_OpenSubgroupSmoothInduction.txt}

\textbf{Direct project imports:} \nolinkurl{PaperXI.SmoothDegreeQuotient}.

\begin{longtable}{@{}p{0.07\textwidth}p{0.12\textwidth}p{0.75\textwidth}@{}}\toprule Line&Kind&Name\\\midrule\endhead

25&def&\nolinkurl{PaperXI.OpenSubgroupSmoothInduction.smoothSubmodule}\\

48&theorem&\nolinkurl{PaperXI.OpenSubgroupSmoothInduction.mem_smoothSubmodule}\\

57&theorem&\nolinkurl{PaperXI.OpenSubgroupSmoothInduction.mk_subgroup}\\

65&theorem&\nolinkurl{PaperXI.OpenSubgroupSmoothInduction.mk_smooth}\\

84&theorem&\nolinkurl{PaperXI.OpenSubgroupSmoothInduction.ind_isSmooth}\\

100&def&\nolinkurl{PaperXI.OpenSubgroupSmoothInduction.functionModel}\\

105&theorem&\nolinkurl{PaperXI.OpenSubgroupSmoothInduction.functionModel_mk}\\

114&theorem&\nolinkurl{PaperXI.OpenSubgroupSmoothInduction.functionModel_mk_self}\\

122&theorem&\nolinkurl{PaperXI.OpenSubgroupSmoothInduction.mk_injective}\\

129&theorem&\nolinkurl{PaperXI.OpenSubgroupSmoothInduction.mk_ne_zero}\\

136&theorem&\nolinkurl{PaperXI.OpenSubgroupSmoothInduction.ind_nontrivial}\\

140&def&\nolinkurl{PaperXI.OpenSubgroupSmoothInduction.finiteCosetSupport}\\

172&theorem&\nolinkurl{PaperXI.OpenSubgroupSmoothInduction.functionModel_finite_support}\\

194&theorem&\nolinkurl{PaperXI.OpenSubgroupSmoothInduction.functionModel_finite_quotient_support}\\

205&theorem&\nolinkurl{PaperXI.OpenSubgroupSmoothInduction.functionModel_compact_quotient_support}\\

212&def&\nolinkurl{PaperXI.OpenSubgroupSmoothInduction.functionIntertwiner}\\

225&def&\nolinkurl{PaperXI.OpenSubgroupSmoothInduction.evaluate}\\

229&theorem&\nolinkurl{PaperXI.OpenSubgroupSmoothInduction.evaluate_mk_self}\\

233&theorem&\nolinkurl{PaperXI.OpenSubgroupSmoothInduction.evaluate_mk_of_not_rel}\\

243&theorem&\nolinkurl{PaperXI.OpenSubgroupSmoothInduction.mk_linearIndependent}\\

264&def&\nolinkurl{PaperXI.OpenSubgroupSmoothInduction.cosetComponent}\\

274&theorem&\nolinkurl{PaperXI.OpenSubgroupSmoothInduction.cosetComponent_mk}\\

279&theorem&\nolinkurl{PaperXI.OpenSubgroupSmoothInduction.cosetComponent_add}\\

286&theorem&\nolinkurl{PaperXI.OpenSubgroupSmoothInduction.cosetComponent_smul}\\

293&theorem&\nolinkurl{PaperXI.OpenSubgroupSmoothInduction.cosetComponent_finite}\\

309&def&\nolinkurl{PaperXI.OpenSubgroupSmoothInduction.reconstruct}\\

322&def&\nolinkurl{PaperXI.OpenSubgroupSmoothInduction.compactModel}\\

326&theorem&\nolinkurl{PaperXI.OpenSubgroupSmoothInduction.reconstruct_compactModel}\\

342&theorem&\nolinkurl{PaperXI.OpenSubgroupSmoothInduction.compactModel_reconstruct}\\

358&def&\nolinkurl{PaperXI.OpenSubgroupSmoothInduction.compactModelEquiv}\\

364&theorem&\nolinkurl{PaperXI.OpenSubgroupSmoothInduction.functionModel_injective}\\

371&theorem&\nolinkurl{PaperXI.OpenSubgroupSmoothInduction.functionModel_range}\\

382&theorem&\nolinkurl{PaperXI.OpenSubgroupSmoothInduction.finite_functions_smooth}\\

\bottomrule\end{longtable}

\section{\texorpdfstring{\nolinkurl{PaperXI.OrbitSpan}}{PaperXI.OrbitSpan}}\label{module:PaperXI.OrbitSpan}

\textbf{Selected layer:} v0.44 parent.\quad\textbf{Source lines:} 86.\par\textbf{SHA-256:} {\scriptsize\nolinkurl{1c23d6d51d10af04e9788052c0601a726639e8d6938f47b653949b20a8bdd0c3}}.

\VerbatimInput[fontsize=\scriptsize,breaklines=true,breakanywhere=true]{excerpts/D_OrbitSpan.txt}

\textbf{Direct project imports:} \nolinkurl{PaperXI.FrobeniusOrbit}.

\begin{longtable}{@{}p{0.07\textwidth}p{0.12\textwidth}p{0.75\textwidth}@{}}\toprule Line&Kind&Name\\\midrule\endhead

12&def&\nolinkurl{PaperXI.OrbitSpan.orbitSpan}\\

16&theorem&\nolinkurl{PaperXI.OrbitSpan.packet_le}\\

22&theorem&\nolinkurl{PaperXI.OrbitSpan.self_packet_le}\\

26&theorem&\nolinkurl{PaperXI.OrbitSpan.map_mem}\\

36&theorem&\nolinkurl{PaperXI.OrbitSpan.inertia_stable}\\

42&theorem&\nolinkurl{PaperXI.OrbitSpan.forward_stable}\\

49&theorem&\nolinkurl{PaperXI.OrbitSpan.backward_stable}\\

56&theorem&\nolinkurl{PaperXI.OrbitSpan.outside_projector_zero}\\

66&theorem&\nolinkurl{PaperXI.OrbitSpan.span_ne_bot}\\

71&theorem&\nolinkurl{PaperXI.OrbitSpan.outside_packet_eq_bot}\\

80&theorem&\nolinkurl{PaperXI.OrbitSpan.related_of_span_top}\\

\bottomrule\end{longtable}

\section{\texorpdfstring{\nolinkurl{PaperXI.OrthogonalPhase}}{PaperXI.OrthogonalPhase}}\label{module:PaperXI.OrthogonalPhase}

\textbf{Selected layer:} v0.44 parent.\quad\textbf{Source lines:} 111.\par\textbf{SHA-256:} {\scriptsize\nolinkurl{f822d2a61741660dd23d656ef3748309c17ca6e982ee0b05bb69f6b06e8edeb0}}.

\VerbatimInput[fontsize=\scriptsize,breaklines=true,breakanywhere=true]{excerpts/D_OrthogonalPhase.txt}

\textbf{Direct project imports:} \nolinkurl{PaperXI.OddPhaseSum}, \nolinkurl{PaperXI.OrthogonalSummands}.

\begin{longtable}{@{}p{0.07\textwidth}p{0.12\textwidth}p{0.75\textwidth}@{}}\toprule Line&Kind&Name\\\midrule\endhead

14&def&\nolinkurl{PaperXI.OrthogonalPhase.decomposition}\\

19&theorem&\nolinkurl{PaperXI.OrthogonalPhase.decomposition_sum}\\

47&def&\nolinkurl{PaperXI.OrthogonalPhase.restricted}\\

52&theorem&\nolinkurl{PaperXI.OrthogonalPhase.restricted_coe}\\

56&theorem&\nolinkurl{PaperXI.OrthogonalPhase.intertwines}\\

68&theorem&\nolinkurl{PaperXI.OrthogonalPhase.phase_product}\\

81&theorem&\nolinkurl{PaperXI.OrthogonalPhase.fixedPointFree}\\

91&theorem&\nolinkurl{PaperXI.OrthogonalPhase.source}\\

\bottomrule\end{longtable}

\section{\texorpdfstring{\nolinkurl{PaperXI.OrthogonalSummands}}{PaperXI.OrthogonalSummands}}\label{module:PaperXI.OrthogonalSummands}

\textbf{Selected layer:} v0.44 parent.\quad\textbf{Source lines:} 120.\par\textbf{SHA-256:} {\scriptsize\nolinkurl{fa9f0303b80e64cd56907157d0dc0c86d1f3e34d1ebd579bf0614a9c89116898}}.

\VerbatimInput[fontsize=\scriptsize,breaklines=true,breakanywhere=true]{excerpts/D_OrthogonalSummands.txt}

\textbf{Direct project imports:} \nolinkurl{PaperXI.OddSymplectic}.

\begin{longtable}{@{}p{0.07\textwidth}p{0.12\textwidth}p{0.75\textwidth}@{}}\toprule Line&Kind&Name\\\midrule\endhead

17&theorem&\nolinkurl{PaperXI.OrthogonalSummands.restricted_isAlt}\\

22&theorem&\nolinkurl{PaperXI.OrthogonalSummands.restricted_form_nondegenerate}\\

58&theorem&\nolinkurl{PaperXI.OrthogonalSummands.iSup_eq_top_of_equiv_sum}\\

71&def&\nolinkurl{PaperXI.OrthogonalSummands.restrictionEquiv}\\

78&theorem&\nolinkurl{PaperXI.OrthogonalSummands.restrictionEquiv_coe}\\

82&theorem&\nolinkurl{PaperXI.OrthogonalSummands.restrictionEquiv_symm_coe}\\

91&theorem&\nolinkurl{PaperXI.OrthogonalSummands.inverse_mem}\\

97&theorem&\nolinkurl{PaperXI.OrthogonalSummands.restrictionEquiv_form}\\

104&theorem&\nolinkurl{PaperXI.OrthogonalSummands.restrictionEquiv_fixed_eq_zero}\\

112&theorem&\nolinkurl{PaperXI.OrthogonalSummands.restricted_even_finrank}\\

\bottomrule\end{longtable}

\section{\texorpdfstring{\nolinkurl{PaperXI.PacketBasis}}{PaperXI.PacketBasis}}\label{module:PaperXI.PacketBasis}

\textbf{Selected layer:} v0.44 parent.\quad\textbf{Source lines:} 97.\par\textbf{SHA-256:} {\scriptsize\nolinkurl{c5f5d83a01759c8aa3984c2b11462888159937d7a2fd4c3f407cc7495a5aee45}}.

\VerbatimInput[fontsize=\scriptsize,breaklines=true,breakanywhere=true]{excerpts/D_PacketBasis.txt}

\textbf{Direct project imports:} \nolinkurl{PaperXI.PacketLines}.

\begin{longtable}{@{}p{0.07\textwidth}p{0.12\textwidth}p{0.75\textwidth}@{}}\toprule Line&Kind&Name\\\midrule\endhead

16&theorem&\nolinkurl{PaperXI.PacketBasis.cycleVector_ne_zero}\\

22&theorem&\nolinkurl{PaperXI.PacketBasis.cycle_linearIndependent}\\

39&def&\nolinkurl{PaperXI.PacketBasis.orbitBasis}\\

52&theorem&\nolinkurl{PaperXI.PacketBasis.orbitBasis_apply}\\

63&theorem&\nolinkurl{PaperXI.PacketBasis.finrank_eq_period_of_return}\\

74&theorem&\nolinkurl{PaperXI.PacketBasis.finrank_eq_period}\\

86&theorem&\nolinkurl{PaperXI.PacketBasis.finrank_eq_orbit_card}\\

\bottomrule\end{longtable}

\section{\texorpdfstring{\nolinkurl{PaperXI.PacketCycle}}{PaperXI.PacketCycle}}\label{module:PaperXI.PacketCycle}

\textbf{Selected layer:} v0.44 parent.\quad\textbf{Source lines:} 76.\par\textbf{SHA-256:} {\scriptsize\nolinkurl{55a96ecf51c156847ccfe2269e6caecfdfb45dbb4792983af8cb809e3a438bc8}}.

\VerbatimInput[fontsize=\scriptsize,breaklines=true,breakanywhere=true]{excerpts/D_PacketCycle.txt}

\textbf{Direct project imports:} \nolinkurl{PaperXI.CycleCoordinates}, \nolinkurl{PaperXI.PacketBasis}.

\begin{longtable}{@{}p{0.07\textwidth}p{0.12\textwidth}p{0.75\textwidth}@{}}\toprule Line&Kind&Name\\\midrule\endhead

13&theorem&\nolinkurl{PaperXI.PacketCycle.linear_power}\\

32&theorem&\nolinkurl{PaperXI.PacketCycle.charpoly}\\

43&theorem&\nolinkurl{PaperXI.PacketCycle.det}\\

54&theorem&\nolinkurl{PaperXI.PacketCycle.ambient_polynomial}\\

66&theorem&\nolinkurl{PaperXI.PacketCycle.exists_cycle_invariants}\\

\bottomrule\end{longtable}

\section{\texorpdfstring{\nolinkurl{PaperXI.PacketLines}}{PaperXI.PacketLines}}\label{module:PaperXI.PacketLines}

\textbf{Selected layer:} v0.44 parent.\quad\textbf{Source lines:} 147.\par\textbf{SHA-256:} {\scriptsize\nolinkurl{a7c34f355ec41735f4fe274cde83c52d9d91c22db3eabb38d6d7de1c7222a4dd}}.

\VerbatimInput[fontsize=\scriptsize,breaklines=true,breakanywhere=true]{excerpts/D_PacketLines.txt}

\textbf{Direct project imports:} \nolinkurl{PaperXI.CycleSpan}.

\begin{longtable}{@{}p{0.07\textwidth}p{0.12\textwidth}p{0.75\textwidth}@{}}\toprule Line&Kind&Name\\\midrule\endhead

13&def&\nolinkurl{PaperXI.PacketLines.orbitLabel}\\

16&theorem&\nolinkurl{PaperXI.PacketLines.orbitLabel_injective}\\

22&theorem&\nolinkurl{PaperXI.PacketLines.cycle_mem_packet}\\

29&theorem&\nolinkurl{PaperXI.PacketLines.inertia_cycle}\\

38&theorem&\nolinkurl{PaperXI.PacketLines.inertia_stable}\\

49&theorem&\nolinkurl{PaperXI.PacketLines.projector_cycle}\\

66&theorem&\nolinkurl{PaperXI.PacketLines.projection_to_line}\\

83&theorem&\nolinkurl{PaperXI.PacketLines.packet_eq_line_of_span_top}\\

98&theorem&\nolinkurl{PaperXI.PacketLines.generators_preserve_cycleSpan}\\

118&theorem&\nolinkurl{PaperXI.PacketLines.cycleSpan_eq_top}\\

131&theorem&\nolinkurl{PaperXI.PacketLines.nonzero_packet_finrank_one}\\

\bottomrule\end{longtable}

\section{\texorpdfstring{\nolinkurl{PaperXI.PacketReturn}}{PaperXI.PacketReturn}}\label{module:PaperXI.PacketReturn}

\textbf{Selected layer:} v0.44 parent.\quad\textbf{Source lines:} 61.\par\textbf{SHA-256:} {\scriptsize\nolinkurl{54669679788cd0095e0d26b16343de68a5196641db0e60dad21ec97090f52942}}.

\VerbatimInput[fontsize=\scriptsize,breaklines=true,breakanywhere=true]{excerpts/D_PacketReturn.txt}

\textbf{Direct project imports:} \nolinkurl{PaperXI.FrobeniusPeriod}.

\begin{longtable}{@{}p{0.07\textwidth}p{0.12\textwidth}p{0.75\textwidth}@{}}\toprule Line&Kind&Name\\\midrule\endhead

13&theorem&\nolinkurl{PaperXI.PacketReturn.power_mem}\\

23&theorem&\nolinkurl{PaperXI.PacketReturn.return_mem}\\

30&def&\nolinkurl{PaperXI.PacketReturn.closingOperator}\\

40&theorem&\nolinkurl{PaperXI.PacketReturn.exists_closing_eigenvector}\\

\bottomrule\end{longtable}

\section{\texorpdfstring{\nolinkurl{PaperXI.PacketSupport}}{PaperXI.PacketSupport}}\label{module:PaperXI.PacketSupport}

\textbf{Selected layer:} v0.44 parent.\quad\textbf{Source lines:} 105.\par\textbf{SHA-256:} {\scriptsize\nolinkurl{3265ba26afef8a86567d0e4ecc83f6448bf3458790ff26d376ae84f7834b78f1}}.

\VerbatimInput[fontsize=\scriptsize,breaklines=true,breakanywhere=true]{excerpts/D_PacketSupport.txt}

\textbf{Direct project imports:} \nolinkurl{PaperXI.TameOrbitTransitivity}.

\begin{longtable}{@{}p{0.07\textwidth}p{0.12\textwidth}p{0.75\textwidth}@{}}\toprule Line&Kind&Name\\\midrule\endhead

12&theorem&\nolinkurl{PaperXI.PacketSupport.nonzero_step_iff}\\

19&theorem&\nolinkurl{PaperXI.PacketSupport.nonzero_back_iff}\\

26&theorem&\nolinkurl{PaperXI.PacketSupport.nonzero_zpow_iff}\\

39&theorem&\nolinkurl{PaperXI.PacketSupport.finrank_zpow}\\

68&theorem&\nolinkurl{PaperXI.PacketSupport.support_iff_orbit}\\

82&theorem&\nolinkurl{PaperXI.PacketSupport.support_is_one_orbit}\\

92&theorem&\nolinkurl{PaperXI.PacketSupport.equal_packet_dimensions}\\

\bottomrule\end{longtable}

\section{\texorpdfstring{\nolinkurl{PaperXI.PacketTransport}}{PaperXI.PacketTransport}}\label{module:PaperXI.PacketTransport}

\textbf{Selected layer:} v0.44 parent.\quad\textbf{Source lines:} 81.\par\textbf{SHA-256:} {\scriptsize\nolinkurl{dc7282e77c935ca61b8807fe87f4fe2fa372d720a28d0a1c8982c6c749e1a910}}.

\VerbatimInput[fontsize=\scriptsize,breaklines=true,breakanywhere=true]{excerpts/D_PacketTransport.txt}

\textbf{Direct project imports:} \nolinkurl{PaperXI.TameProjectors}.

\begin{longtable}{@{}p{0.07\textwidth}p{0.12\textwidth}p{0.75\textwidth}@{}}\toprule Line&Kind&Name\\\midrule\endhead

11&abbrev&\nolinkurl{PaperXI.PacketTransport.packet}\\

14&abbrev&\nolinkurl{PaperXI.PacketTransport.projector}\\

17&theorem&\nolinkurl{PaperXI.PacketTransport.mem_iff_projector}\\

25&theorem&\nolinkurl{PaperXI.PacketTransport.projector_on_other_packet}\\

32&theorem&\nolinkurl{PaperXI.PacketTransport.inertia_mem}\\

38&theorem&\nolinkurl{PaperXI.PacketTransport.transport_of_conjugation}\\

46&theorem&\nolinkurl{PaperXI.PacketTransport.forward_mem}\\

53&theorem&\nolinkurl{PaperXI.PacketTransport.backward_mem}\\

62&def&\nolinkurl{PaperXI.PacketTransport.packetEquiv}\\

75&theorem&\nolinkurl{PaperXI.PacketTransport.packet_finrank_step}\\

\bottomrule\end{longtable}

\section{\texorpdfstring{\nolinkurl{PaperXI.PacketTypeAttachment70}}{PaperXI.PacketTypeAttachment70}}\label{module:PaperXI.PacketTypeAttachment70}

\textbf{Selected layer:} base108 arithmetic branch.\quad\textbf{Source lines:} 85.\par\textbf{SHA-256:} {\scriptsize\nolinkurl{ed82a000e951e6dca19a854165a2e9f92be60a3bcc1ed876f0fae8a9978b8a79}}.

\VerbatimInput[fontsize=\scriptsize,breaklines=true,breakanywhere=true]{excerpts/D_PacketTypeAttachment70.txt}

\textbf{Direct project imports:} \nolinkurl{PaperXI.TensorTypeOperators70}.

\begin{longtable}{@{}p{0.07\textwidth}p{0.12\textwidth}p{0.75\textwidth}@{}}\toprule Line&Kind&Name\\\midrule\endhead

17&def&\nolinkurl{PaperXI.PacketTypeAttachment70.typeAction}\\

21&theorem&\nolinkurl{PaperXI.PacketTypeAttachment70.ambient_action}\\

32&theorem&\nolinkurl{PaperXI.PacketTypeAttachment70.first_factor_action}\\

47&theorem&\nolinkurl{PaperXI.PacketTypeAttachment70.first_factor_average}\\

61&theorem&\nolinkurl{PaperXI.PacketTypeAttachment70.generator_action}\\

71&theorem&\nolinkurl{PaperXI.PacketTypeAttachment70.edge_compatibility_iff}\\

79&theorem&\nolinkurl{PaperXI.PacketTypeAttachment70.edge_residual_iff}\\

\bottomrule\end{longtable}

\section{\texorpdfstring{\nolinkurl{PaperXI.PairIndex110}}{PaperXI.PairIndex110}}\label{module:PaperXI.PairIndex110}

\textbf{Selected layer:} Continuation110.\quad\textbf{Source lines:} 122.\par\textbf{SHA-256:} {\scriptsize\nolinkurl{453d18235bd27cd82e48a7cf367e7b3433288906412d7e7154e80b955025408b}}.

\VerbatimInput[fontsize=\scriptsize,breaklines=true,breakanywhere=true]{excerpts/D_PairIndex110.txt}

\textbf{Direct project imports:} \nolinkurl{PaperXI.IntersectionModel108}.

\begin{longtable}{@{}p{0.07\textwidth}p{0.12\textwidth}p{0.75\textwidth}@{}}\toprule Line&Kind&Name\\\midrule\endhead

19&def&\nolinkurl{PaperXI.PairIndex110.intersection}\\

23&def&\nolinkurl{PaperXI.PairIndex110.scalarIntersection}\\

26&theorem&\nolinkurl{PaperXI.PairIndex110.scalarIntersection_le}\\

29&theorem&\nolinkurl{PaperXI.PairIndex110.scalarIntersection_le_intersection}\\

35&theorem&\nolinkurl{PaperXI.PairIndex110.intersection_image}\\

47&theorem&\nolinkurl{PaperXI.PairIndex110.scalarIntersection_image}\\

58&theorem&\nolinkurl{PaperXI.PairIndex110.left_mem_intersection}\\

62&theorem&\nolinkurl{PaperXI.PairIndex110.left_range}\\

71&theorem&\nolinkurl{PaperXI.PairIndex110.left_injective}\\

82&def&\nolinkurl{PaperXI.PairIndex110.pairEquivIntersection}\\

91&theorem&\nolinkurl{PaperXI.PairIndex110.pairSupport_eq_comap}\\

106&theorem&\nolinkurl{PaperXI.PairIndex110.pairSupport_index}\\

113&theorem&\nolinkurl{PaperXI.PairIndex110.pairSupport_index_ambient}\\

\bottomrule\end{longtable}

\section{\texorpdfstring{\nolinkurl{PaperXI.PeriodicLatticeReduction60}}{PaperXI.PeriodicLatticeReduction60}}\label{module:PaperXI.PeriodicLatticeReduction60}

\textbf{Selected layer:} base108 arithmetic branch.\quad\textbf{Source lines:} 86.\par\textbf{SHA-256:} {\scriptsize\nolinkurl{f70d40f2ca38dfa0635a4232d830d16dbbdc2f58c4472d1c82c4a25950e93a02}}.

\VerbatimInput[fontsize=\scriptsize,breaklines=true,breakanywhere=true]{excerpts/D_PeriodicLatticeReduction60.txt}

\textbf{Direct project imports:} \nolinkurl{PaperXI.ActualELatticeNormalizer59}.

\begin{longtable}{@{}p{0.07\textwidth}p{0.12\textwidth}p{0.75\textwidth}@{}}\toprule Line&Kind&Name\\\midrule\endhead

17&def&\nolinkurl{PaperXI.PeriodicLatticeReduction60.levelCondition}\\

21&theorem&\nolinkurl{PaperXI.PeriodicLatticeReduction60.levelCondition_shift}\\

53&theorem&\nolinkurl{PaperXI.PeriodicLatticeReduction60.levelCondition_periodic}\\

60&theorem&\nolinkurl{PaperXI.PeriodicLatticeReduction60.mem_firstCongruence_iff_fin}\\

\bottomrule\end{longtable}

\section{\texorpdfstring{\nolinkurl{PaperXI.PiPolynomial}}{PaperXI.PiPolynomial}}\label{module:PaperXI.PiPolynomial}

\textbf{Selected layer:} v0.44 parent.\quad\textbf{Source lines:} 58.\par\textbf{SHA-256:} {\scriptsize\nolinkurl{3aaecfa5d206cd6d476431f7e9f3f823ec7b57ec488a1edb5d6170a3c973496e}}.

\VerbatimInput[fontsize=\scriptsize,breaklines=true,breakanywhere=true]{excerpts/D_PiPolynomial.txt}

\textbf{Direct project imports:} \nolinkurl{PaperXI.InvariantPolynomial}.

\begin{longtable}{@{}p{0.07\textwidth}p{0.12\textwidth}p{0.75\textwidth}@{}}\toprule Line&Kind&Name\\\midrule\endhead

13&theorem&\nolinkurl{PaperXI.PiPolynomial.charpoly_pi}\\

46&theorem&\nolinkurl{PaperXI.PiPolynomial.operatorPolynomial_pi}\\

\bottomrule\end{longtable}

\section{\texorpdfstring{\nolinkurl{PaperXI.PolynomialAssembly}}{PaperXI.PolynomialAssembly}}\label{module:PaperXI.PolynomialAssembly}

\textbf{Selected layer:} v0.44 parent.\quad\textbf{Source lines:} 83.\par\textbf{SHA-256:} {\scriptsize\nolinkurl{538480496e4eb1bceb57fd376c890da010064e5ca4d64edc62ceb4394e77663f}}.

\VerbatimInput[fontsize=\scriptsize,breaklines=true,breakanywhere=true]{excerpts/D_PolynomialAssembly.txt}

\textbf{Direct project imports:} \nolinkurl{PaperXI.InvariantAssembly}, \nolinkurl{PaperXI.PiPolynomial}.

\begin{longtable}{@{}p{0.07\textwidth}p{0.12\textwidth}p{0.75\textwidth}@{}}\toprule Line&Kind&Name\\\midrule\endhead

14&def&\nolinkurl{PaperXI.PolynomialAssembly.restricted}\\

20&theorem&\nolinkurl{PaperXI.PolynomialAssembly.restricted_apply}\\

24&theorem&\nolinkurl{PaperXI.PolynomialAssembly.stable_of_decomposition}\\

37&theorem&\nolinkurl{PaperXI.PolynomialAssembly.restricted_conj}\\

53&theorem&\nolinkurl{PaperXI.PolynomialAssembly.polynomial_product}\\

67&def&\nolinkurl{PaperXI.PolynomialAssembly.localFactor}\\

71&theorem&\nolinkurl{PaperXI.PolynomialAssembly.localFactor_product}\\

\bottomrule\end{longtable}

\section{\texorpdfstring{\nolinkurl{PaperXI.PrescribedFieldType70}}{PaperXI.PrescribedFieldType70}}\label{module:PaperXI.PrescribedFieldType70}

\textbf{Selected layer:} base108 arithmetic branch.\quad\textbf{Source lines:} 87.\par\textbf{SHA-256:} {\scriptsize\nolinkurl{6cd430fa9358f4bcd875eaede181bed47500c5d2d71f94e4d937f10231a5d88f}}.

\VerbatimInput[fontsize=\scriptsize,breaklines=true,breakanywhere=true]{excerpts/D_PrescribedFieldType70.txt}

\textbf{Direct project imports:} \nolinkurl{PaperXI.ActualCompactTypeCompletion58}, \nolinkurl{PaperXI.ActualFieldTypeQuotient69}.

\begin{longtable}{@{}p{0.07\textwidth}p{0.12\textwidth}p{0.75\textwidth}@{}}\toprule Line&Kind&Name\\\midrule\endhead

28&def&\nolinkurl{PaperXI.PrescribedFieldType70.compactInclusion}\\

31&def&\nolinkurl{PaperXI.PrescribedFieldType70.fieldGenerator}\\

37&def&\nolinkurl{PaperXI.PrescribedFieldType70.domainEquiv}\\

44&theorem&\nolinkurl{PaperXI.PrescribedFieldType70.exists_prescribed_extension}\\

\bottomrule\end{longtable}

\section{\texorpdfstring{\nolinkurl{PaperXI.PrimitiveCharacterBlock}}{PaperXI.PrimitiveCharacterBlock}}\label{module:PaperXI.PrimitiveCharacterBlock}

\textbf{Selected layer:} v0.44 parent.\quad\textbf{Source lines:} 128.\par\textbf{SHA-256:} {\scriptsize\nolinkurl{5950f0f65ec50d3dc6a4c36c40237c087ccc4a4c0eeb35b1fe536e897bcdf6e5}}.

\VerbatimInput[fontsize=\scriptsize,breaklines=true,breakanywhere=true]{excerpts/D_PrimitiveCharacterBlock.txt}

\textbf{Direct project imports:} \nolinkurl{PaperXI.SemisimpleDetection}.

\begin{longtable}{@{}p{0.07\textwidth}p{0.12\textwidth}p{0.75\textwidth}@{}}\toprule Line&Kind&Name\\\midrule\endhead

17&theorem&\nolinkurl{PaperXI.PrimitiveCharacterBlock.equiv_algebra_action}\\

26&theorem&\nolinkurl{PaperXI.PrimitiveCharacterBlock.nonzero}\\

35&theorem&\nolinkurl{PaperXI.PrimitiveCharacterBlock.idempotent}\\

43&theorem&\nolinkurl{PaperXI.PrimitiveCharacterBlock.supported_eq_zero}\\

60&theorem&\nolinkurl{PaperXI.PrimitiveCharacterBlock.block_action_injective}\\

69&def&\nolinkurl{PaperXI.PrimitiveCharacterBlock.IsPrimitiveCentral}\\

74&theorem&\nolinkurl{PaperXI.PrimitiveCharacterBlock.primitive_central}\\

99&theorem&\nolinkurl{PaperXI.PrimitiveCharacterBlock.no_nontrivial_central_split}\\

112&theorem&\nolinkurl{PaperXI.PrimitiveCharacterBlock.orthogonal}\\

123&theorem&\nolinkurl{PaperXI.PrimitiveCharacterBlock.presentation_invariant}\\

\bottomrule\end{longtable}

\section{\texorpdfstring{\nolinkurl{PaperXI.ProductNormalizerQuotient59}}{PaperXI.ProductNormalizerQuotient59}}\label{module:PaperXI.ProductNormalizerQuotient59}

\textbf{Selected layer:} base108 arithmetic branch.\quad\textbf{Source lines:} 95.\par\textbf{SHA-256:} {\scriptsize\nolinkurl{253347996905c91f7605f14539fa7abaf854e23c86fc22008272369cf3bd30bd}}.

\VerbatimInput[fontsize=\scriptsize,breaklines=true,breakanywhere=true]{excerpts/D_ProductNormalizerQuotient59.txt}

\textbf{Direct project imports:} \nolinkurl{PaperXI.NormalizerExtensionInducing58}.

\begin{longtable}{@{}p{0.07\textwidth}p{0.12\textwidth}p{0.75\textwidth}@{}}\toprule Line&Kind&Name\\\midrule\endhead

25&theorem&\nolinkurl{PaperXI.ProductNormalizerQuotient59.normalizes_sup}\\

32&def&\nolinkurl{PaperXI.ProductNormalizerQuotient59.conjugation}\\

35&theorem&\nolinkurl{PaperXI.ProductNormalizerQuotient59.conjugation_coe}\\

40&theorem&\nolinkurl{PaperXI.ProductNormalizerQuotient59.map_conjugation}\\

63&theorem&\nolinkurl{PaperXI.ProductNormalizerQuotient59.quotient_conjugation}\\

73&def&\nolinkurl{PaperXI.ProductNormalizerQuotient59.productQuotientEquiv}\\

83&def&\nolinkurl{PaperXI.ProductNormalizerQuotient59.quotientEquivOfKernel}\\

\bottomrule\end{longtable}

\section{\texorpdfstring{\nolinkurl{PaperXI.ProjectiveCocycle}}{PaperXI.ProjectiveCocycle}}\label{module:PaperXI.ProjectiveCocycle}

\textbf{Selected layer:} v0.44 parent.\quad\textbf{Source lines:} 171.\par\textbf{SHA-256:} {\scriptsize\nolinkurl{d2eb7d728fef4b48932df07f3dc13a003f39995cc477c3728f0677302794dcf2}}.

\VerbatimInput[fontsize=\scriptsize,breaklines=true,breakanywhere=true]{excerpts/D_ProjectiveCocycle.txt}

\textbf{Direct project imports:} \nolinkurl{PaperXI.CocycleCommutator}, \nolinkurl{PaperXI.SchurScalar}.

\begin{longtable}{@{}p{0.07\textwidth}p{0.12\textwidth}p{0.75\textwidth}@{}}\toprule Line&Kind&Name\\\midrule\endhead

14&theorem&\nolinkurl{PaperXI.ProjectiveCocycle.scalar_injective}\\

21&def&\nolinkurl{PaperXI.ProjectiveCocycle.cocycle}\\

32&theorem&\nolinkurl{PaperXI.ProjectiveCocycle.normalized}\\

45&theorem&\nolinkurl{PaperXI.ProjectiveCocycle.coefficient_unique}\\

53&theorem&\nolinkurl{PaperXI.ProjectiveCocycle.rescale_coefficients}\\

72&theorem&\nolinkurl{PaperXI.ProjectiveCocycle.commutator_rescale_operators}\\

89&theorem&\nolinkurl{PaperXI.ProjectiveCocycle.cocycle_rescale_operators}\\

101&theorem&\nolinkurl{PaperXI.ProjectiveCocycle.cochain_one}\\

109&theorem&\nolinkurl{PaperXI.ProjectiveCocycle.rescaledFamily_mul}\\

123&theorem&\nolinkurl{PaperXI.ProjectiveCocycle.conjugate_mul}\\

133&theorem&\nolinkurl{PaperXI.ProjectiveCocycle.conjugate_rescale_coefficients}\\

149&theorem&\nolinkurl{PaperXI.ProjectiveCocycle.conjugate_rescale_commutator}\\

\bottomrule\end{longtable}

\section{\texorpdfstring{\nolinkurl{PaperXI.ProjectiveHolonomy}}{PaperXI.ProjectiveHolonomy}}\label{module:PaperXI.ProjectiveHolonomy}

\textbf{Selected layer:} v0.44 parent.\quad\textbf{Source lines:} 173.\par\textbf{SHA-256:} {\scriptsize\nolinkurl{dbc14aa36dc916019628345b811c2d1904c4652d23cddbe71792a93751a8db05}}.

\VerbatimInput[fontsize=\scriptsize,breaklines=true,breakanywhere=true]{excerpts/D_ProjectiveHolonomy.txt}

\textbf{Direct project imports:} \nolinkurl{PaperXI.ScalarHolonomy}, \nolinkurl{PaperXI.GraphTransport}, \nolinkurl{PaperXI.GraphLoopRepresentation}.

\begin{longtable}{@{}p{0.07\textwidth}p{0.12\textwidth}p{0.75\textwidth}@{}}\toprule Line&Kind&Name\\\midrule\endhead

20&def&\nolinkurl{PaperXI.ProjectiveHolonomy.rescale}\\

24&theorem&\nolinkurl{PaperXI.ProjectiveHolonomy.rescale_apply}\\

31&theorem&\nolinkurl{PaperXI.ProjectiveHolonomy.rescale_gauge}\\

39&theorem&\nolinkurl{PaperXI.ProjectiveHolonomy.rescale_path}\\

50&theorem&\nolinkurl{PaperXI.ProjectiveHolonomy.rescale_path_apply}\\

59&theorem&\nolinkurl{PaperXI.ProjectiveHolonomy.rescale_loop}\\

66&theorem&\nolinkurl{PaperXI.ProjectiveHolonomy.rescale_holonomy}\\

75&theorem&\nolinkurl{PaperXI.ProjectiveHolonomy.rescale_isScalar}\\

81&theorem&\nolinkurl{PaperXI.ProjectiveHolonomy.isScalar_iff_paths}\\

93&def&\nolinkurl{PaperXI.ProjectiveHolonomy.loopCharacter}\\

97&theorem&\nolinkurl{PaperXI.ProjectiveHolonomy.loopCharacter_spec}\\

105&theorem&\nolinkurl{PaperXI.ProjectiveHolonomy.rescale_character}\\

116&def&\nolinkurl{PaperXI.ProjectiveHolonomy.ratioH1}\\

119&theorem&\nolinkurl{PaperXI.ProjectiveHolonomy.ratioH1_class}\\

123&theorem&\nolinkurl{PaperXI.ProjectiveHolonomy.ratioH1_unique}\\

129&theorem&\nolinkurl{PaperXI.ProjectiveHolonomy.ratioH1_eq_one_iff}\\

145&def&\nolinkurl{PaperXI.ProjectiveHolonomy.transportRatioH1}\\

152&theorem&\nolinkurl{PaperXI.ProjectiveHolonomy.transportRatioH1_class}\\

162&theorem&\nolinkurl{PaperXI.ProjectiveHolonomy.transportRatioH1_rescale}\\

\bottomrule\end{longtable}

\section{\texorpdfstring{\nolinkurl{PaperXI.ProjectorTransport}}{PaperXI.ProjectorTransport}}\label{module:PaperXI.ProjectorTransport}

\textbf{Selected layer:} v0.44 parent.\quad\textbf{Source lines:} 27.\par\textbf{SHA-256:} {\scriptsize\nolinkurl{c88eba65d49511600eaef67165160b9de7d6f1dcb5c6f56476feb8a5259b3592}}.

\VerbatimInput[fontsize=\scriptsize,breaklines=true,breakanywhere=true]{excerpts/D_ProjectorTransport.txt}

\textbf{Direct project imports:} none.

\begin{longtable}{@{}p{0.07\textwidth}p{0.12\textwidth}p{0.75\textwidth}@{}}\toprule Line&Kind&Name\\\midrule\endhead

14&theorem&\nolinkurl{PaperXI.ProjectorTransport.weighted_sum_conjugation}\\

21&theorem&\nolinkurl{PaperXI.ProjectorTransport.normalized_sum_conjugation}\\

\bottomrule\end{longtable}

\section{\texorpdfstring{\nolinkurl{PaperXI.QuotientRepresentation}}{PaperXI.QuotientRepresentation}}\label{module:PaperXI.QuotientRepresentation}

\textbf{Selected layer:} v0.44 parent.\quad\textbf{Source lines:} 72.\par\textbf{SHA-256:} {\scriptsize\nolinkurl{4439d25c13ba1ce6c73f18c0d67c0a8ce5f96caa0c1befb335168c02a01fffe5}}.

\VerbatimInput[fontsize=\scriptsize,breaklines=true,breakanywhere=true]{excerpts/D_QuotientRepresentation.txt}

\textbf{Direct project imports:} none.

\begin{longtable}{@{}p{0.07\textwidth}p{0.12\textwidth}p{0.75\textwidth}@{}}\toprule Line&Kind&Name\\\midrule\endhead

14&def&\nolinkurl{PaperXI.QuotientRepresentation.subrepresentationEquiv}\\

25&theorem&\nolinkurl{PaperXI.QuotientRepresentation.irreducible_comp_iff}\\

30&theorem&\nolinkurl{PaperXI.QuotientRepresentation.intertwines_comp_iff}\\

42&def&\nolinkurl{PaperXI.QuotientRepresentation.descend}\\

46&theorem&\nolinkurl{PaperXI.QuotientRepresentation.descend_apply}\\

49&theorem&\nolinkurl{PaperXI.QuotientRepresentation.descend_comp}\\

52&theorem&\nolinkurl{PaperXI.QuotientRepresentation.descend_irreducible_iff}\\

58&theorem&\nolinkurl{PaperXI.QuotientRepresentation.descend_irreducible}\\

64&theorem&\nolinkurl{PaperXI.QuotientRepresentation.descend_unique}\\

\bottomrule\end{longtable}

\section{\texorpdfstring{\nolinkurl{PaperXI.RadicalOnePlus111}}{PaperXI.RadicalOnePlus111}}\label{module:PaperXI.RadicalOnePlus111}

\textbf{Selected layer:} Continuation111.\quad\textbf{Source lines:} 58.\par\textbf{SHA-256:} {\scriptsize\nolinkurl{3ebe8185adccd13a0c9d9102e437f22fb910637ef074bfb50331602f7eac41a7}}.

\VerbatimInput[fontsize=\scriptsize,breaklines=true,breakanywhere=true]{excerpts/D_RadicalOnePlus111.txt}

\textbf{Direct project imports:} \nolinkurl{PaperXI.OnePlusIntersection110}.

\begin{longtable}{@{}p{0.07\textwidth}p{0.12\textwidth}p{0.75\textwidth}@{}}\toprule Line&Kind&Name\\\midrule\endhead

16&theorem&\nolinkurl{PaperXI.RadicalOnePlus111.isUnit_of_sub_one_mem_jacobson}\\

34&def&\nolinkurl{PaperXI.RadicalOnePlus111.onePlus}\\

36&theorem&\nolinkurl{PaperXI.RadicalOnePlus111.onePlus_offsets}\\

40&theorem&\nolinkurl{PaperXI.RadicalOnePlus111.onePlus_surjective}\\

53&def&\nolinkurl{PaperXI.RadicalOnePlus111.offsetEquiv}\\

\bottomrule\end{longtable}

\section{\texorpdfstring{\nolinkurl{PaperXI.RamificationPolyphase}}{PaperXI.RamificationPolyphase}}\label{module:PaperXI.RamificationPolyphase}

\textbf{Selected layer:} v0.44 parent.\quad\textbf{Source lines:} 326.\par\textbf{SHA-256:} {\scriptsize\nolinkurl{2f34bd694def457221d6f00b874c8b54547dcc95ee249e3afcf34d71b5b64549}}.

\VerbatimInput[fontsize=\scriptsize,breaklines=true,breakanywhere=true]{excerpts/D_RamificationPolyphase.txt}

\textbf{Direct project imports:} \nolinkurl{PaperXI.ValuationConductor}, \nolinkurl{PaperXI.FiniteAverage}.

\begin{longtable}{@{}p{0.07\textwidth}p{0.12\textwidth}p{0.75\textwidth}@{}}\toprule Line&Kind&Name\\\midrule\endhead

29&theorem&\nolinkurl{PaperXI.RamificationPolyphase.average_comp_operator}\\

40&theorem&\nolinkurl{PaperXI.RamificationPolyphase.average_comp_restricted_average}\\

53&theorem&\nolinkurl{PaperXI.RamificationPolyphase.restricted_average_comp_average}\\

69&def&\nolinkurl{PaperXI.RamificationPolyphase.lowerInclusion}\\

75&theorem&\nolinkurl{PaperXI.RamificationPolyphase.lowerInclusion_val}\\

80&theorem&\nolinkurl{PaperXI.RamificationPolyphase.lift_lowerInclusion}\\

84&def&\nolinkurl{PaperXI.RamificationPolyphase.depthProjector}\\

88&theorem&\nolinkurl{PaperXI.RamificationPolyphase.depthProjector_range}\\

93&theorem&\nolinkurl{PaperXI.RamificationPolyphase.depthProjector_idempotent}\\

100&theorem&\nolinkurl{PaperXI.RamificationPolyphase.depthProjector_comp_of_le}\\

114&theorem&\nolinkurl{PaperXI.RamificationPolyphase.depthProjector_comp_of_le_rev}\\

127&def&\nolinkurl{PaperXI.RamificationPolyphase.detailProjector}\\

130&theorem&\nolinkurl{PaperXI.RamificationPolyphase.detailProjector_apply}\\

135&theorem&\nolinkurl{PaperXI.RamificationPolyphase.detailProjector_idempotent}\\

167&theorem&\nolinkurl{PaperXI.RamificationPolyphase.depthProjector_detailProjector}\\

186&theorem&\nolinkurl{PaperXI.RamificationPolyphase.depthProjector_succ_detailProjector}\\

206&theorem&\nolinkurl{PaperXI.RamificationPolyphase.detailProjector_mem_next_fixed}\\

215&theorem&\nolinkurl{PaperXI.RamificationPolyphase.detailProjector_range}\\

232&theorem&\nolinkurl{PaperXI.RamificationPolyphase.detailProjector_comp_eq_zero}\\

265&theorem&\nolinkurl{PaperXI.RamificationPolyphase.detailProjector_comp_eq_zero_rev}\\

299&theorem&\nolinkurl{PaperXI.RamificationPolyphase.depthProjector_eq_zero_add_sum}\\

311&theorem&\nolinkurl{PaperXI.RamificationPolyphase.identity_eq_inertia_add_details}\\

\bottomrule\end{longtable}

\section{\texorpdfstring{\nolinkurl{PaperXI.RamificationPolyphaseBaseChange}}{PaperXI.RamificationPolyphaseBaseChange}}\label{module:PaperXI.RamificationPolyphaseBaseChange}

\textbf{Selected layer:} v0.44 parent.\quad\textbf{Source lines:} 115.\par\textbf{SHA-256:} {\scriptsize\nolinkurl{eb07cacbdaafae8606ca94ef6790ab0e3c7df865ff062943f1757f3ca826daa2}}.

\VerbatimInput[fontsize=\scriptsize,breaklines=true,breakanywhere=true]{excerpts/D_RamificationPolyphaseBaseChange.txt}

\textbf{Direct project imports:} \nolinkurl{PaperXI.RamificationPolyphase}, \nolinkurl{PaperXI.UnramifiedConductor}.

\begin{longtable}{@{}p{0.07\textwidth}p{0.12\textwidth}p{0.75\textwidth}@{}}\toprule Line&Kind&Name\\\midrule\endhead

36&theorem&\nolinkurl{PaperXI.RamificationPolyphaseBaseChange.depthProjector_eq}\\

59&theorem&\nolinkurl{PaperXI.RamificationPolyphaseBaseChange.detailProjector_eq}\\

69&theorem&\nolinkurl{PaperXI.RamificationPolyphaseBaseChange.detailRange_eq}\\

78&theorem&\nolinkurl{PaperXI.RamificationPolyphaseBaseChange.finiteReconstruction_eq}\\

96&theorem&\nolinkurl{PaperXI.RamificationPolyphaseBaseChange.conductorSummand_eq}\\

103&theorem&\nolinkurl{PaperXI.RamificationPolyphaseBaseChange.swan_eq}\\

110&theorem&\nolinkurl{PaperXI.RamificationPolyphaseBaseChange.artin_eq}\\

\bottomrule\end{longtable}

\section{\texorpdfstring{\nolinkurl{PaperXI.RamificationPolyphaseDimensions}}{PaperXI.RamificationPolyphaseDimensions}}\label{module:PaperXI.RamificationPolyphaseDimensions}

\textbf{Selected layer:} v0.44 parent.\quad\textbf{Source lines:} 180.\par\textbf{SHA-256:} {\scriptsize\nolinkurl{f38214cac70955d55434afad4ca9162cd9145d6a627a88311444bac1906e871d}}.

\VerbatimInput[fontsize=\scriptsize,breaklines=true,breakanywhere=true]{excerpts/D_RamificationPolyphaseDimensions.txt}

\textbf{Direct project imports:} \nolinkurl{PaperXI.RamificationPolyphase}, \nolinkurl{PaperXI.RamificationSeparation}.

\begin{longtable}{@{}p{0.07\textwidth}p{0.12\textwidth}p{0.75\textwidth}@{}}\toprule Line&Kind&Name\\\midrule\endhead

31&def&\nolinkurl{PaperXI.RamificationPolyphaseDimensions.detailSpace}\\

37&theorem&\nolinkurl{PaperXI.RamificationPolyphaseDimensions.fixed_mono}\\

47&theorem&\nolinkurl{PaperXI.RamificationPolyphaseDimensions.fixed_succ_eq_sup_detail}\\

76&theorem&\nolinkurl{PaperXI.RamificationPolyphaseDimensions.fixed_disjoint_detail}\\

97&theorem&\nolinkurl{PaperXI.RamificationPolyphaseDimensions.finrank_fixed_succ}\\

108&theorem&\nolinkurl{PaperXI.RamificationPolyphaseDimensions.finrank_fixed_eq_base_add_details}\\

119&theorem&\nolinkurl{PaperXI.RamificationPolyphaseDimensions.finrank_eq_inertia_add_details}\\

135&theorem&\nolinkurl{PaperXI.RamificationPolyphaseDimensions.exists_identity_eq_inertia_add_details}\\

146&theorem&\nolinkurl{PaperXI.RamificationPolyphaseDimensions.exists_finrank_eq_inertia_add_details}\\

157&theorem&\nolinkurl{PaperXI.RamificationPolyphaseDimensions.codimension_eq_detail_tail}\\

170&theorem&\nolinkurl{PaperXI.RamificationPolyphaseDimensions.summand_eq_weight_mul_detail_tail}\\

\bottomrule\end{longtable}

\section{\texorpdfstring{\nolinkurl{PaperXI.RamificationSeparation}}{PaperXI.RamificationSeparation}}\label{module:PaperXI.RamificationSeparation}

\textbf{Selected layer:} v0.44 parent.\quad\textbf{Source lines:} 78.\par\textbf{SHA-256:} {\scriptsize\nolinkurl{f4b0880d3b76829d7b9274fe0f6817a408aa233aed0a8530687b2ec95f7c6094}}.

\VerbatimInput[fontsize=\scriptsize,breaklines=true,breakanywhere=true]{excerpts/D_RamificationSeparation.txt}

\textbf{Direct project imports:} \nolinkurl{PaperXI.ValuationRamification}.

\begin{longtable}{@{}p{0.07\textwidth}p{0.12\textwidth}p{0.75\textwidth}@{}}\toprule Line&Kind&Name\\\midrule\endhead

17&theorem&\nolinkurl{PaperXI.RamificationSeparation.mem_all_lower_iff}\\

37&theorem&\nolinkurl{PaperXI.RamificationSeparation.eventually_trivial}\\

59&instance&\nolinkurl{PaperXI.RamificationSeparation.valuation_faithful}\\

72&theorem&\nolinkurl{PaperXI.RamificationSeparation.valuation_eventually_trivial}\\

\bottomrule\end{longtable}

\section{\texorpdfstring{\nolinkurl{PaperXI.RamifiedUnramifiedQuotient}}{PaperXI.RamifiedUnramifiedQuotient}}\label{module:PaperXI.RamifiedUnramifiedQuotient}

\textbf{Selected layer:} v0.44 parent.\quad\textbf{Source lines:} 322.\par\textbf{SHA-256:} {\scriptsize\nolinkurl{755d457b0f11dfbac94eaee10326fa290bd5b6a331d8883bb8ce26811f00a747}}.

\VerbatimInput[fontsize=\scriptsize,breaklines=true,breakanywhere=true]{excerpts/D_RamifiedUnramifiedQuotient.txt}

\textbf{Direct project imports:} \nolinkurl{PaperXI.RamificationPolyphase}, \nolinkurl{PaperXI.FiniteUnramifiedWeil}.

\begin{longtable}{@{}p{0.07\textwidth}p{0.12\textwidth}p{0.75\textwidth}@{}}\toprule Line&Kind&Name\\\midrule\endhead

38&abbrev&\nolinkurl{PaperXI.RamifiedUnramifiedQuotient.Inertia}\\

42&def&\nolinkurl{PaperXI.RamifiedUnramifiedQuotient.quotientEquivDegreeRange}\\

47&theorem&\nolinkurl{PaperXI.RamifiedUnramifiedQuotient.degree_surjective_of_degree_one}\\

59&theorem&\nolinkurl{PaperXI.RamifiedUnramifiedQuotient.degree_surjective_of_frobenius_mem_range}\\

66&def&\nolinkurl{PaperXI.RamifiedUnramifiedQuotient.quotientEquivIntegerDegree}\\

72&theorem&\nolinkurl{PaperXI.RamifiedUnramifiedQuotient.quotientEquivIntegerDegree_mk}\\

79&theorem&\nolinkurl{PaperXI.RamifiedUnramifiedQuotient.frobenius_lift_ratio_mem_inertia}\\

102&def&\nolinkurl{PaperXI.RamifiedUnramifiedQuotient.degreeModelMap}\\

107&theorem&\nolinkurl{PaperXI.RamifiedUnramifiedQuotient.degreeModelMap_degree}\\

112&theorem&\nolinkurl{PaperXI.RamifiedUnramifiedQuotient.degreeModelMap_ker}\\

119&theorem&\nolinkurl{PaperXI.RamifiedUnramifiedQuotient.degreeModelMap_surjective}\\

130&def&\nolinkurl{PaperXI.RamifiedUnramifiedQuotient.quotientEquivUnramified}\\

136&theorem&\nolinkurl{PaperXI.RamifiedUnramifiedQuotient.quotientEquivUnramified_mk_degree}\\

145&theorem&\nolinkurl{PaperXI.RamifiedUnramifiedQuotient.degreeModelMap_eq_geometricLift_of_degree_one}\\

158&def&\nolinkurl{PaperXI.RamifiedUnramifiedQuotient.pullbackDegreeRepresentation}\\

164&theorem&\nolinkurl{PaperXI.RamifiedUnramifiedQuotient.inertia_le_pullbackDegreeRepresentation_ker}\\

175&def&\nolinkurl{PaperXI.RamifiedUnramifiedQuotient.quotientRepresentation}\\

183&def&\nolinkurl{PaperXI.RamifiedUnramifiedQuotient.descendRepresentation}\\

192&theorem&\nolinkurl{PaperXI.RamifiedUnramifiedQuotient.descendRepresentation_comp_degreeModelMap}\\

215&theorem&\nolinkurl{PaperXI.RamifiedUnramifiedQuotient.descendRepresentation_unique}\\

250&abbrev&\nolinkurl{PaperXI.RamifiedUnramifiedQuotient.degreeOnlyRepresentation}\\

257&theorem&\nolinkurl{PaperXI.RamifiedUnramifiedQuotient.degreeOnlyRepresentation_lift_eq_one}\\

271&theorem&\nolinkurl{PaperXI.RamifiedUnramifiedQuotient.degreeOnly_fixed_eq_top}\\

280&theorem&\nolinkurl{PaperXI.RamifiedUnramifiedQuotient.degreeOnly_depthProjector_eq_id}\\

292&theorem&\nolinkurl{PaperXI.RamifiedUnramifiedQuotient.degreeOnly_detailProjector_eq_zero}\\

301&theorem&\nolinkurl{PaperXI.RamifiedUnramifiedQuotient.degreeOnly_summand_eq_zero}\\

308&theorem&\nolinkurl{PaperXI.RamifiedUnramifiedQuotient.degreeOnly_swan_eq_zero}\\

314&theorem&\nolinkurl{PaperXI.RamifiedUnramifiedQuotient.degreeOnly_artin_eq_zero}\\

\bottomrule\end{longtable}

\section{\texorpdfstring{\nolinkurl{PaperXI.RankOneEulerFactor}}{PaperXI.RankOneEulerFactor}}\label{module:PaperXI.RankOneEulerFactor}

\textbf{Selected layer:} v0.44 parent.\quad\textbf{Source lines:} 113.\par\textbf{SHA-256:} {\scriptsize\nolinkurl{a512466f456be1c2a7b5c0f49b6b7cfad6bbda69ec1bc75283bb51b8140b95c3}}.

\VerbatimInput[fontsize=\scriptsize,breaklines=true,breakanywhere=true]{excerpts/D_RankOneEulerFactor.txt}

\textbf{Direct project imports:} \nolinkurl{PaperXI.UnramifiedConductorZero}.

\begin{longtable}{@{}p{0.07\textwidth}p{0.12\textwidth}p{0.75\textwidth}@{}}\toprule Line&Kind&Name\\\midrule\endhead

27&theorem&\nolinkurl{PaperXI.FiniteUnramifiedWeil.reverse_X_sub_C}\\

38&theorem&\nolinkurl{PaperXI.FiniteUnramifiedWeil.reverse_charpoly_smul_id}\\

59&def&\nolinkurl{PaperXI.FiniteUnramifiedWeil.frobeniusEigenvalue}\\

65&theorem&\nolinkurl{PaperXI.FiniteUnramifiedWeil.geometricFrobenius_action_eq_smul_id}\\

73&theorem&\nolinkurl{PaperXI.FiniteUnramifiedWeil.frobeniusEigenvalue_ne_zero}\\

92&theorem&\nolinkurl{PaperXI.FiniteUnramifiedWeil.frobeniusPolynomial_eq_linear}\\

100&def&\nolinkurl{PaperXI.FiniteUnramifiedWeil.rankOneEulerFactor}\\

105&theorem&\nolinkurl{PaperXI.FiniteUnramifiedWeil.unramifiedLocalFactor_eq_rankOneEulerFactor}\\

\bottomrule\end{longtable}

\section{\texorpdfstring{\nolinkurl{PaperXI.RankOneSourceCandidate110}}{PaperXI.RankOneSourceCandidate110}}\label{module:PaperXI.RankOneSourceCandidate110}

\textbf{Selected layer:} Continuation110.\quad\textbf{Source lines:} 77.\par\textbf{SHA-256:} {\scriptsize\nolinkurl{7f57ede141636d625efea3d9b67e905b1f34bcd5dd6c475f31a7e2ec2a7b2928}}.

\VerbatimInput[fontsize=\scriptsize,breaklines=true,breakanywhere=true]{excerpts/D_RankOneSourceCandidate110.txt}

\textbf{Direct project imports:} \nolinkurl{PaperXI.ArithmeticAttachment110}, \nolinkurl{PaperXI.NormalizerDictionary110}.

\begin{longtable}{@{}p{0.07\textwidth}p{0.12\textwidth}p{0.75\textwidth}@{}}\toprule Line&Kind&Name\\\midrule\endhead

29&theorem&\nolinkurl{PaperXI.RankOneSourceCandidate110.candidate_exists}\\

\bottomrule\end{longtable}

\section{\texorpdfstring{\nolinkurl{PaperXI.RegularCyclic}}{PaperXI.RegularCyclic}}\label{module:PaperXI.RegularCyclic}

\textbf{Selected layer:} v0.44 parent.\quad\textbf{Source lines:} 90.\par\textbf{SHA-256:} {\scriptsize\nolinkurl{2ed6fa5d0f002e4a087b67badeca9c39702efe1de6d6feade2a7eaf34fc5bccf}}.

\VerbatimInput[fontsize=\scriptsize,breaklines=true,breakanywhere=true]{excerpts/D_RegularCyclic.txt}

\textbf{Direct project imports:} \nolinkurl{PaperXI.TameProjectors}.

\begin{longtable}{@{}p{0.07\textwidth}p{0.12\textwidth}p{0.75\textwidth}@{}}\toprule Line&Kind&Name\\\midrule\endhead

12&def&\nolinkurl{PaperXI.RegularCyclic.translation}\\

17&def&\nolinkurl{PaperXI.RegularCyclic.regular}\\

30&theorem&\nolinkurl{PaperXI.RegularCyclic.operator_apply}\\

33&def&\nolinkurl{PaperXI.RegularCyclic.frobenius}\\

41&theorem&\nolinkurl{PaperXI.RegularCyclic.generator_relation}\\

50&theorem&\nolinkurl{PaperXI.RegularCyclic.character_vector_mem}\\

59&theorem&\nolinkurl{PaperXI.RegularCyclic.projector_nonzero}\\

67&theorem&\nolinkurl{PaperXI.RegularCyclic.projectors_distinct}\\

76&def&\nolinkurl{PaperXI.RegularCyclic.twoModFive}\\

79&theorem&\nolinkurl{PaperXI.RegularCyclic.wrong_forward_label}\\

\bottomrule\end{longtable}

\section{\texorpdfstring{\nolinkurl{PaperXI.RepresentationGraph}}{PaperXI.RepresentationGraph}}\label{module:PaperXI.RepresentationGraph}

\textbf{Selected layer:} v0.44 parent.\quad\textbf{Source lines:} 106.\par\textbf{SHA-256:} {\scriptsize\nolinkurl{5919599a901246a3def38a1fc9ae5480fabe1dc0a2f1c09be274d25317788d6d}}.

\VerbatimInput[fontsize=\scriptsize,breaklines=true,breakanywhere=true]{excerpts/D_RepresentationGraph.txt}

\textbf{Direct project imports:} \nolinkurl{PaperXI.GraphGauge}.

\begin{longtable}{@{}p{0.07\textwidth}p{0.12\textwidth}p{0.75\textwidth}@{}}\toprule Line&Kind&Name\\\midrule\endhead

19&def&\nolinkurl{PaperXI.RepresentationGraph.EdgeEquivariant}\\

25&def&\nolinkurl{PaperXI.RepresentationGraph.Intertwines}\\

34&theorem&\nolinkurl{PaperXI.RepresentationGraph.path_equivariant}\\

42&theorem&\nolinkurl{PaperXI.RepresentationGraph.rootGauge_intertwines}\\

55&theorem&\nolinkurl{PaperXI.RepresentationGraph.based_iff_chords}\\

71&theorem&\nolinkurl{PaperXI.RepresentationGraph.connected_iff_chords}\\

81&theorem&\nolinkurl{PaperXI.RepresentationGraph.based_iff_loops}\\

91&theorem&\nolinkurl{PaperXI.RepresentationGraph.unbased_iff_conjugates}\\

\bottomrule\end{longtable}

\section{\texorpdfstring{\nolinkurl{PaperXI.RepresentationStability}}{PaperXI.RepresentationStability}}\label{module:PaperXI.RepresentationStability}

\textbf{Selected layer:} v0.44 parent.\quad\textbf{Source lines:} 47.\par\textbf{SHA-256:} {\scriptsize\nolinkurl{ea48bcbf4e471e1d436ca33befd6c1aa3dcec5b7e765b15b407087b7c26cd33c}}.

\VerbatimInput[fontsize=\scriptsize,breaklines=true,breakanywhere=true]{excerpts/D_RepresentationStability.txt}

\textbf{Direct project imports:} \nolinkurl{PaperXI.GeneratedIrreducibility}.

\begin{longtable}{@{}p{0.07\textwidth}p{0.12\textwidth}p{0.75\textwidth}@{}}\toprule Line&Kind&Name\\\midrule\endhead

13&def&\nolinkurl{PaperXI.RepresentationStability.Stable}\\

17&theorem&\nolinkurl{PaperXI.RepresentationStability.eq_bot_or_top}\\

25&theorem&\nolinkurl{PaperXI.RepresentationStability.irreducible_of_stable}\\

36&theorem&\nolinkurl{PaperXI.RepresentationStability.irreducible_iff}\\

\bottomrule\end{longtable}

\section{\texorpdfstring{\nolinkurl{PaperXI.ResidueDegreeBaseChange}}{PaperXI.ResidueDegreeBaseChange}}\label{module:PaperXI.ResidueDegreeBaseChange}

\textbf{Selected layer:} v0.44 parent.\quad\textbf{Source lines:} 148.\par\textbf{SHA-256:} {\scriptsize\nolinkurl{c581b1a747917d5c92d930d8b817fd79fb5bf2c3de951682b786377850705bee}}.

\VerbatimInput[fontsize=\scriptsize,breaklines=true,breakanywhere=true]{excerpts/D_ResidueDegreeBaseChange.txt}

\textbf{Direct project imports:} \nolinkurl{PaperXI.ResidueWeil}.

\begin{longtable}{@{}p{0.07\textwidth}p{0.12\textwidth}p{0.75\textwidth}@{}}\toprule Line&Kind&Name\\\midrule\endhead

15&def&\nolinkurl{PaperXI.ResidueDegreeBaseChange.degreeMod}\\

24&def&\nolinkurl{PaperXI.ResidueDegreeBaseChange.periodSubgroup}\\

27&instance&\nolinkurl{PaperXI.ResidueDegreeBaseChange.period_normal}\\

30&theorem&\nolinkurl{PaperXI.ResidueDegreeBaseChange.mem_periodSubgroup}\\

36&def&\nolinkurl{PaperXI.ResidueDegreeBaseChange.inclusion}\\

50&theorem&\nolinkurl{PaperXI.ResidueDegreeBaseChange.inclusion_projection}\\

53&theorem&\nolinkurl{PaperXI.ResidueDegreeBaseChange.inclusion_degree}\\

57&theorem&\nolinkurl{PaperXI.ResidueDegreeBaseChange.inclusion_injective}\\

69&theorem&\nolinkurl{PaperXI.ResidueDegreeBaseChange.inclusion_range}\\

89&def&\nolinkurl{PaperXI.ResidueDegreeBaseChange.periodEquiv}\\

95&theorem&\nolinkurl{PaperXI.ResidueDegreeBaseChange.degree_surjective}\\

105&theorem&\nolinkurl{PaperXI.ResidueDegreeBaseChange.degreeMod_surjective}\\

116&def&\nolinkurl{PaperXI.ResidueDegreeBaseChange.quotientEquiv}\\

121&theorem&\nolinkurl{PaperXI.ResidueDegreeBaseChange.quotient_card}\\

127&theorem&\nolinkurl{PaperXI.ResidueDegreeBaseChange.inertia_preimage}\\

141&theorem&\nolinkurl{PaperXI.ResidueDegreeBaseChange.tower_degree}\\

\bottomrule\end{longtable}

\section{\texorpdfstring{\nolinkurl{PaperXI.ResidueExtensionRealization}}{PaperXI.ResidueExtensionRealization}}\label{module:PaperXI.ResidueExtensionRealization}

\textbf{Selected layer:} v0.44 parent.\quad\textbf{Source lines:} 97.\par\textbf{SHA-256:} {\scriptsize\nolinkurl{b24752381c6a6353645a53ffe2613d10e783ecf9086a2cdfb5876dba34bc9ccf}}.

\VerbatimInput[fontsize=\scriptsize,breaklines=true,breakanywhere=true]{excerpts/D_ResidueExtensionRealization.txt}

\textbf{Direct project imports:} \nolinkurl{PaperXI.UnramifiedPolynomial}.

\begin{longtable}{@{}p{0.07\textwidth}p{0.12\textwidth}p{0.75\textwidth}@{}}\toprule Line&Kind&Name\\\midrule\endhead

16&def&\nolinkurl{PaperXI.ResidueExtensionRealization.residueBasis}\\

19&def&\nolinkurl{PaperXI.ResidueExtensionRealization.residuePolynomial}\\

22&theorem&\nolinkurl{PaperXI.ResidueExtensionRealization.residuePolynomial_monic}\\

25&theorem&\nolinkurl{PaperXI.ResidueExtensionRealization.residuePolynomial_irreducible}\\

28&theorem&\nolinkurl{PaperXI.ResidueExtensionRealization.residuePolynomial_degree}\\

33&theorem&\nolinkurl{PaperXI.ResidueExtensionRealization.exists_monic_lift}\\

43&def&\nolinkurl{PaperXI.ResidueExtensionRealization.liftedPolynomial}\\

45&theorem&\nolinkurl{PaperXI.ResidueExtensionRealization.liftedPolynomial_map}\\

49&theorem&\nolinkurl{PaperXI.ResidueExtensionRealization.liftedPolynomial_degree}\\

53&instance&\nolinkurl{PaperXI.ResidueExtensionRealization.monic}\\

56&instance&\nolinkurl{PaperXI.ResidueExtensionRealization.residueIrreducible}\\

60&abbrev&\nolinkurl{PaperXI.ResidueExtensionRealization.ring}\\

61&abbrev&\nolinkurl{PaperXI.ResidueExtensionRealization.field}\\

64&def&\nolinkurl{PaperXI.ResidueExtensionRealization.residueEquiv}\\

72&theorem&\nolinkurl{PaperXI.ResidueExtensionRealization.rank}\\

76&theorem&\nolinkurl{PaperXI.ResidueExtensionRealization.field_degree}\\

82&theorem&\nolinkurl{PaperXI.ResidueExtensionRealization.etale}\\

85&theorem&\nolinkurl{PaperXI.ResidueExtensionRealization.ramificationIdx_one}\\

88&theorem&\nolinkurl{PaperXI.ResidueExtensionRealization.source}\\

\bottomrule\end{longtable}

\section{\texorpdfstring{\nolinkurl{PaperXI.ResidueFrobeniusAttachment}}{PaperXI.ResidueFrobeniusAttachment}}\label{module:PaperXI.ResidueFrobeniusAttachment}

\textbf{Selected layer:} v0.44 parent.\quad\textbf{Source lines:} 144.\par\textbf{SHA-256:} {\scriptsize\nolinkurl{f8bdcb103727ce29ab4fa9647716c21624920c75aff28d23c3a398cee0f59f7c}}.

\VerbatimInput[fontsize=\scriptsize,breaklines=true,breakanywhere=true]{excerpts/D_ResidueFrobeniusAttachment.txt}

\textbf{Direct project imports:} \nolinkurl{PaperXI.IntegerRingComparison}.

\begin{longtable}{@{}p{0.07\textwidth}p{0.12\textwidth}p{0.75\textwidth}@{}}\toprule Line&Kind&Name\\\midrule\endhead

23&abbrev&\nolinkurl{PaperXI.FiniteUnramifiedWeil.UpperField}\\

24&abbrev&\nolinkurl{PaperXI.FiniteUnramifiedWeil.BaseResidue}\\

25&abbrev&\nolinkurl{PaperXI.FiniteUnramifiedWeil.UpperResidue}\\

26&abbrev&\nolinkurl{PaperXI.FiniteUnramifiedWeil.ResidueGalois}\\

27&abbrev&\nolinkurl{PaperXI.FiniteUnramifiedWeil.ResidueDegree}\\

36&def&\nolinkurl{PaperXI.FiniteUnramifiedWeil.arithmeticFrobenius}\\

40&def&\nolinkurl{PaperXI.FiniteUnramifiedWeil.geometricFrobenius}\\

44&theorem&\nolinkurl{PaperXI.FiniteUnramifiedWeil.arithmeticFrobenius_apply}\\

49&theorem&\nolinkurl{PaperXI.FiniteUnramifiedWeil.geometricFrobenius_inverse_power}\\

55&theorem&\nolinkurl{PaperXI.FiniteUnramifiedWeil.arithmeticFrobenius_order}\\

61&theorem&\nolinkurl{PaperXI.FiniteUnramifiedWeil.geometricFrobenius_order}\\

67&theorem&\nolinkurl{PaperXI.FiniteUnramifiedWeil.arithmeticFrobenius_order_eq_fieldDegree}\\

73&theorem&\nolinkurl{PaperXI.FiniteUnramifiedWeil.geometricFrobenius_order_eq_fieldDegree}\\

79&theorem&\nolinkurl{PaperXI.FiniteUnramifiedWeil.natCard_residueGalois}\\

85&theorem&\nolinkurl{PaperXI.FiniteUnramifiedWeil.arithmeticFrobenius_mem_zpowers}\\

95&theorem&\nolinkurl{PaperXI.FiniteUnramifiedWeil.arithmeticFrobenius_zpowers_eq_top}\\

101&theorem&\nolinkurl{PaperXI.FiniteUnramifiedWeil.geometricFrobenius_zpowers_eq_top}\\

107&def&\nolinkurl{PaperXI.FiniteUnramifiedWeil.geometricCycleEquiv}\\

116&theorem&\nolinkurl{PaperXI.FiniteUnramifiedWeil.geometricCycleEquiv_one}\\

122&def&\nolinkurl{PaperXI.FiniteUnramifiedWeil.prescribedResidueGaloisEquiv}\\

127&theorem&\nolinkurl{PaperXI.FiniteUnramifiedWeil.prescribedResidueGaloisEquiv_arithmetic}\\

138&theorem&\nolinkurl{PaperXI.FiniteUnramifiedWeil.prescribedResidueGaloisEquiv_geometric}\\

\bottomrule\end{longtable}

\section{\texorpdfstring{\nolinkurl{PaperXI.ResidueRestriction107}}{PaperXI.ResidueRestriction107}}\label{module:PaperXI.ResidueRestriction107}

\textbf{Selected layer:} base108 arithmetic branch.\quad\textbf{Source lines:} 75.\par\textbf{SHA-256:} {\scriptsize\nolinkurl{d1008a38fcb166fa8167e9591321de444cc9d86d3ae176b13511e4639d0ab405}}.

\VerbatimInput[fontsize=\scriptsize,breaklines=true,breakanywhere=true]{excerpts/D_ResidueRestriction107.txt}

\textbf{Direct project imports:} \nolinkurl{PaperXI.ActualCuspidalCandidate105}.

\begin{longtable}{@{}p{0.07\textwidth}p{0.12\textwidth}p{0.75\textwidth}@{}}\toprule Line&Kind&Name\\\midrule\endhead

12&theorem&\nolinkurl{PaperXI.ResidueRestriction107.quotient_inclusion}\\

44&def&\nolinkurl{PaperXI.ResidueRestriction107.integralInclusion}\\

49&theorem&\nolinkurl{PaperXI.ResidueRestriction107.integralInclusion_coe}\\

56&theorem&\nolinkurl{PaperXI.ResidueRestriction107.reduction_integral}\\

\bottomrule\end{longtable}

\section{\texorpdfstring{\nolinkurl{PaperXI.ResidueWeil}}{PaperXI.ResidueWeil}}\label{module:PaperXI.ResidueWeil}

\textbf{Selected layer:} v0.44 parent.\quad\textbf{Source lines:} 114.\par\textbf{SHA-256:} {\scriptsize\nolinkurl{1eee4bb854ff6f99731ba944d6c0540ef7595df6fb8fa862ae030132fcd38166}}.

\VerbatimInput[fontsize=\scriptsize,breaklines=true,breakanywhere=true]{excerpts/D_ResidueWeil.txt}

\textbf{Direct project imports:} \nolinkurl{PaperXI.NormalInertia}.

\begin{longtable}{@{}p{0.07\textwidth}p{0.12\textwidth}p{0.75\textwidth}@{}}\toprule Line&Kind&Name\\\midrule\endhead

18&def&\nolinkurl{PaperXI.ResidueWeil.subgroup}\\

30&abbrev&\nolinkurl{PaperXI.ResidueWeil.Pullback}\\

32&def&\nolinkurl{PaperXI.ResidueWeil.projection}\\

35&def&\nolinkurl{PaperXI.ResidueWeil.degree}\\

38&theorem&\nolinkurl{PaperXI.ResidueWeil.projection_apply}\\

41&theorem&\nolinkurl{PaperXI.ResidueWeil.degree_apply}\\

45&def&\nolinkurl{PaperXI.ResidueWeil.inertiaEquiv}\\

65&instance&\nolinkurl{PaperXI.ResidueWeil.degreeQuotientCommutative}\\

73&theorem&\nolinkurl{PaperXI.ResidueWeil.exists_degree_one_iff}\\

90&abbrev&\nolinkurl{PaperXI.ResidueWeil.residueAction}\\

94&abbrev&\nolinkurl{PaperXI.ResidueWeil.valuationGroup}\\

98&def&\nolinkurl{PaperXI.ResidueWeil.valuationInertiaEquiv}\\

106&theorem&\nolinkurl{PaperXI.ResidueWeil.valuation_invariants_eq_bot}\\

\bottomrule\end{longtable}

\section{\texorpdfstring{\nolinkurl{PaperXI.RestrictedHeisenbergBridge}}{PaperXI.RestrictedHeisenbergBridge}}\label{module:PaperXI.RestrictedHeisenbergBridge}

\textbf{Selected layer:} v0.44 parent.\quad\textbf{Source lines:} 84.\par\textbf{SHA-256:} {\scriptsize\nolinkurl{4385767418858294f763e1ef3b8a60700679a6399d58874349e83ce8533e5af6}}.

\VerbatimInput[fontsize=\scriptsize,breaklines=true,breakanywhere=true]{excerpts/D_RestrictedHeisenbergBridge.txt}

\textbf{Direct project imports:} \nolinkurl{PaperXI.CentralExtensionComparison}, \nolinkurl{PaperXI.AbstractHeisenbergExtension}, \nolinkurl{PaperXI.SymplecticCentralExtension}.

\begin{longtable}{@{}p{0.07\textwidth}p{0.12\textwidth}p{0.75\textwidth}@{}}\toprule Line&Kind&Name\\\midrule\endhead

24&theorem&\nolinkurl{PaperXI.RestrictedHeisenbergBridge.odd_source}\\

58&theorem&\nolinkurl{PaperXI.RestrictedHeisenbergBridge.dyadic_source}\\

\bottomrule\end{longtable}

\section{\texorpdfstring{\nolinkurl{PaperXI.ScalarCocycle112}}{PaperXI.ScalarCocycle112}}\label{module:PaperXI.ScalarCocycle112}

\textbf{Selected layer:} Continuation112.\quad\textbf{Source lines:} 195.\par\textbf{SHA-256:} {\scriptsize\nolinkurl{351ac3d1e21afeb0a8897f6eea1ba99f83bb63eaeee390eebefa44c1dfbb0f4e}}.

\VerbatimInput[fontsize=\scriptsize,breaklines=true,breakanywhere=true]{excerpts/D_ScalarCocycle112.txt}

\textbf{Direct project imports:} none.

\begin{longtable}{@{}p{0.07\textwidth}p{0.12\textwidth}p{0.75\textwidth}@{}}\toprule Line&Kind&Name\\\midrule\endhead

16&structure&\nolinkurl{PaperXI.ScalarCocycle112.Cocycle}\\

23&def&\nolinkurl{PaperXI.ScalarCocycle112.commutator}\\

26&structure&\nolinkurl{PaperXI.ScalarCocycle112.Extension}\\

33&theorem&\nolinkurl{PaperXI.ScalarCocycle112.Extension.ext}\\

42&theorem&\nolinkurl{PaperXI.ScalarCocycle112.Extension.mul_scalar}\\

44&theorem&\nolinkurl{PaperXI.ScalarCocycle112.Extension.mul_base}\\

45&theorem&\nolinkurl{PaperXI.ScalarCocycle112.Extension.one_scalar}\\

46&theorem&\nolinkurl{PaperXI.ScalarCocycle112.Extension.one_base}\\

69&def&\nolinkurl{PaperXI.ScalarCocycle112.Extension.commGroup}\\

78&def&\nolinkurl{PaperXI.ScalarCocycle112.Extension.scalarInclusion}\\

83&theorem&\nolinkurl{PaperXI.ScalarCocycle112.Extension.scalarInclusion_scalar}\\

85&theorem&\nolinkurl{PaperXI.ScalarCocycle112.Extension.scalarInclusion_base}\\

88&theorem&\nolinkurl{PaperXI.ScalarCocycle112.Extension.scalarInclusion_injective}\\

92&def&\nolinkurl{PaperXI.ScalarCocycle112.Extension.sectionOf}\\

94&theorem&\nolinkurl{PaperXI.ScalarCocycle112.Extension.sectionOf_one}\\

96&theorem&\nolinkurl{PaperXI.ScalarCocycle112.Extension.sectionOf_mul}\\

104&theorem&\nolinkurl{PaperXI.ScalarCocycle112.symmetric_exists_cochain}\\

123&def&\nolinkurl{PaperXI.ScalarCocycle112.ratio}\\

131&theorem&\nolinkurl{PaperXI.ScalarCocycle112.ratio_symmetric}\\

141&theorem&\nolinkurl{PaperXI.ScalarCocycle112.same_commutator_exists_cochain_of_divisible}\\

150&theorem&\nolinkurl{PaperXI.ScalarCocycle112.complex_units_exists_pow}\\

162&def&\nolinkurl{PaperXI.ScalarCocycle112.complexUnitsDivisible}\\

171&theorem&\nolinkurl{PaperXI.ScalarCocycle112.same_commutator_exists_cochain}\\

181&theorem&\nolinkurl{PaperXI.ScalarCocycle112.same_commutator_exists_cochain_raw}\\

\bottomrule\end{longtable}

\section{\texorpdfstring{\nolinkurl{PaperXI.ScalarCovariance}}{PaperXI.ScalarCovariance}}\label{module:PaperXI.ScalarCovariance}

\textbf{Selected layer:} v0.44 parent.\quad\textbf{Source lines:} 52.\par\textbf{SHA-256:} {\scriptsize\nolinkurl{4859f25f802820daa9b4d9f45b627eb9f51c9be6756a872134890a6a94da4e06}}.

\VerbatimInput[fontsize=\scriptsize,breaklines=true,breakanywhere=true]{excerpts/D_ScalarCovariance.txt}

\textbf{Direct project imports:} \nolinkurl{PaperXI.MonodromyChain}.

\begin{longtable}{@{}p{0.07\textwidth}p{0.12\textwidth}p{0.75\textwidth}@{}}\toprule Line&Kind&Name\\\midrule\endhead

13&def&\nolinkurl{PaperXI.ScalarCovariance.frobenius}\\

25&theorem&\nolinkurl{PaperXI.ScalarCovariance.frobenius_apply}\\

29&theorem&\nolinkurl{PaperXI.ScalarCovariance.covariance}\\

41&theorem&\nolinkurl{PaperXI.ScalarCovariance.conjugation}\\

48&theorem&\nolinkurl{PaperXI.ScalarCovariance.scalar_model}\\

\bottomrule\end{longtable}

\section{\texorpdfstring{\nolinkurl{PaperXI.ScalarGluing}}{PaperXI.ScalarGluing}}\label{module:PaperXI.ScalarGluing}

\textbf{Selected layer:} v0.44 parent.\quad\textbf{Source lines:} 167.\par\textbf{SHA-256:} {\scriptsize\nolinkurl{f32ed302e2af8ac3b71237f1ac455195516c0932bb67a8d9f2c363d241b8321c}}.

\VerbatimInput[fontsize=\scriptsize,breaklines=true,breakanywhere=true]{excerpts/D_ScalarGluing.txt}

\textbf{Direct project imports:} \nolinkurl{PaperXI.GraphGauge}, \nolinkurl{PaperXI.ProjectiveHolonomy}.

\begin{longtable}{@{}p{0.07\textwidth}p{0.12\textwidth}p{0.75\textwidth}@{}}\toprule Line&Kind&Name\\\midrule\endhead

25&theorem&\nolinkurl{PaperXI.ScalarGluing.cycle_operator_iff}\\

37&theorem&\nolinkurl{PaperXI.ScalarGluing.cycle_independent}\\

46&theorem&\nolinkurl{PaperXI.ScalarGluing.based_iff_scalar_operators}\\

58&theorem&\nolinkurl{PaperXI.ScalarGluing.based_root_independent}\\

66&theorem&\nolinkurl{PaperXI.ScalarGluing.cycle_units_iff}\\

78&theorem&\nolinkurl{PaperXI.ScalarGluing.based_iff_scalar_units}\\

89&theorem&\nolinkurl{PaperXI.ScalarGluing.unbased_iff_scalar_units}\\

104&theorem&\nolinkurl{PaperXI.ScalarGluing.loops_iff_characters}\\

117&theorem&\nolinkurl{PaperXI.ScalarGluing.based_iff_ratioH1_trivial}\\

127&theorem&\nolinkurl{PaperXI.ScalarGluing.zero_root_gluing}\\

135&theorem&\nolinkurl{PaperXI.ScalarGluing.zero_finrank_gluing}\\

144&theorem&\nolinkurl{PaperXI.ScalarGluing.connected_scalar_criterion}\\

159&theorem&\nolinkurl{PaperXI.ScalarGluing.connected_zero_dimensional}\\

\bottomrule\end{longtable}

\section{\texorpdfstring{\nolinkurl{PaperXI.ScalarHolonomy}}{PaperXI.ScalarHolonomy}}\label{module:PaperXI.ScalarHolonomy}

\textbf{Selected layer:} v0.44 parent.\quad\textbf{Source lines:} 152.\par\textbf{SHA-256:} {\scriptsize\nolinkurl{f3309937e8a28dc3211c81306280c5e8a7c74c5c2cd944bfc884b0ad5e8cdb8a}}.

\VerbatimInput[fontsize=\scriptsize,breaklines=true,breakanywhere=true]{excerpts/D_ScalarHolonomy.txt}

\textbf{Direct project imports:} none.

\begin{longtable}{@{}p{0.07\textwidth}p{0.12\textwidth}p{0.75\textwidth}@{}}\toprule Line&Kind&Name\\\midrule\endhead

15&def&\nolinkurl{PaperXI.ScalarHolonomy.scalarHom}\\

20&theorem&\nolinkurl{PaperXI.ScalarHolonomy.scalarHom_apply}\\

24&theorem&\nolinkurl{PaperXI.ScalarHolonomy.scalar_conjugate}\\

30&theorem&\nolinkurl{PaperXI.ScalarHolonomy.scalar_conjugate_independent}\\

37&theorem&\nolinkurl{PaperXI.ScalarHolonomy.scalarHom_injective}\\

45&def&\nolinkurl{PaperXI.ScalarHolonomy.IsScalar}\\

49&def&\nolinkurl{PaperXI.ScalarHolonomy.character}\\

59&theorem&\nolinkurl{PaperXI.ScalarHolonomy.character_spec}\\

63&theorem&\nolinkurl{PaperXI.ScalarHolonomy.character_apply}\\

69&theorem&\nolinkurl{PaperXI.ScalarHolonomy.character_unique}\\

78&theorem&\nolinkurl{PaperXI.ScalarHolonomy.intertwines_iff}\\

92&theorem&\nolinkurl{PaperXI.ScalarHolonomy.character_intertwines_iff}\\

107&def&\nolinkurl{PaperXI.ScalarHolonomy.ratio}\\

109&theorem&\nolinkurl{PaperXI.ScalarHolonomy.ratio_apply}\\

112&theorem&\nolinkurl{PaperXI.ScalarHolonomy.ratio_eq_one_iff}\\

124&theorem&\nolinkurl{PaperXI.ScalarHolonomy.ratio_apply_eq_one_iff}\\

131&theorem&\nolinkurl{PaperXI.ScalarHolonomy.zero_space_operator_eq}\\

136&theorem&\nolinkurl{PaperXI.ScalarHolonomy.zero_space_intertwines}\\

141&theorem&\nolinkurl{PaperXI.ScalarHolonomy.zero_space_scalar_eq}\\

145&theorem&\nolinkurl{PaperXI.ScalarHolonomy.zero_space_complex_not_injective}\\

\bottomrule\end{longtable}

\section{\texorpdfstring{\nolinkurl{PaperXI.ScalarRepresentation112}}{PaperXI.ScalarRepresentation112}}\label{module:PaperXI.ScalarRepresentation112}

\textbf{Selected layer:} Continuation112.\quad\textbf{Source lines:} 373.\par\textbf{SHA-256:} {\scriptsize\nolinkurl{f8bad889bfb84cf03242cd2fd13dbd620d15080c126ade4b5f9bbfdb2ae7807c}}.

\VerbatimInput[fontsize=\scriptsize,breaklines=true,breakanywhere=true]{excerpts/D_ScalarRepresentation112.txt}

\textbf{Direct project imports:} \nolinkurl{PaperXI.GlobalHeisenbergBridge112}, \nolinkurl{PaperXI.RepresentationStability}, \nolinkurl{PaperXI.ScalarCocycle112}.

\begin{longtable}{@{}p{0.07\textwidth}p{0.12\textwidth}p{0.75\textwidth}@{}}\toprule Line&Kind&Name\\\midrule\endhead

23&def&\nolinkurl{PaperXI.ScalarRepresentation112.sectionResidue}\\

30&def&\nolinkurl{PaperXI.ScalarRepresentation112.sectionScalar}\\

33&def&\nolinkurl{PaperXI.ScalarRepresentation112.sectionCocycle}\\

36&theorem&\nolinkurl{PaperXI.ScalarRepresentation112.sectionResidue_mul}\\

47&theorem&\nolinkurl{PaperXI.ScalarRepresentation112.sectionScalar_mul}\\

56&theorem&\nolinkurl{PaperXI.ScalarRepresentation112.sectionResidue_section}\\

61&theorem&\nolinkurl{PaperXI.ScalarRepresentation112.sectionScalar_section}\\

66&theorem&\nolinkurl{PaperXI.ScalarRepresentation112.sectionScalar_one}\\

71&theorem&\nolinkurl{PaperXI.ScalarRepresentation112.sectionCocycle_one_left}\\

76&theorem&\nolinkurl{PaperXI.ScalarRepresentation112.sectionCocycle_one_right}\\

82&theorem&\nolinkurl{PaperXI.ScalarRepresentation112.sectionCocycle_assoc}\\

91&theorem&\nolinkurl{PaperXI.ScalarRepresentation112.sectionScalar_on_subgroup}\\

111&def&\nolinkurl{PaperXI.ScalarRepresentation112.packetOperator}\\

117&theorem&\nolinkurl{PaperXI.ScalarRepresentation112.packetOperator_mul}\\

150&def&\nolinkurl{PaperXI.ScalarRepresentation112.scalarTransport}\\

178&theorem&\nolinkurl{PaperXI.ScalarRepresentation112.scalarTransport_on_subgroup}\\

193&theorem&\nolinkurl{PaperXI.ScalarRepresentation112.scalarTransport_section}\\

204&theorem&\nolinkurl{PaperXI.ScalarRepresentation112.packet_normalForm}\\

219&theorem&\nolinkurl{PaperXI.ScalarRepresentation112.scalarTransport_irreducible}\\

238&theorem&\nolinkurl{PaperXI.ScalarRepresentation112.scalarTransport_contains}\\

254&def&\nolinkurl{PaperXI.ScalarRepresentation112.actualScalarCocycle}\\

266&def&\nolinkurl{PaperXI.ScalarRepresentation112.sectionCommutator}\\

277&theorem&\nolinkurl{PaperXI.ScalarRepresentation112.sectionCocycle_commutator}\\

297&theorem&\nolinkurl{PaperXI.ScalarRepresentation112.exists_actual_action_of_equal_scalar_commutators}\\

334&theorem&\nolinkurl{PaperXI.ScalarRepresentation112.bridge_to_prescribed_eta_of_equal_commutators}\\

\bottomrule\end{longtable}

\section{\texorpdfstring{\nolinkurl{PaperXI.SchrodingerConsequences}}{PaperXI.SchrodingerConsequences}}\label{module:PaperXI.SchrodingerConsequences}

\textbf{Selected layer:} v0.44 parent.\quad\textbf{Source lines:} 151.\par\textbf{SHA-256:} {\scriptsize\nolinkurl{f0f3bb3edc8a1eeb75ebeb2d70adc089b5e324fb2647841978788425a34c3c19}}.

\VerbatimInput[fontsize=\scriptsize,breaklines=true,breakanywhere=true]{excerpts/D_SchrodingerConsequences.txt}

\textbf{Direct project imports:} \nolinkurl{PaperXI.SchrodingerUniversal}, \nolinkurl{PaperXI.HeisenbergSelector}.

\begin{longtable}{@{}p{0.07\textwidth}p{0.12\textwidth}p{0.75\textwidth}@{}}\toprule Line&Kind&Name\\\midrule\endhead

15&def&\nolinkurl{PaperXI.SchrodingerConsequences.restriction}\\

21&def&\nolinkurl{PaperXI.SchrodingerConsequences.projector}\\

26&theorem&\nolinkurl{PaperXI.SchrodingerConsequences.projector_idempotent}\\

33&theorem&\nolinkurl{PaperXI.SchrodingerConsequences.range_eq_invariants}\\

41&def&\nolinkurl{PaperXI.SchrodingerConsequences.subgroupProjector}\\

47&theorem&\nolinkurl{PaperXI.SchrodingerConsequences.projector_eq_subgroupProjector}\\

68&theorem&\nolinkurl{PaperXI.SchrodingerConsequences.conjugates}\\

79&theorem&\nolinkurl{PaperXI.SchrodingerConsequences.trace_eq_model}\\

85&theorem&\nolinkurl{PaperXI.SchrodingerConsequences.character_support}\\

93&theorem&\nolinkurl{PaperXI.SchrodingerConsequences.conjugates_projector}\\

106&theorem&\nolinkurl{PaperXI.SchrodingerConsequences.trace_projector}\\

114&theorem&\nolinkurl{PaperXI.SchrodingerConsequences.rank_projector}\\

124&theorem&\nolinkurl{PaperXI.SchrodingerConsequences.maximal_trace_one}\\

132&theorem&\nolinkurl{PaperXI.SchrodingerConsequences.maximal_rank_one}\\

144&theorem&\nolinkurl{PaperXI.SchrodingerConsequences.maximal_fixed_line}\\

\bottomrule\end{longtable}

\section{\texorpdfstring{\nolinkurl{PaperXI.SchrodingerFourier}}{PaperXI.SchrodingerFourier}}\label{module:PaperXI.SchrodingerFourier}

\textbf{Selected layer:} v0.44 parent.\quad\textbf{Source lines:} 77.\par\textbf{SHA-256:} {\scriptsize\nolinkurl{1c9a3b9a2264e75afda601089ef6b0665e2e1f3533f7146575d747946f361125}}.

\VerbatimInput[fontsize=\scriptsize,breaklines=true,breakanywhere=true]{excerpts/D_SchrodingerFourier.txt}

\textbf{Direct project imports:} none.

\begin{longtable}{@{}p{0.07\textwidth}p{0.12\textwidth}p{0.75\textwidth}@{}}\toprule Line&Kind&Name\\\midrule\endhead

16&theorem&\nolinkurl{PaperXI.SchrodingerFourier.finite_elimination}\\

49&theorem&\nolinkurl{PaperXI.SchrodingerFourier.coordinate_projection_mem}\\

65&theorem&\nolinkurl{PaperXI.SchrodingerFourier.coordinate_unit_mem}\\

\bottomrule\end{longtable}

\section{\texorpdfstring{\nolinkurl{PaperXI.SchrodingerIrreducible}}{PaperXI.SchrodingerIrreducible}}\label{module:PaperXI.SchrodingerIrreducible}

\textbf{Selected layer:} v0.44 parent.\quad\textbf{Source lines:} 88.\par\textbf{SHA-256:} {\scriptsize\nolinkurl{b45be66566d924a65606b03dc54678b2bb795070e1099de79205c24a0bff555c}}.

\VerbatimInput[fontsize=\scriptsize,breaklines=true,breakanywhere=true]{excerpts/D_SchrodingerIrreducible.txt}

\textbf{Direct project imports:} \nolinkurl{PaperXI.SchrodingerFourier}, \nolinkurl{PaperXI.SchrodingerSeparation}.

\begin{longtable}{@{}p{0.07\textwidth}p{0.12\textwidth}p{0.75\textwidth}@{}}\toprule Line&Kind&Name\\\midrule\endhead

16&theorem&\nolinkurl{PaperXI.SchrodingerIrreducible.invariant_eq_bot_or_top}\\

61&theorem&\nolinkurl{PaperXI.SchrodingerIrreducible.irreducible_of_actions}\\

\bottomrule\end{longtable}

\section{\texorpdfstring{\nolinkurl{PaperXI.SchrodingerOperator}}{PaperXI.SchrodingerOperator}}\label{module:PaperXI.SchrodingerOperator}

\textbf{Selected layer:} v0.44 parent.\quad\textbf{Source lines:} 93.\par\textbf{SHA-256:} {\scriptsize\nolinkurl{2d1478c4d0de8e6ea366020e41530e8967ad47f25516f5cd4d8cbb544d08616d}}.

\VerbatimInput[fontsize=\scriptsize,breaklines=true,breakanywhere=true]{excerpts/D_SchrodingerOperator.txt}

\textbf{Direct project imports:} none.

\begin{longtable}{@{}p{0.07\textwidth}p{0.12\textwidth}p{0.75\textwidth}@{}}\toprule Line&Kind&Name\\\midrule\endhead

13&def&\nolinkurl{PaperXI.Schrodinger.operator}\\

25&theorem&\nolinkurl{PaperXI.Schrodinger.operator_apply}\\

29&def&\nolinkurl{PaperXI.Schrodinger.translation}\\

34&theorem&\nolinkurl{PaperXI.Schrodinger.translation_apply}\\

37&def&\nolinkurl{PaperXI.Schrodinger.modulation}\\

49&theorem&\nolinkurl{PaperXI.Schrodinger.modulation_apply}\\

52&theorem&\nolinkurl{PaperXI.Schrodinger.operator_translation}\\

57&theorem&\nolinkurl{PaperXI.Schrodinger.operator_modulation}\\

62&theorem&\nolinkurl{PaperXI.Schrodinger.operator_central}\\

67&theorem&\nolinkurl{PaperXI.Schrodinger.operator_one}\\

73&theorem&\nolinkurl{PaperXI.Schrodinger.phase_add}\\

83&theorem&\nolinkurl{PaperXI.Schrodinger.operator_mul}\\

\bottomrule\end{longtable}

\section{\texorpdfstring{\nolinkurl{PaperXI.SchrodingerSeparation}}{PaperXI.SchrodingerSeparation}}\label{module:PaperXI.SchrodingerSeparation}

\textbf{Selected layer:} v0.44 parent.\quad\textbf{Source lines:} 32.\par\textbf{SHA-256:} {\scriptsize\nolinkurl{ea554103afc7f9badccc857e5421b5b4749f1bea3df4043f0aae400783cb8b8b}}.

\VerbatimInput[fontsize=\scriptsize,breaklines=true,breakanywhere=true]{excerpts/D_SchrodingerSeparation.txt}

\textbf{Direct project imports:} none.

\begin{longtable}{@{}p{0.07\textwidth}p{0.12\textwidth}p{0.75\textwidth}@{}}\toprule Line&Kind&Name\\\midrule\endhead

15&theorem&\nolinkurl{PaperXI.SchrodingerSeparation.exists_dual_character_ne}\\

\bottomrule\end{longtable}

\section{\texorpdfstring{\nolinkurl{PaperXI.SchrodingerTrace}}{PaperXI.SchrodingerTrace}}\label{module:PaperXI.SchrodingerTrace}

\textbf{Selected layer:} v0.44 parent.\quad\textbf{Source lines:} 88.\par\textbf{SHA-256:} {\scriptsize\nolinkurl{0c523e1824a16155d9072d81a47308e828eef34cdf178595eff2ccbd8dbe7498}}.

\VerbatimInput[fontsize=\scriptsize,breaklines=true,breakanywhere=true]{excerpts/D_SchrodingerTrace.txt}

\textbf{Direct project imports:} \nolinkurl{PaperXI.SchrodingerOperator}.

\begin{longtable}{@{}p{0.07\textwidth}p{0.12\textwidth}p{0.75\textwidth}@{}}\toprule Line&Kind&Name\\\midrule\endhead

16&theorem&\nolinkurl{PaperXI.SchrodingerTrace.functional_surjective}\\

25&theorem&\nolinkurl{PaperXI.SchrodingerTrace.functional_character_ne_one}\\

38&theorem&\nolinkurl{PaperXI.SchrodingerTrace.sum_functional_character}\\

43&theorem&\nolinkurl{PaperXI.SchrodingerTrace.operator_diagonal}\\

58&theorem&\nolinkurl{PaperXI.SchrodingerTrace.trace_of_translation_ne_zero}\\

66&theorem&\nolinkurl{PaperXI.SchrodingerTrace.trace_of_translation_zero}\\

74&theorem&\nolinkurl{PaperXI.SchrodingerTrace.character_support}\\

\bottomrule\end{longtable}

\section{\texorpdfstring{\nolinkurl{PaperXI.SchrodingerUnitary}}{PaperXI.SchrodingerUnitary}}\label{module:PaperXI.SchrodingerUnitary}

\textbf{Selected layer:} v0.44 parent.\quad\textbf{Source lines:} 45.\par\textbf{SHA-256:} {\scriptsize\nolinkurl{90bc68223666d16ea0fd539caf1719dc4d7a73e4774eadd64d5136507b4e40cb}}.

\VerbatimInput[fontsize=\scriptsize,breaklines=true,breakanywhere=true]{excerpts/D_SchrodingerUnitary.txt}

\textbf{Direct project imports:} \nolinkurl{PaperXI.SchrodingerOperator}.

\begin{longtable}{@{}p{0.07\textwidth}p{0.12\textwidth}p{0.75\textwidth}@{}}\toprule Line&Kind&Name\\\midrule\endhead

15&theorem&\nolinkurl{PaperXI.SchrodingerUnitary.conjugate_character}\\

21&theorem&\nolinkurl{PaperXI.SchrodingerUnitary.conjugate_mul_character}\\

29&theorem&\nolinkurl{PaperXI.SchrodingerUnitary.pairing_invariant}\\

\bottomrule\end{longtable}

\section{\texorpdfstring{\nolinkurl{PaperXI.SchrodingerUniversal}}{PaperXI.SchrodingerUniversal}}\label{module:PaperXI.SchrodingerUniversal}

\textbf{Selected layer:} v0.44 parent.\quad\textbf{Source lines:} 217.\par\textbf{SHA-256:} {\scriptsize\nolinkurl{3a45fcb361ce2045666670d9fe67cd2b7406dc7840b50f9a9cbe11c00f77c608}}.

\VerbatimInput[fontsize=\scriptsize,breaklines=true,breakanywhere=true]{excerpts/D_SchrodingerUniversal.txt}

\textbf{Direct project imports:} \nolinkurl{PaperXI.HeisenbergNormalForm}, \nolinkurl{PaperXI.HeisenbergFrequency}, \nolinkurl{PaperXI.HeisenbergDimension}, \nolinkurl{PaperXI.HeisenbergRepresentation}, \nolinkurl{PaperXI.SchurScalar}.

\begin{longtable}{@{}p{0.07\textwidth}p{0.12\textwidth}p{0.75\textwidth}@{}}\toprule Line&Kind&Name\\\midrule\endhead

18&def&\nolinkurl{PaperXI.SchrodingerUniversal.basis_map}\\

22&theorem&\nolinkurl{PaperXI.SchrodingerUniversal.basis_map_basis}\\

28&theorem&\nolinkurl{PaperXI.SchrodingerUniversal.translation_basis}\\

34&theorem&\nolinkurl{PaperXI.SchrodingerUniversal.modulation_basis}\\

43&theorem&\nolinkurl{PaperXI.SchrodingerUniversal.basis_map_translation}\\

56&theorem&\nolinkurl{PaperXI.SchrodingerUniversal.basis_map_modulation}\\

72&theorem&\nolinkurl{PaperXI.SchrodingerUniversal.basis_map_central}\\

83&theorem&\nolinkurl{PaperXI.SchrodingerUniversal.intertwines}\\

103&def&\nolinkurl{PaperXI.SchrodingerUniversal.intertwiningMap}\\

111&theorem&\nolinkurl{PaperXI.SchrodingerUniversal.intertwiningMap_ne_zero}\\

122&def&\nolinkurl{PaperXI.SchrodingerUniversal.equiv_of_fixed}\\

138&theorem&\nolinkurl{PaperXI.SchrodingerUniversal.exists_equiv}\\

148&def&\nolinkurl{PaperXI.SchrodingerUniversal.equiv}\\

155&theorem&\nolinkurl{PaperXI.SchrodingerUniversal.equiv_intertwines}\\

166&theorem&\nolinkurl{PaperXI.SchrodingerUniversal.finiteDimensional}\\

174&theorem&\nolinkurl{PaperXI.SchrodingerUniversal.finrank_eq}\\

182&theorem&\nolinkurl{PaperXI.SchrodingerUniversal.finrank_eq_sqrt}\\

191&theorem&\nolinkurl{PaperXI.SchrodingerUniversal.unique_scalar}\\

206&theorem&\nolinkurl{PaperXI.SchrodingerUniversal.exists_equiv_finiteField}\\

\bottomrule\end{longtable}

\section{\texorpdfstring{\nolinkurl{PaperXI.SchurHom}}{PaperXI.SchurHom}}\label{module:PaperXI.SchurHom}

\textbf{Selected layer:} v0.44 parent.\quad\textbf{Source lines:} 68.\par\textbf{SHA-256:} {\scriptsize\nolinkurl{b86b71d721ec121989b4959337263985da723a313f102af82e8b05a8775c35bf}}.

\VerbatimInput[fontsize=\scriptsize,breaklines=true,breakanywhere=true]{excerpts/D_SchurHom.txt}

\textbf{Direct project imports:} \nolinkurl{PaperXI.DualityLinear}, \nolinkurl{PaperXI.SchurScalar}.

\begin{longtable}{@{}p{0.07\textwidth}p{0.12\textwidth}p{0.75\textwidth}@{}}\toprule Line&Kind&Name\\\midrule\endhead

14&def&\nolinkurl{PaperXI.SchurHom.postcomposeEquiv}\\

25&theorem&\nolinkurl{PaperXI.SchurHom.finrank_eq_indicator}\\

40&theorem&\nolinkurl{PaperXI.SchurHom.finrank_zero_or_one}\\

49&theorem&\nolinkurl{PaperXI.SchurHom.finrank_one_iff_equiv}\\

59&theorem&\nolinkurl{PaperXI.SchurHom.tensor_invariants_zero_or_one}\\

\bottomrule\end{longtable}

\section{\texorpdfstring{\nolinkurl{PaperXI.SchurScalar}}{PaperXI.SchurScalar}}\label{module:PaperXI.SchurScalar}

\textbf{Selected layer:} v0.44 parent.\quad\textbf{Source lines:} 66.\par\textbf{SHA-256:} {\scriptsize\nolinkurl{ab78b7e45ea049accb02eacca96ad12dd74812d23ea41c9186645a9b5700ee92}}.

\VerbatimInput[fontsize=\scriptsize,breaklines=true,breakanywhere=true]{excerpts/D_SchurScalar.txt}

\textbf{Direct project imports:} \nolinkurl{PaperXI.GeneratedIrreducibility}.

\begin{longtable}{@{}p{0.07\textwidth}p{0.12\textwidth}p{0.75\textwidth}@{}}\toprule Line&Kind&Name\\\midrule\endhead

13&theorem&\nolinkurl{PaperXI.SchurScalar.commuting_scalar}\\

25&theorem&\nolinkurl{PaperXI.SchurScalar.scalar_unique}\\

33&theorem&\nolinkurl{PaperXI.SchurScalar.commuting_unit}\\

46&theorem&\nolinkurl{PaperXI.SchurScalar.bridge_unique}\\

\bottomrule\end{longtable}

\section{\texorpdfstring{\nolinkurl{PaperXI.SemisimpleDetection}}{PaperXI.SemisimpleDetection}}\label{module:PaperXI.SemisimpleDetection}

\textbf{Selected layer:} v0.44 parent.\quad\textbf{Source lines:} 71.\par\textbf{SHA-256:} {\scriptsize\nolinkurl{b2713fc088d1ea21bffebfd8dd05a2b93962db392825add54d953db52b00c8c8}}.

\VerbatimInput[fontsize=\scriptsize,breaklines=true,breakanywhere=true]{excerpts/D_SemisimpleDetection.txt}

\textbf{Direct project imports:} \nolinkurl{PaperXI.CharacterProjector}.

\begin{longtable}{@{}p{0.07\textwidth}p{0.12\textwidth}p{0.75\textwidth}@{}}\toprule Line&Kind&Name\\\midrule\endhead

14&theorem&\nolinkurl{PaperXI.SemisimpleDetection.eq_zero_of_simple_actions}\\

41&theorem&\nolinkurl{PaperXI.SemisimpleDetection.eq_zero_of_irreducible_actions}\\

62&theorem&\nolinkurl{PaperXI.SemisimpleDetection.eq_of_irreducible_actions}\\

\bottomrule\end{longtable}

\section{\texorpdfstring{\nolinkurl{PaperXI.SemisimpleProjection}}{PaperXI.SemisimpleProjection}}\label{module:PaperXI.SemisimpleProjection}

\textbf{Selected layer:} base108 arithmetic branch.\quad\textbf{Source lines:} 99.\par\textbf{SHA-256:} {\scriptsize\nolinkurl{dab8a3dd87a59b725c83dd9c32b853b2de9db98afc7461311017ed27ec4b21ea}}.

\VerbatimInput[fontsize=\scriptsize,breaklines=true,breakanywhere=true]{excerpts/D_SemisimpleProjection.txt}

\textbf{Direct project imports:} \nolinkurl{PaperXI.RepresentationStability}.

\begin{longtable}{@{}p{0.07\textwidth}p{0.12\textwidth}p{0.75\textwidth}@{}}\toprule Line&Kind&Name\\\midrule\endhead

19&theorem&\nolinkurl{PaperXI.SemisimpleProjection.module_range_le_of_nonzero}\\

43&theorem&\nolinkurl{PaperXI.SemisimpleProjection.module_projection_mem}\\

66&theorem&\nolinkurl{PaperXI.SemisimpleProjection.range_le_of_nonzero}\\

87&theorem&\nolinkurl{PaperXI.SemisimpleProjection.projection_mem}\\

\bottomrule\end{longtable}

\section{\texorpdfstring{\nolinkurl{PaperXI.SimpleModuleClassification}}{PaperXI.SimpleModuleClassification}}\label{module:PaperXI.SimpleModuleClassification}

\textbf{Selected layer:} v0.44 parent.\quad\textbf{Source lines:} 106.\par\textbf{SHA-256:} {\scriptsize\nolinkurl{b28f26a821af96bffc99660b774cfedfbafa942365b515c8a346a771ce3b4bf3}}.

\VerbatimInput[fontsize=\scriptsize,breaklines=true,breakanywhere=true]{excerpts/D_SimpleModuleClassification.txt}

\textbf{Direct project imports:} none.

\begin{longtable}{@{}p{0.07\textwidth}p{0.12\textwidth}p{0.75\textwidth}@{}}\toprule Line&Kind&Name\\\midrule\endhead

17&theorem&\nolinkurl{PaperXI.SimpleModuleClassification.equivalent}\\

27&theorem&\nolinkurl{PaperXI.SimpleModuleClassification.regular_decomposition}\\

37&theorem&\nolinkurl{PaperXI.SimpleModuleClassification.exists_simple_ideal}\\

46&theorem&\nolinkurl{PaperXI.SimpleModuleClassification.equivalent_complex}\\

55&theorem&\nolinkurl{PaperXI.SimpleModuleClassification.finiteDimensional}\\

64&theorem&\nolinkurl{PaperXI.SimpleModuleClassification.finrank_pos}\\

70&def&\nolinkurl{PaperXI.SimpleModuleClassification.scalarEndEquiv}\\

76&theorem&\nolinkurl{PaperXI.SimpleModuleClassification.finrank_end}\\

81&theorem&\nolinkurl{PaperXI.SimpleModuleClassification.finrank_sq}\\

103&theorem&\nolinkurl{PaperXI.SimpleModuleClassification.finrank_sqrt}\\

\bottomrule\end{longtable}

\section{\texorpdfstring{\nolinkurl{PaperXI.SmithMatrix106}}{PaperXI.SmithMatrix106}}\label{module:PaperXI.SmithMatrix106}

\textbf{Selected layer:} base108 arithmetic branch.\quad\textbf{Source lines:} 57.\par\textbf{SHA-256:} {\scriptsize\nolinkurl{9104a6467adaf608eeb4b3700c5d827828b603184ed83c55516445df5f23a948}}.

\VerbatimInput[fontsize=\scriptsize,breaklines=true,breakanywhere=true]{excerpts/D_SmithMatrix106.txt}

\textbf{Direct project imports:} none.

\begin{longtable}{@{}p{0.07\textwidth}p{0.12\textwidth}p{0.75\textwidth}@{}}\toprule Line&Kind&Name\\\midrule\endhead

16&def&\nolinkurl{PaperXI.SmithMatrix106.basisGL}\\

23&theorem&\nolinkurl{PaperXI.SmithMatrix106.basisGL_apply}\\

27&theorem&\nolinkurl{PaperXI.SmithMatrix106.exists_diagonal}\\

\bottomrule\end{longtable}

\section{\texorpdfstring{\nolinkurl{PaperXI.SmoothDegreeQuotient}}{PaperXI.SmoothDegreeQuotient}}\label{module:PaperXI.SmoothDegreeQuotient}

\textbf{Selected layer:} base108 arithmetic branch.\quad\textbf{Source lines:} 192.\par\textbf{SHA-256:} {\scriptsize\nolinkurl{77138b22ce63047224ba2ad5c14fc85e25acce77ad9ea79b7e3404e4fb684a3d}}.

\VerbatimInput[fontsize=\scriptsize,breaklines=true,breakanywhere=true]{excerpts/D_SmoothDegreeQuotient.txt}

\textbf{Direct project imports:} \nolinkurl{PaperXI.QuotientRepresentation}, \nolinkurl{PaperXI.TameDegree}.

\begin{longtable}{@{}p{0.07\textwidth}p{0.12\textwidth}p{0.75\textwidth}@{}}\toprule Line&Kind&Name\\\midrule\endhead

21&def&\nolinkurl{PaperXI.SmoothDegreeQuotient.vectorStabilizer}\\

33&theorem&\nolinkurl{PaperXI.SmoothDegreeQuotient.mem_vectorStabilizer}\\

40&def&\nolinkurl{PaperXI.SmoothDegreeQuotient.IsSmooth}\\

44&theorem&\nolinkurl{PaperXI.SmoothDegreeQuotient.isOpen_kernel_of_smooth}\\

67&def&\nolinkurl{PaperXI.SmoothDegreeQuotient.jointKernel}\\

69&instance&\nolinkurl{PaperXI.SmoothDegreeQuotient.jointKernel_normal}\\

74&abbrev&\nolinkurl{PaperXI.SmoothDegreeQuotient.DegreeQuotient}\\

77&abbrev&\nolinkurl{PaperXI.SmoothDegreeQuotient.quotientMap}\\

81&def&\nolinkurl{PaperXI.SmoothDegreeQuotient.descendedDegree}\\

85&def&\nolinkurl{PaperXI.SmoothDegreeQuotient.descendedRepresentation}\\

88&theorem&\nolinkurl{PaperXI.SmoothDegreeQuotient.descendedDegree_mk}\\

91&theorem&\nolinkurl{PaperXI.SmoothDegreeQuotient.descendedRepresentation_mk}\\

94&theorem&\nolinkurl{PaperXI.SmoothDegreeQuotient.descendedDegree_comp}\\

97&theorem&\nolinkurl{PaperXI.SmoothDegreeQuotient.descendedRepresentation_comp}\\

101&theorem&\nolinkurl{PaperXI.SmoothDegreeQuotient.descended_irreducible}\\

106&theorem&\nolinkurl{PaperXI.SmoothDegreeQuotient.degree_surjective_iff}\\

118&theorem&\nolinkurl{PaperXI.SmoothDegreeQuotient.degreeSubgroup_comap}\\

126&theorem&\nolinkurl{PaperXI.SmoothDegreeQuotient.jointKernel_descended}\\

135&theorem&\nolinkurl{PaperXI.SmoothDegreeQuotient.degree_descends_iff}\\

150&theorem&\nolinkurl{PaperXI.SmoothDegreeQuotient.inertia_image}\\

165&theorem&\nolinkurl{PaperXI.SmoothDegreeQuotient.isOpen_jointKernel}\\

171&theorem&\nolinkurl{PaperXI.SmoothDegreeQuotient.quotient_discrete}\\

177&theorem&\nolinkurl{PaperXI.SmoothDegreeQuotient.quotientMap_continuous}\\

182&theorem&\nolinkurl{PaperXI.SmoothDegreeQuotient.finite_descended_inertia}\\

\bottomrule\end{longtable}

\section{\texorpdfstring{\nolinkurl{PaperXI.SmoothHomOrbit68}}{PaperXI.SmoothHomOrbit68}}\label{module:PaperXI.SmoothHomOrbit68}

\textbf{Selected layer:} base108 arithmetic branch.\quad\textbf{Source lines:} 95.\par\textbf{SHA-256:} {\scriptsize\nolinkurl{87d24ea9fbb25f0394cace54953a216a8158ddded33b1f0e68f620c6a0e5c2d6}}.

\VerbatimInput[fontsize=\scriptsize,breaklines=true,breakanywhere=true]{excerpts/D_SmoothHomOrbit68.txt}

\textbf{Direct project imports:} \nolinkurl{PaperXI.CompactSmoothSemisimple}.

\begin{longtable}{@{}p{0.07\textwidth}p{0.12\textwidth}p{0.75\textwidth}@{}}\toprule Line&Kind&Name\\\midrule\endhead

22&theorem&\nolinkurl{PaperXI.SmoothHomOrbit68.linHom_smooth}\\

48&theorem&\nolinkurl{PaperXI.SmoothHomOrbit68.descend_smooth}\\

56&theorem&\nolinkurl{PaperXI.SmoothHomOrbit68.orbit_finite_of_compact_quotient}\\

69&theorem&\nolinkurl{PaperXI.SmoothHomOrbit68.scalar_subgroup_le_linHom_kernel}\\

84&theorem&\nolinkurl{PaperXI.SmoothHomOrbit68.linHom_orbit_finite}\\

\bottomrule\end{longtable}

\section{\texorpdfstring{\nolinkurl{PaperXI.SpectralGeometricFrobenius}}{PaperXI.SpectralGeometricFrobenius}}\label{module:PaperXI.SpectralGeometricFrobenius}

\textbf{Selected layer:} normalized Frobenius.\quad\textbf{Source lines:} 184.\par\textbf{SHA-256:} {\scriptsize\nolinkurl{e437a81f9af9b976a3c272dbb0dbb4824e47f138852a60adbb5dc5ac9236685e}}.

\VerbatimInput[fontsize=\scriptsize,breaklines=true,breakanywhere=true]{excerpts/D_SpectralGeometricFrobenius.txt}

\textbf{Direct project imports:} \nolinkurl{PaperXI.HenselianGaloisResidue}, \nolinkurl{PaperXI.HenselianLiftTower}, \nolinkurl{PaperXI.FiniteResidueBaseChange}.

\begin{longtable}{@{}p{0.07\textwidth}p{0.12\textwidth}p{0.75\textwidth}@{}}\toprule Line&Kind&Name\\\midrule\endhead

28&def&\nolinkurl{PaperXI.SpectralGeometricFrobenius.SpectralExtension}\\

30&instance&\nolinkurl{PaperXI.SpectralGeometricFrobenius.spectralField}\\

31&instance&\nolinkurl{PaperXI.SpectralGeometricFrobenius.spectralAlgebra}\\

32&instance&\nolinkurl{PaperXI.SpectralGeometricFrobenius.spectralFinite}\\

34&instance&\nolinkurl{PaperXI.SpectralGeometricFrobenius.spectralGalois}\\

38&def&\nolinkurl{PaperXI.SpectralGeometricFrobenius.originalFieldEquiv}\\

73&def&\nolinkurl{PaperXI.SpectralGeometricFrobenius.integers}\\

76&def&\nolinkurl{PaperXI.SpectralGeometricFrobenius.baseIntegers}\\

80&theorem&\nolinkurl{PaperXI.SpectralGeometricFrobenius.mem_integers}\\

92&theorem&\nolinkurl{PaperXI.SpectralGeometricFrobenius.integers_comap}\\

102&theorem&\nolinkurl{PaperXI.SpectralGeometricFrobenius.automorphism_stable}\\

116&instance&\nolinkurl{PaperXI.SpectralGeometricFrobenius.integersDVR}\\

120&instance&\nolinkurl{PaperXI.SpectralGeometricFrobenius.integersFiniteResidue}\\

124&instance&\nolinkurl{PaperXI.SpectralGeometricFrobenius.integersHenselian}\\

128&abbrev&\nolinkurl{PaperXI.SpectralGeometricFrobenius.BaseResidue}\\

129&abbrev&\nolinkurl{PaperXI.SpectralGeometricFrobenius.ExtensionResidue}\\

140&def&\nolinkurl{PaperXI.SpectralGeometricFrobenius.geometricFrobenius}\\

144&def&\nolinkurl{PaperXI.SpectralGeometricFrobenius.arithmeticFrobenius}\\

147&theorem&\nolinkurl{PaperXI.SpectralGeometricFrobenius.arithmetic_apply}\\

152&theorem&\nolinkurl{PaperXI.SpectralGeometricFrobenius.geometric_inverse_power}\\

158&def&\nolinkurl{PaperXI.SpectralGeometricFrobenius.residueAction}\\

164&theorem&\nolinkurl{PaperXI.SpectralGeometricFrobenius.exists_geometric_lift}\\

172&def&\nolinkurl{PaperXI.SpectralGeometricFrobenius.geometricLift}\\

174&theorem&\nolinkurl{PaperXI.SpectralGeometricFrobenius.residue_geometricLift}\\

179&theorem&\nolinkurl{PaperXI.SpectralGeometricFrobenius.geometricLift_normalized}\\

\bottomrule\end{longtable}

\section{\texorpdfstring{\nolinkurl{PaperXI.SpectralLocalField}}{PaperXI.SpectralLocalField}}\label{module:PaperXI.SpectralLocalField}

\textbf{Selected layer:} v0.44 parent.\quad\textbf{Source lines:} 175.\par\textbf{SHA-256:} {\scriptsize\nolinkurl{96cc4d62e4fc37f3550c6800e3715c60d191bf454547ea9a67d4aeae6376f9ab}}.

\VerbatimInput[fontsize=\scriptsize,breaklines=true,breakanywhere=true]{excerpts/D_SpectralLocalField.txt}

\textbf{Direct project imports:} \nolinkurl{PaperXI.CanonicalDVRValuation}.

\begin{longtable}{@{}p{0.07\textwidth}p{0.12\textwidth}p{0.75\textwidth}@{}}\toprule Line&Kind&Name\\\midrule\endhead

45&abbrev&\nolinkurl{PaperXI.SpectralLocalField.UpperField}\\

48&instance&\nolinkurl{PaperXI.SpectralLocalField.upperNormedField}\\

52&instance&\nolinkurl{PaperXI.SpectralLocalField.upperNontriviallyNormedField}\\

57&instance&\nolinkurl{PaperXI.SpectralLocalField.upperNormedSpace}\\

61&instance&\nolinkurl{PaperXI.SpectralLocalField.upperNormedAlgebra}\\

65&theorem&\nolinkurl{PaperXI.SpectralLocalField.norm_algebraMap}\\

70&theorem&\nolinkurl{PaperXI.SpectralLocalField.norm_eq_spectralNorm}\\

74&theorem&\nolinkurl{PaperXI.SpectralLocalField.algebraMap_isometry}\\

81&instance&\nolinkurl{PaperXI.SpectralLocalField.upperUltrametric}\\

93&instance&\nolinkurl{PaperXI.SpectralLocalField.upperComplete}\\

97&instance&\nolinkurl{PaperXI.SpectralLocalField.upperProper}\\

101&instance&\nolinkurl{PaperXI.SpectralLocalField.upperLocallyCompact}\\

105&instance&\nolinkurl{PaperXI.SpectralLocalField.upperNormValued}\\

109&instance&\nolinkurl{PaperXI.SpectralLocalField.upperValuativeRel}\\

112&instance&\nolinkurl{PaperXI.SpectralLocalField.upperNormValuationCompatible}\\

117&instance&\nolinkurl{PaperXI.SpectralLocalField.upperIsValuativeTopology}\\

124&instance&\nolinkurl{PaperXI.SpectralLocalField.upperValuativeNontrivial}\\

131&theorem&\nolinkurl{PaperXI.SpectralLocalField.relation_integer_eq_norm_integer}\\

143&theorem&\nolinkurl{PaperXI.SpectralLocalField.mem_integer_iff_norm_le_one}\\

150&instance&\nolinkurl{PaperXI.SpectralLocalField.upperIsNonarchimedeanLocalField}\\

157&instance&\nolinkurl{PaperXI.SpectralLocalField.upperFiniteResidue}\\

161&instance&\nolinkurl{PaperXI.SpectralLocalField.upperIntegerIsDVR}\\

166&def&\nolinkurl{PaperXI.SpectralLocalField.upperValueGroupIsoInt}\\

171&theorem&\nolinkurl{PaperXI.SpectralLocalField.finrank_eq_residue_finrank}\\

\bottomrule\end{longtable}

\section{\texorpdfstring{\nolinkurl{PaperXI.SplitRadical105}}{PaperXI.SplitRadical105}}\label{module:PaperXI.SplitRadical105}

\textbf{Selected layer:} base108 arithmetic branch.\quad\textbf{Source lines:} 152.\par\textbf{SHA-256:} {\scriptsize\nolinkurl{bbe4cd3ed2985df636cc326b7fceac83a5c48fec8b393ba3ae4a8da6e7fd5ca7}}.

\VerbatimInput[fontsize=\scriptsize,breaklines=true,breakanywhere=true]{excerpts/D_SplitRadical105.txt}

\textbf{Direct project imports:} none.

\begin{longtable}{@{}p{0.07\textwidth}p{0.12\textwidth}p{0.75\textwidth}@{}}\toprule Line&Kind&Name\\\midrule\endhead

17&def&\nolinkurl{PaperXI.SplitRadical105.rect}\\

20&theorem&\nolinkurl{PaperXI.SplitRadical105.rect_mul}\\

28&theorem&\nolinkurl{PaperXI.SplitRadical105.rect_add}\\

34&def&\nolinkurl{PaperXI.SplitRadical105.unip}\\

40&theorem&\nolinkurl{PaperXI.SplitRadical105.unip_add}\\

48&theorem&\nolinkurl{PaperXI.SplitRadical105.unip_zero}\\

53&def&\nolinkurl{PaperXI.SplitRadical105.radicalMap}\\

58&theorem&\nolinkurl{PaperXI.SplitRadical105.map_unip}\\

67&def&\nolinkurl{PaperXI.SplitRadical105.diagonalGL}\\

79&def&\nolinkurl{PaperXI.SplitRadical105.cartan}\\

82&def&\nolinkurl{PaperXI.SplitRadical105.zeroSet}\\

85&theorem&\nolinkurl{PaperXI.SplitRadical105.mem_zeroSet}\\

89&theorem&\nolinkurl{PaperXI.SplitRadical105.cartan_conjugation}\\

115&theorem&\nolinkurl{PaperXI.SplitRadical105.contracted_reduction}\\

129&theorem&\nolinkurl{PaperXI.SplitRadical105.exists_contracted_lift}\\

146&theorem&\nolinkurl{PaperXI.SplitRadical105.cartan_zero}\\

\bottomrule\end{longtable}

\section{\texorpdfstring{\nolinkurl{PaperXI.SymplecticAssembly}}{PaperXI.SymplecticAssembly}}\label{module:PaperXI.SymplecticAssembly}

\textbf{Selected layer:} v0.44 parent.\quad\textbf{Source lines:} 51.\par\textbf{SHA-256:} {\scriptsize\nolinkurl{67819fec0a0969ee2375d7ee93287a9f1e91d60e5da87e7a12aee6013e3e0d94}}.

\VerbatimInput[fontsize=\scriptsize,breaklines=true,breakanywhere=true]{excerpts/D_SymplecticAssembly.txt}

\textbf{Direct project imports:} \nolinkurl{PaperXI.DyadicCoordinates}.

\begin{longtable}{@{}p{0.07\textwidth}p{0.12\textwidth}p{0.75\textwidth}@{}}\toprule Line&Kind&Name\\\midrule\endhead

11&def&\nolinkurl{PaperXI.SymplecticAssembly.prepend}\\

26&theorem&\nolinkurl{PaperXI.SymplecticAssembly.prepend_form}\\

35&def&\nolinkurl{PaperXI.SymplecticAssembly.assemble}\\

40&theorem&\nolinkurl{PaperXI.SymplecticAssembly.assemble_form}\\

\bottomrule\end{longtable}

\section{\texorpdfstring{\nolinkurl{PaperXI.SymplecticCentralExtension}}{PaperXI.SymplecticCentralExtension}}\label{module:PaperXI.SymplecticCentralExtension}

\textbf{Selected layer:} v0.44 parent.\quad\textbf{Source lines:} 75.\par\textbf{SHA-256:} {\scriptsize\nolinkurl{b54d1f3d499676466c42f50e7b1fff895c0dd532d855652a51347f5bff12c61d}}.

\VerbatimInput[fontsize=\scriptsize,breaklines=true,breakanywhere=true]{excerpts/D_SymplecticCentralExtension.txt}

\textbf{Direct project imports:} \nolinkurl{PaperXI.CentralExtensionStructure}.

\begin{longtable}{@{}p{0.07\textwidth}p{0.12\textwidth}p{0.75\textwidth}@{}}\toprule Line&Kind&Name\\\midrule\endhead

17&def&\nolinkurl{PaperXI.SymplecticCentralExtension.scalarCharacter}\\

20&theorem&\nolinkurl{PaperXI.SymplecticCentralExtension.scalarCharacter_apply}\\

23&theorem&\nolinkurl{PaperXI.SymplecticCentralExtension.scalarCharacter_ofAdd}\\

28&theorem&\nolinkurl{PaperXI.SymplecticCentralExtension.character_separating_left}\\

45&theorem&\nolinkurl{PaperXI.SymplecticCentralExtension.radical_eq_bot}\\

66&theorem&\nolinkurl{PaperXI.SymplecticCentralExtension.radical_eq_bot_scalarCharacter}\\

\bottomrule\end{longtable}

\section{\texorpdfstring{\nolinkurl{PaperXI.SymplecticCoordinates}}{PaperXI.SymplecticCoordinates}}\label{module:PaperXI.SymplecticCoordinates}

\textbf{Selected layer:} v0.44 parent.\quad\textbf{Source lines:} 70.\par\textbf{SHA-256:} {\scriptsize\nolinkurl{b9b5ab88baf7b682acac5698d8cd558de3acf406eaec991f76a8b16e63500ecf}}.

\VerbatimInput[fontsize=\scriptsize,breaklines=true,breakanywhere=true]{excerpts/D_SymplecticCoordinates.txt}

\textbf{Direct project imports:} \nolinkurl{PaperXI.SymplecticSplit}, \nolinkurl{PaperXI.SymplecticAssembly}, \nolinkurl{PaperXI.DyadicDimension}.

\begin{longtable}{@{}p{0.07\textwidth}p{0.12\textwidth}p{0.75\textwidth}@{}}\toprule Line&Kind&Name\\\midrule\endhead

17&theorem&\nolinkurl{PaperXI.SymplecticCoordinates.exists_coordinates}\\

45&theorem&\nolinkurl{PaperXI.SymplecticCoordinates.exists_coordinates_dimension}\\

54&theorem&\nolinkurl{PaperXI.SymplecticCoordinates.even_finrank}\\

61&theorem&\nolinkurl{PaperXI.SymplecticCoordinates.exists_half_coordinates}\\

\bottomrule\end{longtable}

\section{\texorpdfstring{\nolinkurl{PaperXI.SymplecticHeisenbergUniqueness}}{PaperXI.SymplecticHeisenbergUniqueness}}\label{module:PaperXI.SymplecticHeisenbergUniqueness}

\textbf{Selected layer:} v0.44 parent.\quad\textbf{Source lines:} 85.\par\textbf{SHA-256:} {\scriptsize\nolinkurl{e54d8758376e52f844dd630abfdd19a6f940709729fabf3449bfa5f755195104}}.

\VerbatimInput[fontsize=\scriptsize,breaklines=true,breakanywhere=true]{excerpts/D_SymplecticHeisenbergUniqueness.txt}

\textbf{Direct project imports:} \nolinkurl{PaperXI.SymplecticCentralExtension}, \nolinkurl{PaperXI.CentralExtensionExistence}.

\begin{longtable}{@{}p{0.07\textwidth}p{0.12\textwidth}p{0.75\textwidth}@{}}\toprule Line&Kind&Name\\\midrule\endhead

19&theorem&\nolinkurl{PaperXI.SymplecticHeisenbergUniqueness.central_action}\\

35&theorem&\nolinkurl{PaperXI.SymplecticHeisenbergUniqueness.finiteDimensional}\\

42&theorem&\nolinkurl{PaperXI.SymplecticHeisenbergUniqueness.finrank_sqrt}\\

51&theorem&\nolinkurl{PaperXI.SymplecticHeisenbergUniqueness.exists_equiv}\\

61&theorem&\nolinkurl{PaperXI.SymplecticHeisenbergUniqueness.card_eq_square}\\

69&theorem&\nolinkurl{PaperXI.SymplecticHeisenbergUniqueness.exists_irreducible}\\

\bottomrule\end{longtable}

\section{\texorpdfstring{\nolinkurl{PaperXI.SymplecticSplit}}{PaperXI.SymplecticSplit}}\label{module:PaperXI.SymplecticSplit}

\textbf{Selected layer:} v0.44 parent.\quad\textbf{Source lines:} 126.\par\textbf{SHA-256:} {\scriptsize\nolinkurl{66521ce6aa9e9ca4cbf9bc05aad9112ad95e36069f2d3d718f214beb6515bc05}}.

\VerbatimInput[fontsize=\scriptsize,breaklines=true,breakanywhere=true]{excerpts/D_SymplecticSplit.txt}

\textbf{Direct project imports:} \nolinkurl{PaperXI.DyadicGauss}.

\begin{longtable}{@{}p{0.07\textwidth}p{0.12\textwidth}p{0.75\textwidth}@{}}\toprule Line&Kind&Name\\\midrule\endhead

17&theorem&\nolinkurl{PaperXI.SymplecticSplit.alternating_symmetric}\\

24&theorem&\nolinkurl{PaperXI.SymplecticSplit.exists_pair}\\

35&theorem&\nolinkurl{PaperXI.SymplecticSplit.split_of_pair}\\

113&theorem&\nolinkurl{PaperXI.SymplecticSplit.exists_split}\\

\bottomrule\end{longtable}

\section{\texorpdfstring{\nolinkurl{PaperXI.TameConductor}}{PaperXI.TameConductor}}\label{module:PaperXI.TameConductor}

\textbf{Selected layer:} v0.44 parent.\quad\textbf{Source lines:} 113.\par\textbf{SHA-256:} {\scriptsize\nolinkurl{808eb443d2a292fc79cd70ce323c369acd62c38e9393ec4c19269dc06d575dcf}}.

\VerbatimInput[fontsize=\scriptsize,breaklines=true,breakanywhere=true]{excerpts/D_TameConductor.txt}

\textbf{Direct project imports:} \nolinkurl{PaperXI.PacketBasis}, \nolinkurl{PaperXI.InertiaFixed}, \nolinkurl{PaperXI.InertiaQuotient}.

\begin{longtable}{@{}p{0.07\textwidth}p{0.12\textwidth}p{0.75\textwidth}@{}}\toprule Line&Kind&Name\\\midrule\endhead

21&theorem&\nolinkurl{PaperXI.TameConductor.zero_label_period}\\

26&theorem&\nolinkurl{PaperXI.TameConductor.inertia_invariants_eq_bot}\\

40&theorem&\nolinkurl{PaperXI.TameConductor.inertia_finrank_zero}\\

51&def&\nolinkurl{PaperXI.TameConductor.fixedCodimension}\\

56&def&\nolinkurl{PaperXI.TameConductor.finiteSwan}\\

61&theorem&\nolinkurl{PaperXI.TameConductor.invariants_eq_top_of_trivial_on}\\

68&theorem&\nolinkurl{PaperXI.TameConductor.fixedCodimension_zero_of_trivial_on}\\

76&theorem&\nolinkurl{PaperXI.TameConductor.finiteSwan_zero}\\

88&def&\nolinkurl{PaperXI.TameConductor.finiteArtin}\\

95&theorem&\nolinkurl{PaperXI.TameConductor.conductor_of_nontrivial_tame_atom}\\

\bottomrule\end{longtable}

\section{\texorpdfstring{\nolinkurl{PaperXI.TameDegree}}{PaperXI.TameDegree}}\label{module:PaperXI.TameDegree}

\textbf{Selected layer:} v0.44 parent.\quad\textbf{Source lines:} 170.\par\textbf{SHA-256:} {\scriptsize\nolinkurl{868e2f79b22a754bc070c23805f852b2c90a71cb802703bea2eaad248f39ad92}}.

\VerbatimInput[fontsize=\scriptsize,breaklines=true,breakanywhere=true]{excerpts/D_TameDegree.txt}

\textbf{Direct project imports:} \nolinkurl{PaperXI.TameTensorOrbits}.

\begin{longtable}{@{}p{0.07\textwidth}p{0.12\textwidth}p{0.75\textwidth}@{}}\toprule Line&Kind&Name\\\midrule\endhead

13&def&\nolinkurl{PaperXI.TameDegree.degree}\\

15&theorem&\nolinkurl{PaperXI.TameDegree.degree_one}\\

17&theorem&\nolinkurl{PaperXI.TameDegree.degree_mul}\\

20&theorem&\nolinkurl{PaperXI.TameDegree.degree_zpow}\\

23&theorem&\nolinkurl{PaperXI.TameDegree.degree_inv}\\

26&theorem&\nolinkurl{PaperXI.TameDegree.degree_conjugate}\\

30&def&\nolinkurl{PaperXI.TameDegree.degreeMod}\\

35&theorem&\nolinkurl{PaperXI.TameDegree.degreeMod_apply}\\

39&def&\nolinkurl{PaperXI.TameDegree.degreeSubgroup}\\

41&theorem&\nolinkurl{PaperXI.TameDegree.mem_degreeSubgroup}\\

46&theorem&\nolinkurl{PaperXI.TameDegree.degreeSubgroup_inter}\\

54&theorem&\nolinkurl{PaperXI.TameDegree.degreeSubgroup_le}\\

61&def&\nolinkurl{PaperXI.TameDegree.conjugationInto}\\

69&theorem&\nolinkurl{PaperXI.TameDegree.conjugationInto_apply}\\

73&def&\nolinkurl{PaperXI.TameDegree.translationAction}\\

83&theorem&\nolinkurl{PaperXI.TameDegree.translationAction_apply}\\

86&theorem&\nolinkurl{PaperXI.TameDegree.translationAction_fixed_iff}\\

92&def&\nolinkurl{PaperXI.TameDegree.pairAction}\\

98&theorem&\nolinkurl{PaperXI.TameDegree.pairAction_apply}\\

102&theorem&\nolinkurl{PaperXI.TameDegree.pairAction_fixed_iff}\\

107&def&\nolinkurl{PaperXI.TameDegree.blockAction}\\

113&theorem&\nolinkurl{PaperXI.TameDegree.blockAction_apply}\\

117&theorem&\nolinkurl{PaperXI.TameDegree.orbitMap_intertwines}\\

123&theorem&\nolinkurl{PaperXI.TameDegree.orbitEquiv_intertwines}\\

132&theorem&\nolinkurl{PaperXI.TameDegree.degree_phi}\\

135&theorem&\nolinkurl{PaperXI.TameDegree.degree_phi_zpow}\\

139&theorem&\nolinkurl{PaperXI.TameDegree.degreeMod_phi_zpow}\\

144&theorem&\nolinkurl{PaperXI.TameDegree.degreeMod_surjective}\\

152&theorem&\nolinkurl{PaperXI.TameDegree.translationAction_phi_zpow}\\

157&theorem&\nolinkurl{PaperXI.TameDegree.pairAction_phi_zpow}\\

163&theorem&\nolinkurl{PaperXI.TameDegree.fiber_transitive}\\

\bottomrule\end{longtable}

\section{\texorpdfstring{\nolinkurl{PaperXI.TameExponent}}{PaperXI.TameExponent}}\label{module:PaperXI.TameExponent}

\textbf{Selected layer:} v0.44 parent.\quad\textbf{Source lines:} 117.\par\textbf{SHA-256:} {\scriptsize\nolinkurl{738e3f3fe21ccb970422c7dc2b0de7f7780e588ba208153af7e04e70415b767e}}.

\VerbatimInput[fontsize=\scriptsize,breaklines=true,breakanywhere=true]{excerpts/D_TameExponent.txt}

\textbf{Direct project imports:} \nolinkurl{PaperXI.TameNorm}, \nolinkurl{PaperXI.ElementaryEnvelope}.

\begin{longtable}{@{}p{0.07\textwidth}p{0.12\textwidth}p{0.75\textwidth}@{}}\toprule Line&Kind&Name\\\midrule\endhead

11&def&\nolinkurl{PaperXI.TameExponent.rootEmbedding}\\

22&theorem&\nolinkurl{PaperXI.TameExponent.rootEmbedding_injective}\\

32&def&\nolinkurl{PaperXI.TameExponent.faithfulCharacter}\\

36&theorem&\nolinkurl{PaperXI.TameExponent.faithfulCharacter_injective}\\

44&theorem&\nolinkurl{PaperXI.TameExponent.pow_eq_one_iff}\\

63&theorem&\nolinkurl{PaperXI.TameExponent.pow_add_card}\\

73&theorem&\nolinkurl{PaperXI.TameExponent.product}\\

77&theorem&\nolinkurl{PaperXI.TameExponent.exists_exponent}\\

92&theorem&\nolinkurl{PaperXI.TameExponent.pow_eq_pow_iff}\\

102&def&\nolinkurl{PaperXI.TameExponent.coordinate}\\

105&theorem&\nolinkurl{PaperXI.TameExponent.coordinate_bijective}\\

\bottomrule\end{longtable}

\section{\texorpdfstring{\nolinkurl{PaperXI.TameGaussNaturality}}{PaperXI.TameGaussNaturality}}\label{module:PaperXI.TameGaussNaturality}

\textbf{Selected layer:} v0.44 parent.\quad\textbf{Source lines:} 79.\par\textbf{SHA-256:} {\scriptsize\nolinkurl{85bd64090e77888f06f233d6ea04a3cee2a539e3341988dd95273b8cdf5c1f7b}}.

\VerbatimInput[fontsize=\scriptsize,breaklines=true,breakanywhere=true]{excerpts/D_TameGaussNaturality.txt}

\textbf{Direct project imports:} \nolinkurl{PaperXI.GaussCharacters}.

\begin{longtable}{@{}p{0.07\textwidth}p{0.12\textwidth}p{0.75\textwidth}@{}}\toprule Line&Kind&Name\\\midrule\endhead

14&def&\nolinkurl{PaperXI.TameGaussNaturality.pullMul}\\

18&def&\nolinkurl{PaperXI.TameGaussNaturality.pullAdd}\\

22&theorem&\nolinkurl{PaperXI.TameGaussNaturality.sum_transport}\\

27&theorem&\nolinkurl{PaperXI.TameGaussNaturality.phase_transport}\\

34&theorem&\nolinkurl{PaperXI.TameGaussNaturality.phase_of_stable_additive}\\

40&theorem&\nolinkurl{PaperXI.TameGaussNaturality.phase_mulShift}\\

49&theorem&\nolinkurl{PaperXI.TameGaussNaturality.phase_trivial}\\

59&theorem&\nolinkurl{PaperXI.TameGaussNaturality.phase_norm_cases}\\

71&theorem&\nolinkurl{PaperXI.TameGaussNaturality.phase_ne_zero}\\

\bottomrule\end{longtable}

\section{\texorpdfstring{\nolinkurl{PaperXI.TameHomCriterion}}{PaperXI.TameHomCriterion}}\label{module:PaperXI.TameHomCriterion}

\textbf{Selected layer:} v0.44 parent.\quad\textbf{Source lines:} 66.\par\textbf{SHA-256:} {\scriptsize\nolinkurl{d125abe9165991cb809454306c36d0fa17a486c198612b7ddb05b4dd3cb35d35}}.

\VerbatimInput[fontsize=\scriptsize,breaklines=true,breakanywhere=true]{excerpts/D_TameHomCriterion.txt}

\textbf{Direct project imports:} \nolinkurl{PaperXI.SchurHom}, \nolinkurl{PaperXI.TensorFixed}.

\begin{longtable}{@{}p{0.07\textwidth}p{0.12\textwidth}p{0.75\textwidth}@{}}\toprule Line&Kind&Name\\\midrule\endhead

20&theorem&\nolinkurl{PaperXI.TameHomCriterion.hom_finrank_eq_trivial_count}\\

40&theorem&\nolinkurl{PaperXI.TameHomCriterion.tame_hom_test}\\

\bottomrule\end{longtable}

\section{\texorpdfstring{\nolinkurl{PaperXI.TameNorm}}{PaperXI.TameNorm}}\label{module:PaperXI.TameNorm}

\textbf{Selected layer:} v0.44 parent.\quad\textbf{Source lines:} 102.\par\textbf{SHA-256:} {\scriptsize\nolinkurl{249939314ef4f5eff217c783fa013cb8fe51f659337e9f7db4ec9a22d2a11c05}}.

\VerbatimInput[fontsize=\scriptsize,breaklines=true,breakanywhere=true]{excerpts/D_TameNorm.txt}

\textbf{Direct project imports:} none.

\begin{longtable}{@{}p{0.07\textwidth}p{0.12\textwidth}p{0.75\textwidth}@{}}\toprule Line&Kind&Name\\\midrule\endhead

15&def&\nolinkurl{PaperXI.TameNorm.embedding}\\

18&theorem&\nolinkurl{PaperXI.TameNorm.embedding_injective}\\

24&def&\nolinkurl{PaperXI.TameNorm.norm}\\

27&def&\nolinkurl{PaperXI.TameNorm.normExponent}\\

29&theorem&\nolinkurl{PaperXI.TameNorm.embedding_norm}\\

35&theorem&\nolinkurl{PaperXI.TameNorm.norm_lift}\\

44&theorem&\nolinkurl{PaperXI.TameNorm.norm_lift_eq_one_iff}\\

58&def&\nolinkurl{PaperXI.TameNorm.frobenius}\\

61&theorem&\nolinkurl{PaperXI.TameNorm.frobenius_apply}\\

68&theorem&\nolinkurl{PaperXI.TameNorm.frobenius_character}\\

76&theorem&\nolinkurl{PaperXI.TameNorm.shifted_norm_lift}\\

88&theorem&\nolinkurl{PaperXI.TameNorm.norm_frobenius}\\

93&theorem&\nolinkurl{PaperXI.TameNorm.shifted_lower_norm_lift}\\

\bottomrule\end{longtable}

\section{\texorpdfstring{\nolinkurl{PaperXI.TameNormPacket}}{PaperXI.TameNormPacket}}\label{module:PaperXI.TameNormPacket}

\textbf{Selected layer:} v0.44 parent.\quad\textbf{Source lines:} 118.\par\textbf{SHA-256:} {\scriptsize\nolinkurl{d19c32f222a6fd3f565bce086da6baf95a562f7266ec7a0105afefeb54a3a071}}.

\VerbatimInput[fontsize=\scriptsize,breaklines=true,breakanywhere=true]{excerpts/D_TameNormPacket.txt}

\textbf{Direct project imports:} \nolinkurl{PaperXI.TameExponent}.

\begin{longtable}{@{}p{0.07\textwidth}p{0.12\textwidth}p{0.75\textwidth}@{}}\toprule Line&Kind&Name\\\midrule\endhead

17&def&\nolinkurl{PaperXI.TameNormPacket.interaction}\\

23&theorem&\nolinkurl{PaperXI.TameNormPacket.refined_exponent}\\

35&theorem&\nolinkurl{PaperXI.TameNormPacket.refined_trivial_iff}\\

46&theorem&\nolinkurl{PaperXI.TameNormPacket.source_formula}\\

58&theorem&\nolinkurl{PaperXI.TameNormPacket.source_trivial_iff}\\

71&theorem&\nolinkurl{PaperXI.TameNormPacket.restricted_faithful}\\

80&theorem&\nolinkurl{PaperXI.TameNormPacket.exists_compatible_exponents}\\

93&theorem&\nolinkurl{PaperXI.TameNormPacket.source}\\

\bottomrule\end{longtable}

\section{\texorpdfstring{\nolinkurl{PaperXI.TameOrbitTransitivity}}{PaperXI.TameOrbitTransitivity}}\label{module:PaperXI.TameOrbitTransitivity}

\textbf{Selected layer:} v0.44 parent.\quad\textbf{Source lines:} 110.\par\textbf{SHA-256:} {\scriptsize\nolinkurl{e2f06470f27d1f2d54ee8573c71fdb9e4dd3a4f5526cf551c225e087f94f7971}}.

\VerbatimInput[fontsize=\scriptsize,breaklines=true,breakanywhere=true]{excerpts/D_TameOrbitTransitivity.txt}

\textbf{Direct project imports:} \nolinkurl{PaperXI.OrbitSpan}, \nolinkurl{PaperXI.GeneratedIrreducibility}.

\begin{longtable}{@{}p{0.07\textwidth}p{0.12\textwidth}p{0.75\textwidth}@{}}\toprule Line&Kind&Name\\\midrule\endhead

15&abbrev&\nolinkurl{PaperXI.TameOrbitTransitivity.tameRestriction}\\

18&def&\nolinkurl{PaperXI.TameOrbitTransitivity.generators}\\

23&theorem&\nolinkurl{PaperXI.TameOrbitTransitivity.geometric_relation}\\

38&theorem&\nolinkurl{PaperXI.TameOrbitTransitivity.generators_preserve_orbitSpan}\\

58&def&\nolinkurl{PaperXI.TameOrbitTransitivity.orbitSubrepresentation}\\

68&theorem&\nolinkurl{PaperXI.TameOrbitTransitivity.orbitSpan_eq_top}\\

80&theorem&\nolinkurl{PaperXI.TameOrbitTransitivity.nonzero_packets_one_orbit}\\

92&theorem&\nolinkurl{PaperXI.TameOrbitTransitivity.exists_nonzero_packet}\\

\bottomrule\end{longtable}

\section{\texorpdfstring{\nolinkurl{PaperXI.TameProjectors}}{PaperXI.TameProjectors}}\label{module:PaperXI.TameProjectors}

\textbf{Selected layer:} v0.44 parent.\quad\textbf{Source lines:} 86.\par\textbf{SHA-256:} {\scriptsize\nolinkurl{d5881299706cbe60a5cbefc9dc60bb9bdffbc1365228529c181b9a5120b71d8e}}.

\VerbatimInput[fontsize=\scriptsize,breaklines=true,breakanywhere=true]{excerpts/D_TameProjectors.txt}

\textbf{Direct project imports:} \nolinkurl{PaperXI.FourierPackets}.

\begin{longtable}{@{}p{0.07\textwidth}p{0.12\textwidth}p{0.75\textwidth}@{}}\toprule Line&Kind&Name\\\midrule\endhead

14&def&\nolinkurl{PaperXI.TameProjectors.projector}\\

17&theorem&\nolinkurl{PaperXI.TameProjectors.idempotent}\\

22&theorem&\nolinkurl{PaperXI.TameProjectors.orthogonal}\\

27&theorem&\nolinkurl{PaperXI.TameProjectors.complete}\\

32&theorem&\nolinkurl{PaperXI.TameProjectors.range_eq_character_space}\\

37&theorem&\nolinkurl{PaperXI.TameProjectors.forward_generator_covariance}\\

45&theorem&\nolinkurl{PaperXI.TameProjectors.geometric_covariance}\\

54&theorem&\nolinkurl{PaperXI.TameProjectors.character_weight}\\

61&theorem&\nolinkurl{PaperXI.TameProjectors.standard_root}\\

66&theorem&\nolinkurl{PaperXI.TameProjectors.manuscript_formula}\\

77&theorem&\nolinkurl{PaperXI.TameProjectors.tame_projector_theorem}\\

\bottomrule\end{longtable}

\section{\texorpdfstring{\nolinkurl{PaperXI.TameTensorInduction}}{PaperXI.TameTensorInduction}}\label{module:PaperXI.TameTensorInduction}

\textbf{Selected layer:} v0.44 parent.\quad\textbf{Source lines:} 111.\par\textbf{SHA-256:} {\scriptsize\nolinkurl{148b5e2d70afcc38a07c5b9fe8167fd690d2f687def771b465c1a8b81fbb799e}}.

\VerbatimInput[fontsize=\scriptsize,breaklines=true,breakanywhere=true]{excerpts/D_TameTensorInduction.txt}

\textbf{Direct project imports:} \nolinkurl{PaperXI.TameTensorPackets}, \nolinkurl{PaperXI.MonomialInductionSubgroup}.

\begin{longtable}{@{}p{0.07\textwidth}p{0.12\textwidth}p{0.75\textwidth}@{}}\toprule Line&Kind&Name\\\midrule\endhead

20&def&\nolinkurl{PaperXI.TameTensorInduction.character}\\

39&theorem&\nolinkurl{PaperXI.TameTensorInduction.character_apply}\\

48&theorem&\nolinkurl{PaperXI.TameTensorInduction.cycle_transitive}\\

55&def&\nolinkurl{PaperXI.TameTensorInduction.packetEquivalence}\\

70&def&\nolinkurl{PaperXI.TameTensorInduction.directSumEquiv}\\

86&def&\nolinkurl{PaperXI.TameTensorInduction.decomposition}\\

95&def&\nolinkurl{PaperXI.TameTensorInduction.decompositionWithCharacters}\\

\bottomrule\end{longtable}

\section{\texorpdfstring{\nolinkurl{PaperXI.TameTensorOrbits}}{PaperXI.TameTensorOrbits}}\label{module:PaperXI.TameTensorOrbits}

\textbf{Selected layer:} v0.44 parent.\quad\textbf{Source lines:} 210.\par\textbf{SHA-256:} {\scriptsize\nolinkurl{9485733d69ff7a1333197a690fb0ae38421460f342371510df530db8047753ec}}.

\VerbatimInput[fontsize=\scriptsize,breaklines=true,breakanywhere=true]{excerpts/D_TameTensorOrbits.txt}

\textbf{Direct project imports:} none.

\begin{longtable}{@{}p{0.07\textwidth}p{0.12\textwidth}p{0.75\textwidth}@{}}\toprule Line&Kind&Name\\\midrule\endhead

13&def&\nolinkurl{PaperXI.TameTensorOrbits.diagonal}\\

16&theorem&\nolinkurl{PaperXI.TameTensorOrbits.diagonal_injective}\\

19&theorem&\nolinkurl{PaperXI.TameTensorOrbits.diagonal_intCast}\\

24&theorem&\nolinkurl{PaperXI.TameTensorOrbits.diagonal_natCast}\\

30&def&\nolinkurl{PaperXI.TameTensorOrbits.difference}\\

34&theorem&\nolinkurl{PaperXI.TameTensorOrbits.difference_diagonal}\\

39&theorem&\nolinkurl{PaperXI.TameTensorOrbits.difference_add}\\

45&def&\nolinkurl{PaperXI.TameTensorOrbits.orbitMap}\\

48&theorem&\nolinkurl{PaperXI.TameTensorOrbits.orbitMap_apply}\\

53&theorem&\nolinkurl{PaperXI.TameTensorOrbits.orbitMap_zero}\\

57&theorem&\nolinkurl{PaperXI.TameTensorOrbits.difference_orbitMap}\\

62&theorem&\nolinkurl{PaperXI.TameTensorOrbits.orbitMap_injective}\\

77&instance&\nolinkurl{PaperXI.TameTensorOrbits.lcm_neZero}\\

80&instance&\nolinkurl{PaperXI.TameTensorOrbits.gcd_neZero}\\

83&theorem&\nolinkurl{PaperXI.TameTensorOrbits.orbitMap_bijective}\\

90&def&\nolinkurl{PaperXI.TameTensorOrbits.orbitEquiv}\\

93&theorem&\nolinkurl{PaperXI.TameTensorOrbits.orbitEquiv_apply}\\

97&theorem&\nolinkurl{PaperXI.TameTensorOrbits.orbitMap_translate}\\

103&def&\nolinkurl{PaperXI.TameTensorOrbits.translate}\\

107&theorem&\nolinkurl{PaperXI.TameTensorOrbits.orbitMap_translate_int}\\

116&theorem&\nolinkurl{PaperXI.TameTensorOrbits.translate_eq_iff}\\

124&theorem&\nolinkurl{PaperXI.TameTensorOrbits.interaction_count}\\

127&theorem&\nolinkurl{PaperXI.TameTensorOrbits.interaction_length}\\

131&theorem&\nolinkurl{PaperXI.TameTensorOrbits.total_dimension}\\

135&theorem&\nolinkurl{PaperXI.TameTensorOrbits.difference_translate}\\

140&theorem&\nolinkurl{PaperXI.TameTensorOrbits.sameOrbit_iff}\\

159&theorem&\nolinkurl{PaperXI.TameTensorOrbits.representatives_sameOrbit_iff}\\

173&def&\nolinkurl{PaperXI.TameTensorOrbits.fiberEquiv}\\

191&theorem&\nolinkurl{PaperXI.TameTensorOrbits.fiber_card}\\

196&theorem&\nolinkurl{PaperXI.TameTensorOrbits.fiberEquiv_apply}\\

200&theorem&\nolinkurl{PaperXI.TameTensorOrbits.orbitEquiv_symm_translate}\\

\bottomrule\end{longtable}

\section{\texorpdfstring{\nolinkurl{PaperXI.TameTensorPackets}}{PaperXI.TameTensorPackets}}\label{module:PaperXI.TameTensorPackets}

\textbf{Selected layer:} v0.44 parent.\quad\textbf{Source lines:} 81.\par\textbf{SHA-256:} {\scriptsize\nolinkurl{507e433f56fdae07efb7824aa8823c182dd07364260ea7b9d594f003e2514334}}.

\VerbatimInput[fontsize=\scriptsize,breaklines=true,breakanywhere=true]{excerpts/D_TameTensorPackets.txt}

\textbf{Direct project imports:} \nolinkurl{PaperXI.MonomialOrbit}, \nolinkurl{PaperXI.TameDegree}.

\begin{longtable}{@{}p{0.07\textwidth}p{0.12\textwidth}p{0.75\textwidth}@{}}\toprule Line&Kind&Name\\\midrule\endhead

21&def&\nolinkurl{PaperXI.TameTensorPackets.coefficient}\\

26&theorem&\nolinkurl{PaperXI.TameTensorPackets.tensor_action}\\

33&theorem&\nolinkurl{PaperXI.TameTensorPackets.orbit_covariance}\\

39&def&\nolinkurl{PaperXI.TameTensorPackets.packet}\\

46&def&\nolinkurl{PaperXI.TameTensorPackets.decomposition}\\

54&theorem&\nolinkurl{PaperXI.TameTensorPackets.packet_basis}\\

61&theorem&\nolinkurl{PaperXI.TameTensorPackets.packet_zero}\\

70&theorem&\nolinkurl{PaperXI.TameTensorPackets.packet_count}\\

73&theorem&\nolinkurl{PaperXI.TameTensorPackets.packet_dimension}\\

76&theorem&\nolinkurl{PaperXI.TameTensorPackets.dimension_sum}\\

\bottomrule\end{longtable}

\section{\texorpdfstring{\nolinkurl{PaperXI.TameTensorSource}}{PaperXI.TameTensorSource}}\label{module:PaperXI.TameTensorSource}

\textbf{Selected layer:} v0.44 parent.\quad\textbf{Source lines:} 122.\par\textbf{SHA-256:} {\scriptsize\nolinkurl{b52fef0aac236edae7a22103dcf267e147fa13606783529fa64e3c4c4e954d3e}}.

\VerbatimInput[fontsize=\scriptsize,breaklines=true,breakanywhere=true]{excerpts/D_TameTensorSource.txt}

\textbf{Direct project imports:} \nolinkurl{PaperXI.CyclicInducedCharacter}, \nolinkurl{PaperXI.TameTensorInduction}.

\begin{longtable}{@{}p{0.07\textwidth}p{0.12\textwidth}p{0.75\textwidth}@{}}\toprule Line&Kind&Name\\\midrule\endhead

18&def&\nolinkurl{PaperXI.TameTensorSource.interactionCharacter}\\

25&theorem&\nolinkurl{PaperXI.TameTensorSource.interactionCharacter_apply}\\

33&theorem&\nolinkurl{PaperXI.TameTensorSource.interactionCharacter_literal}\\

51&theorem&\nolinkurl{PaperXI.TameTensorSource.representative_val}\\

56&theorem&\nolinkurl{PaperXI.TameTensorSource.interactionCharacter_weight}\\

66&def&\nolinkurl{PaperXI.TameTensorSource.decomposition}\\

81&theorem&\nolinkurl{PaperXI.TameTensorSource.decomposition_intertwines}\\

93&theorem&\nolinkurl{PaperXI.TameTensorSource.interaction_dimension}\\

101&theorem&\nolinkurl{PaperXI.TameTensorSource.interaction_count}\\

106&theorem&\nolinkurl{PaperXI.TameTensorSource.interaction_dimension_sum}\\

116&theorem&\nolinkurl{PaperXI.TameTensorSource.tensor_dimension}\\

\bottomrule\end{longtable}

\section{\texorpdfstring{\nolinkurl{PaperXI.TensorFixed}}{PaperXI.TensorFixed}}\label{module:PaperXI.TensorFixed}

\textbf{Selected layer:} v0.44 parent.\quad\textbf{Source lines:} 34.\par\textbf{SHA-256:} {\scriptsize\nolinkurl{26c3d415523844b3389e5d7054e76b719d6d67574ac4fda6f67934a3be0c1f41}}.

\VerbatimInput[fontsize=\scriptsize,breaklines=true,breakanywhere=true]{excerpts/D_TensorFixed.txt}

\textbf{Direct project imports:} \nolinkurl{PaperXI.InvariantAssembly}.

\begin{longtable}{@{}p{0.07\textwidth}p{0.12\textwidth}p{0.75\textwidth}@{}}\toprule Line&Kind&Name\\\midrule\endhead

15&theorem&\nolinkurl{PaperXI.TensorFixed.count_trivial_packets}\\

\bottomrule\end{longtable}

\section{\texorpdfstring{\nolinkurl{PaperXI.TensorIntertwining107}}{PaperXI.TensorIntertwining107}}\label{module:PaperXI.TensorIntertwining107}

\textbf{Selected layer:} base108 arithmetic branch.\quad\textbf{Source lines:} 106.\par\textbf{SHA-256:} {\scriptsize\nolinkurl{4bb9b42f6f58022b82ae78cc0fe2d83d0aa6a153c7fc03aed66dd2ba7ea6f095}}.

\VerbatimInput[fontsize=\scriptsize,breaklines=true,breakanywhere=true]{excerpts/D_TensorIntertwining107.txt}

\textbf{Direct project imports:} \nolinkurl{PaperXI.TensorSlices107}, \nolinkurl{PaperXI.IntertwiningAlong105}.

\begin{longtable}{@{}p{0.07\textwidth}p{0.12\textwidth}p{0.75\textwidth}@{}}\toprule Line&Kind&Name\\\midrule\endhead

20&theorem&\nolinkurl{PaperXI.TensorIntertwining107.slice_along}\\

36&theorem&\nolinkurl{PaperXI.TensorIntertwining107.restrict_support}\\

45&theorem&\nolinkurl{PaperXI.TensorIntertwining107.second_along}\\

62&theorem&\nolinkurl{PaperXI.TensorIntertwining107.map_along}\\

76&theorem&\nolinkurl{PaperXI.TensorIntertwining107.factor_along}\\

88&theorem&\nolinkurl{PaperXI.TensorIntertwining107.existsAt_iff}\\

\bottomrule\end{longtable}

\section{\texorpdfstring{\nolinkurl{PaperXI.TensorIrreducible}}{PaperXI.TensorIrreducible}}\label{module:PaperXI.TensorIrreducible}

\textbf{Selected layer:} v0.44 parent.\quad\textbf{Source lines:} 78.\par\textbf{SHA-256:} {\scriptsize\nolinkurl{f75b3acfdfeba7e36d31c8d0085caffec072b5c712c4b76f407d8a0b65d5f721}}.

\VerbatimInput[fontsize=\scriptsize,breaklines=true,breakanywhere=true]{excerpts/D_TensorIrreducible.txt}

\textbf{Direct project imports:} \nolinkurl{PaperXI.CentralExtensionUniqueness}.

\begin{longtable}{@{}p{0.07\textwidth}p{0.12\textwidth}p{0.75\textwidth}@{}}\toprule Line&Kind&Name\\\midrule\endhead

18&theorem&\nolinkurl{PaperXI.TensorIrreducible.finrank_sq_of_irreducible}\\

32&theorem&\nolinkurl{PaperXI.TensorIrreducible.action_injective}\\

42&theorem&\nolinkurl{PaperXI.TensorIrreducible.action_surjective}\\

58&theorem&\nolinkurl{PaperXI.TensorIrreducible.irreducible}\\

\bottomrule\end{longtable}

\section{\texorpdfstring{\nolinkurl{PaperXI.TensorPolynomial}}{PaperXI.TensorPolynomial}}\label{module:PaperXI.TensorPolynomial}

\textbf{Selected layer:} v0.44 parent.\quad\textbf{Source lines:} 89.\par\textbf{SHA-256:} {\scriptsize\nolinkurl{45345b4d384aa9e69de8e139f5f1cd75cc18a97728cf3ba6d5d6f53413956ca7}}.

\VerbatimInput[fontsize=\scriptsize,breaklines=true,breakanywhere=true]{excerpts/D_TensorPolynomial.txt}

\textbf{Direct project imports:} \nolinkurl{PaperXI.PolynomialAssembly}.

\begin{longtable}{@{}p{0.07\textwidth}p{0.12\textwidth}p{0.75\textwidth}@{}}\toprule Line&Kind&Name\\\midrule\endhead

13&theorem&\nolinkurl{PaperXI.TensorPolynomial.quotient_stable}\\

22&theorem&\nolinkurl{PaperXI.TensorPolynomial.restricted_eq_factored}\\

33&theorem&\nolinkurl{PaperXI.TensorPolynomial.cyclePolynomial_ne_zero}\\

56&theorem&\nolinkurl{PaperXI.TensorPolynomial.total_stable}\\

63&theorem&\nolinkurl{PaperXI.TensorPolynomial.polynomial_packets}\\

75&theorem&\nolinkurl{PaperXI.TensorPolynomial.determinant_packets}\\

82&theorem&\nolinkurl{PaperXI.TensorPolynomial.localFactor_packets}\\

\bottomrule\end{longtable}

\section{\texorpdfstring{\nolinkurl{PaperXI.TensorSlices107}}{PaperXI.TensorSlices107}}\label{module:PaperXI.TensorSlices107}

\textbf{Selected layer:} base108 arithmetic branch.\quad\textbf{Source lines:} 115.\par\textbf{SHA-256:} {\scriptsize\nolinkurl{0176d323e1f5d380a8aac25c784710e8b877b7d4aca3fa718eb098f4f08b954b}}.

\VerbatimInput[fontsize=\scriptsize,breaklines=true,breakanywhere=true]{excerpts/D_TensorSlices107.txt}

\textbf{Direct project imports:} none.

\begin{longtable}{@{}p{0.07\textwidth}p{0.12\textwidth}p{0.75\textwidth}@{}}\toprule Line&Kind&Name\\\midrule\endhead

14&def&\nolinkurl{PaperXI.TensorSlices107.right}\\

17&def&\nolinkurl{PaperXI.TensorSlices107.left}\\

20&theorem&\nolinkurl{PaperXI.TensorSlices107.right_tmul}\\

23&theorem&\nolinkurl{PaperXI.TensorSlices107.left_tmul}\\

26&theorem&\nolinkurl{PaperXI.TensorSlices107.right_ext}\\

35&theorem&\nolinkurl{PaperXI.TensorSlices107.contractions_commute}\\

42&theorem&\nolinkurl{PaperXI.TensorSlices107.right_map}\\

49&def&\nolinkurl{PaperXI.TensorSlices107.slice}\\

53&theorem&\nolinkurl{PaperXI.TensorSlices107.slice_apply}\\

57&theorem&\nolinkurl{PaperXI.TensorSlices107.exists_nonzero_slice}\\

71&theorem&\nolinkurl{PaperXI.TensorSlices107.factor_of_slice_line}\\

97&theorem&\nolinkurl{PaperXI.TensorSlices107.tmul_left_injective}\\

104&theorem&\nolinkurl{PaperXI.TensorSlices107.map_left_injective}\\

\bottomrule\end{longtable}

\section{\texorpdfstring{\nolinkurl{PaperXI.TensorSubspace}}{PaperXI.TensorSubspace}}\label{module:PaperXI.TensorSubspace}

\textbf{Selected layer:} v0.44 parent.\quad\textbf{Source lines:} 76.\par\textbf{SHA-256:} {\scriptsize\nolinkurl{e65d7a7aea5e82ccccba7c606a3a314c89834d52f16441aa2d64c184362a5038}}.

\VerbatimInput[fontsize=\scriptsize,breaklines=true,breakanywhere=true]{excerpts/D_TensorSubspace.txt}

\textbf{Direct project imports:} none.

\begin{longtable}{@{}p{0.07\textwidth}p{0.12\textwidth}p{0.75\textwidth}@{}}\toprule Line&Kind&Name\\\midrule\endhead

18&def&\nolinkurl{PaperXI.TensorSubspace.image}\\

21&theorem&\nolinkurl{PaperXI.TensorSubspace.inclusion_injective}\\

25&theorem&\nolinkurl{PaperXI.TensorSubspace.map_injective}\\

29&theorem&\nolinkurl{PaperXI.TensorSubspace.tmul_mem}\\

34&theorem&\nolinkurl{PaperXI.TensorSubspace.image_stable}\\

48&theorem&\nolinkurl{PaperXI.TensorSubspace.finrank_image}\\

55&theorem&\nolinkurl{PaperXI.TensorSubspace.image_eq_bot_iff}\\

61&theorem&\nolinkurl{PaperXI.TensorSubspace.image_eq_top_iff}\\

\bottomrule\end{longtable}

\section{\texorpdfstring{\nolinkurl{PaperXI.TensorTypeOperators70}}{PaperXI.TensorTypeOperators70}}\label{module:PaperXI.TensorTypeOperators70}

\textbf{Selected layer:} base108 arithmetic branch.\quad\textbf{Source lines:} 94.\par\textbf{SHA-256:} {\scriptsize\nolinkurl{e82f0e46bab6ff2d1096ab07053553ddb6218bf13ff00afc296a84e2af2b270c}}.

\VerbatimInput[fontsize=\scriptsize,breaklines=true,breakanywhere=true]{excerpts/D_TensorTypeOperators70.txt}

\textbf{Direct project imports:} \nolinkurl{PaperXI.InducedTypeSpace68}.

\begin{longtable}{@{}p{0.07\textwidth}p{0.12\textwidth}p{0.75\textwidth}@{}}\toprule Line&Kind&Name\\\midrule\endhead

17&def&\nolinkurl{PaperXI.TensorTypeOperators70.amplify}\\

20&theorem&\nolinkurl{PaperXI.TensorTypeOperators70.amplify_tmul}\\

24&theorem&\nolinkurl{PaperXI.TensorTypeOperators70.amplify_injective}\\

37&def&\nolinkurl{PaperXI.TensorTypeOperators70.realize}\\

41&theorem&\nolinkurl{PaperXI.TensorTypeOperators70.realize_tmul}\\

50&theorem&\nolinkurl{PaperXI.TensorTypeOperators70.realize_injective}\\

53&theorem&\nolinkurl{PaperXI.TensorTypeOperators70.realize_eq_iff}\\

56&theorem&\nolinkurl{PaperXI.TensorTypeOperators70.realize_zero_iff}\\

61&theorem&\nolinkurl{PaperXI.TensorTypeOperators70.residual_zero_iff}\\

66&theorem&\nolinkurl{PaperXI.TensorTypeOperators70.projector_iff}\\

71&theorem&\nolinkurl{PaperXI.TensorTypeOperators70.intertwining_iff}\\

77&theorem&\nolinkurl{PaperXI.TensorTypeOperators70.word_product}\\

84&theorem&\nolinkurl{PaperXI.TensorTypeOperators70.holonomy_iff}\\

89&theorem&\nolinkurl{PaperXI.TensorTypeOperators70.weighted_sum}\\

\bottomrule\end{longtable}

\section{\texorpdfstring{\nolinkurl{PaperXI.TraceDiscrepancy}}{PaperXI.TraceDiscrepancy}}\label{module:PaperXI.TraceDiscrepancy}

\textbf{Selected layer:} v0.44 parent.\quad\textbf{Source lines:} 84.\par\textbf{SHA-256:} {\scriptsize\nolinkurl{0bdba459861ec52105d2a28d2acdd6d4f0f5ac5d7a55cd93ddff6186ed822e15}}.

\VerbatimInput[fontsize=\scriptsize,breaklines=true,breakanywhere=true]{excerpts/D_TraceDiscrepancy.txt}

\textbf{Direct project imports:} \nolinkurl{PaperXI.CharacterGenerators}.

\begin{longtable}{@{}p{0.07\textwidth}p{0.12\textwidth}p{0.75\textwidth}@{}}\toprule Line&Kind&Name\\\midrule\endhead

14&theorem&\nolinkurl{PaperXI.TraceDiscrepancy.conjugate_eq}\\

20&theorem&\nolinkurl{PaperXI.TraceDiscrepancy.trace_eq}\\

26&theorem&\nolinkurl{PaperXI.TraceDiscrepancy.scalar_eq_one}\\

35&theorem&\nolinkurl{PaperXI.TraceDiscrepancy.normalized_trace_eq}\\

41&theorem&\nolinkurl{PaperXI.TraceDiscrepancy.scalar_eq_one_of_normalized}\\

51&theorem&\nolinkurl{PaperXI.TraceDiscrepancy.character_trivial}\\

63&theorem&\nolinkurl{PaperXI.TraceDiscrepancy.zero_trace_cannot_detect_scalar}\\

67&theorem&\nolinkurl{PaperXI.TraceDiscrepancy.bridge_from_traces}\\

\bottomrule\end{longtable}

\section{\texorpdfstring{\nolinkurl{PaperXI.TranslationAverage}}{PaperXI.TranslationAverage}}\label{module:PaperXI.TranslationAverage}

\textbf{Selected layer:} v0.44 parent.\quad\textbf{Source lines:} 102.\par\textbf{SHA-256:} {\scriptsize\nolinkurl{7b04c80800e59fc6ced892041eb4fe9306e01e08e18e61e79d27c07acefe9d1e}}.

\VerbatimInput[fontsize=\scriptsize,breaklines=true,breakanywhere=true]{excerpts/D_TranslationAverage.txt}

\textbf{Direct project imports:} \nolinkurl{PaperXI.FiniteAverage}.

\begin{longtable}{@{}p{0.07\textwidth}p{0.12\textwidth}p{0.75\textwidth}@{}}\toprule Line&Kind&Name\\\midrule\endhead

13&def&\nolinkurl{PaperXI.TranslationAverage.translation}\\

22&theorem&\nolinkurl{PaperXI.TranslationAverage.translation_apply}\\

25&def&\nolinkurl{PaperXI.TranslationAverage.mean}\\

27&def&\nolinkurl{PaperXI.TranslationAverage.constantSubmodule}\\

31&theorem&\nolinkurl{PaperXI.TranslationAverage.mem_constantSubmodule}\\

39&theorem&\nolinkurl{PaperXI.TranslationAverage.average_apply}\\

48&theorem&\nolinkurl{PaperXI.TranslationAverage.range_eq}\\

61&theorem&\nolinkurl{PaperXI.TranslationAverage.finrank_range}\\

71&def&\nolinkurl{PaperXI.TranslationAverage.pairing}\\

73&theorem&\nolinkurl{PaperXI.TranslationAverage.pairing_translation}\\

79&theorem&\nolinkurl{PaperXI.TranslationAverage.pairing_self_adjoint}\\

91&theorem&\nolinkurl{PaperXI.TranslationAverage.residual_orthogonal}\\

\bottomrule\end{longtable}

\section{\texorpdfstring{\nolinkurl{PaperXI.TwistedAction}}{PaperXI.TwistedAction}}\label{module:PaperXI.TwistedAction}

\textbf{Selected layer:} v0.44 parent.\quad\textbf{Source lines:} 79.\par\textbf{SHA-256:} {\scriptsize\nolinkurl{97ed5ee03bec3fc432018d895c823299a7cc73ddf65cb1b712417d9948a5eb71}}.

\VerbatimInput[fontsize=\scriptsize,breaklines=true,breakanywhere=true]{excerpts/D_TwistedAction.txt}

\textbf{Direct project imports:} \nolinkurl{PaperXI.TwistedAlgebra}.

\begin{longtable}{@{}p{0.07\textwidth}p{0.12\textwidth}p{0.75\textwidth}@{}}\toprule Line&Kind&Name\\\midrule\endhead

15&def&\nolinkurl{PaperXI.TwistedAction.linearExtension}\\

19&theorem&\nolinkurl{PaperXI.TwistedAction.linearExtension_basis}\\

24&theorem&\nolinkurl{PaperXI.TwistedAction.linearExtension_one}\\

29&theorem&\nolinkurl{PaperXI.TwistedAction.linearExtension_mul}\\

44&def&\nolinkurl{PaperXI.TwistedAction.action}\\

51&theorem&\nolinkurl{PaperXI.TwistedAction.action_basis}\\

58&theorem&\nolinkurl{PaperXI.TwistedAction.action_unique}\\

70&theorem&\nolinkurl{PaperXI.TwistedAction.existsUnique_action}\\

\bottomrule\end{longtable}

\section{\texorpdfstring{\nolinkurl{PaperXI.TwistedActionEquivalence}}{PaperXI.TwistedActionEquivalence}}\label{module:PaperXI.TwistedActionEquivalence}

\textbf{Selected layer:} v0.44 parent.\quad\textbf{Source lines:} 107.\par\textbf{SHA-256:} {\scriptsize\nolinkurl{6a14cac5426c7fb2c44d091685da2d07440c4b883f7409f7caff5c2c38919e74}}.

\VerbatimInput[fontsize=\scriptsize,breaklines=true,breakanywhere=true]{excerpts/D_TwistedActionEquivalence.txt}

\textbf{Direct project imports:} \nolinkurl{PaperXI.TwistedAction}, \nolinkurl{PaperXI.TwistedModules}, \nolinkurl{PaperXI.AlgebraActionModule}.

\begin{longtable}{@{}p{0.07\textwidth}p{0.12\textwidth}p{0.75\textwidth}@{}}\toprule Line&Kind&Name\\\midrule\endhead

21&theorem&\nolinkurl{PaperXI.TwistedActionEquivalence.stable_iff}\\

34&theorem&\nolinkurl{PaperXI.TwistedActionEquivalence.simpleModule}\\

47&theorem&\nolinkurl{PaperXI.TwistedActionEquivalence.finiteDimensional}\\

58&theorem&\nolinkurl{PaperXI.TwistedActionEquivalence.finrank_sq}\\

66&theorem&\nolinkurl{PaperXI.TwistedActionEquivalence.finrank_sqrt}\\

71&theorem&\nolinkurl{PaperXI.TwistedActionEquivalence.action_injective}\\

77&theorem&\nolinkurl{PaperXI.TwistedActionEquivalence.action_surjective}\\

92&theorem&\nolinkurl{PaperXI.TwistedActionEquivalence.action_bijective}\\

96&def&\nolinkurl{PaperXI.TwistedActionEquivalence.actionEquiv}\\

100&theorem&\nolinkurl{PaperXI.TwistedActionEquivalence.actionEquiv_apply}\\

103&theorem&\nolinkurl{PaperXI.TwistedActionEquivalence.actionEquiv_basis}\\

\bottomrule\end{longtable}

\section{\texorpdfstring{\nolinkurl{PaperXI.TwistedAlgebra}}{PaperXI.TwistedAlgebra}}\label{module:PaperXI.TwistedAlgebra}

\textbf{Selected layer:} v0.44 parent.\quad\textbf{Source lines:} 109.\par\textbf{SHA-256:} {\scriptsize\nolinkurl{578da8f15a3a094d50dbc2d501fc66ea1c8ad354c2a29fb8acb5fd3f4c7b2020}}.

\VerbatimInput[fontsize=\scriptsize,breaklines=true,breakanywhere=true]{excerpts/D_TwistedAlgebra.txt}

\textbf{Direct project imports:} \nolinkurl{PaperXI.TwistedRegular}.

\begin{longtable}{@{}p{0.07\textwidth}p{0.12\textwidth}p{0.75\textwidth}@{}}\toprule Line&Kind&Name\\\midrule\endhead

16&def&\nolinkurl{PaperXI.TwistedAlgebra.operatorSpan}\\

19&theorem&\nolinkurl{PaperXI.TwistedAlgebra.regular_mem}\\

22&theorem&\nolinkurl{PaperXI.TwistedAlgebra.one_mem}\\

27&theorem&\nolinkurl{PaperXI.TwistedAlgebra.mul_mem}\\

45&def&\nolinkurl{PaperXI.TwistedAlgebra.subalgebra}\\

57&abbrev&\nolinkurl{PaperXI.TwistedAlgebra.Algebra}\\

62&def&\nolinkurl{PaperXI.TwistedAlgebra.basis}\\

65&theorem&\nolinkurl{PaperXI.TwistedAlgebra.coe_basis}\\

71&theorem&\nolinkurl{PaperXI.TwistedAlgebra.basis_mul}\\

78&theorem&\nolinkurl{PaperXI.TwistedAlgebra.basis_one}\\

84&def&\nolinkurl{PaperXI.TwistedAlgebra.basisUnit}\\

96&theorem&\nolinkurl{PaperXI.TwistedAlgebra.coe_basisUnit}\\

99&theorem&\nolinkurl{PaperXI.TwistedAlgebra.basis_isUnit}\\

102&instance&\nolinkurl{PaperXI.TwistedAlgebra.finiteDimensional}\\

105&theorem&\nolinkurl{PaperXI.TwistedAlgebra.finrank_eq_card}\\

\bottomrule\end{longtable}

\section{\texorpdfstring{\nolinkurl{PaperXI.TwistedMatrix}}{PaperXI.TwistedMatrix}}\label{module:PaperXI.TwistedMatrix}

\textbf{Selected layer:} v0.44 parent.\quad\textbf{Source lines:} 54.\par\textbf{SHA-256:} {\scriptsize\nolinkurl{99153ad911c9dab5fd807658f17f2d4d1583b71edd10538aad5c08bfb601e959}}.

\VerbatimInput[fontsize=\scriptsize,breaklines=true,breakanywhere=true]{excerpts/D_TwistedMatrix.txt}

\textbf{Direct project imports:} \nolinkurl{PaperXI.TwistedSimple}.

\begin{longtable}{@{}p{0.07\textwidth}p{0.12\textwidth}p{0.75\textwidth}@{}}\toprule Line&Kind&Name\\\midrule\endhead

18&theorem&\nolinkurl{PaperXI.TwistedMatrix.exists_matrix}\\

27&theorem&\nolinkurl{PaperXI.TwistedMatrix.matrix_sqrt}\\

39&theorem&\nolinkurl{PaperXI.TwistedMatrix.order_square}\\

46&theorem&\nolinkurl{PaperXI.TwistedMatrix.classification}\\

\bottomrule\end{longtable}

\section{\texorpdfstring{\nolinkurl{PaperXI.TwistedModules}}{PaperXI.TwistedModules}}\label{module:PaperXI.TwistedModules}

\textbf{Selected layer:} v0.44 parent.\quad\textbf{Source lines:} 54.\par\textbf{SHA-256:} {\scriptsize\nolinkurl{75c7bf4253db1b19641a24e2fbeaa7317eac3362bc5c6427038f89481a072edb}}.

\VerbatimInput[fontsize=\scriptsize,breaklines=true,breakanywhere=true]{excerpts/D_TwistedModules.txt}

\textbf{Direct project imports:} \nolinkurl{PaperXI.TwistedMatrix}, \nolinkurl{PaperXI.SimpleModuleClassification}.

\begin{longtable}{@{}p{0.07\textwidth}p{0.12\textwidth}p{0.75\textwidth}@{}}\toprule Line&Kind&Name\\\midrule\endhead

18&theorem&\nolinkurl{PaperXI.TwistedModules.exists_simple_ideal}\\

26&theorem&\nolinkurl{PaperXI.TwistedModules.equivalent}\\

40&theorem&\nolinkurl{PaperXI.TwistedModules.finiteDimensional}\\

45&theorem&\nolinkurl{PaperXI.TwistedModules.finrank_sq}\\

51&theorem&\nolinkurl{PaperXI.TwistedModules.finrank_sqrt}\\

\bottomrule\end{longtable}

\section{\texorpdfstring{\nolinkurl{PaperXI.TwistedRegular}}{PaperXI.TwistedRegular}}\label{module:PaperXI.TwistedRegular}

\textbf{Selected layer:} v0.44 parent.\quad\textbf{Source lines:} 81.\par\textbf{SHA-256:} {\scriptsize\nolinkurl{542e57f3ada86dab99247efc614f2e360c4ed62aa6278ac5d318f7cad41242b0}}.

\VerbatimInput[fontsize=\scriptsize,breaklines=true,breakanywhere=true]{excerpts/D_TwistedRegular.txt}

\textbf{Direct project imports:} \nolinkurl{PaperXI.CocycleCommutator}.

\begin{longtable}{@{}p{0.07\textwidth}p{0.12\textwidth}p{0.75\textwidth}@{}}\toprule Line&Kind&Name\\\midrule\endhead

18&def&\nolinkurl{PaperXI.TwistedRegular.regular}\\

23&theorem&\nolinkurl{PaperXI.TwistedRegular.regular_apply}\\

26&theorem&\nolinkurl{PaperXI.TwistedRegular.regular_one}\\

31&theorem&\nolinkurl{PaperXI.TwistedRegular.regular_mul}\\

46&def&\nolinkurl{PaperXI.TwistedRegular.deltaOne}\\

50&theorem&\nolinkurl{PaperXI.TwistedRegular.regular_deltaOne}\\

66&def&\nolinkurl{PaperXI.TwistedRegular.evaluateOne}\\

71&theorem&\nolinkurl{PaperXI.TwistedRegular.regular_linearIndependent}\\

\bottomrule\end{longtable}

\section{\texorpdfstring{\nolinkurl{PaperXI.TwistedRescaling}}{PaperXI.TwistedRescaling}}\label{module:PaperXI.TwistedRescaling}

\textbf{Selected layer:} v0.44 parent.\quad\textbf{Source lines:} 158.\par\textbf{SHA-256:} {\scriptsize\nolinkurl{1bb12f819e8232e254d6147b5fac209667d032836c4aa02fe72c472015d40e0d}}.

\VerbatimInput[fontsize=\scriptsize,breaklines=true,breakanywhere=true]{excerpts/D_TwistedRescaling.txt}

\textbf{Direct project imports:} \nolinkurl{PaperXI.TwistedAction}, \nolinkurl{PaperXI.CocycleBasis}.

\begin{longtable}{@{}p{0.07\textwidth}p{0.12\textwidth}p{0.75\textwidth}@{}}\toprule Line&Kind&Name\\\midrule\endhead

15&def&\nolinkurl{PaperXI.TwistedRescaling.congrEquiv}\\

20&theorem&\nolinkurl{PaperXI.TwistedRescaling.congrEquiv_basis}\\

29&def&\nolinkurl{PaperXI.TwistedRescaling.linearEquiv}\\

36&theorem&\nolinkurl{PaperXI.TwistedRescaling.linearEquiv_basis}\\

46&theorem&\nolinkurl{PaperXI.TwistedRescaling.scaled_basis_mul}\\

56&theorem&\nolinkurl{PaperXI.TwistedRescaling.linearEquiv_one}\\

61&theorem&\nolinkurl{PaperXI.TwistedRescaling.linearEquiv_mul}\\

88&def&\nolinkurl{PaperXI.TwistedRescaling.rescaleEquiv}\\

95&theorem&\nolinkurl{PaperXI.TwistedRescaling.rescaleEquiv_basis}\\

103&theorem&\nolinkurl{PaperXI.TwistedRescaling.rescaleEquiv_symm_basis}\\

116&def&\nolinkurl{PaperXI.TwistedRescaling.transport}\\

122&theorem&\nolinkurl{PaperXI.TwistedRescaling.transport_basis}\\

131&theorem&\nolinkurl{PaperXI.TwistedRescaling.transport_unique}\\

145&def&\nolinkurl{PaperXI.TwistedRescaling.actionEquiv}\\

\bottomrule\end{longtable}

\section{\texorpdfstring{\nolinkurl{PaperXI.TwistedSimple}}{PaperXI.TwistedSimple}}\label{module:PaperXI.TwistedSimple}

\textbf{Selected layer:} v0.44 parent.\quad\textbf{Source lines:} 205.\par\textbf{SHA-256:} {\scriptsize\nolinkurl{f420d05bd41765009035ee02338ae4721a2914e0946fb52c572fcee8ca333c60}}.

\VerbatimInput[fontsize=\scriptsize,breaklines=true,breakanywhere=true]{excerpts/D_TwistedSimple.txt}

\textbf{Direct project imports:} \nolinkurl{PaperXI.CocycleBasis}, \nolinkurl{PaperXI.TwistedAlgebra}, \nolinkurl{PaperXI.SchrodingerFourier}.

\begin{longtable}{@{}p{0.07\textwidth}p{0.12\textwidth}p{0.75\textwidth}@{}}\toprule Line&Kind&Name\\\midrule\endhead

19&theorem&\nolinkurl{PaperXI.TwistedSimple.commutator_separates}\\

34&theorem&\nolinkurl{PaperXI.TwistedSimple.basis_commute}\\

44&theorem&\nolinkurl{PaperXI.TwistedSimple.ideal_smul_mem}\\

52&theorem&\nolinkurl{PaperXI.TwistedSimple.simple_of_basis}\\

117&theorem&\nolinkurl{PaperXI.TwistedSimple.isSimpleRing}\\

129&theorem&\nolinkurl{PaperXI.TwistedSimple.radical_eq_bot_iff}\\

135&theorem&\nolinkurl{PaperXI.TwistedSimple.isSimpleRing_of_nondegenerate}\\

144&theorem&\nolinkurl{PaperXI.TwistedSimple.diagonal_mul_basis}\\

157&theorem&\nolinkurl{PaperXI.TwistedSimple.center_scalar_of_basis}\\

189&theorem&\nolinkurl{PaperXI.TwistedSimple.center_eq_scalars}\\

\bottomrule\end{longtable}

\section{\texorpdfstring{\nolinkurl{PaperXI.TypeNormalizerAction}}{PaperXI.TypeNormalizerAction}}\label{module:PaperXI.TypeNormalizerAction}

\textbf{Selected layer:} v0.44 parent.\quad\textbf{Source lines:} 108.\par\textbf{SHA-256:} {\scriptsize\nolinkurl{623d5fefcb3b3c856476bbd194fd5a1101347efa5511dbdd0c6a68ffc6e7c2f2}}.

\VerbatimInput[fontsize=\scriptsize,breaklines=true,breakanywhere=true]{excerpts/D_TypeNormalizerAction.txt}

\textbf{Direct project imports:} none.

\begin{longtable}{@{}p{0.07\textwidth}p{0.12\textwidth}p{0.75\textwidth}@{}}\toprule Line&Kind&Name\\\midrule\endhead

15&def&\nolinkurl{PaperXI.TypeNormalizerAction.conjugation}\\

25&theorem&\nolinkurl{PaperXI.TypeNormalizerAction.conjugation_coe}\\

28&theorem&\nolinkurl{PaperXI.TypeNormalizerAction.conjugation_inner}\\

34&def&\nolinkurl{PaperXI.TypeNormalizerAction.Stable}\\

37&theorem&\nolinkurl{PaperXI.TypeNormalizerAction.stable_map}\\

48&def&\nolinkurl{PaperXI.TypeNormalizerAction.action}\\

65&theorem&\nolinkurl{PaperXI.TypeNormalizerAction.action_mk}\\

72&theorem&\nolinkurl{PaperXI.TypeNormalizerAction.subgroup_le_kernel}\\

85&def&\nolinkurl{PaperXI.TypeNormalizerAction.quotientAction}\\

90&theorem&\nolinkurl{PaperXI.TypeNormalizerAction.quotientAction_mk}\\

96&theorem&\nolinkurl{PaperXI.TypeNormalizerAction.quotientAction_unique}\\

\bottomrule\end{longtable}

\section{\texorpdfstring{\nolinkurl{PaperXI.TypeNormalizerForm}}{PaperXI.TypeNormalizerForm}}\label{module:PaperXI.TypeNormalizerForm}

\textbf{Selected layer:} v0.44 parent.\quad\textbf{Source lines:} 85.\par\textbf{SHA-256:} {\scriptsize\nolinkurl{1f05cc7cd75bdb3f15a7b10c3d2153b865baf51da1f73e0a4fe87998c5d6969c}}.

\VerbatimInput[fontsize=\scriptsize,breaklines=true,breakanywhere=true]{excerpts/D_TypeNormalizerForm.txt}

\textbf{Direct project imports:} \nolinkurl{PaperXI.TypeNormalizerPairing}.

\begin{longtable}{@{}p{0.07\textwidth}p{0.12\textwidth}p{0.75\textwidth}@{}}\toprule Line&Kind&Name\\\midrule\endhead

15&def&\nolinkurl{PaperXI.TypeNormalizerForm.coordinateAction}\\

26&theorem&\nolinkurl{PaperXI.TypeNormalizerForm.coordinateAction_apply}\\

33&def&\nolinkurl{PaperXI.TypeNormalizerForm.linearAut}\\

44&def&\nolinkurl{PaperXI.TypeNormalizerForm.vectorAction}\\

66&theorem&\nolinkurl{PaperXI.TypeNormalizerForm.source_form_preserved}\\

\bottomrule\end{longtable}

\section{\texorpdfstring{\nolinkurl{PaperXI.TypeNormalizerPairing}}{PaperXI.TypeNormalizerPairing}}\label{module:PaperXI.TypeNormalizerPairing}

\textbf{Selected layer:} v0.44 parent.\quad\textbf{Source lines:} 154.\par\textbf{SHA-256:} {\scriptsize\nolinkurl{47cff0132a4d65fdb6193671065f95af0c45479300e676253cef05a50957f105}}.

\VerbatimInput[fontsize=\scriptsize,breaklines=true,breakanywhere=true]{excerpts/D_TypeNormalizerPairing.txt}

\textbf{Direct project imports:} \nolinkurl{PaperXI.TypeNormalizerAction}, \nolinkurl{PaperXI.TypeSideExtension}.

\begin{longtable}{@{}p{0.07\textwidth}p{0.12\textwidth}p{0.75\textwidth}@{}}\toprule Line&Kind&Name\\\midrule\endhead

17&def&\nolinkurl{PaperXI.TypeNormalizerPairing.subgroupAut}\\

26&def&\nolinkurl{PaperXI.TypeNormalizerPairing.CharacterInvariant}\\

32&theorem&\nolinkurl{PaperXI.TypeNormalizerPairing.inertiaInvariant}\\

43&theorem&\nolinkurl{PaperXI.TypeNormalizerPairing.characterCommutator_preserved}\\

61&def&\nolinkurl{PaperXI.TypeNormalizerPairing.pairing}\\

69&theorem&\nolinkurl{PaperXI.TypeNormalizerPairing.pairing_on_representatives}\\

77&theorem&\nolinkurl{PaperXI.TypeNormalizerPairing.quotientAction_preserves_pairing}\\

93&theorem&\nolinkurl{PaperXI.TypeNormalizerPairing.pairing_alternating}\\

98&theorem&\nolinkurl{PaperXI.TypeNormalizerPairing.pairing_mul_left}\\

106&theorem&\nolinkurl{PaperXI.TypeNormalizerPairing.pairing_mul_right}\\

119&theorem&\nolinkurl{PaperXI.TypeNormalizerPairing.bilinear_character_separates}\\

142&theorem&\nolinkurl{PaperXI.TypeNormalizerPairing.fieldForm_preserved}\\

\bottomrule\end{longtable}

\section{\texorpdfstring{\nolinkurl{PaperXI.TypeProjectiveImplementation}}{PaperXI.TypeProjectiveImplementation}}\label{module:PaperXI.TypeProjectiveImplementation}

\textbf{Selected layer:} v0.44 parent.\quad\textbf{Source lines:} 157.\par\textbf{SHA-256:} {\scriptsize\nolinkurl{3621d11e4ad48c3a7d4986dc55e21778d7bf926a4a008b6de91a38ccc21db617}}.

\VerbatimInput[fontsize=\scriptsize,breaklines=true,breakanywhere=true]{excerpts/D_TypeProjectiveImplementation.txt}

\textbf{Direct project imports:} \nolinkurl{PaperXI.TypeNormalizerPairing}, \nolinkurl{PaperXI.SchurScalar}, \nolinkurl{PaperXI.CliffordDatum}.

\begin{longtable}{@{}p{0.07\textwidth}p{0.12\textwidth}p{0.75\textwidth}@{}}\toprule Line&Kind&Name\\\midrule\endhead

17&def&\nolinkurl{PaperXI.TypeProjectiveImplementation.operator}\\

31&theorem&\nolinkurl{PaperXI.TypeProjectiveImplementation.operator_apply}\\

36&theorem&\nolinkurl{PaperXI.TypeProjectiveImplementation.implements_conjugation}\\

50&theorem&\nolinkurl{PaperXI.TypeProjectiveImplementation.implementer_unique}\\

60&theorem&\nolinkurl{PaperXI.TypeProjectiveImplementation.two_implementers}\\

73&theorem&\nolinkurl{PaperXI.TypeProjectiveImplementation.section_operators_mul}\\

81&theorem&\nolinkurl{PaperXI.TypeProjectiveImplementation.representative_change}\\

86&def&\nolinkurl{PaperXI.TypeProjectiveImplementation.ProjectiveDescends}\\

91&theorem&\nolinkurl{PaperXI.TypeProjectiveImplementation.projectiveDescends_iff}\\

113&theorem&\nolinkurl{PaperXI.TypeProjectiveImplementation.scalar_element_commutator}\\

127&theorem&\nolinkurl{PaperXI.TypeProjectiveImplementation.source_operator_commutator}\\

138&theorem&\nolinkurl{PaperXI.TypeProjectiveImplementation.nontrivial_pairing_obstructs_descent}\\

\bottomrule\end{longtable}

\section{\texorpdfstring{\nolinkurl{PaperXI.TypeSideDescent}}{PaperXI.TypeSideDescent}}\label{module:PaperXI.TypeSideDescent}

\textbf{Selected layer:} v0.44 parent.\quad\textbf{Source lines:} 293.\par\textbf{SHA-256:} {\scriptsize\nolinkurl{e066d77666e2671d02c454f4bd06e3650750cfcb143a6c4635339d8b6a0adb2f}}.

\VerbatimInput[fontsize=\scriptsize,breaklines=true,breakanywhere=true]{excerpts/D_TypeSideDescent.txt}

\textbf{Direct project imports:} \nolinkurl{PaperXI.QuotientRepresentation}, \nolinkurl{PaperXI.GeneratedIrreducibility}.

\begin{longtable}{@{}p{0.07\textwidth}p{0.12\textwidth}p{0.75\textwidth}@{}}\toprule Line&Kind&Name\\\midrule\endhead

18&def&\nolinkurl{PaperXI.TypeSideDescent.conjugate}\\

22&def&\nolinkurl{PaperXI.TypeSideDescent.Invariant}\\

26&def&\nolinkurl{PaperXI.TypeSideDescent.characterSpace}\\

39&def&\nolinkurl{PaperXI.TypeSideDescent.Contains}\\

43&def&\nolinkurl{PaperXI.TypeSideDescent.characterRepresentation}\\

53&theorem&\nolinkurl{PaperXI.TypeSideDescent.contains_iff_nonzero_intertwiner}\\

86&theorem&\nolinkurl{PaperXI.TypeSideDescent.characterSpace_stable}\\

99&def&\nolinkurl{PaperXI.TypeSideDescent.characterSubrepresentation}\\

105&theorem&\nolinkurl{PaperXI.TypeSideDescent.scalar_restriction}\\

121&def&\nolinkurl{PaperXI.TypeSideDescent.characterKernel}\\

124&theorem&\nolinkurl{PaperXI.TypeSideDescent.characterKernel_le}\\

129&theorem&\nolinkurl{PaperXI.TypeSideDescent.mem_characterKernel}\\

139&instance&\nolinkurl{PaperXI.TypeSideDescent.characterKernel_normal}\\

147&theorem&\nolinkurl{PaperXI.TypeSideDescent.characterKernel_le_ker}\\

157&def&\nolinkurl{PaperXI.TypeSideDescent.representation}\\

163&theorem&\nolinkurl{PaperXI.TypeSideDescent.representation_apply}\\

169&theorem&\nolinkurl{PaperXI.TypeSideDescent.representation_irreducible}\\

177&theorem&\nolinkurl{PaperXI.TypeSideDescent.finite_quotient}\\

188&def&\nolinkurl{PaperXI.TypeSideDescent.quotientProjection}\\

193&theorem&\nolinkurl{PaperXI.TypeSideDescent.quotientProjection_apply}\\

198&theorem&\nolinkurl{PaperXI.TypeSideDescent.quotientProjection_surjective}\\

203&theorem&\nolinkurl{PaperXI.TypeSideDescent.quotientProjection_ker}\\

210&theorem&\nolinkurl{PaperXI.TypeSideDescent.central_image}\\

229&def&\nolinkurl{PaperXI.TypeSideDescent.centralCharacter}\\

232&theorem&\nolinkurl{PaperXI.TypeSideDescent.centralCharacter_apply}\\

235&theorem&\nolinkurl{PaperXI.TypeSideDescent.centralCharacter_injective}\\

240&def&\nolinkurl{PaperXI.TypeSideDescent.centerInclusion}\\

248&theorem&\nolinkurl{PaperXI.TypeSideDescent.centerInclusion_apply}\\

253&theorem&\nolinkurl{PaperXI.TypeSideDescent.centerInclusion_injective}\\

260&theorem&\nolinkurl{PaperXI.TypeSideDescent.centerInclusion_range}\\

272&theorem&\nolinkurl{PaperXI.TypeSideDescent.exact_center_quotient}\\

277&theorem&\nolinkurl{PaperXI.TypeSideDescent.centerInclusion_central}\\

285&theorem&\nolinkurl{PaperXI.TypeSideDescent.central_action}\\

\bottomrule\end{longtable}

\section{\texorpdfstring{\nolinkurl{PaperXI.TypeSideExtension}}{PaperXI.TypeSideExtension}}\label{module:PaperXI.TypeSideExtension}

\textbf{Selected layer:} v0.44 parent.\quad\textbf{Source lines:} 170.\par\textbf{SHA-256:} {\scriptsize\nolinkurl{f00297a9c260398048b04cca886e47dbd0ac912e029104cee4892063e5aab11c}}.

\VerbatimInput[fontsize=\scriptsize,breaklines=true,breakanywhere=true]{excerpts/D_TypeSideExtension.txt}

\textbf{Direct project imports:} \nolinkurl{PaperXI.TypeSideDescent}, \nolinkurl{PaperXI.CentralExtensionCocycle}.

\begin{longtable}{@{}p{0.07\textwidth}p{0.12\textwidth}p{0.75\textwidth}@{}}\toprule Line&Kind&Name\\\midrule\endhead

20&instance&\nolinkurl{PaperXI.TypeSideExtension.centralQuotient_isMulCommutative}\\

24&theorem&\nolinkurl{PaperXI.TypeSideExtension.finite_center}\\

27&theorem&\nolinkurl{PaperXI.TypeSideExtension.centralCharacter_ne_one}\\

36&def&\nolinkurl{PaperXI.TypeSideExtension.datum}\\

48&theorem&\nolinkurl{PaperXI.TypeSideExtension.commutator_mem}\\

54&def&\nolinkurl{PaperXI.TypeSideExtension.characterCommutator}\\

60&theorem&\nolinkurl{PaperXI.TypeSideExtension.scalar_commutator_source}\\

81&theorem&\nolinkurl{PaperXI.TypeSideExtension.scalar_commutator_eq_pairing}\\

94&theorem&\nolinkurl{PaperXI.TypeSideExtension.nondegenerate_of_pairing}\\

117&theorem&\nolinkurl{PaperXI.TypeSideExtension.bilinear_character_nondegenerate}\\

138&theorem&\nolinkurl{PaperXI.TypeSideExtension.nondegenerate_of_bilinear_pairing}\\

162&theorem&\nolinkurl{PaperXI.TypeSideExtension.representation_central_action}\\

\bottomrule\end{longtable}

\section{\texorpdfstring{\nolinkurl{PaperXI.TypeSideUniqueness}}{PaperXI.TypeSideUniqueness}}\label{module:PaperXI.TypeSideUniqueness}

\textbf{Selected layer:} v0.44 parent.\quad\textbf{Source lines:} 96.\par\textbf{SHA-256:} {\scriptsize\nolinkurl{3a4c2a7589f7cccf49eec7161e463a3c6d21f217693a596505efa5d339330a63}}.

\VerbatimInput[fontsize=\scriptsize,breaklines=true,breakanywhere=true]{excerpts/D_TypeSideUniqueness.txt}

\textbf{Direct project imports:} \nolinkurl{PaperXI.TypeSideExtension}, \nolinkurl{PaperXI.CentralExtensionUniqueness}.

\begin{longtable}{@{}p{0.07\textwidth}p{0.12\textwidth}p{0.75\textwidth}@{}}\toprule Line&Kind&Name\\\midrule\endhead

21&theorem&\nolinkurl{PaperXI.TypeSideUniqueness.finiteDimensional}\\

34&theorem&\nolinkurl{PaperXI.TypeSideUniqueness.finrank_sqrt}\\

50&theorem&\nolinkurl{PaperXI.TypeSideUniqueness.exists_equiv}\\

71&theorem&\nolinkurl{PaperXI.TypeSideUniqueness.exists_unique_intertwiner}\\

\bottomrule\end{longtable}

\section{\texorpdfstring{\nolinkurl{PaperXI.UnramifiedBaseChangeFactor}}{PaperXI.UnramifiedBaseChangeFactor}}\label{module:PaperXI.UnramifiedBaseChangeFactor}

\textbf{Selected layer:} v0.44 parent.\quad\textbf{Source lines:} 159.\par\textbf{SHA-256:} {\scriptsize\nolinkurl{0f6e188bb24bf02711cdd8d676182d365a8d0fac73456675899fa0f4564a7340}}.

\VerbatimInput[fontsize=\scriptsize,breaklines=true,breakanywhere=true]{excerpts/D_UnramifiedBaseChangeFactor.txt}

\textbf{Direct project imports:} \nolinkurl{PaperXI.RankOneEulerFactor}.

\begin{longtable}{@{}p{0.07\textwidth}p{0.12\textwidth}p{0.75\textwidth}@{}}\toprule Line&Kind&Name\\\midrule\endhead

30&def&\nolinkurl{PaperXI.FiniteUnramifiedWeil.integerDegreeDilation}\\

38&def&\nolinkurl{PaperXI.FiniteUnramifiedWeil.degreeDilation}\\

43&theorem&\nolinkurl{PaperXI.FiniteUnramifiedWeil.degreeEquiv_degreeDilation}\\

49&theorem&\nolinkurl{PaperXI.FiniteUnramifiedWeil.degree_degreeDilation}\\

61&theorem&\nolinkurl{PaperXI.FiniteUnramifiedWeil.degreeDilation_geometricLift}\\

71&theorem&\nolinkurl{PaperXI.FiniteUnramifiedWeil.degreeDilation_comp}\\

84&def&\nolinkurl{PaperXI.FiniteUnramifiedWeil.baseChangeRepresentation}\\

91&theorem&\nolinkurl{PaperXI.FiniteUnramifiedWeil.baseChange_geometricFrobenius}\\

100&theorem&\nolinkurl{PaperXI.FiniteUnramifiedWeil.baseChangeRepresentation_comp}\\

113&theorem&\nolinkurl{PaperXI.FiniteUnramifiedWeil.baseChange_frobeniusPolynomial}\\

121&theorem&\nolinkurl{PaperXI.FiniteUnramifiedWeil.baseChange_localFactor}\\

131&theorem&\nolinkurl{PaperXI.FiniteUnramifiedWeil.irreducible_baseChange_frobeniusPolynomial}\\

151&theorem&\nolinkurl{PaperXI.FiniteUnramifiedWeil.irreducible_baseChange_localFactor}\\

\bottomrule\end{longtable}

\section{\texorpdfstring{\nolinkurl{PaperXI.UnramifiedConductor}}{PaperXI.UnramifiedConductor}}\label{module:PaperXI.UnramifiedConductor}

\textbf{Selected layer:} v0.44 parent.\quad\textbf{Source lines:} 131.\par\textbf{SHA-256:} {\scriptsize\nolinkurl{1a940bc977e7fb0565338370fce87b395fc254c3cbbc39c220d61eef0a6de682}}.

\VerbatimInput[fontsize=\scriptsize,breaklines=true,breakanywhere=true]{excerpts/D_UnramifiedConductor.txt}

\textbf{Direct project imports:} \nolinkurl{PaperXI.ResidueDegreeBaseChange}, \nolinkurl{PaperXI.ValuationConductor}.

\begin{longtable}{@{}p{0.07\textwidth}p{0.12\textwidth}p{0.75\textwidth}@{}}\toprule Line&Kind&Name\\\midrule\endhead

18&theorem&\nolinkurl{PaperXI.UnramifiedConductor.lower_eq_comap}\\

21&def&\nolinkurl{PaperXI.UnramifiedConductor.lowerEquiv}\\

28&theorem&\nolinkurl{PaperXI.UnramifiedConductor.lower_card}\\

33&def&\nolinkurl{PaperXI.UnramifiedConductor.embedding}\\

47&theorem&\nolinkurl{PaperXI.UnramifiedConductor.embedding_degree}\\

52&theorem&\nolinkurl{PaperXI.UnramifiedConductor.embedding_injective}\\

65&theorem&\nolinkurl{PaperXI.UnramifiedConductor.embedding_lift}\\

77&theorem&\nolinkurl{PaperXI.UnramifiedConductor.fixed_eq}\\

95&theorem&\nolinkurl{PaperXI.UnramifiedConductor.summand_eq}\\

103&theorem&\nolinkurl{PaperXI.UnramifiedConductor.swan_eq}\\

112&theorem&\nolinkurl{PaperXI.UnramifiedConductor.artin_eq}\\

120&theorem&\nolinkurl{PaperXI.UnramifiedConductor.summable_restriction}\\

\bottomrule\end{longtable}

\section{\texorpdfstring{\nolinkurl{PaperXI.UnramifiedConductorZero}}{PaperXI.UnramifiedConductorZero}}\label{module:PaperXI.UnramifiedConductorZero}

\textbf{Selected layer:} v0.44 parent.\quad\textbf{Source lines:} 109.\par\textbf{SHA-256:} {\scriptsize\nolinkurl{35a3f791e30467fdb145456cea233a2a85f3fc669b8cc9a5216db7b5a705f19d}}.

\VerbatimInput[fontsize=\scriptsize,breaklines=true,breakanywhere=true]{excerpts/D_UnramifiedConductorZero.txt}

\textbf{Direct project imports:} \nolinkurl{PaperXI.UnramifiedLocalFactor}, \nolinkurl{PaperXI.ValuationConductor}.

\begin{longtable}{@{}p{0.07\textwidth}p{0.12\textwidth}p{0.75\textwidth}@{}}\toprule Line&Kind&Name\\\midrule\endhead

34&theorem&\nolinkurl{PaperXI.ValuationConductor.artin_zero_of_degree_kernel_bot}\\

73&def&\nolinkurl{PaperXI.FiniteUnramifiedWeil.inertiaCodimension}\\

80&def&\nolinkurl{PaperXI.FiniteUnramifiedWeil.swanConductorExponent}\\

84&def&\nolinkurl{PaperXI.FiniteUnramifiedWeil.artinConductorExponent}\\

89&theorem&\nolinkurl{PaperXI.FiniteUnramifiedWeil.inertiaCodimension_eq_zero}\\

97&theorem&\nolinkurl{PaperXI.FiniteUnramifiedWeil.swanConductorExponent_eq_zero}\\

103&theorem&\nolinkurl{PaperXI.FiniteUnramifiedWeil.artinConductorExponent_eq_zero}\\

\bottomrule\end{longtable}

\section{\texorpdfstring{\nolinkurl{PaperXI.UnramifiedFieldTower}}{PaperXI.UnramifiedFieldTower}}\label{module:PaperXI.UnramifiedFieldTower}

\textbf{Selected layer:} v0.44 parent.\quad\textbf{Source lines:} 270.\par\textbf{SHA-256:} {\scriptsize\nolinkurl{79699a103ad55a1254dbaa1ea632d59bd77ce2f4ac03367b50d40d1b35f287ee}}.

\VerbatimInput[fontsize=\scriptsize,breaklines=true,breakanywhere=true]{excerpts/D_UnramifiedFieldTower.txt}

\textbf{Direct project imports:} \nolinkurl{PaperXI.UnramifiedBaseChangeFactor}.

\begin{longtable}{@{}p{0.07\textwidth}p{0.12\textwidth}p{0.75\textwidth}@{}}\toprule Line&Kind&Name\\\midrule\endhead

32&abbrev&\nolinkurl{PaperXI.UnramifiedFieldTower.LowerField}\\

36&def&\nolinkurl{PaperXI.UnramifiedFieldTower.residueTowerMap}\\

53&abbrev&\nolinkurl{PaperXI.UnramifiedFieldTower.UpperField}\\

69&def&\nolinkurl{PaperXI.UnramifiedFieldTower.lowerToUpper}\\

75&theorem&\nolinkurl{PaperXI.UnramifiedFieldTower.residue_finrank}\\

83&theorem&\nolinkurl{PaperXI.UnramifiedFieldTower.relative_degree}\\

90&theorem&\nolinkurl{PaperXI.UnramifiedFieldTower.total_degree}\\

98&def&\nolinkurl{PaperXI.UnramifiedFieldTower.upperResidueAlgEquiv}\\

104&theorem&\nolinkurl{PaperXI.UnramifiedFieldTower.relative_ramificationIdx_one}\\

111&def&\nolinkurl{PaperXI.UnramifiedFieldTower.restriction}\\

120&theorem&\nolinkurl{PaperXI.UnramifiedFieldTower.restriction_degree_square}\\

127&theorem&\nolinkurl{PaperXI.UnramifiedFieldTower.restriction_degree}\\

136&theorem&\nolinkurl{PaperXI.UnramifiedFieldTower.restriction_geometricLift}\\

159&theorem&\nolinkurl{PaperXI.UnramifiedFieldTower.restriction_injective}\\

187&theorem&\nolinkurl{PaperXI.UnramifiedFieldTower.restriction_range}\\

217&def&\nolinkurl{PaperXI.UnramifiedFieldTower.restrictionEquivPeriod}\\

228&theorem&\nolinkurl{PaperXI.UnramifiedFieldTower.restriction_unique}\\

246&def&\nolinkurl{PaperXI.UnramifiedFieldTower.restrictRepresentation}\\

253&theorem&\nolinkurl{PaperXI.UnramifiedFieldTower.restrictRepresentation_geometricLift}\\

261&theorem&\nolinkurl{PaperXI.UnramifiedFieldTower.restriction_frobeniusPolynomial}\\

\bottomrule\end{longtable}

\section{\texorpdfstring{\nolinkurl{PaperXI.UnramifiedLocalFactor}}{PaperXI.UnramifiedLocalFactor}}\label{module:PaperXI.UnramifiedLocalFactor}

\textbf{Selected layer:} v0.44 parent.\quad\textbf{Source lines:} 168.\par\textbf{SHA-256:} {\scriptsize\nolinkurl{ca0d3e6fba2b3caa2549ac65cd50e4fe66b15dc4cf96c9398531f826b265ebbd}}.

\VerbatimInput[fontsize=\scriptsize,breaklines=true,breakanywhere=true]{excerpts/D_UnramifiedLocalFactor.txt}

\textbf{Direct project imports:} \nolinkurl{PaperXI.FiniteUnramifiedWeil}.

\begin{longtable}{@{}p{0.07\textwidth}p{0.12\textwidth}p{0.75\textwidth}@{}}\toprule Line&Kind&Name\\\midrule\endhead

28&abbrev&\nolinkurl{PaperXI.FiniteUnramifiedWeil.Inertia}\\

33&theorem&\nolinkurl{PaperXI.FiniteUnramifiedWeil.inertia_le_kernel}\\

39&theorem&\nolinkurl{PaperXI.FiniteUnramifiedWeil.inertia_invariants_eq_top}\\

47&def&\nolinkurl{PaperXI.FiniteUnramifiedWeil.frobeniusPolynomial}\\

52&def&\nolinkurl{PaperXI.FiniteUnramifiedWeil.unramifiedLocalFactor}\\

59&def&\nolinkurl{PaperXI.FiniteUnramifiedWeil.inertiaFixedEquiv}\\

64&def&\nolinkurl{PaperXI.FiniteUnramifiedWeil.fixedGeometricFrobenius}\\

70&theorem&\nolinkurl{PaperXI.FiniteUnramifiedWeil.fixedGeometricFrobenius_charpoly}\\

77&def&\nolinkurl{PaperXI.FiniteUnramifiedWeil.inertiaFrobeniusPolynomial}\\

83&theorem&\nolinkurl{PaperXI.FiniteUnramifiedWeil.inertiaFrobeniusPolynomial_eq}\\

90&def&\nolinkurl{PaperXI.FiniteUnramifiedWeil.inertiaLocalFactor}\\

95&theorem&\nolinkurl{PaperXI.FiniteUnramifiedWeil.inertiaLocalFactor_eq}\\

103&theorem&\nolinkurl{PaperXI.FiniteUnramifiedWeil.frobeniusPolynomial_eq_det}\\

114&theorem&\nolinkurl{PaperXI.FiniteUnramifiedWeil.frobeniusPolynomial_coeff_zero}\\

121&theorem&\nolinkurl{PaperXI.FiniteUnramifiedWeil.frobeniusPolynomial_eval_zero}\\

127&theorem&\nolinkurl{PaperXI.FiniteUnramifiedWeil.frobeniusPolynomial_ne_zero}\\

136&theorem&\nolinkurl{PaperXI.FiniteUnramifiedWeil.mappedFrobeniusPolynomial_ne_zero}\\

143&theorem&\nolinkurl{PaperXI.FiniteUnramifiedWeil.unramifiedLocalFactor_mul_polynomial}\\

151&theorem&\nolinkurl{PaperXI.FiniteUnramifiedWeil.polynomial_mul_unramifiedLocalFactor}\\

159&theorem&\nolinkurl{PaperXI.FiniteUnramifiedWeil.unramifiedLocalFactor_ne_zero}\\

\bottomrule\end{longtable}

\section{\texorpdfstring{\nolinkurl{PaperXI.UnramifiedPolynomial}}{PaperXI.UnramifiedPolynomial}}\label{module:PaperXI.UnramifiedPolynomial}

\textbf{Selected layer:} v0.44 parent.\quad\textbf{Source lines:} 90.\par\textbf{SHA-256:} {\scriptsize\nolinkurl{2dfea600d7050f62c8adde4b0242ca02748ec7427f610c1876bc40b1d7a2f64b}}.

\VerbatimInput[fontsize=\scriptsize,breaklines=true,breakanywhere=true]{excerpts/D_UnramifiedPolynomial.txt}

\textbf{Direct project imports:} \nolinkurl{PaperXI.MonicResidueAlgebra}.

\begin{longtable}{@{}p{0.07\textwidth}p{0.12\textwidth}p{0.75\textwidth}@{}}\toprule Line&Kind&Name\\\midrule\endhead

21&theorem&\nolinkurl{PaperXI.UnramifiedPolynomial.irreducible}\\

24&instance&\nolinkurl{PaperXI.UnramifiedPolynomial.domain}\\

27&instance&\nolinkurl{PaperXI.UnramifiedPolynomial.faithful}\\

30&theorem&\nolinkurl{PaperXI.UnramifiedPolynomial.not_field}\\

37&instance&\nolinkurl{PaperXI.UnramifiedPolynomial.dvr}\\

44&theorem&\nolinkurl{PaperXI.UnramifiedPolynomial.uniformizer}\\

50&abbrev&\nolinkurl{PaperXI.UnramifiedPolynomial.field}\\

52&instance&\nolinkurl{PaperXI.UnramifiedPolynomial.baseFieldAlgebra}\\

55&instance&\nolinkurl{PaperXI.UnramifiedPolynomial.baseFieldTower}\\

58&theorem&\nolinkurl{PaperXI.UnramifiedPolynomial.natDegree_pos_of_residue_irreducible}\\

62&theorem&\nolinkurl{PaperXI.UnramifiedPolynomial.field_degree}\\

66&instance&\nolinkurl{PaperXI.UnramifiedPolynomial.finiteFieldExtension}\\

71&theorem&\nolinkurl{PaperXI.UnramifiedPolynomial.integralClosure}\\

78&theorem&\nolinkurl{PaperXI.UnramifiedPolynomial.etale}\\

85&theorem&\nolinkurl{PaperXI.UnramifiedPolynomial.ramificationIdx_one}\\

\bottomrule\end{longtable}

\section{\texorpdfstring{\nolinkurl{PaperXI.UnramifiedValuation}}{PaperXI.UnramifiedValuation}}\label{module:PaperXI.UnramifiedValuation}

\textbf{Selected layer:} v0.44 parent.\quad\textbf{Source lines:} 145.\par\textbf{SHA-256:} {\scriptsize\nolinkurl{001406e3abfac1991c15f4006373c9a8673cd0e8e16df3d80a9ef6c276a50054}}.

\VerbatimInput[fontsize=\scriptsize,breaklines=true,breakanywhere=true]{excerpts/D_UnramifiedValuation.txt}

\textbf{Direct project imports:} \nolinkurl{PaperXI.FiniteResidueSubgroup}, \nolinkurl{PaperXI.UnramifiedConductor}, \nolinkurl{PaperXI.ValuationTameEndpoint}.

\begin{longtable}{@{}p{0.07\textwidth}p{0.12\textwidth}p{0.75\textwidth}@{}}\toprule Line&Kind&Name\\\midrule\endhead

22&abbrev&\nolinkurl{PaperXI.UnramifiedValuation.fixingGroup}\\

27&theorem&\nolinkurl{PaperXI.UnramifiedValuation.frobenius_agrees}\\

32&theorem&\nolinkurl{PaperXI.UnramifiedValuation.contains_inertia}\\

37&def&\nolinkurl{PaperXI.UnramifiedValuation.embedding}\\

45&theorem&\nolinkurl{PaperXI.UnramifiedValuation.degree_formula}\\

54&def&\nolinkurl{PaperXI.UnramifiedValuation.degreeSubgroupEquiv}\\

61&theorem&\nolinkurl{PaperXI.UnramifiedValuation.embedding_injective}\\

66&theorem&\nolinkurl{PaperXI.UnramifiedValuation.embedding_range}\\

93&theorem&\nolinkurl{PaperXI.UnramifiedValuation.degree_index}\\

103&theorem&\nolinkurl{PaperXI.UnramifiedValuation.norm_compatibility}\\

118&theorem&\nolinkurl{PaperXI.UnramifiedValuation.source}\\

\bottomrule\end{longtable}

\section{\texorpdfstring{\nolinkurl{PaperXI.ValuationConductor}}{PaperXI.ValuationConductor}}\label{module:PaperXI.ValuationConductor}

\textbf{Selected layer:} v0.44 parent.\quad\textbf{Source lines:} 167.\par\textbf{SHA-256:} {\scriptsize\nolinkurl{3040c48c4558f4796813146cac1147f5d7bf87a20ff05cabecf53a37fce73cb0}}.

\VerbatimInput[fontsize=\scriptsize,breaklines=true,breakanywhere=true]{excerpts/D_ValuationConductor.txt}

\textbf{Direct project imports:} \nolinkurl{PaperXI.ResidueWeil}, \nolinkurl{PaperXI.RamificationSeparation}.

\begin{longtable}{@{}p{0.07\textwidth}p{0.12\textwidth}p{0.75\textwidth}@{}}\toprule Line&Kind&Name\\\midrule\endhead

19&abbrev&\nolinkurl{PaperXI.ValuationConductor.residue}\\

23&def&\nolinkurl{PaperXI.ValuationConductor.lift}\\

32&theorem&\nolinkurl{PaperXI.ValuationConductor.lift_degree}\\

38&def&\nolinkurl{PaperXI.ValuationConductor.fixed}\\

43&theorem&\nolinkurl{PaperXI.ValuationConductor.fixed_zero}\\

61&def&\nolinkurl{PaperXI.ValuationConductor.codimension}\\

66&def&\nolinkurl{PaperXI.ValuationConductor.summand}\\

70&def&\nolinkurl{PaperXI.ValuationConductor.swan}\\

74&def&\nolinkurl{PaperXI.ValuationConductor.artin}\\

79&theorem&\nolinkurl{PaperXI.ValuationConductor.positive_fixed_top}\\

88&theorem&\nolinkurl{PaperXI.ValuationConductor.tame_summand_zero}\\

96&theorem&\nolinkurl{PaperXI.ValuationConductor.tame_swan_zero}\\

102&theorem&\nolinkurl{PaperXI.ValuationConductor.tame_irreducible_dimension}\\

116&theorem&\nolinkurl{PaperXI.ValuationConductor.tame_artin_dichotomy}\\

127&theorem&\nolinkurl{PaperXI.ValuationConductor.fixed_top_of_lower_bot}\\

138&theorem&\nolinkurl{PaperXI.ValuationConductor.summable_summands}\\

151&theorem&\nolinkurl{PaperXI.ValuationConductor.zeroth_summand}\\

160&theorem&\nolinkurl{PaperXI.ValuationConductor.artin_eq_full_sum}\\

\bottomrule\end{longtable}

\section{\texorpdfstring{\nolinkurl{PaperXI.ValuationFixingField}}{PaperXI.ValuationFixingField}}\label{module:PaperXI.ValuationFixingField}

\textbf{Selected layer:} v0.44 parent.\quad\textbf{Source lines:} 41.\par\textbf{SHA-256:} {\scriptsize\nolinkurl{fdaeae5757fb111b416bf264d6b76e43cead77d0d338c4bc16110191bedaeb06}}.

\VerbatimInput[fontsize=\scriptsize,breaklines=true,breakanywhere=true]{excerpts/D_ValuationFixingField.txt}

\textbf{Direct project imports:} \nolinkurl{PaperXI.HenselianFieldLift}.

\begin{longtable}{@{}p{0.07\textwidth}p{0.12\textwidth}p{0.75\textwidth}@{}}\toprule Line&Kind&Name\\\midrule\endhead

12&def&\nolinkurl{PaperXI.ValuationFixingField.decompositionEquiv}\\

27&theorem&\nolinkurl{PaperXI.ValuationFixingField.decompositionEquiv_apply}\\

31&theorem&\nolinkurl{PaperXI.ValuationFixingField.ring_action}\\

35&theorem&\nolinkurl{PaperXI.ValuationFixingField.residue_action}\\

\bottomrule\end{longtable}

\section{\texorpdfstring{\nolinkurl{PaperXI.ValuationIntegralModel}}{PaperXI.ValuationIntegralModel}}\label{module:PaperXI.ValuationIntegralModel}

\textbf{Selected layer:} v0.44 parent.\quad\textbf{Source lines:} 72.\par\textbf{SHA-256:} {\scriptsize\nolinkurl{c126174463d0f9039e8425c76709f0ab71b7babdb8e5ed6b19aa18a6169ffd3c}}.

\VerbatimInput[fontsize=\scriptsize,breaklines=true,breakanywhere=true]{excerpts/D_ValuationIntegralModel.txt}

\textbf{Direct project imports:} \nolinkurl{PaperXI.HenselianGaloisResidue}.

\begin{longtable}{@{}p{0.07\textwidth}p{0.12\textwidth}p{0.75\textwidth}@{}}\toprule Line&Kind&Name\\\midrule\endhead

13&abbrev&\nolinkurl{PaperXI.ValuationIntegralModel.base}\\

15&instance&\nolinkurl{PaperXI.ValuationIntegralModel.baseAlgebra}\\

17&instance&\nolinkurl{PaperXI.ValuationIntegralModel.baseTower}\\

19&instance&\nolinkurl{PaperXI.ValuationIntegralModel.baseFieldTower}\\

21&instance&\nolinkurl{PaperXI.ValuationIntegralModel.baseLocal}\\

25&def&\nolinkurl{PaperXI.ValuationIntegralModel.ringAction}\\

31&theorem&\nolinkurl{PaperXI.ValuationIntegralModel.invariant}\\

50&theorem&\nolinkurl{PaperXI.ValuationIntegralModel.integral}\\

56&theorem&\nolinkurl{PaperXI.ValuationIntegralModel.integralClosure}\\

62&theorem&\nolinkurl{PaperXI.ValuationIntegralModel.finite}\\

67&theorem&\nolinkurl{PaperXI.ValuationIntegralModel.rank}\\

\bottomrule\end{longtable}

\section{\texorpdfstring{\nolinkurl{PaperXI.ValuationRamification}}{PaperXI.ValuationRamification}}\label{module:PaperXI.ValuationRamification}

\textbf{Selected layer:} v0.44 parent.\quad\textbf{Source lines:} 96.\par\textbf{SHA-256:} {\scriptsize\nolinkurl{92ddb509dc327389c2168f44ddfd5f7e16f4315af16835a54398077fe81146ca}}.

\VerbatimInput[fontsize=\scriptsize,breaklines=true,breakanywhere=true]{excerpts/D_ValuationRamification.txt}

\textbf{Direct project imports:} none.

\begin{longtable}{@{}p{0.07\textwidth}p{0.12\textwidth}p{0.75\textwidth}@{}}\toprule Line&Kind&Name\\\midrule\endhead

15&def&\nolinkurl{PaperXI.ValuationRamification.lower}\\

30&theorem&\nolinkurl{PaperXI.ValuationRamification.mem_lower}\\

34&theorem&\nolinkurl{PaperXI.ValuationRamification.smul_mem_power}\\

48&instance&\nolinkurl{PaperXI.ValuationRamification.lower_normal}\\

55&theorem&\nolinkurl{PaperXI.ValuationRamification.antitone}\\

60&theorem&\nolinkurl{PaperXI.ValuationRamification.lower_zero_eq_residue_kernel}\\

80&theorem&\nolinkurl{PaperXI.ValuationRamification.lower_le_inertia}\\

85&theorem&\nolinkurl{PaperXI.ValuationRamification.positive_le_one}\\

91&theorem&\nolinkurl{PaperXI.ValuationRamification.valuation_lower_zero}\\

\bottomrule\end{longtable}

\section{\texorpdfstring{\nolinkurl{PaperXI.ValuationResidueEmbedding}}{PaperXI.ValuationResidueEmbedding}}\label{module:PaperXI.ValuationResidueEmbedding}

\textbf{Selected layer:} v0.44 parent.\quad\textbf{Source lines:} 48.\par\textbf{SHA-256:} {\scriptsize\nolinkurl{ee803ae38b207f0084f4fe3d41d3b6348f38ade5e33a9ca430804b2e35a4c98f}}.

\VerbatimInput[fontsize=\scriptsize,breaklines=true,breakanywhere=true]{excerpts/D_ValuationResidueEmbedding.txt}

\textbf{Direct project imports:} \nolinkurl{PaperXI.HenselianFieldLift}.

\begin{longtable}{@{}p{0.07\textwidth}p{0.12\textwidth}p{0.75\textwidth}@{}}\toprule Line&Kind&Name\\\midrule\endhead

10&def&\nolinkurl{PaperXI.ValuationResidueEmbedding.inclusion}\\

16&theorem&\nolinkurl{PaperXI.ValuationResidueEmbedding.inclusion_val}\\

18&instance&\nolinkurl{PaperXI.ValuationResidueEmbedding.inclusion_local}\\

33&def&\nolinkurl{PaperXI.ValuationResidueEmbedding.residueMap}\\

35&theorem&\nolinkurl{PaperXI.ValuationResidueEmbedding.residueMap_residue}\\

37&theorem&\nolinkurl{PaperXI.ValuationResidueEmbedding.residueMap_injective}\\

39&theorem&\nolinkurl{PaperXI.ValuationResidueEmbedding.mem_range_iff}\\

\bottomrule\end{longtable}

\section{\texorpdfstring{\nolinkurl{PaperXI.ValuationTameEndpoint}}{PaperXI.ValuationTameEndpoint}}\label{module:PaperXI.ValuationTameEndpoint}

\textbf{Selected layer:} v0.44 parent.\quad\textbf{Source lines:} 52.\par\textbf{SHA-256:} {\scriptsize\nolinkurl{613a021b32c397cfc2425e8ee5c9cb796856bc8643069f24fde8b8f634df23d9}}.

\VerbatimInput[fontsize=\scriptsize,breaklines=true,breakanywhere=true]{excerpts/D_ValuationTameEndpoint.txt}

\textbf{Direct project imports:} \nolinkurl{PaperXI.ValuationConductor}.

\begin{longtable}{@{}p{0.07\textwidth}p{0.12\textwidth}p{0.75\textwidth}@{}}\toprule Line&Kind&Name\\\midrule\endhead

20&def&\nolinkurl{PaperXI.ValuationTameEndpoint.geometricFrobenius}\\

24&theorem&\nolinkurl{PaperXI.ValuationTameEndpoint.geometric_power}\\

28&abbrev&\nolinkurl{PaperXI.ValuationTameEndpoint.WeilPullback}\\

30&def&\nolinkurl{PaperXI.ValuationTameEndpoint.inertiaIdentification}\\

38&theorem&\nolinkurl{PaperXI.ValuationTameEndpoint.conductor}\\

\bottomrule\end{longtable}

\chapter{Navigation across the original 66 obligations}

This navigation table extracts the original 66 proof-obligation records from the historical register and gives an entry point into each corresponding source family. Entry modules lead to the definitions and related results for the whole section. They are not assertions that one module alone proves every displayed claim. The full source hypotheses and relevant ordinary mathematical attachments remain decisive.

The historical register contains 102 numbered records, including definitions, imported facts, examples and remarks. Original obligation 66 is printed Theorem 16.2, all-records row 95. Its historical ``proved internally'' flag is superseded in this reading edition by the conjectural-realization and conditional-comparison split. The other historical record types are excluded from this 66-row table. The printed numbers in this historical column now link to the corresponding current monograph destinations. Historical PDF page numbers remain only in the unchanged JSON crosswalk. In the last row, historical 16.2 links to Conjecture 66. The separate conditional comparison is \href{https://komorebi.invalid/b07/current_comparison}{the comparison of one realized candidate}. These are reading links, not new theorem-equivalence or compiler certificates.

\begin{longtable}{@{}p{0.05\textwidth}p{0.06\textwidth}p{0.09\textwidth}>{\raggedright\arraybackslash}p{0.34\textwidth}p{0.32\textwidth}@{}}\toprule Obl.&Row&Printed&Original title&Source family entry\\\midrule\endhead

1&3&\href{https://komorebi.invalid/b07/obligation_01}{2.3}&Exact block reconstruction&\nolinkurl{PaperXI.AtomicReconstruction}\\

2&4&\href{https://komorebi.invalid/b07/obligation_02}{2.4}&Scalar covariance model&\nolinkurl{PaperXI.AtomicReconstruction}\\

3&5&\href{https://komorebi.invalid/b07/obligation_03}{2.5}&Formal atomic reconstruction&\nolinkurl{PaperXI.AtomicReconstruction}\\

4&7&\href{https://komorebi.invalid/b07/obligation_04}{3.1}&Tame character projectors&\nolinkurl{PaperXI.TameProjectors}\\

5&8&\href{https://komorebi.invalid/b07/obligation_05}{3.2}&Tame orbit induction&\nolinkurl{PaperXI.TameProjectors}\\

6&9&\href{https://komorebi.invalid/b07/obligation_06}{3.3}&Dimension and regularity&\nolinkurl{PaperXI.TameProjectors}\\

7&10&\href{https://komorebi.invalid/b07/obligation_07}{3.4}&Dual and self dual tame packets&\nolinkurl{PaperXI.TameProjectors}\\

8&11&\href{https://komorebi.invalid/b07/obligation_08}{3.5}&Unramified twist&\nolinkurl{PaperXI.TameProjectors}\\

9&12&\href{https://komorebi.invalid/b07/obligation_09}{3.6}&Möbius orbit count&\nolinkurl{PaperXI.TameProjectors}\\

10&15&\href{https://komorebi.invalid/b07/obligation_10}{4.2}&Primitive cuspidal idempotent&\nolinkurl{PaperXI.CharacterProjector}\\

11&16&\href{https://komorebi.invalid/b07/obligation_11}{4.3}&Parabolic annihilation&\nolinkurl{PaperXI.CharacterProjector}\\

12&17&\href{https://komorebi.invalid/b07/obligation_12}{4.4}&Frobenius cycle polynomial&\nolinkurl{PaperXI.CharacterProjector}\\

13&18&\href{https://komorebi.invalid/b07/obligation_13}{4.5}&Unramified orientation character&\nolinkurl{PaperXI.CharacterProjector}\\

14&19&\href{https://komorebi.invalid/b07/obligation_14}{5.1}&Greatest common divisor and least common multiple packets&\nolinkurl{PaperXI.TameTensorSource}\\

15&20&\href{https://komorebi.invalid/b07/obligation_15}{5.2}&Packet count and depth&\nolinkurl{PaperXI.TameTensorSource}\\

16&21&\href{https://komorebi.invalid/b07/obligation_16}{5.3}&Refined exponent formula&\nolinkurl{PaperXI.TameTensorSource}\\

17&22&\href{https://komorebi.invalid/b07/obligation_17}{5.4}&Local polynomial of one tame packet&\nolinkurl{PaperXI.TameTensorSource}\\

18&23&\href{https://komorebi.invalid/b07/obligation_18}{5.5}&Tame tensor local factor&\nolinkurl{PaperXI.TameTensorSource}\\

19&24&\href{https://komorebi.invalid/b07/obligation_19}{5.6}&Trivial tensor packet criterion&\nolinkurl{PaperXI.TameTensorSource}\\

20&25&\href{https://komorebi.invalid/b07/obligation_20}{5.7}&Tame Hom test&\nolinkurl{PaperXI.TameTensorSource}\\

21&26&\href{https://komorebi.invalid/b07/obligation_21}{5.8}&Conductor of a nontrivial tame atom&\nolinkurl{PaperXI.TameTensorSource}\\

22&27&\href{https://komorebi.invalid/b07/obligation_22}{6.1}&Gauss magnitude&\nolinkurl{PaperXI.GaussCharacters}\\

23&29&\href{https://komorebi.invalid/b07/obligation_23}{6.3}&Duality of Gauss phases&\nolinkurl{PaperXI.GaussCharacters}\\

24&33&\href{https://komorebi.invalid/b07/obligation_24}{7.2}&Finite Schrödinger representation&\nolinkurl{PaperXI.HeisenbergRepresentation}\\

25&34&\href{https://komorebi.invalid/b07/obligation_25}{7.3}&Character support&\nolinkurl{PaperXI.HeisenbergRepresentation}\\

26&36&\href{https://komorebi.invalid/b07/obligation_26}{7.5}&Isotropic averaging projector&\nolinkurl{PaperXI.HeisenbergRepresentation}\\

27&37&\href{https://komorebi.invalid/b07/obligation_27}{7.6}&Primitive line selector&\nolinkurl{PaperXI.HeisenbergRepresentation}\\

28&38&\href{https://komorebi.invalid/b07/obligation_28}{8.1}&Scalar choice invariance&\nolinkurl{PaperXI.TwistedMatrix}\\

29&39&\href{https://komorebi.invalid/b07/obligation_29}{8.2}&Alternating commutator form&\nolinkurl{PaperXI.TwistedMatrix}\\

30&40&\href{https://komorebi.invalid/b07/obligation_30}{8.3}&Nondegenerate twisted group algebra&\nolinkurl{PaperXI.TwistedMatrix}\\

31&41&\href{https://komorebi.invalid/b07/obligation_31}{8.4}&Radical free envelope&\nolinkurl{PaperXI.TwistedMatrix}\\

32&44&\href{https://komorebi.invalid/b07/obligation_32}{9.2}&Clifford multiplicity action&\nolinkurl{PaperXI.CliffordSource}\\

33&46&\href{https://komorebi.invalid/b07/obligation_33}{9.4}&Intrinsic commutator pairing&\nolinkurl{PaperXI.CliffordSource}\\

34&48&\href{https://komorebi.invalid/b07/obligation_34}{9.6}&Nondegenerate one break packet&\nolinkurl{PaperXI.CliffordSource}\\

35&49&\href{https://komorebi.invalid/b07/obligation_35}{9.7}&Isomorphism invariance&\nolinkurl{PaperXI.CliffordSource}\\

36&50&\href{https://komorebi.invalid/b07/obligation_36}{9.8}&Tame stabilizer action&\nolinkurl{PaperXI.CliffordSource}\\

37&51&\href{https://komorebi.invalid/b07/obligation_37}{9.9}&Heisenberg representation attached to one break&\nolinkurl{PaperXI.CliffordSource}\\

38&55&\href{https://komorebi.invalid/b07/obligation_38}{10.3}&Type side projective commutator&\nolinkurl{PaperXI.TypeNormalizerAction}\\

39&57&\href{https://komorebi.invalid/b07/obligation_39}{11.2}&Restricted Heisenberg bridge&\nolinkurl{PaperXI.CyclicBridge}\\

40&58&\href{https://komorebi.invalid/b07/obligation_40}{11.3}&Cyclic discrepancy character&\nolinkurl{PaperXI.CyclicBridge}\\

41&59&\href{https://komorebi.invalid/b07/obligation_41}{11.4}&Linearization over a cyclic stabilizer&\nolinkurl{PaperXI.CyclicBridge}\\

42&60&\href{https://komorebi.invalid/b07/obligation_42}{11.5}&One phase closes the cyclic bridge&\nolinkurl{PaperXI.CyclicBridge}\\

43&61&\href{https://komorebi.invalid/b07/obligation_43}{11.6}&Bridge uniqueness&\nolinkurl{PaperXI.CyclicBridge}\\

44&65&\href{https://komorebi.invalid/b07/obligation_44}{12.3}&Symplectic invariance&\nolinkurl{PaperXI.OddPhaseCertificate}\\

45&66&\href{https://komorebi.invalid/b07/obligation_45}{12.4}&Orthogonal multiplicativity&\nolinkurl{PaperXI.OddPhaseCertificate}\\

46&67&\href{https://komorebi.invalid/b07/obligation_46}{12.5}&Characteristic polynomial certificate&\nolinkurl{PaperXI.OddPhaseCertificate}\\

47&68&\href{https://komorebi.invalid/b07/obligation_47}{12.6}&Fixed and moving decomposition&\nolinkurl{PaperXI.OddPhaseCertificate}\\

48&70&\href{https://komorebi.invalid/b07/obligation_48}{12.8}&Hyperbolic phase formula&\nolinkurl{PaperXI.OddPhaseCertificate}\\

49&71&\href{https://komorebi.invalid/b07/obligation_49}{12.9}&Phase only bridge test&\nolinkurl{PaperXI.OddPhaseCertificate}\\

50&74&\href{https://komorebi.invalid/b07/obligation_50}{13.2}&Affine space of refinements&\nolinkurl{PaperXI.DyadicAffine}\\

51&77&\href{https://komorebi.invalid/b07/obligation_51}{13.5}&Dyadic Gauss sign&\nolinkurl{PaperXI.DyadicAffine}\\

52&78&\href{https://komorebi.invalid/b07/obligation_52}{13.6}&Orthogonal dyadic multiplicativity&\nolinkurl{PaperXI.DyadicAffine}\\

53&79&\href{https://komorebi.invalid/b07/obligation_53}{13.7}&The form alone does not determine the phase&\nolinkurl{PaperXI.DyadicAffine}\\

54&80&\href{https://komorebi.invalid/b07/obligation_54}{13.8}&Translation of quadratic refinements&\nolinkurl{PaperXI.DyadicAffine}\\

55&81&\href{https://komorebi.invalid/b07/obligation_55}{13.9}&Number of dyadic refinements&\nolinkurl{PaperXI.DyadicAffine}\\

56&83&\href{https://komorebi.invalid/b07/obligation_56}{14.2}&Matrix holonomy criterion&\nolinkurl{PaperXI.GraphGauge}\\

57&84&\href{https://komorebi.invalid/b07/obligation_57}{14.3}&Unbased gauge classification&\nolinkurl{PaperXI.GraphGauge}\\

58&85&\href{https://komorebi.invalid/b07/obligation_58}{14.4}&Scalar and projective gluing&\nolinkurl{PaperXI.GraphGauge}\\

59&86&\href{https://komorebi.invalid/b07/obligation_59}{14.5}&Projective cocycle form&\nolinkurl{PaperXI.GraphGauge}\\

60&87&\href{https://komorebi.invalid/b07/obligation_60}{14.6}&Exact edge residual from cycle defect&\nolinkurl{PaperXI.GraphGauge}\\

61&88&\href{https://komorebi.invalid/b07/obligation_61}{14.7}&Quantitative finite gluing&\nolinkurl{PaperXI.GraphGauge}\\

62&91&\href{https://komorebi.invalid/b07/obligation_62}{15.2}&Heisenberg tensor factorization&\nolinkurl{PaperXI.FiniteBreakAssembly}\\

63&92&\href{https://komorebi.invalid/b07/obligation_63}{15.3}&Breakwise phase product&\nolinkurl{PaperXI.FiniteBreakAssembly}\\

64&93&\href{https://komorebi.invalid/b07/obligation_64}{15.4}&Breakwise restricted bridge&\nolinkurl{PaperXI.FiniteBreakAssembly}\\

65&94&\href{https://komorebi.invalid/b07/obligation_65}{16.1}&Finite tame comparison package&\nolinkurl{PaperXI.FiniteTameComparison}\\

66&95&\href{https://komorebi.invalid/b07/obligation_66}{16.2}&Conditional atomic comparison&\nolinkurl{PaperXI.ConditionalAtomicComparison}\\

\bottomrule\end{longtable}

\chapter{Excerpt and claim crosswalk}

The following index maps teaching excerpts to the preserved formal source. Original proof obligation 66 corresponds to printed Theorem 16.2 and all-records row 95 in the historical 102-record register. Historical register row 66 is a different record, Proposition 12.4. The historical row 95 status ``proved internally'' is preserved as historical metadata and is superseded by this edition's conjectural-realization and conditional-comparison split. The \texttt{ConditionalAtomicComparison.source} interface is a proved conditional implication, while its general arithmetic-realization premise is the conjectural publication endpoint. Proposition 16.1's finite package appears in \texttt{FiniteTameComparison}; an arithmetic local-factor interpretation is a separate attachment.

\begin{longtable}{@{}p{0.06\textwidth}p{0.66\textwidth}p{0.12\textwidth}p{0.08\textwidth}@{}}\toprule ID&Module&Source lines&Page\\\midrule\endhead

E01&\nolinkurl{PaperXI.FourierProjectors}&13--15&\pageref{guide:E01}\\

E02&\nolinkurl{PaperXI.FourierProjectors}&68--77&\pageref{guide:E02}\\

E03&\nolinkurl{PaperXI.PacketTransport}&62--66&\pageref{guide:E03}\\

E04&\nolinkurl{PaperXI.TameProjectors}&45--52&\pageref{guide:E04}\\

E05&\nolinkurl{PaperXI.InducedCycleInvariants}&33--33&\pageref{guide:E05}\\

E06&\nolinkurl{PaperXI.InducedCycleInvariants}&75--77&\pageref{guide:E06}\\

E07&\nolinkurl{PaperXI.TameTensorSource}&66--69&\pageref{guide:E07}\\

E08&\nolinkurl{PaperXI.TameConductor}&56--58&\pageref{guide:E08}\\

E09&\nolinkurl{PaperXI.FiniteTameComparison}&93--107&\pageref{guide:E09}\\

E10&\nolinkurl{PaperXI.HeisenbergRepresentation}&14--16&\pageref{guide:E10}\\

E11&\nolinkurl{PaperXI.HeisenbergRepresentation}&58--59&\pageref{guide:E11}\\

E12&\nolinkurl{PaperXI.CentralProduct}&39--41&\pageref{guide:E12}\\

E13&\nolinkurl{PaperXI.FiniteBreakAssembly}&22--37&\pageref{guide:E13}\\

E14&\nolinkurl{PaperXI.FiniteBreakAssembly}&111--115&\pageref{guide:E14}\\

E15&\nolinkurl{PaperXI.OddPhaseCertificate}&46--50&\pageref{guide:E15}\\

E16&\nolinkurl{PaperXI.DyadicAffine}&74--75&\pageref{guide:E16}\\

E17&\nolinkurl{PaperXI.DyadicGauss}&51--58&\pageref{guide:E17}\\

E18&\nolinkurl{PaperXI.DyadicRefinementCounts}&42--44&\pageref{guide:E18}\\

E19&\nolinkurl{PaperXI.GraphGauge}&155--159&\pageref{guide:E19}\\

E20&\nolinkurl{PaperXI.GluedPacketSections}&22--26&\pageref{guide:E20}\\

E21&\nolinkurl{PaperXI.GluedPacketSections}&138--144&\pageref{guide:E21}\\

E22&\nolinkurl{PaperXI.MonodromyAtom}&64--66&\pageref{guide:E22}\\

E23&\nolinkurl{PaperXI.AtomicReconstruction}&17--21&\pageref{guide:E23}\\

E24&\nolinkurl{PaperXI.AtomicReconstruction}&143--144&\pageref{guide:E24}\\

E25&\nolinkurl{PaperXI.ValuationRamification}&15--28&\pageref{guide:E25}\\

E26&\nolinkurl{PaperXI.RamificationPolyphase}&311--313&\pageref{guide:E26}\\

E27&\nolinkurl{PaperXI.RamifiedUnramifiedQuotient}&42--43&\pageref{guide:E27}\\

E28&\nolinkurl{PaperXI.NormalizedFrobeniusQuotient}&79--79&\pageref{guide:E28}\\

E29&\nolinkurl{PaperXI.ActualProfileCandidate110}&34--88&\pageref{guide:E29}\\

E30&\nolinkurl{PaperXI.RankOneSourceCandidate110}&29--60&\pageref{guide:E30}\\

E31&\nolinkurl{PaperXI.MoritaActualCartan110}&137--152&\pageref{guide:E31}\\

E32&\nolinkurl{PaperXI.OnePlusIntersection110}&179--189&\pageref{guide:E32}\\

E33&\nolinkurl{PaperXI.ActualSourceCandidate111}&40--84&\pageref{guide:E33}\\

E34&\nolinkurl{PaperXI.ScalarCocycle112}&171--174&\pageref{guide:E34}\\

E35&\nolinkurl{PaperXI.EquivariantRepresentation113}&67--91&\pageref{guide:E35}\\

E36&\nolinkurl{PaperXI.CyclicNormalization113}&167--173&\pageref{guide:E36}\\

E37&\nolinkurl{PaperXI.LocalFactorIdentification}&14--18&\pageref{guide:E37}\\

E38&\nolinkurl{PaperXI.LocalFactorIdentification}&25--28&\pageref{guide:E38}\\

E39&\nolinkurl{PaperXI.LocalFactorIdentification}&49--55&\pageref{guide:E39}\\

E40&\nolinkurl{PaperXI.ConditionalAtomicComparison}&127--143&\pageref{guide:E40}\\

\bottomrule\end{longtable}

\end{document}
